% =====================================================================
%  GrafosCompleto.tex
%  ========================
%
%
%  Descripción:
%  ------------
%  Versión extendida de la parte de grafos de la libreta, con los
%  ejercicios resueltos de cada sección. Es el respaldo completo del
%  texto antes de resumirlo; la libreta lleva la versión corta.
%
%
%  Metadata:
%  ----------
%   * Author : zxxz6 (Bryan Violante Arriaga)
%   * Version: 1.0.0
%
%
%  History:
%  ------------
%  Author      Date            Description
%  zxxz6       08/11/2026      Creation
%
%
% =====================================================================
\documentclass[11pt,letterpaper]{article}

\newcommand{\materia}{Análisis y Diseño de Algoritmos}
\newcommand{\profesor}{Dr. José Martínez Carranza}
\newcommand{\periodo}{Otoño 2026}
\newcommand{\alumno}{Bryan Ezequiel Violante Arriaga}

% =====================================================================
%  Preambulo.tex
%  ========================
%
%
%  Descripción:
%  ------------
%  Preámbulo compartido de los apuntes de clase: paquetes, estilo de
%  página, entornos de teorema, las cuatro cajas de color, el estilo de
%  listings y el pseudocódigo en español. Se carga con \input desde
%  cada documento en lugar de copiarlo, para que un arreglo aquí llegue
%  a todas las libretas de una vez.
%
%
%  Considerations:
%  ------------
%  - Se carga con \input y ruta relativa, no con \usepackage, porque
%    \usepackage solo busca en la carpeta del documento y en el árbol
%    de TeX. Un .sty instalado en TEXMFHOME evitaría la ruta, pero vive
%    fuera del repo y una clona nueva no compilaría
%  - Los cuatro datos de la materia se declaran en el documento ANTES
%    del \input. hyperref expande \materia al cargarse, así que si no
%    existe todavía la compilación truena
%  - Cuando la materia la imparte más de un profesor, los nombres van
%    en un solo \profesor separados por \\. La portadilla los imprime
%    dentro de un center, que es donde \\ sí corta renglón. Repetir
%    \newcommand{\profesor} no es opción: truena con "already defined"
%  - algorithmicx no tiene opción de idioma y babel no lo traduce, así
%    que las palabras clave del pseudocódigo se renombran a mano
%  - Los nombres de los entornos van sin acento (\begin{definicion})
%    aunque el rótulo impreso sí lo lleve
%  - babel se carga con es-noshorthands porque en español el " queda
%    activo y se come el texto que le sigue, además de romper listings
%
%
%  Metadata:
%  ----------
%   * Author : zxxz6 (Bryan Violante Arriaga)
%   * Version: 1.7.0
%
%
%  History:
%  ------------
%  Author      Date            Description
%  zxxz6       24/09/2026      Encabezado de una hoja de un ejercicio
%  zxxz6       23/09/2026      Macros de los documentos de soluciones
%  zxxz6       30/08/2026      Macros de lim, limsup y liminf sobre n
%  zxxz6       25/08/2026      Macros de Omega, Theta, o y omega chicas
%  zxxz6       20/08/2026      Cargue pgfplots para graficas de datos
%  zxxz6       20/08/2026      Cargue tikz para diagramas de conjuntos
%  zxxz6       19/08/2026      Documente \profesor con varios nombres
%  zxxz6       19/08/2026      Agregue la caja Idea Random
%  zxxz6       18/08/2026      Creation
%
%
% =====================================================================
%  USO
%  ---
%  \documentclass[11pt,letterpaper]{article}
%  \newcommand{\materia}{...}      % los cuatro datos van arriba
%  \newcommand{\profesor}{...}
%  \newcommand{\periodo}{...}
%  \newcommand{\alumno}{...}
%  \input{../../templates/Preambulo.tex}
%  \begin{document}
%  \portadilla
%
%  Con dos o mas profesores se declara UN solo \profesor y los nombres
%  se separan con \\, uno por renglon:
%
%  \newcommand{\profesor}{Dra. Fulana de Tal \\
%                         Dr. Mengano de Cual}
% =====================================================================


% ==================== IDIOMA Y CODIFICACIÓN ====================
\usepackage[utf8]{inputenc}
\usepackage[T1]{fontenc}
\usepackage[spanish,es-tabla,es-noshorthands]{babel}
\usepackage{lmodern}
\usepackage{microtype}

% ==================== PÁGINA ====================
\usepackage[letterpaper,margin=2.3cm,headheight=14pt]{geometry}
\usepackage{parskip}
\usepackage{fancyhdr}
\usepackage{enumitem}

% ==================== MATEMÁTICAS ====================
\usepackage{amsmath,amssymb,amsthm,mathtools}

% ==================== FIGURAS, TABLAS, CÓDIGO ====================
\usepackage{graphicx}
\usepackage{booktabs}
\usepackage{array}
% float da el especificador [H], que clava una tabla donde se
% escribio. Una traza que LaTeX manda dos paginas adelante deja
% de leerse junto al paso que la produjo
\usepackage{float}
\usepackage{xcolor}
% tikz ya viene arrastrado por tcolorbox[most], pero se carga
% explicito para no depender de ese detalle; la libreria babel
% protege a tikz de los caracteres activos del español
\usepackage{tikz}
\usetikzlibrary{babel}
% pgfplots dibuja graficas de datos (ejes, escalas log) sobre tikz
\usepackage{pgfplots}
\pgfplotsset{compat=1.18}
\usepackage{listings}
\usepackage[ruled,section]{algorithm}
\usepackage{algpseudocode}
\usepackage[most]{tcolorbox}
% soul da \hl, el unico resaltado que parte renglon; \colorbox no
\usepackage{soul}
\usepackage{hyperref}

% =====================================================================
%  ESTILO
% =====================================================================

% --- encabezado y pie de cada página ---
\pagestyle{fancy}
\fancyhf{}
\fancyhead[L]{\small\textsc{\materia}}
\fancyhead[R]{\small\periodo}
\fancyfoot[C]{\small\thepage}
\renewcommand{\headrulewidth}{0.4pt}

% --- enlaces ---
\hypersetup{
  colorlinks=true,
  linkcolor=black,
  urlcolor=blue,
  citecolor=black,
  pdftitle={Apuntes de \materia},
  pdfauthor={\alumno}
}

% --- paleta ---
\definecolor{acento}{HTML}{1F4E79}
\definecolor{fondoIdea}{HTML}{EAF2FB}
\definecolor{fondoOjo}{HTML}{FDF0E6}
\definecolor{fondoPend}{HTML}{F2F0F7}
\definecolor{comentario}{HTML}{5A7D5A}
\definecolor{cadena}{HTML}{A03030}
\definecolor{fondoCodigo}{HTML}{F7F7F7}
\definecolor{fondoResultado}{HTML}{EAF5EC}
\definecolor{marcador}{HTML}{FFF3A3}

% --- entornos de teorema; numeran por sección, o sea por clase ---
\theoremstyle{definition}
\newtheorem{definicion}{Definición}[section]
\newtheorem{ejemplo}[definicion]{Ejemplo}
\newtheorem{ejercicio}[definicion]{Ejercicio}

\theoremstyle{plain}
\newtheorem{teorema}[definicion]{Teorema}
\newtheorem{lema}[definicion]{Lema}
\newtheorem{corolario}[definicion]{Corolario}
\newtheorem{proposicion}[definicion]{Proposición}

\theoremstyle{remark}
\newtheorem*{observacion}{Observación}
\newtheorem*{quick_sum}{Resumen rápido}

% --- cajas de color para lo que se relee antes del examen ---
\newtcolorbox{idea}[1][]{
  colback=fondoIdea, colframe=acento, boxrule=0.6pt,
  arc=2pt, left=6pt, right=6pt, top=5pt, bottom=5pt,
  title={Idea clave}, fonttitle=\bfseries\small, #1
}
\newtcolorbox{ojo}[1][]{
  colback=fondoOjo, colframe=orange!70!black, boxrule=0.6pt,
  arc=2pt, left=6pt, right=6pt, top=5pt, bottom=5pt,
  title={Ojito...}, fonttitle=\bfseries\small, #1
}
\newtcolorbox{pendiente}[1][]{
  colback=fondoPend, colframe=violet!60!black, boxrule=0.6pt,
  arc=2pt, left=6pt, right=6pt, top=5pt, bottom=5pt,
  title={Pendiente por resolver}, fonttitle=\bfseries\small, #1
}
\newtcolorbox{idear}[1][]{
  colback=fondoPend, colframe=pink!60!black, boxrule=0.6pt,
  arc=2pt, left=6pt, right=6pt, top=5pt, bottom=5pt,
  title={Idea Random}, fonttitle=\bfseries\small, #1
}

% --- el resultado al que llego un ejercicio, para que se vea de
%     lejos al hojear la hoja antes del examen ---
\newtcolorbox{resultado}[1][]{
  colback=fondoResultado, colframe=green!45!black, boxrule=0.6pt,
  arc=2pt, left=6pt, right=6pt, top=5pt, bottom=5pt,
  title={Resultado}, fonttitle=\bfseries\small, #1
}

% --- código ---
\lstdefinestyle{codigo}{
  backgroundcolor=\color{fondoCodigo},
  basicstyle=\ttfamily\footnotesize,
  keywordstyle=\color{acento}\bfseries,
  commentstyle=\color{comentario}\itshape,
  stringstyle=\color{cadena},
  numbers=left, numberstyle=\tiny\color{gray}, numbersep=8pt,
  frame=single, rulecolor=\color{gray!40},
  breaklines=true, showstringspaces=false,
  tabsize=4, captionpos=b,
  extendedchars=true, inputencoding=utf8,
  literate={á}{{\'a}}1 {é}{{\'e}}1 {í}{{\'i}}1
           {ó}{{\'o}}1 {ú}{{\'u}}1 {ñ}{{\~n}}1
           {Ü}{{\"U}}1 {ü}{{\"u}}1 {¿}{{?`}}1
           {¡}{{!`}}1
}
\lstset{style=codigo}
\renewcommand{\lstlistingname}{Código}
\renewcommand{\lstlistlistingname}{Índice de código}
% --- pseudocódigo; algorithmicx arma el cuerpo y algorithm el
%     flotante que lo numera y lo pone en el índice. Las palabras
%     clave vienen en inglés de fábrica y babel no las toca, así que
%     se renombran una por una ---
\floatname{algorithm}{Algoritmo}
\renewcommand{\listalgorithmname}{Índice de algoritmos}

\algrenewcommand\algorithmicrequire{\textbf{Entrada:}}
\algrenewcommand\algorithmicensure{\textbf{Salida:}}
\algrenewcommand\algorithmicend{\textbf{fin}}
\algrenewcommand\algorithmicif{\textbf{si}}
\algrenewcommand\algorithmicthen{\textbf{entonces}}
\algrenewcommand\algorithmicelse{\textbf{si no}}
\algrenewcommand\algorithmicfor{\textbf{para}}
\algrenewcommand\algorithmicforall{\textbf{para todo}}
\algrenewcommand\algorithmicdo{\textbf{hacer}}
\algrenewcommand\algorithmicwhile{\textbf{mientras}}
\algrenewcommand\algorithmicrepeat{\textbf{repetir}}
\algrenewcommand\algorithmicuntil{\textbf{hasta que}}
\algrenewcommand\algorithmicloop{\textbf{ciclo}}
\algrenewcommand\algorithmicprocedure{\textbf{procedimiento}}
\algrenewcommand\algorithmicfunction{\textbf{función}}
\algrenewcommand\algorithmicreturn{\textbf{devolver}}
\algrenewcomment[1]{\hfill{\color{comentario}\itshape // #1}}

% Las dos que algorithmicx no trae. El \ del final es obligatorio: sin
% el TeX se come el espacio y sale "hastan-1" en vez de "hasta n-1"
\algnewcommand\Hasta{\textbf{hasta}\ }
\algnewcommand\Intercambiar{\textbf{intercambiar}\ }

% --- abreviaturas que se usan a cada rato ---
\newcommand{\N}{\mathbb{N}}
\newcommand{\Z}{\mathbb{Z}}
\newcommand{\Q}{\mathbb{Q}}
\newcommand{\R}{\mathbb{R}}
\newcommand{\C}{\mathbb{C}}
\newcommand{\abs}[1]{\left\lvert #1 \right\rvert}
\newcommand{\norm}[1]{\left\lVert #1 \right\rVert}
\newcommand{\set}[1]{\left\{ #1 \right\}}
\newcommand{\Ord}[1]{\mathcal{O}\!\left(#1\right)}
% Mayuscula la notacion grande, minuscula la chica
\newcommand{\Omg}[1]{\Omega\!\left(#1\right)}
\newcommand{\Tht}[1]{\Theta\!\left(#1\right)}
\newcommand{\ord}[1]{o\!\left(#1\right)}
\newcommand{\omg}[1]{\omega\!\left(#1\right)}
% Los tres limites sobre n, que salen en cada renglon del criterio
% del cociente para la notacion asintotica
\newcommand{\LM}{\lim_{n \to \infty}}
\newcommand{\LS}{\limsup_{n \to \infty}}
\newcommand{\LI}{\liminf_{n \to \infty}}

% --- encabezado de cada sesión: título a la izquierda, fecha a la
%     derecha. El argumento corto es el que sale en el índice ---
\newcommand{\clase}[2]{%
  \section[#1]{#1 \hfill \normalfont\small\textit{#2}}%
}

% --- portadilla y tabla de contenidos. Se llama una vez, justo
%     despues de \begin{document}. El renglon del profesor va dentro
%     del center, asi que \profesor puede traer varios nombres
%     separados por \\ y salen apilados y centrados sin tocar nada ---
\newcommand{\portadilla}{%
  \begin{center}
    {\LARGE\bfseries\materia\par}
    \vspace{0.4em}
    {\large \profesor \par}
    \vspace{0.2em}
    {\small \periodo\quad-\quad\alumno\par}
    \vspace{0.3em}
    \rule{\linewidth}{0.5pt}
  \end{center}
  \vspace{0.5em}
  \tableofcontents
  \vspace{1em}
  \rule{\linewidth}{0.5pt}%
}

% =====================================================================
%  DOCUMENTOS DE SOLUCIONES
%  Lo que usa una tanda de ejercicios resueltos. Una libreta no toca
%  nada de aqui, pero vive en el preambulo compartido para no tener
%  dos archivos que arreglar cuando cambie el estilo.
% =====================================================================

% --- encabezado compacto: titulo, subtitulo y regla. Sustituye a
%     \portadilla cuando el documento es una tanda de ejercicios: sin
%     indice, porque se hojea de corrido, no se navega ---
\newcommand{\encabezadoSoluciones}[2]{%
  % El paquete algorithm se carga con la opcion [section] para que en
  % una libreta el pseudocodigo se numere por clase. Aqui no hay
  % secciones, asi que esa numeracion sale como "Algoritmo 0.1". Se
  % corrige al vuelo, y solo en los documentos que llaman a este
  % encabezado: la libreta no se entera.
  \renewcommand{\thealgorithm}{\arabic{algorithm}}%
  \begin{center}
    {\LARGE\bfseries #1\par}
    \vspace{0.3em}
    {\small\itshape #2\par}
    \vspace{0.5em}
    \rule{\linewidth}{0.5pt}
  \end{center}
  \vspace{0.3em}%
}

% --- hoja de un solo ejercicio: un renglon chico arriba con el
%     titulo y su numero, y ya. Sin titulo grande y sin pie de pagina,
%     porque en una hoja de dos paginas cada centimetro que se va en
%     adornos es un paso del desarrollo que no cabe ---
\newcommand{\encabezadoEjercicio}[3]{%
  % Misma correccion que en \encabezadoSoluciones: sin secciones, el
  % pseudocodigo saldria numerado "0.1"
  \renewcommand{\thealgorithm}{\arabic{algorithm}}%
  \fancyhf{}%
  \fancyhead[L]{\small\textsc{#1} (ejercicio #2 de #3)}%
  \renewcommand{\headrulewidth}{0.4pt}%
  \pagestyle{fancy}%
  \thispagestyle{fancy}%
}

% --- bloque tematico, numerado con romanos: I, II, III. Agrupa los
%     ejercicios que comparten tema. No dibuja regla: la del
%     \problema que sigue ya separa el titulo del primer ejercicio,
%     y dos rayas seguidas dejan un hueco feo ---
\newcounter{parte}
\renewcommand{\theparte}{\Roman{parte}}
\newcommand{\parte}[1]{%
  \bigskip
  \refstepcounter{parte}%
  \noindent{\large\bfseries \theparte. #1\par}%
}

% --- un ejercicio resuelto. Se numera solo y corrido a lo largo del
%     documento, sin reiniciar en cada parte, para poder decir "el 14"
%     sin aclarar de que bloque ---
\newcounter{problema}
\newcommand{\problema}[1]{%
  \bigskip
  \noindent\rule{\linewidth}{0.4pt}\par
  \vspace{0.3em}
  \refstepcounter{problema}%
  \noindent{\bfseries Ejercicio \theproblema\ --- #1\par}
  \vspace{0.3em}%
}

% --- rotulo del paso que sigue: Caso base, Hipotesis inductiva, Paso
%     inductivo. El texto sigue en el mismo renglon ---
\newcommand{\paso}[1]{\par\vspace{0.3em}\noindent\textbf{#1}\ }

% --- la respuesta de un ejercicio, centrada y en un cuadro. El
%     argumento va en modo matematico, sin los delimitadores ---
\newcommand{\respuesta}[1]{%
  \[
    \boxed{\;#1\;}
  \]
}

% --- resaltado inline, como pasarle el plumon a media frase. No
%     admite matematicas adentro: soul rompe con $...$, asi que una
%     formula que hay que destacar va en \respuesta ---
\sethlcolor{marcador}
\newcommand{\resaltar}[1]{\hl{#1}}

% --- justificacion al final de un renglon de \align, en chiquito y
%     entre parentesis, para decir de donde salio el paso ---
\newcommand{\razon}[1]{\quad\text{\small\itshape (#1)}}

%############################ END OF PREAMBULO.TEX #############################
%###############################################################################


\begin{document}

\portadilla

\newpage

% =====================================================================
%  Clase 22 Oct - Grafos
% =====================================================================

\clase{Grafos}{22 de octubre de 2026}
\label{sec:grafos}

% Estilos de los grafos de toda la parte de grafos de la libreta. Este
% bloque reemplaza al \tikzset de la clase del 22 de octubre y se
% define una sola vez, al principio de esa clase. Ninguna seccion
% agrega estilos globales: si una figura necesita algo propio, lo
% declara en las opciones de su tikzpicture o de un scope.
\tikzset{
  vert/.style={circle, draw=acento, fill=fondoIdea, font=\small,
               minimum size=0.62cm, inner sep=0pt},
  vertOk/.style={vert, draw=green!45!black, fill=fondoResultado},
  vertGris/.style={vert, draw=black!60, fill=gray!30},
  vertNegro/.style={vert, draw=black, fill=black!75, text=white},
  vertMarca/.style={vert, draw=orange!80!black, fill=orange!15,
                    line width=1pt},
  vertAlerta/.style={vert, draw=red!70!black, fill=red!10,
                     line width=1pt},
  arista/.style={draw=gray!70, thick},
  aristaArbol/.style={draw=acento, line width=1.6pt},
  aristaFuera/.style={draw=orange!80!black, thick, dashed},
  aristaResalte/.style={draw=orange!85!black, line width=1.8pt},
  aristaAlerta/.style={draw=red!70!black, line width=1.6pt},
  dirigida/.style={->, >=stealth, thick, draw=gray!70},
  dirigidaArbol/.style={->, >=stealth, draw=acento, line width=1.6pt},
  dirigidaFuera/.style={->, >=stealth, thick, dashed,
                        draw=orange!80!black},
  dirigidaResalte/.style={->, >=stealth, draw=orange!85!black,
                          line width=1.8pt},
  peso/.style={font=\scriptsize, fill=white, inner sep=1.2pt},
  nota/.style={font=\scriptsize, text=gray},
  zona/.style={draw=acento!40, fill=acento!8, rounded corners=8pt},
  corte/.style={draw=red!60!black, thick, dashed},
}

\begin{quick_sum}
  Un grafo es un par de conjuntos, vértices y aristas, y modela
  cualquier cosa hecha de objetos y relaciones entre pares de ellos:
  ciudades y carreteras, páginas y ligas, estados de un juego y
  jugadas. Se representa con matriz o con listas de adyacencia, y se
  clasifica por sus aristas: dirigido o no, pesado o no, cíclico o
  acíclico, simple o no, encajado o topológico, implícito o
  explícito, etiquetado o no. Se vieron los dos recorridos, búsqueda
  en anchura (con cola) y en profundidad (con recursión o pila), y
  tres problemas: camino más corto con Dijkstra y Floyd, detección de
  ciclos con BFS y árbol de expansión mínima con Kruskal. Todos los
  ejemplos usan los mismos tres grafos.
\end{quick_sum}

Esta clase es la base de la parte de grafos: presenta cada tema con
un ejemplo y la idea de por qué funciona. Las pruebas, las variantes y
los temas que no cupieron están en las secciones de complemento que
siguen, cada una sobre un bloque:
\begin{itemize}[itemsep=0.1em]
  \item Sección \ref{sec:teoria-grafos}: terminología completa
        (caminatas, senderos, caminos, circuitos y ciclos), conexidad
        y recorridos de Euler.
  \item Sección \ref{sec:hamilton-arboles}: ciclos de Hamilton,
        grafos bipartitos, cliques y árboles.
  \item Sección \ref{sec:busquedas}: las representaciones y BFS y DFS
        a fondo: colores, tiempos, pruebas de corrección y DFS
        iterativo.
  \item Sección \ref{sec:aplicaciones-dfs}: orden topológico, ciclos
        dirigidos y componentes fuertemente conexas con DFS.
  \item Sección \ref{sec:mst}: árboles de expansión mínima: la
        propiedad del corte, Kruskal y Prim.
  \item Sección \ref{sec:caminos-fuente}: caminos más cortos desde una
        fuente: relajación, Bellman-Ford, caminos en DAG y Dijkstra
        con su prueba.
  \item Sección \ref{sec:caminos-pruebas}: las demostraciones de las
        propiedades de la relajación y problemas de caminos más
        cortos.
  \item Sección \ref{sec:todos-pares}: caminos más cortos entre todos
        los pares: Floyd-Warshall con predecesores, cerradura
        transitiva y Johnson.
  \item Sección \ref{sec:flujo}: flujo máximo con Ford-Fulkerson y
        Edmonds-Karp.
  \item Sección \ref{sec:emparejamiento}: emparejamiento bipartito
        máximo.
  \item Sección \ref{sec:push-relabel}: empujar y reetiquetar, la otra
        familia de algoritmos de flujo máximo.
\end{itemize}

\subsection*{Definición formal}

\begin{definicion}[Grafo]
  Un \textbf{grafo} es un par $G = (V, E)$ donde $V$ es un conjunto
  finito de \textbf{vértices} y $E$ un conjunto de \textbf{aristas}.
  En un grafo no dirigido cada arista es un par no ordenado de
  vértices distintos,
  \[
    E \subseteq \set{\set{u, v} : u, v \in V,\ u \neq v} .
  \]
  Se escribe $n = \abs{V}$ y $m = \abs{E}$.
\end{definicion}

Si $\set{u, v} \in E$, los vértices $u$ y $v$ son
\textbf{adyacentes} y la arista es \textbf{incidente} en los dos. El
\textbf{grado} de un vértice, $\text{gr}(v)$, es el número de aristas
incidentes en él. Otras fuentes, como Cormen, escriben $\deg(v)$.

\begin{ejemplo}[El grafo $G$ de esta clase]
  \[
    V = \set{A, B, C, D, E, F, G} , \qquad
    E = \set{\set{A,B}, \set{A,C}, \set{B,D}, \set{B,E},
             \set{C,F}, \set{E,F}, \set{D,G}} .
  \]
  \begin{center}
  \begin{tikzpicture}[x=1cm, y=1cm]
    \node[vert] (A) at (0, 0) {A};
    \node[vert] (B) at (-1.6, -1.2) {B};
    \node[vert] (C) at (1.6, -1.2) {C};
    \node[vert] (D) at (-2.6, -2.4) {D};
    \node[vert] (E) at (-0.6, -2.4) {E};
    \node[vert] (F) at (1.6, -2.4) {F};
    \node[vert] (G) at (-2.6, -3.6) {G};
    \foreach \a/\b in {A/B, A/C, B/D, B/E, C/F, E/F, D/G}
      \draw[arista] (\a) -- (\b);
  \end{tikzpicture}
  \end{center}
  $n = 7$ y $m = 7$. Los grados son $\text{gr}(B) = 3$,
  $\text{gr}(G) = 1$ y $2$ para los demás. Este grafo se usa en toda
  la clase.
\end{ejemplo}

\begin{observacion}
  Cada arista suma $1$ al grado de sus dos extremos, así que
  $\sum_{v \in V} \text{gr}(v) = 2m$. En $G$:
  $2 + 3 + 2 + 2 + 2 + 2 + 1 = 14 = 2 \cdot 7$. Y como hay a lo más
  un par por cada dos vértices, $0 \leq m \leq n(n-1)/2$.
\end{observacion}

% =====================================================================

\subsection*{Matriz y listas de adyacencia}

\begin{definicion}[Matriz de adyacencia]
  Con los vértices numerados de $1$ a $n$, la matriz de adyacencia es
  la matriz $A$ de $n \times n$ con
  \[
    A[i][j] =
    \begin{cases}
      1 & \text{si } \set{i, j} \in E , \\
      0 & \text{si no.}
    \end{cases}
  \]
\end{definicion}

Para $G$, en el orden $A, \dots, G$:
\[
  \begin{array}{c|ccccccc}
      & A & B & C & D & E & F & G \\ \hline
    A & 0 & 1 & 1 & 0 & 0 & 0 & 0 \\
    B & 1 & 0 & 0 & 1 & 1 & 0 & 0 \\
    C & 1 & 0 & 0 & 0 & 0 & 1 & 0 \\
    D & 0 & 1 & 0 & 0 & 0 & 0 & 1 \\
    E & 0 & 1 & 0 & 0 & 0 & 1 & 0 \\
    F & 0 & 0 & 1 & 0 & 1 & 0 & 0 \\
    G & 0 & 0 & 0 & 1 & 0 & 0 & 0
  \end{array}
\]
En un grafo no dirigido la matriz es simétrica y la suma de cada
renglón es el grado del vértice. La diagonal es cero porque la
definición no admite lazos. En un grafo dirigido las dos cosas cambian,
como se ve más abajo.

La alternativa son las \textbf{listas de adyacencia}: para cada vértice,
la lista de sus vecinos.
\[
  \begin{array}{ll}
    A: B, C & \quad E: B, F \\
    B: A, D, E & \quad F: C, E \\
    C: A, F & \quad G: D \\
    D: B, G &
  \end{array}
\]

\begin{center}
\small
\begin{tabular}{lcc}
  \toprule
  \textbf{Operación} & \textbf{Matriz} & \textbf{Listas} \\
  \midrule
  Memoria & $\Tht{n^2}$ & $\Tht{n + m}$ \\
  ¿Es $\set{u, v}$ arista? & $\Ord{1}$ & $\Ord{\text{gr}(u)}$ \\
  Recorrer los vecinos de $u$ & $\Tht{n}$ & $\Tht{\text{gr}(u)}$ \\
  Recorrer todo el grafo & $\Tht{n^2}$ & $\Tht{n + m}$ \\
  \bottomrule
\end{tabular}
\end{center}

\begin{idea}
  Casi todos los grafos reales son \textbf{dispersos}: $m$ es mucho
  menor que $n^2$. Un mapa de carreteras con un millón de cruces tiene
  unos cuantos millones de tramos, no $10^{12}$. Ahí la matriz gasta
  memoria en ceros y las listas ganan. La matriz conviene en grafos
  \textbf{densos} o cuando se pregunta mucho si una arista existe.
  Las complejidades de esta clase suponen listas.
\end{idea}

% =====================================================================

\subsection*{Tipos de grafos}

La clasificación es la de Skiena. Cada par describe una propiedad de
las aristas o de cómo se guarda el grafo, y cambia qué algoritmo se
puede usar.

\subsubsection*{Dirigido o no dirigido}

En un \textbf{grafo dirigido} las aristas son pares ordenados,
$E \subseteq V \times V$: la arista $(u, v)$ va de $u$ a $v$ y no al
revés. Cada vértice tiene \textbf{grado de entrada}
$\text{gr}^{-}(v)$, las aristas que llegan a él, y \textbf{grado de
salida} $\text{gr}^{+}(v)$, las que salen, y la matriz de adyacencia
deja de ser simétrica. Una calle de un solo sentido o \textit{A sigue
a B} en una red social son dirigidas; una amistad no.

\begin{center}
\begin{tikzpicture}[baseline=(current bounding box.center)]
  \node[vert] (n1) at (0, 0) {1};
  \node[vert] (n2) at (2.6, 0) {2};
  \node[vert] (n3) at (2.6, -2.2) {3};
  \node[vert] (n4) at (0, -2.2) {4};
  \draw[dirigida] (n1) to[bend left=15] (n2);
  \draw[dirigida] (n2) to[bend left=15] (n1);
  \draw[dirigida] (n1) to[bend right=15] (n4);
  \draw[dirigida] (n4) to[bend right=15] (n1);
  \draw[dirigida] (n2) -- (n3);
  \draw[dirigida] (n3) -- (n1);
  \draw[dirigida] (n3) -- (n4);
\end{tikzpicture}
\hspace{1.2cm}
$\begin{array}{c|cccc}
    & 1 & 2 & 3 & 4 \\ \hline
  1 & 0 & 1 & 0 & 1 \\
  2 & 1 & 0 & 1 & 0 \\
  3 & 1 & 0 & 0 & 1 \\
  4 & 1 & 0 & 0 & 0
\end{array}$
\end{center}

$A[2][3] = 1$ pero $A[3][2] = 0$: de $2$ se llega a $3$ en un paso,
de $3$ a $2$ no. El renglón de un vértice suma su grado de salida y
su columna, el de entrada: el vértice $1$ tiene $\text{gr}^{+}(1) = 2$
y $\text{gr}^{-}(1) = 3$. Este grafo se usa en Floyd.

\subsubsection*{Pesado o no pesado}

En un \textbf{grafo pesado} cada arista tiene un número,
$w : E \to \R$: distancia, tiempo, costo. En la matriz se guarda el
peso en lugar del $1$, y $\infty$ donde no hay arista. El $0$ de la
diagonal de la matriz de abajo no es parte de la representación: es
la convención de Floyd, para quien ir de un vértice a sí mismo cuesta
$0$. Otros problemas llenan distinto las casillas vacías; en flujo,
por ejemplo, una arista ausente es una capacidad $0$.

\begin{center}
\begin{tikzpicture}[baseline=(current bounding box.center)]
  \node[vert] (A) at (0, 0) {A};
  \node[vert] (B) at (2, 1.4) {B};
  \node[vert] (C) at (2, -1.4) {C};
  \node[vert] (D) at (4.5, 1.4) {D};
  \node[vert] (E) at (4.5, -1.4) {E};
  \node[vert] (F) at (6.5, 0) {F};
  \foreach \a/\b/\w in {A/B/4, A/C/2, B/C/1, B/D/5, C/D/8, C/E/10,
                        D/E/2, D/F/6, E/F/3}
    \draw[arista] (\a) -- node[peso] {\w} (\b);
  \draw[arista] (A) to[out=75, in=150] node[peso] {7} (D);
\end{tikzpicture}
\hspace{0.8cm}
$\begin{array}{c|cccccc}
    & A & B & C & D & E & F \\ \hline
  A & 0 & 4 & 2 & 7 & \infty & \infty \\
  B & 4 & 0 & 1 & 5 & \infty & \infty \\
  C & 2 & 1 & 0 & 8 & 10 & \infty \\
  D & 7 & 5 & 8 & 0 & 2 & 6 \\
  E & \infty & \infty & 10 & 2 & 0 & 3 \\
  F & \infty & \infty & \infty & 6 & 3 & 0
\end{array}$
\end{center}

En un grafo no pesado el camino más corto es el de menos aristas, y
lo encuentra BFS. En uno pesado el de menos aristas puede no ser el
más barato: de A a D la arista directa pesa $7$, y A, C, B, D pesa
$2 + 1 + 5 = 8$. Hacen falta Dijkstra o Floyd. Este grafo se usa en
Dijkstra y en Kruskal.

\subsubsection*{Cíclico o acíclico}

Una \textbf{caminata} es una sucesión de vértices en la que cada uno
es adyacente al siguiente, y puede repetir lo que sea. Si no repite
aristas es un \textbf{sendero}; si no repite vértices, un
\textbf{camino}, que es el sentido de ``camino'' en toda esta clase.
Un sendero que regresa a su vértice de inicio es un
\textbf{circuito}, y un \textbf{ciclo} es un circuito que tampoco
repite vértices, salvo el primero, que es también el último. $G$
tiene uno: A, B, E, F, C, A. La diferencia importa: un ocho que pasa
dos veces por el mismo vértice es circuito pero no ciclo, y en un
grafo simple no dirigido todo ciclo tiene al menos tres aristas. En
un grafo dirigido el ciclo sigue el sentido de las flechas. La
sección \ref{sec:teoria-grafos} desarrolla esta terminología.

Un grafo sin ciclos es \textbf{acíclico}. Un no dirigido, conexo y
acíclico es un \textbf{árbol}; un dirigido acíclico es un
\textbf{DAG} (\textit{directed acyclic graph}), como las seriaciones
de un plan de estudios.

\begin{center}
\begin{tikzpicture}[
  materia/.style={draw=acento, rounded corners=2pt, fill=fondoIdea,
                  font=\scriptsize, inner sep=3pt}]
  \node[materia] (cal) at (0, 0) {Cálculo};
  \node[materia] (pro) at (2.6, 0) {Probabilidad};
  \node[materia] (alg) at (0, -1.1) {Álgebra lineal};
  \node[materia] (prg) at (0, -2.2) {Programación};
  \node[materia] (ada) at (2.6, -2.2) {ADA};
  \node[materia] (ml) at (5.2, -1.1) {Aprendizaje};
  \draw[dirigida] (cal) -- (pro);
  \draw[dirigida] (pro) -- (ml);
  \draw[dirigida] (alg) -- (ml);
  \draw[dirigida] (prg) -- (ada);
  \draw[dirigida] (ada) -- (ml);
  \node[nota] at (2.6, 0.7) {DAG: ninguna flecha regresa};
  \begin{scope}[xshift=8.6cm, yshift=0.2cm]
    \node[vert] (r) at (0, 0) {};
    \node[vert] (h1) at (-1, -1.1) {};
    \node[vert] (h2) at (1, -1.1) {};
    \node[vert] (h3) at (-1.6, -2.2) {};
    \node[vert] (h4) at (-0.4, -2.2) {};
    \node[vert] (h5) at (1, -2.2) {};
    \foreach \a/\b in {r/h1, r/h2, h1/h3, h1/h4, h2/h5}
      \draw[arista] (\a) -- (\b);
    \node[nota] at (0, 0.5) {árbol: $m = n - 1$};
  \end{scope}
\end{tikzpicture}
\end{center}

Un DAG siempre tiene un \textbf{orden topológico}: una forma de poner
los vértices en fila con todas las flechas hacia adelante, el orden en
que se pueden cursar las materias. Un grafo con un ciclo dirigido no
lo tiene.

\subsubsection*{Simple o no simple}

Un grafo es \textbf{simple} si no tiene \textbf{lazos} (aristas de un
vértice a sí mismo) ni \textbf{aristas múltiples} (dos aristas entre
el mismo par). La definición no dirigida de arriba, con conjuntos,
solo admite grafos simples: $\set{u, u}$ no es un par y un conjunto no
repite elementos. La dirigida, $E \subseteq V \times V$, sí admite
lazos $(u, u)$, aunque tampoco aristas múltiples; un lazo pone un $1$
en la diagonal de la matriz y suma $1$ a cada grado de su vértice, al
de entrada y al de salida. Para tener aristas múltiples $E$ se vuelve
un multiconjunto.

\begin{center}
\begin{tikzpicture}
  \node[vert] (a) at (0, 0) {};
  \node[vert] (b) at (1.6, 0) {};
  \node[vert] (c) at (0.8, -1.3) {};
  \draw[arista] (a) -- (b) -- (c) -- (a);
  \node[nota] at (0.8, 0.6) {simple};
  \begin{scope}[xshift=4.2cm]
    \node[vert] (u) at (0, 0) {$u$};
    \node[vert] (v) at (1.8, 0) {$v$};
    \node[vert] (x) at (0.9, -1.3) {$w$};
    \draw[arista] (u) to[bend left=25] (v);
    \draw[arista] (u) to[bend right=25] (v);
    \draw[arista] (v) -- (x);
    \draw[arista] (x) to[loop below, looseness=6] (x);
    \node[nota] at (0.9, 0.7) {no simple: arista múltiple y lazo};
  \end{scope}
\end{tikzpicture}
\end{center}

Una arista múltiple aparece, por ejemplo, con dos vuelos distintos
entre las mismas ciudades. Casi todos los algoritmos suponen grafos
simples: con aristas múltiples pesadas basta quedarse con la más
barata.

\subsubsection*{Encajado o topológico}

Un grafo es \textbf{encajado} (\textit{embedded}) si a sus vértices y
aristas se les asignan posiciones geométricas: un punto en el plano
para cada vértice y una curva para cada arista. Un grafo es
\textbf{topológico} si solo importa quién está unido con quién. Los
dos dibujos de abajo son el mismo grafo topológico, $K_4$, con cuatro
vértices todos unidos entre sí, en dos encajes distintos:

\begin{center}
\begin{tikzpicture}
  \node[vert] (p1) at (0, 0) {1};
  \node[vert] (p2) at (1.8, 0) {2};
  \node[vert] (p3) at (1.8, -1.8) {3};
  \node[vert] (p4) at (0, -1.8) {4};
  \foreach \a/\b in {p1/p2, p2/p3, p3/p4, p4/p1, p1/p3, p2/p4}
    \draw[arista] (\a) -- (\b);
  \begin{scope}[xshift=5cm]
    \node[vert] (q1) at (0, 0.2) {1};
    \node[vert] (q2) at (-1.2, -1.9) {2};
    \node[vert] (q3) at (1.2, -1.9) {3};
    \node[vert] (q4) at (0, -1.15) {4};
    \foreach \a/\b in {q1/q2, q2/q3, q3/q1, q4/q1, q4/q2, q4/q3}
      \draw[arista] (\a) -- (\b);
  \end{scope}
\end{tikzpicture}
\end{center}

A la izquierda dos aristas se cruzan; a la derecha ninguna. Como grafo
topológico es lo mismo; como grafo encajado no. En un mapa la
geometría es parte del problema: la distancia entre dos cruces sale de
sus coordenadas, y se puede saber si dos calles se cruzan.

\subsubsection*{Implícito o explícito}

Un grafo \textbf{explícito} se construye completo y se guarda en una
matriz o en listas antes de recorrerlo. Uno \textbf{implícito} no se
guarda nunca: los vecinos de un vértice se generan cuando se
necesitan, con una regla.

\begin{ejemplo}[Permutaciones como grafo implícito]
  Los vértices son las permutaciones de $1, 2, 3$, y dos son vecinas
  si una sale de la otra intercambiando dos elementos contiguos:
  \begin{center}
  \begin{tikzpicture}
    \foreach \p/\ang in {123/90, 213/30, 231/-30, 321/-90, 312/-150,
                         132/150}
      \node[vert, minimum size=0.85cm, font=\scriptsize] (p\p)
        at (\ang:1.6) {\p};
    \foreach \a/\b in {123/213, 213/231, 231/321, 321/312, 312/132,
                       132/123}
      \draw[arista] (p\a) -- (p\b);
  \end{tikzpicture}
  \end{center}
  Los vecinos de $123$ son $213$ y $132$, y se obtienen con la regla
  sin haber guardado nada. Con $k$ elementos el grafo tiene $k!$
  vértices: con $k = 10$ son $3\,628\,800$, y solo se generan los que
  el recorrido alcanza a visitar.
\end{ejemplo}

Los grafos implícitos aparecen en búsqueda: los estados de un
rompecabezas, las posiciones de un juego, las soluciones parciales de
un problema de backtracking.

\subsubsection*{Etiquetado o no etiquetado}

En un grafo \textbf{etiquetado} cada vértice tiene un nombre propio,
como A a G en $G$, y dos vértices con nombres distintos son
distintos. En uno \textbf{no etiquetado} los nombres no importan,
solo la estructura. La pregunta natural es si dos grafos son el mismo
salvo los nombres: si existe una biyección $f : V_1 \to V_2$ con
$\set{u, v} \in E_1$ si y solo si $\set{f(u), f(v)} \in E_2$. Es la
\textbf{prueba de isomorfismo}. No se conoce un algoritmo polinomial
para decidirla; en la práctica se hace con búsqueda con retroceso, que
descarta pronto las asignaciones imposibles comparando invariantes
como el número de aristas o la lista de grados. Los dos dibujos de
$K_4$ de arriba tienen las mismas etiquetas y las mismas aristas, así
que son el mismo grafo aun como grafos etiquetados; el isomorfismo
hace falta cuando los nombres no coinciden.

% =====================================================================

\subsection*{Búsqueda en anchura (BFS)}

Recorrer un grafo es visitar todos los vértices alcanzables desde uno
inicial $s$, cada uno una sola vez. La \textbf{búsqueda en anchura}
(\textit{breadth-first search}) visita por capas: primero $s$, luego
todos sus vecinos, luego los vecinos de esos que no se han visto, y
así. Para respetar ese orden usa una \textbf{cola}: el primer vértice
descubierto es el primero que se procesa.

\begin{algorithm}[H]
\caption{\textsc{BFS}$(G, s)$}
\begin{algorithmic}[1]
\ForAll{$v \in V$}
  \State $\text{visto}[v] \gets$ falso, \quad $d[v] \gets \infty$,
         \quad $\pi[v] \gets$ nulo
\EndFor
\State $\text{visto}[s] \gets$ verdadero, \quad $d[s] \gets 0$
\State $Q \gets$ cola vacía; \ \textsc{Encolar}$(Q, s)$
\While{$Q$ no está vacía}
  \State $u \gets \textsc{Desencolar}(Q)$
  \ForAll{$v$ en la lista de adyacencia de $u$}
    \If{\textbf{no} $\text{visto}[v]$}
      \State $\text{visto}[v] \gets$ verdadero, \quad
             $d[v] \gets d[u] + 1$, \quad $\pi[v] \gets u$
      \State \textsc{Encolar}$(Q, v)$
    \EndIf
  \EndFor
\EndWhile
\end{algorithmic}
\end{algorithm}

$d[v]$ es la distancia en aristas desde $s$ y $\pi[v]$ el vértice desde
el que se descubrió $v$, su \textbf{padre}.

\begin{observacion}
  Cormen escribe los atributos con punto, $v.d$ y $v.\pi$, y
  \textsc{nil} para el predecesor vacío; aquí son $d[v]$, $\pi[v]$ y
  nulo. En DFS Cormen usa $u.d$ y $u.f$ para los tiempos de
  descubrimiento y de terminación; aquí son $\text{entra}[u]$ y
  $\text{sale}[u]$, porque $d[v]$ ya es la distancia de BFS. Dentro de
  la notación asintótica Cormen escribe $V$ y $E$ por $\abs{V}$ y
  $\abs{E}$: su $O(V + E)$ es el $\Ord{n + m}$ de aquí.
\end{observacion}

\begin{ejemplo}[BFS sobre $G$ desde A]
  Los vecinos se recorren en orden alfabético.
  \begin{center}
  \small
  \begin{tabular}{cclc}
    \toprule
    \textbf{Paso} & \textbf{Sale} & \textbf{Descubre}
      & \textbf{Cola después} \\
    \midrule
    1 & A & B, C & [B, C] \\
    2 & B & D, E (A ya visto) & [C, D, E] \\
    3 & C & F (A ya visto) & [D, E, F] \\
    4 & D & G (B ya visto) & [E, F, G] \\
    5 & E & nada: B y F ya vistos & [F, G] \\
    6 & F & nada: C y E ya vistos & [G] \\
    7 & G & nada & [\,] \\
    \bottomrule
  \end{tabular}
  \end{center}
  Orden de visita: A, B, C, D, E, F, G. Las aristas por las que se
  descubrió cada vértice forman el \textbf{árbol BFS}:
  \begin{center}
  \begin{tikzpicture}
    \node[vert] (A) at (0, 0) {A};
    \node[vert] (B) at (-1.6, -1.2) {B};
    \node[vert] (C) at (1.6, -1.2) {C};
    \node[vert] (D) at (-2.6, -2.4) {D};
    \node[vert] (E) at (-0.6, -2.4) {E};
    \node[vert] (F) at (1.6, -2.4) {F};
    \node[vert] (G) at (-2.6, -3.6) {G};
    \foreach \a/\b in {A/B, A/C, B/D, B/E, C/F, D/G}
      \draw[aristaArbol] (\a) -- (\b);
    \draw[aristaFuera] (E) -- (F);
    \foreach \y/\d in {0/0, -1.2/1, -2.4/2, -3.6/3}
      \node[nota, anchor=west] at (3.0, \y) {$d = \d$};
  \end{tikzpicture}
  \end{center}
  Cada renglón es una capa: todos los vértices a la misma distancia de
  A. La arista $\set{E, F}$, punteada, no está en el árbol: cuando se
  procesó E, F ya estaba visto.
\end{ejemplo}

La inicialización cuesta $\Tht{n}$. Después, cada vértice alcanzable
entra y sale de la cola una sola vez, porque se marca al descubrirlo,
y su lista se recorre una vez cuando sale. Las listas suman $2m$
entradas en un grafo no dirigido, porque cada arista aparece en las
listas de sus dos extremos, y $m$ en uno dirigido. En total BFS cuesta
$\Ord{n + m}$ con listas. Es $\Tht{n + m}$ cuando $s$ alcanza a todos,
como en $G$; si no, puede ser menos, porque las listas de los
vértices inalcanzables nunca se recorren. Con matriz, encontrar los
vecinos de cada vértice que sale cuesta $\Tht{n}$, y el total es
$\Ord{n^2}$.

\begin{idea}
  En un grafo no pesado, $d[v]$ es la longitud del camino más corto de
  $s$ a $v$, y siguiendo $\pi$ de $v$ hacia atrás se reconstruye ese
  camino. BFS ya resuelve el camino más corto cuando todas las aristas
  valen lo mismo. La prueba está en la sección \ref{sec:busquedas}.
\end{idea}

\begin{ojo}
  En clase se dijo que la complejidad depende del problema, y que en un
  grafo completamente conectado es factorial. Hay que separar dos
  cosas. BFS procesa cada vértice una vez y cada arista desde sus dos
  extremos: en el grafo completo $K_n$, con $m = n(n-1)/2$, cuesta
  $\Tht{n^2}$, no factorial. Lo factorial aparece cuando el problema
  exige recorrer \textbf{caminos} y no vértices, como probar todas las
  rutas posibles: desde un vértice de $K_5$ salen $65$ caminos,
  contando el trivial de longitud cero, y desde uno de $K_{10}$,
  $986\,410$, alrededor de $e \cdot (n-1)!$. O cuando el grafo es
  implícito y sus vértices son permutaciones, como el ejemplo de
  arriba: BFS sigue siendo lineal en el tamaño del grafo, pero el
  grafo mide $k!$.
\end{ojo}

% =====================================================================

\subsection*{Búsqueda en profundidad (DFS)}

La \textbf{búsqueda en profundidad} (\textit{depth-first search}) se
mete lo más hondo que puede por un camino y solo regresa cuando se
atora: cuando el vértice actual ya no tiene vecinos sin visitar. Se
escribe de forma natural con recursión; la pila de llamadas hace el
papel que la cola hace en BFS. La versión iterativa usa una pila
explícita, con dos trampas que se ven en la sección
\ref{sec:busquedas}.

\begin{algorithm}[H]
\caption{\textsc{DFS-Visitar}$(G, u)$}
\begin{algorithmic}[1]
\State $t \gets t + 1$, \quad $\text{entra}[u] \gets t$, \quad
       $\text{visto}[u] \gets$ verdadero
\ForAll{$v$ en la lista de adyacencia de $u$}
  \If{\textbf{no} $\text{visto}[v]$}
    \State $\pi[v] \gets u$
    \State \textsc{DFS-Visitar}$(G, v)$
  \EndIf
\EndFor
\State $t \gets t + 1$, \quad $\text{sale}[u] \gets t$
\end{algorithmic}
\end{algorithm}

$t$ es un reloj global que empieza en $0$. $\text{entra}[u]$ marca
cuándo se descubrió $u$ y $\text{sale}[u]$ cuándo se terminó con todo
lo que cuelga de él. Antes de la primera llamada se pone
$\text{visto}[v] \gets$ falso para todo $v$ y $t \gets 0$. Una llamada
desde $u$ solo alcanza lo que es accesible desde $u$; para recorrer un
grafo desconectado se lanza \textsc{DFS-Visitar} desde cada vértice
que siga sin visitar, y el resultado es un bosque. Ese procedimiento
externo, \textsc{DFS}$(G)$, está en la sección \ref{sec:busquedas}.

\begin{ejemplo}[DFS sobre $G$ desde A]
  Con los vecinos en orden alfabético, como en BFS:
  \begin{itemize}[itemsep=0.1em]
    \item A $\to$ B $\to$ D $\to$ G. G no tiene vecinos nuevos: se
          termina, y se regresa a D, que tampoco tiene; se regresa a B.
    \item B $\to$ E $\to$ F $\to$ C. Los vecinos de C son A y F, los dos
          vistos: se termina C, luego F, E, B y A.
  \end{itemize}
  Orden de visita: A, B, D, G, E, F, C. El de BFS era A, B, C, D, E, F,
  G.
  \begin{center}
  \begin{tikzpicture}
    \node[vert] (A) at (0, 0) {A};
    \node[vert] (B) at (0, -1.1) {B};
    \node[vert] (D) at (-1.3, -2.2) {D};
    \node[vert] (G) at (-1.3, -3.3) {G};
    \node[vert] (E) at (1.3, -2.2) {E};
    \node[vert] (F) at (1.3, -3.3) {F};
    \node[vert] (C) at (1.3, -4.4) {C};
    \foreach \a/\b in {A/B, B/D, D/G, B/E, E/F, F/C}
      \draw[aristaArbol] (\a) -- (\b);
    \draw[aristaFuera] (C) to[out=0, in=0, looseness=1.4] (A);
    \foreach \v/\t/\lado in {A/1--14/left, B/2--13/left, D/3--6/left,
                             G/4--5/left, E/7--12/left, F/8--11/left,
                             C/9--10/left}
      \node[nota, \lado=4pt] at (\v) {\t};
  \end{tikzpicture}
  \end{center}
  Junto a cada vértice, $\text{entra}$--$\text{sale}$. Este árbol DFS,
  el de este orden de las listas, es largo y angosto; el BFS, ancho y
  bajo. La arista
  $\set{A, C}$, punteada, es una \textbf{arista de retroceso}: se
  explora primero desde C, cuando A sigue abierto, y va de C a un
  antecesor suyo en el árbol.
\end{ejemplo}

DFS también cuesta $\Ord{n + m}$ con listas, pero la cuenta es otra,
porque no hay cola: \textsc{DFS-Visitar} se llama una vez por vértice,
ya que lo primero que hace es marcarlo, y cada llamada recorre su
lista una vez. Recorrer todo el grafo cuesta $\Tht{n + m}$.

\begin{center}
\small
\begin{tabular}{lll}
  \toprule
  & \textbf{BFS} & \textbf{DFS} \\
  \midrule
  Estructura & cola & pila o recursión \\
  Orden & por capas, de cerca a lejos & un camino hasta el fondo \\
  Árbol & ancho y bajo & largo y angosto \\
  Útil para & camino más corto sin pesos & ciclos dirigidos, orden
    topológico, componentes \\
  Costo con listas & $\Ord{n + m}$ & $\Ord{n + m}$ \\
  \bottomrule
\end{tabular}
\end{center}

Las aplicaciones de DFS de la tabla se usan aquí sin prueba: un grafo
dirigido tiene ciclo si y solo si DFS encuentra una arista de
retroceso; en un DAG, ordenar por $\text{sale}$ decreciente da un
orden topológico; en uno no dirigido, los árboles DFS son las
componentes conexas (secciones \ref{sec:aplicaciones-dfs} y
\ref{sec:busquedas}).

% =====================================================================

\subsection*{Camino más corto: Dijkstra}

Dado un grafo pesado con pesos \textbf{no negativos} y un vértice $s$,
el algoritmo de Dijkstra (1959) encuentra la distancia $\delta(s, v)$
de $s$ a todos los demás: el peso del camino más barato. Mantiene una
estimación $d[v]$ que empieza en $\infty$ y nunca baja de
$\delta(s, v)$, porque siempre es el peso de una caminata de $s$ a
$v$. En cada paso fija el vértice no
fijado con la $d$ más chica: esa ya es su distancia definitiva. Luego
\textbf{relaja} las aristas que salen de él: si llegar a un vecino
pasando por $u$ es más barato que lo que se tenía, se actualiza.

\begin{algorithm}[H]
\caption{\textsc{Dijkstra}$(G, w, s)$}
\begin{algorithmic}[1]
\ForAll{$v \in V$}
  \State $d[v] \gets \infty$, \quad $\pi[v] \gets$ nulo
\EndFor
\State $d[s] \gets 0$
\State $Q \gets$ montículo de mínimos con todos los vértices, con
       llave $d$
\While{$Q$ no está vacío}
  \State $u \gets \textsc{Extraer-Min}(Q)$ \Comment{$d[u]$ queda fija}
  \ForAll{$v$ adyacente a $u$}
    \If{$d[u] + w(u, v) < d[v]$} \Comment{relajar}
      \State $d[v] \gets d[u] + w(u, v)$, \quad $\pi[v] \gets u$
      \State \textsc{Decrementar-Llave}$(Q, v, d[v])$
    \EndIf
  \EndFor
\EndWhile
\end{algorithmic}
\end{algorithm}

\begin{ejemplo}[Dijkstra sobre el grafo pesado desde A]
  Cada renglón muestra las $d$ después de procesar el vértice que
  sale; en negritas, las ya fijas.
  \begin{center}
  \small
  \begin{tabular}{ccccccccl}
    \toprule
    \textbf{Paso} & \textbf{Sale} & A & B & C & D & E & F
      & \textbf{Relajaciones} \\
    \midrule
    0 & & $0$ & $\infty$ & $\infty$ & $\infty$ & $\infty$ & $\infty$
      & \\
    1 & A & $\mathbf{0}$ & $4$ & $2$ & $7$ & $\infty$ & $\infty$
      & B, C, D desde A \\
    2 & C & $\mathbf{0}$ & $3$ & $\mathbf{2}$ & $7$ & $12$ & $\infty$
      & B: $2 + 1 = 3 < 4$; D: $10$, no \\
    3 & B & $\mathbf{0}$ & $\mathbf{3}$ & $\mathbf{2}$ & $7$ & $12$
      & $\infty$ & D: $3 + 5 = 8$, no mejora \\
    4 & D & $\mathbf{0}$ & $\mathbf{3}$ & $\mathbf{2}$ & $\mathbf{7}$
      & $9$ & $13$ & E: $7 + 2 = 9 < 12$ \\
    5 & E & $\mathbf{0}$ & $\mathbf{3}$ & $\mathbf{2}$ & $\mathbf{7}$
      & $\mathbf{9}$ & $12$ & F: $9 + 3 = 12 < 13$ \\
    6 & F & $\mathbf{0}$ & $\mathbf{3}$ & $\mathbf{2}$ & $\mathbf{7}$
      & $\mathbf{9}$ & $\mathbf{12}$ & \\
    \bottomrule
  \end{tabular}
  \end{center}
  B empezó en $4$ por la arista directa y bajó a $3$ pasando por C. Los
  padres quedan $\pi[B] = C$, $\pi[C] = A$, $\pi[D] = A$,
  $\pi[E] = D$, $\pi[F] = E$, y el camino a F se lee de atrás hacia
  adelante: F $\leftarrow$ E $\leftarrow$ D $\leftarrow$ A, o sea
  A, D, E, F, con $7 + 2 + 3 = 12$.
  \begin{center}
  \begin{tikzpicture}
    \node[vert] (A) at (0, 0) {A};
    \node[vert] (B) at (2, 1.4) {B};
    \node[vert] (C) at (2, -1.4) {C};
    \node[vert] (D) at (4.5, 1.4) {D};
    \node[vert] (E) at (4.5, -1.4) {E};
    \node[vert] (F) at (6.5, 0) {F};
    \foreach \a/\b/\w in {A/B/4, B/D/5, C/D/8, C/E/10, D/F/6}
      \draw[arista] (\a) -- node[peso] {\w} (\b);
    \foreach \a/\b/\w in {A/C/2, C/B/1, D/E/2, E/F/3}
      \draw[aristaArbol] (\a) -- node[peso] {\w} (\b);
    \draw[aristaArbol] (A) to[out=75, in=150] node[peso] {7} (D);
    \foreach \v/\d/\lado in {A/0/left, B/3/below right, C/2/below,
                             D/7/above right, E/9/below, F/12/right}
      \node[nota, \lado=3pt] at (\v) {$d = \d$};
  \end{tikzpicture}
  \end{center}
  En azul, el \textbf{árbol de caminos más cortos}: las aristas
  $\set{\pi[v], v}$. Que el camino de A a cada vértice dentro de este
  árbol sea uno más corto en el grafo se demuestra en la sección
  \ref{sec:caminos-pruebas}.
\end{ejemplo}

\textbf{Por qué funciona.} Supóngase que todos los vértices fijados
antes que $u$ tienen ya su distancia correcta, y tómese un camino más
corto $P$ de $s$ a $u$. (Si $u = s$, o si no hay camino a $u$, no hay
nada que probar.) Sea $y$ el primer vértice de $P$ que no está fijo,
que puede ser el propio $u$, y $x$ el anterior, que sí lo está. Cuando
$x$ se fijó tenía $d[x] = \delta(s, x)$ y relajó la arista hacia $y$;
como esa arista está en un camino más corto, desde entonces
$d[y] = \delta(s, y)$. Esto es lo esencial: no basta con que $y$ tenga
$d[y] \geq d[u]$, hace falta que su $d$ ya sea exacta. El tramo de $P$
que va de $y$ a $u$ no resta, porque los pesos son no negativos, así
que
\[
  d[y] = \delta(s, y) \leq \delta(s, u) \leq d[u] ,
\]
y como el algoritmo eligió $u$, también $d[u] \leq d[y]$. Todo es
igual y $d[u] = \delta(s, u)$. Es un algoritmo voraz: la elección del
momento nunca se revisa, y con pesos no negativos no hace falta. La
prueba completa, con las propiedades de la relajación que usa, está en
la sección \ref{sec:caminos-fuente}.

Dos consecuencias. Con $w \geq 0$ la condición del \textbf{si} nunca
se cumple para un vértice que ya salió: si $v$ está fijo,
$d[v] = \delta(s, v) \leq \delta(s, u) + w(u, v) = d[u] + w(u, v)$.
Así que \textsc{Decrementar-Llave} siempre actúa sobre un vértice que
sigue en el montículo. Y si solo interesa un destino $t$, se puede
parar en cuanto $t$ sale del montículo, no cuando apenas se descubre,
porque su $d$ todavía puede bajar, como la de B en el ejemplo.

\textbf{Costo.} Cada vértice sale del montículo una sola vez, y solo
entonces se recorre su lista. Las listas suman $2m$ entradas ($m$ en
un dirigido), así que hay $\Ord{m}$ relajaciones y cada una hace a lo
más un \textsc{Decrementar-Llave}. Con montículo binario, construirlo
cuesta $\Ord{n}$ y las $n$ extracciones y los $\Ord{m}$ decrementos
cuestan $\Ord{\log n}$ cada uno: $\Ord{(n + m) \log n}$ en total. Es
el montículo de la clase del 17 de septiembre, con un detalle: para
decrementar la llave de $v$ hay que saber dónde está $v$ dentro del
arreglo, así que cada vértice guarda su posición y el montículo la
actualiza en cada intercambio. Con un arreglo indexado por vértice,
decrementar es una asignación pero buscar el mínimo cuesta $\Tht{n}$,
y el total es $\Tht{n^2}$. El montículo gana cuando
$m = \ord{n^2 / \log n}$; en un grafo denso, $m = \Tht{n^2}$, cuesta
$\Ord{n^2 \log n}$ y gana el arreglo. Con montículo de Fibonacci se
llega a $\Ord{n \log n + m}$; esa y las demás colas están en la
sección \ref{sec:caminos-fuente}.

\begin{ojo}
  Con un peso negativo Dijkstra deja de ser confiable. Con las aristas
  dirigidas $(A, B)$ de peso $2$, $(A, C)$ de peso $3$, $(C, B)$ de
  peso $-2$ y $(B, D)$ de peso $2$, desde A se fija B con $d = 2$
  antes de procesar C, y B relaja $(B, D)$: $d[D] = 4$. Al procesar
  C, la relajación de $(C, B)$ baja $d[B]$ a $3 - 2 = 1$, porque el
  pseudocódigo de arriba relaja aunque el vecino ya esté fijo, pero
  nadie vuelve a relajar $(B, D)$. D termina con $4$, y el camino A,
  C, B, D cuesta $3$. Sin la arista $(B, D)$ el resultado final es
  correcto gracias a esa relajación tardía; lo que falla es que B se
  dio por fija con un valor que luego cambió, y que se llamó
  \textsc{Decrementar-Llave} sobre un vértice que ya no estaba en el
  montículo. El argumento de arriba suponía que el tramo de $y$ a $u$
  no resta; aquí, cuando sale B, ese tramo es la arista $(C, B)$. Para
  pesos negativos están Bellman-Ford (sección \ref{sec:caminos-fuente})
  y, entre todos los pares, Floyd y Johnson (sección
  \ref{sec:todos-pares}).
\end{ojo}

% =====================================================================

\subsection*{Camino más corto entre todos los pares: Floyd}

Dijkstra da las distancias desde un vértice. Para tener las de todos
contra todos se puede correr $n$ veces, una por fuente: con montículo
binario son $\Ord{n (n + m) \log n}$. Floyd (Floyd-Warshall, 1962) lo
resuelve con programación dinámica en tres ciclos, en $\Tht{n^3}$, y
admite pesos negativos siempre que no haya ciclos de peso negativo.
Ninguno gana siempre. En un grafo disperso, $m = \Ord{n}$, las $n$
corridas de Dijkstra cuestan $\Ord{n^2 \log n}$ y le ganan a Floyd; en
uno denso, $m = \Tht{n^2}$, cuestan $\Ord{n^3 \log n}$, y Floyd gana
además por sencillez.

Con los vértices numerados de $1$ a $n$, un vértice
\textbf{intermedio} de un camino es cualquiera de sus vértices salvo
los dos extremos. Sea $d^{(k)}_{ij}$ el \textbf{peso} de un camino más
corto de $i$ a $j$ cuyos intermedios están todos en
$\set{1, \dots, k}$. Con $k = 0$ no se permite ningún intermedio y
solo queda la arista directa. Para pasar de $k - 1$ a $k$ hay dos
opciones: el camino no usa $k$, y queda igual; o lo usa. Como no hay
ciclos negativos, ese camino se puede tomar sin vértices repetidos,
así que $k$ aparece una sola vez y lo parte en dos tramos, de $i$ a
$k$ y de $k$ a $j$, cuyos intermedios están en
$\set{1, \dots, k - 1}$. Los dos ya están calculados:
\[
  d^{(k)}_{ij} = \min\left( d^{(k-1)}_{ij},\
                 d^{(k-1)}_{ik} + d^{(k-1)}_{kj} \right) ,
\]
con $d^{(0)}_{ij} = w(i, j)$, $0$ en la diagonal e $\infty$ si no hay
arista. La respuesta es $D^{(n)}$, porque con $k = n$ ya se permite
cualquier intermedio. La hipótesis de no tener ciclos negativos se usa
justo en el segundo caso: si un ciclo negativo pasa por $k$, cada
vuelta abarata el camino, no hay camino más corto y la recurrencia ya
no describe nada. Floyd lo delata: al terminar, algún $D[i][i]$ es
negativo (sección \ref{sec:todos-pares}).

\begin{algorithm}[H]
\caption{\textsc{Floyd}$(W, n)$}
\begin{algorithmic}[1]
\State $D \gets W$
\For{$k \gets 1$ \Hasta $n$}
  \For{$i \gets 1$ \Hasta $n$}
    \For{$j \gets 1$ \Hasta $n$}
      \State $D[i][j] \gets \min(D[i][j],\ D[i][k] + D[k][j])$
    \EndFor
  \EndFor
\EndFor
\State \Return $D$
\end{algorithmic}
\end{algorithm}

Son $\Tht{n^3}$ operaciones y $\Tht{n^2}$ de memoria, porque la misma
matriz se puede actualizar en su lugar. En la vuelta $k$ el renglón
$k$ y la columna $k$ no cambian: sin ciclos negativos $D[k][k] = 0$,
y entonces $D[i][k] + D[k][k] = D[i][k]$. Como $D[i][k]$ y $D[k][j]$
son las únicas casillas que se leen además de $D[i][j]$, da igual
leerlas antes o después de actualizarlas.

\begin{ejemplo}[Floyd sobre el grafo dirigido]
  Los pesos de las aristas del grafo dirigido de arriba:
  \begin{center}
  \begin{tikzpicture}
    \node[vert] (n1) at (0, 0) {1};
    \node[vert] (n2) at (2.8, 0) {2};
    \node[vert] (n3) at (2.8, -2.4) {3};
    \node[vert] (n4) at (0, -2.4) {4};
    \draw[dirigida] (n1) to[bend left=15] node[peso] {3} (n2);
    \draw[dirigida] (n2) to[bend left=15] node[peso] {8} (n1);
    \draw[dirigida] (n1) to[bend right=15] node[peso] {7} (n4);
    \draw[dirigida] (n4) to[bend right=15] node[peso] {2} (n1);
    \draw[dirigida] (n2) -- node[peso] {2} (n3);
    \draw[dirigida] (n3) -- node[peso] {5} (n1);
    \draw[dirigida] (n3) -- node[peso] {1} (n4);
  \end{tikzpicture}
  \end{center}
  En negritas, las casillas que cambian en cada paso:
  \[
    D^{(0)} = \begin{pmatrix}
      0 & 3 & \infty & 7 \\
      8 & 0 & 2 & \infty \\
      5 & \infty & 0 & 1 \\
      2 & \infty & \infty & 0
    \end{pmatrix}
    \qquad
    D^{(1)} = \begin{pmatrix}
      0 & 3 & \infty & 7 \\
      8 & 0 & 2 & \mathbf{15} \\
      5 & \mathbf{8} & 0 & 1 \\
      2 & \mathbf{5} & \infty & 0
    \end{pmatrix}
    \qquad
    D^{(2)} = \begin{pmatrix}
      0 & 3 & \mathbf{5} & 7 \\
      8 & 0 & 2 & 15 \\
      5 & 8 & 0 & 1 \\
      2 & 5 & \mathbf{7} & 0
    \end{pmatrix}
  \]
  \[
    D^{(3)} = \begin{pmatrix}
      0 & 3 & 5 & \mathbf{6} \\
      \mathbf{7} & 0 & 2 & \mathbf{3} \\
      5 & 8 & 0 & 1 \\
      2 & 5 & 7 & 0
    \end{pmatrix}
    \qquad
    D^{(4)} = \begin{pmatrix}
      0 & 3 & 5 & 6 \\
      \mathbf{5} & 0 & 2 & 3 \\
      \mathbf{3} & \mathbf{6} & 0 & 1 \\
      2 & 5 & 7 & 0
    \end{pmatrix}
  \]
  Algunos cambios, uno por paso:
  \begin{itemize}[itemsep=0.1em]
    \item $k = 1$: de $3$ a $2$ no había arista; pasando por $1$,
          $d_{31} + d_{12} = 5 + 3 = 8$.
    \item $k = 2$: de $1$ a $3$, $d_{12} + d_{23} = 3 + 2 = 5$.
    \item $k = 3$: de $2$ a $4$ bajó de $15$ a $d_{23} + d_{34} =
          2 + 1 = 3$.
    \item $k = 4$: de $3$ a $1$ bajó de $5$ a $d_{34} + d_{41} =
          1 + 2 = 3$.
  \end{itemize}
  $D^{(4)}$ tiene todas las distancias. Por ejemplo
  $d_{32} = 6$ es el camino $3 \to 4 \to 1 \to 2$, con
  $1 + 2 + 3 = 6$. Ese camino se lee del dibujo: el pseudocódigo solo
  calcula pesos. Para que el algoritmo lo entregue se guarda además
  una matriz de predecesores, que está en la sección
  \ref{sec:todos-pares}.
\end{ejemplo}

\begin{center}
\small
\begin{tabular}{lll}
  \toprule
  & \textbf{Dijkstra} & \textbf{Floyd} \\
  \midrule
  Responde & desde un vértice & todos los pares \\
  Pesos negativos & no & sí, sin ciclos negativos \\
  Detecta ciclos negativos & no & sí, con la diagonal \\
  Técnica & voraz & programación dinámica \\
  Costo de una corrida & $\Ord{(n + m)\log n}$ & $\Tht{n^3}$ \\
  Costo para todos los pares & $n$ corridas:
    $\Ord{n (n + m)\log n}$ & $\Tht{n^3}$ \\
  Conviene en & grafos dispersos & grafos densos \\
  \bottomrule
\end{tabular}
\end{center}

% =====================================================================

\subsection*{Encontrar ciclos con BFS}

En un grafo no dirigido, BFS encuentra un ciclo en cuanto ve una
arista que no es del árbol: al procesar $u$, un vecino $v$ que ya fue
visto y que no es el padre de $u$. El padre no cuenta porque la arista
$\set{u, \pi[u]}$ es justo por la que se llegó a $u$. Basta cambiar el
\textbf{si} de BFS:

\begin{algorithm}[H]
\caption{Cambio a \textsc{BFS} para detectar un ciclo}
\begin{algorithmic}[1]
\ForAll{$v$ en la lista de adyacencia de $u$}
  \If{\textbf{no} $\text{visto}[v]$}
    \State \ldots igual que en BFS \ldots
  \ElsIf{$v \neq \pi[u]$}
    \State \Return hay un ciclo que pasa por la arista $\set{u, v}$
  \EndIf
\EndFor
\end{algorithmic}
\end{algorithm}

\begin{ejemplo}[El ciclo de $G$]
  En el paso 5 de BFS sale E y su vecino F ya está visto, pero
  $\pi[E] = B \neq F$. Hay ciclo. Para recuperarlo se sube por los
  padres desde los dos extremos hasta que se encuentran:
  \[
    E \to B \to A , \qquad F \to C \to A .
  \]
  Se juntan en A, y el ciclo es A, B, E, F, C, A: la arista punteada
  del árbol BFS cierra la vuelta.
\end{ejemplo}

\begin{ojo}
  En un grafo dirigido esta regla da falsos positivos. Con las aristas
  $(A, B)$, $(A, C)$ y $(B, C)$, BFS desde A llega a C dos veces y la
  regla dice ciclo, pero no hay ningún ciclo dirigido. Para grafos
  dirigidos se usa DFS: hay un ciclo si y solo si DFS encuentra una
  arista de retroceso, una que llega a un vértice todavía abierto, que
  ya entró y no ha salido. No basta $\text{visto}$: en el mismo
  ejemplo, al regresar a A y revisar $(A, C)$, C está visto pero ya
  terminado. Por eso DFS usa tres colores (sección
  \ref{sec:busquedas}); la prueba está en la sección
  \ref{sec:aplicaciones-dfs}. La otra opción es el algoritmo
  de Kahn: con una cola se van quitando los vértices de grado de
  entrada $0$; si alguno nunca llega a $0$, hay ciclo. En el grafo de
  Floyd ningún vértice empieza con $\text{gr}^{-}(v) = 0$, así que la
  cola nace vacía: tiene ciclos, como $1 \to 2 \to 1$.
\end{ojo}

% =====================================================================

\subsection*{Árbol de expansión mínima: Kruskal}

\begin{definicion}[Árbol de expansión mínima]
  En un grafo no dirigido, conexo y pesado, un \textbf{árbol de
  expansión} es un subconjunto de aristas que conecta todos los
  vértices sin formar ciclos; tiene exactamente $n - 1$ aristas. El
  de \textbf{expansión mínima} (MST, \textit{minimum spanning tree})
  es el de menor peso total.
\end{definicion}

Es el problema de conectar $n$ ciudades con cable gastando lo menos
posible. El algoritmo de Kruskal (1956) es voraz: ordena las aristas de
la más ligera a la más pesada y las agrega en ese orden, saltándose las
que cerrarían un ciclo.

\begin{algorithm}[H]
\caption{\textsc{Kruskal}$(G, w)$}
\begin{algorithmic}[1]
\State $T \gets \emptyset$
\ForAll{$v \in V$}
  \State \textsc{Crear-Conjunto}$(v)$
\EndFor
\State ordenar $E$ por peso, de menor a mayor
\ForAll{$\set{u, v} \in E$, en ese orden}
  \If{$\textsc{Encontrar}(u) \neq \textsc{Encontrar}(v)$}
    \State $T \gets T \cup \set{\set{u, v}}$
    \State \textsc{Unir}$(u, v)$
  \EndIf
\EndFor
\State \Return $T$
\end{algorithmic}
\end{algorithm}

Para saber si una arista cierra un ciclo se usa una estructura de
\textbf{conjuntos disjuntos} (\textit{union-find}): cada vértice
empieza en su propio conjunto, \textsc{Encontrar} dice en qué conjunto
está un vértice y \textsc{Unir} junta dos. Si $u$ y $v$ están en el
mismo conjunto, hay un camino entre ellos y la arista cerraría un
ciclo.

\begin{ejemplo}[Kruskal sobre el grafo pesado]
  Las aristas en orden de peso:
  \begin{center}
  \small
  \begin{tabular}{cclll}
    \toprule
    \textbf{Arista} & \textbf{Peso} & \textbf{¿Entra?}
      & \textbf{Componentes después} \\
    \midrule
    $\set{B, C}$ & $1$ & sí & \{A\} \{B, C\} \{D\} \{E\} \{F\} \\
    $\set{A, C}$ & $2$ & sí & \{A, B, C\} \{D\} \{E\} \{F\} \\
    $\set{D, E}$ & $2$ & sí & \{A, B, C\} \{D, E\} \{F\} \\
    $\set{E, F}$ & $3$ & sí & \{A, B, C\} \{D, E, F\} \\
    $\set{A, B}$ & $4$ & no: A y B ya están juntos & \\
    $\set{B, D}$ & $5$ & sí & \{A, B, C, D, E, F\} \\
    \bottomrule
  \end{tabular}
  \end{center}
  Con $n - 1 = 5$ aristas el árbol está completo, y las que faltan
  ($\set{D, F}$, $\set{A, D}$, $\set{C, D}$, $\set{C, E}$) cerrarían
  ciclos. Peso total: $1 + 2 + 2 + 3 + 5 = 13$.
  \begin{center}
  \begin{tikzpicture}[scale=0.82, transform shape]
    \foreach \pre/\x in {s/0, k/8.2} {
      \begin{scope}[xshift=\x cm, name prefix=\pre-]
        \node[vert] (A) at (0, 0) {A};
        \node[vert] (B) at (2, 1.4) {B};
        \node[vert] (C) at (2, -1.4) {C};
        \node[vert] (D) at (4.5, 1.4) {D};
        \node[vert] (E) at (4.5, -1.4) {E};
        \node[vert] (F) at (6.5, 0) {F};
        \foreach \a/\b in {A/B, B/C, B/D, C/D, C/E, D/E, D/F, E/F, A/C}
          \draw[arista] (\a) -- (\b);
        \draw[arista] (A) to[out=75, in=150] (D);
      \end{scope}
    }
    \foreach \a/\b in {A/C, C/B, D/E, E/F}
      \draw[aristaArbol] (s-\a) -- (s-\b);
    \draw[aristaArbol] (s-A) to[out=75, in=150] (s-D);
    \foreach \a/\b in {B/C, A/C, D/E, E/F, B/D}
      \draw[aristaArbol] (k-\a) -- (k-\b);
    \node[font=\small] at (3.25, -2.3)
      {Caminos más cortos desde A: peso 15};
    \node[font=\small] at (11.45, -2.3) {Expansión mínima: peso 13};
  \end{tikzpicture}
  \end{center}
\end{ejemplo}

\begin{ojo}
  El árbol de caminos más cortos de Dijkstra y el árbol de expansión
  mínima no son lo mismo, aunque los dos tengan $n - 1$ aristas. El de
  Dijkstra minimiza la distancia de A a cada vértice; el MST, el peso
  total. En el MST el camino de A a D es A, C, B, D, con $8$; en el de
  Dijkstra es la arista directa, con $7$. Pero el MST pesa $13$ y el de
  Dijkstra $15$.
\end{ojo}

\textbf{Por qué funciona.} Se parte $V$ en dos lados de modo que
ninguna arista ya aceptada cruce de un lado al otro. La arista más
ligera que cruza, $e$, es \textbf{segura}: si las aceptadas caben en
un MST $T$, también caben junto con $e$ en uno. Si $T$ no usa $e$,
agregarla cierra un ciclo que cruza el corte por otra arista $e'$, que
no es de las aceptadas; cambiar $e'$ por $e$ da un árbol de expansión
que pesa $w(T) - w(e') + w(e) \leq w(T)$. El intercambio no empeora,
y con empates empata en lugar de mejorar; así que el árbol nuevo
también es mínimo y contiene a $e$ y a las aceptadas. Pedir que el
corte no separe aristas aceptadas es lo que permite encadenar los
pasos: no basta con que $e$ esté en algún MST, tiene que estar en el
mismo que las anteriores. Cada arista que Kruskal acepta es la más
ligera que sale del componente de uno de sus extremos, porque las más
ligeras ya se revisaron y quedaron dentro de algún componente; así que
es segura. Es la misma idea voraz de la clase del 8 de octubre, y la
sección \ref{sec:mst} la demuestra completa como la propiedad del
corte.

\textbf{Costo.} Ordenar las aristas cuesta $\Ord{m \log m}$. Con unión
por rango y compresión de caminos, la sucesión completa de operaciones
de conjuntos disjuntos ($n$ de \textsc{Crear-Conjunto}, $2m$ de
\textsc{Encontrar} y $n - 1$ de \textsc{Unir}) cuesta
$\Ord{(n + m)\,\alpha(n)}$, donde $\alpha$ es la inversa de la función
de Ackermann: crece tan despacio que $\alpha(n) \leq 4$ para cualquier
$n$ concebible en la práctica (Cormen, sección 21.4). Como $G$ es
conexo, $m \geq n - 1$ y eso es $\Ord{m\,\alpha(n)}$. Domina el
ordenamiento: $\Ord{m \log m} = \Ord{m \log n}$, porque $m < n^2$
implica $\log m < 2 \log n$.

\begin{observacion}
  El otro algoritmo clásico es el de Prim: hace crecer un solo árbol
  desde un vértice, agregando siempre la arista más ligera que sale de
  él. Es casi idéntico a Dijkstra: la llave de $v$ pasa de
  $d[u] + w(u, v)$ a solo $w(u, v)$, y se actualiza solo si $v$ sigue
  en $Q$. Esa condición sobra en Dijkstra, donde con pesos no
  negativos la relajación nunca dispara hacia un vértice fijo, pero en
  Prim es indispensable. Sin ella, en el grafo pesado desde A, al
  sacar B (llave $1$) se compara $w(B, C) = 1$ con la llave $2$ de C,
  que ya está en el árbol, y se cambia $\pi[C]$ a B cuando
  $\pi[B] = C$. Al final los padres forman dos ciclos, B con C y D con
  E, y ningún vértice cuelga de A. Con montículo binario Prim cuesta
  $\Ord{m \log n}$, como Kruskal; está en la sección \ref{sec:mst}.
\end{observacion}

% =====================================================================

\subsection*{Resumen}

\begin{center}
\small
\begin{tabular}{l>{\raggedright\arraybackslash}p{4.4cm}%
                 >{\raggedright\arraybackslash}p{3.2cm}l}
  \toprule
  \textbf{Algoritmo} & \textbf{Resuelve} & \textbf{Estructura}
    & \textbf{Costo} \\
  \midrule
  BFS & recorrido por capas; camino más corto sin pesos; ciclos en no
        dirigidos & cola & $\Ord{n + m}$ \\
  \addlinespace
  DFS & recorrido en profundidad; ciclos dirigidos; orden topológico
      & pila o recursión & $\Ord{n + m}$ \\
  \addlinespace
  Dijkstra & caminos más cortos desde un vértice, pesos $\geq 0$
           & montículo binario & $\Ord{(n + m) \log n}$ \\
           & & arreglo & $\Tht{n^2}$ \\
  \addlinespace
  Floyd & caminos más cortos entre todos los pares, sin ciclos
          negativos & matriz & $\Tht{n^3}$ \\
  \addlinespace
  Kruskal & árbol de expansión mínima & conjuntos disjuntos
          & $\Ord{m \log n}$ \\
  \addlinespace
  Prim & árbol de expansión mínima & montículo binario
       & $\Ord{m \log n}$ \\
  \bottomrule
\end{tabular}
\end{center}

BFS y DFS cuestan $\Tht{n + m}$ cuando recorren todo el grafo. Floyd
usa $\Tht{n^2}$ de memoria. El resumen de toda la parte de grafos, con
lo que agregan las secciones de complemento, está al final de la
sección \ref{sec:push-relabel}.

\newpage

% =====================================================================
%  Grafos - Teoria de grafos
% =====================================================================

\clase{Teoría de grafos: terminología, conexidad y Euler}%
  {Complemento: Rosen, cap.\ 10}
\label{sec:teoria-grafos}

\begin{quick_sum}
  Esta sección arma el vocabulario que el resto de la parte de grafos
  da por sabido. Parte del problema de los puentes de Königsberg y
  nombra lo que la clase del 22 de octubre dejó sin nombre:
  multigrafos y pseudografos, el grado con lazos y el teorema del
  apretón de manos, subgrafos, la matriz de incidencia y el
  isomorfismo. Fija la convención de caminata, sendero, camino,
  circuito y ciclo que usan todas las secciones siguientes, define
  conexidad, componentes, puentes y vértices de corte, y cierra con
  los grafos eulerianos: el teorema de Euler con su prueba completa,
  el caso dirigido y el algoritmo de Hierholzer, que construye el
  recorrido en $\Tht{n + m}$. La idea que se lleva el lector: decidir
  si existe un recorrido que use cada arista una vez, y construirlo,
  cuesta lo mismo que leer el grafo.
\end{quick_sum}

La materia prima son las diapositivas \textit{Grafos 1} de la
Dra.\ Kelsey Ramírez, que siguen de cerca a Rosen (cap.\ 10) y
recomiendan además a Grimaldi (cap.\ 11). Son teoría pura: no traen
pseudocódigo ni complejidades. Las pruebas de suficiencia, el caso
dirigido y el algoritmo de Hierholzer son agregados, para que la
libreta responda con un procedimiento lo que la clase decide mirando
el dibujo. Lo que ya está en la sección~\ref{sec:grafos} (la
definición formal, adyacencia, incidencia y grado en grafos simples,
la matriz y las listas de adyacencia, la clasificación de Skiena) no
se repite: se cita y se complementa.

% =====================================================================

\subsection*{Motivación: los siete puentes de Königsberg}

En el siglo \textsc{xviii} la ciudad de Königsberg (hoy Kaliningrado)
estaba partida por el río Pregel en cuatro regiones: la orilla norte
$N$, la orilla sur $S$, la isla Kneiphof $I$ y la región del este
$E$, entre los dos brazos del río. Siete puentes las unían. La
pregunta era si se podía dar un paseo que cruzara cada puente
exactamente una vez y regresara al punto de partida.

Euler (1736) notó que la respuesta no depende de la forma de las
islas, de las distancias ni del tamaño de los puentes: solo importa
qué regiones une cada puente. Cada región se vuelve un vértice y cada
puente una arista. Es el primer grafo de la historia, y la figura
muestra las dos versiones: a la izquierda el mapa, redibujado de forma
esquemática; a la derecha el grafo, con el grado de cada vértice.

% Konigsberg: mapa esquematico y su grafo (diapositivas 5, 112 y 125)
\begin{center}
\begin{tikzpicture}[x=1cm, y=1cm,
  tierra/.style={fill=green!8, draw=gray!60},
  puente/.style={draw=brown!70!black, line width=3.2pt}]
  \fill[green!8] (0, 1.3) rectangle (6.4, 2.0);
  \fill[green!8] (0, -2.0) rectangle (6.4, -1.3);
  \fill[cyan!22] (0, -1.3) rectangle (6.4, 1.3);
  \draw[tierra, rounded corners=6pt] (0.9, -0.55) rectangle (2.9, 0.55);
  \draw[tierra, rounded corners=6pt] (3.6, -0.55) -- (3.6, 0.55)
    -- (6.4, 0.75) -- (6.4, -0.75) -- cycle;
  \foreach \x in {1.4, 2.4} {
    \draw[puente] (\x, 0.55) -- (\x, 1.3);
    \draw[puente] (\x, -0.55) -- (\x, -1.3);
  }
  \draw[puente] (2.9, 0) -- (3.6, 0);
  \draw[puente] (5.0, 0.65) -- (5.0, 1.3);
  \draw[puente] (5.0, -0.65) -- (5.0, -1.3);
  \node[font=\small\bfseries] at (3.2, 1.65) {N};
  \node[font=\small\bfseries] at (3.2, -1.65) {S};
  \node[font=\small\bfseries] at (1.9, 0) {I};
  \node[font=\small\bfseries] at (5.6, 0) {E};
  \node[nota] at (0.5, -0.92) {Pregel};
  \begin{scope}[xshift=8.6cm]
    \node[vert] (I) at (0, 0) {I};
    \node[vert] (N) at (2, 1.5) {N};
    \node[vert] (S) at (2, -1.5) {S};
    \node[vert] (E) at (4, 0) {E};
    \draw[arista] (I) to[bend left=22] (N);
    \draw[arista] (I) to[bend right=22] (N);
    \draw[arista] (I) to[bend left=22] (S);
    \draw[arista] (I) to[bend right=22] (S);
    \draw[arista] (I) -- (E);
    \draw[arista] (N) -- (E);
    \draw[arista] (S) -- (E);
    \node[nota, left=3pt] at (I.west) {5};
    \node[nota, above=3pt] at (N.north) {3};
    \node[nota, below=3pt] at (S.south) {3};
    \node[nota, right=3pt] at (E.east) {3};
  \end{scope}
\end{tikzpicture}
\end{center}

El grafo tiene $n = 4$ y $m = 7$, y no es simple: entre $I$ y $N$ hay
dos puentes, y otros dos entre $I$ y $S$. Sus aristas forman el
multiconjunto
\[
  E(K) = \set{\set{I,N}, \set{I,N}, \set{I,S}, \set{I,S}, \set{I,E},
              \set{N,E}, \set{S,E}} .
\]
La $E$ de la región este es la de las diapositivas; para no
confundirla con el conjunto de aristas, a este se le llama $E(K)$. La
pregunta de Euler se vuelve: ¿hay un recorrido cerrado que use cada
arista exactamente una vez? La respuesta, al final de la sección, sale
de los números junto a cada vértice.

% =====================================================================

\subsection*{Multigrafos, pseudografos y grado}

La sección~\ref{sec:grafos} define el grafo simple, donde $E$ es un
conjunto de pares no ordenados de vértices distintos, y nota que un
lazo o dos aristas entre el mismo par no caben en esa definición. Para
admitirlos, $E$ se vuelve un multiconjunto. Rosen les pone nombre
según tres preguntas: ¿las aristas tienen dirección?, ¿puede haber
varias entre el mismo par de vértices?, ¿puede haber lazos?

\begin{center}
\small
\begin{tabular}{lccc}
  \toprule
  \textbf{Tipo} & \textbf{Dirigido} & \textbf{Aristas múltiples}
    & \textbf{Lazos} \\
  \midrule
  Grafo simple & no & no & no \\
  Multigrafo & no & sí & no \\
  Pseudografo & no & sí & sí \\
  Grafo dirigido simple & sí & no & no \\
  Grafo dirigido & sí & no & sí \\
  Multigrafo dirigido & sí & sí & sí \\
  \bottomrule
\end{tabular}
\end{center}

\begin{definicion}[Multigrafo y pseudografo]
  \label{def:pseudografo}
  Un \textbf{multigrafo} es un par $G = (V, E)$ en el que $E$ es un
  multiconjunto de pares $\set{u, v}$ con $u \neq v$: dos aristas
  distintas pueden tener los mismos extremos, y se llaman
  \textbf{múltiples} o \textbf{paralelas}. Un \textbf{pseudografo}
  admite además \textbf{lazos}, aristas $\set{v, v}$ cuyos dos extremos
  son el mismo vértice.
\end{definicion}

Como dos aristas paralelas no se distinguen por sus extremos, en un
multigrafo se les pone nombre, $e_1, \dots, e_m$, y cada una tiene sus
dos extremos. Todo grafo simple es multigrafo y todo multigrafo es
pseudografo, así que un teorema probado para pseudografos vale para
los tres. Rosen agrega el \textbf{grafo mixto}, con aristas dirigidas y
no dirigidas a la vez; no aparece en esta parte.

\begin{ojo}
  En un grafo dirigido, los arcos $(c, d)$ y $(d, c)$ \textbf{no} son
  múltiples: tienen sentidos opuestos y son arcos distintos, como las
  dos flechas entre $1$ y $2$ del grafo dirigido de la
  sección~\ref{sec:grafos}. Múltiples son dos arcos con el mismo origen
  y el mismo destino. La tabla de las diapositivas dice que un grafo
  dirigido ``puede tener'' aristas múltiples; eso contradice a Rosen y
  al propio ejercicio 7 de abajo, cuya respuesta es grafo dirigido y no
  multigrafo dirigido.
\end{ojo}

\begin{ejercicio}[Rosen, \S 10.1, ejercicios 3, 5 y 7]
  Decir qué tipo de grafo es cada uno de los tres dibujos.
  \begin{center}
  \begin{tikzpicture}[x=1cm, y=1cm]
    \node[vert] (a) at (0, 2) {a};  \node[vert] (b) at (2, 2) {b};
    \node[vert] (c) at (0, 0) {c};  \node[vert] (d) at (2, 0) {d};
    \foreach \u/\v in {a/b, a/c, b/c, b/d} \draw[arista] (\u) -- (\v);
    \node[nota] at (1, -1.15) {ejercicio 3};
    \begin{scope}[xshift=4.2cm]
      \node[vert] (a) at (0, 2) {a};  \node[vert] (b) at (2, 2) {b};
      \node[vert] (c) at (0, 0) {c};  \node[vert] (d) at (2, 0) {d};
      \draw[arista] (a) to[bend left=14] (b);
      \draw[arista] (a) to[bend right=14] (b);
      \draw[arista] (b) to[bend left=14] (d);
      \draw[arista] (b) to[bend right=14] (d);
      \draw[arista] (c) to[bend left=14] (d);
      \draw[arista] (c) to[bend right=14] (d);
      \draw[arista] (a) -- (c);
      \draw[arista] (a) to[out=150, in=210, looseness=7] (a);
      \draw[arista] (b) to[out=30, in=330, looseness=7] (b);
      \draw[arista] (d) to[out=30, in=330, looseness=7] (d);
      \node[nota] at (1, -1.15) {ejercicio 5};
    \end{scope}
    \begin{scope}[xshift=8.6cm]
      \node[vert] (a) at (0, 0) {a};    \node[vert] (b) at (1.5, 2) {b};
      \node[vert] (e) at (1.5, 0) {e};  \node[vert] (c) at (3.2, 2) {c};
      \node[vert] (d) at (3.2, 0) {d};
      \draw[dirigida] (a) -- (b);
      \draw[dirigida] (b) -- (e);
      \draw[dirigida] (e) -- (a);
      \draw[dirigida] (e) -- (d);
      \draw[dirigida] (c) -- (b);
      \draw[dirigida] (c) to[bend left=25] (d);
      \draw[dirigida] (d) to[bend left=25] (c);
      \draw[dirigida] (e) to[out=300, in=240, looseness=7] (e);
      \draw[dirigida] (c) to[out=60, in=0, looseness=7] (c);
      \node[nota] at (1.6, -1.15) {ejercicio 7};
    \end{scope}
  \end{tikzpicture}
  \end{center}
  \textit{Solución.} Se contestan las tres preguntas en orden. El 3 no
  es dirigido, no tiene múltiples ni lazos: grafo simple. El 5 tiene
  aristas dobles entre $a$ y $b$, $b$ y $d$, $c$ y $d$, y lazos en
  $a$, $b$ y $d$: pseudografo. El 7 es dirigido y tiene lazos en $c$ y
  en $e$; los dos arcos entre $c$ y $d$ son antiparalelos, no
  múltiples: grafo dirigido.
\end{ejercicio}

Modelar es decidir qué se guarda. El ejemplo de las rutas aéreas de
las diapositivas lo muestra con los mismos datos y dos grafos
distintos.

\begin{ejemplo}[Rutas aéreas, Rosen \S 10.1, ejercicio 1]
  Una aerolínea tiene estos vuelos diarios: Boston a Newark 4, Newark a
  Boston 2, Newark a Miami 3, Miami a Newark 2, Newark a Detroit 1,
  Detroit a Newark 2, Newark a Washington 3, Washington a Newark 2,
  Washington a Miami 1. Si solo interesa qué ciudades tienen vuelo
  entre sí, en cualquier sentido, el modelo es un grafo simple no
  dirigido con $5$ aristas (izquierda). Si importa el sentido, es un
  grafo dirigido simple con $9$ arcos (derecha).
  \begin{center}
  \begin{tikzpicture}[x=1cm, y=1cm,
    lugar/.style={draw=acento, fill=fondoIdea, rounded corners=2pt,
                  font=\scriptsize, inner sep=2.5pt}]
    \node[lugar] (Bo) at (0, 0) {Boston};
    \node[lugar] (Ne) at (2.1, 0) {Newark};
    \node[lugar] (De) at (2.1, 1.3) {Detroit};
    \node[lugar] (Wa) at (4.3, 0.8) {Washington};
    \node[lugar] (Mi) at (4.3, -0.8) {Miami};
    \foreach \u/\v in {Bo/Ne, Ne/De, Ne/Wa, Ne/Mi, Wa/Mi}
      \draw[arista] (\u) -- (\v);
    \node[nota] at (2.1, -1.4) {simple no dirigido, $m = 5$};
    \begin{scope}[xshift=7.6cm]
      \node[lugar] (Bo) at (0, 0) {Boston};
      \node[lugar] (Ne) at (2.1, 0) {Newark};
      \node[lugar] (De) at (2.1, 1.3) {Detroit};
      \node[lugar] (Wa) at (4.3, 0.8) {Washington};
      \node[lugar] (Mi) at (4.3, -0.8) {Miami};
      \foreach \u/\v in {Bo/Ne, Ne/De, Ne/Wa, Ne/Mi}{
        \draw[dirigida] (\u) to[bend left=14] (\v);
        \draw[dirigida] (\v) to[bend left=14] (\u);
      }
      \draw[dirigida] (Wa) -- (Mi);
      \node[nota] at (2.1, -1.4) {dirigido simple, $9$ arcos};
    \end{scope}
  \end{tikzpicture}
  \end{center}
  El de la derecha responde que hay vuelo de Washington a Miami pero no
  de regreso; el de la izquierda no puede. Ninguno de los dos sabe
  cuántos vuelos hay: para eso haría falta un multigrafo con una arista
  por vuelo, $20$ en total, o un grafo pesado con el número de vuelos
  como peso.
\end{ejemplo}

\subsubsection*{El grado y el teorema del apretón de manos}

En un pseudografo el grado necesita una precisión: un lazo toca a su
vértice por los dos extremos.

\begin{definicion}[Grado]
  \label{def:grado-lazos}
  El \textbf{grado} $\text{gr}(v)$ de un vértice de un pseudografo es
  el número de extremos de aristas que caen en $v$: cada arista
  $\set{u, v}$ con $u \neq v$ suma $1$ al grado de cada extremo, y cada
  lazo en $v$ suma $2$ al grado de $v$. Un vértice de grado $0$ es
  \textbf{aislado} y uno de grado $1$ es una \textbf{hoja}.
\end{definicion}

Rosen, Cormen y las diapositivas escriben $\deg(v)$; aquí se conserva
el $\text{gr}(v)$ de la sección~\ref{sec:grafos}. Ahí mismo aparece,
como observación, la suma de los grados en un grafo simple. Vale en
general y tiene nombre.

\begin{teorema}[Apretón de manos]
  \label{teo:apreton-manos}
  En todo pseudografo $G = (V, E)$,
  \[
    \sum_{v \in V} \text{gr}(v) = 2m .
  \]
\end{teorema}

\begin{proof}
  Se cuentan de dos maneras los pares (vértice, extremo de arista que
  cae en él). Contando por vértices, cada $v$ aporta $\text{gr}(v)$
  pares, y el total es la suma de los grados. Contando por aristas,
  cada arista tiene exactamente dos extremos, sea o no un lazo, y
  aporta $2$. Las dos cuentas cuentan lo mismo, así que coinciden.
\end{proof}

El nombre viene de una fiesta: si cada quien anota cuántas manos
estrechó, la suma de las anotaciones es el doble del número de
apretones. En el pseudografo del ejercicio 5,
$\text{gr}(a) = 2 + 2 + 1 = 5$ (el lazo, la arista doble a $b$ y la
arista a $c$), $\text{gr}(b) = 6$, $\text{gr}(c) = 3$ y
$\text{gr}(d) = 6$; suman $20 = 2 \cdot 10$, y en efecto hay tres
lazos, tres pares de aristas dobles y una sencilla.

\begin{ojo}
  Las diapositivas definen el grado como ``el número de aristas
  incidentes''. Leída al pie de la letra, esa definición cuenta un lazo
  una sola vez, y entonces el teorema falla: en el ejercicio 5 la suma
  daría $17$ y no $20$. La versión de Rosen, con el lazo contado dos
  veces, es la que hace funcionar todo lo que sigue.
\end{ojo}

\begin{corolario}
  \label{cor:impares-par}
  En todo pseudografo, el número de vértices de grado impar es par.
\end{corolario}

\begin{proof}
  Se parte la suma de los grados en la de los vértices de grado par y
  la de los de grado impar. El total es $2m$, par, y la primera parte
  es par por ser suma de pares; luego la suma de los grados impares
  también es par. Una suma de números impares es par solo si tiene un
  número par de sumandos.
\end{proof}

El corolario parece una curiosidad. Es lo que hace exhaustiva la
clasificación de los grafos eulerianos al final de la sección.

\begin{ejercicio}
  \label{ej:siete-personas}
  ¿Puede haber una reunión de siete personas en la que cada una
  estreche la mano de exactamente otras tres? ¿Hay un grafo simple con
  grados $3, 3, 3, 1$?

  \textit{Solución.} Lo primero no: serían siete vértices de grado
  impar, y por el corolario~\ref{cor:impares-par} eso es imposible (la
  suma sería $21$). Lo segundo tampoco, aunque la suma, $10$, sea par.
  Los tres vértices de grado $3$ tienen solo tres vecinos posibles
  cada uno, así que cada uno es adyacente a todos los demás, incluido
  el de grado $1$, que terminaría con grado $3$. La paridad es
  necesaria, no suficiente. En un multigrafo la secuencia sí existe:
  basta duplicar la arista entre dos de los vértices de grado $3$.
\end{ejercicio}

\subsubsection*{Grado de entrada y de salida}

En un grafo dirigido, $\text{gr}^{-}(v)$ es el \textbf{grado de
entrada}, el número de arcos que llegan a $v$, y $\text{gr}^{+}(v)$ el
\textbf{grado de salida}, el número de arcos que salen de $v$. Un lazo
$(v, v)$ suma $1$ a cada uno.

\begin{proposicion}
  \label{prop:grados-dirigidos}
  En todo grafo dirigido,
  $\sum_{v} \text{gr}^{-}(v) = \sum_{v} \text{gr}^{+}(v) = m$.
\end{proposicion}

\begin{proof}
  Cada arco tiene exactamente un origen y un destino: aporta $1$ a la
  suma de grados de salida y $1$ a la de entrada.
\end{proof}

En el ejercicio 7 los grados de entrada y salida son $a{:}\ 1, 1$;
$b{:}\ 2, 1$; $c{:}\ 2, 3$; $d{:}\ 2, 1$; $e{:}\ 2, 3$. Las dos columnas
suman $9$, el número de arcos, contando cada lazo una vez en cada lado.

\subsubsection*{Familias con nombre}

Algunos grafos aparecen tanto que tienen nombre propio. Las
diapositivas usan cuatro, y las secciones siguientes también.

\begin{ejemplo}[Familias de grafos]
  \label{ej:familias-teoria}
  El \textbf{camino} $P_n$ tiene $n$ vértices en fila y $n - 1$
  aristas; el \textbf{ciclo} $C_n$, con $n \geq 3$, cierra la fila con
  una arista más; el \textbf{completo} $K_n$ une cada par de vértices,
  y la \textbf{estrella} une un centro con $k$ hojas.
  \begin{center}
  \begin{tikzpicture}[x=1cm, y=1cm]
    \foreach \k/\tit in {0/{$P_4$}, 1/{$C_4$}, 2/{$K_4$}} {
      \begin{scope}[xshift=\k*3.6cm]
        \node[vert] (a\k) at (0, 1.5) {a};
        \node[vert] (b\k) at (1.5, 1.5) {b};
        \node[vert] (d\k) at (0, 0) {d};
        \node[vert] (c\k) at (1.5, 0) {c};
        \node[nota] at (0.75, -0.7) {\tit};
      \end{scope}
    }
    \foreach \u/\v in {a/b, b/c, c/d} \draw[arista] (\u0) -- (\v0);
    \foreach \u/\v in {a/b, b/c, c/d, d/a} \draw[arista] (\u1) -- (\v1);
    \foreach \u/\v in {a/b, b/c, c/d, d/a, a/c, b/d}
      \draw[arista] (\u2) -- (\v2);
    \begin{scope}[xshift=11.55cm, yshift=0.6cm]
      \node[vert] (o) at (0, 0) {};
      \foreach \ang/\nom in {90/h1, 210/h2, 330/h3}
        \node[vert] (\nom) at (\ang:1.0) {};
      \foreach \nom in {h1, h2, h3} \draw[arista] (o) -- (\nom);
      \node[nota] at (0, -1.3) {estrella, $k = 3$};
    \end{scope}
  \end{tikzpicture}
  \end{center}
  En $P_n$ los extremos tienen grado $1$ y los demás $2$; en $C_n$
  todos tienen grado $2$ y $m = n$; en $K_n$ todos tienen grado $n - 1$
  y $m = n(n-1)/2$, el máximo posible en un grafo simple; en la
  estrella el centro tiene grado $k$ y $m = k$. La
  sección~\ref{sec:hamilton-arboles} vuelve a la estrella como el
  bipartito completo $K_{1,k}$.
\end{ejemplo}

\subsubsection*{Subgrafos}

\begin{definicion}[Subgrafo]
  \label{def:subgrafo}
  $H = (W, F)$ es un \textbf{subgrafo} de $G = (V, E)$ si
  $W \subseteq V$, $F \subseteq E$ y cada arista de $F$ tiene sus dos
  extremos en $W$. Es el \textbf{subgrafo inducido} por $W$, y se
  escribe $G[W]$, si $F$ contiene \emph{todas} las aristas de $E$ con
  los dos extremos en $W$. Es un \textbf{subgrafo generador} (o de
  expansión) si $W = V$.
\end{definicion}

Dos subgrafos se usan a cada rato. $G - e$ es el grafo sin la arista
$e$, con todos los vértices. $G - v$ es el grafo sin el vértice $v$
\emph{y sin todas las aristas que inciden en $v$}, porque una arista
sin uno de sus extremos deja de ser arista.

\begin{ojo}
  La diapositiva 46 define subgrafo con $W \subseteq V$ y
  $F \subseteq E$, sin la tercera condición. Sin ella, $W = \set{A}$ y
  $F = \set{\set{A, B}}$ cumplirían la definición, y eso no es un
  grafo: la arista tiene un extremo que no está.
\end{ojo}

\begin{ejemplo}[Subgrafos del grafo $G$ de la libreta]
  En el grafo $G$ de la sección~\ref{sec:grafos}, el subgrafo inducido
  por $\set{A, B, C, E, F}$ tiene las aristas $\set{A,B}$,
  $\set{A,C}$, $\set{B,E}$, $\set{C,F}$ y $\set{E,F}$: es el ciclo
  $C_5$ que la clase encuentra con BFS. El subgrafo $G - \set{E, F}$
  es generador, porque conserva los siete vértices, y es exactamente el
  árbol BFS de esa clase. En general, los árboles de BFS, de DFS y de
  expansión mínima son subgrafos generadores.
\end{ejemplo}

% =====================================================================

\subsection*{Matriz de incidencia}

La matriz de adyacencia de la sección~\ref{sec:grafos} relaciona
vértices con vértices. La de incidencia relaciona vértices con
aristas, y por eso admite aristas paralelas sin perder información.

\begin{definicion}[Matriz de incidencia]
  \label{def:matriz-incidencia}
  Con los vértices $v_1, \dots, v_n$ y las aristas $e_1, \dots, e_m$
  numerados, la matriz de incidencia es la matriz $M$ de $n \times m$
  con $M_{ij} = 1$ si $v_i$ es extremo de $e_j$, y $M_{ij} = 0$ si no.
\end{definicion}

\begin{ejemplo}[Diapositivas 41 y 44]
  \label{ej:incidencia-teoria}
  A la izquierda, un grafo simple con cinco vértices y seis aristas; a
  la derecha, un pseudografo con lazos $e_1$ y $e_8$ y aristas
  paralelas $e_2$ y $e_3$.
  \begin{center}
  \begin{tikzpicture}[x=1cm, y=1cm]
    \node[vert] (v1) at (0, 1.6) {$v_1$};
    \node[vert] (v2) at (1.6, 1.6) {$v_2$};
    \node[vert] (v3) at (3.2, 1.6) {$v_3$};
    \node[vert] (v4) at (0.8, 0) {$v_4$};
    \node[vert] (v5) at (2.4, 0) {$v_5$};
    \draw[arista] (v1) -- node[peso, pos=0.5] {$e_1$} (v4);
    \draw[arista] (v1) -- node[peso, pos=0.22] {$e_2$} (v5);
    \draw[arista] (v2) -- node[peso, pos=0.78] {$e_3$} (v4);
    \draw[arista] (v2) -- node[peso, pos=0.5] {$e_4$} (v5);
    \draw[arista] (v3) -- node[peso, pos=0.5] {$e_5$} (v5);
    \draw[arista] (v2) -- node[peso, pos=0.5] {$e_6$} (v3);
    \begin{scope}[xshift=6.2cm]
      \node[vert] (v1) at (0, 1.6) {$v_1$};
      \node[vert] (v2) at (1.8, 1.6) {$v_2$};
      \node[vert] (v3) at (3.8, 1.6) {$v_3$};
      \node[vert] (v4) at (0.4, 0) {$v_4$};
      \node[vert] (v5) at (2.8, 0) {$v_5$};
      \draw[arista] (v1) to[out=200, in=160, looseness=7]
        node[peso, left] {$e_1$} (v1);
      \draw[arista] (v1) to[bend left=25] node[peso] {$e_2$} (v2);
      \draw[arista] (v1) to[bend right=25] node[peso] {$e_3$} (v2);
      \draw[arista] (v2) -- node[peso] {$e_4$} (v3);
      \draw[arista] (v3) -- node[peso] {$e_5$} (v5);
      \draw[arista] (v2) -- node[peso] {$e_6$} (v5);
      \draw[arista] (v2) -- node[peso] {$e_7$} (v4);
      \draw[arista] (v4) to[out=290, in=250, looseness=7]
        node[peso, below] {$e_8$} (v4);
    \end{scope}
  \end{tikzpicture}
  \end{center}
  Para el grafo simple, su matriz de incidencia $M$ y el producto
  $M M^{T}$:
  \[
    M = \begin{array}{c|cccccc}
        & e_1 & e_2 & e_3 & e_4 & e_5 & e_6 \\ \hline
      v_1 & 1 & 1 & 0 & 0 & 0 & 0 \\
      v_2 & 0 & 0 & 1 & 1 & 0 & 1 \\
      v_3 & 0 & 0 & 0 & 0 & 1 & 1 \\
      v_4 & 1 & 0 & 1 & 0 & 0 & 0 \\
      v_5 & 0 & 1 & 0 & 1 & 1 & 0
    \end{array}
    \qquad
    M M^{T} = \begin{array}{c|ccccc}
        & v_1 & v_2 & v_3 & v_4 & v_5 \\ \hline
      v_1 & 2 & 0 & 0 & 1 & 1 \\
      v_2 & 0 & 3 & 1 & 1 & 1 \\
      v_3 & 0 & 1 & 2 & 0 & 1 \\
      v_4 & 1 & 1 & 0 & 2 & 0 \\
      v_5 & 1 & 1 & 1 & 0 & 3
    \end{array}
  \]
  Para el pseudografo:
  \[
    M = \begin{array}{c|cccccccc}
        & e_1 & e_2 & e_3 & e_4 & e_5 & e_6 & e_7 & e_8 \\ \hline
      v_1 & 1 & 1 & 1 & 0 & 0 & 0 & 0 & 0 \\
      v_2 & 0 & 1 & 1 & 1 & 0 & 1 & 1 & 0 \\
      v_3 & 0 & 0 & 0 & 1 & 1 & 0 & 0 & 0 \\
      v_4 & 0 & 0 & 0 & 0 & 0 & 0 & 1 & 1 \\
      v_5 & 0 & 0 & 0 & 0 & 1 & 1 & 0 & 0
    \end{array}
  \]
  Las columnas de los lazos tienen un solo $1$, y las de las aristas
  paralelas $e_2$ y $e_3$ son iguales.
\end{ejemplo}

\begin{proposicion}
  \label{prop:incidencia}
  Sea $M$ la matriz de incidencia de un multigrafo sin lazos y $A$ su
  matriz de adyacencia, con $A_{ij}$ igual al número de aristas entre
  $v_i$ y $v_j$. Entonces cada columna de $M$ tiene exactamente dos
  unos, la suma del renglón $i$ es $\text{gr}(v_i)$, y
  \[
    M M^{T} = D + A ,
  \]
  donde $D$ es la matriz diagonal con los grados.
\end{proposicion}

\begin{proof}
  Las dos primeras afirmaciones son la definición leída por columnas y
  por renglones. Para la tercera,
  $(M M^{T})_{ij} = \sum_{k} M_{ik} M_{jk}$, y el sumando $k$ vale $1$
  justo cuando $e_k$ incide a la vez en $v_i$ y en $v_j$. Con
  $i = j$ eso cuenta las aristas que inciden en $v_i$, su grado; con
  $i \neq j$ cuenta las aristas entre $v_i$ y $v_j$, que es $A_{ij}$.
\end{proof}

En el ejemplo, la diagonal de $M M^{T}$ es $2, 3, 2, 2, 3$, los grados,
y fuera de ella está la matriz de adyacencia. Con lazos la suma del
renglón deja de ser el grado: en el pseudografo el renglón de $v_1$
suma $3$ y $\text{gr}(v_1) = 4$, y el de $v_4$ suma $2$ con
$\text{gr}(v_4) = 3$. Algunos textos escriben un $2$ en la entrada del
lazo para que la suma sí dé el grado.

\begin{idea}
  La matriz de incidencia ocupa $\Tht{nm}$ de memoria, más que la de
  adyacencia en cuanto $m > n$, y preguntar si $\set{u, v}$ es arista
  cuesta $\Tht{m}$: hay que recorrer el renglón de $u$ buscando una
  columna que también tenga $1$ en $v$. Por eso casi ningún algoritmo
  la usa. Su valor es teórico: identidades como $M M^{T} = D + A$, y
  que se generaliza sin cambios a aristas con más de dos extremos. La
  versión dirigida, con $-1$ en el origen y $1$ en el destino, está en
  la sección~\ref{sec:busquedas} (Cormen, ejercicio 22.1-7).
\end{idea}

% =====================================================================

\subsection*{Isomorfismo}

La sección~\ref{sec:grafos} distingue el grafo encajado del
topológico: dos dibujos distintos son el mismo grafo topológico si unen
a los mismos pares. Cuando además los vértices tienen nombres
distintos, ``el mismo grafo'' se precisa con una biyección.

\begin{definicion}[Isomorfismo]
  \label{def:isomorfismo}
  Dos grafos simples $G_1 = (V_1, E_1)$ y $G_2 = (V_2, E_2)$ son
  \textbf{isomorfos}, $G_1 \cong G_2$, si existe una biyección
  $f : V_1 \to V_2$ tal que
  \[
    \set{u, v} \in E_1 \iff \set{f(u), f(v)} \in E_2 .
  \]
  Una propiedad que conserva todo isomorfismo es un
  \textbf{invariante}.
\end{definicion}

El número de vértices, el de aristas y la secuencia de grados son
invariantes, y también lo son la conexidad, el número de ciclos de
cada longitud y el grado de los vecinos de cada vértice: la biyección
manda vecinos en vecinos. Esto da la regla para trabajar. Si dos grafos
difieren en un invariante, no son isomorfos. Si coinciden en todos los
que se revisaron, no se sabe nada: para afirmar que lo son hay que
exhibir la biyección.

\begin{ejemplo}[De la matriz al grafo, diapositivas 36 a 38]
  La diapositiva 36 da una matriz de adyacencia en el orden
  $a, b, c, d$ que, leída renglón por renglón, tiene las aristas
  $\set{a,b}$, $\set{a,c}$, $\set{b,d}$, $\set{c,d}$. La diapositiva
  38 la dibuja con dos aristas cruzadas (derecha), y parece otra cosa
  que el cuadrado $C_4$ (izquierda).
  \begin{center}
  \begin{tikzpicture}[x=1cm, y=1cm]
    \node[vert] (A) at (0, 1.6) {A};  \node[vert] (B) at (1.6, 1.6) {B};
    \node[vert] (C) at (0, 0) {C};    \node[vert] (D) at (1.6, 0) {D};
    \foreach \u/\v in {A/B, A/C, B/D, C/D} \draw[arista] (\u) -- (\v);
    \node[nota] at (0.8, -0.7) {$C_4$};
    \begin{scope}[xshift=4.5cm]
      \node[vert] (a) at (0, 1.6) {a};  \node[vert] (b) at (1.6, 1.6) {b};
      \node[vert] (d) at (0, 0) {d};    \node[vert] (c) at (1.6, 0) {c};
      \foreach \u/\v in {a/b, a/c, b/d, d/c} \draw[arista] (\u) -- (\v);
      \node[nota] at (0.8, -0.7) {diapositiva 38};
    \end{scope}
  \end{tikzpicture}
  \end{center}
  La biyección $f(A) = a$, $f(B) = b$, $f(C) = c$, $f(D) = d$ manda
  las cuatro aristas de la izquierda en las cuatro de la derecha: son
  isomorfos. El cruce es del dibujo, no del grafo.
\end{ejemplo}

\begin{ejemplo}[Dos grafos que se parecen, diapositivas 48 a 51]
  \label{ej:isomorfismo-teoria}
  $G$ y $H$ tienen ocho vértices cada uno, un cuadrado exterior y uno
  interior unidos por dos aristas:
  \[
    E_G = \set{ab, bc, cd, da, ef, fg, gh, he, bf, dh} , \qquad
    E_H = \set{st, tu, uv, vs, wx, xy, yz, zw, sw, vz} .
  \]
  En verde, los vértices de grado $2$; punteadas, las aristas de $H$
  que unen dos vértices de grado $2$.
  \begin{center}
  \begin{tikzpicture}[x=1cm, y=1cm]
    \foreach \nom/\x/\y in {a/0/2.6, b/3.2/2.6, c/3.2/0, d/0/0,
                            e/1/1.8, f/2.2/1.8, g/2.2/0.8, h/1/0.8}
      \node[vert] (\nom) at (\x, \y) {\nom};
    \foreach \nom in {a, c, e, g} \node[vertOk] at (\nom) {\nom};
    \foreach \u/\v in {a/b, b/c, c/d, d/a, e/f, f/g, g/h, h/e, b/f, d/h}
      \draw[arista] (\u) -- (\v);
    \node[font=\small] at (1.6, -0.7) {$G$};
    \begin{scope}[xshift=5.6cm]
      \foreach \nom/\x/\y in {s/0/2.6, t/3.2/2.6, u/3.2/0, v/0/0,
                              w/1/1.8, x/2.2/1.8, y/2.2/0.8, z/1/0.8}
        \node[vert] (\nom) at (\x, \y) {\nom};
      \foreach \nom in {t, u, x, y} \node[vertOk] at (\nom) {\nom};
      \foreach \u/\v in {s/t, u/v, v/s, w/x, y/z, z/w, s/w, v/z}
        \draw[arista] (\u) -- (\v);
      \draw[aristaFuera] (t) -- (u);
      \draw[aristaFuera] (x) -- (y);
      \node[font=\small] at (1.6, -0.7) {$H$};
    \end{scope}
  \end{tikzpicture}
  \end{center}
  Los invariantes de la diapositiva 49 coinciden: ocho vértices, diez
  aristas y la misma secuencia de grados, $2, 2, 2, 2, 3, 3, 3, 3$. Los
  dos últimos renglones de la tabla no coinciden.
  \begin{center}
  \small
  \begin{tabular}{lcc}
    \toprule
    \textbf{Invariante} & $G$ & $H$ \\
    \midrule
    Vértices, aristas & $8$, $10$ & $8$, $10$ \\
    Secuencia de grados & $2,2,2,2,3,3,3,3$ & $2,2,2,2,3,3,3,3$ \\
    Vértices de grado $2$ & $a, c, e, g$ & $t, u, x, y$ \\
    Aristas entre vértices de grado $2$ & $0$ & $2$ ($tu$, $xy$) \\
    Ciclos de longitud $4$ & $2$ & $3$ ($stuv$, $wxyz$, $swzv$) \\
    \bottomrule
  \end{tabular}
  \end{center}
  El argumento de la clase va por vértices. Un isomorfismo conserva el
  grado y la adyacencia, así que $a$, de grado $2$, tendría que ir a un
  vértice de grado $2$ de $H$ cuyos dos vecinos tengan grado $3$, como
  los de $a$. En $H$ no hay ninguno: cada vértice de grado $2$ tiene un
  vecino de grado $2$. No existe a dónde mandar $a$, así que $G$ y $H$
  \textbf{no son isomorfos}.
\end{ejemplo}

\begin{ojo}
  La diapositiva 51 termina en la observación sobre los vecinos y no
  escribe la conclusión; la de arriba se confirmó probando las
  $8! = 40\,320$ biyecciones. De paso, la diapositiva 52 que sigue es
  un ejemplo de grado fuera de lugar, no la respuesta. Mismo dibujo no
  implica mismo grafo, ni al revés: la diapositiva 38 es $C_4$ con un
  cruce, y $G$ y $H$ se parecen mucho sin ser isomorfos.
\end{ojo}

Probar las $n!$ biyecciones cuesta $\Omg{n!}$. No se conoce un
algoritmo polinomial para decidir isomorfismo en general, ni se sabe
que el problema sea NP-completo; el mejor resultado, de Babai (2015),
es cuasipolinomial. En la práctica se descarta con invariantes baratos
y solo se busca la biyección cuando todos coinciden.

% =====================================================================

\subsection*{Caminatas, senderos, caminos, circuitos y ciclos}

Moverse dentro de un grafo tiene cinco nombres, cada uno con una
restricción más que el anterior. La terminología cambia de un libro a
otro, así que esta subsección fija la que usa toda la parte de grafos
de la libreta.

\begin{definicion}[Caminata]
  \label{def:caminata}
  En un pseudografo, una \textbf{caminata} de longitud $k$ de $u$ a $v$
  es una secuencia alternada de vértices y aristas
  \[
    v_0, e_1, v_1, e_2, \dots, e_k, v_k
  \]
  con $v_0 = u$, $v_k = v$ y cada $e_i$ una arista de extremos
  $v_{i-1}$ y $v_i$. Es \textbf{cerrada} si $v_0 = v_k$. Un solo
  vértice es una caminata de longitud $0$.
\end{definicion}

En un grafo simple cada arista queda determinada por sus extremos, y
la caminata se escribe solo con vértices, $v_0, v_1, \dots, v_k$. Las
diapositivas usan $n$ para la longitud; aquí es $k$, porque $n$ es el
número de vértices.

\begin{definicion}[Sendero, camino, circuito y ciclo]
  \label{def:sendero-ciclo}
  Una caminata es
  \begin{itemize}[itemsep=0.1em]
    \item un \textbf{sendero} si no repite aristas;
    \item un \textbf{camino} si no repite vértices;
    \item un \textbf{circuito} si es un sendero cerrado de longitud al
          menos $1$;
    \item un \textbf{ciclo} si es un circuito en el que
          $v_0, v_1, \dots, v_{k-1}$ son distintos: un camino cerrado,
          que no repite vértices salvo el primero y el último.
  \end{itemize}
\end{definicion}

\begin{center}
\small
\begin{tabular}{lccc}
  \toprule
  & \textbf{¿Repite vértices?} & \textbf{¿Repite aristas?}
    & \textbf{¿Cerrada?} \\
  \midrule
  Caminata & puede & puede & puede \\
  Sendero & puede & no & puede \\
  Camino & no & no & solo con $k = 0$ \\
  Circuito & puede & no & sí \\
  Ciclo & solo $v_0 = v_k$ & no & sí \\
  \bottomrule
\end{tabular}
\end{center}

Todo camino es sendero: si no se repiten vértices, una arista no puede
repetirse, porque sus extremos se repetirían. Por lo mismo, todo ciclo
es circuito. El recíproco falla: en dos triángulos pegados por un
vértice (la ``mariposa'' de la diapositiva 119, más abajo) el
recorrido $c, a, b, c, d, e, c$ es un circuito que pasa dos
veces por $c$, y no es ciclo.

\begin{observacion}
  En un grafo simple todo ciclo tiene longitud al menos $3$: uno de
  longitud $1$ sería un lazo, y uno de longitud $2$, como $a, b, a$,
  usaría dos veces la misma arista. En un multigrafo dos aristas
  paralelas $a, e_1, b, e_2, a$ sí forman un ciclo de longitud $2$, y
  en un pseudografo un lazo es un ciclo de longitud $1$.
\end{observacion}

La definición de ciclo de la sección~\ref{sec:grafos} y el ``ciclo''
que Kruskal evita cerrar son ciclos en este sentido estricto.

\begin{ojo}
  La profesora llama \textbf{circuito} a cualquier caminata cerrada
  (diapositiva 81), aunque repita aristas, y lo advierte: ``depende de
  la convención''. Grimaldi, que ella misma recomienda, exige que el
  circuito no repita aristas, y esta libreta sigue a Grimaldi. Hay dos
  razones. Así los nombres van en pares, abierto y cerrado:
  sendero y circuito, camino y ciclo. Y así un circuito euleriano es
  literalmente un circuito. El precio es que una respuesta de la
  actividad de la clase cambia (la secuencia 3 del
  ejercicio~\ref{ej:actividad-recorridos}). Con los libros en inglés
  hay que traducir con cuidado: el \textit{simple path} de Rosen es un
  sendero y el de Cormen es un camino, y los ``caminos simples'' de la
  sección~\ref{sec:grafos} son caminos a secas.\par
  \vspace{6pt}
  {\centering\small
  \begin{tabular}{lllll}
    \toprule
    \textbf{Libreta} & \textbf{Diapositivas} & \textbf{Grimaldi}
      & \textbf{Rosen} & \textbf{Cormen} \\
    \midrule
    caminata & caminata & \textit{walk} & \textit{path}
      & \textit{path} \\
    sendero & sendero & \textit{trail} & \textit{simple path}
      & sin nombre \\
    camino & camino, trayectoria & \textit{path} & sin nombre
      & \textit{simple path} \\
    caminata cerrada & circuito & \textit{closed walk}
      & \textit{circuit} & \textit{cycle} (dirigido) \\
    circuito & circuito & \textit{circuit} & \textit{simple circuit}
      & sin nombre \\
    ciclo & ciclo & \textit{cycle} & sin nombre
      & \textit{simple cycle} \\
    \bottomrule
  \end{tabular}\par}
\end{ojo}

\begin{ejercicio}[Actividad de la diapositiva 89]
  \label{ej:actividad-recorridos}
  En el grafo
  \[
    V = \set{a, b, c, d, e, f} , \qquad
    E = \set{ab, bc, de, ef, ad, be, cf, ae, bf, cd} ,
  \]
  clasificar las secuencias $a,d,c,f,e$; $\ a,b,e,d,a$;
  $\ a,b,e,a,d,a$; $\ d,e,c,a$; $\ a,b,e,d,a,b$.

  \textit{Solución.} Se revisa en este orden: que cada par consecutivo
  sea arista, que no se repitan aristas, que no se repitan vértices y
  si la secuencia cierra. En la figura, el número sobre cada arista es
  el paso en que se usa; la punteada no existe.
  \begin{center}
  \begin{tikzpicture}[x=0.9cm, y=0.9cm]
    \foreach \k/\dx/\dy/\tit in {
        0/0/0/{1. $a,d,c,f,e$},
        1/6.4/0/{3. $a,b,e,a,d,a$},
        2/0/-3.3/{4. $d,e,c,a$},
        3/6.4/-3.3/{5. $a,b,e,d,a,b$}} {
      \begin{scope}[xshift=\dx cm, yshift=\dy cm]
        \foreach \nom/\x/\y in {a/0/1.8, b/2/1.8, c/4/1.8,
                                d/0/0, e/2/0, f/4/0}
          \node[vert] (\nom\k) at (\x, \y) {\nom};
        \foreach \u/\v in {a/b, b/c, d/e, e/f, a/d, b/e, c/f, a/e, b/f,
                           d/c}
          \draw[arista] (\u\k) -- (\v\k);
        \node[nota] at (2, -0.8) {\tit};
      \end{scope}
    }
    \draw[aristaArbol] (a0) -- node[peso] {1} (d0);
    \draw[aristaArbol] (d0) -- node[peso, pos=0.2] {2} (c0);
    \draw[aristaArbol] (c0) -- node[peso] {3} (f0);
    \draw[aristaArbol] (f0) -- node[peso] {4} (e0);
    \draw[aristaArbol] (a1) -- node[peso] {1} (b1);
    \draw[aristaArbol] (b1) -- node[peso] {2} (e1);
    \draw[aristaArbol] (e1) -- node[peso, pos=0.4] {3} (a1);
    \draw[aristaArbol] (a1) -- node[peso] {4,5} (d1);
    \draw[aristaArbol] (d2) -- node[peso] {1} (e2);
    \draw[aristaFuera] (e2) -- node[peso] {?} (c2);
    \draw[aristaArbol] (a3) -- node[peso] {1,5} (b3);
    \draw[aristaArbol] (b3) -- node[peso] {2} (e3);
    \draw[aristaArbol] (e3) -- node[peso] {3} (d3);
    \draw[aristaArbol] (d3) -- node[peso] {4} (a3);
  \end{tikzpicture}
  \end{center}
  \begin{center}
  \small
  \begin{tabular}{llll}
    \toprule
    \textbf{Secuencia} & \textbf{Qué pasa} & \textbf{Libreta}
      & \textbf{Diapositivas} \\
    \midrule
    $a,d,c,f,e$ & nada se repite, abierta & camino, $k = 4$
      & camino \\
    $a,b,e,d,a$ & solo $a$ al cerrar & ciclo, $k = 4$
      & circuito y ciclo \\
    $a,b,e,a,d,a$ & repite $\set{a,d}$, cerrada
      & caminata cerrada & circuito, no ciclo \\
    $d,e,c,a$ & $\set{e,c}$ no es arista & no es caminata
      & no es caminata \\
    $a,b,e,d,a,b$ & repite $\set{a,b}$, abierta
      & caminata, $k = 5$ & caminata \\
    \bottomrule
  \end{tabular}
  \end{center}
  La tercera es la que cambia de nombre: repite la arista $\set{a,d}$,
  así que no es sendero ni circuito, solo una caminata cerrada. Las
  diapositivas 77 y 80 dan otras dos: $a, b, e, a, d, c, f$ es un
  sendero de longitud $6$ que no es camino (repite $a$), y
  $a, b, e, a, d$, que la diapositiva llama ``caminata que no es
  camino'', es más precisamente un sendero.
\end{ejercicio}

Repetir vértices nunca ayuda a llegar: siempre se puede quitar la
vuelta.

\begin{lema}
  \label{lem:caminata-camino}
  Si hay una caminata de $u$ a $v$, hay un camino de $u$ a $v$, y su
  longitud no es mayor.
\end{lema}

\begin{proof}
  Si la caminata repite un vértice $x$, se borra el tramo entre la
  primera y la última aparición de $x$, que es una vuelta cerrada. Lo
  que queda sigue siendo una caminata de $u$ a $v$ y es más corta. Como
  la longitud no puede bajar para siempre, el proceso termina en una
  caminata sin vértices repetidos: un camino.
\end{proof}

En particular, la caminata más corta de $u$ a $v$ es un camino. Por
eso la $d[v]$ de BFS, que mide la caminata más corta, es la longitud
de un camino, y por eso con pesos no negativos siempre hay un camino
más corto que no repite vértices: quitar una vuelta no aumenta el
peso.

La matriz de adyacencia también cuenta caminatas, y eso enlaza esta
subsección con la representación.

\begin{proposicion}
  \label{prop:potencias-adyacencia}
  Si $A$ es la matriz de adyacencia de un multigrafo, con $A_{ij}$
  igual al número de aristas entre $v_i$ y $v_j$, entonces
  $(A^{k})_{ij}$ es el número de caminatas de longitud $k$ de $v_i$ a
  $v_j$.
\end{proposicion}

\begin{proof}
  Por inducción sobre $k$. Con $k = 1$ es la definición de $A$. Una
  caminata de longitud $k + 1$ de $v_i$ a $v_j$ es una caminata de
  longitud $k$ de $v_i$ a algún $v_l$ seguida de una arista de $v_l$ a
  $v_j$; por hipótesis hay $(A^{k})_{il}$ de las primeras y $A_{lj}$
  de las segundas, y sumando sobre $l$ queda
  $\sum_{l} (A^{k})_{il} A_{lj} = (A^{k+1})_{ij}$.
\end{proof}

En el grafo del ejercicio~\ref{ej:actividad-recorridos},
$(A^{2})_{ac} = 2$, por las caminatas $a, b, c$ y $a, d, c$, y la
diagonal de $A^{2}$ es $3, 4, 3, 3, 4, 3$, los grados: una caminata
cerrada de longitud $2$ es ir y volver por una arista. La
sección~\ref{sec:todos-pares} usa la misma idea con pesos, cambiando
la suma por un mínimo y el producto por una suma.

En un grafo dirigido las definiciones son las mismas con arcos: cada
paso va de $v_{i-1}$ a $v_i$ siguiendo la flecha, $(v_{i-1}, v_i) \in
E$. En el grafo dirigido de la sección~\ref{sec:grafos},
$3 \to 4 \to 1 \to 2$ es un camino dirigido (el que da $d_{32} = 6$
en Floyd) y $2 \to 3 \to 1 \to 2$ es un ciclo dirigido de longitud
$3$. Ahí $1 \to 2 \to 1$ es un ciclo de longitud $2$: los arcos
$(1, 2)$ y $(2, 1)$ son distintos, así que no se repite ninguno.

% =====================================================================

\subsection*{Conexidad}

\begin{definicion}[Conexo]
  \label{def:conexo}
  Un grafo no dirigido es \textbf{conexo} si hay un camino entre cada
  par de vértices, y \textbf{desconectado} si no.
\end{definicion}

Por el lema~\ref{lem:caminata-camino} da lo mismo pedir caminos o
caminatas. La sección~\ref{sec:grafos} ya usa la palabra: un árbol es
un grafo no dirigido, conexo y acíclico, y el árbol de expansión
mínima supone un grafo conexo. Ahora las dos palabras tienen
definición, y los árboles se estudian en la
sección~\ref{sec:hamilton-arboles}.

\begin{proposicion}
  \label{prop:conectados-equivalencia}
  La relación ``$u$ y $v$ están conectados'', es decir, hay un camino
  de $u$ a $v$, es de equivalencia en $V$.
\end{proposicion}

\begin{proof}
  Es reflexiva por el camino de longitud $0$, y simétrica porque un
  camino se recorre al revés. Es transitiva porque un camino de $u$ a
  $v$ seguido de uno de $v$ a $w$ es una caminata de $u$ a $w$, y el
  lema~\ref{lem:caminata-camino} la convierte en camino.
\end{proof}

\begin{definicion}[Componente conexa]
  \label{def:componente-conexa}
  Las \textbf{componentes conexas} de $G$ son los subgrafos inducidos
  por las clases de esa relación. Equivalentemente, son los subgrafos
  conexos maximales: conexos, y sin poder agregar un vértice más sin
  dejar de serlo.
\end{definicion}

Las componentes parten a $V$, cada arista queda dentro de una sola
componente, y $G$ es conexo si y solo si tiene una sola. Un vértice
aislado es una componente de un solo vértice.

\begin{ejemplo}[Diapositivas 95 y 100]
  \label{ej:componentes-teoria}
  A la izquierda, un grafo conexo: $a$ y $e$ no son adyacentes, pero
  el camino $a, c, e$ los une. A la derecha, un grafo con tres
  componentes, $\set{a,b,c}$, $\set{d,e}$ y $\set{f,g}$: no hay camino
  de $a$ a $d$ ni de $e$ a $g$.
  \begin{center}
  \begin{tikzpicture}[x=1cm, y=1cm]
    \node[vert] (a) at (0, 1.2) {a};    \node[vert] (b) at (1.4, 1.2) {b};
    \node[vert] (d) at (2.8, 1.2) {d};  \node[vert] (c) at (0.7, 0) {c};
    \node[vert] (e) at (2.8, 0) {e};
    \foreach \u/\v in {a/b, a/c, b/c, b/d, d/e, c/e}
      \draw[arista] (\u) -- (\v);
    \node[nota] at (1.4, -0.75) {conexo};
    \begin{scope}[xshift=5cm]
      \foreach \x/\w in {-0.45/2.1, 1.75/1.5, 3.35/1.5}
        \draw[zona] (\x, -0.45) rectangle ++(\w, 2.1);
      \node[vert] (a) at (0, 1.2) {a};    \node[vert] (b) at (1.2, 1.2) {b};
      \node[vert] (c) at (0.6, 0) {c};
      \node[vert] (d) at (2.2, 1.2) {d};  \node[vert] (e) at (2.8, 0) {e};
      \node[vert] (f) at (3.8, 1.2) {f};  \node[vert] (g) at (4.4, 0) {g};
      \foreach \u/\v in {a/b, a/c, b/c, d/e, f/g}
        \draw[arista] (\u) -- (\v);
      \node[nota] at (2.2, -0.8) {tres componentes};
    \end{scope}
  \end{tikzpicture}
  \end{center}
\end{ejemplo}

Las componentes se calculan con los recorridos de la
sección~\ref{sec:grafos}. Una BFS desde $s$ visita exactamente la
componente de $s$: todo lo que alcanza está conectado con $s$, y todo
lo conectado con $s$ se alcanza. Repetida desde cada vértice que siga
sin visitar, etiqueta todas las componentes en $\Tht{n + m}$, porque
cada vértice entra a la cola una sola vez en total y cada lista se
recorre una sola vez, no una vez por componente. La
sección~\ref{sec:busquedas} hace lo mismo con DFS.

En un grafo dirigido hay dos nociones. Es \textbf{débilmente conexo}
si, olvidando el sentido de los arcos, el grafo no dirigido que queda
es conexo; es \textbf{fuertemente conexo} si hay un camino dirigido de
cada vértice a cada otro. El grafo dirigido de Floyd en la
sección~\ref{sec:grafos} es fuertemente conexo. El DAG de seriaciones
es débilmente conexo y no fuertemente: ninguna flecha regresa. Las
componentes fuertemente conexas y su algoritmo están en la
sección~\ref{sec:aplicaciones-dfs}.

\subsubsection*{Puentes y vértices de corte}

Hay aristas y vértices de los que depende la conexidad: si fallan, la
red se parte.

\begin{definicion}[Puente y vértice de corte]
  \label{def:puente-corte}
  Una arista $e$ es un \textbf{puente} (o arista de corte) si $G - e$
  tiene más componentes conexas que $G$. Un vértice $v$ es un
  \textbf{vértice de corte} (o punto de articulación) si $G - v$ tiene
  más componentes conexas que $G$.
\end{definicion}

Recuérdese que $G - v$ quita también las aristas de $v$. En el grafo
$G$ de la sección~\ref{sec:grafos}, $\set{B, D}$ y $\set{D, G}$ son
puentes y $B$ y $D$ son vértices de corte; ninguna arista del ciclo
$A, B, E, F, C, A$ es puente. Esa última observación es general.

\begin{teorema}
  \label{teo:puente-ciclo}
  Una arista es puente si y solo si no pertenece a ningún ciclo.
\end{teorema}

\begin{proof}
  Sea $e$ una arista de extremos $u$ y $v$. Si $e$ está en un ciclo
  $u, v, x_1, \dots, x_r, u$, el resto del ciclo,
  $v, x_1, \dots, x_r, u$, es un camino de $v$ a $u$ que no usa $e$.
  Cualquier caminata de $G$ que pase por $e$ se desvía por ese camino,
  así que $G - e$ tiene las mismas componentes que $G$, y $e$ no es
  puente. Esto cubre los lazos, que son ciclos de longitud $1$.

  Si $e$ no está en ningún ciclo, $u \neq v$ y $G - e$ no tiene camino
  de $u$ a $v$: ese camino, seguido de $e$, cerraría un ciclo que
  contiene a $e$ (de longitud $2$ si el camino es una arista paralela a
  $e$). Entonces la componente de $G$ que contiene a $u$ y a $v$ se
  parte en dos al quitar $e$, las demás no cambian, y $G - e$ tiene una
  componente más.
\end{proof}

Para vértices no hay una regla así: estar en un ciclo no protege a un
vértice. Lo que sí se tiene es un puente entre la definición de puente
y la de vértice de corte.

\begin{corolario}
  \label{cor:puente-corte}
  Si $\set{u, v}$ es un puente y $u$ tiene otro vecino $w \neq v$,
  entonces $u$ es vértice de corte.
\end{corolario}

\begin{proof}
  Sea $e = \set{u, v}$. En $G - e$, el vértice $w$ sigue unido a $u$ y
  queda en una componente distinta de la de $v$. Un camino de $w$ a
  $v$ en $G - u$ no pasaría por $u$, y por lo tanto tampoco por $e$:
  sería un camino en $G - e$, que no existe.
  Así que al quitar $u$ la componente de $u$ se parte al menos en la
  parte de $w$ y la de $v$, y las demás no cambian.
\end{proof}

\begin{ejercicio}[Actividad de la diapositiva 109]
  \label{ej:puentes-teoria}
  Encontrar los puentes y los vértices de corte del grafo con
  $V = \set{a, \dots, g}$ y
  $E = \set{ab, ac, bc, bd, de, df, ef, eg}$.

  \textit{Solución.} Los ciclos son los triángulos $a, b, c$ y
  $d, e, f$; por el teorema~\ref{teo:puente-ciclo}, los puentes son las
  dos aristas que no están en ellos, $\set{b,d}$ y $\set{e,g}$. Por el
  corolario~\ref{cor:puente-corte}, $b$, $d$ y $e$ son de corte, porque
  cada uno es extremo de un puente y tiene otro vecino; $g$ no, porque
  es una hoja. Falta revisar los demás: quitar $a$, $c$ o $f$ deja un
  triángulo roto que sigue conectado por los otros dos vértices.
  \begin{center}
  \begin{tikzpicture}[x=1cm, y=1cm]
    \node[vert] (a) at (0, 1.3) {a};
    \node[vertAlerta] (b) at (1.5, 1.3) {b};
    \node[vert] (c) at (0.75, 0) {c};
    \node[vertAlerta] (d) at (3.0, 1.3) {d};
    \node[vert] (f) at (3.0, 0) {f};
    \node[vertAlerta] (e) at (4.0, 0) {e};
    \node[vert] (g) at (5.3, 0.6) {g};
    \foreach \u/\v in {a/b, a/c, b/c, d/f, d/e, f/e}
      \draw[arista] (\u) -- (\v);
    \draw[aristaResalte] (b) -- (d);
    \draw[aristaResalte] (e) -- (g);
    \node[nota] at (2.6, -0.75)
      {puentes en naranja, vértices de corte en rojo};
  \end{tikzpicture}
  \end{center}
  El caso fino está en $d$. La arista $\set{d, e}$ no es puente,
  porque está en el triángulo $d, e, f$: si se quita, $e$ y $d$ siguen
  unidos por $f$. Pero el vértice $d$ sí es de corte: quitarlo rompe
  el triángulo y separa $\set{e, f, g}$ de $\set{a, b, c}$.
\end{ejercicio}

\begin{ojo}
  Tres detalles que las diapositivas dejan pasar. Primero, la
  diapositiva 104 dice que una arista en un ciclo ``no suele ser''
  crítica: por el teorema~\ref{teo:puente-ciclo}, nunca lo es. Segundo,
  la diapositiva 102 señala un puente, $\set{b, d}$, pero
  $\set{d, e}$ también lo es: toda arista que llega a una hoja es
  puente. Del mismo modo, en la diapositiva 105 también $d$ es vértice
  de corte, además de $b$. Tercero, dos aristas paralelas nunca son
  puentes, porque juntas forman un ciclo de longitud $2$; eso importa
  en multigrafos como el de Königsberg, que no tiene ningún puente.
\end{ojo}

Por la definición, saber si una arista es puente cuesta una BFS sobre
$G - e$, $\Tht{n + m}$; encontrar todos los puentes así cuesta
$\Tht{m(n + m)}$, y todos los vértices de corte $\Tht{n(n + m)}$. Una
sola DFS con un valor más por vértice, $\text{low}[v]$, encuentra
todos en $\Tht{n + m}$; ese algoritmo está en la
sección~\ref{sec:aplicaciones-dfs}.

% =====================================================================

\subsection*{Grafos eulerianos}

La diapositiva 111 plantea dos preguntas gemelas: ¿hay un recorrido
que pase por cada \emph{arista} exactamente una vez?, ¿hay uno que pase
por cada \emph{vértice} exactamente una vez? La primera es la de
Euler y se resuelve aquí. La segunda es la de Hamilton, y está en la
sección~\ref{sec:hamilton-arboles}; se parecen tanto que uno esperaría
resolverlas igual, y no.

\begin{definicion}[Recorridos eulerianos]
  \label{def:euleriano}
  Un \textbf{sendero euleriano} de $G$ es un sendero que usa cada
  arista de $G$ exactamente una vez; un \textbf{circuito euleriano} es
  un sendero euleriano cerrado. $G$ es \textbf{euleriano} si tiene
  circuito euleriano, \textbf{semi-euleriano} si tiene sendero
  euleriano pero no circuito, y \textbf{no euleriano} si no tiene
  ninguno.
\end{definicion}

\begin{ojo}
  Las diapositivas y Rosen dicen ``camino euleriano'' (\textit{Euler
  path}). Con la convención de esta libreta no es un camino, porque casi
  siempre repite vértices: el recorrido $b, a, d, c, b, d$ de abajo
  pasa dos veces por $b$ y por $d$. Grimaldi lo llama \textit{Euler
  trail}, sendero euleriano, y es el nombre que se usa aquí. Cormen
  llama \textit{Euler tour} al circuito euleriano de un grafo dirigido
  y lo describe como un \textit{cycle} que puede repetir vértices
  (Cormen, problema 22-3): para esta libreta es un circuito.
\end{ojo}

La intuición está en un vértice intermedio. Cada vez que el recorrido
pasa por él usa dos de sus aristas, una para entrar y otra para salir.
Si el recorrido usa todas las aristas una vez, las de cada vértice
quedan agrupadas en pares, salvo en los extremos de un recorrido
abierto.

\subsubsection*{El teorema de Euler}

\begin{teorema}[Euler]
  \label{teo:euler-circuito}
  Sea $G$ un pseudografo con al menos una arista, conexo salvo vértices
  aislados. $G$ tiene un circuito euleriano si y solo si todos sus
  vértices tienen grado par.
\end{teorema}

``Conexo salvo vértices aislados'' significa que todas las aristas
están en una misma componente. Los vértices aislados no estorban: no
tienen aristas que recorrer.

\begin{proof}
  \paso{Necesidad.} Sea $v_0, e_1, v_1, \dots, e_m, v_m$ un circuito
  euleriano, con $v_m = v_0$, y sea $e_{m+1} = e_1$. Cada índice
  $i \in \set{1, \dots, m}$ empareja dos extremos que caen en $v_i$:
  el final de $e_i$, por donde se entra, y el inicio de $e_{i+1}$, por
  donde se sale. Para $i = m$ el par es la última llegada a $v_0$ y la
  primera salida. Como cada arista aparece una sola vez, cada extremo
  que cae en un vértice $v$ queda en exactamente un par, así que
  $\text{gr}(v)$ es dos veces el número de apariciones de $v$ entre
  $v_1, \dots, v_m$. Un lazo en $v$ también encaja: se sale y se
  regresa a $v$ en el mismo paso, y sus dos extremos quedan en pares
  distintos.

  \paso{Un sendero solo se atora en su inicio.} Supóngase que todos los
  grados de un pseudografo son pares, y que desde un vértice $s$ se
  camina usando cada vez una arista todavía no usada, mientras la haya.
  Estando en un $v \neq s$, el sendero ha entrado a $v$ una vez más de
  las que ha salido, así que ha usado un número impar de los extremos
  de $v$ (un lazo, entrado y salido en un paso, usa dos). Como
  $\text{gr}(v)$ es par, queda un extremo libre y se puede seguir. Las
  aristas son finitas, así que la caminata se detiene, y solo puede
  hacerlo en $s$: el resultado es un circuito.

  \paso{Suficiencia.} Sea $s$ un vértice con aristas. Por el paso
  anterior, desde $s$ sale un circuito $C$. Si $C$ usa todas las
  aristas, terminamos. Si no, alguna arista sin usar toca un vértice
  $x$ de $C$. En efecto, si una arista sin usar $\set{y, z}$ no toca a
  $C$, se toma un camino de $s$ a $y$, que existe porque todas las
  aristas están en una componente; las aristas de $C$ unen vértices de
  $C$ y $y$ no es uno de ellos, así que la primera arista del camino
  que no está en $C$ sale de un vértice $x$ de $C$ y está sin usar.

  Sea $R$ el grafo de las aristas sin usar. En $R$ todos los grados
  siguen siendo pares, porque $C$ usa en cada vértice un número par de
  extremos (es la necesidad aplicada a $C$). Desde $x$ sale entonces
  un circuito $C'$ en $R$, que regresa a $x$. Se \textbf{empalma}
  $C'$ en $C$: en una aparición de $x$ en $C$ se inserta $C'$
  completo. El resultado es un circuito, porque $C$ y $C'$ no
  comparten aristas, y tiene más aristas que $C$. Se repite; como $m$
  es finito, al final el circuito usa todas las aristas.
\end{proof}

La prueba de la suficiencia es constructiva: camina, se atora, busca
un vértice con aristas libres y empalma. Es exactamente el algoritmo
de Hierholzer (1873), que se ve más abajo. Las diapositivas solo dan
la necesidad.

\begin{teorema}[Sendero euleriano]
  \label{teo:euler-sendero}
  Sea $G$ un pseudografo conexo salvo vértices aislados. $G$ tiene un
  sendero euleriano que no es cerrado si y solo si tiene exactamente
  dos vértices de grado impar, y en ese caso todo sendero euleriano
  empieza en uno de ellos y termina en el otro.
\end{teorema}

\begin{proof}
  Si hay un sendero euleriano de $u$ a $v$ con $u \neq v$, las
  apariciones intermedias de cada vértice emparejan sus extremos como
  en la prueba anterior; en $u$ sobra la primera salida y en $v$ la
  última llegada. Así que $u$ y $v$ tienen grado impar y los demás
  par.

  Al revés, si $u$ y $v$ son los únicos vértices de grado impar, se
  agrega una arista nueva $e^{*} = \set{u, v}$, que puede quedar
  paralela a otra; por eso el teorema se enuncia para multigrafos. Ahora
  todos los grados son pares y, por el teorema~\ref{teo:euler-circuito},
  hay un circuito euleriano. Como es cerrado, se puede empezar a leer
  justo después de $e^{*}$: queda $u, \dots, v, e^{*}, u$. Borrando
  $e^{*}$ queda un sendero de $u$ a $v$ que usa cada arista original
  una vez.
\end{proof}

Con los dos teoremas y el corolario~\ref{cor:impares-par}, la
clasificación de un grafo conexo salvo aislados se lee contando
vértices de grado impar:

\begin{center}
\small
\begin{tabular}{cll}
  \toprule
  \textbf{Vértices impares} & \textbf{Tipo} & \textbf{Recorrido} \\
  \midrule
  $0$ & euleriano & circuito desde cualquier vértice con aristas \\
  $2$ & semi-euleriano & sendero de un impar al otro \\
  $4, 6, \dots$ & no euleriano & ninguno \\
  $1, 3, 5, \dots$ & imposible & por el corolario~\ref{cor:impares-par}
    \\
  \bottomrule
\end{tabular}
\end{center}

La tabla de la diapositiva 124 dice ``más de $2$''; el último renglón
explica por qué eso significa $4$ o más y por qué la tabla no deja
casos fuera.

\begin{ejemplo}[Diapositivas 115 y 119]
  \label{ej:euler-teoria}
  A la izquierda, el grafo con $E = \set{ab, ad, bd, bc, dc}$ tiene
  grados $\text{gr}(a) = \text{gr}(c) = 2$ y
  $\text{gr}(b) = \text{gr}(d) = 3$: es semi-euleriano, y el sendero
  $b, a, d, c, b, d$ va de un impar (en verde) al otro. A la derecha,
  la mariposa con $E = \set{ab, ac, bc, cd, ce, de}$ tiene
  $\text{gr}(c) = 4$ y $2$ en los demás: es euleriana, y
  $c, a, b, c, d, e, c$ es un circuito euleriano. El número sobre cada
  arista es el orden en que se usa.
  \begin{center}
  \begin{tikzpicture}[x=1cm, y=1cm]
    \node[vertOk] (b) at (1.5, 1.3) {b};  \node[vert] (a) at (0, 1.3) {a};
    \node[vertOk] (d) at (0, 0) {d};      \node[vert] (c) at (3.0, 0) {c};
    \draw[aristaArbol] (b) -- node[peso] {1} (a);
    \draw[aristaArbol] (a) -- node[peso] {2} (d);
    \draw[aristaArbol] (d) -- node[peso] {3} (c);
    \draw[aristaArbol] (c) -- node[peso] {4} (b);
    \draw[aristaArbol] (b) -- node[peso] {5} (d);
    \node[nota] at (1.5, -0.7) {sendero $b,a,d,c,b,d$};
    \begin{scope}[xshift=5.4cm]
      \node[vert] (b) at (0, 0) {b};      \node[vert] (c) at (1.4, 0) {c};
      \node[vert] (d) at (2.8, 0) {d};    \node[vert] (a) at (1.4, 1.4) {a};
      \node[vert] (e) at (3.1, 1.4) {e};
      \draw[aristaArbol] (c) -- node[peso] {1} (a);
      \draw[aristaArbol] (a) -- node[peso] {2} (b);
      \draw[aristaArbol] (b) -- node[peso] {3} (c);
      \draw[aristaArbol] (c) -- node[peso] {4} (d);
      \draw[aristaArbol] (d) -- node[peso] {5} (e);
      \draw[aristaArbol] (e) -- node[peso] {6} (c);
      \node[nota] at (1.5, -0.7) {circuito $c,a,b,c,d,e,c$};
    \end{scope}
  \end{tikzpicture}
  \end{center}
  El circuito de la mariposa pasa dos veces por $c$: es circuito, no
  ciclo.
\end{ejemplo}

\begin{ojo}
  La hipótesis de conexidad no sobra, pero solo se pide entre los
  vértices que tienen aristas. Dos triángulos ajenos tienen todos sus
  grados pares y ningún circuito euleriano: no hay forma de pasar de
  uno al otro. En cambio, un triángulo más un vértice suelto sí tiene
  circuito euleriano.
\end{ojo}

\begin{ejercicio}[Actividad de la diapositiva 128]
  Clasificar $P_4$, $C_4$ y $K_4$, del ejemplo~\ref{ej:familias-teoria},
  en euleriano, semi-euleriano o no euleriano. ¿Y la estrella con tres
  hojas?

  \textit{Solución.} Se cuentan los vértices de grado impar. $C_4$ no
  tiene ninguno: euleriano. $P_4$ tiene dos, sus extremos $a$ y $d$:
  semi-euleriano, con el sendero $a, b, c, d$. $K_4$ tiene cuatro
  vértices de grado $3$: no euleriano. La estrella tiene el centro de
  grado $3$ y tres hojas de grado $1$, cuatro impares: tampoco. En
  general, para $n \geq 3$, $K_n$ es euleriano si y solo si $n$ es
  impar, porque todos sus grados valen $n - 1$.
\end{ejercicio}

Ya se puede contestar la pregunta de Euler. En Königsberg,
$\text{gr}(I) = 5$ y $\text{gr}(N) = \text{gr}(S) = \text{gr}(E) = 3$.

\begin{resultado}
  Los cuatro vértices del grafo de Königsberg tienen grado impar. No
  hay circuito euleriano ni sendero euleriano: no se pueden cruzar los
  siete puentes una sola vez cada uno, ni regresando al punto de
  partida ni sin regresar.
\end{resultado}

\begin{ejercicio}
  \label{ej:konigsberg-puentes}
  ¿Cuántos puentes nuevos necesita Königsberg para que exista un
  sendero euleriano? ¿Y para un circuito? ¿Qué pasa con el grafo $G$
  de la sección~\ref{sec:grafos} y con el del
  ejercicio~\ref{ej:actividad-recorridos}?

  \textit{Solución.} Cada arista nueva cambia la paridad de sus dos
  extremos. Con un puente $\set{N, S}$, $N$ y $S$ pasan a grado $4$ y
  quedan dos impares, $I$ y $E$: hay un sendero euleriano, por ejemplo
  $I, E, N, I, N, S, I, S, E$. Con un segundo puente $\set{I, E}$ todos
  los grados son pares, y hay un circuito, por ejemplo
  $N, E, I, E, S,\allowbreak I, N, I, S, N$. En general, con $2k$ vértices
  impares bastan $k$ aristas nuevas que los emparejen para tener
  circuito, y $k - 1$ para tener sendero. El grafo $G$ tiene dos
  impares, $B$ (grado $3$) y $G$ (grado $1$): es semi-euleriano, con el
  sendero $B, A, C, F, E, B, D, G$. El de la actividad tiene cuatro,
  $a$, $c$, $d$ y $f$: no es euleriano, y con una arista nueva
  $\set{a, c}$ pasa a semi-euleriano, de $d$ a $f$.
\end{ejercicio}

\subsubsection*{El caso dirigido}

En un grafo dirigido el recorrido tiene que seguir las flechas, y el
par entrada-salida de cada vértice intermedio se vuelve un arco que
llega y uno que sale. La paridad se cambia por el balance.

\begin{teorema}[Euler, caso dirigido]
  \label{teo:euler-dirigido}
  Sea $G$ un grafo dirigido con al menos un arco, débilmente conexo
  salvo vértices aislados.
  \begin{enumerate}[itemsep=0.1em]
    \item $G$ tiene un circuito euleriano si y solo si
          $\text{gr}^{-}(v) = \text{gr}^{+}(v)$ para todo $v$.
    \item $G$ tiene un sendero euleriano no cerrado si y solo si hay un
          vértice $s$ con $\text{gr}^{+}(s) - \text{gr}^{-}(s) = 1$, uno
          $t$ con $\text{gr}^{-}(t) - \text{gr}^{+}(t) = 1$, y todos
          los demás balanceados. El sendero va de $s$ a $t$.
  \end{enumerate}
\end{teorema}

\begin{proof}
  Es la misma prueba con dos ajustes. En la necesidad, cada paso por
  un vértice usa un arco que llega y uno que sale, así que en un
  circuito cada vértice queda balanceado; en un sendero abierto, el
  inicio tiene una salida de más y el final una llegada de más. En el
  paso de atorarse, estando en $v \neq s$ se ha llegado a $v$ una vez
  más de las que se ha salido, y como $v$ está balanceado le queda un
  arco de salida libre. En el empalme, la conexidad débil da un arco
  libre que toca un vértice $x$ del circuito, pero podría llegar a $x$
  en lugar de salir; como en los arcos libres todo sigue balanceado,
  si a $x$ le llega un arco libre también le sale uno. El caso del
  sendero se reduce al del circuito agregando el arco $(t, s)$.
\end{proof}

Cormen enuncia la parte 1 para grafos fuertemente conexos
(problema 22-3). No es una hipótesis más fuerte en el fondo: si el
grafo está balanceado y es débilmente conexo, el circuito euleriano
pasa por todos los vértices con arcos, y siguiéndolo se llega de
cualquiera de ellos a cualquier otro. Balanceado y débilmente conexo
implica fuertemente conexo.

\begin{ejemplo}[Una cadena de De Bruijn]
  \label{ej:de-bruijn}
  Se busca la cadena binaria cíclica más corta que contenga, como
  ventana de tres símbolos consecutivos, cada una de las ocho cadenas
  de 3 bits. Una cerradura de dos botones que se abre con los últimos
  tres oprimidos cae tecleando esa cadena: $10$ pulsaciones en lugar
  de $24$. El grafo tiene un vértice por
  cada cadena de 2 bits y, por cada cadena $xyz$ de 3 bits, un arco de
  $xy$ a $yz$, con lazos para $000$ y $111$.
  \begin{center}
  \begin{tikzpicture}[x=1cm, y=1cm]
    \node[vert] (p) at (0, 0) {00};
    \node[vert] (q) at (2.4, 1.1) {01};
    \node[vert] (r) at (2.4, -1.1) {10};
    \node[vert] (t) at (4.8, 0) {11};
    \draw[dirigida] (p) to[out=150, in=210, looseness=6]
      node[peso, left] {000} (p);
    \draw[dirigida] (p) -- node[peso] {001} (q);
    \draw[dirigida] (q) -- node[peso] {011} (t);
    \draw[dirigida] (t) to[out=30, in=330, looseness=6]
      node[peso, right] {111} (t);
    \draw[dirigida] (t) -- node[peso] {110} (r);
    \draw[dirigida] (r) -- node[peso] {100} (p);
    \draw[dirigida] (q) to[bend left=22] node[peso] {010} (r);
    \draw[dirigida] (r) to[bend left=22] node[peso] {101} (q);
  \end{tikzpicture}
  \end{center}
  Cada vértice tiene $\text{gr}^{-} = \text{gr}^{+} = 2$ y el grafo es
  conexo: hay circuito euleriano. Desde $00$, con las listas en orden,
  el algoritmo de Hierholzer de abajo da
  $00, 00, 01, 10, 01, 11, 11, 10, 00$. Empezando con $00$ y agregando
  el último bit de cada arco sale $0001011100$; como el circuito
  regresa a $00$, los dos últimos símbolos repiten los dos primeros, y
  la cadena cíclica es $00010111$. Sus ocho ventanas, dando la vuelta
  al final, son $000, 001, 010, 101, 011, 111, 110, 100$, cada cadena
  una vez. No hay una más corta, porque ocho ventanas necesitan ocho
  posiciones. La misma construcción, con fragmentos de ADN en lugar de
  bits, es la base de varios ensambladores de genomas.
\end{ejemplo}

\subsubsection*{El algoritmo de Hierholzer}

Decidir el tipo cuesta lo mismo que contar grados y correr una BFS.

\begin{algorithm}[H]
\caption{\textsc{Tipo-Euler}$(G)$}
\label{alg:tipo-euler}
\begin{algorithmic}[1]
\Procedure{Tipo-Euler}{$G$}
  \State $I \gets \set{v \in V : \text{gr}(v) \text{ es impar}}$
  \State correr BFS desde un vértice $s$ con $\text{gr}(s) > 0$
  \If{algún $v$ con $\text{gr}(v) > 0$ quedó sin visitar}
    \State \Return no euleriano \Comment{aristas en dos componentes}
  \ElsIf{$\abs{I} = 0$}
    \State \Return euleriano: circuito desde cualquier $s$ con
           $\text{gr}(s) > 0$
  \ElsIf{$\abs{I} = 2$}
    \State \Return semi-euleriano: sendero de un impar al otro
  \Else
    \State \Return no euleriano \Comment{$\abs{I} \geq 4$}
  \EndIf
\EndProcedure
\end{algorithmic}
\end{algorithm}

Los grados salen de las longitudes de las listas en $\Tht{n + m}$, y
la BFS cuesta $\Tht{n + m}$. La rama final es la única que queda,
porque $\abs{I}$ es par. Si $m = 0$ no hay nada que recorrer y no
hace falta llamar al algoritmo.

Construir el recorrido también es lineal. La prueba del
teorema~\ref{teo:euler-circuito} empalma subcircuitos; Hierholzer lo
organiza con una pila para que el empalme salga solo, sin manipular
listas.

\begin{algorithm}[H]
\caption{\textsc{Hierholzer}$(G, s)$}
\label{alg:hierholzer}
\begin{algorithmic}[1]
\Procedure{Hierholzer}{$G, s$}
  \ForAll{$e \in E$}
    \State $\text{usada}[e] \gets$ falso
  \EndFor
  \ForAll{$v \in V$}
    \State $\text{sig}[v] \gets$ primera entrada de la lista de $v$
  \EndFor
  \State $P \gets$ pila vacía; \ \textsc{Apilar}$(P, s)$
  \State $C \gets$ lista vacía
  \While{$P$ no está vacía}
    \State $u \gets \textsc{Tope}(P)$
    \While{$\text{sig}[u]$ apunta a una arista usada}
      \State avanzar $\text{sig}[u]$
    \EndWhile
    \If{$\text{sig}[u]$ apunta a una arista $e = \set{u, w}$}
      \State $\text{usada}[e] \gets$ verdadero
      \State \textsc{Apilar}$(P, w)$ \Comment{seguir caminando}
    \Else \Comment{$u$ ya no tiene aristas libres}
      \State \textsc{Desapilar}$(P)$; \ agregar $u$ al frente de $C$
    \EndIf
  \EndWhile
  \State \Return $C$
\EndProcedure
\end{algorithmic}
\end{algorithm}

Si el grafo es euleriano se llama con cualquier $s$ de grado positivo;
si es semi-euleriano, \textbf{obligatoriamente} con uno de los dos
impares. Las líneas que no son obvias:
\begin{itemize}[itemsep=0.1em]
  \item Líneas 12 a 14: $\text{sig}[u]$ es un apuntador que avanza
        sobre la lista de $u$ y nunca retrocede. Cada arista
        $\set{u, w}$ aparece en la lista de $u$ y en la de $w$, con el
        mismo identificador; al usarla desde $u$ se marca, y la copia
        en la lista de $w$ se salta cuando el apuntador de $w$ pase por
        ella.
  \item Líneas 15 a 17: mientras el tope tenga aristas libres se
        camina, y la pila guarda la caminata actual.
  \item Línea 19: un vértice sin aristas libres ya no puede extender
        nada, y pasa al frente de $C$. Al desapilar se regresa por la
        caminata; si aparece un vértice con aristas libres, desde él
        sale un subcircuito, cuyos vértices se desapilarán más tarde y
        quedarán en $C$ justo delante de ese vértice. Ese es el empalme
        de la prueba.
\end{itemize}
Por qué el subcircuito regresa a su vértice de arranque $x$: en ese
momento las aristas usadas forman un sendero del inicio a $x$ (la
pila) y otro de $x$ al final (lo que ya está en $C$). Por la paridad
de la prueba, en las aristas libres todos los grados son pares, y el
paso de atorarse dice que la nueva caminata solo se detiene en $x$.
La primera caminata es la excepción en el caso semi-euleriano: sale
de un impar y se atora en el otro.
$C$ termina con $m + 1$ vértices, en el orden del recorrido.

\begin{ejemplo}[Hierholzer sobre la mariposa desde $a$]
  \label{ej:hierholzer-traza}
  Desde $c$, como en la diapositiva, la primera caminata ya usa todo y
  no se ve el empalme; desde $a$ sí. Las listas, en orden alfabético:
  $a{:}\ b, c$; $\ b{:}\ a, c$; $\ c{:}\ a, b, d, e$; $\ d{:}\ c, e$;
  $\ e{:}\ c, d$.
  \begin{center}
  \small
  \begin{tabular}{cclll}
    \toprule
    \textbf{Paso} & \textbf{Tope} & \textbf{Acción} & \textbf{Pila $P$}
      & \textbf{$C$} \\
    \midrule
    1 & $a$ & usa $\set{a,b}$ y apila $b$ & $a, b$ & \\
    2 & $b$ & usa $\set{b,c}$ y apila $c$ & $a, b, c$ & \\
    3 & $c$ & usa $\set{c,a}$ y apila $a$ & $a, b, c, a$ & \\
    4 & $a$ & sin aristas libres: desapila & $a, b, c$ & $a$ \\
    5 & $c$ & usa $\set{c,d}$ y apila $d$ & $a, b, c, d$ & $a$ \\
    6 & $d$ & usa $\set{d,e}$ y apila $e$ & $a, b, c, d, e$ & $a$ \\
    7 & $e$ & usa $\set{e,c}$ y apila $c$ & $a, b, c, d, e, c$ & $a$ \\
    8 & $c$ & sin aristas libres: desapila & $a, b, c, d, e$ & $c, a$ \\
    9 & $e$ & desapila & $a, b, c, d$ & $e, c, a$ \\
    10 & $d$ & desapila & $a, b, c$ & $d, e, c, a$ \\
    11 & $c$ & desapila & $a, b$ & $c, d, e, c, a$ \\
    12 & $b$ & desapila & $a$ & $b, c, d, e, c, a$ \\
    13 & $a$ & desapila & vacía & $a, b, c, d, e, c, a$ \\
    \bottomrule
  \end{tabular}
  \end{center}
  En el paso 4 la caminata se atora en $a = s$, como predice la
  paridad, después de cerrar $a, b, c, a$ y con $\set{c,d}$,
  $\set{d,e}$, $\set{e,c}$ sin usar. Al retroceder, $c$ todavía tiene
  aristas libres, y de él sale el subcircuito $c, d, e, c$, que queda
  empalmado en $c$. Las tres instantáneas: en azul lo recorrido antes,
  en naranja el subcircuito.
  \begin{center}
  \begin{tikzpicture}[x=1cm, y=1cm]
    \foreach \k/\tit in {0/{se atora en $a$: cierra $a,b,c,a$},
                         1/{desde $c$ sale $c,d,e,c$},
                         2/{empalme: $a,b,c,d,e,c,a$}} {
      \begin{scope}[xshift=\k*5cm]
        \node[vert] (b\k) at (0, 0) {b};
        \node[vert] (c\k) at (1.4, 0) {c};
        \node[vert] (d\k) at (2.8, 0) {d};
        \node[vert] (a\k) at (1.4, 1.4) {a};
        \node[vert] (e\k) at (3.1, 1.4) {e};
        \node[nota] at (1.5, -0.7) {\tit};
      \end{scope}
    }
    \draw[aristaArbol] (a0) -- node[peso] {1} (b0);
    \draw[aristaArbol] (b0) -- node[peso] {2} (c0);
    \draw[aristaArbol] (c0) -- node[peso] {3} (a0);
    \foreach \u/\v in {c/d, d/e, e/c} \draw[arista] (\u0) -- (\v0);
    \foreach \u/\v in {a/b, b/c, c/a} \draw[aristaArbol] (\u1) -- (\v1);
    \draw[aristaResalte] (c1) -- node[peso] {4} (d1);
    \draw[aristaResalte] (d1) -- node[peso] {5} (e1);
    \draw[aristaResalte] (e1) -- node[peso] {6} (c1);
    \draw[aristaArbol] (a2) -- node[peso] {1} (b2);
    \draw[aristaArbol] (b2) -- node[peso] {2} (c2);
    \draw[aristaResalte] (c2) -- node[peso] {3} (d2);
    \draw[aristaResalte] (d2) -- node[peso] {4} (e2);
    \draw[aristaResalte] (e2) -- node[peso] {5} (c2);
    \draw[aristaArbol] (c2) -- node[peso] {6} (a2);
    \node[vertOk] at (c2) {c};
  \end{tikzpicture}
  \end{center}
  El resultado, $a, b, c, d, e, c, a$, usa las seis aristas una vez
  cada una.
\end{ejemplo}

\textbf{Costo.} Cada arista se marca una sola vez y provoca un solo
\textsc{Apilar}: hay $m + 1$ apilados contando $s$, y como cada vértice
apilado se desapila una vez, $m + 1$ desapilados. El trabajo de las
líneas 12 a 14 se paga con el apuntador: $\text{sig}[u]$ recorre la
lista de $u$ una sola vez en todo el algoritmo, no una vez por visita
a $u$, y las listas suman $2m$ entradas. Inicializar cuesta
$\Tht{n + m}$. En total, $\Tht{n + m}$ con listas de adyacencia. Con
matriz, guardando multiplicidades, el apuntador de cada renglón barre
$n$ columnas y el total es $\Tht{n^{2} + m}$.

En un grafo dirigido el algoritmo es el mismo con la lista de arcos
que salen de cada vértice. Cada arco aparece en una sola lista, la de
su origen, así que ni siquiera hacen falta las marcas: basta avanzar
$\text{sig}[u]$ cada vez que se usa un arco. Así se obtuvo el circuito
del ejemplo~\ref{ej:de-bruijn}.

\begin{ojo}
  Hierholzer no comprueba nada: supone que el recorrido existe y que
  $s$ es un inicio válido. Primero se clasifica con
  \textsc{Tipo-Euler} y luego se elige $s$. En el grafo semi-euleriano
  del ejemplo~\ref{ej:euler-teoria}, desde $b$ (impar) devuelve
  $b, a, d, b, c, d$, un sendero euleriano distinto del de la
  diapositiva e igual de válido. Desde $a$ (par) devuelve
  $a, b, c, d, b, a$, que repite $\set{a, b}$ y omite $\set{a, d}$: ni
  siquiera es un sendero. La caminata de un vértice par se atora en el
  otro impar, no en su inicio, y el empalme queda mal hecho.
\end{ojo}

Varios textos presentan en su lugar el algoritmo de Fleury: caminar
sin cruzar un puente del grafo restante salvo que no haya otra opción.
Es correcto, pero cada paso tiene que averiguar si una arista es
puente, y con la prueba por definición eso cuesta $\Ord{m(n + m)}$ en
total. Hierholzer evita la pregunta: no le importa atorarse, porque
sabe dónde se va a atorar.

% =====================================================================

\subsection*{Resumen}

\begin{center}
\small
\begin{tabular}{>{\raggedright\arraybackslash}p{4.3cm}%
                >{\raggedright\arraybackslash}p{5.4cm}%
                >{\raggedright\arraybackslash}p{3.6cm}}
  \toprule
  \textbf{Pregunta} & \textbf{Cómo se decide} & \textbf{Costo} \\
  \midrule
  ¿Es conexo?, ¿cuáles son sus componentes? & BFS o DFS desde cada
    vértice sin visitar & $\Tht{n + m}$ \\
  \addlinespace
  ¿Es $e$ un puente?, ¿es $v$ de corte? & quitarlo y contar
    componentes; todos a la vez con $\text{low}$
    (sección~\ref{sec:aplicaciones-dfs}) & $\Tht{n + m}$ \\
  \addlinespace
  ¿Hay circuito o sendero euleriano? & contar vértices impares y
    revisar conexidad & $\Tht{n + m}$ \\
  \addlinespace
  Construir el recorrido & Hierholzer & $\Tht{n + m}$ \\
  \addlinespace
  ¿Son isomorfos? & descartar con invariantes; confirmar con una
    biyección & sin algoritmo polinomial conocido \\
  \addlinespace
  ¿Hay ciclo hamiltoniano? & sección~\ref{sec:hamilton-arboles}
    & NP-completo \\
  \bottomrule
\end{tabular}
\end{center}

\begin{idea}
  Euler contra Hamilton es el contraste que importa en un curso de
  algoritmos. Recorrer cada arista una vez se decide localmente,
  vértice por vértice, contando paridades, y se construye en tiempo
  lineal. Recorrer cada vértice una vez no tiene una condición local
  equivalente, y es NP-completo. Dos preguntas que suenan iguales
  quedan en los dos extremos de la escala de dificultad.
\end{idea}

\newpage

% =====================================================================
%  Grafos - Hamilton, bipartitos, cliques y arboles
% =====================================================================

\clase{Hamilton, bipartitos, cliques y árboles}%
      {Complemento: Rosen, caps.\ 10 y 11}
\label{sec:hamilton-arboles}

\begin{quick_sum}
  La sección \ref{sec:teoria-grafos} buscó recorridos que usan cada
  arista una vez. Esta empieza con la pregunta gemela, pasar por cada
  vértice una vez, y muestra que el parecido engaña: Euler se decide
  en tiempo lineal y Hamilton es NP-completo. Se ven las condiciones
  necesarias, la suficiente de Dirac, la búsqueda con retroceso y el
  agente viajero con la programación dinámica de Held y Karp, que
  baja de $(n-1)!/2$ rutas a $\Tht{n^2 2^n}$ pasos. Luego los grafos
  bipartitos, que son exactamente los que no tienen ciclos impares,
  con un algoritmo de BFS que en $\Tht{n + m}$ entrega la partición o
  un ciclo impar como testigo. Las cliques dan el contraste entre
  maximal y máxima. La segunda mitad es de árboles: sus seis
  caracterizaciones equivalentes, árboles con raíz, árboles $k$-arios
  llenos, casi completos y perfectos, las fórmulas de conteo, el
  balanceo, los tres recorridos y la evaluación de expresiones
  posfijas con una pila. La materia prima son las diapositivas
  \textit{Grafos 2} de la Dra.\ Kelsey Ramírez, que siguen a Rosen;
  los algoritmos y sus costos son agregados.
\end{quick_sum}

\subsection*{Caminos y ciclos hamiltonianos}

Un recorrido euleriano pasa por cada \textbf{arista} exactamente una
vez. Si se cambia arista por \textbf{vértice} sale otra pregunta, que
lleva el nombre de W.\ R.\ Hamilton por un juego suyo de 1857: dar la
vuelta a los $20$ vértices de un dodecaedro sin repetir ninguno.

\begin{definicion}[Camino y ciclo hamiltonianos]
  \label{def:hamilton-camino}
  Sea $G = (V, E)$ con $n$ vértices. Un \textbf{camino hamiltoniano}
  es un camino $v_1, v_2, \dots, v_n$ que pasa por todos los
  vértices, cada uno exactamente una vez. Un \textbf{ciclo
  hamiltoniano} es un ciclo $v_1, v_2, \dots, v_n, v_1$ que pasa por
  todos; el único repetido es $v_1$, por ser inicio y final. Un grafo
  con ciclo hamiltoniano es un \textbf{grafo hamiltoniano}.
\end{definicion}

Camino y ciclo van en el sentido de la sección
\ref{sec:teoria-grafos}: no repiten vértices. En un grafo simple el
ciclo hamiltoniano exige $n \ge 3$, porque en $K_2$ la sucesión
$a, b, a$ usaría dos veces la misma arista. Quitarle una arista a un
ciclo hamiltoniano deja un camino hamiltoniano; al revés no siempre
se puede cerrar.

\begin{ejemplo}[Hamilton no usa todas las aristas]
  \label{ej:hamilton-pentagono}
  Junto a cada vértice de la figura está su grado; las aristas son
  \[
    E = \set{ab, bc, cd, de, ea, ac, ce, bd} .
  \]
  \begin{center}
  \begin{tikzpicture}
    \foreach \v/\ang in {a/90, b/18, c/-54, d/-126, e/162}
      \node[vert] (\v) at (\ang:1.4) {\v};
    \foreach \x/\y in {a/b, b/c, c/d, d/e, e/a}
      \draw[aristaArbol] (\x) -- (\y);
    \foreach \x/\y in {a/c, c/e, b/d}
      \draw[arista] (\x) -- (\y);
    \foreach \v/\g/\lado/\ancla in {a/3/above/north, b/3/right/east,
        c/4/below right/south east, d/3/below left/south west,
        e/3/left/west}
      \node[nota, \lado=1pt] at (\v.\ancla) {\g};
  \end{tikzpicture}
  \end{center}
  El ciclo a, b, c, d, e, a, en azul, visita los cinco vértices y deja
  fuera las tres diagonales. Hay otros tres ciclos hamiltonianos
  distintos (cuatro en total, sin contar sentido ni punto de partida).
  Un ciclo hamiltoniano usa exactamente $n$ aristas, tenga el grafo
  las que tenga; un circuito euleriano las usa todas.
\end{ejemplo}

\subsubsection*{Condiciones necesarias}

Lo primero que se revisa antes de buscar un ciclo es lo que lo hace
imposible.

\begin{teorema}[Condiciones de grado]
  \label{teo:hamilton-grado}
  Sea $G$ un grafo simple con $n \ge 3$ vértices.
  \begin{enumerate}[label=(\alph*), itemsep=0pt, parsep=0pt]
    \item Si $G$ tiene ciclo hamiltoniano, $\text{gr}(v) \ge 2$ para
          todo $v$.
    \item Si $G$ tiene más de dos vértices de grado $1$, no tiene
          camino hamiltoniano.
  \end{enumerate}
\end{teorema}

\begin{proof}
  (a) En el ciclo cada vértice tiene un antecesor y un sucesor
  distintos, y las aristas que lo unen con ellos son dos aristas
  distintas. (b) En un camino, cada vértice interior usa dos aristas
  del camino, así que uno de grado $1$ solo puede ser extremo, y un
  camino tiene dos extremos.
\end{proof}

\begin{teorema}[Vértices de corte]
  \label{teo:hamilton-corte}
  Si $G$ tiene ciclo hamiltoniano, quitar cualquier conjunto $S$ de
  $k \ge 1$ vértices deja a lo más $k$ componentes conexas. En
  particular, un grafo con un vértice de corte no es hamiltoniano.
\end{teorema}

\begin{proof}
  Quitar $k$ vértices de un ciclo lo parte en a lo más $k$ tramos, y
  cada tramo es un camino, que es conexo. Los tramos cubren todos los
  vértices que quedan y sus aristas siguen en $G - S$, así que
  $G - S$ tiene a lo más $k$ componentes. Con $k = 1$ sale lo
  segundo: el ciclo sin un vértice es un camino que toca a todos los
  demás.
\end{proof}

Para caminos hamiltonianos la cuenta es la misma con un tramo más:
quitar $k$ vértices deja a lo más $k + 1$ componentes.

\begin{ejemplo}[Tres grafos que no pasan]
  \label{ej:hamilton-no-pasan}
  Los tres grafos de las diapositivas que fallan alguna condición;
  en rojo, los vértices que estorban.
  \begin{center}
  \begin{tikzpicture}
    \node[vertAlerta] (a) at (0, 0.6) {a};
    \node[vert] (b) at (1.6, 0.6) {b};
    \node[vert] (c) at (1.6, -1) {c};
    \node[vert] (d) at (0, -1) {d};
    \foreach \x/\y in {a/b, b/c, c/d}
      \draw[aristaArbol] (\x) -- (\y);
    \draw[arista] (d) -- (b);
    \node[nota] at (0.8, -1.8) {(1) camino a, b, c, d; sin ciclo};
    \begin{scope}[xshift=5.2cm, yshift=0.2cm]
      \node[vertAlerta] (a) at (-0.9, 1.3) {a};
      \node[vert] (g) at (0.9, 1.5) {g};
      \node[vert] (b) at (0, 0.7) {b};
      \node[vert] (c) at (0, -0.5) {c};
      \node[vert] (e) at (1.6, -0.5) {e};
      \node[vertAlerta] (d) at (-0.9, -1.2) {d};
      \node[vertAlerta] (f) at (2.5, -1.2) {f};
      \foreach \x/\y in {a/b, b/g, b/c, c/d, c/e, g/e, e/f}
        \draw[arista] (\x) -- (\y);
      \node[nota] at (0.8, -2.0) {(2) tres vértices de grado 1};
    \end{scope}
    \begin{scope}[xshift=11cm, yshift=0.2cm]
      \node[vert] (v) at (0, 0) {v};
      \foreach \h/\ang in {a/90, c/18, e/-54, d/-126, b/162}
        \node[vertAlerta] (\h) at (\ang:1.2) {\h};
      \foreach \h in {a, b, c, d, e}
        \draw[arista] (v) -- (\h);
      \node[nota] at (0, -2.0) {(3) la estrella $K_{1,5}$};
    \end{scope}
  \end{tikzpicture}
  \end{center}
  (1) Con $E = \set{ab, bc, cd, db}$ hay camino hamiltoniano, a, b, c,
  d (y también a, b, d, c), pero no ciclo: $\text{gr}(a) = 1$. Además
  b es vértice de corte. (2) Con
  $E = \set{ab, bg, bc, cd, ce, ge, ef}$, los vértices a, d y f tienen
  grado $1$: ni siquiera hay camino. (3) La estrella, un centro unido
  a cinco hojas, no tiene ni ciclo ni camino: las hojas tienen grado
  $1$ y un camino admite a lo más dos como extremos.
\end{ejemplo}

\begin{ojo}
  Dos detalles de las diapositivas. Para el grafo (1) y para el
  camino $P_5$ (a, b, c, d, e) justifican que no hay ciclo ``porque no
  existe la arista $\set{d, a}$'' (o $\set{e, a}$). Eso solo descarta
  el ciclo
  que cerraría esa ruta en particular; podría haber otra ruta que sí
  cierre. El argumento que sí funciona es el del teorema
  \ref{teo:hamilton-grado}: a tiene grado $1$ y en un ciclo todo
  vértice usa dos aristas. El otro detalle: el texto llama a la
  estrella $K_{1,4}$ con ``cuatro hojas'', pero la figura dibuja cinco
  (a, b, c, d, e), así que es $K_{1,5}$. La conclusión no cambia; con
  tres hojas ya basta.
\end{ojo}

\subsubsection*{Una condición suficiente: Dirac}

Las condiciones de arriba solo descartan. La que garantiza un ciclo
pide muchas aristas en cada vértice.

\begin{teorema}[Dirac, 1952]
  \label{teo:hamilton-dirac}
  Si $G$ es simple, $n \ge 3$ y $\text{gr}(v) \ge n/2$ para todo
  $v$, entonces $G$ tiene ciclo hamiltoniano.
\end{teorema}

Conviene probar algo más fuerte, el teorema de Ore (1960): si
$\text{gr}(u) + \text{gr}(v) \ge n$ para cada par $u, v$ no
adyacente, $G$ tiene ciclo hamiltoniano. Dirac es el caso particular
en que cada sumando ya vale al menos $n/2$.

\begin{proof}[Prueba de Ore]
  Supóngase que hay un contraejemplo con $n$ vértices. Agregar
  aristas solo sube los grados, así que la condición se conserva. Se
  le agregan aristas, una por una, mientras haya alguna que no lo
  vuelva hamiltoniano. Se llega a un $G$ que cumple la condición, no
  es hamiltoniano, y se vuelve hamiltoniano al agregarle
  \textbf{cualquier} arista que le falte; como $K_n$ sí es
  hamiltoniano, a $G$ le falta alguna. Como $G \ne K_n$, hay $u, v$
  no adyacentes. El ciclo hamiltoniano de $G + \set{u, v}$ usa la
  arista nueva; sin ella es un camino hamiltoniano de $G$,
  \[
    u = v_1, v_2, \dots, v_n = v .
  \]
  Sean $S = \set{i : u \text{ es adyacente a } v_{i+1}}$ y
  $T = \set{i : v \text{ es adyacente a } v_i}$, ambos dentro de
  $\set{1, \dots, n-1}$. Cada vecino de $u$ aporta un índice a $S$ y
  cada vecino de $v$ uno a $T$, así que
  $\abs{S} + \abs{T} = \text{gr}(u) + \text{gr}(v) \ge n > n - 1$, y
  algún $i$ está en los dos. Para ese $i$,
  \[
    v_1, \dots, v_i, \; v_n, v_{n-1}, \dots, v_{i+1}, \; v_1
  \]
  es un ciclo hamiltoniano de $G$: va por el camino hasta $v_i$,
  salta a $v_n$ por la arista $\set{v_i, v}$, regresa por el camino
  hasta $v_{i+1}$ y cierra con $\set{v_{i+1}, u}$. Contradicción.
\end{proof}

\begin{ejemplo}[Suficiente, no necesaria]
  En el grafo del ejemplo \ref{ej:hamilton-pentagono}, $n = 5$ y el
  grado mínimo es $3 \ge 5/2$: Dirac garantiza el ciclo sin
  buscarlo. Las diapositivas usan ese grafo antes del teorema y no lo
  notan, pero es el contraste natural con el ciclo $C_5$: ahí cada
  vértice tiene grado $2 < 5/2$, Dirac no dice nada, y el grafo es él
  mismo un ciclo hamiltoniano. La condición de Dirac es suficiente,
  no necesaria.
\end{ejemplo}

% =====================================================================

\subsubsection*{Euler contra Hamilton}

Recordatorio de la sección \ref{sec:teoria-grafos}: un grafo conexo
(sin contar vértices aislados) tiene circuito euleriano si y solo si
todos sus grados son pares, y camino euleriano si y solo si tiene $0$
o $2$ vértices de grado impar. Decidirlo cuesta $\Tht{n + m}$: contar
grados impares y correr una BFS. Para Hamilton no se conoce ningún
algoritmo polinomial. Las dos preguntas tienen la misma forma,
recorrer todo sin repetir, y solo cambia qué no se repite. Tres
hechos explican la diferencia.

\textbf{La condición de Euler es local.} Como los vértices sí se
pueden repetir, cada vértice solo necesita emparejar las aristas por
las que se entra con las aristas por las que se sale, y eso se puede
exactamente cuando su grado es par. Lo único global es la conexidad,
y una BFS la revisa. Además los pedazos se pegan: dos circuitos sin
aristas en común que comparten un vértice se funden en uno, pasando
dos veces por ese vértice. Por eso Hierholzer nunca deshace una
decisión.

\textbf{La condición de Hamilton es global.} En cada vértice hay que
escoger cuáles dos de sus aristas usar, y lo que se escoge en uno
restringe a todos los demás. Los pedazos no se pegan: fundir dos
ciclos que comparten un vértice obliga a pasar dos veces por él, que
es justo lo prohibido. Una mala elección temprana no se repara
localmente; hay que regresar, y el árbol de búsqueda crece como
$(n-1)!$. Ninguna condición de grados cierra el hueco entre las
necesarias y la de Dirac: en medio quedan grafos como $C_5$,
hamiltoniano sin cumplir Dirac, y el de Petersen, que cumple todas
las necesarias y no es hamiltoniano.

\textbf{La prueba formal es la NP-completitud.} Un ciclo propuesto se
verifica con $n$ consultas a la matriz, así que el problema está en
NP; y cualquier problema de NP se reduce a él (Karp, 1972; Cormen,
teorema 34.13, lo hace desde cubierta de vértices). Un algoritmo
polinomial para Hamilton daría uno para todo NP, es decir $P = NP$.
Lo mejor que se conoce es exponencial: la programación dinámica del
agente viajero, más abajo, decide Hamilton en $\Tht{n^2 2^n}$.

\begin{idea}
  Euler se lee en los grados. Hamilton no se lee en nada local: un
  ciclo propuesto se verifica en tiempo lineal, pero encontrarlo,
  hasta donde se sabe, cuesta tiempo exponencial. Las dos propiedades
  son independientes, como muestra la tabla.
\end{idea}

\begin{center}
\small
\begin{tabular}{lcccccc}
  \toprule
  \textbf{Grafo} & $n$ & $m$ & \textbf{Impares} & \textbf{Euler}
    & \textbf{Camino H.} & \textbf{Ciclo H.} \\
  \midrule
  $C_5$ & 5 & 5 & 0 & circuito & sí & sí \\
  Ejemplo \ref{ej:hamilton-pentagono} & 5 & 8 & 4 & no & sí & sí \\
  Corbata & 5 & 6 & 0 & circuito & sí & no \\
  Camino $P_5$ & 5 & 4 & 2 & camino & sí & no \\
  Estrella $K_{1,5}$ & 6 & 5 & 6 & no & no & no \\
  Petersen & 10 & 15 & 10 & no & sí & no \\
  \bottomrule
\end{tabular}
\end{center}

La ``corbata'' son dos triángulos que comparten el vértice c: tiene
circuito euleriano, a, b, c, d, e, c, a, pero c es vértice de corte y
el teorema \ref{teo:hamilton-corte} descarta el ciclo hamiltoniano.

\begin{center}
\begin{tikzpicture}
  \node[vert] (a) at (-1.5, 0.9) {a};
  \node[vert] (b) at (-1.5, -0.9) {b};
  \node[vertAlerta] (c) at (0, 0) {c};
  \node[vert] (d) at (1.5, 0.9) {d};
  \node[vert] (e) at (1.5, -0.9) {e};
  \foreach \x/\y in {a/b, b/c, c/a, c/d, d/e, e/c}
    \draw[aristaArbol] (\x) -- (\y);
  \node[nota, align=center] at (0, -2.1)
    {corbata: euleriano sí,\\ hamiltoniano no (c es de corte)};
  \begin{scope}[xshift=6.8cm]
    \foreach \i in {0, ..., 4} {
      \node[vert] (o\i) at (90 + 72 * \i:1.6)
        {\the\numexpr\i+1\relax};
      \node[vert] (i\i) at (90 + 72 * \i:0.8)
        {\the\numexpr\i+6\relax};
    }
    \foreach \i/\j in {0/1, 1/2, 2/3, 3/4, 4/0}
      \draw[arista] (o\i) -- (o\j);
    \foreach \i/\j in {0/2, 2/4, 4/1, 1/3, 3/0}
      \draw[arista] (i\i) -- (i\j);
    \foreach \i in {0, ..., 4}
      \draw[arista] (o\i) -- (i\i);
    \node[nota, align=center] at (0, -2.1)
      {Petersen: $3$-regular, sin vértices de corte,\\
       sin ciclo hamiltoniano};
  \end{scope}
\end{tikzpicture}
\end{center}

El \textbf{grafo de Petersen} es el caso clásico. Tiene $n = 10$ y
$m = 15$, todos sus vértices tienen grado $3$, es conexo y quitar
cualquier vértice lo deja conexo: pasa los teoremas
\ref{teo:hamilton-grado} y \ref{teo:hamilton-corte}. Dirac pediría
grado $5$ y no aplica. No tiene ciclo hamiltoniano, aunque sí camino:
$1, 6, 8, 10, 5, 4, 9, 7, 2, 3$ con la numeración de la figura. Que
no tiene ciclo se puede probar a mano.

\begin{ejercicio}[Petersen no es hamiltoniano]
  \label{ej:hamilton-petersen}
  Llame $o_j$ al vértice exterior $j + 1$ e $i_j$ al interior
  $j + 6$, con $j = 0, \dots, 4$ e índices módulo $5$. Las aristas son
  las del pentágono, $o_j o_{j+1}$; las de la estrella,
  $i_j i_{j+2}$; y los rayos $o_j i_j$. Demuestre que no hay ciclo
  hamiltoniano.

  \textit{Solución.} Supóngase que $H$ es un ciclo hamiltoniano; tiene
  $10$ aristas. Cada vez que $H$ pasa del pentágono a la estrella usa
  un rayo, y como regresa a donde empezó, cruza un número par de
  veces; cero no puede ser, porque visita los dos lados. Usa $2$ o
  $4$ rayos.

  \paso{Dos rayos.} En $H$ cada vértice tiene grado $2$. Los dos
  exteriores con rayo usan una arista del pentágono y los otros tres
  usan dos: $8$ extremos, $4$ aristas. Cuatro aristas del pentágono
  son el pentágono menos una, un camino por los cinco exteriores
  cuyos extremos son los de la arista quitada: los dos con rayo son
  vecinos, $o_j$ y $o_{j+1}$. Por lo mismo, las aristas interiores de
  $H$ forman un camino por los cinco interiores que va de $i_j$ a
  $i_{j+1}$. Pero la estrella es el ciclo
  $i_0, i_2, i_4, i_1, i_3, i_0$, y en un ciclo de cinco vértices solo
  hay camino por todos entre dos vértices vecinos. $i_j$ e $i_{j+1}$
  no lo son.

  \paso{Cuatro rayos.} Falta uno; por la simetría de rotación, sea
  $o_0 i_0$. Entonces $o_0$ usa sus dos aristas del pentágono,
  $o_4 o_0$ y $o_0 o_1$, e $i_0$ sus dos de la estrella, $i_2 i_0$ e
  $i_0 i_3$. Cada uno de los otros ocho usa su rayo y una sola arista
  de su lado. $o_1$ y $o_4$ ya tienen la suya, así que la tercera
  arista exterior es $o_2 o_3$; $i_2$ e $i_3$ ya tienen la suya, así
  que la tercera interior es la única que queda, $i_4 i_1$. Pero
  entonces $o_1, o_0, o_4, i_4, i_1, o_1$ es un ciclo de cinco
  vértices dentro de $H$, que debía ser un solo ciclo de diez.
\end{ejercicio}

% =====================================================================

\subsubsection*{Búsqueda con retroceso}

Sin una caracterización, lo que queda es buscar. La búsqueda con
retroceso (\textit{backtracking}) hace crecer un camino desde un
vértice fijo, y cuando se atora deshace el último paso y prueba el
siguiente vecino.

\begin{algorithm}[H]
\caption{\textsc{Ciclo-Hamiltoniano}$(G)$: búsqueda con retroceso}
\label{alg:hamilton-retroceso}
\begin{algorithmic}[1]
\Procedure{Ciclo-Hamiltoniano}{$G$}
  \ForAll{$v \in V$}
    \State $\text{usado}[v] \gets$ falso
  \EndFor
  \State $s \gets$ cualquier vértice; \ $P \gets [s]$; \
         $\text{usado}[s] \gets$ verdadero
  \State \Return \textsc{Extender}$(G, P)$
\EndProcedure
\Statex
\Procedure{Extender}{$G, P$}
  \State $u \gets$ último vértice de $P$
  \If{$\abs{P} = n$}
    \State \Return $\set{u, P[1]} \in E$ \Comment{¿cierra el ciclo?}
  \EndIf
  \ForAll{$v$ adyacente a $u$ con \textbf{no} $\text{usado}[v]$}
    \State agregar $v$ al final de $P$; \
           $\text{usado}[v] \gets$ verdadero
    \If{\textsc{Extender}$(G, P)$}
      \State \Return verdadero
    \EndIf
    \State quitar $v$ de $P$; \ $\text{usado}[v] \gets$ falso
           \Comment{retroceder}
  \EndFor
  \State \Return falso
\EndProcedure
\end{algorithmic}
\end{algorithm}

Fijar el inicio $s$ no pierde nada: todo ciclo hamiltoniano pasa por
$s$. Para camino hamiltoniano sí hay que probar todos los inicios, o
agregar un vértice nuevo unido a todos los demás: $G$ tiene camino
hamiltoniano si y solo si $G$ más ese vértice tiene ciclo.

\textbf{Costo.} Con el arreglo $\text{usado}$, cada consulta es
$\Ord{1}$. Cada llamada a \textsc{Extender} corresponde a un camino
que empieza en $s$, distinto para cada llamada, y el grafo con más
caminos así es $K_n$:
\[
  \sum_{k=0}^{n-1} \frac{(n-1)!}{(n-1-k)!}
  = (n-1)! \sum_{j=0}^{n-1} \frac{1}{j!} \le e \cdot (n-1)! .
\]
Para $K_5$ son $65$, el mismo número del ojito sobre los caminos que
salen de un vértice, en la sección \ref{sec:grafos}: es aquí donde ese
factorial aparece de verdad. Cada llamada trabaja $\Ord{n}$ fuera de
sus llamadas recursivas, así que el total es
$\Ord{n \cdot (n-1)!} = \Ord{n!}$. Las podas de los teoremas
\ref{teo:hamilton-grado} y
\ref{teo:hamilton-corte} se revisan antes en $\Tht{n + m}$ y en la
práctica ahorran mucho, pero no cambian el peor caso.

\begin{ejemplo}[Retroceso sobre el grafo pesado, sin pesos]
  \label{ej:hamilton-traza}
  Se busca un ciclo en el grafo pesado de la sección \ref{sec:grafos}
  ignorando los pesos, desde A y con los vecinos en orden
  alfabético. Cada nodo del árbol de búsqueda es el vértice que se
  agrega; las ramas se exploran de arriba hacia abajo.
  \begin{center}
  \begin{tikzpicture}[x=1.3cm, y=0.95cm]
    \node[vertOk] (n0) at (0, 0) {A};
    \node[vertOk] (n1) at (1, 0) {B};
    \node[vertOk] (n2) at (2, 0) {C};
    \node[vert] (n3) at (3, 1) {D};
    \node[vertOk] (n4) at (3, -1) {E};
    \node[vert] (n5) at (4, 1.5) {E};
    \node[vert] (n6) at (4, 0.5) {F};
    \node[vert] (n7) at (4, -0.5) {D};
    \node[vertOk] (n8) at (4, -1.5) {F};
    \node[vertAlerta] (n9) at (5, 1.5) {F};
    \node[vertAlerta] (n10) at (5, 0.5) {E};
    \node[vertAlerta] (n11) at (5, -0.5) {F};
    \node[vertOk] (n12) at (5, -1.5) {D};
    \foreach \a/\b in {n2/n3, n3/n5, n3/n6, n4/n7, n5/n9, n6/n10,
                       n7/n11}
      \draw[arista] (\a) -- (\b);
    \foreach \a/\b in {n0/n1, n1/n2, n2/n4, n4/n8, n8/n12}
      \draw[aristaArbol] (\a) -- (\b);
    \foreach \y/\t in {1.5/{A, B, C, D, E, F: F--A no es arista},
                       0.5/{A, B, C, D, F, E: E--A no es arista},
                       -0.5/{A, B, C, E, D, F: F--A no es arista},
                       -1.5/{A, B, C, E, F, D: D--A sí, cierra}}
      \node[nota, anchor=west] at (5.35, \y) {\t};
  \end{tikzpicture}
  \end{center}
  De A el primer vecino libre es B, de B es C, y de C se prueba
  primero D. Las dos ramas bajo D llegan a usar los seis vértices sin
  poder cerrar, porque A solo es adyacente a B, C y D. Se regresa a C
  y se prueba E: bajo E, la rama de D falla igual y la de F termina en
  D, que sí es vecino de A. El ciclo es A, B, C, E, F, D, A, tras $13$
  llamadas a \textsc{Extender}.
\end{ejemplo}

% =====================================================================

\subsubsection*{El agente viajero}

\begin{definicion}[Problema del agente viajero]
  Dadas $n$ ciudades y el costo $w(u, v) \ge 0$ de ir de una a otra,
  el \textbf{problema del agente viajero} (\textit{traveling salesman
  problem}, TSP) pide la ruta que visita cada ciudad una vez, regresa
  al inicio y tiene el menor costo total: un ciclo hamiltoniano de
  costo mínimo en un grafo completo y pesado.
\end{definicion}

Si el grafo no es completo, una arista que falta vale $w = \infty$.
Es al menos tan difícil como Hamilton: dado $G$, se arma el completo
con peso $0$ en las aristas de $G$ y $1$ en las demás, y $G$ es
hamiltoniano si y solo si hay una ruta de costo $0$. Por eso la
versión de decisión también es NP-completa (Cormen, teorema 34.14).

\textbf{Fuerza bruta.} Se fija la ciudad de salida y se permutan las
otras $n - 1$. Una ruta y su reversa cuestan lo mismo, así que hay
$(n-1)!/2$ rutas distintas, cada una evaluada en $\Tht{n}$: en total
$\Tht{n!}$. Para $n = 5$ son $12$ rutas; para $n = 20$, unas
$6.1 \times 10^{16}$.

\textbf{Programación dinámica (Bellman; Held y Karp, 1962).} Con las
ciudades numeradas $1, \dots, n$ y salida en $1$, sea $g[S][j]$ el
costo del camino más barato que sale de $1$, visita exactamente los
vértices de $S \subseteq \set{2, \dots, n}$ y termina en $j \in S$.
El penúltimo vértice de ese camino es algún $i \in S$ distinto de
$j$, y lo que viene antes es un camino de $1$ a $i$ por
$S \setminus \set{j}$, que debe ser el más barato. Eso da
\[
  g[\set{j}][j] = w(1, j) , \qquad
  g[S][j] = \min_{i \in S \setminus \set{j}}
            \bigl( g[S \setminus \set{j}][i] + w(i, j) \bigr) ,
\]
y la ruta óptima cuesta
$\min_{j} \bigl( g[\set{2, \dots, n}][j] + w(j, 1) \bigr)$. Es la
subestructura óptima de la clase de programación dinámica: la mejor
manera de terminar depende de qué falta y de dónde se está, no del
orden en que se visitó $S$. Esa observación es la que ahorra: hay
$(n-1)!$ órdenes, pero solo $2^{n-1}$ conjuntos.

\textbf{Costo.} Recorrer los conjuntos por tamaño creciente, como en
el algoritmo de abajo, garantiza que $g[S \setminus \set{j}][\cdot]$
ya está calculado. Hay $(n-1)\, 2^{n-2}$ pares $(S, j)$ con
$j \in S$, y cada uno toma un mínimo sobre a lo más $n - 2$
términos: $\Tht{n^2 2^n}$ en tiempo y $\Tht{n 2^n}$ en memoria,
porque la tabla se guarda completa. Con $S$ guardado como máscara de
bits, cada consulta es $\Ord{1}$. Para
$n = 20$ son unos $4.2 \times 10^8$ pasos contra $6.1 \times 10^{16}$
rutas. Sigue siendo exponencial, como se espera de un problema
NP-difícil, y en la práctica la memoria es el cuello de botella.

\begin{algorithm}[H]
\caption{\textsc{Agente-Viajero}$(w, n)$: programación dinámica de
         Held y Karp}
\label{alg:hamilton-held-karp}
\begin{algorithmic}[1]
\Procedure{Agente-Viajero}{$w, n$}
  \State $g[\set{j}][j] \gets w(1, j)$ para $j = 2, \dots, n$
  \For{$k \gets 2$ \Hasta $n - 1$} \Comment{por tamaño creciente}
    \ForAll{$S \subseteq \set{2, \dots, n}$ con $\abs{S} = k$}
      \ForAll{$j \in S$}
        \State $g[S][j] \gets \min_{i \in S \setminus \set{j}}
               \bigl( g[S \setminus \set{j}][i] + w(i, j) \bigr)$
      \EndFor
    \EndFor
  \EndFor
  \State \Return $\min_{2 \le j \le n}
         \bigl( g[\set{2, \dots, n}][j] + w(j, 1) \bigr)$
\EndProcedure
\end{algorithmic}
\end{algorithm}

\begin{ejemplo}[El agente viajero de las diapositivas]
  \label{ej:hamilton-tsp}
  Cinco ciudades con $w(A,B) = 4$, $w(A,C) = 7$, $w(A,E) = 6$,
  $w(B,C) = 3$, $w(B,D) = 8$, $w(C,D) = 5$, $w(C,E) = 4$ y
  $w(D,E) = 2$.
  \begin{center}
  \begin{tikzpicture}[x=1.1cm, y=1.1cm]
    \node[vert] (A) at (0, 1.6) {A};
    \node[vert] (B) at (1.9, 0.9) {B};
    \node[vert] (C) at (1.9, -1.0) {C};
    \node[vert] (D) at (0, -1.7) {D};
    \node[vert] (E) at (-2.0, 0) {E};
    \draw[aristaArbol] (A) -- node[peso] {4} (B);
    \draw[aristaArbol] (B) -- node[peso] {3} (C);
    \draw[aristaArbol] (C) -- node[peso] {5} (D);
    \draw[aristaArbol] (D) -- node[peso] {2} (E);
    \draw[aristaArbol] (E) -- node[peso] {6} (A);
    \draw[arista] (A) -- node[peso, pos=0.25] {7} (C);
    \draw[arista] (B) -- node[peso, pos=0.75] {8} (D);
    \draw[arista] (E) -- node[peso, pos=0.3] {4} (C);
    \node[nota] at (0, -2.45) {óptimo: A, B, C, D, E, A, con costo 20};
  \end{tikzpicture}
  \end{center}
  Faltan A--D y B--E, que valen $\infty$. De las $12$ rutas de $K_5$
  solo existen cuatro:
  \begin{center}
  \small
  \begin{tabular}{ll}
    \toprule
    \textbf{Ruta} & \textbf{Costo} \\
    \midrule
    A, B, C, D, E, A & $4 + 3 + 5 + 2 + 6 = 20$ \\
    A, B, D, E, C, A & $4 + 8 + 2 + 4 + 7 = 25$ \\
    A, C, B, D, E, A & $7 + 3 + 8 + 2 + 6 = 26$ \\
    A, B, D, C, E, A & $4 + 8 + 5 + 4 + 6 = 27$ \\
    \bottomrule
  \end{tabular}
  \end{center}
  La tabla de Held y Karp, con salida en A, da lo mismo sin enumerar
  rutas. Cada renglón es un conjunto $S$ y cada columna el vértice
  final $j$; un guion indica $j \notin S$.
  \begin{center}
  \small
  \begin{tabular}[t]{lcccc}
    \toprule
    $S$ & $j = B$ & $j = C$ & $j = D$ & $j = E$ \\
    \midrule
    $\set{B}$ & $4$ & -- & -- & -- \\
    $\set{C}$ & -- & $7$ & -- & -- \\
    $\set{D}$ & -- & -- & $\infty$ & -- \\
    $\set{E}$ & -- & -- & -- & $6$ \\
    \addlinespace
    $\set{B, C}$ & $10$ & $7$ & -- & -- \\
    $\set{B, D}$ & $\infty$ & -- & $12$ & -- \\
    $\set{B, E}$ & $\infty$ & -- & -- & $\infty$ \\
    $\set{C, D}$ & -- & $\infty$ & $12$ & -- \\
    $\set{C, E}$ & -- & $10$ & -- & $11$ \\
    $\set{D, E}$ & -- & -- & $8$ & $\infty$ \\
    \bottomrule
  \end{tabular}
  \qquad
  \begin{tabular}[t]{lcccc}
    \toprule
    $S$ & $j = B$ & $j = C$ & $j = D$ & $j = E$ \\
    \midrule
    $\set{B, C, D}$ & $20$ & $17$ & $12$ & -- \\
    $\set{B, C, E}$ & $13$ & $\infty$ & -- & $11$ \\
    $\set{B, D, E}$ & $16$ & -- & $\infty$ & $14$ \\
    $\set{C, D, E}$ & -- & $13$ & $13$ & $14$ \\
    \addlinespace
    $\set{B, C, D, E}$ & $16$ & $18$ & $13$ & $14$ \\
    \midrule
    $+\,w(j, A)$ & $20$ & $25$ & $\infty$ & $20$ \\
    \bottomrule
  \end{tabular}
  \end{center}
  Un par de entradas a mano: $g[\set{B, C}][C] = g[\set{B}][B] +
  w(B, C) = 4 + 3 = 7$, porque la única opción es A, B, C, y
  $g[\set{B, C}][B] = g[\set{C}][C] + w(C, B) = 7 + 3 = 10$. Y
  $g[\set{B, C, D, E}][E] = 14$ es
  A, B, C, D, E, con $4 + 3 + 5 + 2$; cerrar con $w(E, A) = 6$ da
  $20$. Por la columna de B también sale $20$: es la misma ruta al
  revés. Por D sale $\infty$ porque no hay arista D--A.
\end{ejemplo}

\begin{ojo}
  El grafo de las diapositivas no es completo y la diapositiva no da
  la respuesta. En el TSP el grafo se supone completo; una arista que
  falta se modela con costo $\infty$, y entonces una ruta ``existe''
  solo si su costo es finito. Ojo también con la frase ``la ruta más
  corta'': no es el camino más corto entre dos ciudades de las
  secciones \ref{sec:caminos-fuente} a \ref{sec:todos-pares}, que se
  resuelve en tiempo polinomial, sino un ciclo que toca a todas.
\end{ojo}

% =====================================================================

\subsection*{Grafos bipartitos}

\begin{definicion}[Grafo bipartito]
  \label{def:bipartito}
  Un grafo no dirigido $G = (V, E)$ es \textbf{bipartito} si $V$ se
  parte en dos conjuntos disjuntos, $V = V_1 \cup V_2$, de modo que
  toda arista tiene un extremo en $V_1$ y el otro en $V_2$. El
  \textbf{bipartito completo} $K_{p,q}$ tiene $\abs{V_1} = p$,
  $\abs{V_2} = q$ y cada vértice de $V_1$ unido con todos los de
  $V_2$: tiene $pq$ aristas.
\end{definicion}

Modela cualquier relación entre dos tipos de objetos: estudiantes y
proyectos, trabajadores y tareas, usuarios y películas. Asignar a
cada estudiante un proyecto distinto es el emparejamiento de la
sección \ref{sec:emparejamiento}. Cormen llama a los lados $L$ y $R$.

\begin{center}
\begin{tikzpicture}
  \foreach \i/\y in {1/-0.6, 2/-1.6}
    \node[vert] (u\i) at (0, \y) {$u_{\i}$};
  \foreach \j/\y in {1/-0.1, 2/-1.1, 3/-2.1}
    \node[vertOk] (v\j) at (2.6, \y) {$v_{\j}$};
  \foreach \i in {1, 2}
    \foreach \j in {1, 2, 3}
      \draw[arista] (u\i) -- (v\j);
  \node[nota] at (1.3, -2.8) {$K_{2,3}$: $2 \cdot 3 = 6$ aristas};
  \begin{scope}[xshift=5.6cm]
    \foreach \i/\y in {1/-0.1, 2/-1.1, 3/-2.1} {
      \node[vert] (p\i) at (0, \y) {$u_{\i}$};
      \node[vertOk] (q\i) at (2.6, \y) {$v_{\i}$};
    }
    \foreach \i in {1, 2, 3}
      \foreach \j in {1, 2, 3}
        \draw[arista] (p\i) -- (q\j);
    \node[nota] at (1.3, -2.8) {$K_{3,3}$: $3 \cdot 3 = 9$ aristas};
  \end{scope}
\end{tikzpicture}
\end{center}

\begin{observacion}
  Las diapositivas y Rosen escriben $K_{m,n}$. Aquí $n = \abs{V}$ y
  $m = \abs{E}$, y con esa lectura ``$K_{m,n}$ tiene $mn$ aristas''
  diría $m = mn$. Por eso los lados se llaman $p$ y $q$.
\end{observacion}

\subsubsection*{La caracterización por ciclos impares}

Un ciclo en un bipartito va de un lado al otro en cada paso, así que
para volver al lado de salida necesita un número par de pasos. Lo
notable es que eso basta: lo único que impide partir un grafo en dos
lados es un ciclo impar.

\begin{teorema}[Caracterización de los bipartitos]
  \label{teo:bipartito-ciclos}
  Un grafo no dirigido es bipartito si y solo si no contiene ciclos
  de longitud impar.
\end{teorema}

\begin{proof}
  ($\Rightarrow$) Sean $V_1, V_2$ los lados y
  $v_0, v_1, \dots, v_k = v_0$ un ciclo. Cada arista cambia de lado,
  así que $v_i$ está en el lado de $v_0$ exactamente cuando $i$ es
  par. Como $v_k = v_0$, $k$ es par.

  ($\Leftarrow$) Basta partir cada componente conexa por separado: no
  hay aristas entre componentes, así que la unión de las particiones
  sirve. Sea entonces $G$ conexo y sin ciclos impares. Se corre BFS
  desde cualquier $s$ y se pone
  \[
    V_1 = \set{v : d[v] \text{ es par}} , \qquad
    V_2 = \set{v : d[v] \text{ es impar}} .
  \]
  Falta ver que ninguna arista queda dentro de un lado. Supóngase que
  $\set{u, v} \in E$ con $d[u]$ y $d[v]$ de la misma paridad.

  \paso{Primero, $d[u] = d[v]$.} En BFS, $d[v] = \delta(s, v)$, el
  número de aristas del camino más corto (sección \ref{sec:grafos};
  la prueba está en la \ref{sec:busquedas}). Un camino más corto a $u$
  seguido de la arista $\set{u, v}$ llega a $v$, así que
  $\delta(s, v) \le \delta(s, u) + 1$, y por simetría
  $\abs{d[u] - d[v]} \le 1$. Con la misma paridad, la diferencia es
  $0$.

  \paso{Luego, el ciclo.} Sea $w$ el ancestro común más profundo de
  $u$ y $v$ en el árbol BFS; existe, porque $s$ es ancestro de todos.
  Cada arista del árbol baja exactamente un nivel, así que el camino
  del árbol de $w$ a $u$ mide $d[u] - d[w]$, y el de $w$ a $v$ mide
  $d[v] - d[w] = d[u] - d[w]$. Por ser $w$ el más profundo, los dos
  caminos solo comparten a $w$; además $w \ne u$, porque $u$ no puede
  ser ancestro de $v$ estando en el mismo nivel. La arista
  $\set{u, v}$ no es del árbol, porque une dos vértices del mismo
  nivel. Los dos caminos y la arista forman entonces un ciclo de
  longitud
  \[
    2\bigl(d[u] - d[w]\bigr) + 1 ,
  \]
  impar. Contradicción.
\end{proof}

\begin{observacion}
  Da igual hablar de ciclos impares o de caminatas cerradas impares:
  una caminata cerrada de longitud impar siempre contiene un ciclo
  impar. Si la caminata repite un vértice, además de cerrar en el
  inicial, se parte ahí en dos
  caminatas cerradas más cortas cuyas longitudes suman la original;
  una de las dos es impar, y se repite el argumento hasta llegar a
  una que no repite vértices.
\end{observacion}

\subsubsection*{El algoritmo: dos lados con BFS}

La prueba es un algoritmo. Se corre BFS, se asigna a cada vértice el
lado que dice la paridad de su distancia, y se revisa cada arista. Si
ninguna une dos vértices del mismo lado, los lados son la partición;
si alguna lo hace, se reconstruye el ciclo impar subiendo por los
padres desde sus dos extremos, como en la prueba. Para no chocar con
$\text{color}[v]$ de la sección \ref{sec:busquedas}, que es blanco,
gris o negro, el lado se guarda en $\text{lado}[v] \in \set{0, 1}$.
Las secciones \ref{sec:busquedas} (el problema de los luchadores) y
\ref{sec:emparejamiento} (emparejamiento) usan este algoritmo tal
cual.

Primero la pieza que reconstruye el testigo. Se llama solo en un
conflicto, cuando $u$ y $v$ son adyacentes y están en el mismo
lado, y sube por $\pi$ desde los dos hasta que se encuentran.

\begin{algorithm}[H]
\caption{\textsc{Ciclo-Impar}$(u, v)$: el testigo de que no es
         bipartito}
\label{alg:bipartito-ciclo-impar}
\begin{algorithmic}[1]
\Procedure{Ciclo-Impar}{$u, v$} \Comment{al llamarse, $d[u] = d[v]$}
  \State $x \gets u$, \quad $y \gets v$, \quad $P \gets [u]$, \quad
         $R \gets [v]$
  \While{$x \neq y$} \Comment{suben a la par}
    \State $x \gets \pi[x]$; \ agregar $x$ al final de $P$
    \State $y \gets \pi[y]$; \ agregar $y$ al final de $R$
  \EndWhile
  \State \Return $P$ seguida de $R$ al revés, sin repetir $x$
\EndProcedure
\end{algorithmic}
\end{algorithm}

El procedimiento principal es la BFS de la clase con un lado por
vértice.

\begin{algorithm}[H]
\caption{\textsc{Bipartito}$(G)$: dos lados con BFS}
\label{alg:bipartito-bfs}
\begin{algorithmic}[1]
\Procedure{Bipartito}{$G$}
  \ForAll{$v \in V$}
    \State $\text{lado}[v] \gets$ nulo, \quad $\pi[v] \gets$ nulo
  \EndFor
  \ForAll{$s \in V$} \Comment{una BFS por componente}
    \If{$\text{lado}[s] =$ nulo}
      \State $\text{lado}[s] \gets 0$, \quad $d[s] \gets 0$
      \State $Q \gets$ cola vacía; \ \textsc{Encolar}$(Q, s)$
      \While{$Q$ no está vacía}
        \State $u \gets \textsc{Desencolar}(Q)$
        \ForAll{$v$ en la lista de adyacencia de $u$}
          \If{$\text{lado}[v] =$ nulo}
            \State $\text{lado}[v] \gets 1 - \text{lado}[u]$, \quad
                   $d[v] \gets d[u] + 1$, \quad $\pi[v] \gets u$
            \State \textsc{Encolar}$(Q, v)$
          \ElsIf{$\text{lado}[v] = \text{lado}[u]$}
            \State \Return falso, \textsc{Ciclo-Impar}$(u, v)$
          \EndIf
        \EndFor
      \EndWhile
    \EndIf
  \EndFor
  \State \Return verdadero,
         $V_1 = \set{v : \text{lado}[v] = 0}$,
         $V_2 = \set{v : \text{lado}[v] = 1}$
\EndProcedure
\end{algorithmic}
\end{algorithm}

Detalles que no son obvios:
\begin{itemize}[itemsep=0.1em]
  \item El ciclo externo arranca una BFS en cada componente. Los lados
        de componentes distintas son independientes, y la partición
        final es la unión.
  \item El lado es la paridad de $d[v]$. Para decidir no hace falta
        $d$; se guarda para que \textsc{Ciclo-Impar} sepa que los dos
        extremos están en el mismo nivel y puedan subir a la par.
  \item A diferencia de la detección de ciclos de la sección
        \ref{sec:grafos}, no hace falta excluir al padre: el padre
        está del otro lado y nunca dispara la condición. Una arista
        fuera del árbol entre niveles consecutivos cierra un ciclo
        \textbf{par} y es inofensiva; solo una arista dentro de un
        mismo nivel delata un ciclo impar.
  \item \textsc{Ciclo-Impar} devuelve $u, \dots, w, \dots, v$, donde
        $w$ es el ancestro común; la arista $\set{v, u}$ lo cierra.
\end{itemize}

\textbf{Costo.} La inicialización es $\Tht{n}$. Cada vértice recibe
lado una sola vez, justo antes de encolarse, y nunca cambia: $n$
operaciones de cola de $\Ord{1}$. La lista de cada vértice se recorre
una vez, cuando sale de la cola; las listas suman $2m$ entradas, cada
una con una comparación de $\Ord{1}$. El ciclo externo agrega
$\Tht{n}$. Total, $\Tht{n + m}$ con listas, y $\Tht{n^2}$ con matriz.
\textsc{Ciclo-Impar} sube un nivel por iteración: $\Ord{n}$. Decidir
si un grafo es bipartito no cuesta más que recorrerlo.

\begin{idea}
  El algoritmo siempre entrega un testigo que se revisa a mano: la
  partición, si es bipartito, o un ciclo impar, si no. Es la
  diferencia con Hamilton, donde un ``no'' no viene con ningún
  certificado corto conocido.
\end{idea}

\begin{ejemplo}[El grafo $G$ de la clase no es bipartito]
  \label{ej:bipartito-traza-G}
  Se corre \textsc{Bipartito} sobre el grafo $G$ de la sección
  \ref{sec:grafos}, con los vecinos en orden alfabético. En la
  figura, azul es el lado $0$ y verde el lado $1$.
  \begin{center}
  \small
  \begin{tabular}{cclc}
    \toprule
    \textbf{Paso} & \textbf{Sale} & \textbf{Asigna o revisa}
      & \textbf{Cola después} \\
    \midrule
    0 & & A: 0 & [A] \\
    1 & A & B: 1, C: 1 & [B, C] \\
    2 & B & A bien; D: 0, E: 0 & [C, D, E] \\
    3 & C & A bien; F: 0 & [D, E, F] \\
    4 & D & B bien; G: 1 & [E, F, G] \\
    5 & E & B bien; F tiene 0, \textbf{igual} que E & conflicto \\
    \bottomrule
  \end{tabular}
  \end{center}
  \begin{center}
  \begin{tikzpicture}[scale=0.72, transform shape,
    sinLado/.style={vert, draw=gray!60, fill=white}]
    \foreach \pre/\x in {p/0, q/6.4, r/12.8} {
      \begin{scope}[xshift=\x cm, name prefix=\pre-]
        \coordinate (A) at (0, 0);
        \coordinate (B) at (-1.6, -1.2);
        \coordinate (C) at (1.6, -1.2);
        \coordinate (D) at (-2.6, -2.4);
        \coordinate (E) at (-0.6, -2.4);
        \coordinate (F) at (1.6, -2.4);
        \coordinate (G) at (-2.6, -3.6);
      \end{scope}
    }
    % instantanea 1: sale A
    \foreach \a/\b in {B/D, B/E, C/F, E/F, D/G}
      \draw[arista] (p-\a) -- (p-\b);
    \foreach \a/\b in {A/B, A/C}
      \draw[aristaArbol] (p-\a) -- (p-\b);
    \node[vert] at (p-A) {A};
    \foreach \v in {B, C} \node[vertOk] at (p-\v) {\v};
    \foreach \v in {D, E, F, G} \node[sinLado] at (p-\v) {\v};
    \node[font=\small] at (0, -4.5) {1. sale A};
    % instantanea 2: sale D
    \draw[arista] (q-E) -- (q-F);
    \foreach \a/\b in {A/B, A/C, B/D, B/E, C/F, D/G}
      \draw[aristaArbol] (q-\a) -- (q-\b);
    \foreach \v in {A, D, E, F} \node[vert] at (q-\v) {\v};
    \foreach \v in {B, C, G} \node[vertOk] at (q-\v) {\v};
    \node[font=\small] at (6.4, -4.5) {4. sale D};
    % instantanea 3: sale E, conflicto en E--F
    \foreach \a/\b in {B/D, D/G}
      \draw[aristaArbol] (r-\a) -- (r-\b);
    \foreach \a/\b in {A/B, A/C, B/E, C/F}
      \draw[aristaAlerta] (r-\a) -- (r-\b);
    \draw[aristaAlerta, dashed] (r-E) -- (r-F);
    \foreach \v in {A, D, E, F} \node[vert] at (r-\v) {\v};
    \foreach \v in {B, C, G} \node[vertOk] at (r-\v) {\v};
    \node[font=\small] at (12.8, -4.5) {5. sale E: E y F en el lado 0};
    \foreach \y/\d in {0/0, -1.2/1, -2.4/2, -3.6/3}
      \node[nota, anchor=west] at (15.2, \y) {$d = \d$};
  \end{tikzpicture}
  \end{center}
  En el paso 5 sale E, del lado $0$, y su vecino F también está en el
  lado $0$: $d[E] = d[F] = 2$, la arista E--F queda dentro de un
  nivel. \textsc{Ciclo-Impar}$(E, F)$ sube $E \to B \to A$ y
  $F \to C \to A$, se encuentra en A y devuelve E, B, A, C, F; con la
  arista F--E es el ciclo A, B, E, F, C, A, de longitud $5$, en rojo.
  Es el mismo ciclo que la clase encontró con BFS en ``Encontrar
  ciclos'', y ahora se sabe además que es impar. $G$ no es bipartito.
\end{ejemplo}

\begin{ejemplo}[Un bipartito disfrazado]
  \label{ej:bipartito-traza-c6}
  El ejemplo 6 de bipartitos de las diapositivas tiene
  $V = \set{a, b, c, d, e, f}$ y
  $E = \set{ae, af, bd, bf, cd, ce}$.

  \begin{minipage}[c]{0.47\linewidth}
  \centering
  \small
  \begin{tabular}{cclc}
    \toprule
    \textbf{Paso} & \textbf{Sale} & \textbf{Asigna o revisa}
      & \textbf{Cola} \\
    \midrule
    0 & & a: 0 & [a] \\
    1 & a & e: 1, f: 1 & [e, f] \\
    2 & e & a bien; c: 0 & [f, c] \\
    3 & f & a bien; b: 0 & [c, b] \\
    4 & c & d: 1; e bien & [b, d] \\
    5 & b & d bien, f bien & [d] \\
    6 & d & b bien, c bien & [\,] \\
    \bottomrule
  \end{tabular}
  \end{minipage}\hfill
  \begin{minipage}[c]{0.5\linewidth}
  \centering
  \begin{tikzpicture}
    \foreach \v/\x in {a/0, b/1.4, c/2.8}
      \node[vert] (\v) at (\x, 1.6) {\v};
    \foreach \v/\x in {d/0, e/1.4, f/2.8}
      \node[vertOk] (\v) at (\x, 0) {\v};
    \foreach \x/\y in {a/e, a/f, b/d, b/f, c/d, c/e}
      \draw[arista] (\x) -- (\y);
    \node[nota] at (1.4, -0.7) {como en la diapositiva};
    \begin{scope}[xshift=5.4cm, yshift=0.8cm]
      \foreach \v/\ang/\s in {a/90/vert, e/30/vertOk, c/-30/vert,
                              d/-90/vertOk, b/-150/vert, f/150/vertOk}
        \node[\s] (\v) at (\ang:1.15) {\v};
      \foreach \x/\y in {a/e, e/c, c/d, d/b, b/f, f/a}
        \draw[arista] (\x) -- (\y);
      \node[nota] at (0, -1.5) {el mismo grafo: $C_6$};
    \end{scope}
  \end{tikzpicture}
  \end{minipage}

  Nunca hay conflicto: $V_1 = \set{a, b, c}$, $V_2 = \set{d, e, f}$.
  El dibujo de la diapositiva esconde que el grafo es el hexágono
  $C_6$: a, e, c, d, b, f, a. Son el mismo grafo dibujado de dos
  maneras, isomorfos en el sentido de la sección
  \ref{sec:teoria-grafos}. También es $K_{3,3}$ sin las aristas a--d,
  b--e y c--f.
\end{ejemplo}

\begin{observacion}
  En un bipartito conexo la partición es única salvo intercambiar los
  lados: el lado de $s$ fija el de todos los demás. Con $c$
  componentes hay $2^c$ maneras de asignar los lados, porque cada
  componente se puede voltear por separado. Y ser bipartito ignora el
  sentido de las aristas: en un dirigido se pregunta por el no
  dirigido que queda al borrar las flechas.
\end{observacion}

\subsubsection*{Bipartitos y Hamilton}

El mismo argumento de alternancia da una poda para Hamilton.

\begin{corolario}
  \label{cor:bipartito-hamilton}
  Sea $G$ bipartito con lados $V_1$ y $V_2$. Si $G$ tiene ciclo
  hamiltoniano, $\abs{V_1} = \abs{V_2}$. Si tiene camino
  hamiltoniano, $\bigl| \abs{V_1} - \abs{V_2} \bigr| \le 1$.
\end{corolario}

\begin{proof}
  El ciclo alterna de lado, así que tiene tantos vértices de un lado
  como del otro, y pasa por todos. El camino también alterna: con un
  número par de vértices los dos lados quedan iguales, y con uno impar
  empieza y termina en el mismo lado, que lleva uno de más.
\end{proof}

Esto liquida de un golpe la estrella $K_{1,5}$ (lados de $1$ y $5$:
ni camino) y muestra que $K_{2,3}$ no es hamiltoniano aunque tenga el
camino $v_1, u_1, v_2, u_2, v_3$. $K_{3,3}$ sí lo es: tiene seis
ciclos hamiltonianos distintos.

\begin{ejercicio}[Determine si son bipartitos]
  \label{ej:bipartito-determine}
  (a) El grafo pesado de la sección \ref{sec:grafos}, sin pesos.
  (b) El grafo de Petersen. (c) Cualquier árbol. (d) El ciclo $C_k$.
  (e) El grafo con $V_1 = \set{a, b, c, d}$, $V_2 = \set{1, 2, 3}$ y
  $E = \set{a1, a2, b2, c1, c3, d3}$: ¿tiene camino o ciclo
  hamiltoniano?

  \textit{Solución.} (a) No: A, B, C es un triángulo.
  \textsc{Bipartito} desde A pone B y C en el lado $1$, y al sacar B
  encuentra que su vecino C está en su mismo lado;
  \textsc{Ciclo-Impar}$(B, C)$ devuelve B, A, C. (b) No: el pentágono
  exterior es un ciclo de longitud $5$. (c) Sí: no tiene ciclos, así
  que tampoco impares. En el algoritmo se ve directo: todas las
  aristas son del árbol BFS y unen niveles consecutivos. (d) Si y
  solo si $k$ es par. (e) Por construcción es bipartito, con lados de
  $4$ y $3$: el corolario descarta el ciclo pero no el camino, y el
  camino existe: b, 2, a, 1, c, 3, d.
\end{ejercicio}

% =====================================================================

\subsection*{Cliques}

\begin{definicion}[Clique]
  \label{def:clique}
  Una \textbf{clique} de $G$ es un conjunto $C \subseteq V$ de
  vértices adyacentes dos a dos: si $u, v \in C$ y $u \ne v$,
  entonces $\set{u, v} \in E$. Su tamaño es $\abs{C}$. El
  \textbf{número de clique} $\omega(G)$ es el tamaño de la clique más
  grande. Una clique es \textbf{maximal} si no se le puede agregar
  ningún vértice, porque ninguno fuera de ella es adyacente a todos
  los de adentro; es \textbf{máxima} si tiene tamaño $\omega(G)$.
\end{definicion}

El \textbf{subgrafo inducido} $G[C]$ tiene los vértices de $C$ y todas
las aristas de $G$ entre ellos. $C$ es clique exactamente cuando
$G[C]$ es isomorfo al completo $K_{\abs{C}}$, que es como lo escriben
las diapositivas: $G[C] \cong K_k$. Toda clique máxima es maximal; al
revés no. En redes sociales una clique es un grupo en el que todos se
conocen, y buscarlas es la forma más estricta de detectar grupos
densos.

\begin{ejemplo}[Las cliques del grafo pesado]
  \label{ej:clique-pesado}
  En el grafo pesado de la sección \ref{sec:grafos}, sin pesos, los
  cuatro vértices A, B, C, D están unidos todos con todos.
  \begin{center}
  \begin{tikzpicture}
    \node[vertOk] (A) at (0, 0) {A};
    \node[vertOk] (B) at (2, 1.4) {B};
    \node[vertOk] (C) at (2, -1.4) {C};
    \node[vertOk] (D) at (4.5, 1.4) {D};
    \node[vertMarca] (E) at (4.5, -1.4) {E};
    \node[vertMarca] (F) at (6.5, 0) {F};
    \draw[arista] (C) -- (E);
    \foreach \a/\b in {D/E, D/F, E/F}
      \draw[aristaResalte] (\a) -- (\b);
    \foreach \a/\b in {A/B, A/C, B/C, B/D, C/D}
      \draw[aristaArbol] (\a) -- (\b);
    \draw[aristaArbol] (A) to[out=75, in=150] (D);
    \node[nota, align=left, anchor=west] at (7.2, 0)
      {azul: $\set{A, B, C, D}$, máxima\\
       naranja: $\set{D, E, F}$, maximal\\
       no máxima};
  \end{tikzpicture}
  \end{center}
  $\set{A, B, C, D}$ es una clique de tamaño $4$, así que
  $\omega(G) \ge 4$. Para la igualdad hay que descartar una de tamaño
  $5$: una $K_5$ necesita cinco vértices de grado al menos $4$, y los
  grados son $3, 3, 4, 5, 3, 2$ para A a F; solo C y D llegan. Así
  que $\omega(G) = 4$. Las cliques maximales son tres:
  $\set{A, B, C, D}$, $\set{C, D, E}$ y $\set{D, E, F}$. La última es
  maximal, porque ningún otro vértice es adyacente a D, E y F a la
  vez, pero no máxima.
\end{ejemplo}

\begin{ojo}
  Exhibir una clique de tamaño $k$ solo prueba $\omega(G) \ge k$; para
  la igualdad falta un argumento de que no hay una más grande, como
  el de grados de arriba. Las diapositivas se quedan en la primera
  mitad. Y una clique es un \textbf{conjunto de vértices}, no un
  grafo: la frase ``el grafo completo $K_4$ es una clique de tamaño
  4'' mezcla las dos cosas. La clique es $\set{a, b, c, d}$; su
  subgrafo inducido es isomorfo a $K_4$.
\end{ojo}

La diferencia entre maximal y máxima es algorítmica. Una clique
maximal se encuentra con un barrido voraz: se recorren los vértices
en cualquier orden y se agrega cada uno que sea adyacente a todos los
ya escogidos. Con matriz, cada revisión cuesta $\Ord{n}$ y el total
es $\Ord{n^2}$. El resultado depende del orden: en el grafo pesado,
el orden A, B, \dots, F da $\set{A, B, C, D}$, la máxima, y el orden
F, E, \dots, A da $\set{F, E, D}$, que es maximal y nada más. Encontrar
la máxima es otra historia: decidir si $\omega(G) \ge k$ es el
problema CLIQUE, NP-completo (Cormen, teorema 34.11), y la fuerza
bruta revisa los $2^n$ subconjuntos. Es la misma brecha que entre
Euler y Hamilton.

\begin{observacion}
  Si $G$ es bipartito y tiene al menos una arista, $\omega(G) = 2$,
  porque un triángulo es un ciclo impar. Al revés no: el grafo $G$ de
  la clase no tiene triángulos, así que $\omega(G) = 2$, y aun así no
  es bipartito por su ciclo de longitud $5$.
\end{observacion}

\begin{ejercicio}[Número de clique]
  \label{ej:clique-hexagono}
  Encuentre $\omega(G)$ para
  $E = \set{ab, bc, cd, de, ef, fa, ac, bf, bd}$.

  \begin{minipage}[c]{0.28\linewidth}
  \centering
  \begin{tikzpicture}
    \node[vertOk] (a) at (-0.8, 1.4) {a};
    \node[vertOk] (b) at (0.8, 1.4) {b};
    \node[vertOk] (c) at (1.6, 0) {c};
    \node[vert] (d) at (0.8, -1.4) {d};
    \node[vert] (e) at (-0.8, -1.4) {e};
    \node[vert] (f) at (-1.6, 0) {f};
    \foreach \x/\y in {c/d, d/e, e/f, f/a, b/f, b/d}
      \draw[arista] (\x) -- (\y);
    \foreach \x/\y in {a/b, b/c, a/c}
      \draw[aristaArbol] (\x) -- (\y);
  \end{tikzpicture}
  \end{minipage}\hfill
  \begin{minipage}[c]{0.68\linewidth}
  \textit{Solución.} Hay tres triángulos, $\set{a, b, c}$,
  $\set{a, b, f}$ y $\set{b, c, d}$, así que $\omega(G) \ge 3$. Una
  clique de tamaño $4$ necesita cuatro vértices de grado al menos
  $3$; los candidatos son a, b, c, d y f, y cada una de sus cinco
  cuartetas contiene un par no adyacente, a--d o c--f. Por tanto
  $\omega(G) = 3$.
  \end{minipage}
\end{ejercicio}

% =====================================================================

\subsection*{Árboles}

\subsubsection*{Definición y caracterizaciones}

\begin{definicion}[Árbol y bosque]
  \label{def:arbol}
  Un \textbf{árbol} es un grafo no dirigido, conexo y sin ciclos. Un
  \textbf{bosque} es un grafo sin ciclos, no necesariamente conexo:
  cada una de sus componentes es un árbol. Una \textbf{hoja} de un
  árbol es un vértice de grado $1$; las diapositivas también le dicen
  vértice terminal o colgante.
\end{definicion}

Cormen dice \textit{free tree} para distinguirlo del árbol con raíz.
Todo lo que sigue sale de dos lemas.

\begin{lema}[Hojas]
  \label{lem:arbol-hojas}
  Todo árbol con $n \ge 2$ vértices tiene al menos dos hojas.
\end{lema}

\begin{proof}
  Sea $v_0, v_1, \dots, v_k$ un camino de longitud máxima; $k \ge 1$
  porque un árbol conexo con dos o más vértices tiene aristas. Si
  $v_0$ tuviera un vecino distinto de $v_1$, o ese vecino está en el
  camino y cierra un ciclo, o no está y alarga el camino. Ninguna de
  las dos se puede, así que $\text{gr}(v_0) = 1$; lo mismo para
  $v_k$.
\end{proof}

\begin{lema}[Aristas]
  \label{lem:arbol-aristas}
  Todo árbol con $n$ vértices tiene $n - 1$ aristas. Un bosque con
  $n$ vértices y $c$ componentes tiene $n - c$.
\end{lema}

\begin{proof}
  Inducción sobre $n$. Con $n = 1$ no hay aristas. Con $n \ge 2$ se
  quita una hoja $v$ y su arista. Lo que queda sigue sin ciclos y
  sigue conexo, porque ningún camino entre otros dos vértices pasa
  por $v$: entraría y saldría por su única arista. Es un árbol con
  $n - 1$ vértices, que por hipótesis tiene $n - 2$ aristas; con la
  de $v$ son $n - 1$. Para el bosque se suma sobre las componentes:
  $\sum (n_i - 1) = n - c$.
\end{proof}

\begin{teorema}[Caracterizaciones de un árbol]
  \label{teo:arbol-caracterizaciones}
  Para un grafo no dirigido $G$ con $n$ vértices y $m$ aristas, son
  equivalentes:
  \begin{enumerate}[label=(\arabic*), itemsep=0pt, parsep=0pt]
    \item $G$ es un árbol: conexo y sin ciclos.
    \item Entre cada par de vértices hay un único camino.
    \item $G$ es conexo y $m = n - 1$.
    \item $G$ no tiene ciclos y $m = n - 1$.
    \item $G$ es conexo y cada arista es un puente.
    \item $G$ no tiene ciclos y agregar cualquier arista nueva crea
          uno.
  \end{enumerate}
\end{teorema}

\begin{proof}
  (1)$\Leftrightarrow$(2). Conexo dice que hay al menos un camino. Si
  hubiera dos caminos distintos $P$ y $P'$ de $u$ a $v$, sea $x$ el
  último vértice común antes de que se separen y $y$ el primer vértice
  de $P$, después de $x$, que vuelve a estar en $P'$. Los tramos de
  $P$ y de $P'$ entre $x$ e $y$ solo comparten sus extremos y juntos
  cierran un ciclo. Al revés, si hubiera un ciclo, dos vértices
  vecinos en él estarían unidos por la arista y por el resto del
  ciclo: dos caminos.

  (1)$\Rightarrow$(3) y (1)$\Rightarrow$(4) son el lema
  \ref{lem:arbol-aristas}.

  (3)$\Rightarrow$(1). Si $G$ conexo tuviera un ciclo, quitar una
  arista del ciclo no lo desconecta, porque sus extremos siguen unidos
  por el resto del ciclo. Repitiendo hasta que no queden ciclos se
  llega a un árbol con $n$ vértices y menos de $n - 1$ aristas,
  contra el lema.

  (4)$\Rightarrow$(1). Sin ciclos es un bosque; si tiene $c$
  componentes, $m = n - c$, y con $m = n - 1$ sale $c = 1$.

  (1)$\Leftrightarrow$(5). En un árbol, quitar $\set{u, v}$ desconecta
  a $u$ de $v$: el único camino entre ellos era esa arista. Al revés,
  una arista de un ciclo nunca es puente, así que si todas lo son no
  hay ciclos.

  (1)$\Leftrightarrow$(6). En un árbol, si $\set{u, v} \notin E$, el
  único camino de $u$ a $v$ más la arista nueva es un ciclo. Al
  revés, si $G$ sin ciclos no fuera conexo, una arista entre dos
  componentes no crearía ningún ciclo.
\end{proof}

\begin{idea}
  (5) y (6) dicen lo mismo desde los dos lados: un árbol es un grafo
  conexo con el mínimo de aristas y un grafo acíclico con el máximo.
  Kruskal vive de eso: por (6), una arista entre dos vértices que ya
  están unidos cerraría un ciclo, y por (4), se detiene a las $n - 1$
  aristas (sección \ref{sec:mst}).
\end{idea}

\begin{ejemplo}[Árbol, ni árbol ni bosque, bosque]
  \label{ej:arbol-tres}
  Tres grafos sobre los mismos seis vértices.
  \begin{center}
  \begin{tikzpicture}
    \foreach \s/\x in {1/0, 2/3.4, 3/6.8} {
      \begin{scope}[xshift=\x cm, name prefix=g\s-]
        \node[vert] (a) at (-0.7, 1.3) {a};
        \node[vert] (b) at (0.7, 1.3) {b};
        \node[vert] (c) at (0, 0.4) {c};
        \node[vert] (d) at (0, -0.6) {d};
        \node[vert] (e) at (-0.7, -1.5) {e};
        \node[vert] (f) at (0.7, -1.5) {f};
        \node[font=\small] at (0, -2.2) {$G_{\s}$};
      \end{scope}
    }
    \foreach \x/\y in {a/c, b/c, c/d, d/e, d/f}
      \draw[arista] (g1-\x) -- (g1-\y);
    \foreach \x/\y in {a/c, b/c, c/d, d/e, d/f}
      \draw[arista] (g2-\x) -- (g2-\y);
    \draw[aristaAlerta] (g2-a) -- (g2-b);
    \foreach \x/\y in {a/b, c/d, d/e, d/f}
      \draw[arista] (g3-\x) -- (g3-\y);
  \end{tikzpicture}
  \end{center}
  $G_1$ es árbol: conexo, con $6$ vértices y $5$ aristas, por (3).
  $G_2$ le agrega a--b y tiene el ciclo a, b, c, a: conexo con $6$
  aristas, no es árbol ni bosque. $G_3$ no tiene ciclos pero tampoco
  es conexo: es un bosque de dos árboles, y en efecto
  $m = 4 = 6 - 2$.
\end{ejemplo}

\subsubsection*{Árboles con raíz}

\begin{definicion}[Árbol con raíz]
  \label{def:arbol-raiz}
  Un \textbf{árbol con raíz} (\textit{rooted tree}) es un árbol con un
  vértice distinguido, la \textbf{raíz}; las aristas se leen de arriba
  hacia abajo. Para un vértice $v$:
  \begin{itemize}[itemsep=0.1em]
    \item si $v$ no es la raíz, su \textbf{padre} es el vecino que
          está en el camino de $v$ a la raíz; $v$ es \textbf{hijo} de
          su padre, y dos vértices con el mismo padre son
          \textbf{hermanos};
    \item sus \textbf{ancestros} son los vértices del camino de la
          raíz a $v$, sin contar a $v$; sus \textbf{descendientes},
          los vértices que tienen a $v$ como ancestro;
    \item una \textbf{hoja} es un vértice sin hijos y un vértice
          \textbf{interno} uno con al menos un hijo;
    \item el \textbf{subárbol con raíz $v$} tiene a $v$, sus
          descendientes y las aristas entre ellos.
  \end{itemize}
  El \textbf{nivel} de $v$ es la longitud del camino de la raíz a $v$;
  la raíz está en el nivel $0$. La \textbf{altura} del árbol es el
  mayor nivel de sus hojas.
\end{definicion}

\begin{ejemplo}[El vocabulario sobre un árbol]
  \label{ej:arbol-vocabulario}
  El árbol de las diapositivas, con raíz r; en verde, las hojas.
  \begin{center}
  \begin{tikzpicture}
    \node[vert] (r) at (0, 0) {r};
    \node[vert] (a) at (-2.8, -1) {a};
    \node[vert] (b) at (0, -1) {b};
    \node[vert] (c) at (2.8, -1) {c};
    \node[vertOk] (d) at (-3.6, -2) {d};
    \node[vertOk] (e) at (-2.0, -2) {e};
    \node[vertOk] (f) at (-0.8, -2) {f};
    \node[vert] (g) at (0.8, -2) {g};
    \node[vertOk] (h) at (2.0, -2) {h};
    \node[vertOk] (i) at (3.6, -2) {i};
    \node[vertOk] (j) at (0.8, -3) {j};
    \foreach \x/\y in {a/d, a/e, b/f, c/h, c/i, r/a, r/c}
      \draw[arista] (\x) -- (\y);
    \foreach \x/\y in {r/b, b/g, g/j}
      \draw[aristaArbol] (\x) -- (\y);
    \draw[nota, dashed, rounded corners] (0.3, -1.6) rectangle
      (1.3, -3.4);
    \node[nota, anchor=east] at (0.25, -3.25) {subárbol de g};
    \foreach \y/\l in {0/0, -1/1, -2/2, -3/3}
      \node[nota, anchor=west] at (4.3, \y) {nivel \l};
  \end{tikzpicture}
  \end{center}
  El padre de g es b; los hijos de b son f y g; el hermano de h es i.
  Las hojas son d, e, f, h, i, j y los internos r, a, b, c, g. Los
  ancestros de j son r, b, g, el camino en azul; los descendientes de
  c son h e i. El subárbol con raíz g tiene los vértices g, j y la
  arista $\set{g, j}$. j está en el nivel $3$ y la altura es $3$.
\end{ejemplo}

Los árboles de los recorridos ya son árboles con raíz: en el árbol BFS
de la sección \ref{sec:grafos}, $\pi[v]$ es el padre de $v$, $d[v]$ es
su nivel y la altura es la mayor distancia desde $s$. Un árbol es
\textbf{ordenado} si los hijos de cada vértice tienen un orden fijo
de izquierda a derecha; dos árboles con los mismos padres pero con los
hijos de r intercambiados son iguales como árboles con raíz y
distintos como ordenados. En un binario ordenado se habla de hijo
izquierdo e hijo derecho.

\begin{ojo}
  ``Hoja'' no significa lo mismo en el árbol libre y en el árbol con
  raíz. En el libre es grado $1$; en el árbol con raíz, sin hijos.
  Difieren en la raíz: una raíz con un solo hijo tiene grado $1$ y no
  es hoja. Con los ancestros pasa algo parecido: las diapositivas
  excluyen al propio vértice, y Cormen lo incluye y dice ``ancestro
  propio'' para excluirlo.
\end{ojo}

\subsubsection*{Árboles $k$-arios: lleno, casi completo y perfecto}

\begin{definicion}[Árboles $k$-arios]
  \label{def:arbol-kario}
  Un árbol con raíz es \textbf{$k$-ario} si cada vértice tiene a lo más
  $k$ hijos; \textbf{binario} si $k = 2$ y \textbf{ternario} si
  $k = 3$. Un árbol $k$-ario es:
  \begin{itemize}[itemsep=0.1em]
    \item \textbf{lleno} si cada vértice interno tiene exactamente $k$
          hijos;
    \item \textbf{casi completo} si todos sus niveles están llenos
          salvo tal vez el último, que se llena de izquierda a
          derecha: es la forma del heap;
    \item \textbf{perfecto} si es lleno y todas sus hojas están en el
          mismo nivel.
  \end{itemize}
\end{definicion}

\begin{center}
\begin{tikzpicture}
  \node[vert] (r) at (0, 0) {r};
  \node[vertOk] (a) at (-0.9, -1) {a};
  \node[vert] (b) at (0.9, -1) {b};
  \node[vertOk] (c) at (0.4, -2) {c};
  \node[vertOk] (d) at (1.4, -2) {d};
  \foreach \x/\y in {r/a, r/b, b/c, b/d}
    \draw[arista] (\x) -- (\y);
  \node[nota, align=center] at (0.2, -2.8)
    {lleno, no casi completo};
  \begin{scope}[xshift=4.6cm]
    \node[vert] (r) at (0, 0) {r};
    \node[vert] (a) at (-0.9, -1) {a};
    \node[vertOk] (b) at (0.9, -1) {b};
    \node[vertOk] (c) at (-1.4, -2) {c};
    \foreach \x/\y in {r/a, r/b, a/c}
      \draw[arista] (\x) -- (\y);
    \node[nota, align=center] at (0, -2.8)
      {casi completo, no lleno};
  \end{scope}
  \begin{scope}[xshift=9.4cm]
    \node[vert] (r) at (0, 0) {r};
    \node[vert] (a) at (-1.1, -1) {a};
    \node[vert] (b) at (1.1, -1) {b};
    \foreach \v/\x in {c/-1.6, d/-0.6, e/0.6, f/1.6}
      \node[vertOk] (\v) at (\x, -2) {\v};
    \foreach \x/\y in {r/a, r/b, a/c, a/d, b/e, b/f}
      \draw[arista] (\x) -- (\y);
    \node[nota, align=center] at (0, -2.8)
      {perfecto: lleno y casi completo};
  \end{scope}
\end{tikzpicture}
\end{center}

A la izquierda, cada interno tiene dos hijos, pero el nivel $2$ se
llenó por la derecha. En medio, el heap de cuatro elementos: el
nivel $2$ está a la izquierda, pero a tiene un solo hijo. A la
derecha, el perfecto de altura $2$, que es las dos cosas.

\begin{ojo}
  ``Completo'' tiene tres significados según la fuente, y aquí se
  evita a propósito. Las diapositivas, siguiendo a Rosen
  (\textit{full $m$-ary tree}), llaman árbol $n$-ario completo al que
  aquí es \textbf{lleno}. La clase de heaps (17 de septiembre) llama
  árbol binario completo al que aquí es \textbf{casi completo}, que
  es el \textit{nearly complete} de Cormen. Y Cormen llama
  \textit{complete $k$-ary tree} al que aquí es \textbf{perfecto}.
  Leer ``completo'' sin saber de qué libro viene es como se confunde
  un heap con un árbol lleno, y las fórmulas de abajo valen para uno
  y no para el otro.
\end{ojo}

\begin{teorema}[Conteo en árboles llenos]
  \label{teo:arbol-conteo}
  Un árbol $k$-ario lleno con $i$ vértices internos tiene
  \[
    n = ki + 1 \text{ vértices} \qquad \text{y} \qquad
    \ell = (k - 1)i + 1 \text{ hojas.}
  \]
\end{teorema}

\begin{proof}
  Todo vértice salvo la raíz es hijo de exactamente un interno, y cada
  interno tiene exactamente $k$ hijos: $n - 1 = ki$. Las hojas son el
  resto, $\ell = n - i = (k - 1)i + 1$.
\end{proof}

Despejando, cualquiera de $n$, $i$ y $\ell$ da los otros dos; por
ejemplo, $i = (\ell - 1)/(k - 1)$ y $n = (k\ell - 1)/(k - 1)$. Las
diapositivas escriben $m = ni + 1$, con $n$ la aridad; aquí la aridad
es $k$ para no chocar con $n = \abs{V}$.

\begin{definicion}[Árbol balanceado]
  Un árbol con raíz de altura $h$ es \textbf{balanceado} si todas sus
  hojas están en los niveles $h$ o $h - 1$.
\end{definicion}

El del ejemplo \ref{ej:arbol-vocabulario} es balanceado: altura $3$,
hojas en los niveles $2$ y $3$. Todo casi completo es balanceado,
porque sus niveles están llenos hasta el penúltimo, y todo perfecto
también. Un árbol con raíz r, hijos a y b, y una cadena b, c, d no lo
es: tiene altura $3$ y la hoja a está en el nivel $1$.

\begin{ojo}
  Balanceado, en este sentido, no es AVL. La definición de las
  diapositivas es la de Rosen y solo mira dónde están las hojas. Un
  árbol AVL pide que en \textbf{cada} vértice las alturas de los dos
  subárboles difieran en a lo más $1$. Hay árboles que cumplen una y
  no la otra; que coincidan en un ejemplo no las hace la misma
  definición.
\end{ojo}

\begin{teorema}[Altura contra hojas]
  \label{teo:arbol-altura}
  (a) Un árbol $k$-ario de altura $h$ tiene a lo más $k^h$ hojas;
  por tanto $h \ge \lceil \log_k \ell \rceil$. (b) Si además es lleno
  y balanceado, $h = \lceil \log_k \ell \rceil$.
\end{teorema}

\begin{proof}
  (a) Inducción sobre $h$. Con $h = 0$ hay una hoja, $k^0$. Con
  $h \ge 1$, la raíz tiene a lo más $k$ hijos y cada uno es raíz de un
  subárbol de altura a lo más $h - 1$, con a lo más $k^{h-1}$ hojas.
  Como $h$ es entero y $\log_k \ell \le h$, sale la cota con techo.

  (b) Con $h \ge 1$. Por ser balanceado no hay hojas antes del nivel
  $h - 1$, y por ser lleno cada vértice de los niveles $0$ a $h - 2$
  tiene $k$ hijos: el nivel $h - 1$ tiene exactamente $k^{h-1}$
  vértices. De ellos, $I \ge 1$ son internos (alguna hoja está en el
  nivel $h$) y cada uno aporta $k$ hojas en el nivel $h$; los demás
  son hojas. Entonces
  \[
    \ell = (k^{h-1} - I) + kI = k^{h-1} + (k - 1)I > k^{h-1} ,
  \]
  y con (a), $k^{h-1} < \ell \le k^h$, o sea
  $h - 1 < \log_k \ell \le h$.
\end{proof}

\begin{observacion}
  La parte (a) da la cota inferior $\Omg{n \lg n}$ de ordenar por
  comparaciones. Un algoritmo que solo compara se describe con un
  árbol binario de decisión: cada vértice interno es una comparación y
  cada hoja un orden de salida. Para ordenar bien necesita al menos
  $n!$ hojas, así que su altura, el número de comparaciones en el peor
  caso, es al menos $\lg n! = \Tht{n \lg n}$.
\end{observacion}

\begin{ejemplo}[Torneo de eliminación directa]
  \label{ej:arbol-torneo}
  ¿Cuántos partidos hay en un torneo de eliminación directa con $16$
  jugadores? El torneo es un árbol binario lleno: las hojas son los
  jugadores, cada vértice interno es un partido y la raíz es la
  final. Con $k = 2$ y $\ell = 16$, el teorema \ref{teo:arbol-conteo}
  da $i = \ell - 1 = 15$ partidos, y como el cuadro es balanceado, el
  teorema \ref{teo:arbol-altura} da $\lceil \lg 16 \rceil = 4$ rondas.
  El argumento directo es más general: cada partido elimina a uno y
  todos salvo el campeón quedan eliminados una vez. Con $13$ jugadores
  y pases libres hay $12$ partidos y $\lceil \lg 13 \rceil = 4$
  rondas.
\end{ejemplo}

\begin{ejercicio}[Determine si existe]
  \label{ej:arbol-existe}
  (a) ¿Cuántas hojas tiene un árbol ternario lleno con $100$
  vértices? (b) ¿Existe un árbol binario lleno con $10$ vértices?
  (c) ¿Cuántos vértices y hojas tiene el binario perfecto de altura
  $3$?

  \textit{Solución.} (a) $100 = 3i + 1$ da $i = 33$ internos y
  $\ell = 2 \cdot 33 + 1 = 67$ hojas. (b) No: en un binario lleno
  $n = 2i + 1$ es impar. (c) Es lleno y balanceado con todas las hojas
  en el nivel $3$, así que $\ell = 2^3 = 8$, $i = \ell - 1 = 7$ y
  $n = 15$. En general el $k$-ario perfecto de altura $h$ tiene
  $k^h$ hojas y $(k^{h+1} - 1)/(k - 1)$ vértices.
\end{ejercicio}

% =====================================================================

\subsubsection*{Recorridos}

Un \textbf{recorrido} de un árbol visita todos sus vértices de manera
sistemática. En un árbol binario ordenado hay tres naturales, según
cuándo se visita la raíz respecto a sus dos subárboles:
\textbf{preorden} (raíz, izquierdo, derecho), \textbf{inorden}
(izquierdo, raíz, derecho) y \textbf{postorden} (izquierdo, derecho,
raíz).

\begin{algorithm}[H]
\caption{\textsc{Recorrer}$(x)$: los tres recorridos de un árbol
         binario}
\label{alg:arbol-recorridos}
\begin{algorithmic}[1]
\Procedure{Recorrer}{$x$}
  \If{$x \neq$ nulo}
    \State visitar $x$ \Comment{solo en preorden}
    \State \textsc{Recorrer}$(\text{izq}[x])$
    \State visitar $x$ \Comment{solo en inorden}
    \State \textsc{Recorrer}$(\text{der}[x])$
    \State visitar $x$ \Comment{solo en postorden}
  \EndIf
\EndProcedure
\end{algorithmic}
\end{algorithm}

Cada recorrido conserva una sola de las tres líneas que visitan:
\textsc{Preorden} es \textsc{Recorrer} con la línea 3,
\textsc{Inorden} con la 5 y \textsc{Postorden} con la 7, y las
llamadas recursivas se renombran igual. Lo único que cambia es en qué
momento se visita la raíz.

\textbf{Costo.} Cada vértice se visita una vez, y cada llamada hace
trabajo constante más dos llamadas recursivas. Con $j$ vértices en el
subárbol izquierdo,
\[
  T(n) = T(j) + T(n - j - 1) + d , \qquad T(0) = c ,
\]
y por sustitución $T(n) = (c + d)n + c$: la hipótesis da
$(c + d)j + c + (c + d)(n - j - 1) + c + d = (c + d)n + c$. Los tres
son $\Tht{n}$, sin importar la forma del árbol (Cormen, teorema
12.1).

\begin{ejemplo}[Los tres recorridos y el reloj de DFS]
  \label{ej:arbol-recorridos}
  El árbol de las diapositivas, con $\text{entra}$--$\text{sale}$ de
  \textsc{DFS-Visitar} junto a cada vértice, visitando primero el hijo
  izquierdo.
  \begin{center}
  \begin{tikzpicture}
    \node[vert] (a) at (0, 0) {a};
    \node[vert] (b) at (-2, -1) {b};
    \node[vert] (c) at (2, -1) {c};
    \node[vert] (d) at (-3, -2) {d};
    \node[vert] (e) at (-1, -2) {e};
    \node[vert] (f) at (1, -2) {f};
    \node[vert] (g) at (3, -2) {g};
    \node[vert] (h) at (-1.6, -3) {h};
    \node[vert] (i) at (-0.4, -3) {i};
    \foreach \x/\y in {a/b, a/c, b/d, b/e, c/f, c/g, e/h, e/i}
      \draw[arista] (\x) -- (\y);
    \foreach \v/\t/\lado/\ancla in {a/1--18/right/east,
        b/2--11/left/west, c/12--17/right/east, d/3--4/left/west,
        e/5--10/right/east, f/13--14/below/south,
        g/15--16/below/south, h/6--7/below/south,
        i/8--9/below/south}
      \node[nota, \lado=1pt] at (\v.\ancla) {\t};
  \end{tikzpicture}
  \end{center}
  \begin{center}
  \small
  \begin{tabular}{lll}
    \toprule
    \textbf{Recorrido} & \textbf{Resultado}
      & \textbf{Cómo sale del reloj} \\
    \midrule
    Preorden & a, b, d, e, h, i, c, f, g
      & orden creciente de $\text{entra}$ \\
    Inorden & d, b, h, e, i, a, f, c, g & (solo binarios) \\
    Postorden & d, h, i, e, b, f, g, c, a
      & orden creciente de $\text{sale}$ \\
    Por niveles & a, b, c, d, e, f, g, h, i & BFS desde a \\
    \bottomrule
  \end{tabular}
  \end{center}
  El preorden es el orden en que DFS descubre los vértices y el
  postorden el orden en que los termina. El inorden solo tiene sentido
  en binarios, porque necesita un ``en medio'' entre dos hijos. El
  recorrido por niveles es BFS.
\end{ejemplo}

\subsubsection*{Árboles de expresiones y notación posfija}

Una expresión aritmética con operadores binarios es un árbol binario
lleno: cada operador es un vértice interno con sus dos operandos como
hijos, y los números y variables son las hojas. El árbol de
$(x + y)^2 + (x - 4)/3$, con $\wedge$ para la potencia:
\begin{center}
\begin{tikzpicture}[op/.style={vert, font=\normalsize}]
  \node[op] (s) at (0, 0) {$+$};
  \node[op] (p) at (-2.2, -1) {$\wedge$};
  \node[op] (q) at (2.2, -1) {$/$};
  \node[op] (s2) at (-3.2, -2) {$+$};
  \node[vertOk] (two) at (-1.2, -2) {$2$};
  \node[op] (m) at (1.2, -2) {$-$};
  \node[vertOk] (three) at (3.2, -2) {$3$};
  \node[vertOk] (x1) at (-3.8, -3) {$x$};
  \node[vertOk] (y) at (-2.6, -3) {$y$};
  \node[vertOk] (x2) at (0.6, -3) {$x$};
  \node[vertOk] (four) at (1.8, -3) {$4$};
  \foreach \x/\y in {s/p, s/q, p/s2, p/two, q/m, q/three, s2/x1,
                     s2/y, m/x2, m/four}
    \draw[arista] (\x) -- (\y);
\end{tikzpicture}
\end{center}
Los tres recorridos dan tres notaciones; las diapositivas escriben
``postfija'' por posfija.
\begin{center}
\small
\begin{tabular}{lll}
  \toprule
  \textbf{Recorrido} & \textbf{Notación} & \textbf{Resultado} \\
  \midrule
  Preorden & \textbf{prefija} (polaca)
    & ${+}\ {\wedge}\ {+}\ x\ y\ 2\ {/}\ {-}\ x\ 4\ 3$ \\
  Postorden & \textbf{posfija} (polaca inversa)
    & $x\ y\ {+}\ 2\ {\wedge}\ x\ 4\ {-}\ 3\ {/}\ {+}$ \\
  Inorden & \textbf{infija}, con paréntesis
    & $((x + y) \wedge 2) + ((x - 4) / 3)$ \\
  \bottomrule
\end{tabular}
\end{center}

\begin{ojo}
  El inorden sin paréntesis, $x + y \wedge 2 + x - 4 / 3$, es ambiguo:
  con la precedencia usual se leería $x + y^2 + x - 4/3$, que es otra
  expresión. La infija necesita paréntesis para fijar la forma del
  árbol. La prefija y la posfija no: cada operador sabe que lleva
  exactamente dos operandos, así que la cadena determina el árbol.
  Las diapositivas no escriben la infija.
\end{ojo}

La posfija se evalúa con una pila, de izquierda a derecha y sin mirar
adelante: los operandos se apilan, y cada operador desapila sus dos
operandos y apila el resultado.

\begin{algorithm}[H]
\caption{\textsc{Evaluar-Posfija}$(t_1, \dots, t_k)$}
\label{alg:arbol-evaluar-posfija}
\begin{algorithmic}[1]
\Procedure{Evaluar-Posfija}{$t_1, \dots, t_k$}
  \State $P \gets$ pila vacía
  \For{$i \gets 1$ \Hasta $k$}
    \If{$t_i$ es un número o una variable}
      \State \textsc{Apilar}$(P, \text{valor}(t_i))$
    \Else \Comment{$t_i$ es un operador binario $\circ$}
      \State $b \gets \textsc{Desapilar}(P)$; \
             $a \gets \textsc{Desapilar}(P)$
      \State \textsc{Apilar}$(P, a \circ b)$
    \EndIf
  \EndFor
  \State \Return \textsc{Desapilar}$(P)$
\EndProcedure
\end{algorithmic}
\end{algorithm}

\begin{ejemplo}[Evaluar la posfija con $x = 1$, $y = 2$]
  \label{ej:arbol-posfija}
  Contenido de la pila después de cada símbolo, con el fondo arriba:
  \begin{center}
  \small
  \begin{tabular}{*{11}{c}}
    \toprule
    $x$ & $y$ & $+$ & $2$ & $\wedge$ & $x$ & $4$ & $-$ & $3$ & $/$
      & $+$ \\
    \midrule
    $1$ & $1$ & $3$ & $3$ & $9$ & $9$ & $9$ & $9$ & $9$ & $9$ & $8$ \\
        & $2$ &     & $2$ &     & $1$ & $1$ & $-3$ & $-3$ & $-1$ & \\
        &     &     &     &     &     & $4$ &      & $3$  &      & \\
    \bottomrule
  \end{tabular}
  \end{center}
  El resultado es $8$, y directamente
  $(1 + 2)^2 + (1 - 4)/3 = 9 - 1 = 8$.
\end{ejemplo}

El orden de los desapilamientos importa: el primero que sale es el
operando \textbf{derecho}. En el $-$ del ejemplo sale primero el $4$
y luego el $1$, y el resultado es $1 - 4 = -3$, no $4 - 1$; con $-$,
$/$ o $\wedge$, invertirlos cambia el resultado. Hay un apilamiento
por símbolo y dos desapilamientos por operador: $\Tht{k}$. Funciona
porque el postorden escribe cada operador justo después de sus dos
subárboles completos, así que cuando llega el operador sus dos
valores son los dos de arriba de la pila.

% =====================================================================

\subsection*{Resumen}

\begin{center}
\small
\begin{tabular}{>{\raggedright\arraybackslash}p{5.2cm}%
                >{\raggedright\arraybackslash}p{5.4cm}l}
  \toprule
  \textbf{Problema} & \textbf{Método} & \textbf{Costo} \\
  \midrule
  ¿Hay recorrido euleriano? & grados impares y una BFS
    & $\Tht{n + m}$ \\
  \addlinespace
  ¿Hay ciclo hamiltoniano? & retroceso; NP-completo
    & $\Ord{n!}$ \\
  \addlinespace
  Agente viajero & fuerza bruta & $\Tht{n!}$ \\
  & Held y Karp & $\Tht{n^2 2^n}$ \\
  \addlinespace
  ¿Es bipartito? & BFS con dos lados; da partición o ciclo impar
    & $\Tht{n + m}$ \\
  \addlinespace
  Clique maximal & voraz con matriz & $\Ord{n^2}$ \\
  Clique máxima & NP-difícil; subconjuntos & $\Ord{2^n n^2}$ \\
  \addlinespace
  Recorrer un árbol & preorden, inorden, postorden & $\Tht{n}$ \\
  \addlinespace
  Evaluar una posfija de $k$ símbolos & pila & $\Tht{k}$ \\
  \bottomrule
\end{tabular}
\end{center}

\newpage

% =====================================================================
%  Grafos - Representaciones, BFS y DFS
% =====================================================================

\clase{Representaciones, BFS y DFS a fondo}{Complemento: Cormen,
  cap.\ 22.1 a 22.3}
\label{sec:busquedas}

\begin{quick_sum}
  La clase del 22 de octubre presentó las dos representaciones y los
  dos recorridos con un ejemplo cada uno. Esta sección los retoma con
  el rigor de Cormen. Primero, el costo de cada operación con listas y
  con matriz, y seis algoritmos chicos que manipulan representaciones:
  grados, transpuesto, cuadrado, sumidero universal en $\Ord{n}$,
  matriz de incidencia y simplificación de un multigrafo. Después, BFS
  con colores: por qué su costo es $\Ord{n + m}$ y no $\Tht{n + m}$, la
  prueba de que calcula distancias, el árbol BFS y el diámetro de un
  árbol. Al final, DFS con colores y tiempos: el bosque DFS, los
  teoremas del paréntesis y del camino blanco, los cuatro tipos de
  arista, la versión iterativa con sus dos trampas y las componentes
  conexas. Lo que queda es la maquinaria de pruebas sobre la que se
  apoyan las aplicaciones de DFS y los caminos más cortos de las
  secciones que siguen.
\end{quick_sum}

% Macros de esta seccion. \busqListaAdy{y}{etiqueta}{v1,v2,...} dibuja
% un renglon del arreglo Adj con su lista ligada, al estilo de las
% figuras 22.1 y 22.2 de Cormen.
% \busqBaseBFS dibuja el grafo de la figura 22.3 con todo en blanco.
\newcommand{\busqListaAdy}[3]{%
  \node[anchor=east, font=\small] at (-0.3, #1) {#2};
  \node[draw=gray!80!black, fill=fondoIdea, minimum size=0.5cm,
        inner sep=0pt] at (0, #1) {};
  \foreach \busqElem [count=\busqPos] in {#3} {
    \pgfmathsetmacro{\busqX}{0.85 + (\busqPos - 1)*1.25}
    \node[draw=gray!80!black, minimum size=0.5cm, inner sep=0pt,
          font=\small] at (\busqX, #1) {\busqElem};
    \node[draw=gray!80!black, minimum width=0.3cm,
          minimum height=0.5cm, inner sep=0pt]
      at (\busqX + 0.4, #1) {};
    \fill (\busqX - 0.85, #1) circle (1.1pt);
    \draw[->, >=stealth] (\busqX - 0.85, #1) -- (\busqX - 0.25, #1);
    \xdef\busqUltimoX{\busqX}
  }
  \draw (\busqUltimoX + 0.25, #1 - 0.25)
    -- (\busqUltimoX + 0.55, #1 + 0.25);
}
\newcommand{\busqBaseBFS}{%
  \foreach \busqV/\busqX/\busqY in {r/0/1, s/1/1, t/2/1, u/3/1,
                                     v/0/0, w/1/0, x/2/0, y/3/0}
    \coordinate (\busqV) at (\busqX, \busqY);
  \foreach \busqA/\busqB in {r/s, r/v, s/w, w/t, w/x, t/x, t/u, x/u,
                             x/y, u/y}
    \draw[arista] (\busqA) -- (\busqB);
  \foreach \busqV in {r, s, t, u}
    \node[vert, fill=white, label={[nota, text=black]above:$\busqV$}]
      at (\busqV) {$\infty$};
  \foreach \busqV in {v, w, x, y}
    \node[vert, fill=white, label={[nota, text=black]below:$\busqV$}]
      at (\busqV) {$\infty$};
  \node[font=\small] at (-0.15, -1.15) {$Q$};
}

La notación de atributos es la de la sección~\ref{sec:grafos}, que
también da su equivalencia con la de Cormen: $d[v]$, $\pi[v]$,
$\text{color}[v]$, y $\text{entra}[v]$ y $\text{sale}[v]$ para los
tiempos de DFS.

% =====================================================================

\subsection*{Listas y matriz: el costo de cada operación}

La clase del 22 de octubre definió las dos representaciones y comparó
cuatro operaciones. Aquí se fija la notación y se completa la cuenta.
Las \textbf{listas de adyacencia} son un arreglo $\mathit{Adj}$ de $n$
listas, una por vértice: $\mathit{Adj}[u]$ contiene los $v$ con
$(u, v) \in E$ en un dirigido, o con $\set{u, v} \in E$ en uno no
dirigido, en un orden arbitrario. El orden no significa nada para el
grafo, pero sí decide qué vecino descubren primero BFS y DFS; por eso
cada ejemplo de esta sección dice qué orden usa.

\begin{lema}[Longitud total de las listas]
  \label{lem:busq-longitud-listas}
  En un grafo dirigido $\sum_{u \in V} \abs{\mathit{Adj}[u]} = m$. En
  uno no dirigido la suma es $2m$.
\end{lema}

\begin{proof}
  En un dirigido la arista $(u, v)$ aparece una sola vez, como $v$ en
  $\mathit{Adj}[u]$. En uno no dirigido $\set{u, v}$ aparece dos
  veces: como $v$ en $\mathit{Adj}[u]$ y como $u$ en $\mathit{Adj}[v]$.
\end{proof}

Es la misma cuenta que $\sum_v \text{gr}(v) = 2m$ de la clase. En un
dirigido se lee así: la suma de los grados de salida es $m$, y la de
los grados de entrada también. Escribimos $\text{gr}^{+}(u)$ para el
\textbf{grado de salida} (el número de aristas $(u, v)$) y
$\text{gr}^{-}(u)$ para el \textbf{grado de entrada} (el de aristas
$(v, u)$). De este lema sale el $\Tht{n + m}$ de memoria de las
listas: $n$ cabezas más $m$ o $2m$ nodos. El término $n$ no se puede
quitar, porque un grafo con vértices aislados sigue pagando su
arreglo de cabezas.

\begin{center}
\small
\begin{tabular}{l>{\raggedright\arraybackslash}p{6.1cm}%
                >{\raggedright\arraybackslash}p{3.6cm}}
  \toprule
  \textbf{Operación} & \textbf{Listas} & \textbf{Matriz} \\
  \midrule
  Memoria & $\Tht{n + m}$ & $\Tht{n^2}$ \\
  ¿Es $(u, v)$ arista? & $\Ord{\text{gr}^{+}(u)}$, peor caso
    $\Ord{n}$ & $\Tht{1}$ \\
  Recorrer los vecinos de $u$ & $\Tht{1 + \text{gr}^{+}(u)}$
    & $\Tht{n}$ \\
  Recorrer todas las aristas & $\Tht{n + m}$ & $\Tht{n^2}$ \\
  Grado de salida de $u$ & $\Tht{\text{gr}^{+}(u)}$; $\Tht{1}$ si la
    lista guarda su longitud & $\Tht{n}$ \\
  Grados de entrada de todos & $\Tht{n + m}$ & $\Tht{n^2}$ \\
  Agregar una arista & $\Ord{1}$, al frente y sin revisar duplicados
    & $\Tht{1}$ \\
  Borrar la arista $(u, v)$ & $\Ord{\text{gr}^{+}(u)}$; en no
    dirigido, también $\Ord{\text{gr}(v)}$ & $\Tht{1}$ \\
  Grafo transpuesto & $\Tht{n + m}$ & $\Tht{n^2}$ \\
  Cuadrado $G^2$ & $\Ord{nm}$ & $\Ord{n^3}$, u
    $\Ord{n^{\lg 7}}$ con Strassen \\
  \bottomrule
\end{tabular}
\end{center}

En un no dirigido, $\text{gr}^{+}(u)$ es simplemente $\text{gr}(u)$.
Las cuatro filas de la clase siguen ahí; las demás las justifican los
algoritmos de esta subsección.

\textbf{Grafos pesados y atributos.} En listas, el peso $w(u, v)$
viaja en el nodo de $v$ dentro de $\mathit{Adj}[u]$; para el grafo
pesado de la clase, $\mathit{Adj}[A] = (B, 4), (C, 2), (D, 7)$. En la
matriz va en la casilla $(u, v)$, y en las casillas sin arista se
guarda lo que pida el problema: $\infty$ para caminos más cortos, $0$
para capacidades, nulo si hay que distinguir. El $0$ en la diagonal
de la matriz pesada de la clase es la convención de Floyd, no parte de
la representación. Los atributos de vértice, como $d[v]$ o
$\text{color}[v]$, viven en arreglos indexados por vértice, con acceso
$\Ord{1}$; los de arista, como un flujo $f(u, v)$, en el nodo de la
lista o en una segunda matriz.

\begin{ojo}
  \begin{itemize}[leftmargin=1.2em, itemsep=0.15em]
    \item $\Tht{n + m}$ no es $\Tht{m}$. Las dos coinciden cuando el
      grafo es no dirigido y conexo, porque entonces $m \geq n - 1$,
      pero no en general.
    \item Cada arista no dirigida vive dos veces en las listas.
      Borrarla o cambiarle el peso obliga a tocar $\mathit{Adj}[u]$ y
      $\mathit{Adj}[v]$; olvidar una copia deja el grafo
      inconsistente.
    \item La matriz cobra $\Tht{n^2}$ desde que se crea, porque
      inicializarla ya cuesta eso. Un algoritmo que promete
      $\Ord{n + m}$ no puede construir una matriz por dentro, ni
      siquiera para consultar aristas rápido.
    \item ``La suma de cada renglón es el grado'' y ``la diagonal es
      cero'' valen en grafos no dirigidos sin lazos. En un dirigido
      el renglón $u$ suma $\text{gr}^{+}(u)$ y la columna $u$ suma
      $\text{gr}^{-}(u)$; un lazo $(u, u)$ pone un $1$ en la diagonal
      y suma $1$ a cada uno de los dos grados.
    \item Consultar una arista con listas no es $\Ord{1}$. En un no
      dirigido conviene buscar en la lista más corta:
      $\Ord{\min(\text{gr}(u), \text{gr}(v))}$.
  \end{itemize}
\end{ojo}

\subsubsection*{Grados de entrada y de salida}

\begin{algorithm}[H]
\caption{\textsc{Grados}$(G)$}
\label{alg:busq-grados}
\begin{algorithmic}[1]
\Procedure{Grados}{$G$}
  \ForAll{$u \in V$}
    \State $\text{gr}^{+}[u] \gets 0$, \quad $\text{gr}^{-}[u] \gets 0$
  \EndFor
  \ForAll{$u \in V$}
    \ForAll{$v \in \mathit{Adj}[u]$}
      \State $\text{gr}^{+}[u] \gets \text{gr}^{+}[u] + 1$
      \State $\text{gr}^{-}[v] \gets \text{gr}^{-}[v] + 1$
    \EndFor
  \EndFor
  \State \Return $\text{gr}^{+}, \text{gr}^{-}$
\EndProcedure
\end{algorithmic}
\end{algorithm}

Las líneas 2 a 4 cuestan $\Tht{n}$. El ciclo de la línea 5 da $n$
vueltas, cada una de costo $\Tht{1}$ más la longitud de
$\mathit{Adj}[u]$, y por el lema~\ref{lem:busq-longitud-listas} las
vueltas internas suman $m$: total $\Tht{n + m}$. Si solo interesan los
grados de salida y cada lista guarda su longitud, basta leerla, en
$\Tht{n}$. Los de entrada no admiten atajo: cualquier entrada de
cualquier lista puede sumarle a cualquier vértice, así que hay que
verlas todas. Con matriz, $\text{gr}^{+}(u)$ es la suma del renglón $u$
y $\text{gr}^{-}(u)$ la de la columna $u$: $\Tht{n}$ cada uno y
$\Tht{n^2}$ para todos.

\subsubsection*{Grafo transpuesto}

\begin{definicion}[Grafo transpuesto]
  El \textbf{transpuesto} de un grafo dirigido $G = (V, E)$ es
  $G^{\mathsf{T}} = (V, E^{\mathsf{T}})$ con
  $E^{\mathsf{T}} = \set{(v, u) : (u, v) \in E}$: el mismo grafo con
  todas las aristas invertidas. Su matriz es $A^{\mathsf{T}}$.
\end{definicion}

\begin{algorithm}[H]
\caption{\textsc{Transponer-Listas}$(G)$}
\label{alg:busq-transponer}
\begin{algorithmic}[1]
\Procedure{Transponer-Listas}{$G$}
  \ForAll{$u \in V$}
    \State $\mathit{Adj}^{\mathsf{T}}[u] \gets$ lista vacía
  \EndFor
  \ForAll{$u \in V$, en orden creciente}
    \ForAll{$v \in \mathit{Adj}[u]$}
      \State agregar $u$ al final de $\mathit{Adj}^{\mathsf{T}}[v]$
        \Comment{$(u, v)$ se vuelve $(v, u)$}
    \EndFor
  \EndFor
  \State \Return $\mathit{Adj}^{\mathsf{T}}$
\EndProcedure
\end{algorithmic}
\end{algorithm}

Cada inserción cuesta $\Ord{1}$ si la lista guarda un apuntador a su
último nodo, y hay $m$ de ellas: $\Tht{n + m}$, que es también el
tamaño de la salida. Hay un efecto secundario útil: como $u$ se
recorre en orden creciente y se agrega al final, cada
$\mathit{Adj}^{\mathsf{T}}[v]$ sale ordenada. Transponer dos veces
deja ordenadas las listas de $G$ en $\Tht{n + m}$, sin una sola
comparación; es un ordenamiento por cubetas disfrazado. Con matriz,
transponer es intercambiar $A[i][j]$ con $A[j][i]$ arriba de la
diagonal, $(n^2 - n)/2$ intercambios en su lugar, $\Tht{n^2}$. El
transpuesto reaparece en las componentes fuertemente conexas de la
sección~\ref{sec:aplicaciones-dfs}, que corren DFS sobre $G$ y luego
sobre $G^{\mathsf{T}}$ y necesitan transponer en tiempo lineal.

\begin{ejemplo}[Un grafo dirigido, sus grados y su transpuesto]
  \label{ej:busq-dirigido-cormen}
  El grafo de la figura 22.2 de Cormen tiene seis vértices y ocho
  aristas: $(1,2)$, $(1,4)$, $(2,5)$, $(3,5)$, $(3,6)$, $(4,2)$,
  $(5,4)$ y el lazo $(6,6)$. A la izquierda el dibujo, al centro el
  arreglo $\mathit{Adj}$ con sus listas ligadas y a la derecha la
  matriz.
  \begin{center}
  \begin{tikzpicture}[baseline=(current bounding box.center)]
    \node[vert] (n1) at (0, 1.3) {1};
    \node[vert] (n2) at (1.4, 1.3) {2};
    \node[vert] (n3) at (2.8, 1.3) {3};
    \node[vert] (n4) at (0, 0) {4};
    \node[vert] (n5) at (1.4, 0) {5};
    \node[vert] (n6) at (2.8, 0) {6};
    \foreach \a/\b in {n1/n2, n1/n4, n2/n5, n3/n6, n3/n5, n4/n2,
                       n5/n4}
      \draw[dirigida] (\a) -- (\b);
    \draw[dirigida] (n6) to[loop right, looseness=6] (n6);
  \end{tikzpicture}
  \hspace{0.7cm}
  \begin{tikzpicture}[baseline=(current bounding box.center)]
    \busqListaAdy{0}{1}{2,4}
    \busqListaAdy{-0.6}{2}{5}
    \busqListaAdy{-1.2}{3}{6,5}
    \busqListaAdy{-1.8}{4}{2}
    \busqListaAdy{-2.4}{5}{4}
    \busqListaAdy{-3.0}{6}{6}
  \end{tikzpicture}
  \hspace{0.7cm}
  $\begin{array}{c|cccccc}
      & 1 & 2 & 3 & 4 & 5 & 6 \\ \hline
    1 & 0 & 1 & 0 & 1 & 0 & 0 \\
    2 & 0 & 0 & 0 & 0 & 1 & 0 \\
    3 & 0 & 0 & 0 & 0 & 1 & 1 \\
    4 & 0 & 1 & 0 & 0 & 0 & 0 \\
    5 & 0 & 0 & 0 & 1 & 0 & 0 \\
    6 & 0 & 0 & 0 & 0 & 0 & 1
  \end{array}$
  \end{center}
  Cada arista aparece una sola vez en las listas, que suman
  $8 = m$. La matriz no es simétrica ($a_{12} = 1$, $a_{21} = 0$) y el
  lazo pone un $1$ en $a_{66}$. \textsc{Grados} y
  \textsc{Transponer-Listas} recorren las listas en el mismo orden y se
  pueden correr en una sola pasada. Los grados de salida son
  $2, 1, 2, 1, 1, 1$, las longitudes de las listas o las sumas de los
  renglones, y los de entrada $0, 2, 0, 2, 2, 2$, las sumas de las
  columnas; las dos listas suman $8$. Los vértices $1$ y $3$
  no reciben ninguna arista, así que en $G^{\mathsf{T}}$ sus listas
  quedan vacías. Las demás salen
  $\mathit{Adj}^{\mathsf{T}}[2] = 1, 4$,
  $\mathit{Adj}^{\mathsf{T}}[4] = 1, 5$,
  $\mathit{Adj}^{\mathsf{T}}[5] = 2, 3$ y
  $\mathit{Adj}^{\mathsf{T}}[6] = 3, 6$, todas en orden creciente.
\end{ejemplo}

\subsubsection*{Cuadrado de un grafo}

\begin{definicion}[Cuadrado de un grafo]
  El \textbf{cuadrado} de un grafo dirigido $G = (V, E)$ es
  $G^2 = (V, E^2)$, donde $(u, v) \in E^2$ si y solo si $G$ tiene una
  caminata de una o de dos aristas de $u$ a $v$.
\end{definicion}

Cormen dice ``un camino de a lo más dos aristas''. Leído al pie de la
letra eso incluye el camino de cero aristas de $u$ a $u$, y entonces
todo vértice tendría lazo en $G^2$. La lectura usual, y la de los
algoritmos de abajo, es la de la definición: $(u, u) \in E^2$ solo si
$u$ tiene lazo o hay una ida y vuelta $u \to w \to u$.

\begin{algorithm}[H]
\caption{\textsc{Cuadrado-Listas}$(G)$}
\label{alg:busq-cuadrado}
\begin{algorithmic}[1]
\Procedure{Cuadrado-Listas}{$G$}
  \ForAll{$u \in V$}
    \State $\mathit{Adj}^2[u] \gets$ lista vacía, \quad
           $\text{marca}[u] \gets \textsc{nulo}$
  \EndFor
  \ForAll{$u \in V$}
    \ForAll{$w \in \mathit{Adj}[u]$} \Comment{caminatas de una arista}
      \If{$\text{marca}[w] \neq u$}
        \State $\text{marca}[w] \gets u$; \ agregar $w$ a
               $\mathit{Adj}^2[u]$
      \EndIf
      \ForAll{$v \in \mathit{Adj}[w]$} \Comment{de dos aristas}
        \If{$\text{marca}[v] \neq u$}
          \State $\text{marca}[v] \gets u$; \ agregar $v$ a
                 $\mathit{Adj}^2[u]$
        \EndIf
      \EndFor
    \EndFor
  \EndFor
  \State \Return $\mathit{Adj}^2$
\EndProcedure
\end{algorithmic}
\end{algorithm}

$\text{marca}[x] = u$ significa ``$x$ ya está en $\mathit{Adj}^2[u]$''.
Sin la marca, $v$ entraría una vez por cada caminata $u \to w \to v$ y
otra más si además $(u, v) \in E$. La marca no se reinicializa entre
un $u$ y el siguiente: una marca puesta por un vértice anterior nunca
coincide con el $u$ actual. Reinicializarla costaría $\Tht{n}$ por
vértice, $\Tht{n^2}$ en total, y destruiría la cota.

\textbf{Costo.} El ciclo de la línea 6 da $m$ vueltas en total. Dentro
de la vuelta $(u, w)$, el de la línea 10 da $\text{gr}^{+}(w)$
vueltas, y cada $w$ aparece en $\text{gr}^{-}(w)$ listas, así que el
trabajo interno es
\[
  \sum_{u \in V} \sum_{w \in \mathit{Adj}[u]} \text{gr}^{+}(w)
  = \sum_{w \in V} \text{gr}^{-}(w)\, \text{gr}^{+}(w)
  \leq n \sum_{w \in V} \text{gr}^{-}(w) = nm .
\]
El total es $\Ord{n + m + nm} = \Ord{nm}$ para $m \geq 1$. La cota se
alcanza: con $k$ vértices que apuntan a un centro y el centro
apuntando a otros $k$, se tiene $n = 2k + 1$ y $m = 2k$, pero $G^2$
tiene $k^2 + 2k$ aristas. La salida sola ya mide $\Tht{n^2}$ con
$m = \Tht{n}$, así que no existe un algoritmo lineal para $G^2$.

Con matriz, la entrada $(i, j)$ del producto entero $A \cdot A$ es
$\sum_k a_{ik} a_{kj}$, y cada sumando vale $1$ exactamente cuando
$i \to k \to j$ es una caminata. Por inducción, $(A^k)_{ij}$ cuenta
las caminatas de $k$ aristas de $i$ a $j$. Entonces $E^2$ son las
entradas no nulas de $A + A^2$: con tres ciclos anidados cuesta
$\Tht{n^3}$, y con el producto de Strassen $\Ord{n^{\lg 7}}$, con
$\lg 7 \approx 2.81$, más $\Tht{n^2}$ para llevar a $1$ lo no nulo.
Es la semilla de los caminos más cortos por multiplicación de matrices
de la sección~\ref{sec:todos-pares}.

\begin{ejemplo}[El cuadrado del grafo dirigido]
  El grafo es el del ejemplo~\ref{ej:busq-dirigido-cormen}. Para
  $u = 1$, el algoritmo recorre $w = 2$ (agrega $2$) y su vecino $5$
  (agrega $5$), luego $w = 4$ (agrega $4$) y su vecino $2$, que ya
  tiene la marca de $1$ y no se repite. Así con cada $u$; las
  aristas nuevas salen punteadas.
  \begin{center}
  \begin{tikzpicture}
    \node[vert] (n1) at (0, 1.8) {1};
    \node[vert] (n2) at (2, 1.8) {2};
    \node[vert] (n3) at (4, 1.8) {3};
    \node[vert] (n4) at (0, 0) {4};
    \node[vert] (n5) at (2, 0) {5};
    \node[vert] (n6) at (4, 0) {6};
    \draw[dirigida] (n1) -- (n2);
    \draw[dirigida] (n1) -- (n4);
    \draw[dirigida] (n3) -- (n6);
    \draw[dirigida] (n3) -- (n5);
    \draw[dirigida] (n2) to[bend left=15] (n5);
    \draw[dirigida] (n4) to[bend left=15] (n2);
    \draw[dirigida] (n5) to[bend left=15] (n4);
    \draw[dirigida] (n6) to[loop right, looseness=6] (n6);
    \draw[dirigidaFuera] (n5) to[bend left=15] (n2);
    \draw[dirigidaFuera] (n2) to[bend left=15] (n4);
    \draw[dirigidaFuera] (n4) to[bend left=15] (n5);
    \draw[dirigidaFuera] (n1) -- (n5);
    \draw[dirigidaFuera] (n3) to[out=120, in=150, looseness=1.4]
      (n4);
    \node[nota] at (2, -0.75) {$G^{2}$: punteadas las cinco aristas
      nuevas};
  \end{tikzpicture}
  \end{center}
  $\abs{E^2} = 13$: las ocho de $G$ más $(1,5)$, $(2,4)$, $(3,4)$,
  $(4,5)$ y $(5,2)$. El ciclo $2 \to 5 \to 4 \to 2$ produce en
  $G^2$ el mismo ciclo recorrido al revés, porque saltar dos pasos
  hacia adelante en un ciclo de tres es dar uno hacia atrás.
\end{ejemplo}

\subsubsection*{Sumidero universal en $\Ord{n}$}

\begin{definicion}[Sumidero universal]
  En un grafo dirigido, un \textbf{sumidero universal} es un vértice
  con $\text{gr}^{-} = n - 1$ y $\text{gr}^{+} = 0$: todos los demás
  le apuntan y él no apunta a nadie, ni a sí mismo.
\end{definicion}

Con la matriz, casi cualquier pregunta sobre un grafo cuesta
$\Omg{n^2}$. Decidir si hay al menos una arista, por ejemplo, exige
leer todas las casillas en el peor caso: un adversario contesta $0$ a
cada lectura y, si el algoritmo se detiene antes, pone un $1$ en una
casilla no leída. El sumidero universal es la excepción que vale la
pena conocer, y escapa al argumento por este lema.

\begin{lema}
  \label{lem:busq-una-lectura}
  Sean $k \neq j$. Si $a_{kj} = 1$, $k$ no es sumidero universal,
  porque tiene una arista de salida. Si $a_{kj} = 0$, $j$ no lo es,
  porque le falta la arista desde $k$.
\end{lema}

Cada lectura fuera de la diagonal descarta a uno de los dos vértices,
no solo una casilla. Con $n - 1$ lecturas sobrevive un solo candidato.

\begin{algorithm}[H]
\caption{\textsc{Sumidero-Universal}$(A, n)$}
\label{alg:busq-sumidero}
\begin{algorithmic}[1]
\Procedure{Sumidero-Universal}{$A, n$}
  \State $k \gets 1$ \Comment{candidato}
  \For{$j \gets 2$ \Hasta $n$}
    \If{$A[k][j] = 1$} \Comment{$k$ tiene una arista de salida}
      \State $k \gets j$
    \EndIf
  \EndFor
  \For{$j \gets 1$ \Hasta $n$} \Comment{verificar al candidato}
    \If{$A[k][j] = 1$ \textbf{o} ($j \neq k$ \textbf{y}
        $A[j][k] = 0$)}
      \State \Return \textsc{nulo}
    \EndIf
  \EndFor
  \State \Return $k$
\EndProcedure
\end{algorithmic}
\end{algorithm}

Que el algoritmo devuelva un solo vértice no deja nada fuera:

\begin{lema}
  Un grafo dirigido tiene a lo más un sumidero universal.
\end{lema}

\begin{proof}
  Si $k \neq k'$ fueran sumideros, $k'$ recibiría arista de todos los
  demás, en particular $(k, k') \in E$, y eso contradice
  $\text{gr}^{+}(k) = 0$.
\end{proof}

\textbf{Corrección.} Invariante del ciclo de la línea 3: al empezar la
vuelta $j$, ningún vértice de $\set{1, \dots, j - 1}$ distinto de $k$
es sumidero universal. Al inicio el conjunto es $\set{1} = \set{k}$.
En la vuelta $j$ se lee $a_{kj}$ con $k < j$: si vale $1$, $k$ queda
descartado y $j$ es el nuevo candidato; si vale $0$, $j$ queda
descartado (lema~\ref{lem:busq-una-lectura}). Al terminar, el único
sumidero posible es $k$, pero puede no serlo: la primera fase siempre
entrega un candidato, haya sumidero o no. Las líneas 8 a 12 lo
verifican: el renglón $k$ completo en ceros, incluida la diagonal, que
descarta un lazo, y la columna $k$ en unos fuera de la diagonal.

\textbf{Costo.} La primera fase lee $n - 1$ casillas; la segunda, el
renglón ($n$) y la columna sin la diagonal ($n - 1$). A lo más
$3n - 2$ lecturas: $\Ord{n}$, sobre una entrada de tamaño $n^2$.

\begin{ejemplo}[Sumidero universal, con y sin sumidero]
  Seis vértices; todos tienen una arista hacia $4$, y además $(1,2)$,
  $(2,5)$, $(3,1)$, $(3,6)$, $(5,6)$ y $(6,1)$. El vértice $4$ tiene
  $\text{gr}^{-} = 5$ y $\text{gr}^{+} = 0$. En la matriz del centro,
  en amarillo, las casillas que lee la primera fase; a la derecha, en
  verde, el renglón y la columna que revisa la segunda.
  \begin{center}
  \begin{tikzpicture}[baseline=(current bounding box.center)]
    \node[vertOk] (n4) at (0, 0) {4};
    \node[vert] (n1) at (162:1.7) {1};
    \node[vert] (n2) at (90:1.7) {2};
    \node[vert] (n5) at (18:1.7) {5};
    \node[vert] (n6) at (-54:1.7) {6};
    \node[vert] (n3) at (-126:1.7) {3};
    \foreach \a/\b in {n1/n2, n2/n5, n5/n6, n3/n6, n3/n1, n6/n1}
      \draw[dirigida] (\a) -- (\b);
    \foreach \a in {n1, n2, n3, n5, n6}
      \draw[dirigida, draw=green!45!black] (\a) -- (n4);
  \end{tikzpicture}
  \hspace{0.8cm}
  \begin{tikzpicture}[baseline=(current bounding box.center),
      x=0.5cm, y=0.5cm,
      celda/.style={draw=gray!80!black, minimum size=0.5cm,
                    inner sep=0pt, outer sep=0pt, font=\small}]
    \foreach \fil/\col in {1/2, 2/3, 2/4, 4/5, 4/6}
      \fill[marcador] (\col - 0.5, -\fil - 0.5)
        rectangle (\col + 0.5, -\fil + 0.5);
    \foreach \fila [count=\fil] in {{0,1,0,1,0,0}, {0,0,0,1,1,0},
                                  {1,0,0,1,0,1}, {0,0,0,0,0,0},
                                  {0,0,0,1,0,1}, {1,0,0,1,0,0}} {
      \node[font=\scriptsize] at (0, -\fil) {\fil};
      \node[font=\scriptsize] at (\fil, 0) {\fil};
      \foreach \val [count=\col] in \fila
        \node[celda] at (\col, -\fil) {\val};
    }
    \begin{scope}[->, >=stealth, draw=acento, thick]
      \draw (2.2, -1.2) -- (2.8, -1.8);
      \draw (3.3, -2) -- (3.7, -2);
      \draw (4.15, -2.3) -- (4.85, -3.7);
      \draw (5.3, -4) -- (5.7, -4);
    \end{scope}
    \node[nota] at (3.5, -7.1) {fase 1: eliminación};
  \end{tikzpicture}
  \hspace{0.5cm}
  \begin{tikzpicture}[baseline=(current bounding box.center),
      x=0.5cm, y=0.5cm,
      celda/.style={draw=gray!80!black, minimum size=0.5cm,
                    inner sep=0pt, outer sep=0pt, font=\small}]
    \fill[fondoResultado] (0.5, -4.5) rectangle (6.5, -3.5);
    \fill[fondoResultado] (3.5, -6.5) rectangle (4.5, -0.5);
    \foreach \fila [count=\fil] in {{0,1,0,1,0,0}, {0,0,0,1,1,0},
                                  {1,0,0,1,0,1}, {0,0,0,0,0,0},
                                  {0,0,0,1,0,1}, {1,0,0,1,0,0}} {
      \node[font=\scriptsize] at (0, -\fil) {\fil};
      \node[font=\scriptsize] at (\fil, 0) {\fil};
      \foreach \val [count=\col] in \fila
        \node[celda] at (\col, -\fil) {\val};
    }
    \draw[green!45!black, thick] (0.5, -4.5) rectangle (6.5, -3.5);
    \draw[green!45!black, thick] (3.5, -6.5) rectangle (4.5, -0.5);
    \node[nota] at (3.5, -7.1) {fase 2: verificación de $k = 4$};
  \end{tikzpicture}
  \end{center}
  La primera fase lee una casilla por columna, siempre arriba de la
  diagonal: un $0$ avanza a la derecha y un $1$ baja al renglón de esa
  columna.
  \begin{center}
  \small
  \begin{tabular}{ccccl}
    \toprule
    $j$ & $k$ antes & $a_{kj}$ & $k$ después
      & \textbf{Qué se descarta} \\
    \midrule
    2 & 1 & 1 & 2 & $1$, porque tiene una arista de salida \\
    3 & 2 & 0 & 2 & $3$, porque no recibe arista de $2$ \\
    4 & 2 & 1 & 4 & $2$, porque tiene una arista de salida \\
    5 & 4 & 0 & 4 & $5$, porque no recibe arista de $4$ \\
    6 & 4 & 0 & 4 & $6$, porque no recibe arista de $4$ \\
    \bottomrule
  \end{tabular}
  \end{center}
  Queda $k = 4$. La segunda fase lee su renglón, seis ceros, y su
  columna fuera de la diagonal, cinco unos: $4$ es el sumidero. Se
  leyeron $5 + 6 + 5 = 16 = 3 \cdot 6 - 2$ de las $36$ casillas.

  Si se agrega la arista $(4, 6)$, ya no hay sumidero. La eliminación
  es la misma hasta $j = 6$, donde $a_{46} = 1$ deja como candidato a
  $6$, y la verificación lo tumba en la primera lectura, porque
  $a_{61} = 1$. Sin la segunda fase, el algoritmo devolvería $6$, que
  tiene $\text{gr}^{-} = 3$ y $\text{gr}^{+} = 2$.
\end{ejemplo}

\subsubsection*{Matriz de incidencia y $BB^{\mathsf{T}}$}

\begin{definicion}[Matriz de incidencia]
  Para un grafo dirigido sin lazos con las aristas numeradas
  $e_1, \dots, e_m$, la \textbf{matriz de incidencia} es la matriz
  $B$ de $n \times m$ con $b_{ij} = -1$ si $e_j$ sale del vértice $i$,
  $b_{ij} = 1$ si $e_j$ entra al vértice $i$, y $b_{ij} = 0$ si no.
  Los renglones son vértices y las columnas aristas.
\end{definicion}

\begin{proposicion}[Cormen, ejercicio 22.1-7]
  \label{prop:busq-incidencia}
  Para un grafo dirigido sin lazos,
  $(BB^{\mathsf{T}})_{ii} = \text{gr}^{-}(i) + \text{gr}^{+}(i)$, y
  para $i \neq j$, $(BB^{\mathsf{T}})_{ij}$ es menos el número de
  aristas entre $i$ y $j$, en cualquier sentido.
\end{proposicion}

\begin{proof}
  $(BB^{\mathsf{T}})_{ij} = \sum_{e} b_{ie} b_{je}$, el producto punto
  de los renglones $i$ y $j$. Si $i = j$, cada sumando es $b_{ie}^2$,
  que vale $1$ si $e$ toca a $i$ y $0$ si no: la suma cuenta las
  aristas que entran o salen de $i$. Si $i \neq j$, el producto
  $b_{ie} b_{je}$ es distinto de cero solo cuando $e$ une a $i$ con
  $j$; en ese caso uno de los dos vale $-1$ y el otro $1$, y el
  producto es $-1$.
\end{proof}

Dicho de otro modo, $BB^{\mathsf{T}} = D - (A + A^{\mathsf{T}})$, con
$D$ la diagonal de grados totales: es la \textbf{matriz laplaciana} del
grafo no dirigido subyacente, contando multiplicidades. Cada renglón
suma $0$, porque cada columna de $B$ tiene un $-1$ y un $1$. La
laplaciana es la base del teorema de Kirchhoff para contar árboles de
expansión; aquí basta reconocerla. La definición no admite lazos: la
columna de $(u, u)$ tendría que valer $-1$ y $1$ en la misma casilla.

\begin{ejemplo}[Incidencia del grafo dirigido de la clase]
  Es el grafo dirigido de la clase del 22 de octubre, el de Floyd, con
  sus siete aristas numeradas. Sus dos pares de aristas antiparalelas
  producen entradas $-2$.
  \begin{center}
  \begin{tikzpicture}[baseline=(current bounding box.center),
                      scale=0.85]
    \node[vert] (n1) at (0, 0) {1};
    \node[vert] (n2) at (2.6, 0) {2};
    \node[vert] (n3) at (2.6, -2.2) {3};
    \node[vert] (n4) at (0, -2.2) {4};
    \draw[dirigida] (n1) to[bend left=15]
      node[peso, above] {$e_1$} (n2);
    \draw[dirigida] (n2) to[bend left=15]
      node[peso, below] {$e_2$} (n1);
    \draw[dirigida] (n1) to[bend right=15]
      node[peso, left] {$e_3$} (n4);
    \draw[dirigida] (n4) to[bend right=15]
      node[peso, right] {$e_4$} (n1);
    \draw[dirigida] (n2) -- node[peso, right] {$e_5$} (n3);
    \draw[dirigida] (n3) -- node[peso] {$e_6$} (n1);
    \draw[dirigida] (n3) -- node[peso, below] {$e_7$} (n4);
  \end{tikzpicture}
  \hspace{0.4cm}
  $B = \begin{array}{c|rrrrrrr}
      & e_1 & e_2 & e_3 & e_4 & e_5 & e_6 & e_7 \\ \hline
    1 & -1 &  1 & -1 &  1 &  0 &  1 &  0 \\
    2 &  1 & -1 &  0 &  0 & -1 &  0 &  0 \\
    3 &  0 &  0 &  0 &  0 &  1 & -1 & -1 \\
    4 &  0 &  0 &  1 & -1 &  0 &  0 &  1
  \end{array}$
  \hspace{0.3cm}
  $BB^{\mathsf{T}} =
  \begin{pmatrix*}[r]
     5 & -2 & -1 & -2 \\
    -2 &  3 & -1 &  0 \\
    -1 & -1 &  3 & -1 \\
    -2 &  0 & -1 &  3
  \end{pmatrix*}$
  \end{center}
  El $5$ de la esquina son las cinco aristas que tocan a $1$:
  $\text{gr}^{-}(1) = 3$ más $\text{gr}^{+}(1) = 2$. Hay $-2$ en
  $(1,2)$ porque entre $1$ y $2$ hay dos aristas, una en cada
  sentido; $-1$ en $(2,3)$ por la única arista $e_5$; $0$ en $(2,4)$
  porque no hay ninguna.
\end{ejemplo}

\subsubsection*{De multigrafo a grafo simple}

Un multigrafo no dirigido admite aristas repetidas y lazos. Su grafo
simple equivalente $G' = (V, E')$ deja una sola arista por cada grupo
de repetidas y quita los lazos (Cormen, ejercicio 22.1-4). El truco es
el mismo de la marca del cuadrado.

\begin{algorithm}[H]
\caption{\textsc{Simplificar-Multigrafo}$(G)$}
\label{alg:busq-simplificar}
\begin{algorithmic}[1]
\Procedure{Simplificar-Multigrafo}{$G$}
  \ForAll{$u \in V$}
    \State $\mathit{Adj}'[u] \gets$ lista vacía, \quad
           $\text{último}[u] \gets \textsc{nulo}$
  \EndFor
  \ForAll{$u \in V$}
    \ForAll{$v \in \mathit{Adj}[u]$}
      \If{$v \neq u$ \textbf{y} $\text{último}[v] \neq u$}
        \State $\text{último}[v] \gets u$
        \State agregar $v$ a $\mathit{Adj}'[u]$
      \EndIf
    \EndFor
  \EndFor
  \State \Return $\mathit{Adj}'$
\EndProcedure
\end{algorithmic}
\end{algorithm}

En la línea 7, $v \neq u$ quita los lazos y $\text{último}[v] \neq u$
quita las repetidas: $\text{último}[v] = u$ quiere decir que $v$ ya
entró a $\mathit{Adj}'[u]$. La salida es simétrica sin hacer nada
especial, porque si $v$ entra a $\mathit{Adj}'[u]$, al procesar $v$ se
encontrará a $u$ en $\mathit{Adj}[v]$. Por ejemplo, con dos aristas
entre $1$ y $2$, un lazo en $1$ y $\mathit{Adj}[1] = 2, 1, 2, 3$, el
primer $2$ entra y deja $\text{último}[2] = 1$, el lazo se omite, el
segundo $2$ se descarta y el $3$ entra con $\text{último}[3] = 1$.
Al pasar a $u = 2$ con $\mathit{Adj}[2] = 1, 3, 1$, esa marca es
vieja, distinta de $2$, y el $3$ entra sin haber reinicializado nada.
El arreglo se inicializa una sola vez y el doble ciclo hace una vuelta
$\Ord{1}$ por entrada de las listas: $\Tht{n + m}$.

% =====================================================================

\subsection*{Búsqueda en anchura a fondo}

La clase del 22 de octubre (sección~\ref{sec:grafos}) dio BFS con un
arreglo booleano $\text{visto}$ y lo corrió sobre $G$ desde A. Aquí se
usa la versión de Cormen, que pinta cada vértice de uno de tres
colores. El resultado es el mismo; los colores sirven para las
pruebas.

\begin{definicion}[Colores de BFS]
  Un vértice es \textbf{blanco} mientras no se ha descubierto,
  \textbf{gris} desde que se descubre hasta que se termina de recorrer
  su lista, y \textbf{negro} después. Los grises forman la
  \textbf{frontera} entre lo descubierto y lo que falta.
\end{definicion}

\begin{algorithm}[H]
\caption{\textsc{BFS}$(G, s)$ con colores (Cormen, sección 22.2)}
\label{alg:bfs-colores}
\begin{algorithmic}[1]
\Procedure{BFS}{$G, s$}
  \ForAll{$u \in V - \set{s}$}
    \State $\text{color}[u] \gets \textsc{blanco}$, \quad
           $d[u] \gets \infty$, \quad $\pi[u] \gets \textsc{nulo}$
  \EndFor
  \State $\text{color}[s] \gets \textsc{gris}$, \quad $d[s] \gets 0$,
         \quad $\pi[s] \gets \textsc{nulo}$
  \State $Q \gets$ cola vacía; \ \textsc{Encolar}$(Q, s)$
  \While{$Q$ no está vacía}
    \State $u \gets \textsc{Desencolar}(Q)$
    \ForAll{$v \in \mathit{Adj}[u]$}
      \If{$\text{color}[v] = \textsc{blanco}$} \Comment{se descubre $v$}
        \State $\text{color}[v] \gets \textsc{gris}$, \quad
               $d[v] \gets d[u] + 1$, \quad $\pi[v] \gets u$
        \State \textsc{Encolar}$(Q, v)$
      \EndIf
    \EndFor
    \State $\text{color}[u] \gets \textsc{negro}$
      \Comment{lista de $u$ agotada}
  \EndWhile
\EndProcedure
\end{algorithmic}
\end{algorithm}

Tres líneas merecen comentario. La línea 10 solo pregunta si $v$ es
blanco: gris y negro se tratan igual, así que un bit por vértice
basta, y ese bit es el $\text{visto}$ de la clase (Cormen, ejercicio
22.2-3). La línea 11 pinta de gris y fija $d$ y $\pi$ en el instante
en que $v$ se descubre, no cuando sale de la cola; eso es lo que
garantiza que cada vértice entre a $Q$ una sola vez. La línea 15 no
cambia el resultado: existe para el invariante siguiente, que da
sentido a la palabra frontera.

\begin{observacion}
  Al principio de cada vuelta del ciclo de la línea 7, $Q$ es
  exactamente el conjunto de vértices grises. Al inicio, $s$ es el
  único gris y el único en $Q$; cada vez que un vértice se pinta de
  gris (línea 11) también se encola (línea 12), y el que se desencola
  en la línea 8 se pinta de negro en la 15.
\end{observacion}

\subsubsection*{Costo: $\Ord{n + m}$, y por qué no $\Tht{n + m}$}

El análisis es agregado: en lugar de acotar cada vuelta del ciclo, se
cuenta cuántas veces puede pasar cada cosa en toda la ejecución.
\begin{enumerate}[itemsep=0.15em]
  \item \textit{Inicialización} (líneas 2 a 6): $\Tht{n}$.
  \item \textit{Operaciones de cola.} Un vértice entra a $Q$ solo en
    la línea 12, solo si era blanco en la 10, y la 11 lo pinta de gris
    en ese mismo momento; ninguna línea lo regresa a blanco. Así que
    cada vértice se encola a lo más una vez, contando $s$, y como cada
    \textsc{Desencolar} saca algo que entró, también se desencola a lo
    más una vez. Con una cola de arreglo circular o de lista ligada,
    cada operación es $\Ord{1}$: $\Ord{n}$ en total.
  \item \textit{Recorrido de listas.} El ciclo de la línea 9 corre
    para $u$ solo cuando $u$ sale de la cola, a lo más una vez por
    vértice, y da $\abs{\mathit{Adj}[u]}$ vueltas de costo $\Ord{1}$.
    Por el lema~\ref{lem:busq-longitud-listas}, todas las listas
    suman $m$ o $2m$: $\Ord{m}$.
  \item \textit{Suma:} $\Tht{n} + \Ord{n} + \Ord{m} = \Ord{n + m}$,
    lineal en el tamaño de la representación.
\end{enumerate}

Por eso la clase escribe $\Ord{n + m}$ y no $\Tht{n + m}$: el
segundo es el costo del peor caso, que se alcanza cuando todo vértice
es alcanzable desde $s$, pero no describe cada ejecución. Si $s$ no
tiene aristas de salida, BFS inicializa los $n$ vértices, saca a
$s$, mira una lista vacía y termina en $\Tht{n}$, aunque $m$ sea del
orden de $n^2$. El costo exacto de una ejecución es
\[
  \Tht{n + \sum_{u \text{ alcanzable desde } s}
               \abs{\mathit{Adj}[u]}} ,
\]
porque las listas de los vértices inalcanzables nunca se leen. Por eso
Cormen escribe $O(V + E)$: una cota superior que vale para toda
entrada. Afirmar $\Tht{n + m}$ sin más promete también una cota
inferior, y esa falla en cuanto hay vértices inalcanzables con listas
largas. En DFS no pasa, como se verá: su ciclo externo obliga a leer
todas las listas.

Cada término depende de una estructura. El $\Ord{n}$ de la cola supone
operaciones en $\Ord{1}$; con una lista de Python y \texttt{pop(0)},
que cuesta lo que mide la lista, BFS sube a $\Ord{n^2 + m}$. El
$\Ord{m}$ supone listas: con matriz, cada vértice que sale de la cola
lee su renglón completo, $n$ casillas, y el total es $n$ por el número
de alcanzables, $\Ord{n^2}$. Y si la cola se
cambia por un montículo con llave $d$, sale Dijkstra con pesos
unitarios, $\Ord{(n + m) \log n}$: más caro y sin ganancia, porque,
como se prueba abajo, la cola ya entrega los vértices en orden de $d$.

\subsubsection*{Distancias: por qué BFS las calcula}

\begin{definicion}[Distancia de camino más corto]
  $\delta(s, v)$ es el mínimo número de aristas de un camino de $s$ a
  $v$, e $\infty$ si no hay ninguno. Un \textbf{camino más corto} de
  $s$ a $v$ es uno de longitud $\delta(s, v)$.
\end{definicion}

Aquí todas las aristas pesan $1$; la sección~\ref{sec:caminos-fuente}
generaliza a pesos reales. La caja de idea de la clase afirma que BFS
deja $d[v] = \delta(s, v)$. La prueba de Cormen tiene tres pasos: $d$
nunca subestima, la cola está ordenada por $d$, y un vértice con la
$d$ equivocada lleva a una contradicción.

\begin{lema}[Cormen, lema 22.1]
  \label{lem:busq-triangulo}
  Sea $G$ dirigido o no y $s \in V$. Para toda arista $(u, v) \in E$,
  $\delta(s, v) \leq \delta(s, u) + 1$.
\end{lema}

\begin{proof}
  Si $u$ es alcanzable, un camino más corto a $u$ seguido de la arista
  $(u, v)$ llega a $v$ con $\delta(s, u) + 1$ aristas, y el más corto
  a $v$ no puede ser más largo. Si $u$ no es alcanzable, el lado
  derecho es $\infty$.
\end{proof}

En un dirigido la desigualdad va en un solo sentido: una arista
$(f, s)$ puede ir de un vértice lejano a la fuente sin que eso acote
nada en sentido contrario. En uno no dirigido cada arista sirve en
los dos sentidos, y queda
$\abs{\delta(s, u) - \delta(s, v)} \leq 1$ para toda arista
$\set{u, v}$.

\begin{lema}[Cormen, lema 22.2]
  \label{lem:busq-cota-d}
  Al terminar BFS desde $s$, $d[v] \geq \delta(s, v)$ para todo
  $v \in V$.
\end{lema}

\begin{proof}
  Inducción sobre el número de operaciones \textsc{Encolar}, con la
  hipótesis ``$d[v] \geq \delta(s, v)$ para todo $v$''. Tras encolar
  a $s$, $d[s] = 0 = \delta(s, s)$ y los demás valen $\infty$. Si un
  $v$ blanco se descubre desde $u$, entonces
  $d[v] = d[u] + 1 \geq \delta(s, u) + 1 \geq \delta(s, v)$, por la
  hipótesis y el lema~\ref{lem:busq-triangulo}. Como $v$ queda gris,
  la línea 11 no vuelve a tocarlo y su $d$ ya no cambia.
\end{proof}

\begin{lema}[Cormen, lema 22.3]
  \label{lem:busq-cola}
  Si en algún momento $Q = \langle v_1, \dots, v_r \rangle$, con $v_1$
  a la cabeza, entonces $d[v_r] \leq d[v_1] + 1$ y
  $d[v_i] \leq d[v_{i+1}]$ para $i = 1, \dots, r - 1$.
\end{lema}

\begin{proof}[Idea de la prueba]
  Inducción sobre las operaciones de cola. Al desencolar $v_1$, la
  nueva cabeza es $v_2$ y $d[v_r] \leq d[v_1] + 1 \leq d[v_2] + 1$.
  Al encolar $v$ desde $u$, $u$ acaba de salir como cabeza, así que
  $d[u] \leq d[v_1]$ para la cabeza actual $v_1$, y
  $d[v] = d[u] + 1 \leq d[v_1] + 1$; además
  $d[v_r] \leq d[u] + 1 = d[v]$. Si la cola quedó vacía al sacar $u$,
  $v$ queda solo y no hay nada que probar.
\end{proof}

La consecuencia que conviene recordar: la cola contiene a lo más dos
valores distintos de $d$, consecutivos y en orden no decreciente. La
columna ``$d$ en $Q$'' de la traza del
ejemplo~\ref{ej:busq-fig223} lo muestra paso a paso.

\begin{corolario}[Cormen, corolario 22.4]
  \label{cor:busq-orden-cola}
  Si $v_i$ se encola antes que $v_j$, entonces $d[v_i] \leq d[v_j]$
  en el momento en que $v_j$ se encola.
\end{corolario}

Es inmediato del lema~\ref{lem:busq-cola} y de que cada vértice recibe
un valor finito de $d$ una sola vez.

\begin{teorema}[Corrección de BFS; Cormen, teorema 22.5]
  \label{teo:busq-correccion-bfs}
  Sea $G$ dirigido o no y supóngase que BFS corre desde $s$. Entonces
  descubre todo vértice alcanzable desde $s$ y, al terminar,
  $d[v] = \delta(s, v)$ para todo $v \in V$. Además, para todo
  $v \neq s$ alcanzable, un camino más corto de $s$ a $\pi[v]$
  seguido de la arista $(\pi[v], v)$ es un camino más corto de $s$ a
  $v$.
\end{teorema}

\begin{proof}
  Por contradicción. Supóngase que algún vértice termina con
  $d[v] \neq \delta(s, v)$, y de todos ellos tómese uno con
  $\delta(s, v)$ mínima. No es $s$, porque $d[s] = 0 = \delta(s, s)$.
  Por el lema~\ref{lem:busq-cota-d}, $d[v] > \delta(s, v)$, así que
  $\delta(s, v)$ es finita: $v$ es alcanzable. Sea $u$ el vértice
  anterior a $v$ en un camino más corto de $s$ a $v$; entonces
  $\delta(s, v) = \delta(s, u) + 1$, y como $\delta(s, u)$ es menor
  que $\delta(s, v)$, la elección de $v$ obliga a que
  $d[u] = \delta(s, u)$. Juntando,
  \begin{equation}
    d[v] > \delta(s, v) = \delta(s, u) + 1 = d[u] + 1 .
    \tag{$\ast$} \label{eq:busq-contradiccion}
  \end{equation}
  Como $d[u]$ es finita, $u$ pasó por la cola. Considérese el instante
  en que $u$ sale de ella. En ese momento $v$ es blanco, gris o negro.
  \begin{itemize}[itemsep=0.1em]
    \item Blanco: la línea 11 fija $d[v] = d[u] + 1$.
    \item Negro: $v$ ya salió de la cola antes que $u$, así que entró
      antes, y por el corolario~\ref{cor:busq-orden-cola},
      $d[v] \leq d[u]$.
    \item Gris: $v$ se descubrió al sacar algún $w$ que salió antes
      que $u$, así que $d[v] = d[w] + 1$ con $d[w] \leq d[u]$ por el
      mismo corolario, y $d[v] \leq d[u] + 1$.
  \end{itemize}
  Los tres casos contradicen~\eqref{eq:busq-contradiccion}. Entonces
  $d[v] = \delta(s, v)$ para todo $v$. Un alcanzable no descubierto
  tendría $d[v] = \infty > \delta(s, v)$, así que todos se descubren.
  Por último, si $\pi[v] = u$, la línea 11 dio
  $d[v] = d[u] + 1$, es decir, $\delta(s, v) = \delta(s, u) + 1$, y
  un camino más corto a $u$ más la arista $(u, v)$ mide exactamente
  eso.
\end{proof}

El caso gris es el que obliga a la prueba a separar colores, y
ocurre de verdad: en la figura 22.3, cuando $t$ revisa su lista
encuentra a $x$ gris, ya descubierto por $w$.

\subsubsection*{El árbol BFS y la impresión de caminos}

\begin{definicion}[Subgrafo predecesor de BFS]
  El \textbf{subgrafo predecesor} de BFS desde $s$ es
  $G_\pi = (V_\pi, E_\pi)$ con
  $V_\pi = \set{v : \pi[v] \neq \text{nulo}} \cup \set{s}$ y
  $E_\pi = \set{(\pi[v], v) : v \in V_\pi - \set{s}}$. Es un
  \textbf{árbol BFS} si $V_\pi$ son exactamente los vértices
  alcanzables desde $s$ y, para cada $v \in V_\pi$, el único camino de
  $s$ a $v$ en $G_\pi$ es un camino más corto en $G$.
\end{definicion}

\begin{lema}[Cormen, lema 22.6]
  \label{lem:busq-arbol-bfs}
  Aplicado a un grafo dirigido o no, BFS construye $\pi$ de modo que
  $G_\pi$ es un árbol BFS.
\end{lema}

\begin{proof}[Idea de la prueba]
  $\pi[v]$ recibe un valor distinto de nulo exactamente cuando
  $v \neq s$ es alcanzable, y entonces $(\pi[v], v) \in E$; por eso
  $V_\pi$ son los alcanzables. $G_\pi$ es conexo, porque de cada
  vértice se sube por $\pi$ hasta $s$, y tiene
  $\abs{E_\pi} = \abs{V_\pi} - 1$, así que es un árbol. A lo largo del
  único camino de $s$ a $v$ cada arista sube $d$ en exactamente $1$,
  y por el teorema~\ref{teo:busq-correccion-bfs} el camino mide
  $d[v] = \delta(s, v)$.
\end{proof}

Cormen escribe en esa prueba que $\pi[v] = u$ si y solo si
$(u, v) \in E$ y $\delta(s, v) < \infty$. Leída al pie de la letra es
falsa: en la figura 22.3, $(x, u) \in E$ y $u$ es alcanzable, pero
$\pi[u] = t$. Lo que quiere decir es lo que dice la idea de arriba.
Para leer el camino del árbol se sube por $\pi$ y se imprime al
regresar de la recursión:

\begin{algorithm}[H]
\caption{\textsc{Imprimir-Camino}$(G, s, v)$}
\label{alg:busq-imprimir-camino}
\begin{algorithmic}[1]
\Procedure{Imprimir-Camino}{$G, s, v$}
  \If{$v = s$}
    \State imprimir $s$
  \ElsIf{$\pi[v] = \textsc{nulo}$}
    \State imprimir ``no hay camino de $s$ a $v$''
  \Else
    \State \Call{Imprimir-Camino}{$G, s, \pi[v]$}
    \State imprimir $v$
  \EndIf
\EndProcedure
\end{algorithmic}
\end{algorithm}

Supone que BFS ya corrió desde la misma $s$. Si $v$ es alcanzable, la
cadena de padres termina en $s$ sin tocar nulo, así que la rama de la
línea 4 solo puede dispararse en la primera llamada. Hace
$\delta(s, v) + 1$ llamadas de costo $\Ord{1}$. Si los $\pi$ vinieran
de un BFS desde otra fuente $s'$, la cadena terminaría en $s'$ y el
mensaje ``no hay camino'' saldría a mitad de la salida.

\begin{observacion}
  Las distancias no dependen del orden de las listas, porque
  $d[v] = \delta(s, v)$ solo depende de $G$ (Cormen, ejercicio
  22.2-5). El árbol sí puede depender. En la figura 22.3, si la lista
  de $w$ es $s, x, t$ en lugar de $s, t, x$, sale $x$ antes que $t$ y
  descubre a $u$: queda $\pi[u] = x$. De los $41\,472$ órdenes
  posibles de las listas de ese grafo solo salen dos árboles, y solo
  difieren en $\pi[u]$.
\end{observacion}

\begin{ejemplo}[La figura 22.3 de Cormen]
  \label{ej:busq-fig223}
  El grafo es no dirigido, con $V = \set{r, s, t, u, v, w, x, y}$,
  diez aristas y fuente $s$. Cormen no imprime las listas; las que
  reproducen sus colas son $r$: $s, v$; $s$: $w, r$; $t$: $w, x, u$;
  $u$: $t, x, y$; $v$: $r$; $w$: $s, t, x$; $x$: $w, t, u, y$; $y$:
  $x, u$. La de $s$ no es alfabética, porque la instantánea (b) tiene
  $Q = \langle w, r \rangle$. Cada instantánea muestra el estado al
  principio de una vuelta del ciclo: blancos con $\infty$, grises y
  negros con su $d$, y la cola con la $d$ de cada vértice debajo; en
  azul, las aristas del árbol.
  \begin{center}
  \begin{tikzpicture}[x=1.05cm, y=1.05cm,
      cajaCola/.style={draw=gray!70, fill=gray!15, font=\scriptsize,
                       minimum size=0.42cm, inner sep=0pt}]
    \begin{scope}[xshift=0.0cm]
      \busqBaseBFS
      \node[font=\small] at (-0.8, 1.2) {(a)};
      \node[vertGris] at (s) {$0$};
      \node[cajaCola] at (0.45, -1.15) {$s$};
      \node[nota] at (0.45, -1.6) {0};
    \end{scope}
    \begin{scope}[xshift=5.5cm]
      \busqBaseBFS
      \node[font=\small] at (-0.8, 1.2) {(b)};
      \foreach \a/\b in {s/r, s/w}
        \draw[aristaArbol] (\a) -- (\b);
      \foreach \p/\d in {r/1, w/1}
        \node[vertGris] at (\p) {$\d$};
      \node[vertNegro] at (s) {$0$};
      \foreach \p/\d [count=\c] in {w/1, r/1} {
        \node[cajaCola] at (0.45*\c, -1.15) {$\p$};
        \node[nota] at (0.45*\c, -1.6) {\d};
      }
    \end{scope}
    \begin{scope}[xshift=11.0cm]
      \busqBaseBFS
      \node[font=\small] at (-0.8, 1.2) {(c)};
      \foreach \a/\b in {s/r, w/t, s/w, w/x}
        \draw[aristaArbol] (\a) -- (\b);
      \foreach \p/\d in {r/1, t/2, x/2}
        \node[vertGris] at (\p) {$\d$};
      \foreach \p/\d in {s/0, w/1}
        \node[vertNegro] at (\p) {$\d$};
      \foreach \p/\d [count=\c] in {r/1, t/2, x/2} {
        \node[cajaCola] at (0.45*\c, -1.15) {$\p$};
        \node[nota] at (0.45*\c, -1.6) {\d};
      }
    \end{scope}
  \end{tikzpicture}

  \vspace{0.25cm}

  \begin{tikzpicture}[x=1.05cm, y=1.05cm,
      cajaCola/.style={draw=gray!70, fill=gray!15, font=\scriptsize,
                       minimum size=0.42cm, inner sep=0pt}]
    \begin{scope}[xshift=0.0cm]
      \busqBaseBFS
      \node[font=\small] at (-0.8, 1.2) {(d)};
      \foreach \a/\b in {s/r, w/t, r/v, s/w, w/x}
        \draw[aristaArbol] (\a) -- (\b);
      \foreach \p/\d in {t/2, v/2, x/2}
        \node[vertGris] at (\p) {$\d$};
      \foreach \p/\d in {r/1, s/0, w/1}
        \node[vertNegro] at (\p) {$\d$};
      \foreach \p/\d [count=\c] in {t/2, x/2, v/2} {
        \node[cajaCola] at (0.45*\c, -1.15) {$\p$};
        \node[nota] at (0.45*\c, -1.6) {\d};
      }
    \end{scope}
    \begin{scope}[xshift=5.5cm]
      \busqBaseBFS
      \node[font=\small] at (-0.8, 1.2) {(e)};
      \foreach \a/\b in {s/r, w/t, t/u, r/v, s/w, w/x}
        \draw[aristaArbol] (\a) -- (\b);
      \foreach \p/\d in {u/3, v/2, x/2}
        \node[vertGris] at (\p) {$\d$};
      \foreach \p/\d in {r/1, s/0, t/2, w/1}
        \node[vertNegro] at (\p) {$\d$};
      \foreach \p/\d [count=\c] in {x/2, v/2, u/3} {
        \node[cajaCola] at (0.45*\c, -1.15) {$\p$};
        \node[nota] at (0.45*\c, -1.6) {\d};
      }
    \end{scope}
    \begin{scope}[xshift=11.0cm]
      \busqBaseBFS
      \node[font=\small] at (-0.8, 1.2) {(f)};
      \foreach \a/\b in {s/r, w/t, t/u, r/v, s/w, w/x, x/y}
        \draw[aristaArbol] (\a) -- (\b);
      \foreach \p/\d in {u/3, v/2, y/3}
        \node[vertGris] at (\p) {$\d$};
      \foreach \p/\d in {r/1, s/0, t/2, w/1, x/2}
        \node[vertNegro] at (\p) {$\d$};
      \foreach \p/\d [count=\c] in {v/2, u/3, y/3} {
        \node[cajaCola] at (0.45*\c, -1.15) {$\p$};
        \node[nota] at (0.45*\c, -1.6) {\d};
      }
    \end{scope}
  \end{tikzpicture}

  \vspace{0.25cm}

  \begin{tikzpicture}[x=1.05cm, y=1.05cm,
      cajaCola/.style={draw=gray!70, fill=gray!15, font=\scriptsize,
                       minimum size=0.42cm, inner sep=0pt}]
    \begin{scope}[xshift=0.0cm]
      \busqBaseBFS
      \node[font=\small] at (-0.8, 1.2) {(g)};
      \foreach \a/\b in {s/r, w/t, t/u, r/v, s/w, w/x, x/y}
        \draw[aristaArbol] (\a) -- (\b);
      \foreach \p/\d in {u/3, y/3}
        \node[vertGris] at (\p) {$\d$};
      \foreach \p/\d in {r/1, s/0, t/2, v/2, w/1, x/2}
        \node[vertNegro] at (\p) {$\d$};
      \foreach \p/\d [count=\c] in {u/3, y/3} {
        \node[cajaCola] at (0.45*\c, -1.15) {$\p$};
        \node[nota] at (0.45*\c, -1.6) {\d};
      }
    \end{scope}
    \begin{scope}[xshift=5.5cm]
      \busqBaseBFS
      \node[font=\small] at (-0.8, 1.2) {(h)};
      \foreach \a/\b in {s/r, w/t, t/u, r/v, s/w, w/x, x/y}
        \draw[aristaArbol] (\a) -- (\b);
      \node[vertGris] at (y) {$3$};
      \foreach \p/\d in {r/1, s/0, t/2, u/3, v/2, w/1, x/2}
        \node[vertNegro] at (\p) {$\d$};
      \node[cajaCola] at (0.45, -1.15) {$y$};
      \node[nota] at (0.45, -1.6) {3};
    \end{scope}
    \begin{scope}[xshift=11.0cm]
      \busqBaseBFS
      \node[font=\small] at (-0.8, 1.2) {(i)};
      \foreach \a/\b in {s/r, w/t, t/u, r/v, s/w, w/x, x/y}
        \draw[aristaArbol] (\a) -- (\b);
      \foreach \p/\d in {r/1, s/0, t/2, u/3, v/2, w/1, x/2, y/3}
        \node[vertNegro] at (\p) {$\d$};
      \node[font=\small, anchor=west] at (0.15, -1.15)
        {$= \emptyset$};
    \end{scope}
  \end{tikzpicture}
  \end{center}
  En la traza, B, G y N son el color del vecino en el momento en que
  se examina.
  \begin{center}
  \small
  \begin{tabular}{cllcll}
    \toprule
    \textbf{It.} & \textbf{$Q$ al inicio} & \textbf{$d$ en $Q$}
      & \textbf{Sale} & \textbf{Examina} & \textbf{Descubre} \\
    \midrule
    1 & $s$ & $0$ & $s$ & $w$\,B, $r$\,B
      & $w, r$ con $d = 1$, $\pi = s$ \\
    2 & $w, r$ & $1, 1$ & $w$ & $s$\,N, $t$\,B, $x$\,B
      & $t, x$ con $d = 2$, $\pi = w$ \\
    3 & $r, t, x$ & $1, 2, 2$ & $r$ & $s$\,N, $v$\,B
      & $v$ con $d = 2$, $\pi = r$ \\
    4 & $t, x, v$ & $2, 2, 2$ & $t$ & $w$\,N, $x$\,G, $u$\,B
      & $u$ con $d = 3$, $\pi = t$ \\
    5 & $x, v, u$ & $2, 2, 3$ & $x$ & $w$\,N, $t$\,N, $u$\,G, $y$\,B
      & $y$ con $d = 3$, $\pi = x$ \\
    6 & $v, u, y$ & $2, 3, 3$ & $v$ & $r$\,N & nada \\
    7 & $u, y$ & $3, 3$ & $u$ & $t$\,N, $x$\,N, $y$\,G & nada \\
    8 & $y$ & $3$ & $y$ & $x$\,N, $u$\,N & nada \\
    \bottomrule
  \end{tabular}
  \end{center}
  Quedan siete aristas de árbol, $n - 1$, y tres fuera: $t$--$x$,
  $x$--$u$ y $u$--$y$. Las tres unen vértices del mismo nivel o de
  niveles contiguos, como exige
  $\abs{\delta(s, u) - \delta(s, v)} \leq 1$.
  \textsc{Imprimir-Camino}$(G, s, y)$ imprime $s, w, x, y$. En las
  iteraciones 4, 5 y 7 aparece un vecino gris: son los casos que la
  prueba del teorema~\ref{teo:busq-correccion-bfs} trata aparte.
\end{ejemplo}

\begin{ejemplo}[Dirigido, con un vértice inalcanzable]
  \label{ej:busq-bfs-dirigido}
  $V = \set{s, a, b, c, e, f, z}$ y nueve aristas dirigidas:
  $(s,a)$, $(s,b)$, $(a,c)$, $(b,c)$, $(b,e)$, $(c,f)$, $(e,f)$,
  $(f,s)$ y $(z,s)$, con las listas en orden alfabético. No hay
  vértice $d$, para no confundirlo con el atributo. Tres instantáneas:
  tras sacar $s$, tras sacar $b$ y el estado final, con la $d$ junto
  a cada vértice.
  \begin{center}
  \begin{tikzpicture}[
      cajaCola/.style={draw=gray!70, fill=gray!15, font=\scriptsize,
                       minimum size=0.42cm, inner sep=0pt}]
    \foreach \p/\x/\y in {c/2.6/0.8, e/2.6/-0.8, f/3.9/0, z/0/-1.6}
      \node[vert, fill=white] (\p) at (\x, \y) {$\p$};
    \foreach \p/\x/\y in {a/1.3/0.8, b/1.3/-0.8}
      \node[vertGris] (\p) at (\x, \y) {$\p$};
    \node[vertNegro] (s) at (0, 0) {$s$};
    \foreach \a/\b in {a/c, b/c, b/e, c/f, e/f, z/s}
      \draw[dirigida] (\a) -- (\b);
    \draw[dirigida] (f) to[out=105, in=75, looseness=1.4] (s);
    \foreach \a/\b in {s/a, s/b}
      \draw[dirigidaArbol] (\a) -- (\b);
    \node[font=\small] at (1.0, -1.75) {$Q$};
    \foreach \p [count=\c] in {a, b}
      \node[cajaCola] at (1.05 + 0.45*\c, -1.75) {$\p$};
    \node[font=\small] at (1.95, -2.5) {(1) sale $s$};
  \end{tikzpicture}
  \hfill
  \begin{tikzpicture}[
      cajaCola/.style={draw=gray!70, fill=gray!15, font=\scriptsize,
                       minimum size=0.42cm, inner sep=0pt}]
    \foreach \p/\x/\y in {f/3.9/0, z/0/-1.6}
      \node[vert, fill=white] (\p) at (\x, \y) {$\p$};
    \foreach \p/\x/\y in {c/2.6/0.8, e/2.6/-0.8}
      \node[vertGris] (\p) at (\x, \y) {$\p$};
    \foreach \p/\x/\y in {s/0/0, a/1.3/0.8, b/1.3/-0.8}
      \node[vertNegro] (\p) at (\x, \y) {$\p$};
    \foreach \a/\b in {c/f, e/f, z/s}
      \draw[dirigida] (\a) -- (\b);
    \draw[dirigidaFuera] (b) -- (c);
    \draw[dirigida] (f) to[out=105, in=75, looseness=1.4] (s);
    \foreach \a/\b in {s/a, s/b, a/c, b/e}
      \draw[dirigidaArbol] (\a) -- (\b);
    \node[font=\small] at (1.0, -1.75) {$Q$};
    \foreach \p [count=\c] in {c, e}
      \node[cajaCola] at (1.05 + 0.45*\c, -1.75) {$\p$};
    \node[font=\small] at (1.95, -2.5) {(3) sale $b$};
  \end{tikzpicture}
  \hfill
  \begin{tikzpicture}
    \node[vert, fill=white, label={[nota]left:$\infty$}] (z)
      at (0, -1.6) {$z$};
    \foreach \p/\x/\y/\d/\lado in {s/0/0/0/left, a/1.3/0.8/1/above,
                                   b/1.3/-0.8/1/below,
                                   c/2.6/0.8/2/above,
                                   e/2.6/-0.8/2/below,
                                   f/3.9/0/3/right}
      \node[vertNegro, label={[nota]\lado:$\d$}] (\p) at (\x, \y)
        {$\p$};
    \foreach \a/\b in {b/c, e/f, z/s}
      \draw[dirigida] (\a) -- (\b);
    \draw[dirigida] (f) to[out=105, in=75, looseness=1.4] (s);
    \foreach \a/\b in {s/a, s/b, a/c, b/e, c/f}
      \draw[dirigidaArbol] (\a) -- (\b);
    \node[font=\small] at (1.5, -1.75) {$Q = \emptyset$};
    \node[font=\small] at (1.95, -2.5) {(6) final};
  \end{tikzpicture}
  \end{center}
  Tres cosas que este ejemplo enseña y el de la clase no:
  \begin{itemize}[itemsep=0.1em]
    \item $z$ tiene una arista \textit{hacia} $s$, pero ninguna llega
      a $z$: es inalcanzable. Termina blanco, con $d[z] = \infty$ y
      $\pi[z] = \text{nulo}$, y su lista nunca se lee: salen seis
      vértices y se examinan ocho de las nueve entradas de las
      listas. \textsc{Imprimir-Camino}$(G, s, z)$ imprime ``no hay
      camino de $s$ a $z$''.
    \item La arista $(f, s)$ baja del nivel $3$ al $0$. En un dirigido
      una arista puede regresar cualquier número de niveles; solo
      hacia adelante está acotada (lema~\ref{lem:busq-triangulo}).
    \item También aquí aparecen vecinos grises: al sacar $b$, $c$ ya
      está en la cola (punteada en la instantánea 3), y al sacar $e$,
      $f$ también. En la versión no dirigida, $z$ y $f$ quedarían a
      distancia $1$.
  \end{itemize}
\end{ejemplo}

\begin{ejercicio}[Cormen, ejercicio 22.2-6]
  Dar un grafo dirigido y un conjunto $E_\pi$ de aristas tal que, para
  todo $v$, el único camino de $s$ a $v$ en $E_\pi$ sea un camino más
  corto de $G$, pero que BFS no produzca $E_\pi$ con ningún orden de
  las listas.

  \textit{Solución.} Aristas $s \to u$, $s \to v$, $u \to w$,
  $u \to x$, $v \to w$, $v \to x$, y
  $E_\pi = \set{(s,u), (s,v), (u,w), (v,x)}$. Todo camino de $E_\pi$
  es más corto, pero el primero de $u$ y $v$ que sale de la cola
  descubre a $w$ y a $x$ a la vez, y los dos quedan con el mismo
  padre. No todo árbol de caminos más cortos es un árbol BFS.
\end{ejercicio}

\begin{ojo}
  \begin{itemize}[leftmargin=1.2em, itemsep=0.15em]
    \item Marcar al descubrir, no al sacar. Si el vértice se marca
      cuando sale de la cola, puede entrar varias veces, una por cada
      vecino que lo encuentre antes de que salga; la cola crece hasta
      $\Tht{m}$, se repite trabajo y, si $d$ y $\pi$ se sobrescriben
      al reencolar, se pierden los valores correctos.
    \item En un dirigido, alcanzable es en la dirección de las
      flechas: $z$ apunta a $s$ y es inalcanzable. La propiedad
      $\abs{d[u] - d[v]} \leq 1$ en cada arista es exclusiva del caso
      no dirigido.
    \item Un lazo no estorba. En un BFS desde $3$ sobre el grafo del
      ejemplo~\ref{ej:busq-dirigido-cormen}, al sacar a $6$ este se
      examina a sí mismo y se encuentra gris, no negro, porque se
      ennegrece hasta acabar su lista; como solo importa que no sea
      blanco, no pasa nada.
    \item Un BFS solo visita lo alcanzable desde $s$. Para recorrer
      todo el grafo hace falta un ciclo externo sobre los vértices,
      como el del algoritmo bipartito o el de DFS.
  \end{itemize}
\end{ojo}

\subsubsection*{Ciclos impares y grafos bipartitos}

En un grafo no dirigido, el lema~\ref{lem:busq-triangulo} dice que
toda arista une dos niveles iguales o contiguos del árbol BFS. Eso
afina la regla de ciclos de la clase. Si una arista fuera del árbol
une $u$ y $v$ con $d[u] = d[v] = k$ y $x$ es su ancestro común más
bajo, los dos caminos del árbol de $x$ a $u$ y a $v$ más la arista
cierran un ciclo de $2(k - d[x]) + 1$ aristas: impar. Si los niveles
difieren en $1$, el ciclo es par. En el grafo $G$ de la clase, la
arista E--F une dos vértices con $d = 2$, y el ciclo A, B, E, F, C, A
tiene cinco aristas: $G$ no es bipartito. Esta es exactamente la
prueba del algoritmo~\ref{alg:bipartito-bfs} de la
sección~\ref{sec:hamilton-arboles}, que colorea por la paridad de $d$
y también resuelve el problema de los luchadores buenos y malos de
Cormen (ejercicio 22.2-7) en $\Ord{n + r}$, con $r$ el número de
rivalidades.

\subsubsection*{Diámetro de un árbol con dos BFS}

El \textbf{diámetro} de un árbol $T$ es la mayor de todas las
distancias, $\max_{u, v} \delta(u, v)$. Calcularlo par por par
costaría $n$ BFS. En un árbol bastan dos, gracias a esta propiedad:

\begin{proposicion}
  \label{prop:busq-diametro}
  En un árbol, desde cualquier $s$, el vértice $a$ más lejano es
  extremo de algún diámetro.
\end{proposicion}

\begin{proof}
  Sea $p, q$ un par que realiza el diámetro y $P$ el camino entre
  ellos. Sea $x$ el vértice de $P$ más cercano a $s$; en un árbol todo
  camino de $s$ a un vértice de $P$ pasa por $x$. Sea $z$ el último
  vértice común de los caminos de $s$ a $a$ y de $s$ a $q$, el punto
  donde se separan. Hay dos casos.

  Si $z$ está en el tramo de $s$ a $x$, como $a$ es el más lejano,
  $\delta(s, z) + \delta(z, a) = \delta(s, a) \geq \delta(s, p) =
  \delta(s, z) + \delta(z, x) + \delta(x, p)$, así que
  $\delta(z, a) \geq \delta(z, x) + \delta(x, p)$. Entonces
  \[
    \delta(a, q) = \delta(a, z) + \delta(z, x) + \delta(x, q)
    \geq \delta(x, p) + \delta(x, q) = \delta(p, q) ,
  \]
  y el par $a, q$ también realiza el diámetro.

  Si $z$ está en $P$ después de $x$, en el tramo hacia $q$, entonces
  $\delta(s, z) + \delta(z, a) \geq \delta(s, z) + \delta(z, q)$, o
  sea $\delta(z, a) \geq \delta(z, q)$, y
  $\delta(p, a) = \delta(p, z) + \delta(z, a) \geq \delta(p, z) +
  \delta(z, q) = \delta(p, q)$: el par $p, a$ realiza el diámetro.
\end{proof}

Entonces un BFS desde cualquier $s$ da un extremo $a$ de un diámetro,
y un segundo BFS desde $a$ encuentra el otro extremo:

\begin{algorithm}[H]
\caption{\textsc{Diámetro}$(T)$ de un árbol con dos BFS}
\label{alg:busq-diametro}
\begin{algorithmic}[1]
\Procedure{Diámetro}{$T$}
  \State $s \gets$ cualquier vértice de $T$
  \State \Call{BFS}{$T, s$}; \ $a \gets$ un vértice con $d[a]$ máxima
  \State \Call{BFS}{$T, a$}; \ $b \gets$ un vértice con $d[b]$ máxima
  \State \Return $d[b]$
    \Comment{el camino sale de Imprimir-Camino$(T, a, b)$}
\EndProcedure
\end{algorithmic}
\end{algorithm}

Son dos BFS de $\Ord{n + m}$ con $m = n - 1$, más dos recorridos para
hallar el máximo: $\Ord{n}$.

Un solo BFS no basta, y el truco es exclusivo de
árboles: en $K_4$ con vértices $1$ a $4$ sin la arista $\set{3, 4}$,
desde $1$ todos están a distancia $1$; si se toma $2$ como el más
lejano, el segundo BFS vuelve a dar $1$, pero
$\delta(3, 4) = 2$. La prueba usó que el camino entre dos vértices es
único, y eso solo pasa en árboles.

\begin{ejemplo}[Diámetro de un árbol]
  El árbol tiene $V = \set{a, \dots, i}$ y aristas $i$--$a$,
  $a$--$b$, $b$--$c$, $b$--$d$, $d$--$e$, $e$--$f$, $d$--$g$,
  $g$--$h$. Junto a cada vértice, su $d$ en cada BFS; en verde, la
  fuente y el más lejano de cada uno.
  \begin{center}
  \begin{tikzpicture}
    \begin{scope}
      \foreach \p/\x/\y/\d/\lado in {i/0/0/1/above, a/1.1/0/0/above,
          b/2.2/0/1/above, c/2.2/-1/2/right, d/3.3/0/2/above,
          e/4.4/0.6/3/above, f/5.5/0.6/4/above, g/4.4/-0.6/3/below,
          h/5.5/-0.6/4/below}
        \node[vert, label={[nota]\lado:$\d$}] (\p) at (\x, \y) {$\p$};
      \foreach \a/\b in {i/a, a/b, b/c, b/d, d/e, e/f, d/g, g/h}
        \draw[arista] (\a) -- (\b);
      \node[vertOk] at (a) {$a$};
      \node[vertOk] at (f) {$f$};
      \node[font=\small] at (2.75, -1.75)
        {BFS desde $a$: el más lejano es $f$, a $4$};
    \end{scope}
    \begin{scope}[xshift=7.4cm]
      \foreach \p/\x/\y/\d/\lado in {i/0/0/5/above, a/1.1/0/4/above,
          b/2.2/0/3/above, c/2.2/-1/4/right, d/3.3/0/2/above,
          e/4.4/0.6/1/above, f/5.5/0.6/0/above, g/4.4/-0.6/3/below,
          h/5.5/-0.6/4/below}
        \node[vert, label={[nota]\lado:$\d$}] (\p) at (\x, \y) {$\p$};
      \foreach \a/\b in {b/c, d/g, g/h}
        \draw[arista] (\a) -- (\b);
      \foreach \a/\b in {f/e, e/d, d/b, b/a, a/i}
        \draw[aristaArbol] (\a) -- (\b);
      \node[vertOk] at (f) {$f$};
      \node[vertOk] at (i) {$i$};
      \node[font=\small] at (2.75, -1.75)
        {BFS desde $f$: el más lejano es $i$, a $5$};
    \end{scope}
  \end{tikzpicture}
  \end{center}
  Desde $a$ hay empate entre $f$ y $h$ a distancia $4$; cualquiera
  sirve. Desde $f$ el más lejano es $i$, a $5$: el diámetro es $5$ y
  el camino, en azul, es $f, e, d, b, a, i$. Desde $a$ el vértice más
  lejano estaba a $4$, así que un solo BFS se habría quedado corto.
  El procedimiento da $5$ desde cualquiera de los nueve vértices
  iniciales.
\end{ejemplo}

% =====================================================================

\subsection*{Búsqueda en profundidad a fondo}

La clase dio \textsc{DFS-Visitar} con $\text{visto}$, un reloj y los
tiempos $\text{entra}$ y $\text{sale}$, y la corrió sobre $G$ desde A.
Le faltan dos piezas que Cormen pone: un ciclo externo que vuelve a
arrancar desde cada vértice que siga sin descubrir, de modo que el
recorrido cubre todo el grafo y produce un bosque, y los tres colores,
que aquí sí cambian lo que se puede deducir. Un vértice es blanco
antes de $\text{entra}[u]$, gris entre $\text{entra}[u]$ y
$\text{sale}[u]$, y negro después. El reloj $t$ de la clase se llama
aquí $\mathit{tiempo}$, porque uno de los ejemplos tiene un vértice
$t$.

\begin{definicion}[Bosque DFS]
  El subgrafo predecesor de DFS es $G_\pi = (V, E_\pi)$ con
  $E_\pi = \set{(\pi[v], v) : v \in V,\ \pi[v] \neq \text{nulo}}$.
  Es un \textbf{bosque DFS}, formado por \textbf{árboles DFS}, uno
  por cada raíz que elige el ciclo externo; sus aristas son las
  \textbf{aristas de árbol}.
\end{definicion}

A diferencia del de BFS, $G_\pi$ contiene a todo $V$. Reproduce el
árbol de llamadas recursivas: $\pi[v] = u$ si y solo si
\textsc{DFS-Visitar}$(G, v)$ se llamó desde el ciclo de $u$. En cada
instante, los vértices grises son exactamente las llamadas activas,
del más viejo al más reciente, y forman una cadena de descendientes en
el bosque. De ahí sale la llave de casi todo lo que sigue: una arista
que llega a un vértice gris llega a un antecesor.

\begin{algorithm}[H]
\caption{\textsc{DFS}$(G)$ con colores y tiempos (Cormen, sección
  22.3)}
\label{alg:dfs-tiempos}
\begin{algorithmic}[1]
\Procedure{DFS}{$G$}
  \ForAll{$u \in V$}
    \State $\text{color}[u] \gets \textsc{blanco}$, \quad
           $\pi[u] \gets \textsc{nulo}$
  \EndFor
  \State $\mathit{tiempo} \gets 0$
  \ForAll{$u \in V$}
    \If{$\text{color}[u] = \textsc{blanco}$}
      \State \Call{DFS-Visitar}{$G, u$}
        \Comment{$u$ es raíz de un árbol nuevo}
    \EndIf
  \EndFor
\EndProcedure
\end{algorithmic}
\end{algorithm}

\begin{algorithm}[H]
\caption{\textsc{DFS-Visitar}$(G, u)$, la visita recursiva de
  \textsc{DFS}}
\label{alg:busq-dfs-visitar}
\begin{algorithmic}[1]
\Procedure{DFS-Visitar}{$G, u$}
  \State $\mathit{tiempo} \gets \mathit{tiempo} + 1$, \quad
         $\text{entra}[u] \gets \mathit{tiempo}$
  \State $\text{color}[u] \gets \textsc{gris}$
  \ForAll{$v \in \mathit{Adj}[u]$} \Comment{explorar la arista
    $(u, v)$}
    \If{$\text{color}[v] = \textsc{blanco}$}
      \State $\pi[v] \gets u$
      \State \Call{DFS-Visitar}{$G, v$}
    \EndIf
  \EndFor
  \State $\text{color}[u] \gets \textsc{negro}$ \Comment{$u$ terminó}
  \State $\mathit{tiempo} \gets \mathit{tiempo} + 1$, \quad
         $\text{sale}[u] \gets \mathit{tiempo}$
\EndProcedure
\end{algorithmic}
\end{algorithm}

El reloj sube antes de asignar, así que el primer descubrimiento es el
$1$. Hay exactamente un descubrimiento y una finalización por vértice,
de modo que los tiempos son los enteros de $1$ a $2n$, todos
distintos, y siempre $\text{entra}[u] < \text{sale}[u]$. Un vértice
gris o negro no se vuelve a visitar, pero la arista hacia él sí cuenta
como explorada: eso es lo que permite clasificarla más adelante.

\textbf{Costo: aquí sí es $\Tht{n + m}$.} \textsc{DFS-Visitar} se llama
exactamente una vez por vértice. A lo más una, porque solo se llama
sobre un vértice blanco y lo primero que hace es pintarlo de gris,
color del que nunca regresa. Al menos una, porque el ciclo externo
llama sobre todo vértice que siga blanco. Fuera de su ciclo, cada
llamada hace trabajo $\Ord{1}$: $\Tht{n}$. En la llamada sobre $u$, el
ciclo da $\abs{\mathit{Adj}[u]}$ vueltas de costo $\Ord{1}$, sin
contar la llamada recursiva, que ya se contó. Sumando con el
lema~\ref{lem:busq-longitud-listas} sale $\Tht{m}$, y el total es
$\Tht{n + m}$ en toda entrada, porque cada vértice y cada entrada de
cada lista se tocan. La cuenta no es la de BFS, que contaba entradas a
la cola; aquí se cuentan llamadas. Con matriz el total sube a
$\Tht{n^2}$. La memoria extra es $\Tht{n}$ más la pila de recursión,
que llega a profundidad $n$ cuando el árbol es un camino; en grafos
grandes eso desborda la pila del sistema (Python corta en $1000$
llamadas por omisión), y es la razón práctica de la versión
iterativa del final.

\begin{ejemplo}[La figura 22.4 de Cormen, paso a paso]
  \label{ej:busq-fig224}
  $V = \set{u, v, w, x, y, z}$ y las aristas $(u,v)$, $(u,x)$,
  $(v,y)$, $(w,y)$, $(w,z)$, $(x,v)$, $(y,x)$ y el lazo $(z,z)$, con
  listas $u$: $v, x$; $v$: $y$; $w$: $y, z$; $x$: $v$; $y$: $x$; $z$:
  $z$. El ciclo externo toma los vértices en orden alfabético. La
  columna de grises va del fondo al tope y es literalmente la pila de
  llamadas.
  \begin{center}
  \small
  \begin{tabular}{cllll}
    \toprule
    $\mathit{tiempo}$ & \textbf{Evento} & \textbf{Grises}
      & \textbf{Arista explorada} & \textbf{Color de $v$, tipo} \\
    \midrule
    1 & descubre $u$ (raíz) & $u$ & $(u, v)$ & blanco: árbol \\
    2 & descubre $v$ & $u, v$ & $(v, y)$ & blanco: árbol \\
    3 & descubre $y$ & $u, v, y$ & $(y, x)$ & blanco: árbol \\
    4 & descubre $x$ & $u, v, y, x$ & $(x, v)$ & gris: retroceso \\
    5 & termina $x$ & $u, v, y$ & & \\
    6 & termina $y$ & $u, v$ & & \\
    7 & termina $v$ & $u$ & $(u, x)$
      & negro, $1 < 4$: avance \\
    8 & termina $u$ & vacía & & \\
    \addlinespace
    9 & descubre $w$ (raíz) & $w$ & $(w, y)$
      & negro, $9 > 3$: cruce \\
      & & & $(w, z)$ & blanco: árbol \\
    10 & descubre $z$ & $w, z$ & $(z, z)$ & gris: retroceso \\
    11 & termina $z$ & $w$ & & \\
    12 & termina $w$ & vacía & & \\
    \bottomrule
  \end{tabular}
  \end{center}
  Tres instantáneas, en los tiempos $4$, $8$ y $12$; junto a cada
  vértice, $\text{entra}/\text{sale}$. Las aristas grises delgadas
  aún no se exploran; las punteadas llevan su tipo: R de retroceso, A
  de avance, C de cruce.
  \begin{center}
  \begin{tikzpicture}[vert/.append style={fill=white},
      ret/.style={dirigidaFuera, draw=red!70!black},
      ava/.style={dirigidaFuera, draw=green!45!black},
      etR/.style={peso, text=red!70!black},
      etA/.style={peso, text=green!45!black},
      etC/.style={peso, text=orange!80!black}]
    \begin{scope}
      \node[vertGris] (u) at (0, 0) {$u$};
      \node[vertGris] (v) at (1.6, 0) {$v$};
      \node[vert] (w) at (3.2, 0) {$w$};
      \node[vertGris] (x) at (0, -1.5) {$x$};
      \node[vertGris] (y) at (1.6, -1.5) {$y$};
      \node[vert] (z) at (3.2, -1.5) {$z$};
      \draw[dirigidaArbol] (u) -- (v);
      \draw[dirigidaArbol] (v) -- (y);
      \draw[dirigidaArbol] (y) -- (x);
      \draw[ret] (x) -- node[etR] {R} (v);
      \draw[dirigida] (u) -- (x);
      \draw[dirigida] (w) -- (y);
      \draw[dirigida] (w) -- (z);
      \draw[dirigida] (z) to[out=30, in=-30, looseness=7] (z);
      \foreach \n/\t in {u/{1/}, v/{2/}}
        \node[nota, above=1pt] at (\n.north) {\t};
      \foreach \n/\t in {x/{4/}, y/{3/}}
        \node[nota, below=1pt] at (\n.south) {\t};
      \node[font=\small] at (1.6, -2.55) {(a) $\mathit{tiempo} = 4$};
    \end{scope}
    \begin{scope}[xshift=5.2cm]
      \node[vertNegro] (u) at (0, 0) {$u$};
      \node[vertNegro] (v) at (1.6, 0) {$v$};
      \node[vert] (w) at (3.2, 0) {$w$};
      \node[vertNegro] (x) at (0, -1.5) {$x$};
      \node[vertNegro] (y) at (1.6, -1.5) {$y$};
      \node[vert] (z) at (3.2, -1.5) {$z$};
      \draw[dirigidaArbol] (u) -- (v);
      \draw[dirigidaArbol] (v) -- (y);
      \draw[dirigidaArbol] (y) -- (x);
      \draw[ret] (x) -- node[etR] {R} (v);
      \draw[ava] (u) -- node[etA] {A} (x);
      \draw[dirigida] (w) -- (y);
      \draw[dirigida] (w) -- (z);
      \draw[dirigida] (z) to[out=30, in=-30, looseness=7] (z);
      \foreach \n/\t in {u/{1/8}, v/{2/7}}
        \node[nota, above=1pt] at (\n.north) {\t};
      \foreach \n/\t in {x/{4/5}, y/{3/6}}
        \node[nota, below=1pt] at (\n.south) {\t};
      \node[font=\small] at (1.6, -2.55) {(b) $\mathit{tiempo} = 8$};
    \end{scope}
    \begin{scope}[xshift=10.4cm]
      \node[vertNegro] (u) at (0, 0) {$u$};
      \node[vertNegro] (v) at (1.6, 0) {$v$};
      \node[vertNegro] (w) at (3.2, 0) {$w$};
      \node[vertNegro] (x) at (0, -1.5) {$x$};
      \node[vertNegro] (y) at (1.6, -1.5) {$y$};
      \node[vertNegro] (z) at (3.2, -1.5) {$z$};
      \draw[dirigidaArbol] (u) -- (v);
      \draw[dirigidaArbol] (v) -- (y);
      \draw[dirigidaArbol] (y) -- (x);
      \draw[dirigidaArbol] (w) -- (z);
      \draw[ret] (x) -- node[etR] {R} (v);
      \draw[ava] (u) -- node[etA] {A} (x);
      \draw[dirigidaFuera] (w) -- node[etC] {C} (y);
      \draw[ret] (z) to[out=30, in=-30, looseness=7]
        node[etR, right=1pt] {R} (z);
      \foreach \n/\t in {u/{1/8}, v/{2/7}, w/{9/12}}
        \node[nota, above=1pt] at (\n.north) {\t};
      \foreach \n/\t in {x/{4/5}, y/{3/6}, z/{10/11}}
        \node[nota, below=1pt] at (\n.south) {\t};
      \node[font=\small] at (1.6, -2.55) {(c) $\mathit{tiempo} = 12$};
    \end{scope}
  \end{tikzpicture}
  \end{center}
  Quedan dos árboles, $\set{u, v, y, x}$ y $\set{w, z}$. El vértice
  $y$ queda en el primero aunque $w$ también le apunte: cuando $w$
  explora $(w, y)$, $y$ ya es negro. El ejemplo de la clase sobre $G$
  es una corrida del mismo algoritmo con una sola raíz, porque desde A
  se alcanza todo.
\end{ejemplo}

\subsubsection*{El teorema del paréntesis}

Si se escribe ``$(u$'' en el instante $\text{entra}[u]$ y ``$u)$'' en
$\text{sale}[u]$, la sucesión de eventos de una DFS es una expresión
de paréntesis bien balanceada. En el ejemplo~\ref{ej:busq-fig224} es
$(u\ (v\ (y\ (x\ x)\ y)\ v)\ u)\ (w\ (z\ z)\ w)$, y en el de la clase
sobre $G$,
$(A\ (B\ (D\ (G\ G)\ D)\ (E\ (F\ (C\ C)\ F)\ E)\ B)\ A)$. La razón es
que la recursión es una pila: una llamada que empieza dentro de otra
termina antes que ella.

\begin{teorema}[Del paréntesis; Cormen, teorema 22.7]
  \label{teo:busq-parentesis}
  En cualquier DFS de un grafo, dirigido o no, para cualesquiera dos
  vértices $u$ y $v$ se cumple exactamente una de estas tres
  condiciones:
  \begin{enumerate}[itemsep=0.1em]
    \item los intervalos $[\text{entra}[u], \text{sale}[u]]$ y
      $[\text{entra}[v], \text{sale}[v]]$ son disjuntos, y ninguno de
      los dos vértices desciende del otro en el bosque;
    \item el intervalo de $u$ está contenido en el de $v$, y $u$
      desciende de $v$;
    \item el intervalo de $v$ está contenido en el de $u$, y $v$
      desciende de $u$.
  \end{enumerate}
\end{teorema}

\begin{proof}
  Supóngase $\text{entra}[u] < \text{entra}[v]$; el otro caso es
  simétrico. Si $\text{entra}[v] < \text{sale}[u]$, $v$ se descubrió
  mientras $u$ era gris, es decir, durante la llamada sobre $u$, así
  que desciende de $u$. Como $v$ es más reciente, su llamada termina
  antes de que la búsqueda regrese a $u$:
  $\text{sale}[v] < \text{sale}[u]$, y el intervalo de $v$ queda dentro
  del de $u$. Si en cambio $\text{sale}[u] < \text{entra}[v]$, entonces
  $\text{entra}[u] < \text{sale}[u] < \text{entra}[v] < \text{sale}[v]$:
  los intervalos son disjuntos, y ninguno se descubrió con el otro
  gris, así que ninguno desciende del otro.
\end{proof}

Lo que el teorema prohíbe es el traslape parcial: dos intervalos se
anidan o se separan, como los paréntesis.

\begin{corolario}[Anidamiento; Cormen, corolario 22.8]
  \label{cor:busq-anidamiento}
  $v$ es descendiente propio de $u$ si y solo si
  $\text{entra}[u] < \text{entra}[v] < \text{sale}[v] < \text{sale}[u]$.
\end{corolario}

\begin{ejemplo}[La figura 22.5 de Cormen: los paréntesis como
  intervalos]
  \label{ej:busq-fig225}
  $V = \set{s, t, u, v, w, x, y, z}$ y trece aristas: $(s,z)$,
  $(s,w)$, $(t,v)$, $(t,u)$, $(u,t)$, $(u,v)$, $(v,s)$, $(v,w)$,
  $(w,x)$, $(x,z)$, $(y,x)$, $(z,y)$, $(z,w)$. Cormen no da las
  listas; estas reproducen sus tiempos: $s$: $z, w$; $t$: $v, u$;
  $u$: $t, v$; $v$: $s, w$; $w$: $x$; $x$: $z$; $y$: $x$; $z$: $y, w$,
  con $s$ primero y $t$ segundo en el ciclo externo. Cada vértice
  está dibujado como su intervalo, un renglón por nivel de
  profundidad, y abajo van los paréntesis en el instante en que se
  escriben.
  \begin{center}
  \begin{tikzpicture}[x=0.72cm, y=1cm]
    \foreach \k in {1,...,16} {
      \draw[gray!45, densely dotted] (\k, 0.45) -- (\k, -2.15);
      \node[font=\scriptsize] at (\k, -2.45) {\k};
    }
    \foreach \p/\a/\b/\n in {s/1/10/0, z/2/9/1, y/3/6/2, x/4/5/3,
                             w/7/8/2, t/11/16/0, v/12/13/1,
                             u/14/15/1} {
      \draw[draw=acento, fill=fondoIdea]
        (\a, -0.6*\n + 0.2) rectangle (\b, -0.6*\n - 0.2);
      \node[font=\small] at ({(\a + \b)/2}, -0.6*\n) {$\p$};
    }
    \foreach \k/\q in {1/{(s}, 2/{(z}, 3/{(y}, 4/{(x}, 5/{x)}, 6/{y)},
                       7/{(w}, 8/{w)}, 9/{z)}, 10/{s)}, 11/{(t},
                       12/{(v}, 13/{v)}, 14/{(u}, 15/{u)}, 16/{t)}}
      \node[font=\small] at (\k, -2.95) {$\q$};
    \node[nota, anchor=east] at (0.4, -2.45) {$\mathit{tiempo}$};
  \end{tikzpicture}
  \end{center}
  Un rectángulo debajo de otro y dentro de su sombra es un
  descendiente: $x$ dentro de $y$, dentro de $z$, dentro de $s$. Dos
  rectángulos sin traslape no tienen relación: $y$ y $w$ son hermanos,
  y $s$ y $t$ están en árboles distintos. Nunca hay traslape parcial.
  Los árboles son $s, z, y, x, w$, con $\pi[z] = s$,
  $\pi[y] = \pi[w] = z$ y $\pi[x] = y$, y $t, v, u$, con
  $\pi[v] = \pi[u] = t$.
\end{ejemplo}

\subsubsection*{El teorema del camino blanco}

El paréntesis dice cómo reconocer a un descendiente a partir de los
tiempos. El siguiente teorema dice cómo saber, en el instante en que
se descubre $u$, quiénes van a terminar colgando de él.

\begin{teorema}[Del camino blanco; Cormen, teorema 22.9]
  \label{teo:busq-camino-blanco}
  En un bosque DFS de un grafo dirigido o no, $v$ es descendiente de
  $u$ si y solo si, en el instante $\text{entra}[u]$, existe un camino
  de $u$ a $v$ formado solo por vértices blancos.
\end{teorema}

\begin{proof}
  ($\Rightarrow$) Si $v = u$, el camino es solo $u$, que es blanco
  justo antes de que se fije $\text{entra}[u]$. Si $v$ es descendiente
  propio, todos los vértices del camino del árbol de $u$ a $v$
  descienden de $u$, así que por el
  corolario~\ref{cor:busq-anidamiento} se descubren después de
  $\text{entra}[u]$: en ese instante eran blancos.

  ($\Leftarrow$) Supóngase que en $\text{entra}[u]$ hay un camino
  blanco de $u$ a $v$, pero $v$ no desciende de $u$. Sea $x$ el primer
  vértice del camino que no desciende de $u$, y $w$ el anterior a él
  en el camino; $w$ sí desciende de $u$ (puede ser $w = u$), así que
  $\text{sale}[w] \leq \text{sale}[u]$. Como $x$ era blanco en
  $\text{entra}[u]$, se descubre después; y como $w$ explora la arista
  $(w, x)$ antes de terminar, $x$ se descubre a más tardar en ese
  momento. Entonces
  \[
    \text{entra}[u] < \text{entra}[x] < \text{sale}[w] \leq
    \text{sale}[u] .
  \]
  El intervalo de $x$ empieza dentro del de $u$; por el
  teorema~\ref{teo:busq-parentesis} no puede traslaparse a medias, así
  que queda contenido, y por el corolario~\ref{cor:busq-anidamiento}
  $x$ desciende de $u$. Contradicción.
\end{proof}

La lectura útil: ser descendiente no depende de lo que pase después,
sino de qué estaba blanco al descubrir $u$. Un camino de $u$ a $v$ que
pase por un vértice ya gris o negro no obliga a nada.

\begin{ejercicio}[Cormen, ejercicios 22.3-8 y 22.3-9]
  Refutar dos conjeturas: ``si hay camino de $u$ a $v$ y
  $\text{entra}[u] < \text{entra}[v]$, entonces $v$ desciende de
  $u$'', y ``si hay camino de $u$ a $v$, entonces
  $\text{entra}[v] \leq \text{sale}[u]$''.

  \textit{Solución.} Un solo grafo tumba las dos. Aristas
  $(w, u)$, $(u, w)$ y $(w, v)$, listas $w$: $u, v$ y $u$: $w$, y el
  ciclo externo empezando en $w$.
  Los tiempos son $w$ 1/6, $u$ 2/3, $v$ 4/5. Hay camino
  $u \to w \to v$ y $\text{entra}[u] = 2 < \text{entra}[v] = 4$, pero
  $v$ es hermano de $u$, no descendiente. Y
  $\text{entra}[v] = 4 > \text{sale}[u] = 3$. El teorema del camino
  blanco no se aplica porque en el instante $2$ el vértice $w$ ya era
  gris: el camino existe, pero no es blanco.
\end{ejercicio}

\subsubsection*{Clasificación de aristas}

\begin{definicion}[Tipos de arista]
  Respecto del bosque DFS, una arista $(u, v)$ es
  \begin{enumerate}[itemsep=0.1em]
    \item \textbf{de árbol} si está en $E_\pi$: $v$ se descubrió al
      explorarla;
    \item \textbf{de retroceso} si une a $u$ con un antecesor $v$ en
      su árbol; los lazos cuentan como de retroceso;
    \item \textbf{de avance} si no es de árbol y une a $u$ con un
      descendiente $v$;
    \item \textbf{de cruce} en cualquier otro caso: entre vértices del
      mismo árbol sin relación de ancestría, o entre árboles
      distintos.
  \end{enumerate}
\end{definicion}

\begin{proposicion}[Clasificación por el color]
  \label{prop:busq-clasificacion}
  Cuando DFS explora $(u, v)$ por primera vez, el tipo lo decide el
  color de $v$: blanco, árbol; gris, retroceso; negro, avance si
  $\text{entra}[u] < \text{entra}[v]$ y cruce si
  $\text{entra}[u] > \text{entra}[v]$.
\end{proposicion}

\begin{proof}
  Blanco: $v$ se descubre justo ahora, desde $u$, por construcción.
  Gris: los grises son la cadena de llamadas activas y la exploración
  sale siempre del más reciente, que es $u$; cualquier otro gris es
  antecesor de $u$, y si $v = u$ es un lazo. Negro con
  $\text{entra}[u] < \text{entra}[v]$: $v$ se descubrió y terminó
  mientras $u$ era gris, así que desciende de $u$, y no es de árbol
  porque no estaba blanco. Negro con
  $\text{entra}[v] < \text{entra}[u]$: $u$ no puede descender de $v$,
  porque entonces $v$ seguiría gris; por el
  teorema~\ref{teo:busq-parentesis} los intervalos son disjuntos,
  $\text{sale}[v] < \text{entra}[u]$, y no hay relación.
\end{proof}

Para separar avance de cruce se compara $\text{entra}$, no
$\text{sale}$: cuando se explora la arista, $\text{sale}[u]$ todavía
no existe. Para clasificar al vuelo basta cambiar el \textbf{si} del
ciclo de \textsc{DFS-Visitar} por estos cuatro casos; el costo sigue
siendo $\Tht{n + m}$, porque cada arista se clasifica en $\Ord{1}$.

Con los intervalos ya calculados también se clasifica sin repetir la
búsqueda (Cormen, ejercicio 22.3-5).

\begin{proposicion}[Clasificación por intervalos]
  \label{prop:busq-intervalos}
  La arista $(u, v)$ es de árbol o de avance si y solo si
  $\text{entra}[u] < \text{entra}[v] < \text{sale}[v] < \text{sale}[u]$;
  de retroceso si y solo si
  $\text{entra}[v] \leq \text{entra}[u] < \text{sale}[u] \leq
  \text{sale}[v]$, con igualdad solo en un lazo; y de cruce si y solo
  si $\text{entra}[v] < \text{sale}[v] < \text{entra}[u] <
  \text{sale}[u]$. Árbol y avance se separan con $\pi$: es de árbol si
  y solo si $\pi[v] = u$.
\end{proposicion}

Falta a propósito un cuarto caso,
$\text{entra}[u] < \text{sale}[u] < \text{entra}[v]$: ninguna arista
sale de un intervalo hacia otro que empieza después de que el primero
cerró, porque al explorar $(u, v)$ el vértice $v$ habría estado blanco
y se habría vuelto descendiente de $u$. Por la misma razón nunca hay
una arista de un vértice negro a uno blanco: un vértice no termina
mientras le quede un vecino sin descubrir. La consecuencia que más se
usa es esta: \textbf{toda arista que no es de retroceso cumple
$\text{sale}[v] < \text{sale}[u]$}. Es la desigualdad que hace
funcionar el orden topológico de la sección~\ref{sec:aplicaciones-dfs}.
En el ejemplo~\ref{ej:busq-fig225}: $(x, z)$ va de $[4, 5]$ hacia
$[2, 9]$, que lo contiene, y es de retroceso; $(s, w)$ va de $[1, 10]$
hacia $[7, 8]$, contenido, con $\pi[w] = z \neq s$, y es de avance;
$(w, x)$ va de $[7, 8]$ hacia $[4, 5]$, que ya había cerrado, y es de
cruce.

\begin{ejemplo}[Los cuatro tipos; Cormen, figura 22.6 y ejercicio
  22.3-2]
  \label{ej:busq-cuatro-tipos}
  $V = \set{q, r, s, t, u, v, w, x, y, z}$ y catorce aristas: $(q,s)$,
  $(q,t)$, $(q,w)$, $(r,u)$, $(r,y)$, $(s,v)$, $(t,x)$, $(t,y)$,
  $(u,y)$, $(v,w)$, $(w,s)$, $(x,z)$, $(y,q)$, $(z,x)$. Como pide el
  ejercicio, el ciclo externo y todas las listas van en orden
  alfabético.
  \begin{center}
  \small
  \begin{tabular}{cccl@{\hspace{1cm}}cccl}
    \toprule
    \textbf{Arista} & $\mathit{tiempo}$ & \textbf{Color de $v$}
      & \textbf{Tipo} & \textbf{Arista} & $\mathit{tiempo}$
      & \textbf{Color de $v$} & \textbf{Tipo} \\
    \midrule
    $(q, s)$ & 1 & blanco & árbol & $(t, y)$ & 12 & blanco & árbol \\
    $(s, v)$ & 2 & blanco & árbol & $(y, q)$ & 13 & gris
      & retroceso \\
    $(v, w)$ & 3 & blanco & árbol & $(q, w)$ & 15 & negro
      & avance, $1 < 4$ \\
    $(w, s)$ & 4 & gris & retroceso & $(r, u)$ & 17 & blanco
      & árbol \\
    $(q, t)$ & 7 & blanco & árbol & $(u, y)$ & 18 & negro
      & cruce, $18 > 13$ \\
    $(t, x)$ & 8 & blanco & árbol & $(r, y)$ & 19 & negro
      & cruce, $17 > 13$ \\
    $(x, z)$ & 9 & blanco & árbol & & & & \\
    $(z, x)$ & 10 & gris & retroceso & & & & \\
    \bottomrule
  \end{tabular}
  \end{center}
  Tiempos finales: $q$ 1/16, $s$ 2/7, $v$ 3/6, $w$ 4/5, $t$ 8/15, $x$
  9/12, $z$ 10/11, $y$ 13/14, $r$ 17/20, $u$ 18/19. Ocho de árbol,
  tres de retroceso, una de avance y dos de cruce. A la izquierda el
  grafo; a la derecha el mismo grafo redibujado como bosque: las de
  árbol y la de avance bajan, las de retroceso suben y las de cruce
  van de derecha a izquierda, hacia subárboles que ya terminaron.
  \begin{center}
  \begin{tikzpicture}[scale=0.92, transform shape,
      ret/.style={dirigidaFuera, draw=red!70!black},
      ava/.style={dirigidaFuera, draw=green!45!black},
      etR/.style={peso, text=red!70!black},
      etA/.style={peso, text=green!45!black},
      etC/.style={peso, text=orange!80!black}]
    \begin{scope}
      \node[vert] (q) at (2, 0) {$q$};
      \node[vert] (s) at (0.8, -0.9) {$s$};
      \node[vert] (v) at (0, -1.8) {$v$};
      \node[vert] (w) at (1.6, -1.8) {$w$};
      \node[vert] (t) at (3.4, -0.9) {$t$};
      \node[vert] (x) at (3.4, -1.8) {$x$};
      \node[vert] (z) at (3.4, -2.8) {$z$};
      \node[vert] (y) at (5, -1.8) {$y$};
      \node[vert] (r) at (5.9, 0) {$r$};
      \node[vert] (u) at (6.8, -0.9) {$u$};
      \foreach \a/\b in {q/s, q/w, q/t, s/v, v/w, w/s, t/x, t/y, r/u,
                         r/y, u/y}
        \draw[dirigida] (\a) -- (\b);
      \draw[dirigida] (x) to[bend right=25] (z);
      \draw[dirigida] (z) to[bend right=25] (x);
      \draw[dirigida] (y) to[out=95, in=15] (q);
      \node[font=\small] at (3.4, -4.0) {(a) el grafo};
    \end{scope}
    \begin{scope}[xshift=9.8cm]
      \node[vertNegro] (q) at (0, 0) {$q$};
      \node[vertNegro] (s) at (-1.6, -1.0) {$s$};
      \node[vertNegro] (v) at (-2.4, -2.0) {$v$};
      \node[vertNegro] (w) at (-1.6, -3.0) {$w$};
      \node[vertNegro] (t) at (1.4, -1.0) {$t$};
      \node[vertNegro] (x) at (0.7, -2.0) {$x$};
      \node[vertNegro] (z) at (0.7, -3.0) {$z$};
      \node[vertNegro] (y) at (2.2, -2.0) {$y$};
      \node[vertNegro] (r) at (4.0, 0) {$r$};
      \node[vertNegro] (u) at (4.0, -1.0) {$u$};
      \foreach \a/\b in {q/s, s/v, v/w, q/t, t/x, x/z, t/y, r/u}
        \draw[dirigidaArbol] (\a) -- (\b);
      \draw[ret] (w) -- node[etR] {R} (s);
      \draw[ret] (z) to[out=30, in=-30] node[etR] {R} (x);
      \draw[ret] (y) to[out=80, in=10] node[etR] {R} (q);
      \draw[ava] (q) -- node[etA] {A} (w);
      \draw[dirigidaFuera] (u) -- node[etC] {C} (y);
      \draw[dirigidaFuera] (r) to[out=200, in=60] node[etC] {C} (y);
      \foreach \n/\t/\l in {q/{1/16}/above, s/{2/7}/left,
                            v/{3/6}/left, w/{4/5}/left,
                            t/{8/15}/left, x/{9/12}/left,
                            z/{10/11}/left, y/{13/14}/below,
                            r/{17/20}/above, u/{18/19}/right}
        \node[nota, \l=10pt] at (\n) {\t};
      \node[font=\small] at (0.8, -4.0) {(b) el bosque DFS};
    \end{scope}
  \end{tikzpicture}
  \end{center}
  Las tres de retroceso cierran tres ciclos: $s \to v \to w \to s$,
  $x \to z \to x$ y $q \to t \to y \to q$. Para saber que el grafo no
  es acíclico basta una; es el criterio que la
  sección~\ref{sec:aplicaciones-dfs} demuestra con el teorema del
  camino blanco.
\end{ejemplo}

\subsubsection*{Grafos no dirigidos: solo árbol y retroceso}

En un no dirigido, $\set{u, v}$ aparece en las dos listas y se explora
dos veces. Cormen la clasifica con el primero de los cuatro tipos que
le aplique, lo que equivale a clasificarla por la dirección en que se
explora primero (ejercicio 22.3-6). Con esa regla, dos de los tipos
desaparecen.

\begin{teorema}[Cormen, teorema 22.10]
  \label{teo:busq-no-dirigido}
  En una DFS de un grafo no dirigido, toda arista es de árbol o de
  retroceso.
\end{teorema}

\begin{proof}
  Sea $\set{u, v}$ una arista con, sin perder generalidad,
  $\text{entra}[u] < \text{entra}[v]$. Como $v$ está en la lista de
  $u$ y era blanco cuando se descubrió $u$, el teorema del camino
  blanco lo pone debajo de $u$: $v$ se descubre y termina mientras $u$
  es gris. Si la arista se explora primero de $u$ a $v$, $v$ seguía
  blanco en ese momento, porque si ya se hubiera descubierto, la
  llamada sobre $v$ habría explorado la arista desde su lado antes: es
  de árbol. Si se explora primero de $v$ a $u$, $u$ sigue gris: es de
  retroceso.
\end{proof}

El ejemplo de la clase lo ilustra. La arista $\set{A, C}$ se explora
primero desde C, en el tiempo $9$, con A gris: es de retroceso. Cuando
A la ve en el tiempo $13$, C ya es negro y no se reclasifica. Para
implementarlo, al explorar $\set{u, v}$ desde $u$: si $v$ es blanco,
árbol; si es gris y $v \neq \pi[u]$, retroceso; si $v = \pi[u]$, es la
arista de árbol vista desde abajo; si $v$ es negro, la arista ya se
exploró desde $v$. En los dos últimos casos no se hace nada. Es la
misma condición $v \neq \pi[u]$ de la regla de ciclos con BFS de la
clase.

\begin{ojo}
  Con aristas paralelas la condición $v \neq \pi[u]$ falla. En el
  multigrafo con dos aristas entre A y B y una entre B y C, DFS desde
  A baja a B por la primera copia. Desde B, la segunda copia lleva a
  A, que es gris pero es $\pi[B]$, y se descarta como si fuera la de
  árbol; cuando A la ve después, B ya es negro y también se descarta.
  El ciclo A, B, A de dos aristas nunca se reporta. La solución es
  comparar aristas y no vértices: guardar con cada vértice la arista
  por la que se llegó a él, y descartar solo esa. Con esa regla la
  segunda copia sale de retroceso, como debe.
\end{ojo}

\begin{ojo}
  \begin{itemize}[leftmargin=1.2em, itemsep=0.15em]
    \item El resultado depende del orden. En la figura 22.4, si el
      ciclo externo empieza por $w$, sale $w$ 1/10, $y$ 2/7, $x$ 3/6,
      $v$ 4/5, $z$ 8/9, $u$ 11/12, y la de retroceso es $(v, y)$ en
      lugar de $(x, v)$. Lo que no depende del orden es
      \textit{si hay} alguna de retroceso, porque eso equivale a
      tener un ciclo.
    \item Un bit basta para recorrer, no para clasificar. Si se quita
      la línea que pinta de negro, $\text{entra}$, $\text{sale}$ y
      $\pi$ salen idénticos (Cormen, ejercicio 22.3-4), así que el
      $\text{visto}$ de la clase alcanza para el recorrido. Pero sin
      distinguir gris de negro una de retroceso y una de cruce se ven
      igual, y la detección de ciclos dirigidos da falsos positivos,
      como la regla con BFS que la clase advierte.
    \item Un vértice puede quedar solo en su árbol aunque tenga
      aristas de entrada y de salida (ejercicio 22.3-11): con $(u, v)$
      y $(w, u)$ y el ciclo externo en el orden $v, u, w$, salen tres
      árboles de un vértice y dos aristas de cruce.
  \end{itemize}
\end{ojo}

\subsubsection*{DFS iterativo con pila}

La recursión de DFS llega a profundidad $n$, y en un grafo grande eso
desborda la pila del sistema. La clase menciona que la versión
iterativa usa una pila explícita, pero no la escribe, y hay dos formas
naturales de escribirla mal. Primero la versión corta, que sirve
cuando solo interesan el orden de visita y $\pi$.

\begin{algorithm}[H]
\caption{\textsc{DFS-Pila}$(G, s)$: orden de visita y padres}
\label{alg:busq-dfs-pila}
\begin{algorithmic}[1]
\Procedure{DFS-Pila}{$G, s$}
  \ForAll{$v \in V$}
    \State $\text{visto}[v] \gets$ falso
  \EndFor
  \State $P \gets$ pila vacía; \ \textsc{Apilar}$(P, (s, \textsc{nulo}))$
  \While{$P$ no está vacía}
    \State $(u, p) \gets \textsc{Desapilar}(P)$
    \If{\textbf{no} $\text{visto}[u]$} \Comment{si no, es copia vieja}
      \State $\text{visto}[u] \gets$ verdadero, \quad $\pi[u] \gets p$
      \ForAll{$v \in \mathit{Adj}[u]$, del último al primero}
        \If{\textbf{no} $\text{visto}[v]$}
          \State \textsc{Apilar}$(P, (v, u))$
        \EndIf
      \EndFor
    \EndIf
  \EndWhile
\EndProcedure
\end{algorithmic}
\end{algorithm}

Las dos decisiones que la hacen DFS son justo las dos trampas.

\textbf{Trampa 1: marcar al apilar.} Aquí se marca al \textit{sacar}
(línea 9), y un vértice puede quedar varias veces en la pila; las
copias viejas se descartan en la línea 8. La tentación es marcar al
apilar, como en BFS, para que nadie entre dos veces. Eso ya no es DFS:
el vecino se marca cuando lo apila el primero que lo ve, y queda
colgando de él aunque la búsqueda llegue después por otro lado. En el
triángulo $a, b, c$ produce una estrella con centro $a$,
$\pi[b] = \pi[c] = a$, y la arista $\set{b, c}$ sería de cruce, cosa
que el teorema~\ref{teo:busq-no-dirigido} prohíbe. En $G$ es más
traicionero: el orden de visita sale igual al de la recursión, A, B,
D, G, E, F, C, pero C queda con $\pi[C] = A$ en lugar de F, porque A
lo marcó al principio. El árbol es otro, y F--C une dos ramas
distintas.

\textbf{Trampa 2: apilar en el orden de la lista.} La pila invierte el
orden: lo último que se apila es lo primero que sale. Para visitar los
vecinos en el orden de la lista hay que apilarlos del último al
primero (línea 10). Si se apilan en el orden de la lista, el resultado
es una DFS válida, pero la de las listas al revés: en $G$ sale A, C,
F, E, B, D, G, y el árbol ya no es el de la clase.

Además, esta versión no da tiempos de salida, y su pila puede guardar
copias: cada apilado corresponde a una entrada de una lista revisada
desde un vértice recién visitado, así que hay a lo más $m + 1$
apilados en un dirigido y $2m + 1$ en uno no dirigido. Cada vértice se
visita una vez y recorre su lista una vez: $\Ord{n + m}$ de tiempo,
con $\Ord{m}$ de memoria para la pila.

\begin{ejemplo}[\textsc{DFS-Pila} sobre $G$ desde A]
  Con las listas en orden alfabético. En la pila, $C_A$ es el par
  $(C, A)$: el vértice $C$, apilado por $A$. Cada renglón saca el tope.
  \begin{center}
  \small
  \begin{tabular}{lll}
    \toprule
    \textbf{Saca} & \textbf{Acción} & \textbf{Pila después, fondo a
      tope} \\
    \midrule
    $A_{\text{nulo}}$ & visita A; apila C y B & $C_A, B_A$ \\
    $B_A$ & visita B; apila E y D & $C_A, E_B, D_B$ \\
    $D_B$ & visita D; apila G & $C_A, E_B, G_D$ \\
    $G_D$ & visita G; nada que apilar & $C_A, E_B$ \\
    $E_B$ & visita E; apila F & $C_A, F_E$ \\
    $F_E$ & visita F; apila C & $C_A, C_F$ \\
    $C_F$ & visita C, con $\pi[C] = F$ & $C_A$ \\
    $C_A$ & copia vieja: C ya se visitó & vacía \\
    \bottomrule
  \end{tabular}
  \end{center}
  Orden A, B, D, G, E, F, C, y los mismos padres que la recursión. Hubo
  ocho apilados y una copia vieja, $C_A$. Las tres versiones, lado a
  lado:
  \begin{center}
  \begin{tikzpicture}[scale=0.82, transform shape]
    \foreach \pre/\x in {p/0, q/5.6, r/11.2} {
      \begin{scope}[xshift=\x cm]
        \node[vert] (\pre A) at (0, 0) {A};
        \node[vert] (\pre B) at (-1.6, -1.2) {B};
        \node[vert] (\pre C) at (1.6, -1.2) {C};
        \node[vert] (\pre D) at (-2.6, -2.4) {D};
        \node[vert] (\pre E) at (-0.6, -2.4) {E};
        \node[vert] (\pre F) at (1.6, -2.4) {F};
        \node[vert] (\pre G) at (-2.6, -3.6) {G};
      \end{scope}
    }
    \foreach \a/\b in {A/B, B/D, D/G, B/E, E/F, F/C}
      \draw[aristaArbol] (p\a) -- (p\b);
    \draw[aristaFuera] (pC) -- (pA);
    \foreach \a/\b in {A/B, A/C, B/D, D/G, B/E, E/F}
      \draw[aristaArbol] (q\a) -- (q\b);
    \draw[aristaAlerta, dashed] (qF) -- (qC);
    \foreach \a/\b in {A/C, C/F, F/E, E/B, B/D, D/G}
      \draw[aristaArbol] (r\a) -- (r\b);
    \draw[aristaFuera] (rA) -- (rB);
    \node[font=\small, align=center] at (-0.5, -4.5)
      {(a) recursión o \textsc{DFS-Pila}:\\C--A de retroceso};
    \node[font=\small, align=center] at (5.1, -4.5)
      {(b) marcando al apilar:\\F--C sería de cruce};
    \node[font=\small, align=center] at (10.7, -4.5)
      {(c) apilando en orden:\\otra DFS};
  \end{tikzpicture}
  \end{center}
\end{ejemplo}

Para tener tiempos de salida hay que simular la recursión de verdad
(Cormen, ejercicio 22.3-7). La pila guarda solo vértices grises, y
cada uno lleva un apuntador $\text{sig}[u]$ al siguiente nodo por
revisar de su lista, que es justo lo que la recursión guarda sin que
se vea en la variable del ciclo. En cada vuelta se mira el tope $u$
sin sacarlo. Si $\text{sig}[u]$ no es nulo, se toma ese vecino $v$ y
se avanza el apuntador; si $v$ es blanco, recibe $\pi[v] = u$ y su
$\text{entra}[v]$, se pinta de gris y se apila. Si $\text{sig}[u]$ es
nulo, $u$ agotó su lista: se desapila, se pinta de negro y recibe su
$\text{sale}[u]$. El tope avanza una sola arista por vuelta, y el
ciclo externo es el mismo de \textsc{DFS}.

El invariante es que la pila contiene exactamente los vértices grises,
del más viejo en el fondo al más reciente en el tope, igual que la
pila de llamadas; por eso da los mismos $\text{entra}$, $\text{sale}$
y $\pi$ que la versión recursiva. Sobre $G$ desde A reproduce los
tiempos de la clase, de A 1/14 a C 9/10, con la pila llegando a cinco
vértices (A, B, E, F, C). Sin $\text{sig}[u]$, cada vez que $u$
regresa al tope habría que recorrer su lista desde el principio, y el
costo podría subir a $\Ord{nm}$. Con él, cada vértice entra y sale de
la pila una vez y cada vuelta del ciclo avanza un apuntador o saca un
vértice: $\Tht{n + m}$, con $\Ord{n}$ de pila.

\subsubsection*{Componentes conexas con DFS}

En un grafo no dirigido, cada árbol del bosque DFS es exactamente una
componente conexa (Cormen, ejercicio 22.3-12). Basta numerar los
árboles.

\begin{algorithm}[H]
\caption{Cambio a \textsc{DFS} para numerar las componentes}
\label{alg:busq-componentes}
\begin{algorithmic}[1]
\State $k \gets 0$ \Comment{antes del ciclo externo}
\ForAll{$u \in V$}
  \If{$\text{color}[u] = \textsc{blanco}$}
    \State $k \gets k + 1$ \Comment{árbol nuevo, componente nueva}
    \State \Call{DFS-Visitar}{$G, u$}
  \EndIf
\EndFor
\Statex En \textsc{DFS-Visitar}$(G, u)$, junto con pintar a $u$ de
  gris: $\text{cc}[u] \gets k$.
\end{algorithmic}
\end{algorithm}

\textbf{Por qué funciona.} Por inducción sobre los árboles, en el
orden en que se crean. Supóngase que los árboles anteriores son
componentes completas. La nueva raíz $u$ es blanca, así que no está en
ninguna de ellas, y su componente es disjunta de todas: en el
instante $\text{entra}[u]$ toda la componente de $u$ es blanca. Por el
teorema del camino blanco, todo vértice de esa componente termina
debajo de $u$. Y nada más, porque las aristas de árbol son aristas del
grafo y no hay aristas hacia otras componentes. El costo es el de una
DFS con una asignación más por vértice, $\Tht{n + m}$. Con $G$ de la
clase más una arista H--I y un vértice J aislado salen tres árboles y
tres componentes: $\set{A, \dots, G}$, $\set{H, I}$ y $\set{J}$. En un
dirigido los árboles DFS no son componentes de nada, y su número
depende del orden: las componentes fuertemente conexas necesitan el
algoritmo de dos DFS de la sección~\ref{sec:aplicaciones-dfs}.

\newpage

% =====================================================================
%  Grafos - Aplicaciones de DFS
% =====================================================================

\clase{Aplicaciones de DFS}{Complemento: Cormen, cap.\ 22.4, 22.5 y
  problemas}
\label{sec:aplicaciones-dfs}

\begin{quick_sum}
  El DFS de la sección~\ref{sec:busquedas} deja en cada vértice dos
  marcas de tiempo, $\text{entra}[v]$ y $\text{sale}[v]$, y esta
  sección muestra que esas marcas no son contabilidad del recorrido:
  son la estructura del grafo. Con ellas se resuelven, todos en
  $\Tht{n + m}$, cuatro problemas. El orden topológico de un DAG sale
  de listar los vértices por tiempo de salida decreciente, o del
  algoritmo de Kahn, que quita vértices sin predecesores. Las
  componentes fuertemente conexas de un grafo dirigido salen de dos
  DFS, el segundo sobre el grafo transpuesto. Los puntos de
  articulación, los puentes y las componentes biconexas de un grafo no
  dirigido salen de un solo DFS que calcula $\text{low}[v]$. Y la
  etiqueta mínima alcanzable desde cada vértice sale de un DFS sobre el
  transpuesto. Se cierra con el recorrido de Euler dirigido. Cada
  algoritmo trae su prueba, su costo y un ejemplo trabajado; la ropa
  del profesor, la figura 22.9 de Cormen y el grafo $G$ de la libreta
  son los grafos de trabajo.
\end{quick_sum}

\subsection*{Lo que se usa del DFS}

Todo lo que sigue corre sobre el DFS con colores y tiempos de la
sección~\ref{sec:busquedas} (algoritmo~\ref{alg:dfs-tiempos}): un reloj
global $t$ que sube en uno cada vez que un vértice se descubre o se
termina, los tiempos $\text{entra}[u]$ y $\text{sale}[u]$, el padre
$\pi[u]$ y el color, que es blanco antes de $\text{entra}[u]$, gris
entre $\text{entra}[u]$ y $\text{sale}[u]$ y negro después. En cada
instante los vértices grises son exactamente la pila de llamadas
activas, una cadena de antecesores en el árbol DFS. Las pruebas de esta
sección se apoyan en cuatro hechos de la sección~\ref{sec:busquedas}:

\begin{itemize}[itemsep=0.1em]
  \item \textbf{Paréntesis.} Para dos vértices cualesquiera, los
        intervalos $[\text{entra}[u], \text{sale}[u]]$ y
        $[\text{entra}[v], \text{sale}[v]]$ son ajenos o uno contiene
        al otro.
  \item \textbf{Descendientes.} $v$ es descendiente propio de $u$ si y
        solo si $\text{entra}[u] < \text{entra}[v] < \text{sale}[v] <
        \text{sale}[u]$.
  \item \textbf{Camino blanco.} $v$ es descendiente de $u$ si y solo
        si, en el instante $\text{entra}[u]$, hay un camino de $u$ a
        $v$ formado solo por vértices blancos.
  \item \textbf{Clasificación de aristas.} Al explorar $(u, v)$, la
        arista es de árbol si $v$ es blanco; de retroceso si $v$ es
        gris (va a un antecesor de $u$, y los lazos cuentan aquí); de
        avance si $v$ es negro y $\text{entra}[u] < \text{entra}[v]$;
        de cruce si $v$ es negro y $\text{entra}[v] <
        \text{entra}[u]$. En un grafo no dirigido solo hay aristas de
        árbol y de retroceso.
\end{itemize}

Hace falta una notación más. Para un conjunto de vértices $U$ se
escribe $\text{entra}(U) = \min_{u \in U} \text{entra}[u]$ y
$\text{sale}(U) = \max_{u \in U} \text{sale}[u]$: el primer
descubrimiento y la última salida dentro de $U$. Cormen los llama
$d(U)$ y $f(U)$. Con $u \leadsto v$ se indica que hay un camino de $u$
a $v$; el camino de longitud cero cuenta, así que $u \leadsto u$
siempre.

% =====================================================================

\subsection*{Orden topológico}

La clase del 22 de octubre afirma que un DAG siempre tiene un orden
topológico, el orden en que se pueden cursar las materias del plan de
estudios, pero no lo prueba. Aquí se prueba de dos maneras, y cada
prueba trae su algoritmo lineal.

\begin{definicion}[Orden topológico]
  \label{def:orden-topologico}
  Un \textbf{orden topológico} (\textit{topological sort}) de un grafo
  dirigido $G = (V, E)$ es un orden lineal de todos sus vértices tal
  que, si $(u, v) \in E$, entonces $u$ aparece antes que $v$.
\end{definicion}

Dibujados en fila según ese orden, todas las flechas apuntan a la
derecha. El orden no es único: en un grafo sin aristas cualquier
permutación sirve. Si el grafo tiene un ciclo no existe, porque cada
vértice del ciclo tendría que ir antes que su sucesor y, dando la
vuelta, antes que sí mismo. Las aplicaciones son las de siempre:
calendarizar tareas con precedencias, compilar módulos que dependen
unos de otros, recalcular las celdas de una hoja de cálculo.

\subsubsection*{Ciclos y aristas de retroceso}

\begin{lema}[Cormen, lema 22.11]
  \label{lem:aciclico-retroceso}
  Un grafo dirigido $G$ es acíclico si y solo si un DFS de $G$ no
  produce aristas de retroceso.
\end{lema}

\begin{proof}
  Si hay una arista de retroceso $(u, v)$, entonces $v$ es antecesor de
  $u$ en el bosque DFS, o $v = u$ si es un lazo. El camino de árbol de
  $v$ a $u$ seguido de la arista $(u, v)$ es un ciclo.

  Al revés, supóngase que $G$ tiene un ciclo $c$. Sea $v$ el primer
  vértice de $c$ que se descubre y sea $(u, v)$ la arista de $c$ que
  llega a $v$. En el instante $\text{entra}[v]$ los demás vértices de
  $c$ siguen blancos, así que el tramo de $c$ que va de $v$ a $u$ es un
  camino blanco. Por el teorema del camino blanco, $u$ queda como
  descendiente de $v$. Cuando el DFS explora $(u, v)$, desde dentro de
  la llamada de $u$, la llamada de $v$ sigue abierta: $v$ es gris y la
  arista es de retroceso.
\end{proof}

El lema vale para cualquier DFS, con cualquier orden de los vértices en
el ciclo exterior y en las listas. Basta una corrida para decidir si un
grafo dirigido tiene ciclo, y la arista de retroceso, junto con los
padres, entrega el ciclo mismo: se sube por $\pi$ desde $u$ hasta
llegar a $v$. Es la versión dirigida de la regla de ciclos con BFS de
la clase del 22 de octubre, que en dirigidos da falsos positivos.

\begin{ejercicio}[Cormen, ejercicio 22.4-3]
  Decidir si un grafo \emph{no dirigido} tiene un ciclo en tiempo
  $\Ord{n}$, sin importar cuántas aristas tenga.

  \textit{Solución.} Un DFS que se detiene en la primera arista de
  retroceso; la arista hacia el padre no cuenta, igual que en la regla
  de BFS. Mientras no aparezca una arista de retroceso, las aristas
  exploradas son de árbol y forman un bosque, que tiene a lo más
  $n - 1$ aristas. Cada una se examina dos veces, una desde cada
  extremo, así que se revisan a lo más $2(n - 1) + 1$ entradas de las
  listas antes de terminar o de encontrar el ciclo. En particular, si
  $m \ge n$ hay ciclo seguro.
\end{ejercicio}

\subsubsection*{El algoritmo por tiempos de salida}

La idea cabe en una línea: correr DFS y listar los vértices en orden
decreciente de $\text{sale}$. El teorema siguiente es la razón.

\begin{teorema}[Cormen, teorema 22.12]
  \label{teo:topologico-salida}
  Si $G$ es un DAG, en cualquier DFS de $G$ se cumple
  $\text{sale}[v] < \text{sale}[u]$ para toda arista $(u, v)$. En
  consecuencia, los vértices en orden decreciente de $\text{sale}$
  forman un orden topológico.
\end{teorema}

\begin{proof}
  Considérese el momento en que el DFS explora $(u, v)$. En ese momento
  $u$ es gris. Hay tres casos según el color de $v$.
  \begin{itemize}[itemsep=0.1em]
    \item $v$ gris: $v$ sería antecesor de $u$ y la arista sería de
          retroceso, lo que el lema~\ref{lem:aciclico-retroceso}
          descarta en un DAG.
    \item $v$ blanco: $v$ se vuelve descendiente de $u$ y, por la
          caracterización de descendientes,
          $\text{sale}[v] < \text{sale}[u]$.
    \item $v$ negro: $v$ ya terminó y $\text{sale}[v]$ ya tiene valor,
          mientras que $u$ todavía no termina. El reloj solo avanza,
          así que $\text{sale}[v] < \text{sale}[u]$.
  \end{itemize}
  En los dos casos posibles $u$ termina después que $v$, y queda antes
  que $v$ en el orden decreciente.
\end{proof}

No hace falta ordenar nada. Cada vértice que termina se inserta al
frente de una lista: el primero en terminar queda hasta el final y el
último en terminar queda al frente. El
algoritmo~\ref{alg:orden-topologico} es el DFS del
algoritmo~\ref{alg:dfs-tiempos} con dos agregados: esa inserción y la
prueba del vértice gris, que Cormen no pone y que convierte al
algoritmo en un detector de ciclos.

\begin{algorithm}[H]
\caption{\textsc{Orden-Topológico}$(G)$, con DFS y detección de
         ciclos}
\label{alg:orden-topologico}
\begin{algorithmic}[1]
\Procedure{Orden-Topológico}{$G$}
  \ForAll{$u \in V$}
    \State $\text{color}[u] \gets$ blanco, \quad $\pi[u] \gets$ nulo
  \EndFor
  \State $t \gets 0$; \ $L \gets$ lista vacía
  \ForAll{$u \in V$}
    \If{$\text{color}[u] =$ blanco}
      \State \textsc{Visitar-Orden}$(G, u)$
    \EndIf
  \EndFor
  \State \Return $L$
\EndProcedure
\Statex
\Procedure{Visitar-Orden}{$G, u$}
  \State $t \gets t + 1$, \quad $\text{entra}[u] \gets t$, \quad
         $\text{color}[u] \gets$ gris
  \ForAll{$v$ en la lista de adyacencia de $u$}
    \If{$\text{color}[v] =$ blanco}
      \State $\pi[v] \gets u$; \ \textsc{Visitar-Orden}$(G, v)$
    \ElsIf{$\text{color}[v] =$ gris} \Comment{arista de retroceso}
      \State \textbf{error}: hay un ciclo, no existe orden topológico
    \EndIf
  \EndFor
  \State $\text{color}[u] \gets$ negro, \quad $t \gets t + 1$, \quad
         $\text{sale}[u] \gets t$
  \State insertar $u$ al \textbf{frente} de $L$
\EndProcedure
\end{algorithmic}
\end{algorithm}

\begin{itemize}[itemsep=0.1em]
  \item \textbf{Línea 14.} El vértice se pinta de gris al descubrirlo
        y sigue gris mientras su llamada está en la pila.
  \item \textbf{Líneas 18 y 19.} Una arista hacia un vértice gris es
        de retroceso y, por el lema~\ref{lem:aciclico-retroceso}, hay
        ciclo. Sin estas dos líneas el algoritmo devuelve un orden de
        todos modos, sin avisar que es inválido.
  \item \textbf{Línea 23.} Insertar al frente equivale a ordenar por
        $\text{sale}$ decreciente sin ordenar.
  \item Las aristas hacia vértices negros, de avance o de cruce, se
        ignoran: esos vértices ya terminaron y quedarán a la derecha de
        $u$.
\end{itemize}

\textbf{Costo.} La inicialización cuesta $\Tht{n}$.
\textsc{Visitar-Orden} se llama exactamente una vez por vértice,
porque solo se llama sobre vértices blancos y lo primero que hace es
pintarlos de gris. Cada llamada recorre la lista de su vértice una vez;
las listas suman $m$ entradas y cada una cuesta $\Ord{1}$. Insertar al
frente de una lista ligada es $\Ord{1}$. Total: $\Tht{n + m}$ con
listas de adyacencia. Con matriz, encontrar los vecinos de cada vértice
cuesta $\Tht{n}$ y el total sube a $\Tht{n^2}$.

\begin{ejemplo}[La ropa del profesor, con DFS]
  \label{ej:ropa-dfs}
  Es la figura 22.7 de Cormen. Nueve prendas, donde la arista $(u, v)$
  significa que $u$ se pone antes que $v$: calzoncillos $\to$
  pantalón, calzoncillos $\to$ zapatos, pantalón $\to$ zapatos,
  pantalón $\to$ cinturón, camisa $\to$ cinturón, camisa $\to$ corbata,
  corbata $\to$ saco, cinturón $\to$ saco y calcetines $\to$ zapatos.
  El reloj está aislado. $n = 9$ y $m = 9$.

  Cormen no dice en qué orden recorre los vértices; este orden
  reproduce exactamente sus tiempos. Ciclo exterior: camisa, corbata,
  saco, cinturón, reloj, calzoncillos, pantalón, zapatos, calcetines.
  En las listas, camisa tiene a corbata antes que a cinturón, y
  pantalón tiene a zapatos antes que a cinturón.
  \begin{itemize}[itemsep=0.1em]
    \item camisa (1) $\to$ corbata (2) $\to$ saco (3); saco termina (4)
          y corbata (5). De regreso en camisa se descubre cinturón (6),
          cuyo único vecino, saco, ya es negro: cinturón termina (7) y
          camisa (8).
    \item reloj se descubre y termina sin vecinos: 9 y 10.
    \item calzoncillos (11) $\to$ pantalón (12) $\to$ zapatos (13), que
          termina (14). Cinturón ya es negro, así que pantalón termina
          (15). Zapatos ya es negro; calzoncillos termina (16).
    \item calcetines: 17 y 18; su vecino zapatos ya es negro.
  \end{itemize}
  Ninguna arista llega a un vértice gris: es un DAG. La lista $L$
  crece por el frente:
  \begin{center}
  \small
  \begin{tabular}{cll}
    \toprule
    $\text{sale}$ & \textbf{Termina} & \textbf{$L$ después (el frente
      a la izquierda)} \\
    \midrule
    $4$ & saco & saco \\
    $5$ & corbata & corbata, saco \\
    $7$ & cinturón & cinturón, corbata, saco \\
    $8$ & camisa & camisa, cinturón, corbata, saco \\
    $10$ & reloj & reloj, camisa, \ldots \\
    $14$ & zapatos & zapatos, reloj, camisa, \ldots \\
    $15$ & pantalón & pantalón, zapatos, reloj, \ldots \\
    $16$ & calzoncillos & calzoncillos, pantalón, zapatos, \ldots \\
    $18$ & calcetines & calcetines, calzoncillos, pantalón, zapatos,
      reloj, camisa, cinturón, corbata, saco \\
    \bottomrule
  \end{tabular}
  \end{center}
  % La ropa: arriba el DAG con entra/sale, abajo el orden topologico
  \begin{center}
  \begin{tikzpicture}[
    prenda/.style={draw=acento, rounded corners=3pt, fill=fondoIdea,
                   font=\scriptsize, inner sep=2.5pt,
                   minimum height=0.45cm}]
    \node[prenda] (cal) at (0, 0) {calzoncillos};
    \node[prenda] (pan) at (0, -1.1) {pantalón};
    \node[prenda] (cin) at (0, -2.3) {cinturón};
    \node[prenda] (cam) at (2.6, -1.8) {camisa};
    \node[prenda] (cor) at (2.6, -2.8) {corbata};
    \node[prenda] (sac) at (2.6, -3.8) {saco};
    \node[prenda] (cts) at (4.8, 0) {calcetines};
    \node[prenda] (zap) at (4.8, -1.1) {zapatos};
    \node[prenda] (rel) at (6.7, -0.55) {reloj};
    \foreach \a/\b in {cal/pan, cal/zap, pan/zap, pan/cin, cam/cin,
                       cam/cor, cor/sac, cin/sac, cts/zap}
      \draw[dirigida] (\a) -- (\b);
    \foreach \v/\d/\f in {cal/11/16, pan/12/15, cin/6/7}
      \node[nota, left=2pt] at (\v.west) {\d/\f};
    \foreach \v/\d/\f in {cam/1/8, cor/2/5, sac/3/4, cts/17/18,
                          zap/13/14, rel/9/10}
      \node[nota, right=2pt] at (\v.east) {\d/\f};
  \end{tikzpicture}
  \end{center}
  \begin{center}
  \begin{tikzpicture}[
    prenda/.style={draw=acento, rounded corners=3pt, fill=fondoIdea,
                   font=\scriptsize, inner sep=2.5pt,
                   minimum height=0.45cm}]
    \foreach \v/\nombre/\d/\f [count=\i from 0] in
        {cts/calcetines/17/18, cal/calzoncillos/11/16,
         pan/pantalón/12/15, zap/zapatos/13/14, rel/reloj/9/10,
         cam/camisa/1/8, cin/cinturón/6/7, cor/corbata/2/5,
         sac/saco/3/4} {
      \node[prenda] (\v) at (1.8*\i, 0) {\nombre};
      \node[nota, below=1pt] at (\v.south) {\d/\f};
    }
    \foreach \a/\b in {cal/pan, pan/zap, cam/cin, cor/sac}
      \draw[dirigida] (\a) -- (\b);
    \draw[dirigida] (cts.north) to[out=30, in=150]
      ([xshift=-5pt]zap.north);
    \draw[dirigida] (cal.north) to[out=40, in=140]
      ([xshift=5pt]zap.north);
    \draw[dirigida] (pan.north) to[out=25, in=155]
      ([xshift=-5pt]cin.north);
    \draw[dirigida] (cam.north) to[out=40, in=140] (cor.north);
    \draw[dirigida] ([xshift=5pt]cin.north) to[out=40, in=140]
      (sac.north);
  \end{tikzpicture}
  \end{center}
  Arriba, el DAG con $\text{entra}/\text{sale}$ junto a cada prenda.
  Abajo, las mismas prendas en orden decreciente de $\text{sale}$:
  calcetines, calzoncillos, pantalón, zapatos, reloj, camisa, cinturón,
  corbata, saco. Todas las flechas van a la derecha. El
  teorema~\ref{teo:topologico-salida} se comprueba arista por arista;
  por ejemplo, en cinturón $\to$ saco se tiene $7 > 4$, aunque cuando
  se exploró esa arista saco ya era negro.
\end{ejemplo}

\begin{ojo}
  Sin la prueba de gris (líneas 18 y 19) el algoritmo no avisa que hay
  ciclo: devuelve un orden inválido sin quejarse. Si a la ropa se le
  agrega la arista saco $\to$ camisa, absurda pero ilustrativa, el DFS
  con el mismo orden la explora cuando camisa sigue gris: es de
  retroceso, y cierra los ciclos camisa, corbata, saco, camisa y
  camisa, cinturón, saco, camisa. La versión de Cormen, que no tiene
  esa prueba, devuelve exactamente el mismo orden de antes, y ese orden
  pone a camisa antes que saco, en contra de la arista nueva.
\end{ojo}

\subsubsection*{El algoritmo de Kahn}

La otra forma de construir el orden va de adelante hacia atrás. Un
vértice sin aristas de entrada no depende de nada y puede ir primero.
Se saca, se borran sus aristas y se repite con lo que queda. La clase
del 22 de octubre lo menciona en el recuadro de ciclos con BFS, y
Cormen lo deja como ejercicio (22.4-5). En lugar de borrar aristas se
lleva un contador por vértice, $\text{faltan}[v]$: cuántos predecesores
de $v$ faltan por salir. Empieza en el grado de entrada
$\text{gr}^{-}(v)$ y solo baja.

\begin{algorithm}[H]
\caption{\textsc{Kahn}$(G)$, orden topológico por grados de entrada}
\label{alg:kahn}
\begin{algorithmic}[1]
\Procedure{Kahn}{$G$}
  \State $\text{faltan}[v] \gets 0$ para todo $v \in V$
  \ForAll{$u \in V$}
    \ForAll{$v$ en la lista de adyacencia de $u$}
      \State $\text{faltan}[v] \gets \text{faltan}[v] + 1$
    \EndFor
  \EndFor
  \State $Q \gets$ cola con todos los $v$ tales que
         $\text{faltan}[v] = 0$; \ $L \gets$ lista vacía
  \While{$Q$ no está vacía}
    \State $u \gets \textsc{Desencolar}(Q)$; \ agregar $u$ al final
           de $L$
    \ForAll{$v$ en la lista de adyacencia de $u$}
      \State $\text{faltan}[v] \gets \text{faltan}[v] - 1$
             \Comment{se borra la arista $(u, v)$}
      \If{$\text{faltan}[v] = 0$}
        \State \textsc{Encolar}$(Q, v)$
      \EndIf
    \EndFor
  \EndWhile
  \If{$\abs{L} < n$}
    \State \textbf{error}: los vértices que no están en $L$ están en un
           ciclo o debajo de uno
  \EndIf
  \State \Return $L$
\EndProcedure
\end{algorithmic}
\end{algorithm}

\begin{itemize}[itemsep=0.1em]
  \item \textbf{Líneas 2 a 7.} Calculan los grados de entrada
        recorriendo todas las listas.
  \item \textbf{Línea 12.} No se borra nada de la estructura;
        decrementar el contador simula borrar la arista.
  \item \textbf{Líneas 18 y 19.} Si salen menos de $n$ vértices, hay
        ciclo.
\end{itemize}

\begin{proposicion}[Correctitud de Kahn]
  \label{prop:kahn-correcto}
  Si $G$ es un DAG, {\upshape\textsc{Kahn}} devuelve un orden
  topológico de sus $n$ vértices. Si $G$ tiene un ciclo, devuelve
  menos de $n$ vértices, y ninguno de los vértices del ciclo está entre
  ellos.
\end{proposicion}

\begin{proof}
  En todo momento, $\text{faltan}[v]$ es el número de predecesores de
  $v$ que no están en $L$, porque cada vértice que se agrega a $L$
  decrementa una vez el contador de cada uno de sus sucesores. Un
  vértice entra a la cola justo cuando su contador llega a cero, es
  decir, cuando todos sus predecesores ya están en $L$. Por lo tanto,
  si $(u, v) \in E$ y $v$ está en $L$, $u$ está antes que $v$.

  Falta ver cuántos salen. Supóngase que la cola se vacía y queda un
  conjunto no vacío $R$ de vértices fuera de $L$. Ningún $v \in R$
  tiene $\text{faltan}[v] = 0$, porque habría entrado a la cola y
  estaría en $L$; así que todo vértice de $R$ tiene un predecesor en
  $R$. Partiendo de cualquier vértice de $R$ y retrocediendo de
  predecesor en predecesor sin salir de $R$, en a lo más $\abs{R}$
  pasos algún vértice se repite, y eso es un ciclo. En un DAG es
  imposible, así que $R = \emptyset$ y salen los $n$ vértices.

  Si $G$ tiene un ciclo $c$, ningún vértice de $c$ llega a $L$: el
  primero que llegara necesitaría a su predecesor en $c$ ya en $L$.
\end{proof}

La prueba da una segunda demostración de que todo DAG tiene orden
topológico, y de paso un hecho que se usa seguido: todo DAG no vacío
tiene al menos un vértice con grado de entrada cero, una
\textbf{fuente}, y por simetría uno con grado de salida cero, un
\textbf{sumidero}.

\textbf{Costo.} Calcular los grados recorre todas las listas:
$\Tht{n + m}$. Cada vértice entra a la cola a lo más una vez, porque su
contador solo baja y pasa por cero una sola vez; son $\Ord{n}$
operaciones de cola de $\Ord{1}$. Cada vértice que sale recorre su
lista una vez, $\Ord{m}$ en total. Suma: $\Tht{n + m}$. Si se quiere el
orden lexicográficamente menor, la cola se cambia por un montículo de
mínimos y el costo sube a $\Ord{n \log n + m}$.

\begin{ejemplo}[La ropa del profesor, con Kahn]
  \label{ej:ropa-kahn}
  Con cola FIFO y los empates en orden alfabético. Los grados de
  entrada iniciales son: zapatos $3$; cinturón y saco $2$; pantalón y
  corbata $1$; calcetines, calzoncillos, camisa y reloj $0$.
  \begin{center}
  \small
  \begin{tabular}{clll}
    \toprule
    \textbf{Paso} & \textbf{Sale} & \textbf{Contadores que bajan}
      & \textbf{Cola después} \\
    \midrule
    0 & & & calcetines, calzoncillos, camisa, reloj \\
    1 & calcetines & zapatos $3 \to 2$ & calzoncillos, camisa,
      reloj \\
    2 & calzoncillos & pantalón $1 \to \mathbf{0}$, zapatos
      $2 \to 1$ & camisa, reloj, pantalón \\
    3 & camisa & cinturón $2 \to 1$, corbata $1 \to \mathbf{0}$
      & reloj, pantalón, corbata \\
    4 & reloj & & pantalón, corbata \\
    5 & pantalón & cinturón $1 \to \mathbf{0}$, zapatos
      $1 \to \mathbf{0}$ & corbata, cinturón, zapatos \\
    6 & corbata & saco $2 \to 1$ & cinturón, zapatos \\
    7 & cinturón & saco $1 \to \mathbf{0}$ & zapatos, saco \\
    8 & zapatos & & saco \\
    9 & saco & & \\
    \bottomrule
  \end{tabular}
  \end{center}
  El orden de Kahn es calcetines, calzoncillos, camisa, reloj,
  pantalón, corbata, cinturón, zapatos, saco. Es distinto del de DFS y
  también es válido; el grafo tiene $1728$ órdenes topológicos,
  contados enumerando las $9!$ permutaciones. La diferencia de estilo
  se nota: Kahn saca primero todo lo que no depende de nada, las cuatro
  fuentes, mientras que DFS se mete por una cadena completa antes de
  pasar a otra.

  Con la arista saco $\to$ camisa agregada, Kahn saca calcetines,
  calzoncillos, reloj, pantalón y zapatos, y se detiene. Camisa,
  corbata, cinturón y saco nunca llegan a cero, y son justo los
  vértices de los dos ciclos.
\end{ejemplo}

Los dos algoritmos cuestan lo mismo y resuelven lo mismo, pero no son
intercambiables en todo:

\begin{center}
\small
\begin{tabular}{>{\raggedright\arraybackslash}p{3.1cm}
                >{\raggedright\arraybackslash}p{5.6cm}
                >{\raggedright\arraybackslash}p{5.6cm}}
  \toprule
  & \textbf{DFS por $\text{sale}$} & \textbf{Kahn} \\
  \midrule
  Costo & $\Tht{n + m}$ & $\Tht{n + m}$ \\
  Estructura & recursión o pila & cola y arreglo de contadores \\
  Detecta ciclos & solo con la prueba de gris & siempre: salen menos
    de $n$ \\
  Si hay ciclo, reporta & una arista de retroceso; con $\pi$ se
    reconstruye el ciclo & el conjunto de vértices bloqueados \\
  Construye el orden & de atrás hacia adelante & de adelante hacia
    atrás \\
  Recursión & hasta $n$ niveles & ninguna \\
  Lo que da de más & los tiempos que usan las componentes fuertes &
    los niveles: cada ronda es un lote de tareas paralelas \\
  \bottomrule
\end{tabular}
\end{center}

\textbf{Kahn por niveles.} Si en lugar de sacar un vértice a la vez se
saca la ronda completa de vértices con contador cero, y solo después
se decrementan los contadores, cada ronda es un conjunto de tareas que
se pueden hacer en paralelo. El número de rondas es el mínimo posible
para terminarlo todo, y es igual al número de vértices del camino más
largo del DAG. Con el DAG de materias de la clase del 22 de octubre
salen tres semestres:

\begin{center}
\begin{tikzpicture}[
  materia/.style={draw=acento, rounded corners=2pt, fill=fondoIdea,
                  font=\scriptsize, inner sep=3pt}]
  \node[materia] (cal) at (0, 0) {Cálculo};
  \node[materia] (alg) at (0, -1) {Álgebra lineal};
  \node[materia] (prg) at (0, -2) {Programación};
  \node[materia] (pro) at (3.6, 0) {Probabilidad};
  \node[materia] (ada) at (3.6, -2) {ADA};
  \node[materia] (ml) at (7.2, -1) {Aprendizaje};
  \foreach \a/\b in {cal/pro, pro/ml, alg/ml, prg/ada, ada/ml}
    \draw[dirigida] (\a) -- (\b);
  \foreach \x/\s in {0/1, 3.6/2, 7.2/3}
    \node[nota] at (\x, 0.65) {semestre \s};
\end{tikzpicture}
\end{center}

En la primera ronda salen las tres materias sin requisitos. Al
sacarlas, Probabilidad y ADA quedan libres y el contador de
Aprendizaje baja de $3$ a $2$. En la segunda ronda salen esas dos y
Aprendizaje llega a cero. El camino más largo, Cálculo, Probabilidad,
Aprendizaje, tiene tres vértices: no hay manera de terminar en dos
semestres.

\begin{ojo}
  Lo que Kahn deja fuera no es solo el ciclo: son los vértices que
  están en un ciclo \emph{o debajo de uno}, alcanzables desde él. Con
  $a \to b$, $b \to a$ y $b \to c$, ningún vértice tiene grado de
  entrada cero y la cola nace vacía; $c$ se queda fuera sin estar en
  ningún ciclo. Si se quiere el ciclo mismo, Kahn no lo da directo; DFS
  sí, con la arista de retroceso y los padres. Del lado de DFS, el
  riesgo es la recursión: en una cadena de $10^6$ vértices el DFS
  recursivo anida $10^6$ llamadas y revienta la pila (Python se detiene
  en $1000$ por omisión). Kahn no tiene ese problema; un DFS para
  grafos grandes se escribe con pila explícita.
\end{ojo}

\begin{ejercicio}[Cormen, ejercicio 22.4-4]
  En un grafo con ciclos no existe orden topológico, pero el algoritmo
  sin la prueba de gris devuelve algún orden. ¿Minimiza ese orden el
  número de aristas que apuntan hacia atrás?

  \textit{Solución.} No. Con $a \to b$, $b \to a$, $b \to c$,
  $c \to a$ y el DFS desde $a$ en orden alfabético, los tiempos son
  $a$: 1/6, $b$: 2/5, $c$: 3/4, y el orden es $a, b, c$, con dos
  aristas hacia atrás, $b \to a$ y $c \to a$. El orden $b, c, a$ tiene
  solo una, $a \to b$. Minimizar las aristas hacia atrás es el problema
  del conjunto mínimo de arcos de retroalimentación (\textit{feedback
  arc set}), que es NP-difícil; no hay que esperar que un recorrido
  lineal lo resuelva.
\end{ejercicio}

\begin{observacion}
  El orden topológico reaparece en la sección~\ref{sec:caminos-fuente}:
  relajar las aristas de un DAG en orden topológico da las distancias
  desde una fuente en $\Tht{n + m}$, incluso con pesos negativos
  (algoritmo~\ref{alg:caminos-dag}), que es justo donde Dijkstra falla.
  Con los pesos negados da el camino más largo, la ruta crítica de un
  proyecto. Y toda recurrencia de programación dinámica de la libreta
  se evalúa en un orden topológico del DAG de sus subproblemas.
\end{observacion}

% =====================================================================

\subsection*{Componentes fuertemente conexas}

En un grafo dirigido, que de $u$ se llegue a $v$ no dice nada sobre si
de $v$ se llega a $u$. Las componentes fuertemente conexas agrupan los
vértices que sí se alcanzan en los dos sentidos. Al contraer cada una a
un punto queda un DAG, el esqueleto del grafo: muchos algoritmos sobre
grafos dirigidos empiezan por ahí, resuelven cada componente por
separado y luego combinan siguiendo ese DAG en orden topológico.

\begin{definicion}[Componente fuertemente conexa]
  \label{def:componente-fuerte}
  Una \textbf{componente fuertemente conexa} (\textit{strongly
  connected component}) de un grafo dirigido $G = (V, E)$ es un
  conjunto maximal $C \subseteq V$ tal que, para todo $u, v \in C$, se
  tiene $u \leadsto v$ y $v \leadsto u$.
\end{definicion}

``Se alcanzan mutuamente'' es una relación de equivalencia: reflexiva
por el camino de longitud cero, simétrica por definición y transitiva
pegando caminos. Las componentes son sus clases, así que parten a $V$ y
cada vértice está en exactamente una. Un vértice solo forma una
componente aunque no tenga lazo. El grafo es \textbf{fuertemente
conexo} si tiene una sola componente.

\begin{definicion}[Grafo transpuesto y grafo de componentes]
  \label{def:grafo-componentes}
  El \textbf{transpuesto} de $G = (V, E)$ es $G^T = (V, E^T)$, con
  $E^T = \set{(u, v) : (v, u) \in E}$: las mismas aristas, volteadas.
  Si $C_1, \dots, C_k$ son las componentes fuertemente conexas de $G$,
  el \textbf{grafo de componentes} $G^{\text{CFC}}$ tiene un vértice
  $c_i$ por cada $C_i$, y la arista $(c_i, c_j)$, con $i \neq j$, si
  hay alguna arista $(x, y) \in E$ con $x \in C_i$ y $y \in C_j$.
  Cormen lo escribe $G^{\text{SCC}}$.
\end{definicion}

\subsubsection*{El algoritmo de dos DFS}

El algoritmo, de Kosaraju y Sharir, corre dos DFS. El primero, sobre
$G$, solo sirve para calcular los tiempos de salida. El segundo corre
sobre $G^T$ y toma las raíces en orden decreciente de esos tiempos.
Cada árbol del segundo bosque es exactamente una componente.

\begin{algorithm}[H]
\caption{\textsc{Componentes-Fuertes}$(G)$, dos DFS (Kosaraju)}
\label{alg:componentes-fuertes}
\begin{algorithmic}[1]
\Procedure{Componentes-Fuertes}{$G$}
  \State DFS sobre $G$ (algoritmo~\ref{alg:dfs-tiempos}); \ $L \gets$
         vértices en $\text{sale}$ decreciente
  \State construir $G^T$; \ $\text{comp}[u] \gets 0$ para todo
         $u \in V$; \ $k \gets 0$ \Comment{$0$: sin componente}
  \ForAll{$u$ en el orden de $L$} \Comment{$\text{sale}$ decreciente}
    \If{$\text{comp}[u] = 0$}
      \State $k \gets k + 1$; \ \textsc{Marcar}$(u, k)$
    \EndIf
  \EndFor
  \State \Return $\text{comp}$
\EndProcedure
\Statex
\Procedure{Marcar}{$u, k$} \Comment{DFS sobre $G^T$}
  \State $\text{comp}[u] \gets k$
  \ForAll{$v \in \text{Adj}^T[u]$}
    \State si $\text{comp}[v] = 0$, \textsc{Marcar}$(v, k)$
  \EndFor
\EndProcedure
\end{algorithmic}
\end{algorithm}

\begin{itemize}[itemsep=0.1em]
  \item \textbf{Línea 2.} Es el orden por $\text{sale}$ del
        algoritmo~\ref{alg:orden-topologico}, sin la prueba de ciclo
        (aquí los ciclos son el punto): cada vértice que termina se
        inserta al frente de $L$.
  \item \textbf{Línea 3.} Para construir $G^T$ se crea una lista
        vacía $\text{Adj}^T[u]$ por vértice y, por cada arista $(u, v)$
        de $G$, se agrega $u$ a $\text{Adj}^T[v]$.
  \item \textbf{Líneas 4 a 8.} El ciclo exterior del segundo DFS
        recorre $L$, o sea los vértices en $\text{sale}$ decreciente.
        Cada llamada a \textsc{Marcar} desde aquí es la raíz de un
        árbol nuevo, una componente nueva, y $\text{comp}[u] = k$ dice
        a cuál pertenece $u$.
  \item \textsc{Marcar} es un DFS normal sobre $G^T$. No necesita gris,
        porque no importa la clasificación de aristas, solo quién es
        alcanzable; $\text{comp}[v] = 0$ hace de blanco.
\end{itemize}

\textbf{Costo.} El primer DFS cuesta $\Tht{n + m}$. Transponer cuesta
$\Tht{n}$ para crear las listas vacías y $\Tht{m}$ para llenarlas. En
el segundo DFS cada vértice se marca una vez y cada lista de $G^T$, que
suman $m$, se recorre una vez: $\Tht{n + m}$. El orden por
$\text{sale}$ no cuesta un ordenamiento, porque la lista $L$ se arma en
$\Ord{1}$ por vértice; y aunque se ordenara, como
$\text{sale} \in \set{1, \dots, 2n}$, bastaría un ordenamiento por
conteo en $\Tht{n}$. Total: $\Tht{n + m}$, con memoria extra
$\Tht{n + m}$ por las listas de $G^T$.

\subsubsection*{Por qué funciona}

Primero, dos hechos estructurales. $G$ y $G^T$ tienen exactamente las
mismas componentes: voltear todas las aristas voltea todos los caminos,
así que $u$ y $v$ se alcanzan mutuamente en $G$ si y solo si se
alcanzan mutuamente en $G^T$. Y entre dos componentes distintas los
caminos van en un solo sentido:

\begin{lema}[Cormen, lema 22.13]
  \label{lem:componentes-dag}
  Sean $C \neq C'$ componentes, $u, v \in C$ y $u', v' \in C'$. Si hay
  un camino $u \leadsto u'$, no hay camino $v' \leadsto v$. En
  consecuencia, $G^{\text{CFC}}$ es un DAG.
\end{lema}

\begin{proof}
  Si hubiera un camino $v' \leadsto v$, se tendría
  $u \leadsto u' \leadsto v'$ y $v' \leadsto v \leadsto u$, así que $u$
  y $v'$ se alcanzarían mutuamente y estarían en la misma componente.
  Un ciclo en $G^{\text{CFC}}$ daría caminos en los dos sentidos entre
  dos componentes distintas del ciclo, que es justo lo que se acaba de
  descartar.
\end{proof}

Todo lo demás descansa en un lema sobre los tiempos de salida de las
componentes.

\begin{lema}[Cormen, lema 22.14]
  \label{lem:sale-componentes}
  Sean $C \neq C'$ componentes y $(u, v) \in E$ con $u \in C$ y
  $v \in C'$. Entonces $\text{sale}(C) > \text{sale}(C')$.
\end{lema}

\begin{proof}
  Hay dos casos, según cuál de las dos componentes toca primero el
  DFS.

  Si $\text{entra}(C) < \text{entra}(C')$, sea $x$ el primer vértice
  de $C$ que se descubre. En el instante $\text{entra}[x]$ todos los
  vértices de $C \cup C'$ son blancos. Desde $x$ hay un camino blanco a
  cada vértice de $C$ y, pasando por la arista $(u, v)$, a cada vértice
  de $C'$. Por el teorema del camino blanco todos quedan como
  descendientes de $x$, así que $x$ es el último en terminar:
  $\text{sale}[x] = \text{sale}(C) > \text{sale}(C')$.

  Si $\text{entra}(C) > \text{entra}(C')$, sea $y$ el primer vértice
  de $C'$ que se descubre. Por el mismo argumento, todo $C'$ queda
  debajo de $y$ y termina a más tardar en $\text{sale}[y] =
  \text{sale}(C')$. Durante el intervalo de $y$ no se descubre ningún
  vértice de $C$: lo que se descubre ahí es alcanzable desde $y$, y por
  el lema~\ref{lem:componentes-dag} ningún vértice de $C$ es
  alcanzable desde $C'$, porque sí hay camino de $C$ a $C'$. Como
  tampoco se descubrió ninguno antes de $\text{entra}[y]$, todo $C$ se
  descubre después de $\text{sale}(C')$ y, por lo tanto, termina
  después.
\end{proof}

\begin{corolario}[Cormen, corolario 22.15]
  \label{cor:sale-transpuesto}
  Si $(u, v) \in E^T$ con $u \in C$, $v \in C'$ y $C \neq C'$,
  entonces $\text{sale}(C) < \text{sale}(C')$. En palabras: en $G^T$,
  toda arista entre componentes va de la que terminó antes a la que
  terminó después.
\end{corolario}

\begin{proof}
  $(u, v) \in E^T$ significa $(v, u) \in E$; se aplica el
  lema~\ref{lem:sale-componentes} con los papeles invertidos.
\end{proof}

\begin{idea}
  Sea $C$ la componente con $\text{sale}(C)$ máximo, la que contiene al
  vértice que terminó al último en el primer DFS. Por el
  lema~\ref{lem:sale-componentes}, ninguna arista de $G$ entra a $C$
  desde otra componente, porque esa otra tendría $\text{sale}$ mayor:
  $C$ es una \textbf{fuente} de $G^{\text{CFC}}$. En $G^T$ las flechas
  se voltean y $C$ se vuelve un \textbf{sumidero}: un DFS que arranca
  en $C$ recorre toda $C$ y no puede salir. Ese primer árbol es
  exactamente $C$. Después se toma el vértice con $\text{sale}$ máximo
  entre los que quedan; por el corolario~\ref{cor:sale-transpuesto},
  las aristas de $G^T$ que salen de su componente solo pueden ir a
  componentes con $\text{sale}$ mayor, que ya están marcadas. Cada
  árbol queda encerrado porque todas sus salidas en $G^T$ llevan a
  árboles ya construidos. Por eso el orden es \emph{decreciente}.
\end{idea}

\begin{teorema}[Cormen, teorema 22.16]
  \label{teo:componentes-fuertes}
  {\upshape\textsc{Componentes-Fuertes}} calcula correctamente las
  componentes fuertemente conexas de $G$.
\end{teorema}

\begin{proof}
  Por inducción sobre el número de árboles del segundo DFS. La
  hipótesis es que los primeros $k$ árboles son componentes completas;
  con $k = 0$ no hay nada que probar. Sea $u$ la raíz del árbol
  $k + 1$ y $C$ su componente. Ningún vértice de $C$ está marcado,
  porque los árboles anteriores son componentes completas distintas de
  $C$. Por la elección de las raíces, $u$ tiene el mayor $\text{sale}$
  entre los vértices sin marcar, así que
  $\text{sale}[u] = \text{sale}(C) > \text{sale}(C')$ para toda
  componente $C' \neq C$ todavía sin marcar.

  Todos los vértices de $C$ se alcanzan desde $u$ en $G^T$ por caminos
  dentro de $C$, sin marcar, así que quedan en el árbol de $u$. Por
  otro lado, por el corolario~\ref{cor:sale-transpuesto}, toda arista
  de $G^T$ que sale de $C$ llega a una componente con $\text{sale}$
  mayor que $\text{sale}(C)$, que entonces no puede estar sin marcar:
  ya está marcada, y \textsc{Marcar} no entra ahí. Ningún vértice fuera
  de $C$ entra al árbol, y el árbol es exactamente $C$.
\end{proof}

Hay otra forma de leer el resultado: el segundo DFS visita las
componentes en orden topológico de $G^{\text{CFC}}$, de fuente a
sumidero. Las componentes salen ya ordenadas.

\begin{ejemplo}[La figura 22.9 de Cormen]
  \label{ej:componentes-cormen}
  Se usan las letras minúsculas de Cormen para no confundir con el
  grafo $G$ de la libreta. $V = \set{a, b, c, d, e, f, g, h}$, $n = 8$
  y $m = 14$:
  \[
    \begin{aligned}
      E = \{&(a,b), (b,c), (b,e), (b,f), (c,d), (c,g), (d,c),
             (d,h), \\
            &(e,a), (e,f), (f,g), (g,f), (g,h), (h,h)\} .
    \end{aligned}
  \]
  \textbf{Primer DFS.} Para reproducir los tiempos de la figura, el
  ciclo exterior empieza en $c$ y el siguiente blanco es $b$; en la
  lista de $c$ va $g$ antes que $d$, y en la de $g$, $f$ antes que $h$.
  \begin{itemize}[itemsep=0.1em]
    \item $c$ (1) $\to$ $g$ (2) $\to$ $f$ (3); $f \to g$ es de
          retroceso y $f$ termina (4). $g \to h$ (5); el lazo
          $h \to h$ es de retroceso, $h$ termina (6) y $g$ termina (7).
    \item $c \to d$ (8); $d \to c$ es de retroceso y $d \to h$ de
          cruce; $d$ termina (9) y $c$ termina (10).
    \item $b$ (11): $b \to c$ es de cruce; $b \to e$ (12) $\to$ $a$
          (13); $a \to b$ es de retroceso y $a$ termina (14); $e \to f$
          es de cruce y $e$ termina (15); $b \to f$ es de cruce y $b$
          termina (16).
  \end{itemize}
  En $\text{sale}$ decreciente,
  $L = b\,(16), e\,(15), a\,(14), c\,(10), d\,(9), g\,(7), h\,(6),
  f\,(4)$.
  % Figura 22.9: (a) primer DFS sobre G, (b) segundo DFS sobre G^T
  \begin{center}
  \begin{tikzpicture}[scale=0.85, transform shape,
    comp/.style={line width=0.8cm, line cap=round, line join=round,
                 draw=gray!18, fill=gray!18}]
    \draw[comp] (0, 0) -- (2.2, 0) -- (0, -1.8) -- cycle;
    \draw[comp] (4.4, 0) -- (6.6, 0);
    \draw[comp] (2.2, -1.8) -- (4.4, -1.8);
    \fill[gray!18] (6.6, -1.8) circle (0.4);
    \foreach \v/\x/\y in {a/0/0, b/2.2/0, c/4.4/0, d/6.6/0,
                          e/0/-1.8, f/2.2/-1.8, g/4.4/-1.8,
                          h/6.6/-1.8}
      \node[vert] (\v) at (\x, \y) {$\v$};
    \foreach \a/\b in {a/b, b/c, b/f, e/f, d/h}
      \draw[dirigida] (\a) -- (\b);
    \draw[dirigida] (d) to[bend left=18] (c);
    \draw[dirigida] (f) to[bend left=18] (g);
    \draw[dirigida] (h) to[out=-30, in=30, looseness=6] (h);
    \foreach \a/\b in {b/e, e/a, c/g, g/h, c/d}
      \draw[dirigidaArbol] (\a) -- (\b);
    \draw[dirigidaArbol] (g) to[bend left=18] (f);
    \foreach \v/\d/\f in {a/13/14, b/11/16, c/1/10, d/8/9}
      \node[nota, above=3pt] at (\v.north) {\d/\f};
    \foreach \v/\d/\f in {e/12/15, f/3/4, g/2/7, h/5/6}
      \node[nota, below=3pt] at (\v.south) {\d/\f};
    \node[font=\small] at (3.3, -3.1) {(a) $G$ y el primer DFS};
    \begin{scope}[xshift=8.6cm]
      \draw[comp] (0, 0) -- (2.2, 0) -- (0, -1.8) -- cycle;
      \draw[comp] (4.4, 0) -- (6.6, 0);
      \draw[comp] (2.2, -1.8) -- (4.4, -1.8);
      \fill[gray!18] (6.6, -1.8) circle (0.4);
      \foreach \v/\x/\y in {a/0/0, e/0/-1.8, d/6.6/0, f/2.2/-1.8}
        \node[vert] (t\v) at (\x, \y) {$\v$};
      \foreach \v/\x/\y in {b/2.2/0, c/4.4/0, g/4.4/-1.8, h/6.6/-1.8}
        \node[vertMarca] (t\v) at (\x, \y) {$\v$};
      \foreach \a/\b in {c/b, e/b, f/b, f/e, h/d, g/c, h/g}
        \draw[dirigida] (t\a) -- (t\b);
      \draw[dirigida] (td) to[bend left=18] (tc);
      \draw[dirigida] (tf) to[bend left=18] (tg);
      \draw[dirigida] (th) to[out=-30, in=30, looseness=6] (th);
      \foreach \a/\b in {b/a, a/e}
        \draw[dirigidaArbol] (t\a) -- (t\b);
      \draw[dirigidaArbol] (tc) to[bend left=18] (td);
      \draw[dirigidaArbol] (tg) to[bend left=18] (tf);
      \foreach \v/\t in {a/14, b/16, c/10, d/9}
        \node[nota, above=3pt] at (t\v.north) {\t};
      \foreach \v/\t in {e/15, f/4, g/7, h/6}
        \node[nota, below=3pt] at (t\v.south) {\t};
      \node[font=\small] at (3.3, -3.1) {(b) $G^T$ y el segundo DFS};
    \end{scope}
  \end{tikzpicture}
  \end{center}
  En (a), las aristas de árbol del primer DFS van en azul grueso, las
  componentes sombreadas y $\text{entra}/\text{sale}$ junto a cada
  vértice. En (b), $G^T$ con el bosque del segundo DFS que se describe
  enseguida; las raíces $b$, $c$, $g$ y $h$ resaltadas, y junto a cada
  vértice su $\text{sale}$ del primer DFS.


  \textbf{Segundo DFS}, sobre $G^T$ y en el orden de $L$. Las listas
  de $G^T$ son $a{:}\ e$; $b{:}\ a$; $c{:}\ b, d$; $d{:}\ c$;
  $e{:}\ b$; $f{:}\ b, e, g$; $g{:}\ c, f$; $h{:}\ d, g, h$.
  \begin{center}
  \small
  \begin{tabular}{clll}
    \toprule
    \textbf{Árbol} & \textbf{Raíz ($\text{sale}$)} & \textbf{Vértices}
      & \textbf{Aristas de $G^T$ que topan con árboles anteriores} \\
    \midrule
    1 & $b$ ($16$) & $b, a, e$ & ninguna: en $G^T$ nada sale de
      $\set{a, b, e}$ \\
    2 & $c$ ($10$) & $c, d$ & $c \to b$ \\
    3 & $g$ ($7$) & $g, f$ & $g \to c$, \ $f \to b$, \ $f \to e$ \\
    4 & $h$ ($6$) & $h$ & $h \to d$, \ $h \to g$ \\
    \bottomrule
  \end{tabular}
  \end{center}
  $e$, $a$, $d$ y $f$ aparecen en $L$, pero cuando les toca ya están
  marcados. Las componentes son $\set{a, b, e}$, $\set{c, d}$,
  $\set{f, g}$ y $\set{h}$. La última columna es el
  corolario~\ref{cor:sale-transpuesto} en acción: cada arista de $G^T$
  que sale de un árbol choca con un árbol anterior, nunca con uno
  futuro. Al contraer cada componente queda el grafo de componentes:
  \begin{center}
  \begin{tikzpicture}
    \node[vert, minimum size=0.85cm] (abe) at (0, -0.3) {$abe$};
    \node[vert, minimum size=0.85cm] (cd) at (4.2, 0.4) {$cd$};
    \node[vert, minimum size=0.85cm] (fg) at (2.4, -1.4) {$fg$};
    \node[vert, minimum size=0.85cm] (h) at (6.2, -1.4) {$h$};
    \foreach \a/\b in {abe/cd, abe/fg, cd/fg, cd/h, fg/h}
      \draw[dirigida] (\a) -- (\b);
    \node[nota, above=2pt] at (abe.north) {sale $16$};
    \node[nota, above=2pt] at (cd.north) {sale $10$};
    \node[nota, below=2pt] at (fg.south) {sale $7$};
    \node[nota, below=2pt] at (h.south) {sale $6$};
    \node[font=\small] at (3.1, -2.85) {(c) grafo de componentes
      $G^{\text{CFC}}$};
  \end{tikzpicture}
  \end{center}
  En (c), el grafo de componentes con el $\text{sale}$ de cada una: en
  cada arista el $\text{sale}$ baja, que es el
  lema~\ref{lem:sale-componentes}. Las componentes salieron en el orden
  $abe$, $cd$, $fg$, $h$, que es un orden topológico de
  $G^{\text{CFC}}$.
\end{ejemplo}

\begin{ejercicio}[Cormen, ejercicio 22.5-3]
  El profesor Bacon propone simplificar: correr el segundo DFS sobre
  $G$ mismo, sin transponer, tomando las raíces en orden
  \emph{creciente} de $\text{sale}$. Su razonamiento es que el primer
  vértice en terminar debería estar en una componente sumidero.
  ¿Funciona?

  \textit{Solución.} No, y el ejemplo mínimo muestra por qué. Con
  $a \to b$, $b \to a$, $a \to c$ y el DFS desde $a$ visitando primero
  a $b$, los tiempos son $a$: 1/6, $b$: 2/3, $c$: 4/5. El primero en
  terminar es $b$, pero su componente $\set{a, b}$ tiene una arista
  hacia $\set{c}$: no es sumidero. Desde $b$ en $G$ se alcanza todo y
  sale un solo árbol $\set{a, b, c}$ en lugar de $\set{a, b}$ y
  $\set{c}$. La asimetría es real: el vértice que termina al último
  siempre está en una componente fuente
  (lema~\ref{lem:sale-componentes}), pero el que termina primero no
  tiene por qué estar en una componente sumidero.
\end{ejercicio}

\begin{ojo}
  El primer DFS \emph{no} respeta las componentes. En la figura (a), el
  árbol que empieza en $c$ se lleva $\set{c, d}$, $\set{f, g}$ y
  $\set{h}$ juntos. Del primer DFS solo sirven los tiempos de salida;
  sus árboles no significan nada.

  Las dos piezas del segundo DFS son necesarias, el transpuesto y el
  orden decreciente de $\text{sale}$, y los tiempos son siempre los del
  \emph{primer} DFS. En el mismo grafo, quitar cualquiera de las dos
  falla:
  \begin{itemize}[itemsep=0.1em]
    \item sobre $G$, sin transponer, en $\text{sale}$ decreciente: desde
          $b$ se alcanza todo y sale un solo árbol con los ocho
          vértices;
    \item sobre $G^T$ en $\text{sale}$ creciente: desde $f$ salen
          $\set{f, b, a, e, g, c, d}$ y luego $\set{h}$;
    \item sobre $G$ en $\text{sale}$ creciente, la propuesta de Bacon:
          $\set{f, g, h}$, $\set{c, d}$, $\set{a, b, e}$, que junta a
          $\set{f, g}$ con $\set{h}$.
  \end{itemize}
\end{ojo}

\begin{ejercicio}[Contar caminos en un DAG; Cormen, ejercicio 22.4-2]
  Dado un DAG y dos vértices $s$ y $t$, contar los caminos de $s$ a $t$
  en tiempo $\Tht{n + m}$. Aplicarlo al grafo de componentes de la
  figura (c), de $abe$ a $h$.

  \textit{Solución.} En un DAG todo camino dirigido es simple, así que
  basta una recurrencia. Sea $P(u)$ el número de caminos de $u$ a $t$:
  $P(t) = 1$ y $P(u) = \sum_{(u, w) \in E} P(w)$. Cada $P(u)$ depende
  de sus sucesores, así que se evalúa en orden topológico inverso, y
  cada arista se usa una vez. En el ejemplo, $P(h) = 1$, $P(fg) = 1$,
  $P(cd) = P(fg) + P(h) = 2$ y $P(abe) = P(cd) + P(fg) = 3$; los tres
  caminos son $abe, cd, h$; $abe, cd, fg, h$; y $abe, fg, h$. Es
  programación dinámica sobre un DAG, como las recurrencias de monedas
  y mochila de la libreta.
\end{ejercicio}

\begin{ejercicio}[Cormen, ejercicio 22.5-7]
  Un grafo dirigido es \textbf{semiconexo} si para todo par $u, v$ se
  tiene $u \leadsto v$ o $v \leadsto u$. Decidirlo en $\Tht{n + m}$.

  \textit{Solución.} Calcular las componentes, construir
  $G^{\text{CFC}}$ y ordenarlo topológicamente. El grafo es semiconexo
  si y solo si hay una arista entre cada par de componentes
  consecutivas en ese orden. Si falta la arista entre $c_i$ y
  $c_{i + 1}$, ninguna alcanza a la otra: $c_{i + 1}$ no alcanza a
  $c_i$ por el orden, y un camino de $c_i$ a $c_{i + 1}$ tendría que
  pasar por componentes entre las dos, que no existen. El grafo de la
  figura 22.9 es semiconexo: $abe \to cd \to fg \to h$.
\end{ejercicio}

\begin{observacion}
  La tabla BFS contra DFS de la clase del 22 de octubre dice que DFS
  sirve para ``componentes''. En un grafo no dirigido, las componentes
  conexas son simplemente los árboles del bosque DFS. En un dirigido
  no: los árboles de un DFS mezclan componentes, como se acaba de ver,
  y hace falta el algoritmo de dos DFS. El grafo dirigido que la clase
  usa en Floyd es una sola componente fuerte, porque
  $1 \to 2 \to 3 \to 4 \to 1$ es un circuito que pasa por los cuatro
  vértices; su grafo de componentes es un solo punto.
\end{observacion}

\begin{observacion}
  El algoritmo de Tarjan encuentra las mismas componentes con un solo
  DFS, sin construir $G^T$, usando una pila de vértices y un valor muy
  parecido al $\text{low}[v]$ que aparece más abajo. Kosaraju es más
  fácil de explicar y de probar; Tarjan hace un solo recorrido. Una
  aplicación clásica de las componentes fuertes es 2-SAT: una fórmula
  en forma normal conjuntiva con dos literales por cláusula es
  satisfacible si y solo si ninguna variable queda en la misma
  componente que su negación en el grafo de implicaciones.
\end{observacion}

% =====================================================================

\subsection*{Puntos de articulación, puentes y componentes biconexas}

Esta parte es sobre grafos \emph{no dirigidos}. La
sección~\ref{sec:teoria-grafos} define los dos conceptos de fragilidad
de un grafo conexo. Un \textbf{punto de articulación}
(\textit{articulation point}; allá, vértice de corte) es un vértice cuya
eliminación, junto con sus aristas, desconecta el grafo. Un
\textbf{puente} (\textit{bridge}) es una arista cuya eliminación lo
desconecta. Allá también se prueba que una arista es puente si y solo
si no está en ningún ciclo. En una red de computadoras, un punto de
articulación es un servidor cuya caída parte la red en dos. Decidirlo
con la definición cuesta un recorrido por cada vértice,
$\Ord{n(n + m)}$. Aquí se hace todo en un solo DFS (Cormen, problema
22-2).

\begin{definicion}[Componente biconexa]
  \label{def:biconexa}
  Una \textbf{componente biconexa} (\textit{biconnected component}) de
  un grafo no dirigido es un conjunto maximal de aristas tal que dos
  aristas cualesquiera del conjunto están en un ciclo común.
\end{definicion}

La relación ``$e = e'$, o $e$ y $e'$ están en un ciclo común'' es de
equivalencia sobre las aristas. Lo delicado es la transitividad, que se
prueba recombinando dos ciclos que comparten una arista (Cormen,
problema 22-2g). Sus clases son las componentes biconexas, así que
parten a $E$.

\begin{observacion}[Convención para los puentes]
  Un puente no está en ningún ciclo, así que su clase es él solo. En
  esta libreta \textbf{cada puente es una componente biconexa de una
  sola arista}. Es la convención de la teoría de grafos (los
  \textit{bloques}), cumple la definición de arriba por vacuidad y es
  lo que produce el algoritmo de la pila: con ella las componentes
  biconexas parten a todas las aristas, sin excepción, y los puntos de
  articulación son exactamente los vértices que tocan dos o más
  componentes. Cormen, en la figura 22.10 y en el inciso (g) del
  problema, deja los puentes fuera de toda componente biconexa. Para
  obtener su versión basta descartar las componentes de una sola
  arista, que son exactamente los puentes.
\end{observacion}

\subsubsection*{Criterios con el árbol DFS}

Sea $G$ no dirigido y conexo, y sea $T$ el árbol de un DFS desde una
raíz $r$. Todo lo que sigue depende de un hecho de la
sección~\ref{sec:busquedas}: en un no dirigido no hay aristas de avance
ni de cruce, así que toda arista que no es del árbol une a un vértice
con un antecesor suyo. El \textbf{subárbol} de un vértice $v$ son $v$ y
sus descendientes en $T$.

\begin{proposicion}[Regla de la raíz]
  \label{prop:raiz-articulacion}
  La raíz $r$ es punto de articulación si y solo si tiene al menos dos
  hijos en $T$.
\end{proposicion}

\begin{proof}
  Si $r$ tiene dos hijos $s_1$ y $s_2$, toda arista que sale del
  subárbol de $s_1$ llega a un antecesor de $s_1$ fuera del subárbol,
  y el único es $r$. Sin $r$, ninguna arista une el subárbol de $s_1$
  con el de $s_2$, y el grafo se desconecta. Si $r$ tiene un solo hijo
  $s$, el subárbol de $s$ contiene a todos los demás vértices, porque
  el DFS desde $r$ alcanza todo el grafo conexo, y sigue conexo por sus
  aristas de árbol al quitar $r$.
\end{proof}

\begin{proposicion}[Vértices que no son la raíz]
  \label{prop:hijo-articulacion}
  Un vértice $v \neq r$ es punto de articulación si y solo si tiene un
  hijo $s$ tal que ninguna arista de retroceso sale del subárbol de $s$
  hacia un antecesor propio de $v$.
\end{proposicion}

\begin{proof}
  Si existe ese hijo $s$, las aristas que salen del subárbol de $s$
  llegan a antecesores de $s$ fuera del subárbol, o sea a $v$ o a
  antecesores propios de $v$. Por hipótesis ninguna llega a estos
  últimos, así que todas llegan a $v$. Al quitar $v$, el subárbol de
  $s$ queda separado de $r$, que existe porque $v \neq r$.

  Si cada hijo $s$ de $v$ tiene una arista de retroceso desde su
  subárbol hasta un antecesor propio de $v$, al quitar $v$ queda la
  parte de $T$ fuera del subárbol de $v$, que sigue conexa por aristas
  de árbol y contiene a $r$. Cada subárbol de un hijo sigue conexo por
  dentro y queda colgado de esa parte por su arista de retroceso. El
  grafo sigue conexo.
\end{proof}

Revisar si desde cada subárbol hay una arista que escapa por encima de
$v$ se vuelve una comparación de números con el valor siguiente.

\begin{definicion}[El valor low]
  \label{def:low}
  Para cada vértice $v$, $\text{low}[v]$ es el mínimo entre
  $\text{entra}[v]$ y los $\text{entra}[w]$ de todas las aristas de
  retroceso $\set{x, w}$ con $x$ en el subárbol de $v$ y $w$ antecesor
  de $x$. Es el descubrimiento más antiguo al que se puede saltar desde
  el subárbol de $v$ con una sola arista de retroceso. Cormen lo
  escribe $v.\mathit{low}$.
\end{definicion}

Como el subárbol de $v$ es $v$ junto con los subárboles de sus hijos,
\[
  \text{low}[v] = \min\Bigl(\text{entra}[v],\
    \min_{\set{v, w} \text{ de retroceso}} \text{entra}[w],\
    \min_{s \text{ hijo de } v} \text{low}[s]\Bigr) ,
\]
y se calcula en postorden dentro del mismo DFS: al regresar de cada
hijo, su $\text{low}$ ya es definitivo.

\begin{corolario}[Criterio de articulación con low]
  \label{cor:articulacion-low}
  Un vértice $v \neq r$ es punto de articulación si y solo si tiene un
  hijo $s$ con $\text{low}[s] \ge \text{entra}[v]$.
\end{corolario}

\begin{proof}
  Los antecesores propios de $v$ tienen $\text{entra}$ menor que
  $\text{entra}[v]$. Una arista de retroceso desde el subárbol de $s$
  llega a un antecesor de su extremo de abajo, y ese antecesor está
  dentro del subárbol de $s$ (con $\text{entra}$ mayor que
  $\text{entra}[v]$), es $v$, o es un antecesor propio de $v$. Por lo
  tanto, alguna llega a un antecesor propio de $v$ si y solo si
  $\text{low}[s] < \text{entra}[v]$. Se aplica la
  proposición~\ref{prop:hijo-articulacion}.
\end{proof}

\begin{proposicion}[Criterio de puente]
  \label{prop:puente-low}
  Una arista de retroceso nunca es puente. Una arista de árbol
  $\set{u, v}$, con $u = \pi[v]$, es puente si y solo si
  $\text{low}[v] > \text{entra}[u]$, lo que equivale a
  $\text{low}[v] = \text{entra}[v]$.
\end{proposicion}

\begin{proof}
  Una arista de retroceso cierra un ciclo con el camino del árbol entre
  sus extremos, y un puente no está en ningún ciclo. Para la arista de
  árbol $\set{u, v}$: si una arista de retroceso $\set{x, w}$ sale del
  subárbol de $v$ hacia $u$ o más arriba, el camino de árbol de $w$ a
  $u$, la arista $\set{u, v}$, el camino de árbol de $v$ a $x$ y la
  arista $\set{x, w}$ forman un ciclo, y $\set{u, v}$ no es puente. Si
  ninguna sale, la única arista que une el subárbol de $v$ con el resto
  es $\set{u, v}$, y quitarla lo separa. ``Alguna sale hacia $u$ o más
  arriba'' es exactamente $\text{low}[v] \le \text{entra}[u]$.

  La equivalencia final se debe a que $\text{low}[v]$ es el
  $\text{entra}$ de un antecesor de algún vértice del subárbol: o está
  en el subárbol, con $\text{entra} \ge \text{entra}[v]$, o es $u$ o
  algo más arriba, con $\text{entra} \le \text{entra}[u]$. No hay
  valores intermedios.
\end{proof}

Las componentes biconexas salen en el mismo recorrido con una pila de
aristas. Cada arista se apila la primera vez que se explora. Cuando, al
regresar del hijo $v$ a $u$, se cumple
$\text{low}[v] \ge \text{entra}[u]$, nada del subárbol de $v$ salta por
encima de $u$: las aristas apiladas desde $\set{u, v}$ hasta el tope
son las del subárbol de $v$ que todavía no tienen componente, junto con
las de retroceso que llegan a $u$, y forman una componente completa. Se
desapilan juntas.

\subsubsection*{El algoritmo}

El algoritmo~\ref{alg:articulaciones} junta los tres criterios en un
solo DFS. No necesita colores: $\text{entra}[v] = 0$ indica que $v$ no
se ha visto, y en un no dirigido no hace falta distinguir gris de
negro. El reloj avanza también al terminar, como en el
algoritmo~\ref{alg:dfs-tiempos}, para que los tiempos coincidan con los
de la libreta. Cormen solo pide los criterios y la idea en $\Ord{m}$;
el pseudocódigo es de esta libreta.

\begin{algorithm}[H]
\caption{\textsc{Articulaciones}$(G)$: puntos de articulación, puentes
         y componentes biconexas}
\label{alg:articulaciones}
\begin{algorithmic}[1]
\Procedure{Articulaciones}{$G$} \Comment{$G$ no dirigido}
  \State $\text{entra}[u] \gets 0$, \ $\text{hijos}[u] \gets 0$, \
         $\pi[u] \gets$ nulo, para todo $u \in V$
  \State $t \gets 0$; \ $P \gets$ pila de aristas vacía
  \ForAll{$r \in V$}
    \If{$\text{entra}[r] = 0$} \Comment{$r$ es raíz de un árbol nuevo}
      \State \textsc{Visitar-Low}$(r)$
      \State si $\text{hijos}[r] \ge 2$, marcar $r$ como punto de
             articulación \Comment{regla de la raíz}
    \EndIf
  \EndFor
\EndProcedure
\Statex
\Procedure{Visitar-Low}{$u$}
  \State $t \gets t + 1$, \quad $\text{entra}[u] \gets t$, \quad
         $\text{low}[u] \gets t$
  \ForAll{$v$ vecino de $u$}
    \If{$\text{entra}[v] = 0$} \Comment{arista de árbol}
      \State $\pi[v] \gets u$, \quad
             $\text{hijos}[u] \gets \text{hijos}[u] + 1$; \
             apilar $\set{u, v}$ en $P$
      \State \textsc{Visitar-Low}$(v)$
      \State $\text{low}[u] \gets \min(\text{low}[u], \text{low}[v])$
      \If{$\text{low}[v] \ge \text{entra}[u]$}
        \Comment{nada debajo de $v$ salta arriba de $u$}
        \State si $\pi[u] \neq$ nulo, marcar $u$ como punto de
               articulación
        \State desapilar de $P$ hasta $\set{u, v}$ inclusive: es una
               componente biconexa
      \EndIf
      \State si $\text{low}[v] > \text{entra}[u]$, $\set{u, v}$ es
             puente
    \ElsIf{$v \neq \pi[u]$ \textbf{y}
           $\text{entra}[v] < \text{entra}[u]$}
      \Comment{retroceso}
      \State apilar $\set{u, v}$ en $P$
      \State $\text{low}[u] \gets \min(\text{low}[u], \text{entra}[v])$
    \EndIf
  \EndFor
  \State $t \gets t + 1$, \quad $\text{sale}[u] \gets t$
\EndProcedure
\end{algorithmic}
\end{algorithm}

\begin{itemize}[itemsep=0.1em]
  \item \textbf{Líneas 14 a 22.} Arista de árbol. Al regresar de $v$
        ya se conoce $\text{low}[v]$, y la línea 17 lo propaga a $u$.
  \item \textbf{Línea 18.} $\text{low}[v] \ge \text{entra}[u]$ dice que
        nada del subárbol de $v$ salta por encima de $u$. Si $u$ no es
        raíz, eso lo hace articulación
        (corolario~\ref{cor:articulacion-low}). En la raíz la condición
        se cumple siempre, porque nada está por encima de ella; por eso
        la raíz usa la regla de los dos hijos (línea 7). Las
        componentes biconexas sí se desapilan también en la raíz.
  \item \textbf{Línea 22.} Desigualdad estricta. Con igualdad, el
        subárbol de $v$ regresa justo a $u$ por otra arista, así que
        $\set{u, v}$ está en un ciclo.
  \item \textbf{Línea 23.} $v \neq \pi[u]$ excluye la arista por la que
        se llegó, que vista desde $u$ no es de retroceso. La condición
        $\text{entra}[v] < \text{entra}[u]$ hace que cada arista de
        retroceso se apile una sola vez, desde el extremo de abajo:
        cuando el antecesor la ve desde arriba, $v$ ya es un
        descendiente terminado con $\text{entra}$ mayor.
  \item \textbf{Línea 25.} Se usa $\text{entra}[v]$, no
        $\text{low}[v]$. Con $\text{low}[v]$ el algoritmo pierde
        articulaciones, como se ve en el ejemplo siguiente.
\end{itemize}

\textbf{Costo.} Es un DFS sobre un grafo no dirigido: cada vértice se
visita una vez y cada lista se recorre una vez, y las listas suman
$2m$ entradas. Cada arista se apila exactamente una vez, las de árbol
al descubrirlas y las de retroceso solo desde abajo, y se desapila una
vez. Total: $\Tht{n + m}$, que es $\Ord{m}$ en un grafo conexo porque
ahí $n \le m + 1$.

\begin{ejemplo}[Un moño con cola]
  \label{ej:articulaciones-mono}
  El grafo $H$ está hecho para que aparezcan todos los casos.
  $V = \set{A, \dots, H}$, $n = 8$, $m = 10$:
  \[
    E = \set{\set{A,B}, \set{A,C}, \set{B,C}, \set{C,D}, \set{C,E},
             \set{D,E}, \set{E,F}, \set{F,G}, \set{F,H}, \set{G,H}} .
  \]
  Son dos triángulos que comparten el vértice $C$, un moño, más un
  puente $\set{E, F}$ y un tercer triángulo colgado de $F$.

  El DFS empieza en $A$, con los vecinos en orden alfabético. El árbol
  resulta ser un camino, $A, B, C, \dots, H$, así que $\text{entra}$ va
  de $1$ a $8$. Las aristas de retroceso son $\set{C, A}$,
  $\set{E, C}$ y $\set{H, F}$. La tabla está ordenada por
  $\text{entra}$ para seguir el árbol, pero se llena en postorden, de
  $H$ hacia arriba:
  \begin{center}
  \small
  \begin{tabular}{cccclc}
    \toprule
    $u$ & $\text{entra}$ & $\pi$ & \textbf{Retroceso}
      & \textbf{Prueba con el hijo $s$} & $\text{low}$ \\
    \midrule
    $A$ & $1$ & & & raíz con un solo hijo: no & $1$ \\
    $B$ & $2$ & $A$ & & $\text{low}[C] = 1 < 2$: no & $1$ \\
    $C$ & $3$ & $B$ & a $A$, entra $1$ & $\text{low}[D] = 3 \ge 3$:
      \textbf{sí} & $1$ \\
    $D$ & $4$ & $C$ & & $\text{low}[E] = 3 < 4$: no & $3$ \\
    $E$ & $5$ & $D$ & a $C$, entra $3$ & $\text{low}[F] = 6 \ge 5$:
      \textbf{sí}, y $6 > 5$: puente & $3$ \\
    $F$ & $6$ & $E$ & & $\text{low}[G] = 6 \ge 6$: \textbf{sí} & $6$ \\
    $G$ & $7$ & $F$ & & $\text{low}[H] = 6 < 7$: no & $6$ \\
    $H$ & $8$ & $G$ & a $F$, entra $6$ & sin hijos: no & $6$ \\
    \bottomrule
  \end{tabular}
  \end{center}
  % H: el arbol DFS tendido en linea, con entra y low
  \begin{center}
  \begin{tikzpicture}
    \foreach \v [count=\i from 0] in {A, B, C, D, E, F, G, H} {
      \coordinate (p\v) at (1.35*\i, 0);
    }
    \foreach \v in {A, B, D, G, H}
      \node[vert] (\v) at (p\v) {$\v$};
    \foreach \v in {C, E, F}
      \node[vertMarca] (\v) at (p\v) {$\v$};
    \foreach \a/\b in {A/B, B/C, C/D, D/E, F/G, G/H}
      \draw[aristaArbol] (\a) -- (\b);
    \draw[aristaResalte] (E) -- (F);
    \draw[aristaFuera] (C) to[out=125, in=55] (A);
    \draw[aristaFuera] (E) to[out=125, in=55] (C);
    \draw[aristaFuera] (H) to[out=125, in=55] (F);
    \node[nota, anchor=east] at (-0.5, -0.65) {entra};
    \node[nota, anchor=east] at (-0.5, -1.05) {low};
    \foreach \v/\d/\l in {A/1/1, B/2/1, C/3/1, D/4/3, E/5/3, F/6/6,
                          G/7/6, H/8/6} {
      \node[nota] at ([yshift=-0.65cm]p\v) {$\d$};
      \node[font=\scriptsize] at ([yshift=-1.05cm]p\v) {$\l$};
    }
  \end{tikzpicture}
  \end{center}
  El árbol DFS tendido en línea, con las aristas de retroceso punteadas
  y las filas de $\text{entra}$ y $\text{low}$. Resultado: puntos de
  articulación $C$, $E$ y $F$, resaltados, y un solo puente,
  $\set{E, F}$, en naranja. $C$ es el caso fino:
  $\text{low}[D] = \text{entra}[C] = 3$, así que $C$ es articulación
  pero $\set{C, D}$ no es puente, porque la arista $\set{E, C}$ regresa
  justo a $C$.

  Las componentes biconexas salen de la pila en este orden:
  $\set{FG, GH, HF}$ al regresar de $G$ a $F$; $\set{EF}$, el puente,
  al regresar de $F$ a $E$; $\set{CD, DE, EC}$ al regresar de $D$ a
  $C$; y $\set{AB, BC, CA}$ al regresar de $B$ a $A$.
  % H: las cuatro componentes biconexas
  \begin{center}
  \begin{tikzpicture}[
    bloque/.style={line width=0.8cm, line cap=round, line join=round,
                   draw=#1, fill=#1}]
    \draw[bloque=blue!10] (0, 0.8) -- (0, -0.8) -- (1.4, 0) -- cycle;
    \draw[bloque=green!12] (1.4, 0) -- (2.8, 0.8) -- (2.8, -0.8)
      -- cycle;
    \draw[bloque=violet!10] (4.4, -0.8) -- (5.8, 0) -- (5.8, -1.6)
      -- cycle;
    \foreach \v/\x/\y in {A/0/0.8, B/0/-0.8, D/2.8/0.8, G/5.8/0,
                          H/5.8/-1.6}
      \node[vert] (\v) at (\x, \y) {$\v$};
    \foreach \v/\x/\y in {C/1.4/0, E/2.8/-0.8, F/4.4/-0.8}
      \node[vertMarca] (\v) at (\x, \y) {$\v$};
    \foreach \a/\b in {A/B, B/C, C/A, C/D, D/E, E/C, F/G, G/H, H/F}
      \draw[arista] (\a) -- (\b);
    \draw[aristaResalte] (E) -- node[nota, below=3pt] {2} (F);
    \node[nota] at (-0.75, 0) {4};
    \node[nota] at (3.45, 0.35) {3};
    \node[nota] at (6.5, -0.8) {1};
  \end{tikzpicture}
  \end{center}
  Las cuatro componentes, numeradas en el orden en que salen de la
  pila; la 2 es el puente. Las articulaciones son justo los vértices
  que tocan dos componentes. Todo coincide con la fuerza bruta: quitar
  cada vértice y cada arista y revisar si lo que queda sigue conexo.

  Si el DFS arrancara en $C$, entraría en juego la regla de la raíz:
  $C \to A \to B$, y $\set{B, C}$ es de retroceso; de vuelta en $C$,
  $B$ ya está visto y el recorrido sigue
  $C \to D \to E \to F \to G \to H$. La raíz $C$ queda con dos hijos,
  $A$ y $D$, y es articulación por la
  proposición~\ref{prop:raiz-articulacion}. El resultado es el mismo:
  $C$, $E$ y $F$, con el puente $\set{E, F}$.
\end{ejemplo}

\begin{ojo}
  En la línea 25, la arista de retroceso aporta $\text{entra}[v]$, no
  $\text{low}[v]$. Es un error muy común, y en $H$ se nota: cuando $E$
  ve su arista de retroceso hacia $C$, $\text{low}[C]$ ya vale $1$ por
  la arista $\set{C, A}$, aunque $C$ todavía no termina. Con
  $\text{low}[C]$, los valores de $E$ y de $D$ bajan a $1$, la prueba
  en $C$ da $1 < 3$ y el algoritmo \emph{pierde la articulación $C$}.
  El $\text{low}$ de un antecesor gris todavía no es definitivo y,
  peor, describe saltos que no pasan por la arista en cuestión. Para
  puentes el error no se nota; para articulaciones sí: en $3000$ grafos
  conexos aleatorios de $3$ a $9$ vértices, la variante errónea cambió
  las articulaciones en $25$ y los puentes en ninguno.
\end{ojo}

\begin{ejemplo}[El grafo $G$ de la libreta]
  \label{ej:articulaciones-libreta}
  La clase del 22 de octubre ya hizo un DFS sobre $G$ desde $A$, con
  los vecinos en orden alfabético, y obtuvo $A$ 1/14, $B$ 2/13, $D$
  3/6, $G$ 4/5, $E$ 7/12, $F$ 8/11 y $C$ 9/10, con una sola arista de
  retroceso, $\set{C, A}$. Con esa misma traza:
  \begin{itemize}[itemsep=0.1em]
    \item $\text{low}[C] = \text{entra}[A] = 1$, y ese $1$ sube por
          $F$, $E$, $B$ y $A$: todos tienen $\text{low} = 1$.
    \item $G$ no tiene aristas de retroceso, así que
          $\text{low}[G] = \text{entra}[G] = 4$, y
          $\text{low}[D] = \min(3, 4) = 3$.
    \item $B$ es articulación por su hijo $D$:
          $\text{low}[D] = 3 \ge \text{entra}[B] = 2$. $D$ lo es por su
          hijo $G$: $4 \ge 3$. Los demás no: $E$ y $F$ tienen hijos con
          $\text{low} = 1$, y la raíz $A$ tiene un solo hijo.
    \item Puentes: $\set{B, D}$, porque $\text{low}[D] = 3 > 2$, y
          $\set{D, G}$, porque $4 > 3$.
  \end{itemize}
  Las componentes biconexas salen en el orden $\set{DG}$, $\set{BD}$ y
  $\set{AB, BE, EF, FC, CA}$: los dos puentes y el ciclo
  $A, B, E, F, C, A$ que la clase encontró con BFS.
  \begin{center}
  \begin{tikzpicture}[baseline=(current bounding box.north)]
    \node[vert] (A) at (0, 0) {A};
    \node[vertMarca] (B) at (0, -1.1) {B};
    \node[vertMarca] (D) at (-1.3, -2.2) {D};
    \node[vert] (G) at (-1.3, -3.3) {G};
    \node[vert] (E) at (1.3, -2.2) {E};
    \node[vert] (F) at (1.3, -3.3) {F};
    \node[vert] (C) at (1.3, -4.4) {C};
    \foreach \a/\b in {A/B, B/E, E/F, F/C}
      \draw[aristaArbol] (\a) -- (\b);
    \foreach \a/\b in {B/D, D/G}
      \draw[aristaResalte] (\a) -- (\b);
    \draw[aristaFuera] (C) to[out=0, in=0, looseness=1.4] (A);
    \foreach \v/\e/\l in {A/1/1, B/2/1, D/3/3, G/4/4, E/7/1, F/8/1,
                          C/9/1}
      \node[nota, left=2pt] at (\v.west) {$\e/\l$};
  \end{tikzpicture}
  \hspace{1.6cm}
  \begin{tikzpicture}[baseline=(current bounding box.north),
    bloque/.style={line width=0.8cm, line cap=round, line join=round,
                   draw=#1, fill=#1}]
    \draw[bloque=blue!10] (0, 0) -- (-1.6, -1.2) -- (-0.6, -2.4)
      -- (1.6, -2.4) -- (1.6, -1.2) -- cycle;
    \node[vert] (A) at (0, 0) {A};
    \node[vertMarca] (B) at (-1.6, -1.2) {B};
    \node[vert] (C) at (1.6, -1.2) {C};
    \node[vertMarca] (D) at (-2.6, -2.4) {D};
    \node[vert] (E) at (-0.6, -2.4) {E};
    \node[vert] (F) at (1.6, -2.4) {F};
    \node[vert] (G) at (-2.6, -3.6) {G};
    \foreach \a/\b in {A/B, A/C, B/E, C/F, E/F}
      \draw[arista] (\a) -- (\b);
    \foreach \a/\b in {B/D, D/G}
      \draw[aristaResalte] (\a) -- (\b);
  \end{tikzpicture}
  \end{center}
  A la izquierda, el árbol DFS de la clase con
  $\text{entra}/\text{low}$ junto a cada vértice y la arista de
  retroceso punteada. A la derecha, $G$ con sus tres componentes
  biconexas: el ciclo sombreado y los dos puentes en naranja. Los
  puntos de articulación, resaltados, son los vértices donde se tocan.
\end{ejemplo}

\begin{ejercicio}
  ¿Qué devuelve \textsc{Articulaciones} sobre un árbol con $n \ge 3$
  vértices? ¿Y sobre un ciclo?

  \textit{Solución.} En un árbol no hay aristas de retroceso, así que
  $\text{low}[v] = \text{entra}[v]$ para todo $v$: las $n - 1$ aristas
  son puentes, cada una es su propia componente biconexa, y los puntos
  de articulación son los vértices de grado al menos $2$. En un ciclo
  el árbol DFS es un camino y la única arista de retroceso va del
  último vértice a la raíz, así que $\text{low} = 1$ en todos los
  vértices: no hay puentes ni articulaciones, y todas las aristas
  forman una sola componente biconexa.
\end{ejercicio}

\begin{ojo}
  Tres detalles más que rompen el algoritmo si se olvidan.
  \begin{itemize}[itemsep=0.1em]
    \item \textbf{La arista al padre.} Si no se excluye la arista por
          la que se llegó a $u$, siempre
          $\text{low}[u] \le \text{entra}[\pi[u]]$ y no se encuentra
          ningún puente. En un multigrafo hay que excluir \emph{esa
          arista}, por identificador, no al vértice padre: dos aristas
          paralelas entre $u$ y $v$ forman un ciclo y ninguna es
          puente.
    \item \textbf{La raíz.} La prueba $\text{low}[s] \ge
          \text{entra}[r]$ se cumple siempre en la raíz. Aplicarla ahí
          declara articulación a toda raíz; por eso la raíz usa la
          regla de los dos hijos.
    \item \textbf{Grafos no conexos.} El ciclo exterior de
          \textsc{Articulaciones} arranca un DFS en cada componente
          conexa, y cada raíz usa su propia regla. En ese caso
          ``desconecta'' se lee como ``aumenta el número de componentes
          conexas''.
  \end{itemize}
\end{ojo}

% =====================================================================

\subsection*{Recorrido de Euler en grafos dirigidos}

La sección~\ref{sec:teoria-grafos} trata el caso no dirigido: un
multigrafo conexo tiene circuito euleriano si y solo si todos sus
grados son pares, y el algoritmo de Hierholzer lo construye en
$\Tht{n + m}$. El problema 22-3 de Cormen pide la versión dirigida. Un
\textbf{recorrido de Euler} (\textit{Euler tour}) de un grafo dirigido
es un circuito que usa cada arista exactamente una vez. Puede repetir
vértices, así que no es un ciclo en el sentido de esta libreta, aunque
Cormen lo llame así.

\begin{teorema}[Cormen, problema 22-3]
  \label{teo:euler-dirigido-cormen}
  Un grafo dirigido fuertemente conexo tiene recorrido de Euler si y
  solo si $\text{gr}^{-}(v) = \text{gr}^{+}(v)$ para todo vértice $v$.
\end{teorema}

Es el teorema~\ref{teo:euler-dirigido} de la
sección~\ref{sec:teoria-grafos} en la forma de Cormen. Allá basta la
conexidad débil, porque con los grados balanceados un grafo débilmente
conexo resulta fuertemente conexo.

\begin{proof}[Idea de la prueba]
  Cada paso del recorrido por $v$ entra por una arista y sale por otra,
  y cada arista se usa una vez: los grados quedan balanceados. Al
  revés, se camina desde $s$ por aristas no usadas mientras se pueda.
  Al llegar a un $v \neq s$ se ha entrado a $v$ una vez más de las que
  se ha salido, y como sus grados están balanceados le queda una salida
  libre: la caminata solo puede atorarse en $s$, y al hacerlo cierra un
  circuito. Si sobran aristas, por conexidad alguna sale de un vértice
  $x$ del circuito; en lo que sobra los grados siguen balanceados, así
  que desde $x$ sale otro circuito, que se empalma en $x$. Es el
  argumento del caso no dirigido con ``grado par'' cambiado por
  ``entra tanto como sale''.
\end{proof}

Hierholzer funciona igual siguiendo solo las aristas de salida: una
pila guarda el camino actual, se avanza mientras el tope tenga aristas
sin usar, y cuando se atora se desapila e inserta al frente del
recorrido. Cada arista se usa una vez: $\Tht{n + m}$.

\begin{ejemplo}[Un empalme]
  \label{ej:euler-dirigido}
  Con las aristas $a \to b$, $b \to c$, $c \to a$, $b \to d$,
  $d \to e$ y $e \to b$, el vértice $b$ tiene grados de entrada y
  salida $(2, 2)$ y los demás $(1, 1)$, y el grafo es fuertemente
  conexo. En la lista de $b$ va $c$ antes que $d$, para que el primer
  intento se atore antes de usarlo todo.
  \begin{center}
  \begin{tikzpicture}
    \node[vert] (a) at (0, 0) {$a$};
    \node[vert] (b) at (2, 0) {$b$};
    \node[vert] (c) at (1, -1.5) {$c$};
    \node[vert] (d) at (3.6, 0.9) {$d$};
    \node[vert] (e) at (3.6, -0.9) {$e$};
    \draw[dirigidaArbol] (a) -- node[peso] {1} (b);
    \draw[dirigidaArbol] (b) -- node[peso] {5} (c);
    \draw[dirigidaArbol] (c) -- node[peso] {6} (a);
    \draw[dirigidaFuera] (b) -- node[peso] {2} (d);
    \draw[dirigidaFuera] (d) -- node[peso] {3} (e);
    \draw[dirigidaFuera] (e) -- node[peso] {4} (b);
  \end{tikzpicture}
  \end{center}
  Desde $a$, la pila crece $a, b, c, a$ y se atora en $a$ sin haber
  usado $b \to d$. Al desapilar ($a$, luego $c$) se regresa a $b$, que
  todavía tiene aristas, y desde ahí se hace el desvío $b, d, e, b$.
  Luego todo se atora y se desapila. El recorrido final es
  $a, b, d, e, b, c, a$: el desvío quedó empalmado en $b$. En azul, el
  primer circuito que se encontró; en naranja punteado, el desvío; los
  números dan la posición de cada arista en el recorrido final.
\end{ejemplo}

\begin{ojo}
  Sin conexidad fuerte el balance no basta: dos ciclos ajenos están
  balanceados y no tienen recorrido de Euler. La condición general es
  el balance más que todas las aristas estén en una misma componente
  fuertemente conexa; los vértices aislados no importan. Y el grafo
  dirigido de Floyd de la clase del 22 de octubre, aunque es
  fuertemente conexo, no tiene recorrido de Euler: el vértice $1$
  tiene grado de entrada $3$ y de salida $2$.
\end{ojo}

% =====================================================================

\subsection*{Alcanzabilidad mínima}

El problema 22-4 de Cormen es un ejercicio sobre la dirección de las
aristas. Cada vértice $u$ de un grafo dirigido tiene una etiqueta
distinta $L(u) \in \set{1, \dots, n}$, y se quiere, para cada $u$, el
vértice $\text{mín}[u]$ de etiqueta mínima entre los alcanzables desde
$u$, incluido $u$. Correr un recorrido desde cada vértice costaría
$\Ord{n(n + m)}$.

La clave es que $u \leadsto v$ en $G$ si y solo si $v \leadsto u$ en
$G^T$. Se procesan los vértices en orden \emph{creciente} de etiqueta,
y desde cada $v$ todavía sin asignar se corre un DFS en $G^T$ que
asigna $\text{mín}[u] = v$ a todo vértice sin asignar que alcance.

\begin{algorithm}[H]
\caption{\textsc{Mínimos-Alcanzables}$(G, L)$, problema 22-4}
\label{alg:minimos-alcanzables}
\begin{algorithmic}[1]
\Procedure{Mínimos-Alcanzables}{$G, L$}
  \State construir $G^T$
  \ForAll{$u \in V$}
    \State $\text{vert}[L(u)] \gets u$, \quad
           $\text{mín}[u] \gets$ nulo
           \Comment{ordena por $L$ en $\Tht{n}$}
  \EndFor
  \For{$i \gets 1$ \Hasta $n$}
    \State $v \gets \text{vert}[i]$; \ si $\text{mín}[v] =$ nulo,
           \textsc{Asignar}$(v, v)$
  \EndFor
  \State \Return $\text{mín}$
\EndProcedure
\Statex
\Procedure{Asignar}{$u, v$} \Comment{DFS sobre $G^T$}
  \State $\text{mín}[u] \gets v$
  \ForAll{$w \in \text{Adj}^T[u]$}
    \State si $\text{mín}[w] =$ nulo, \textsc{Asignar}$(w, v)$
  \EndFor
\EndProcedure
\end{algorithmic}
\end{algorithm}

\textbf{Por qué funciona.} Si $u$ sigue sin asignar cuando le toca a
$v$, ningún vértice de etiqueta menor es alcanzable desde $u$; si
además $u \leadsto v$, entonces $v$ es su mínimo. Detener el DFS en los
vértices ya asignados no pierde nada: si $w$ fue asignado por $v'$, con
$L(v') < L(v)$, todo lo que alcanza a $w$ alcanza a $v'$ y ya tiene un
mínimo de etiqueta a lo más $L(v')$. Por la misma razón, el camino
$u \leadsto v$ de un $u$ sin asignar pasa solo por vértices sin
asignar, y el DFS lo encuentra.

\textbf{Costo.} $G^T$ cuesta $\Tht{n + m}$. Invertir $L$ cuesta
$\Tht{n}$, porque es una biyección con $\set{1, \dots, n}$ y no hace
falta comparar. \textsc{Asignar} se llama una vez por vértice y cada
lista de $G^T$ se recorre una vez. Total: $\Tht{n + m}$.

\begin{ejemplo}[Mínimos alcanzables]
  \label{ej:minimos-alcanzables}
  Con las aristas $p \to q$, $q \to p$, $q \to r$, $r \to s$,
  $s \to r$, $t \to s$, $t \to u$ y las etiquetas $L(p) = 1$,
  $L(s) = 2$, $L(u) = 3$, $L(r) = 4$, $L(q) = 5$, $L(t) = 6$:
  \begin{center}
  \small
  \begin{tabular}{cccl}
    \toprule
    $L$ & \textbf{Vértice} & \textbf{¿Asignado?}
      & \textbf{El DFS en $G^T$ asigna} \\
    \midrule
    $1$ & $p$ & no & $p, q$ \\
    $2$ & $s$ & no & $s, r, t$ (de $r$ se llegaría a $q$, ya
      asignado) \\
    $3$ & $u$ & no & $u$ (su vecino en $G^T$, $t$, ya está
      asignado) \\
    $4, 5, 6$ & $r, q, t$ & sí & nada \\
    \bottomrule
  \end{tabular}
  \end{center}
  Resultado: $\text{mín}[p] = \text{mín}[q] = p$;
  $\text{mín}[r] = \text{mín}[s] = \text{mín}[t] = s$;
  $\text{mín}[u] = u$. El caso interesante es $t$: alcanza a $s$, $r$
  y $u$, con etiquetas $2$, $4$ y $3$, y su mínimo es $s$ aunque $u$
  esté a un solo paso.
\end{ejemplo}

\begin{ojo}
  Aquí importa la dirección dos veces. El DFS va sobre $G^T$, porque se
  busca quién \emph{alcanza} a $v$, y en orden \emph{creciente} de
  etiqueta. Sobre $G$ se calcularía otra cosa: el mínimo de los
  vértices que alcanzan a cada uno.
\end{ojo}

% =====================================================================

\subsection*{Resumen}

\begin{center}
\small
\begin{tabular}{>{\raggedright\arraybackslash}p{3.4cm}
                >{\raggedright\arraybackslash}p{8.2cm}l}
  \toprule
  \textbf{Problema} & \textbf{Idea} & \textbf{Costo} \\
  \midrule
  Orden topológico & DFS; vértices en $\text{sale}$ decreciente; un
    vértice gris delata un ciclo & $\Tht{n + m}$ \\
  \addlinespace
  Orden topológico & Kahn: sacar vértices con grado de entrada cero;
    si salen menos de $n$, hay ciclo & $\Tht{n + m}$ \\
  \addlinespace
  Componentes fuertes & DFS en $G$ para los tiempos; DFS en $G^T$ con
    raíces en $\text{sale}$ decreciente & $\Tht{n + m}$ \\
  \addlinespace
  Articulaciones, puentes y biconexas & un DFS con $\text{low}$;
    articulación si un hijo cumple $\text{low}[s] \ge
    \text{entra}[v]$, puente si $\text{low}[v] > \text{entra}[u]$;
    pila de aristas para las componentes & $\Tht{n + m}$ \\
  \addlinespace
  Recorrido de Euler dirigido & fuertemente conexo y balanceado;
    Hierholzer con aristas de salida & $\Tht{n + m}$ \\
  \addlinespace
  Mínimo alcanzable & DFS en $G^T$ en orden creciente de etiqueta
    & $\Tht{n + m}$ \\
  \bottomrule
\end{tabular}
\end{center}

\newpage

% =====================================================================
%  Grafos - Arboles de expansion minima
% =====================================================================

\clase{Árboles de expansión mínima}{Complemento: Cormen, cap.\ 23}
\label{sec:mst}

\begin{quick_sum}
  La clase del 22 de octubre presentó Kruskal con un ejemplo y dejó a
  Prim en una observación. Esta sección da la teoría que los sostiene
  y las dos implementaciones completas. Todo descansa en un teorema,
  la propiedad del corte: si se parten los vértices en dos lados sin
  separar ninguna arista ya elegida, la arista más ligera que cruza de
  un lado al otro se puede agregar sin perder la optimalidad. De ahí
  salen un método genérico, la propiedad del ciclo y la unicidad del
  árbol cuando los pesos son distintos. Kruskal y Prim son dos maneras
  de escoger el corte. Kruskal usa conjuntos disjuntos y cuesta
  $\Ord{m \log n}$, dominado por el ordenamiento; Prim usa una cola de
  prioridad y cuesta $\Ord{m \log n}$ con montículo binario,
  $\Ord{m + n \log n}$ con montículo de Fibonacci y $\Tht{n^2}$ con un
  arreglo. Cierran la sección Prim contra Dijkstra sobre el grafo
  pesado de la libreta y unos ejercicios: un corte, un intercambio, el
  segundo mejor árbol y el árbol de cuello de botella.
\end{quick_sum}

\subsection*{El problema y el método genérico}

La clase del 22 de octubre (sección~\ref{sec:grafos}) definió el
problema: en un grafo $G = (V, E)$ no dirigido y conexo, con un peso
$w : E \to \R$ en cada arista, se busca un conjunto $T \subseteq E$
acíclico que conecte todos los vértices y minimice
\[
  w(T) = \sum_{\set{u, v} \in T} w(u, v) .
\]
Un conjunto así es un árbol y tiene $n - 1$ aristas, igual que
cualquier otro árbol de expansión. Lo que se minimiza no es el número
de aristas, que no cambia, sino su peso: el nombre completo es
\textbf{árbol de expansión de peso mínimo}, y se abrevia MST por
\textit{minimum spanning tree}. Los pesos son reales cualesquiera,
negativos incluidos; a diferencia de Dijkstra, aquí nada depende de su
signo.

Dos convenciones para esta sección. Las letras $e$ y $f$ denotan
aristas, y $T - f + e$ abrevia $(T \setminus \set{f}) \cup \set{e}$,
para no apilar llaves dentro de llaves. Los hechos sobre árboles que
usan las pruebas son las caracterizaciones de la
sección~\ref{sec:hamilton-arboles} (Cormen, teorema B.2); el
lema~\ref{lem:hechos-arboles-mst} reúne los tres que hacen falta.

\subsubsection*{Aristas seguras}

Kruskal y Prim son voraces y agregan una arista a la vez. Lo que
tienen en común es un invariante sobre el conjunto $A$ de aristas
elegidas hasta el momento: \emph{antes de cada iteración, $A$ es
subconjunto de algún árbol de expansión mínima.}

\begin{definicion}[Arista segura]
  \label{def:arista-segura-mst}
  Sea $A \subseteq E$ un subconjunto de algún MST de $G$. Una arista
  $e \notin A$ es \textbf{segura} para $A$ (\textit{safe edge}) si
  $A \cup \set{e}$ también es subconjunto de algún MST.
\end{definicion}

Con esa definición el algoritmo cabe en cuatro líneas, aunque todavía
no se pueda ejecutar:

\begin{algorithm}[H]
\caption{\textsc{MST-Genérico}$(G, w)$}
\label{alg:mst-generico}
\begin{algorithmic}[1]
\Procedure{MST-Genérico}{$G, w$}
  \State $A \gets \emptyset$
  \While{$A$ no forma un árbol de expansión}
    \State encontrar una arista $e$ segura para $A$
    \State $A \gets A \cup \set{e}$
  \EndWhile
  \State \Return $A$
\EndProcedure
\end{algorithmic}
\end{algorithm}

La línea 4 no dice cómo encontrar la arista, y es justo donde Kruskal
y Prim se separan. Escribir primero el método genérico tiene una
ventaja: la corrección se prueba una sola vez, aquí, y a cada
algoritmo concreto solo le queda mostrar que la arista que escoge es
segura.

\textbf{Por qué funciona.} Al inicio $A = \emptyset$, que está
contenido en cualquier MST. Cada iteración agrega una arista segura,
que por definición conserva el invariante. Al salir del ciclo $A$ es
un árbol de expansión contenido en un MST $T$; los dos tienen $n - 1$
aristas, así que $A = T$. El ciclo termina porque cada vuelta agrega
una arista a un conjunto que, al vivir dentro de un árbol, no pasa de
$n - 1$ aristas. Falta ver que la línea 4 siempre tiene algo que
encontrar: el invariante da un MST $T$ con $A \subseteq T$, y como $A$
todavía no es árbol de expansión, la contención es estricta. Cualquier
arista de $T \setminus A$ es segura, porque $A$ con ella sigue dentro
de $T$.

Ese argumento muestra que la arista existe, pero la encuentra dentro
de un MST que nadie conoce. Reconocer una arista segura sin conocer el
árbol es el trabajo del teorema central de la sección.

% =====================================================================

\subsection*{Cortes y la propiedad del corte}

\begin{definicion}[Corte, cruce, arista ligera]
  \label{def:corte-mst}
  Un \textbf{corte} $(S, V - S)$ de $G$ (\textit{cut}) es una partición
  de $V$ en dos conjuntos. La arista $\set{u, v}$ \textbf{cruza} el
  corte (\textit{crosses}) si uno de sus extremos está en $S$ y el otro
  en $V - S$. El corte \textbf{respeta} un conjunto de aristas $A$
  (\textit{respects}) si ninguna arista de $A$ lo cruza. Una
  \textbf{arista ligera} que cruza el corte (\textit{light edge}) es
  una de peso mínimo entre las que lo cruzan.
\end{definicion}

Con empates puede haber varias aristas ligeras para el mismo corte, y
el teorema vale para cualquiera de ellas. ``Ligera'' siempre se dice
respecto a un corte; no tiene nada que ver con las aristas más baratas
del grafo, como muestra el ejemplo.

\begin{ejemplo}[Un corte del grafo de Cormen]
  \label{ej:corte-cormen-mst}
  Este ejemplo y casi todos los que siguen usan el grafo de la figura
  23.1 de Cormen, con $n = 9$ vértices y $m = 14$ aristas. Los
  vértices van en minúscula para no confundirlos con los del grafo
  $G$ de la libreta:
  \[
    \begin{array}{l}
      \set{a,b}\,4,\ \set{a,h}\,8,\ \set{b,c}\,8,\ \set{b,h}\,11,\
      \set{c,d}\,7,\ \set{c,f}\,4,\ \set{c,i}\,2, \\
      \set{d,e}\,9,\ \set{d,f}\,14,\ \set{e,f}\,10,\ \set{f,g}\,2,\
      \set{g,h}\,1,\ \set{g,i}\,6,\ \set{h,i}\,7 .
    \end{array}
  \]
  Se toma $S = \set{a, b, d, e}$, en relleno oscuro, y
  $A = \set{\set{a,b}, \set{c,i}, \set{c,f}, \set{g,h}, \set{f,g}}$,
  en azul grueso. Ninguna arista de $A$ tiene un extremo de cada lado,
  así que el corte $(S, V - S)$ respeta $A$. Las que cruzan,
  punteadas, son seis:
  \begin{center}
  \small
  \begin{tabular}{lcccccc}
    \toprule
    \textbf{Arista} & $\set{c,d}$ & $\set{a,h}$ & $\set{b,c}$
      & $\set{e,f}$ & $\set{b,h}$ & $\set{d,f}$ \\
    \textbf{Peso} & $\mathbf{7}$ & $8$ & $8$ & $10$ & $11$ & $14$ \\
    \textbf{Extremo en $S$} & $d$ & $a$ & $b$ & $e$ & $b$ & $d$ \\
    \bottomrule
  \end{tabular}
  \end{center}
  La arista ligera es $\set{c, d}$, con peso $7$, y es la única. El
  teorema de abajo garantiza que $A \cup \set{\set{c, d}}$ cabe en
  algún árbol de expansión mínima, aunque el corte no tenga nada que
  ver con cómo se construyó $A$. Y $\set{c, d}$ es ligera aunque el
  grafo tenga seis aristas más baratas que ella: ninguna de esas seis
  cruza este corte.
  \begin{center}
  \begin{tikzpicture}[x=1.25cm, y=1.1cm,
      vertS/.style={vert, fill=acento, text=white}]
    \foreach \v/\x/\y in {c/3/1, f/5/-1, g/3/-1, h/1/-1, i/2/0}
      \node[vert] (\v) at (\x, \y) {$\v$};
    \foreach \v/\x/\y in {a/0/0, b/1/1, d/5/1, e/6/0}
      \node[vertS] (\v) at (\x, \y) {$\v$};
    \foreach \u/\v in {g/i, h/i}
      \draw[arista] (\u) -- (\v);
    \foreach \u/\v in {a/b, c/i, c/f, g/h, f/g}
      \draw[aristaArbol] (\u) -- (\v);
    \draw[arista] (d) -- (e);
    \foreach \u/\v in {a/h, b/c, b/h, e/f, d/f}
      \draw[aristaFuera] (\u) -- (\v);
    \draw[aristaResalte] (c) -- (d);
    \foreach \u/\v/\w/\lado in {a/b/4/above left, a/h/8/below left,
        b/c/8/above, b/h/11/left, c/d/7/above, c/f/4/above right,
        c/i/2/above left, d/e/9/above right, d/f/14/left,
        e/f/10/below right, f/g/2/below, g/h/1/below,
        g/i/6/above right, h/i/7/above left}
      \path (\u) -- node[peso, \lado=1pt] {\w} (\v);
    \draw[corte] plot[smooth, tension=0.6] coordinates {
      (-0.8,-0.25) (0.3,-0.35) (1.0,0.4) (2.3,1.0) (3.0,1.65)
      (3.7,1.0) (4.5,0.45) (5.0,0.35) (5.65,-0.4) (6.8,-0.4)};
    \node[nota, anchor=south west] at (-0.8,-0.2) {$S$};
    \node[nota, anchor=north west] at (-0.8,-0.3) {$V - S$};
  \end{tikzpicture}
  \end{center}
\end{ejemplo}

Las tres pruebas que siguen manipulan árboles de la misma manera:
quitan una arista, agregan otra y comprueban que lo que queda sigue
siendo árbol. Conviene tener a la mano por qué.

\begin{lema}[Tres hechos sobre árboles de expansión]
  \label{lem:hechos-arboles-mst}
  Sea $T$ un árbol de expansión de $G$.
  \begin{enumerate}[label=(\alph*), itemsep=0.1em]
    \item Entre cada par de vértices hay un único camino dentro de
          $T$.
    \item Quitar una arista $f = \set{x, y}$ de $T$ deja exactamente
          dos componentes, una con $x$ y otra con $y$. Ellas definen
          un corte, y $f$ es la única arista de $T$ que lo cruza.
    \item Si $e = \set{u, v} \notin T$ y $f$ es cualquier arista del
          camino de $u$ a $v$ en $T$, entonces $T - f + e$ es otra vez
          un árbol de expansión.
  \end{enumerate}
\end{lema}

\begin{proof}
  (a) y (b) son caracterizaciones de árbol de la
  sección~\ref{sec:hamilton-arboles}; la última frase de (b) se debe a
  que las demás aristas de $T$ quedan dentro de alguna de las dos
  componentes. Para (c): por (b), $T - f$ tiene dos componentes. Los
  vértices $u$ y $v$ quedan en componentes distintas, porque si
  siguieran conectados en $T - f$ habría en $T$ un segundo camino de
  $u$ a $v$ que no usa $f$, contra (a). La arista $e$ une esas dos
  componentes, así que $T - f + e$ es conexo, y tiene $n - 1$
  aristas: es un árbol.
\end{proof}

\begin{teorema}[Propiedad del corte; Cormen, teorema 23.1]
  \label{teo:propiedad-corte}
  Sea $G = (V, E)$ un grafo no dirigido y conexo con pesos
  $w : E \to \R$. Sea $A \subseteq E$ un subconjunto de algún árbol de
  expansión mínima de $G$, sea $(S, V - S)$ un corte que respeta $A$ y
  sea $e = \set{u, v}$ una arista ligera que cruza $(S, V - S)$.
  Entonces $e$ es segura para $A$.
\end{teorema}

\begin{proof}
  Sea $T$ un MST con $A \subseteq T$. Si $e \in T$, entonces
  $A \cup \set{e} \subseteq T$ y no hay nada que probar. Supongamos
  que $e \notin T$; vamos a fabricar otro MST que sí la contenga.

  \paso{Una arista de $T$ que cruza.} Sin perder generalidad
  $u \in S$ y $v \in V - S$. Sea $p$ el camino de $u$ a $v$ dentro de
  $T$. Empieza en $S$ y termina fuera de $S$, así que al recorrerlo
  desde $u$ hay una primera arista $f = \set{x, y}$ con $x \in S$ y
  $y \in V - S$. Esa arista cruza el corte. No está en $A$, porque el
  corte respeta $A$. Y es distinta de $e$, porque $f \in T$ y
  $e \notin T$.

  \paso{El intercambio.} Por el lema~\ref{lem:hechos-arboles-mst}(c),
  $T' = T - f + e$ es un árbol de expansión, y pesa
  \[
    w(T') = w(T) - w(f) + w(e) .
  \]
  Las dos aristas $e$ y $f$ cruzan el corte y $e$ es ligera, así que
  $w(e) \le w(f)$ y por lo tanto $w(T') \le w(T)$.

  \paso{Los dos casos.} Si $w(e) < w(f)$, entonces $w(T') < w(T)$, lo
  que es imposible porque $T$ es mínimo: este caso no ocurre. Queda
  el empate, $w(e) = w(f)$. Entonces $w(T') = w(T)$ y $T'$ es otro
  MST: el intercambio no mejora, empata.

  \paso{$T'$ contiene lo que hace falta.} Como $A \subseteq T$ y
  $f \notin A$, se tiene $A \subseteq T - f \subseteq T'$; además
  $e \in T'$. Entonces $A \cup \set{e} \subseteq T'$, con $T'$ un MST,
  y $e$ es segura.
\end{proof}

El tercer paso dice algo más de lo que pide el teorema. Si $e$ no está
en un MST $T$ que contiene a $A$, es porque hay otra arista de cruce
empatada con ella. En particular, cuando la arista ligera de un corte
que respeta $A$ es \emph{única}, está en todo MST que contenga a $A$.
Esa observación es la semilla de la unicidad, más adelante.

\begin{ejemplo}[El intercambio de la prueba]
  \label{ej:intercambio-mst}
  Mismo grafo y mismo $A$, ahora con el corte $S = \set{a, b}$, que
  también respeta $A$. Cruzan $\set{a,h}$ con $8$, $\set{b,c}$ con $8$
  y $\set{b,h}$ con $11$: hay \textbf{dos} aristas ligeras. Sea $T$ el
  árbol que construye Prim desde $a$ (ejemplo~\ref{ej:prim-cormen}),
  de peso $37$, que contiene a $A$ y a $\set{b, c}$ pero no a
  $e = \set{a, h}$.
  \begin{enumerate}[itemsep=0.1em]
    \item El camino $p$ de $a$ a $h$ dentro de $T$ es
          $a, b, c, f, g, h$. Con $e$ cierra un ciclo.
    \item La primera arista de $p$ que sale de $S$ es
          $f = \set{b, c}$; no está en $A$, porque el corte respeta $A$.
    \item $T' = T - f + e$ es otro árbol de expansión, de peso
          $37 - 8 + 8 = 37$.
  \end{enumerate}
  El peso no bajó: empató, porque $w(a,h) = w(b,c)$. Es el caso del
  empate de la prueba, y $T'$ resulta ser justamente el árbol que
  devuelve Kruskal (ejemplo~\ref{ej:kruskal-cormen}). Una enumeración
  de los $662$ árboles de expansión del grafo confirma que solo esos
  dos pesan $37$.
  \begin{center}
  \begin{tikzpicture}[x=0.95cm, y=0.95cm,
      vertS/.style={vert, fill=acento, text=white}]
    \def\pesos{a/b/4/above left, a/h/8/below left, b/c/8/above,
      b/h/11/left, c/d/7/above, c/f/4/above right, c/i/2/above left,
      d/e/9/above right, d/f/14/left, e/f/10/below right, f/g/2/below,
      g/h/1/below, g/i/6/above right, h/i/7/above left}
    \foreach \p/\dx/\pie in {t/0/{$T$: árbol de Prim, peso $37$},
        s/8.6/{$T' = T - f + e$: árbol de Kruskal, peso $37$}} {
      \begin{scope}[xshift=\dx cm, name prefix=\p-]
        \node[font=\small] at (3, -2) {\pie};
        \foreach \v/\x/\y in {c/3/1, d/5/1, e/6/0, f/5/-1, g/3/-1,
                              h/1/-1, i/2/0}
          \node[vert] (\v) at (\x, \y) {$\v$};
        \foreach \v/\x/\y in {a/0/0, b/1/1}
          \node[vertS] (\v) at (\x, \y) {$\v$};
        \foreach \u/\v in {b/h, d/f, e/f, g/i, h/i}
          \draw[arista] (\u) -- (\v);
        \foreach \u/\v in {a/b, c/i, c/f, g/h, f/g, c/d, d/e}
          \draw[aristaArbol] (\u) -- (\v);
        \draw[corte] (0.5, 0.5) ellipse [x radius=1.25, y radius=0.62,
                                         rotate=45];
        \node[nota] at (-0.35, 1.0) {$S$};
      \end{scope}
    }
    \draw[aristaArbol] (t-b) -- (t-c);
    \draw[aristaFuera] (t-a) -- (t-h);
    \draw[arista] (s-b) -- (s-c);
    \draw[aristaArbol] (s-a) -- (s-h);
    \foreach \p in {t, s}
      \foreach \u/\v/\w/\lado in \pesos
        \path (\p-\u) -- node[peso, \lado=1pt] {\w} (\p-\v);
    \path (t-b) -- node[nota, above=8pt] {$f$} (t-c);
    \path (t-a) -- node[nota, below left=5pt] {$e$} (t-h);
  \end{tikzpicture}
  \end{center}
\end{ejemplo}

El teorema deja escoger el corte con total libertad, siempre que
respete $A$. Kruskal y Prim usan el mismo caso particular: el corte
que separa uno de los árboles ya construidos del resto del grafo.

\begin{corolario}[Cormen, corolario 23.2]
  \label{cor:mst-componente}
  Sean $G$ y $A$ como en el teorema~\ref{teo:propiedad-corte} y sea
  $C = (V_C, E_C)$ un componente del bosque $G_A = (V, A)$. Si $e$ es
  una arista ligera entre las que unen $C$ con otro componente de
  $G_A$, entonces $e$ es segura para $A$.
\end{corolario}

\begin{proof}
  El corte $(V_C, V - V_C)$ respeta $A$, porque ninguna arista de $A$
  sale de un componente de $G_A$. Las aristas que cruzan ese corte son
  exactamente las que unen $C$ con otro componente, así que $e$ es
  ligera para él. El teorema hace el resto.
\end{proof}

Del teorema salen también tres datos de la forma del método genérico.
$A$ siempre es acíclico, porque vive dentro de un árbol, así que
$G_A = (V, A)$ es un bosque; al inicio tiene $n$ árboles de un
vértice. Toda arista segura une dos componentes distintos de $G_A$:
una arista con los dos extremos en el mismo componente cerraría un
ciclo con $A$, y ningún MST contiene un ciclo. Y como cada arista
segura junta dos árboles en uno, el ciclo del
algoritmo~\ref{alg:mst-generico} da exactamente $n - 1$ vueltas.

\begin{ojo}
  \textbf{Segura no es lo mismo que ``está en algún MST''.} El
  invariante pide que $A \cup \set{e}$ quepa entero en \emph{un mismo}
  MST. En el grafo de Cormen, $\set{a,h}$ está en un MST y
  $\set{b,c}$ en el otro. Si $A$ son las siete aristas comunes a los
  dos MST más $\set{a,h}$, la arista $\set{b,c}$ ya no es segura,
  aunque pertenezca a un MST: con $A$ cierra el ciclo
  $a, b, c, f, g, h, a$. Por eso el teorema exige que el corte respete
  $A$; es lo que garantiza que las decisiones de pasos distintos
  apunten al mismo árbol. Y por eso el argumento informal ``la arista
  más ligera que cruza un corte está en algún MST'' no alcanza para
  encadenar pasos voraces.
\end{ojo}

\begin{observacion}
  La recíproca del teorema es falsa (Cormen, ejercicio 23.1-2): una
  arista segura no tiene que ser ligera para el corte que se escogió.
  En el camino $a, b, c$ con $w(a,b) = 1$ y $w(b,c) = 2$, con
  $A = \emptyset$ y el corte $(\set{a, c}, \set{b})$, cruzan las dos
  aristas. $\set{b, c}$ es segura, porque el grafo es su propio MST, y
  no es ligera. Lo que sí es cierto (ejercicio 23.1-3) es que toda
  arista de un MST es ligera para \emph{algún} corte: el que deja el
  árbol al quitarle esa arista.
\end{observacion}

% =====================================================================

\subsection*{La propiedad del ciclo y la unicidad}

El teorema del corte dice qué aristas se pueden meter. Su gemelo dice
cuáles se pueden dejar fuera.

\begin{lema}[Propiedad del ciclo; Cormen, ejercicio 23.1-5]
  \label{lem:propiedad-ciclo}
  Sea $e$ una arista de peso máximo en algún ciclo $K$ de $G$.
  Entonces hay un MST de $G$ que no contiene a $e$. Si además $e$ es
  la única de peso máximo en $K$, ningún MST contiene a $e$.
\end{lema}

\begin{proof}
  Si ningún MST contiene a $e$ no hay nada que probar. Sea entonces $T$
  un MST con $e = \set{u, v} \in T$. Por el
  lema~\ref{lem:hechos-arboles-mst}(b), quitar $e$ de $T$ deja dos
  componentes, con $u$ y $v$ en lados distintos, que definen un corte
  cruzado por $e$ y por ninguna otra arista de $T$. El resto del ciclo
  $K$ es un camino de $u$ a $v$: empieza de un lado y termina del
  otro, así que alguna arista $f \neq e$ de $K$ cruza ese corte, y
  $f \notin T$. Entonces $T - e + f$ vuelve a conectar las dos
  componentes; es un árbol de expansión de peso
  $w(T) - w(e) + w(f) \le w(T)$, porque $e$ es máxima en $K$. Es un
  MST sin $e$. Si $e$ es la única máxima, $w(f) < w(e)$ y el árbol
  nuevo pesa menos que $T$, lo que contradice que $T$ sea mínimo; así
  que ningún MST contenía a $e$.
\end{proof}

\begin{ojo}
  La versión que circula, ``la arista más pesada de un ciclo no está en
  ningún MST'', es falsa con empates. En el grafo de Cormen el ciclo
  $a, b, c, i, h, a$ tiene pesos $4, 8, 2, 7, 8$. Sus dos aristas de
  peso máximo, $\set{b,c}$ y $\set{a,h}$, están cada una en uno de los
  dos MST. La versión fuerte pide un máximo estricto: $\set{d,f}$, con
  $14$, es la única máxima del ciclo $c, d, f, c$ (pesos $7, 14, 4$) y,
  en efecto, no está en ningún MST.
\end{ojo}

\begin{teorema}[Unicidad; Cormen, ejercicio 23.1-6 y problema 23-1a]
  \label{teo:unicidad-mst}
  Si todo corte de $G$ tiene una sola arista ligera, en particular si
  todos los pesos son distintos, el MST de $G$ es único.
\end{teorema}

\begin{proof}
  Supongamos que hay dos MST distintos, $T$ y $T'$. Como tienen el
  mismo número de aristas, existe $e = \set{u, v} \in T \setminus T'$.
  Quitar $e$ de $T$ define un corte, y $e$ es ligera para él: si otra
  arista $g$ que lo cruza pesara menos, $T - e + g$ sería un árbol de
  expansión más barato que $T$. Por hipótesis $e$ es la \emph{única}
  ligera de ese corte. Ahora se agrega $e$ a $T'$: cierra un ciclo con
  el camino de $u$ a $v$ en $T'$, y ese camino, que va de un lado del
  corte al otro, contiene alguna arista $f \in T'$ que cruza. Como
  $f \neq e$ y la ligera es única, $w(f) > w(e)$. Por el
  lema~\ref{lem:hechos-arboles-mst}(c), $T' - f + e$ es un árbol de
  expansión que pesa menos que $T'$, contradicción.
\end{proof}

La recíproca es falsa. En el triángulo con $w(x,y) = w(y,z) = 1$ y
$w(x,z) = 2$ el MST $\set{\set{x,y}, \set{y,z}}$ es único, pero el
corte $(\set{y}, \set{x, z})$ tiene dos aristas ligeras. El grafo de
Cormen ilustra el otro lado: tiene pesos repetidos y dos MST, que
difieren en $\set{a,h}$ contra $\set{b,c}$. Aun así comparten la lista
ordenada de pesos, $1, 2, 2, 4, 4, 7, 8, 9$; todos los MST de un grafo
la comparten (Cormen, ejercicio 23.1-8).

\begin{observacion}
  Kruskal y Prim nunca suman pesos: solo los comparan. Por eso los
  pesos negativos no estorban. Por la misma razón, sumar una constante
  $c$ a todos los pesos no cambia cuáles árboles son mínimos, porque
  todo árbol de expansión tiene $n - 1$ aristas y todos suben
  $(n - 1)c$. En caminos más cortos esa suma sí cambia la respuesta,
  porque dos caminos pueden tener distinto número de aristas. Un árbol
  de expansión de peso \emph{máximo} se obtiene negando los pesos.
\end{observacion}

% =====================================================================

\subsection*{Conjuntos disjuntos}

Kruskal pregunta una y otra vez si dos vértices están en el mismo
árbol del bosque $G_A$, mientras los árboles se van fusionando. La
clase del 22 de octubre nombró las tres operaciones que necesita:
\textsc{Crear-Conjunto}$(x)$ forma el conjunto $\set{x}$,
\textsc{Encontrar}$(x)$ devuelve un representante del conjunto de $x$
y \textsc{Unir}$(x, y)$ junta los conjuntos de $x$ y de $y$. El costo
de Kruskal depende de cómo se implementen (Cormen, cap.\ 21).

La implementación rápida guarda cada conjunto como un árbol con raíz,
en el que cada nodo apunta a su padre $p[x]$. La raíz es el
representante y se apunta a sí misma. \textsc{Encontrar} sube hasta la
raíz, y \textsc{Unir} cuelga una raíz de la otra. Sin cuidado se
forman cadenas de largo $n$ y cada \textsc{Encontrar} cuesta
$\Ord{n}$. Dos heurísticas lo evitan:

\begin{itemize}[itemsep=0.1em]
  \item \textbf{Unión por rango.} Cada raíz guarda un $\text{rango}$,
        una cota superior de la altura de su árbol. Al unir, la raíz
        de rango menor cuelga de la de rango mayor; si empatan, se
        escoge una y su rango sube en uno.
  \item \textbf{Compresión de caminos.} \textsc{Encontrar}, al
        regresar de la recursión, deja a cada nodo del camino
        apuntando directo a la raíz. El trabajo de subir una vez
        ahorra las subidas siguientes.
\end{itemize}

\begin{algorithm}[H]
\caption{Conjuntos disjuntos con unión por rango y compresión de
         caminos}
\label{alg:conjuntos-disjuntos}
\begin{algorithmic}[1]
\Procedure{Crear-Conjunto}{$x$}
  \State $p[x] \gets x$, \quad $\text{rango}[x] \gets 0$
\EndProcedure
\Procedure{Encontrar}{$x$}
  \If{$x \neq p[x]$}
    \State $p[x] \gets \textsc{Encontrar}(p[x])$
      \Comment{compresión de caminos}
  \EndIf
  \State \Return $p[x]$
\EndProcedure
\Procedure{Unir}{$x, y$}
  \State $a \gets \textsc{Encontrar}(x)$, \quad
         $b \gets \textsc{Encontrar}(y)$
  \If{$\text{rango}[a] < \text{rango}[b]$}
    \State $p[a] \gets b$ \Comment{el árbol bajo cuelga del alto}
  \Else
    \State $p[b] \gets a$
    \If{$\text{rango}[a] = \text{rango}[b]$}
      \State $\text{rango}[a] \gets \text{rango}[a] + 1$
    \EndIf
  \EndIf
\EndProcedure
\end{algorithmic}
\end{algorithm}

\textsc{Unir} recibe vértices cualesquiera, no representantes, y busca
las raíces por su cuenta. La unión por rango sola ya da una cota
logarítmica, con una prueba de tres líneas.

\begin{lema}[Unión por rango]
  \label{lem:rango-conjuntos}
  Con unión por rango y sin compresión de caminos, un árbol cuya raíz
  tiene rango $r$ tiene altura $r$ y al menos $2^r$ nodos. En
  consecuencia la altura nunca pasa de $\log_2 n$ y \textsc{Encontrar}
  cuesta $\Ord{\log n}$.
\end{lema}

\begin{proof}
  Por inducción sobre las uniones. Un conjunto recién creado tiene
  rango $0$, altura $0$ y un nodo. Al unir raíces $a$ y $b$ de rangos
  distintos, digamos $\text{rango}[b] < \text{rango}[a]$, el árbol de
  $b$ queda un nivel abajo de $a$, con altura
  $\text{rango}[b] + 1 \le \text{rango}[a]$: la altura y el rango de
  $a$ no cambian y sus nodos aumentan. Si los dos rangos valen $r$, el
  árbol nuevo tiene altura $r + 1$, que es el rango nuevo, y al menos
  $2^r + 2^r = 2^{r+1}$ nodos. Como ningún árbol tiene más de $n$
  nodos, $2^r \le n$.
\end{proof}

Con las dos heurísticas juntas la cota baja a casi constante por
operación. La prueba ocupa la sección 21.4 de Cormen y no se repite
aquí; basta el enunciado.

\begin{teorema}[Cormen, teorema 21.14]
  \label{teo:cota-conjuntos}
  Con unión por rango y compresión de caminos, una sucesión de $k$
  operaciones, de las cuales $n$ son \textsc{Crear-Conjunto}, cuesta
  $\Ord{k\,\alpha(n)}$ en el peor caso.
\end{teorema}

Cormen llama $m$ al número de operaciones; aquí es $k$, porque $m$ ya
es el número de aristas. La función $\alpha$ es la inversa de la
función de Ackermann y crece tan despacio que $\alpha(n) \le 4$ para
todo $n \le 10^{80}$, más que el número estimado de átomos del
universo observable. En la práctica es una constante, pero no lo es
en teoría: $\alpha(n) \to \infty$. El ejemplo~\ref{ej:conjuntos-libreta},
en la parte de Kruskal, sigue la estructura durante una corrida
completa.


% =====================================================================

\subsection*{Kruskal}

Kruskal toma, entre todas las aristas que unen dos árboles distintos
del bosque $G_A$, la de menor peso, y para eso basta recorrerlas en
orden de peso. Cada conjunto disjunto guarda los vértices de un árbol
del bosque. El pseudocódigo es el de la clase, dentro de un
procedimiento para que los números de línea coincidan con los que se
citan abajo.

\begin{algorithm}[H]
\caption{\textsc{MST-Kruskal}$(G, w)$}
\label{alg:mst-kruskal}
\begin{algorithmic}[1]
\Procedure{MST-Kruskal}{$G, w$}
  \State $A \gets \emptyset$
  \ForAll{$v \in V$}
    \State \textsc{Crear-Conjunto}$(v)$
  \EndFor
  \State ordenar $E$ por peso, de menor a mayor
  \ForAll{$\set{u, v} \in E$, en ese orden}
    \If{$\textsc{Encontrar}(u) \neq \textsc{Encontrar}(v)$}
      \State $A \gets A \cup \set{\set{u, v}}$
        \Comment{une dos árboles distintos}
      \State \textsc{Unir}$(u, v)$
    \EndIf
  \EndFor
  \State \Return $A$
\EndProcedure
\end{algorithmic}
\end{algorithm}

En la línea 6 el orden es no decreciente y los empates se rompen como
sea; esa decisión cambia qué árbol sale, no su peso. La línea 8 es la
que decide, y su justificación es el corolario.

\begin{proposicion}[Corrección de Kruskal]
  \label{prop:kruskal-mst}
  Si $G$ es conexo, \textsc{MST-Kruskal} devuelve un árbol de
  expansión mínima.
\end{proposicion}

\begin{proof}
  Los conjuntos disjuntos son siempre los vértices de los árboles de
  $G_A$: empiezan como $n$ conjuntos de un vértice, que son los árboles
  de $G_\emptyset$, y la línea 10 junta dos conjuntos justo cuando la
  9 junta dos árboles. Basta ver que cada arista aceptada es segura;
  el método genérico hace el resto.

  Sea $e = \set{u, v}$ la arista examinada. Si $u$ y $v$ están en el
  mismo conjunto, ya hay un camino entre ellos dentro de $A$ y $e$
  cerraría un ciclo; se rechaza. Si están en conjuntos distintos, sea
  $C$ el árbol de $u$. Toda arista $g$ que une $C$ con otro árbol pesa
  al menos $w(e)$. Las que vienen después de $e$ en el orden pesan
  $\ge w(e)$. Las que vinieron antes ya se examinaron, y en ese momento
  o se aceptaron o se rechazaron por tener los extremos juntos; en los
  dos casos sus extremos quedaron en un mismo árbol, y como los árboles
  solo se fusionan, siguen juntos y $g$ no sale de $C$. Entonces $e$ es
  ligera entre las aristas que unen $C$ con otro árbol, y el
  corolario~\ref{cor:mst-componente} la hace segura.

  Al final $A$ es un árbol de expansión. Si $G_A$ tuviera dos árboles,
  alguna arista de $G$ los uniría, porque $G$ es conexo. Esa arista se
  examinó, y sus extremos ya estaban entonces en árboles distintos
  (lo están al final y los árboles solo se fusionan), así que se habría
  aceptado. Contradicción.
\end{proof}

\textbf{Costo.} Con la estructura del
algoritmo~\ref{alg:conjuntos-disjuntos}:
\begin{enumerate}[itemsep=0.1em]
  \item Líneas 2 a 5: $n$ llamadas a \textsc{Crear-Conjunto}, $\Ord{n}$.
  \item Línea 6: un ordenamiento por comparaciones (mergesort,
        heapsort), $\Ord{m \log m}$.
  \item Líneas 7 a 12: $2m$ llamadas a \textsc{Encontrar} y $n - 1$ a
        \textsc{Unir}, una por arista aceptada. Junto con las $n$
        \textsc{Crear-Conjunto} son $k = n + 2m + (n - 1)$ operaciones,
        y por el teorema~\ref{teo:cota-conjuntos} cuestan
        $\Ord{(n + m)\,\alpha(n)}$. Como $G$ es conexo,
        $m \ge n - 1$, y esto es $\Ord{m\,\alpha(n)}$.
  \item Total: $\Ord{m \log m + m\,\alpha(n)}$. Como
        $\alpha(n) = \Ord{\log n}$ y $n \le m + 1$, es
        $\Ord{m \log m}$. Y como $m < n^2$, $\log m < 2 \log n$, así
        que Kruskal cuesta $\Ord{m \log n}$.
\end{enumerate}

El cuello de botella es el ordenamiento, no la estructura. De hecho,
para la cota $\Ord{m \log n}$ basta la unión por rango sola: por el
lema~\ref{lem:rango-conjuntos} las $\Ord{m}$ operaciones cuestan
$\Ord{m \log n}$, que el ordenamiento absorbe. La compresión de
caminos importa cuando ordenar es barato: si las aristas ya vienen
ordenadas, o si los pesos son enteros en un rango chico (de $1$ a $n$,
o de $1$ a una constante), counting sort las ordena en $\Ord{n + m}$ y
Kruskal baja a $\Ord{m\,\alpha(n)}$, casi lineal (Cormen, ejercicio
23.2-4). Al revés, sin conjuntos disjuntos, revisando con BFS dentro
de $A$ si $u$ y $v$ ya están conectados, cada prueba cuesta $\Ord{n}$
y el total sube a $\Ord{mn}$.

La clase del 22 de octubre ya corrió Kruskal sobre el grafo pesado de
la libreta. Vale la pena repetir esa corrida mirando dentro de la
estructura de conjuntos.

\begin{ejemplo}[El bosque de conjuntos en el Kruskal de la clase]
  \label{ej:conjuntos-libreta}
  Las operaciones del ejemplo de Kruskal de la clase sobre el grafo
  pesado, con las aristas en el mismo orden.
  \textsc{Unir}$(B, C)$ encuentra rangos iguales: $C$ cuelga de $B$ y
  $\text{rango}[B] = 1$. En \textsc{Unir}$(A, C)$,
  $\textsc{Encontrar}(C) = B$ y $\text{rango}[A] = 0 < 1$, así que $A$
  cuelga de $B$. Del otro lado, $E$ cuelga de $D$, que sube a rango
  $1$, y luego $F$ cuelga de $D$. La arista $\set{A, B}$ se rechaza
  porque $\textsc{Encontrar}(A) = \textsc{Encontrar}(B) = B$. En
  \textsc{Unir}$(B, D)$ empatan los rangos en $1$: $D$ cuelga de $B$,
  que sube a rango $2$. El pseudocódigo sigue revisando aristas, y al
  examinar $\set{D, F}$ la llamada $\textsc{Encontrar}(F)$ sube por
  $F, D, B$ y, al regresar, deja a $F$ colgado directo de $B$
  (en naranja). Las flechas apuntan al padre.
  \begin{center}
  \begin{tikzpicture}[x=1cm, y=0.9cm]
    % Antes de unir los dos arboles
    \node[vert] (b1) at (0.6, 0) {B};
    \node[vert] (c1) at (0, -1.1) {C};
    \node[vert] (a1) at (1.2, -1.1) {A};
    \node[vert] (d1) at (3.0, 0) {D};
    \node[vert] (e1) at (2.4, -1.1) {E};
    \node[vert] (f1) at (3.6, -1.1) {F};
    \foreach \h/\p in {c1/b1, a1/b1, e1/d1, f1/d1}
      \draw[dirigida] (\h) -- (\p);
    \node[nota, above=1pt] at (b1.north) {rango 1};
    \node[nota, above=1pt] at (d1.north) {rango 1};
    \node[font=\small] at (1.8, -2.85) {antes de $\set{A, B}$};
    % Despues de Unir(B, D)
    \node[vert] (b2) at (6.6, 0) {B};
    \node[vert] (c2) at (5.6, -1.1) {C};
    \node[vert] (a2) at (6.6, -1.1) {A};
    \node[vert] (d2) at (7.6, -1.1) {D};
    \node[vert] (e2) at (7.1, -2.2) {E};
    \node[vert] (f2) at (8.1, -2.2) {F};
    \foreach \h/\p in {c2/b2, a2/b2, d2/b2, e2/d2, f2/d2}
      \draw[dirigida] (\h) -- (\p);
    \node[nota, above=1pt] at (b2.north) {rango 2};
    \node[font=\small] at (6.85, -2.85) {después de \textsc{Unir}$(B, D)$};
    % Despues de Encontrar(F)
    \node[vert] (b3) at (12.0, 0) {B};
    \node[vert] (c3) at (10.7, -1.1) {C};
    \node[vert] (a3) at (11.55, -1.1) {A};
    \node[vert] (d3) at (12.45, -1.1) {D};
    \node[vert] (f3) at (13.3, -1.1) {F};
    \node[vert] (e3) at (12.45, -2.2) {E};
    \foreach \h/\p in {c3/b3, a3/b3, d3/b3, e3/d3}
      \draw[dirigida] (\h) -- (\p);
    \draw[dirigidaResalte] (f3) -- (b3);
    \node[nota, above=1pt] at (b3.north) {rango 2};
    \node[font=\small] at (12.0, -2.85)
      {después de \textsc{Encontrar}$(F)$};
  \end{tikzpicture}
  \end{center}
  El rango de $B$ se queda en $2$ aunque la altura real haya bajado:
  después de una compresión el rango es solo una cota. Al examinar
  $\set{C, E}$ le pasa lo mismo a $E$, y al terminar los otros cinco
  vértices cuelgan directo de $B$.
\end{ejemplo}

El grafo de Cormen tiene más aristas y, sobre todo, un empate que
decide el árbol.

\begin{ejemplo}[Kruskal sobre el grafo de Cormen]
  \label{ej:kruskal-cormen}
  Las $14$ aristas en orden de peso, con los empates en el orden de la
  figura 23.4 de Cormen. La última columna es el bosque $G_A$ después
  del paso; en los rechazos no cambia.
  \begin{center}
  \small
  \begin{tabular}{rcclp{7.2cm}}
    \toprule
    \textbf{Paso} & \textbf{Arista} & \textbf{Peso} & \textbf{¿Entra?}
      & \textbf{Árboles del bosque después} \\
    \midrule
    1 & $\set{h,g}$ & $1$ & sí
      & $\{a\}\ \{b\}\ \{c\}\ \{d\}\ \{e\}\ \{f\}\ \{g,h\}\ \{i\}$ \\
    2 & $\set{c,i}$ & $2$ & sí
      & $\{a\}\ \{b\}\ \{c,i\}\ \{d\}\ \{e\}\ \{f\}\ \{g,h\}$ \\
    3 & $\set{g,f}$ & $2$ & sí
      & $\{a\}\ \{b\}\ \{c,i\}\ \{d\}\ \{e\}\ \{f,g,h\}$ \\
    4 & $\set{a,b}$ & $4$ & sí
      & $\{a,b\}\ \{c,i\}\ \{d\}\ \{e\}\ \{f,g,h\}$ \\
    5 & $\set{c,f}$ & $4$ & sí
      & $\{a,b\}\ \{c,f,g,h,i\}\ \{d\}\ \{e\}$ \\
    6 & $\set{i,g}$ & $6$ & no & cerraría el ciclo $i, c, f, g, i$ \\
    7 & $\set{c,d}$ & $7$ & sí & $\{a,b\}\ \{c,d,f,g,h,i\}\ \{e\}$ \\
    8 & $\set{i,h}$ & $7$ & no & cerraría el ciclo $i, c, f, g, h, i$ \\
    9 & $\set{a,h}$ & $8$ & sí & $\{a,b,c,d,f,g,h,i\}\ \{e\}$ \\
    10 & $\set{b,c}$ & $8$ & no
      & cerraría el ciclo $b, a, h, g, f, c, b$ \\
    11 & $\set{d,e}$ & $9$ & sí & $\{a,b,c,d,e,f,g,h,i\}$ \\
    12--14 & \multicolumn{2}{c}{$\set{e,f}$, $\set{b,h}$, $\set{d,f}$}
      & no & ya hay $n - 1 = 8$ aristas \\
    \bottomrule
  \end{tabular}
  \end{center}
  Peso total $1 + 2 + 2 + 4 + 4 + 7 + 8 + 9 = 37$. En el paso 9 el
  empate entre $\set{a,h}$ y $\set{b,c}$, las dos de peso $8$, se
  rompió a favor de $\set{a,h}$ solo porque aparece antes en la lista
  ordenada. Fueron $28$ llamadas a \textsc{Encontrar} en la línea 8 y
  $8$ a \textsc{Unir}. Las tres instantáneas: después del paso 5,
  después del 8 (en naranja los dos rechazos) y al final.
  \begin{center}
  \begin{tikzpicture}[x=0.78cm, y=0.78cm]
    \def\pesos{a/b/4/above left, a/h/8/below left, b/c/8/above,
      b/h/11/left, c/d/7/above, c/f/4/above right, c/i/2/above left,
      d/e/9/above right, d/f/14/left, e/f/10/below right, f/g/2/below,
      g/h/1/below, g/i/6/above right, h/i/7/above left}
    \foreach \p/\dx/\pie in {k1/0/{pasos 1 a 5}, k2/5.6/{pasos 6 a 8},
                             k3/11.2/{final: peso 37}} {
      \begin{scope}[xshift=\dx cm, name prefix=\p-]
        \node[font=\small] at (3, -1.9) {\pie};
        \foreach \v/\x/\y in {a/0/0, b/1/1, c/3/1, d/5/1, e/6/0, f/5/-1,
                              g/3/-1, h/1/-1, i/2/0}
          \node[vert] (\v) at (\x, \y) {$\v$};
        \foreach \u/\v in {a/b, a/h, b/c, b/h, c/d, c/f, c/i, d/e, d/f,
                           e/f, f/g, g/h, g/i, h/i}
          \draw[arista] (\u) -- (\v);
      \end{scope}
    }
    \foreach \u/\v in {g/h, c/i, f/g, a/b, c/f}
      \draw[aristaArbol] (k1-\u) -- (k1-\v);
    \foreach \u/\v in {g/h, c/i, f/g, a/b, c/f, c/d}
      \draw[aristaArbol] (k2-\u) -- (k2-\v);
    \foreach \u/\v in {g/i, h/i}
      \draw[aristaFuera,
        preaction={draw=white, solid, line width=2pt}]
        (k2-\u) -- (k2-\v);
    \foreach \u/\v in {g/h, c/i, f/g, a/b, c/f, c/d, a/h, d/e}
      \draw[aristaArbol] (k3-\u) -- (k3-\v);
    \foreach \p in {k1, k2, k3}
      \foreach \u/\v/\w/\lado in \pesos
        \path (\p-\u) -- node[peso, \lado=1pt] {\w} (\p-\v);
  \end{tikzpicture}
  \end{center}
\end{ejemplo}

\begin{observacion}
  El pseudocódigo de Cormen recorre las $m$ aristas aunque ya tenga
  $n - 1$; en el ejemplo, las tres últimas solo se examinan para
  rechazarlas. Cortar al llegar a $n - 1$ ahorra llamadas a
  \textsc{Encontrar}, pero no el ordenamiento, que domina. Si $G$ no
  es conexo, Kruskal no se queja: devuelve un \textbf{bosque} de
  expansión mínima, con $n - c$ aristas si $G$ tiene $c$ componentes.
  Y detenido cuando quedan $k$ árboles, es exactamente el
  agrupamiento jerárquico de enlace simple con $k$ grupos.
\end{observacion}

% =====================================================================

\subsection*{Prim}

Prim hace crecer un solo árbol desde una raíz $r$ hasta cubrir $V$. En
cada paso agrega la arista más ligera que une el árbol con un vértice
de afuera: es el corolario~\ref{cor:mst-componente} con $C$ igual al
árbol que se está construyendo. Para escoger esa arista rápido, los
vértices de afuera viven en una cola de prioridad $Q$ con la llave
$\text{llave}[v]$: el peso de la arista más barata que une a $v$ con el
árbol, o $\infty$ si no hay ninguna (Cormen escribe $v.key$). El padre
$\pi[v]$ es el extremo de esa arista dentro del árbol. Es la misma
organización que Dijkstra, con otra llave.

\begin{algorithm}[H]
\caption{\textsc{MST-Prim}$(G, w, r)$}
\label{alg:mst-prim}
\begin{algorithmic}[1]
\Procedure{MST-Prim}{$G, w, r$}
  \ForAll{$u \in V$}
    \State $\text{llave}[u] \gets \infty$, \quad $\pi[u] \gets$ nulo
  \EndFor
  \State $\text{llave}[r] \gets 0$
  \State $Q \gets$ montículo de mínimos con todo $V$, ordenado por
         $\text{llave}$
  \While{$Q \neq \emptyset$}
    \State $u \gets \textsc{Extraer-Min}(Q)$
      \Comment{entra $u$ con la arista $\set{u, \pi[u]}$}
    \ForAll{$v$ adyacente a $u$}
      \If{$v \in Q$ \textbf{y} $w(u, v) < \text{llave}[v]$}
        \State $\pi[v] \gets u$
        \State $\text{llave}[v] \gets w(u, v)$
          \Comment{\textsc{Decrementar-Llave}}
      \EndIf
    \EndFor
  \EndWhile
  \State \Return $\set{\set{v, \pi[v]} : v \in V - \set{r}}$
\EndProcedure
\end{algorithmic}
\end{algorithm}

Las líneas no obvias:
\begin{itemize}[itemsep=0.1em]
  \item \textbf{Líneas 2 a 6.} Todas las llaves en $\infty$ salvo la de
        la raíz, que vale $0$ para que salga primero. En la primera
        vuelta $u = r$ y no entra ninguna arista; $r$ es el único
        vértice con $\pi$ nulo al final.
  \item \textbf{Línea 8.} Los vértices del árbol son los de $V - Q$.
        El que sale tiene la arista más barata hacia el árbol, así que
        $\set{u, \pi[u]}$ es una arista ligera que cruza el corte
        $(V - Q, Q)$; entra al árbol de manera implícita.
  \item \textbf{Línea 10.} La prueba $v \in Q$ no es decorativa: sin
        ella el resultado deja de ser un árbol, como muestra el
        ejemplo~\ref{ej:prim-sin-cola}. Se implementa en $\Ord{1}$ con
        un bit por vértice que se apaga al extraerlo.
  \item \textbf{Línea 12.} Es un \textsc{Decrementar-Llave} sobre el
        montículo. Para hacerlo en $\Ord{\log n}$ el montículo tiene
        que saber en qué posición del arreglo está $v$; se guarda un
        índice por vértice que se actualiza en cada intercambio.
\end{itemize}

\begin{proposicion}[Corrección de Prim]
  \label{prop:prim-mst}
  Si $G$ es conexo, \textsc{MST-Prim} devuelve un árbol de expansión
  mínima.
\end{proposicion}

\begin{proof}
  Sea $A = \set{\set{v, \pi[v]} : v \in V - Q - \set{r}}$, las aristas
  con las que entraron los vértices que ya salieron de $Q$. Antes de
  cada vuelta del \textbf{mientras}, salvo la primera, se cumple:
  \begin{enumerate}[label=(\roman*), itemsep=0.1em]
    \item $A$ es subconjunto de algún MST, y sus aristas tienen los
          dos extremos en $V - Q$;
    \item para todo $v \in Q$, $\text{llave}[v]$ es el peso mínimo de
          una arista entre $v$ y $V - Q$, y $\pi[v]$ es el otro extremo
          de esa arista ($\infty$ y nulo si no hay ninguna).
  \end{enumerate}
  Después de la primera vuelta $V - Q = \set{r}$, $A = \emptyset$ y las
  líneas 9 a 13 pusieron $\text{llave}[v] = w(r, v)$ en los vecinos de
  $r$: las dos partes valen.

  Para (ii): cuando $u$ pasa de $Q$ a $V - Q$, el nuevo mínimo de un
  $v \in Q$ vecino de $u$ es el menor entre $\text{llave}[v]$ y
  $w(u, v)$, y el de un $v$ no vecino no cambia. Las líneas 9 a 13
  calculan exactamente eso, y gracias a la prueba $v \in Q$ no tocan a
  nadie más. Para (i): el $u$ que sale tiene la llave mínima de $Q$,
  que por (ii) es el peso de la arista más ligera que cruza el corte
  $(V - Q, Q)$, y esa arista es $\set{u, \pi[u]}$. La llave es finita
  porque $G$ es conexo y los dos lados del corte son no vacíos. El
  corte respeta $A$ por la segunda mitad de (i), así que el
  teorema~\ref{teo:propiedad-corte} hace segura a $\set{u, \pi[u]}$.
  Sus dos extremos quedan en $V - Q$, y $\pi[u]$ ya no vuelve a
  cambiar, otra vez por la prueba $v \in Q$; el conjunto $A$ de la
  fórmula es entonces el que hace crecer el método genérico.

  Al terminar, $Q = \emptyset$ y $A$ tiene $n - 1$ aristas dentro de
  un MST: es ese MST.
\end{proof}

Cormen enuncia un invariante parecido, en tres partes, pero no lo
prueba en detalle. La prueba deja ver dónde trabaja la línea 10: en
las dos partes del invariante.

\begin{ejemplo}[Prim sobre el grafo de Cormen, desde $a$]
  \label{ej:prim-cormen}
  Cada renglón muestra las llaves después de procesar el vértice que
  sale; en negritas, los que ya están en el árbol, cuya llave ya no
  cambia y es el peso de la arista con la que entraron. Los empates en
  la llave mínima se rompen por orden alfabético, como en la figura
  23.5 de Cormen.
  \begin{center}
  \small
  \setlength{\tabcolsep}{4pt}
  \begin{tabular}{cc*{9}{c}>{\raggedright\arraybackslash}p{5.2cm}}
    \toprule
    \textbf{Paso} & \textbf{Sale} & $a$ & $b$ & $c$ & $d$ & $e$ & $f$
      & $g$ & $h$ & $i$ & \textbf{Cambios} \\
    \midrule
    0 & & $0$ & $\infty$ & $\infty$ & $\infty$ & $\infty$ & $\infty$
      & $\infty$ & $\infty$ & $\infty$ & \\
    1 & $a$ & $\mathbf{0}$ & $4$ & $\infty$ & $\infty$ & $\infty$
      & $\infty$ & $\infty$ & $8$ & $\infty$ & $b$, $h$ desde $a$ \\
    2 & $b$ & $\mathbf{0}$ & $\mathbf{4}$ & $8$ & $\infty$ & $\infty$
      & $\infty$ & $\infty$ & $8$ & $\infty$
      & $c$ desde $b$; $h$: $11 \not< 8$ \\
    3 & $c$ & $\mathbf{0}$ & $\mathbf{4}$ & $\mathbf{8}$ & $7$
      & $\infty$ & $4$ & $\infty$ & $8$ & $2$
      & empate $c$, $h$ en $8$; $d$, $f$, $i$ desde $c$ \\
    4 & $i$ & $\mathbf{0}$ & $\mathbf{4}$ & $\mathbf{8}$ & $7$
      & $\infty$ & $4$ & $6$ & $7$ & $\mathbf{2}$
      & $g$ desde $i$; $h$: $7 < 8$ \\
    5 & $f$ & $\mathbf{0}$ & $\mathbf{4}$ & $\mathbf{8}$ & $7$ & $10$
      & $\mathbf{4}$ & $2$ & $7$ & $\mathbf{2}$
      & $e$ desde $f$; $g$: $2 < 6$ \\
    6 & $g$ & $\mathbf{0}$ & $\mathbf{4}$ & $\mathbf{8}$ & $7$ & $10$
      & $\mathbf{4}$ & $\mathbf{2}$ & $1$ & $\mathbf{2}$
      & $h$: $1 < 7$ \\
    7 & $h$ & $\mathbf{0}$ & $\mathbf{4}$ & $\mathbf{8}$ & $7$ & $10$
      & $\mathbf{4}$ & $\mathbf{2}$ & $\mathbf{1}$ & $\mathbf{2}$
      & nada \\
    8 & $d$ & $\mathbf{0}$ & $\mathbf{4}$ & $\mathbf{8}$ & $\mathbf{7}$
      & $9$ & $\mathbf{4}$ & $\mathbf{2}$ & $\mathbf{1}$ & $\mathbf{2}$
      & $e$: $9 < 10$ \\
    9 & $e$ & $\mathbf{0}$ & $\mathbf{4}$ & $\mathbf{8}$ & $\mathbf{7}$
      & $\mathbf{9}$ & $\mathbf{4}$ & $\mathbf{2}$ & $\mathbf{1}$
      & $\mathbf{2}$ & \\
    \bottomrule
  \end{tabular}
  \end{center}
  Orden de entrada: $a, b, c, i, f, g, h, d, e$. Los padres quedan
  $\pi[b] = a$, $\pi[c] = b$, $\pi[i] = c$, $\pi[f] = c$,
  $\pi[g] = f$, $\pi[h] = g$, $\pi[d] = c$, $\pi[e] = d$, y el peso
  es la suma de las llaves finales,
  $0 + 4 + 8 + 7 + 9 + 4 + 2 + 1 + 2 = 37$. La llave de $h$ cambió
  tres veces, $\infty \to 8 \to 7 \to 1$: desde $a$, desde $i$ y desde
  $g$. En las instantáneas, en verde los vértices del árbol y en
  naranja, para cada vértice de afuera con llave finita, la arista
  $\set{v, \pi[v]}$ que le da esa llave: la más corta de ellas es la
  que entra en el paso siguiente.
  \begin{center}
  \begin{tikzpicture}[x=0.78cm, y=0.78cm]
    \def\pesos{a/b/4/above left, a/h/8/below left, b/c/8/above,
      b/h/11/left, c/d/7/above, c/f/4/above right, c/i/2/above left,
      d/e/9/above right, d/f/14/left, e/f/10/below right, f/g/2/below,
      g/h/1/below, g/i/6/above right, h/i/7/above left}
    \foreach \p/\dx/\dentro/\pie in {p1/0/{a,b,c}/{sale $c$},
        p2/5.6/{a,b,c,i,f,g}/{sale $g$},
        p3/11.2/{a,b,c,d,e,f,g,h,i}/{final: peso 37}} {
      \begin{scope}[xshift=\dx cm, name prefix=\p-]
        \node[font=\small] at (3, -1.9) {\pie};
        \foreach \v/\x/\y in {a/0/0, b/1/1, c/3/1, d/5/1, e/6/0, f/5/-1,
                              g/3/-1, h/1/-1, i/2/0}
          \node[vert] (\v) at (\x, \y) {$\v$};
        \foreach \v in \dentro
          \node[vertOk] at (\v) {$\v$};
        \foreach \u/\v in {a/b, a/h, b/c, b/h, c/d, c/f, c/i, d/e, d/f,
                           e/f, f/g, g/h, g/i, h/i}
          \draw[arista] (\u) -- (\v);
      \end{scope}
    }
    \foreach \u/\v in {a/b, b/c}
      \draw[aristaArbol] (p1-\u) -- (p1-\v);
    \foreach \u/\v in {c/d, c/f, c/i, a/h}
      \draw[aristaFuera,
        preaction={draw=white, solid, line width=2pt}]
        (p1-\u) -- (p1-\v);
    \foreach \u/\v in {a/b, b/c, c/i, c/f, f/g}
      \draw[aristaArbol] (p2-\u) -- (p2-\v);
    \foreach \u/\v in {g/h, c/d, e/f}
      \draw[aristaFuera,
        preaction={draw=white, solid, line width=2pt}]
        (p2-\u) -- (p2-\v);
    \foreach \u/\v in {a/b, b/c, c/i, c/f, f/g, g/h, c/d, d/e}
      \draw[aristaArbol] (p3-\u) -- (p3-\v);
    \foreach \p in {p1, p2, p3}
      \foreach \u/\v/\w/\lado in \pesos
        \path (\p-\u) -- node[peso, \lado=1pt] {\w} (\p-\v);
  \end{tikzpicture}
  \end{center}
  Kruskal y Prim llegan a $37$ por caminos distintos: Kruskal tomó
  $\set{g,h}$ primero y $\set{d,e}$ al último; Prim tomó $\set{a,b}$
  primero. Difieren en una arista. Kruskal se quedó con $\set{a,h}$
  porque aparecía antes en la lista ordenada; Prim con $\set{b,c}$
  porque en el paso 3 la llave de $c$ empató con la de $h$ en $8$ y
  el desempate alfabético sacó a $c$. Con cualquiera de las nueve
  raíces Prim llega a $37$.
\end{ejemplo}

\subsubsection*{La prueba $v \in Q$}

Prim se suele describir como Dijkstra con la llave $w(u, v)$ en lugar
de $d[u] + w(u, v)$. Es cierto, con una diferencia que se paga cara:
en Prim la condición $v \in Q$ de la línea 10 es indispensable.

\begin{ejemplo}[Prim sin la prueba $v \in Q$]
  \label{ej:prim-sin-cola}
  Se corre \textsc{MST-Prim} sobre el grafo de Cormen desde $a$,
  borrando ``$v \in Q$ \textbf{y}'' de la línea 10. Los pasos 1 a 3
  son idénticos a los del ejemplo~\ref{ej:prim-cormen}. En el paso 4
  sale $i$, con llave $2$, y revisa a su vecino $c$, que ya está en el
  árbol con llave $8$ y padre $b$. Como $w(i, c) = 2 < 8$, la línea 11
  hace $\pi[c] = i$, mientras que $\pi[i] = c$. Pasa lo mismo dos
  veces más:
  \begin{center}
  \small
  \begin{tabular}{cccccc}
    \toprule
    \textbf{Paso} & \textbf{Sale} & \textbf{Vértice ya en el árbol}
      & \textbf{Padre} & \textbf{Llave} & \textbf{Ahora} \\
    \midrule
    4 & $i$ & $c$ & $b \to i$ & $8 \to 2$ & $\pi[c] = i$, $\pi[i] = c$ \\
    6 & $g$ & $f$ & $c \to g$ & $4 \to 2$ & $\pi[f] = g$ \\
    7 & $h$ & $g$ & $f \to h$ & $2 \to 1$ & $\pi[g] = h$, $\pi[h] = g$ \\
    \bottomrule
  \end{tabular}
  \end{center}
  El orden de extracción no cambia, porque \textsc{Extraer-Min} solo
  mira las llaves de $Q$, y las que se corrompen son de vértices que ya
  salieron. Lo que se rompe es lo que devuelve la línea 16: los pares
  $\set{v, \pi[v]}$ son $\set{b,a}$, $\set{c,i}$, $\set{i,c}$,
  $\set{d,c}$, $\set{e,d}$, $\set{f,g}$, $\set{g,h}$, $\set{h,g}$.
  Ocho pares pero solo seis aristas distintas, repartidas en tres
  componentes, $\set{a,b}$, $\set{c,d,e,i}$ y $\set{f,g,h}$. No es un
  árbol. La suma de las llaves finales da $28$, menos que el óptimo
  de $37$, y eso solo ya delata el problema: ningún árbol de expansión
  pesa menos que un MST. En la figura, cada flecha va de $v$ a
  $\pi[v]$; en rojo, los tres padres reescritos.
  \begin{center}
  \begin{tikzpicture}[x=1.25cm, y=1.1cm,
      dirigidaAlerta/.style={->, >=stealth, draw=red!70!black,
                             line width=1.4pt}]
    \foreach \v/\x/\y in {a/0/0, b/1/1, c/3/1, d/5/1, e/6/0, f/5/-1,
                          g/3/-1, h/1/-1, i/2/0}
      \node[vert] (\v) at (\x, \y) {$\v$};
    \foreach \u/\v in {a/h, b/c, b/h, d/f, e/f, g/i, h/i, c/f}
      \draw[arista, draw=gray!35] (\u) -- (\v);
    \foreach \h/\p in {b/a, d/c, e/d}
      \draw[dirigidaArbol] (\h) -- (\p);
    \draw[dirigidaArbol] (i) to[bend left=18] (c);
    \draw[dirigidaAlerta] (c) to[bend left=18] (i);
    \draw[dirigidaArbol] (h) to[bend left=18] (g);
    \draw[dirigidaAlerta] (g) to[bend left=18] (h);
    \draw[dirigidaAlerta] (f) -- (g);
  \end{tikzpicture}
  \end{center}
  La arista $\set{u, \pi[u]}$ se elige en el momento en que $u$ sale
  de $Q$, y desde ese momento $\pi[u]$ debe quedarse quieto. La línea
  10 es la que lo protege.
\end{ejemplo}

En Dijkstra la misma prueba se puede omitir sin daño: con pesos no
negativos, todo $v$ que ya salió tiene $d[v] \le d[u]$, así que
$d[u] + w(u, v) \ge d[v]$ y la relajación nunca dispara sobre él. Por
eso el Dijkstra de la clase no la escribe. En Prim la comparación
$w(u, v) < \text{llave}[v]$ no guarda ninguna relación con el orden de
salida: $w(i, c) = 2$ es menor que la llave $8$ con la que $c$ entró
al árbol.

\subsubsection*{Cuánto cuesta Prim}

El costo depende de la cola de prioridad. Con un \textbf{montículo
binario de mínimos}, el de la clase del 17 de septiembre en su versión
de mínimos:
\begin{enumerate}[itemsep=0.1em]
  \item Líneas 2 a 6: inicializar es $\Ord{n}$, y construir el
        montículo también, con el mismo análisis que
        \textsc{Build-Max-Heap}.
  \item El \textbf{mientras} da $n$ vueltas, una por vértice, y cada
        \textsc{Extraer-Min} cuesta $\Ord{\log n}$: $\Ord{n \log n}$.
  \item El \textbf{para todo} de la línea 9 recorre la lista de
        adyacencia de cada vértice una vez, $2m$ iteraciones en total.
        Cada una hace la prueba de la línea 10 en $\Ord{1}$ y, a veces,
        un \textsc{Decrementar-Llave} de $\Ord{\log n}$.
  \item Hay a lo más $m$ decrementos: una arista solo puede bajar una
        llave cuando sale el primero de sus extremos, porque cuando
        sale el segundo el otro ya no está en $Q$. En el grafo de
        Cormen hubo $12$ decrementos para $m = 14$; en el pesado de la
        libreta, $9$ para $m = 10$.
  \item Total: $\Ord{n + n \log n + m + m \log n} = \Ord{(n + m) \log n}$,
        que con $G$ conexo es $\Ord{m \log n}$. Asintóticamente igual
        que Kruskal.
\end{enumerate}

Con un \textbf{montículo de Fibonacci} (Cormen, cap.\ 19),
\textsc{Extraer-Min} cuesta $\Ord{\log n}$ amortizado y
\textsc{Decrementar-Llave} $\Ord{1}$ amortizado. ``Amortizado''
significa que la suma de los costos de toda la sucesión está acotada
por la suma de esas cotas, aunque una operación suelta cueste más. Los
términos quedan $n \cdot \Ord{\log n}$ por las extracciones y
$m \cdot \Ord{1}$ por los decrementos: $\Ord{m + n \log n}$. Lo que se
ganó es que el término que multiplica a $m$ ya no lleva $\log n$.

Con un \textbf{arreglo} (Cormen, ejercicio 23.2-2), conveniente cuando
el grafo viene como matriz de adyacencia, la llave se guarda en un
arreglo indexado por vértice. Buscar el mínimo es un barrido de
$\Tht{n}$ y decrementar es una asignación de $\Ord{1}$. Con la matriz,
revisar los vecinos de $u$ es otro barrido de $\Tht{n}$. Son $n$
vueltas de $\Tht{n}$: $\Tht{n^2}$, sin importar $m$. La prueba
$v \in Q$ se vuelve un arreglo booleano $\text{dentro}[v]$, y por la
misma razón no se puede quitar. La tabla~\ref{tab:costos-mst} reúne
las cuatro opciones (Cormen, ejercicio 23.2-3).

\begin{table}[H]
\centering
\small
\begin{tabular}{llll}
  \toprule
  \textbf{Implementación} & \textbf{Costo}
    & \textbf{Disperso, $m = \Tht{n}$} & \textbf{Denso, $m = \Tht{n^2}$} \\
  \midrule
  Kruskal, conjuntos disjuntos & $\Ord{m \log n}$ & $\Ord{n \log n}$
    & $\Ord{n^2 \log n}$ \\
  Prim, montículo binario & $\Ord{m \log n}$ & $\Ord{n \log n}$
    & $\Ord{n^2 \log n}$ \\
  Prim, montículo de Fibonacci & $\Ord{m + n \log n}$
    & $\Ord{n \log n}$ & $\Ord{n^2}$ \\
  Prim, arreglo o matriz & $\Tht{n^2}$ & $\Tht{n^2}$ & $\Tht{n^2}$ \\
  \bottomrule
\end{tabular}
\caption{Costo de los algoritmos de árbol de expansión mínima según la
         densidad del grafo.}
\label{tab:costos-mst}
\end{table}

La tabla se lee así. Fibonacci le gana asintóticamente al binario
exactamente cuando $m = \omg{n}$: en ese caso $m \log n$ domina tanto
a $m$ como a $n \log n$. Cuando $m = \Omg{n \log n}$, Prim con
Fibonacci ya es $\Ord{m}$, lineal. En grafos densos el arreglo empata
con Fibonacci y tiene constantes mucho menores, así que es la opción
práctica. En grafos dispersos todos dan $\Ord{n \log n}$ y deciden las
constantes: Kruskal es el más simple de programar y es casi lineal si
los pesos se ordenan en tiempo lineal; Prim con montículo binario es
lo usual con listas de adyacencia. Para dar una idea de tamaños, con
$n = 1000$ y sin constantes: en un grafo disperso con $m = 3000$,
$m \log_2 n \approx 3 \cdot 10^4$ contra $n^2 = 10^6$, y gana el
montículo; en el grafo completo, con $m = 499\,500$,
$m \log_2 n \approx 5 \cdot 10^6$ contra $n^2 = 10^6$, y gana el
arreglo.

\begin{ojo}
  \textbf{El montículo de la biblioteca no trae
  \textsc{Decrementar-Llave}.} Esto no viene en Cormen, pero es lo
  primero que aparece al programarlo: \texttt{heapq} de Python y
  \texttt{priority\_queue} de C++ no decrementan. La salida usual es la
  versión perezosa. En lugar de decrementar se inserta otra copia
  $(\text{llave}, v)$, y al extraer se descartan los $v$ que ya están
  en el árbol; ese descarte es la prueba $v \in Q$ con otro disfraz. El
  montículo llega a tener $\Ord{m}$ entradas y el costo es
  $\Ord{m \log m} = \Ord{m \log n}$, la misma cota.
\end{ojo}

\begin{ojo}
  \textbf{Grafos desconexos y dirigidos.} Si $G$ no es conexo, Prim
  desde $r$ solo cubre el componente de $r$: los demás vértices salen
  de $Q$ con llave $\infty$ y padre nulo, y el pseudocódigo no lo
  detecta. Basta revisar si la llave extraída es $\infty$. Con grafos
  dirigidos el capítulo no aplica, porque el corte y el ciclo usan que
  la arista sirve en los dos sentidos. Con raíz $r$ y arcos $r \to a$
  de peso $2$, $r \to b$ de $3$, $a \to b$ de $3$ y $b \to a$ de $1$,
  tomar siempre el arco más barato que sale del árbol da $r \to a$ y
  luego un arco de peso $3$ hacia $b$, total $5$; la arborescencia
  óptima es $r \to b$, $b \to a$, con peso $4$. El problema dirigido
  tiene su propio algoritmo (Chu-Liu y Edmonds), fuera de este curso.
\end{ojo}

\subsubsection*{Prim contra Dijkstra}

La clase del 22 de octubre comparó el árbol de Dijkstra desde A
(peso $15$) con el MST (peso $13$) en el grafo pesado. Corriendo los
dos algoritmos lado a lado se ve en qué paso se separan.

\begin{ejemplo}[Prim contra Dijkstra en el grafo pesado de la libreta]
  \label{ej:prim-dijkstra}
  Desde A, con el grafo pesado de la clase. En cada celda,
  $\text{llave}$ de Prim / $d$ de Dijkstra; en negritas, los vértices
  ya fijados. A tiene $0 / 0$ todo el tiempo.
  \begin{center}
  \small
  \begin{tabular}{ccccccc}
    \toprule
    \textbf{Paso} & \textbf{Sale} & B & C & D & E & F \\
    \midrule
    1 & A & $4 / 4$ & $2 / 2$ & $7 / 7$ & $\infty / \infty$
      & $\infty / \infty$ \\
    2 & C & $1 / 3$ & $\mathbf{2 / 2}$ & $7 / 7$ & $10 / 12$
      & $\infty / \infty$ \\
    3 & B & $\mathbf{1 / 3}$ & $\mathbf{2 / 2}$ & $5 / 7$ & $10 / 12$
      & $\infty / \infty$ \\
    4 & D & $\mathbf{1 / 3}$ & $\mathbf{2 / 2}$ & $\mathbf{5 / 7}$
      & $2 / 9$ & $6 / 13$ \\
    5 & E & $\mathbf{1 / 3}$ & $\mathbf{2 / 2}$ & $\mathbf{5 / 7}$
      & $\mathbf{2 / 9}$ & $3 / 12$ \\
    6 & F & $\mathbf{1 / 3}$ & $\mathbf{2 / 2}$ & $\mathbf{5 / 7}$
      & $\mathbf{2 / 9}$ & $\mathbf{3 / 12}$ \\
    \bottomrule
  \end{tabular}
  \end{center}
  Los dos sacan los vértices en el mismo orden, A, C, B, D, E, F, y aun
  así construyen árboles distintos. La diferencia está en el paso 3:
  al salir B, Prim compara $w(B, D) = 5$ contra la llave $7$ y cambia
  el padre de D a B; Dijkstra compara $d[B] + w(B, D) = 8$ contra $7$
  y lo deja colgado de A. Prim termina con el MST de peso
  $0 + 2 + 1 + 5 + 2 + 3 = 13$, el mismo que encontró Kruskal;
  Dijkstra, con el árbol de caminos más cortos de peso $15$. En azul el
  árbol de Prim, punteada la arista que Dijkstra usa en lugar de
  $\set{B, D}$, y junto a cada vértice su llave final y su $d$ final.
  \begin{center}
  \begin{tikzpicture}
    \node[vert] (A) at (0, 0) {A};
    \node[vert] (B) at (2, 1.4) {B};
    \node[vert] (C) at (2, -1.4) {C};
    \node[vert] (D) at (4.5, 1.4) {D};
    \node[vert] (E) at (4.5, -1.4) {E};
    \node[vert] (F) at (6.5, 0) {F};
    \foreach \a/\b/\w in {A/B/4, C/D/8, C/E/10, D/F/6}
      \draw[arista] (\a) -- node[peso] {\w} (\b);
    \foreach \a/\b/\w in {A/C/2, C/B/1, B/D/5, D/E/2, E/F/3}
      \draw[aristaArbol] (\a) -- node[peso] {\w} (\b);
    \draw[aristaFuera] (A) to[out=75, in=150]
      node[peso, pos=0.3] {7} (D);
    \foreach \v/\k/\dd/\ancla/\lado in {A/0/0/west/left,
        B/1/3/south east/below right, C/2/2/south/below,
        D/5/7/north east/above right, E/2/9/south/below,
        F/3/12/east/right}
      \node[nota, \lado=1pt] at (\v.\ancla) {$\k / \dd$};
  \end{tikzpicture}
  \end{center}
\end{ejemplo}

Prim minimiza el peso total y Dijkstra la distancia desde la raíz, así
que sus árboles no tienen por qué parecerse aunque el código sí. Otra
diferencia: el árbol de Dijkstra depende de la raíz por definición; el
peso del de Prim no (en el grafo de Cormen sale $37$ desde cualquiera
de los nueve vértices). Y un dato curioso: el árbol de Dijkstra desde A
es uno de los dos \emph{segundos mejores} árboles de expansión de este
grafo (ejercicio~\ref{ej:segundo-mejor-mst}).

% =====================================================================

\subsection*{Ejercicios}

\begin{ejercicio}[El corte dibujado]
  \label{ej:corte-dibujado-mst}
  En el grafo pesado de la libreta sea
  $X = \set{\set{B, C}, \set{D, E}}$, contenido en el MST de peso
  $13$; hace el papel de $A$, letra que aquí ya nombra un vértice.
  Para cada corte, con $S_1 = \set{A, B, C}$, $S_2 = \set{A, B, D}$ y
  $S_3 = \set{A}$, di si respeta $X$; si lo respeta, lista las aristas
  que cruzan, encuentra la ligera y di qué garantiza el
  teorema~\ref{teo:propiedad-corte}. Dibuja el primero.
\end{ejercicio}

\textbf{Solución.} $S_1$ respeta $X$: $B$ y $C$ quedan adentro, $D$ y
$E$ afuera. Cruzan $\set{B, D}$ con $5$, $\set{A, D}$ con $7$,
$\set{C, D}$ con $8$ y $\set{C, E}$ con $10$. La ligera es
$\set{B, D}$, única, así que $X \cup \set{\set{B, D}}$ está en un MST,
y de hecho en todo MST que contenga a $X$. $S_2$ no respeta $X$:
$\set{B, C}$ lo cruza, y $\set{D, E}$ también. El teorema no dice
nada, y se nota por qué: la ligera de ese corte es $\set{B, C}$, que ya
está en $X$. $S_3$ respeta $X$, porque ninguna arista de $X$ toca al
vértice A. Cruzan $\set{A, B}$ con $4$, $\set{A, C}$ con $2$ y
$\set{A, D}$ con $7$; la ligera es $\set{A, C}$. Es el corte que usa
Prim desde A en su primer paso. El dibujo de $S_1$, con $X$ en azul,
las aristas de cruce punteadas y la ligera resaltada:
\begin{center}
\begin{tikzpicture}[vertS/.style={vert, fill=acento, text=white}]
  \node[vertS] (A) at (0, 0) {A};
  \node[vertS] (B) at (2, 1.4) {B};
  \node[vertS] (C) at (2, -1.4) {C};
  \node[vert] (D) at (4.5, 1.4) {D};
  \node[vert] (E) at (4.5, -1.4) {E};
  \node[vert] (F) at (6.5, 0) {F};
  \foreach \a/\b/\w in {A/B/4, A/C/2, D/F/6, E/F/3}
    \draw[arista] (\a) -- node[peso] {\w} (\b);
  \foreach \a/\b/\w in {B/C/1, D/E/2}
    \draw[aristaArbol] (\a) -- node[peso] {\w} (\b);
  \foreach \a/\b/\w in {C/D/8, C/E/10}
    \draw[aristaFuera] (\a) -- node[peso] {\w} (\b);
  \draw[aristaFuera] (A) to[out=75, in=150]
    node[peso, pos=0.3] {7} (D);
  \draw[aristaResalte] (B) -- node[peso] {5} (D);
  \draw[corte] (2.75, 2.6) -- (2.75, -2.0);
  \node[nota, anchor=east] at (2.65, -1.95) {$S_1$};
  \node[nota, anchor=west] at (2.85, -1.95) {$V - S_1$};
\end{tikzpicture}
\end{center}

\begin{ejercicio}[El intercambio de la prueba]
  \label{ej:intercambio-dijkstra-mst}
  El árbol de Dijkstra desde A en el grafo pesado es
  $T_D = \set{\set{A,C}, \set{A,D}, \set{B,C}, \set{D,E}, \set{E,F}}$,
  de peso $15$, y contiene al $X$ del ejercicio anterior. Aplica a
  $T_D$ los pasos de la prueba del teorema~\ref{teo:propiedad-corte}
  con el corte $S_1$ y su arista ligera $e = \set{B, D}$: encuentra el
  camino $p$, la arista $f$, el árbol $T_D - f + e$ y su peso. ¿Por qué
  no contradice a la prueba que el peso baje?
\end{ejercicio}

\textbf{Solución.} El camino de $B$ a $D$ en $T_D$ es $B, C, A, D$.
Sale de $S_1$ por primera vez en $f = \set{A, D}$, porque $B$, $C$ y
$A$ están en $S_1$ y $D$ no; $f \notin X$, como debe ser. El árbol
$T_D - f + e$ tiene las aristas $\set{A,C}$, $\set{B,C}$, $\set{B,D}$,
$\set{D,E}$ y $\set{E,F}$, y pesa $15 - 7 + 5 = 13$: es el MST. La
prueba suponía que
$T$ era mínimo y de ahí concluía que una baja estricta es imposible.
Aquí la baja ocurre, y lo único que demuestra es que $T_D$ no era
mínimo. Visto así, el intercambio es un paso de mejora que sirve con
cualquier árbol de expansión que contenga a $X$: nunca sube el peso, y
si baja, el árbol de partida no era óptimo.

\begin{ejercicio}[Segundo mejor árbol de expansión; Cormen, problema
  23-1]
  \label{ej:segundo-mejor-mst}
  Supón que $G$ tiene un único MST $T$. Un \textbf{segundo mejor}
  árbol de expansión es uno de peso mínimo entre los árboles de
  expansión distintos de $T$. Para $u, v \in V$ sea $\max[u, v]$ una
  arista de peso máximo en el camino de $u$ a $v$ dentro de $T$.
  \begin{enumerate}[label=(\alph*), itemsep=0.1em]
    \item Muestra con un ejemplo que el segundo mejor no tiene que ser
          único.
    \item Con pesos distintos, prueba que algún segundo mejor se
          obtiene de $T$ con un solo intercambio: agregar una arista
          $e \notin T$ y quitar una arista $f$ del ciclo que cierra.
    \item Calcula $\max[u, v]$ para todos los pares en $\Ord{n^2}$.
    \item Da un algoritmo para el segundo mejor y aplícalo al grafo
          pesado de la libreta.
  \end{enumerate}
\end{ejercicio}

\textbf{Solución.} (b) Sea $T_2$ un segundo mejor y sea $e$ la arista
más ligera de la diferencia simétrica entre $T$ y $T_2$. Si $e$
estuviera en $T_2 \setminus T$, al agregarla a $T$ cerraría un ciclo
con alguna arista $g \in T \setminus T_2$, más pesada que $e$, y
$T - g + e$ sería más barato que $T$: imposible. Entonces
$e \in T \setminus T_2$. Agregada a $T_2$ cierra un ciclo con alguna
$f \in T_2 \setminus T$, más pesada que $e$, y $T_2 - f + e$ pesa
menos que $T_2$. Como $T_2$ es segundo mejor, ese árbol tiene que ser
$T$, y entonces $T_2 = T - e + f$: un solo intercambio. (c) Un BFS
dentro de $T$ desde cada $u$: si $q$ es el padre de $v$ en ese BFS,
$\max[u, v]$ es la más pesada entre $\max[u, q]$ y $\set{q, v}$. Cada
BFS cuesta $\Ord{n}$ porque $T$ tiene $n - 1$ aristas, y son $n$
recorridos. (d) Se calcula $T$, la tabla $\max$, y para cada arista
$\set{x, y} \notin T$ el candidato
$w(T) - w(\max[x, y]) + w(x, y)$; el menor es el segundo mejor. Cuesta
$\Ord{m \log n + n^2}$, o $\Tht{n^2}$ calculando $T$ con Prim sobre
un arreglo. En el grafo pesado, con el MST
$\set{\set{A,C}, \set{B,C}, \set{B,D}, \set{D,E}, \set{E,F}}$ de peso
$13$:
\begin{center}
\small
\begin{tabular}{cccc}
  \toprule
  \textbf{Arista fuera de $T$} & \textbf{Camino en $T$}
    & \textbf{$\max$ en el camino} & \textbf{Candidato} \\
  \midrule
  $\set{A,B}$, $4$ & A, C, B & $\set{A,C}$, $2$ & $13 - 2 + 4 = 15$ \\
  $\set{A,D}$, $7$ & A, C, B, D & $\set{B,D}$, $5$
    & $13 - 5 + 7 = 15$ \\
  $\set{C,D}$, $8$ & C, B, D & $\set{B,D}$, $5$ & $13 - 5 + 8 = 16$ \\
  $\set{D,F}$, $6$ & D, E, F & $\set{E,F}$, $3$ & $13 - 3 + 6 = 16$ \\
  $\set{C,E}$, $10$ & C, B, D, E & $\set{B,D}$, $5$
    & $13 - 5 + 10 = 18$ \\
  \bottomrule
\end{tabular}
\end{center}
El segundo mejor pesa $15$ y hay dos, lo que responde (a). El que
cambia $\set{B,D}$ por $\set{A,D}$ es exactamente el árbol de Dijkstra
desde A. Los pesos de este grafo no son todos distintos ($\set{A,C}$ y
$\set{D,E}$ pesan $2$), pero su MST es único y la enumeración de sus
$104$ árboles de expansión confirma la tabla.

\begin{ejercicio}[Árbol de cuello de botella; Cormen, problema 23-3]
  \label{ej:cuello-botella-mst}
  Un \textbf{árbol de expansión de cuello de botella} es uno cuya
  arista más pesada es lo más ligera posible; su \textbf{valor} es el
  peso de esa arista.
  \begin{enumerate}[label=(\alph*), itemsep=0.1em]
    \item Prueba que todo MST es un árbol de cuello de botella.
    \item Da un algoritmo de $\Ord{n + m}$ que decida si el valor es a
          lo más $b$.
    \item Usa (b) para hallar el valor en tiempo lineal.
    \item ¿Es cierta la recíproca de (a)?
  \end{enumerate}
\end{ejercicio}

\textbf{Solución.} (a) Sea $e$ la arista más pesada de un MST $T$ y
supón que otro árbol $T'$ tiene todas sus aristas más ligeras que $e$.
Quitar $e$ de $T$ define un corte; como $T'$ es conexo, alguna arista
$g \in T'$ lo cruza, con $w(g) < w(e)$, y $T - e + g$ pesa menos que
$T$. Contradicción. (b) Se quedan solo las aristas de peso $\le b$ y se
revisa con BFS o DFS si el grafo sigue conexo: $\Ord{n + m}$.
(c) Es el algoritmo de Camerini. Con selección lineal (Cormen, sección
9.3) se halla la mediana de los pesos. Si las aristas de la mitad
ligera conectan el grafo, se recurre solo sobre ellas; si no, se
contraen sus componentes a supervértices y se recurre sobre la mitad
pesada. Cada nivel cuesta $\Ord{m}$ y deja la mitad de las aristas:
$t(m) = t(m/2) + \Ord{m} = \Ord{m}$. (d) No. En el grafo de Cormen el
valor es $9$, porque el vértice $e$ solo tiene aristas de $9$ y $10$,
y hay $44$ árboles de valor $9$ que no son MST. Por ejemplo, el de
Kruskal cambiando $\set{c,f}$ por $\set{i,g}$: sigue siendo árbol, su
arista más pesada sigue siendo $\set{d,e}$, y pesa $39$.

\newpage

% =====================================================================
%  Grafos - Caminos más cortos desde una fuente
% =====================================================================

\clase{Caminos más cortos desde una fuente}
  {Complemento: Cormen, cap.\ 24.1 a 24.4}
\label{sec:caminos-fuente}

\begin{quick_sum}
  El problema: un grafo dirigido con pesos reales en las aristas, que
  pueden ser negativos, y un vértice fuente $s$; se quiere el peso
  $\delta(s, v)$ de un camino más corto de $s$ a cada vértice, y el
  camino. La sección arranca con los fundamentos: qué vale $\delta$
  cuando hay ciclos de peso negativo, por qué un camino más corto se
  puede tomar sin ciclos y cómo se guardan todos los caminos en un
  árbol. Sigue la herramienta común a todos los algoritmos, la
  relajación, con las seis propiedades que la gobiernan. Sobre ella se
  montan tres algoritmos que solo difieren en el orden en que relajan:
  Bellman-Ford, para pesos cualesquiera, que detecta los ciclos
  negativos en $\Tht{nm}$; el de los grafos acíclicos, que relaja en
  orden topológico en $\Tht{n + m}$ y de paso calcula rutas críticas;
  y Dijkstra, con su prueba completa y su costo según la cola de
  prioridad. Al final, los sistemas de restricciones de diferencia
  muestran que Bellman-Ford también resuelve una familia de programas
  lineales.
\end{quick_sum}

% =====================================================================

\subsection*{El problema}

La clase del 22 de octubre resolvió dos casos particulares del
problema de esta sección: BFS, cuando todas las aristas valen lo
mismo, y Dijkstra, cuando ningún peso es negativo. Aquí va el caso
general, con pesos reales cualesquiera. En toda la sección
$G = (V, E)$ es un grafo dirigido guardado en listas de adyacencia, y
cada entrada de la lista trae el peso de su arista, así que leer
$w(u, v)$ cuesta $\Ord{1}$. El peso es cualquier cantidad que se
acumule sumando a lo largo de un recorrido: distancia, tiempo, costo,
pérdida.

\begin{definicion}[Peso de una caminata y peso de camino más corto]
  \label{def:distancia-fuente}
  El \textbf{peso} de una caminata
  $p = \langle v_0, v_1, \dots, v_k \rangle$ es la suma de los pesos
  de sus aristas,
  \[
    w(p) = \sum_{i=1}^{k} w(v_{i-1}, v_i) .
  \]
  El \textbf{peso de camino más corto} (\textit{shortest-path
  weight}) de $u$ a $v$ es
  \[
    \delta(u, v) =
    \begin{cases}
      \inf \set{w(p) : p \text{ es una caminata de } u
        \text{ a } v} & \text{si hay alguna,} \\
      \infty & \text{si no.}
    \end{cases}
  \]
  Un \textbf{camino más corto} de $u$ a $v$ es un camino $p$ de $u$ a
  $v$ con $w(p) = \delta(u, v)$. No tiene por qué ser único.
\end{definicion}

\begin{observacion}
  Cormen llama \textit{path} a lo que aquí es una caminata, que puede
  repetir vértices, y \textit{simple path} a lo que aquí es un camino.
  El ínfimo se toma sobre caminatas a propósito: si una caminata de
  $u$ a $v$ puede pasar por un ciclo de peso negativo, las que le dan
  más vueltas pesan cada vez menos, y eso tiene que reflejarse en
  $\delta$. Cuando no hay ciclos negativos de por medio, el ínfimo lo
  alcanza un camino (proposición \ref{prop:caminos-simples-fuente}) y
  las dos lecturas coinciden.
\end{observacion}

Con la fuente única se resuelven las otras variantes del problema:
\begin{itemize}[itemsep=0.1em]
  \item \textbf{Fuente única} (\textit{single-source}): de un vértice
        $s$ fijo a todos los demás. Es el problema de esta sección.
  \item \textbf{Destino único}: de todos los vértices a un destino
        $t$ fijo. Se invierte cada arista y se resuelve fuente única
        desde $t$.
  \item \textbf{Par único}: de $u$ a $v$. Se resuelve fuente única
        desde $u$. No se conoce nada asintóticamente mejor en el peor
        caso: saber el destino no ayuda.
  \item \textbf{Todos los pares}: basta correr fuente única desde
        cada vértice, pero hay métodos mejores (sección
        \ref{sec:todos-pares}).
\end{itemize}

\subsubsection*{Subestructura óptima}

\begin{lema}[Los subcaminos de un camino más corto son caminos más
  cortos; Cormen, lema 24.1]\label{lem:subcaminos-cortos}
  Sea $p = \langle v_0, v_1, \dots, v_k\rangle$ un camino más corto
  de $v_0$ a $v_k$. Para todo $0 \le i \le j \le k$, el subcamino
  $p_{ij} = \langle v_i, v_{i+1}, \dots, v_j\rangle$ es un camino más
  corto de $v_i$ a $v_j$.
\end{lema}

\begin{proof}
  Se parte $p$ en tres tramos,
  $v_0 \overset{p_{0i}}{\leadsto} v_i \overset{p_{ij}}{\leadsto} v_j
  \overset{p_{jk}}{\leadsto} v_k$, de modo que
  $w(p) = w(p_{0i}) + w(p_{ij}) + w(p_{jk})$. Si hubiera una caminata
  $p'$ de $v_i$ a $v_j$ con $w(p') < w(p_{ij})$, la caminata que
  recorre $p_{0i}$, luego $p'$ y luego $p_{jk}$ iría de $v_0$ a $v_k$
  con peso
  \[
    w(p_{0i}) + w(p') + w(p_{jk}) < w(p) = \delta(v_0, v_k) ,
  \]
  imposible porque $\delta$ es el ínfimo sobre todas las caminatas.
  Así que ninguna caminata de $v_i$ a $v_j$ pesa menos que $p_{ij}$, y
  $w(p_{ij}) = \delta(v_i, v_j)$.
\end{proof}

Es la \textbf{subestructura óptima} de la programación dinámica: una
solución óptima contiene soluciones óptimas de sus subproblemas.
Cormen apoya en este lema tanto a Dijkstra, que es voraz, como a
Floyd-Warshall, que es programación dinámica.

\subsubsection*{Pesos negativos y ciclos negativos}

Una arista de peso negativo no rompe nada por sí sola. Lo que rompe la
definición es un \textbf{ciclo de peso negativo}
(\textit{negative-weight cycle}) en alguna caminata de $s$ a $v$: cada
vuelta adicional baja el peso, el ínfimo es $\delta(s, v) = -\infty$ y
ningún camino es el más corto. Si $v$ no es alcanzable desde $s$, en
cambio, $\delta(s, v) = \infty$, aunque $v$ esté en un ciclo negativo.
Para que las cuentas valgan con infinitos se fija la aritmética
\[
  a + \infty = \infty \quad (a \neq -\infty) , \qquad
  a + (-\infty) = -\infty \quad (a \neq \infty) .
\]
La segunda regla solo hace falta en las pruebas, cuando hay ciclos
negativos.

\begin{ejemplo}[Los tres comportamientos; Cormen, figura 24.1]
  \label{ej:ciclos-negativos-fuente}
  En la figura, junto a cada vértice está $\delta(s, \cdot)$, y en
  naranja los dos ciclos de peso negativo.
  \begin{center}
  \begin{tikzpicture}[distancia/.style={font=\scriptsize\bfseries,
      text=acento, inner sep=1.5pt}]
    \node[vert, label={[distancia]left:$0$}] (s) at (0, 0) {$s$};
    \node[vert, label={[distancia]above:$3$}] (a) at (1.7, 1.3)
      {$a$};
    \node[vert, label={[distancia]above:$-1$}] (b) at (3.4, 1.3)
      {$b$};
    \node[vert, label={[distancia]above left:$5$}] (c) at (1.7, 0)
      {$c$};
    \node[vert, label={[distancia]above right:$11$}] (d) at (3.4, 0)
      {$d$};
    \node[vert, label={[distancia]below:$-\infty$}] (e)
      at (1.7, -1.3) {$e$};
    \node[vert, label={[distancia]below:$-\infty$}] (f)
      at (3.4, -1.3) {$f$};
    \node[vert, label={[distancia]right:$-\infty$}] (g) at (5.1, 0)
      {$g$};
    \draw[dirigida] (s) -- node[peso] {3} (a);
    \draw[dirigida] (s) -- node[peso] {5} (c);
    \draw[dirigida] (s) -- node[peso] {2} (e);
    \draw[dirigida] (a) -- node[peso] {$-4$} (b);
    \draw[dirigida] (b) -- node[peso] {4} (g);
    \draw[dirigida] (c) to[bend left=25] node[peso] {6} (d);
    \draw[dirigida] (d) to[bend left=25] node[peso] {$-3$} (c);
    \draw[dirigida] (d) -- node[peso] {8} (g);
    \draw[dirigidaResalte] (e) to[bend left=25] node[peso] {3} (f);
    \draw[dirigidaResalte] (f) to[bend left=25] node[peso] {$-6$}
      (e);
    \draw[dirigida] (f) -- node[peso] {7} (g);
    \node[vert, label={[distancia]above:$\infty$}] (h) at (7.2, 0.8)
      {$h$};
    \node[vert, label={[distancia]above:$\infty$}] (i) at (9.0, 0.8)
      {$i$};
    \node[vert, label={[distancia]below:$\infty$}] (j)
      at (8.1, -0.7) {$j$};
    \draw[dirigidaResalte] (h) -- node[peso] {2} (i);
    \draw[dirigidaResalte] (i) -- node[peso] {3} (j);
    \draw[dirigidaResalte] (j) -- node[peso] {$-8$} (h);
    \node[nota] at (8.1, -1.7) {ciclo negativo no alcanzable};
    \node[nota] at (2.55, -2.3)
      {ciclo negativo alcanzable: $3 - 6 = -3$};
  \end{tikzpicture}
  \end{center}
  \begin{itemize}[itemsep=0.1em]
    \item A $a$ y a $b$ se llega de una sola forma:
          $\delta(s, a) = 3$ y $\delta(s, b) = 3 - 4 = -1$. Una
          arista negativa sola no estorba.
    \item A $c$ se llega por $\langle s, c\rangle$,
          $\langle s, c, d, c\rangle$,
          $\langle s, c, d, c, d, c\rangle$, y así sin fin. El ciclo
          $\langle c, d, c\rangle$ pesa $6 - 3 = 3 > 0$, así que cada
          vuelta empeora: $\delta(s, c) = 5$ y $\delta(s, d) = 11$.
    \item El ciclo $\langle e, f, e\rangle$ pesa $3 - 6 = -3$ y es
          alcanzable: $\delta(s, e) = \delta(s, f) = -\infty$. Como
          $g$ es alcanzable desde $f$, también
          $\delta(s, g) = -\infty$, aunque el camino $s, a, b, g$ pese
          $3$.
    \item $h$, $i$ y $j$ forman un ciclo de peso $2 + 3 - 8 = -3$,
          pero desde $s$ no se llega a ninguno de ellos: los tres
          tienen $\delta = \infty$. Un ciclo negativo solo contamina
          lo que está a la vez al alcance de la fuente y río abajo de
          él.
  \end{itemize}
\end{ejemplo}

\subsubsection*{Un camino más corto no tiene ciclos}

\begin{proposicion}\label{prop:caminos-simples-fuente}
  Si desde $s$ no se alcanza ningún ciclo de peso negativo, para todo
  $v$ alcanzable desde $s$ el valor $\delta(s, v)$ es un número real
  y lo alcanza un camino, que tiene a lo más $n - 1$ aristas.
\end{proposicion}

\begin{proof}
  Sea $q = \langle v_0, v_1, \dots, v_k\rangle$ una caminata de
  $s = v_0$ a $v = v_k$ que repite algún vértice. Entre las parejas de
  posiciones $i < j$ con $v_i = v_j$ se toma una con $j - i$ mínimo.
  El tramo $c = \langle v_i, \dots, v_j\rangle$ no repite vértices
  salvo sus extremos, porque otra repetición dentro de él estaría más
  cerca; es un ciclo, y es alcanzable desde $s$, así que
  $w(c) \ge 0$. Quitarlo deja la caminata
  $q' = \langle v_0, \dots, v_i, v_{j+1}, \dots, v_k\rangle$ de $s$ a
  $v$, con menos aristas y $w(q') = w(q) - w(c) \le w(q)$. Repitiendo
  se llega a un camino $p$ con $w(p) \le w(q)$. Entonces el ínfimo
  sobre caminatas es el mínimo sobre caminos, y como hay una cantidad
  finita de caminos, ese mínimo existe y es un número real. Un camino
  no repite vértices, así que tiene a lo más $n$ vértices y $n - 1$
  aristas.
\end{proof}

La prueba dice qué pasa con cada tipo de ciclo. Uno de peso positivo
nunca está en un camino más corto, porque quitarlo baja el peso. Uno
de peso cero se puede quitar sin cambiar nada. Uno de peso negativo
hace que no haya camino más corto. Por eso basta buscar entre
caminos, y ese $n - 1$ es exactamente el número de pasadas de
Bellman-Ford.

\subsubsection*{Cómo se guardan los caminos}

Igual que BFS, los algoritmos de esta sección no guardan los caminos
completos: cada vértice guarda un \textbf{predecesor} $\pi[v]$, que es
otro vértice o nulo. Al terminar, la cadena
$v, \pi[v], \pi[\pi[v]], \dots$ recorre hacia atrás un camino más
corto de $s$ a $v$, y se imprime como en BFS (sección
\ref{sec:busquedas}). A media ejecución los $\pi$ no tienen por qué
describir caminos más cortos, y con un ciclo negativo ni siquiera
tienen que formar un árbol (ejemplo
\ref{ej:bellman-ford-ciclo-negativo}).

\begin{definicion}[Subgrafo predecesor]
  \label{def:subgrafo-pred-fuente}
  El \textbf{subgrafo predecesor} (\textit{predecessor subgraph}) es
  $G_\pi = (V_\pi, E_\pi)$, con
  \[
    V_\pi = \set{v \in V : \pi[v] \neq \text{nulo}} \cup \set{s} ,
    \qquad
    E_\pi = \set{(\pi[v], v) : v \in V_\pi - \set{s}} .
  \]
\end{definicion}

\begin{definicion}[Árbol de caminos más cortos]
  \label{def:arbol-caminos-cortos}
  Si desde $s$ no se alcanza ningún ciclo negativo, un \textbf{árbol
  de caminos más cortos} (\textit{shortest-paths tree}) con raíz $s$
  es un subgrafo $G' = (V', E')$, con $V' \subseteq V$ y
  $E' \subseteq E$, tal que
  \begin{enumerate}[itemsep=0.1em]
    \item $V'$ es el conjunto de vértices alcanzables desde $s$;
    \item $G'$ es un árbol con raíz $s$;
    \item para todo $v \in V'$, el único camino de $s$ a $v$ en $G'$
          es un camino más corto de $s$ a $v$ en $G$.
  \end{enumerate}
\end{definicion}

Es el árbol BFS con pesos en lugar de número de aristas, y es el que
la clase del 22 de octubre dibujó en azul en el ejemplo de Dijkstra.
No es único: cuando dos caminos empatan, cualquiera de los dos sirve.

\begin{ejemplo}[Dos árboles para las mismas distancias]
  \label{ej:dos-arboles-fuente}
  En este grafo $\delta(s, \cdot) = 0, 1, 2, 4$ para $s, a, b, c$. A
  $b$ se llega con peso $2$ directo o pasando por $a$, y a $c$ con
  peso $4$ por $a$ o por $b$.
  \begin{center}
  \begin{tikzpicture}[distancia/.style={font=\scriptsize\bfseries,
      text=acento, inner sep=1.5pt}]
    \foreach \p/\X/\tit in {g/0/{(a) el grafo},
                            b/4.8/{(b) un árbol},
                            c/9.6/{(c) otro árbol}} {
      \begin{scope}[xshift=\X cm]
        \node[vert, label={[distancia]left:$0$}] (\p-s) at (0, 0)
          {$s$};
        \node[vert, label={[distancia]above:$1$}] (\p-a) at (1.6, 1)
          {$a$};
        \node[vert, label={[distancia]below:$2$}] (\p-b)
          at (1.6, -1) {$b$};
        \node[vert, label={[distancia]right:$4$}] (\p-c) at (3.2, 0)
          {$c$};
        \draw[dirigida] (\p-s) -- node[peso] {1} (\p-a);
        \draw[dirigida] (\p-s) -- node[peso] {2} (\p-b);
        \draw[dirigida] (\p-a) -- node[peso] {1} (\p-b);
        \draw[dirigida] (\p-a) -- node[peso] {3} (\p-c);
        \draw[dirigida] (\p-b) -- node[peso] {2} (\p-c);
        \node[font=\small] at (1.6, -2.15) {\tit};
      \end{scope}
    }
    \foreach \u/\v/\w in {s/a/1, s/b/2, a/c/3}
      \draw[dirigidaArbol] (b-\u) -- node[peso] {\w} (b-\v);
    \foreach \u/\v/\w in {s/a/1, a/b/1, b/c/2}
      \draw[dirigidaArbol] (c-\u) -- node[peso] {\w} (c-\v);
  \end{tikzpicture}
  \end{center}
  Los dos son árboles de caminos más cortos válidos (Cormen muestra
  otro caso en su figura 24.2). Un mismo algoritmo puede dar uno u
  otro según el orden en que revise las aristas: Bellman-Ford, que
  viene más adelante, da el de (b) con el orden
  $(s, a), (s, b), (a, b), (a, c), (b, c)$ y el de (c) con
  $(s, a), (a, b), (b, c), (s, b), (a, c)$.
\end{ejemplo}

% =====================================================================

\subsection*{Relajación}

Todos los algoritmos de la sección comparten el mismo esqueleto. Cada
vértice guarda una \textbf{estimación} $d[v]$ del peso de camino más
corto, que siempre es una cota superior de $\delta(s, v)$, y un
predecesor $\pi[v]$. Se arranca con $d[s] = 0$ e $\infty$ en los
demás (\textsc{Initialize-Single-Source} en Cormen):

\begin{algorithm}[H]
\caption{\textsc{Inicializar-Fuente}$(G, s)$}
\label{alg:inicializar-fuente}
\begin{algorithmic}[1]
\Procedure{Inicializar-Fuente}{$G, s$}
  \ForAll{$v \in V$}
    \State $d[v] \gets \infty$, \quad $\pi[v] \gets$ nulo
  \EndFor
  \State $d[s] \gets 0$
\EndProcedure
\end{algorithmic}
\end{algorithm}

Después, la única operación que cambia $d$ y $\pi$ es
\textbf{relajar} una arista $(u, v)$: si llegar a $v$ pasando por $u$
es más barato que lo que se tiene, se actualizan $d[v]$ y $\pi[v]$
(\textsc{Relax} en Cormen).

\begin{algorithm}[H]
\caption{\textsc{Relajar}$(u, v, w)$}
\label{alg:relajar}
\begin{algorithmic}[1]
\Procedure{Relajar}{$u, v, w$}
  \If{$d[v] > d[u] + w(u, v)$}
    \State $d[v] \gets d[u] + w(u, v)$, \quad $\pi[v] \gets u$
  \EndIf
\EndProcedure
\end{algorithmic}
\end{algorithm}

\textsc{Inicializar-Fuente} cuesta $\Tht{n}$ y \textsc{Relajar}
$\Ord{1}$. Es la misma relajación que la clase del 22 de octubre
escribió en línea dentro de Dijkstra. El nombre parece al revés,
porque la operación aprieta una cota, y la razón es histórica: las
distancias finales cumplen $\delta(s, v) \le \delta(s, u) + w(u, v)$,
relajar es forzar la restricción $d[v] \le d[u] + w(u, v)$, y una
restricción que ya se cumple está ``relajada''.

Los algoritmos difieren solo en cuántas veces relajan cada arista y en
qué orden. Dijkstra y el algoritmo para grafos acíclicos relajan cada
arista exactamente una vez; Bellman-Ford, $n - 1$ veces. Que eso
baste lo explican seis propiedades.

\subsubsection*{Las seis propiedades de la relajación}

La primera es un hecho sobre $\delta$. Las otras cinco hablan de
cualquier sucesión de relajaciones, y todas suponen lo mismo: el
grafo se inicializó con \textsc{Inicializar-Fuente}$(G, s)$ y desde
entonces $d$ y $\pi$ solo cambian por llamadas a \textsc{Relajar}.
Aquí van los enunciados con su intuición; las pruebas completas están
en la sección \ref{sec:caminos-pruebas}.

\begin{proposicion}[Desigualdad del triángulo; Cormen, lema 24.10]
  \label{prop:triangulo}
  Para toda arista $(u, v) \in E$,
  $\delta(s, v) \le \delta(s, u) + w(u, v)$.
\end{proposicion}

\textbf{Intuición.} Ir a $u$ de la mejor manera y luego tomar
$(u, v)$ es una forma de llegar a $v$, y la mejor forma no cuesta más
que una forma cualquiera. Vale con infinitos: si $u$ no es
alcanzable, el lado derecho es $\infty$; si $\delta(s, u) = -\infty$,
el ciclo negativo que llega a $u$ también llega a $v$ y
$\delta(s, v) = -\infty$.

\begin{proposicion}[Cota superior; Cormen, lema 24.11]
  \label{prop:cota-superior}
  En todo momento $d[v] \ge \delta(s, v)$ para todo $v \in V$, y en
  cuanto $d[v]$ alcanza $\delta(s, v)$ ya no cambia.
\end{proposicion}

\textbf{Intuición.} Cada valor finito de $d[v]$ es el peso de una
caminata de verdad de $s$ a $v$, la que fueron armando las
relajaciones, y ninguna caminata pesa menos que el ínfimo. Como
relajar solo baja valores, cuando $d[v]$ toca el piso se queda ahí.

\begin{proposicion}[Sin camino; Cormen, corolario 24.12]
  \label{prop:sin-camino}
  Si no hay camino de $s$ a $v$, en todo momento
  $d[v] = \delta(s, v) = \infty$.
\end{proposicion}

\textbf{Intuición.} Es la cota superior con $\delta(s, v) = \infty$:
ninguna relajación baja un vértice al que no llega nada.

\begin{proposicion}[Convergencia; Cormen, lema 24.14]
  \label{prop:convergencia}
  Si $s \leadsto u \to v$ es un camino más corto y
  $d[u] = \delta(s, u)$ en algún momento antes de relajar $(u, v)$,
  entonces $d[v] = \delta(s, v)$ en todo momento después de esa
  relajación.
\end{proposicion}

\textbf{Intuición.} Al relajar queda
$d[v] \le d[u] + w(u, v) = \delta(s, u) + w(u, v) = \delta(s, v)$; la
última igualdad es el lema \ref{lem:subcaminos-cortos} aplicado al
camino más corto. La cota superior impide bajar más. La distancia
correcta de $u$ se hereda a $v$ por la última arista del camino.

\begin{proposicion}[Relajación de caminos; Cormen, lema 24.15]
  \label{prop:relajacion-caminos}
  Sea $p = \langle v_0, v_1, \dots, v_k\rangle$ un camino más corto de
  $s = v_0$ a $v_k$. Si sus aristas se relajan en el orden
  $(v_0, v_1), (v_1, v_2), \dots, (v_{k-1}, v_k)$, entonces
  $d[v_k] = \delta(s, v_k)$ desde la última de esas relajaciones en
  adelante. Vale sin importar qué otras relajaciones se intercalen
  entre ellas.
\end{proposicion}

\textbf{Intuición.} Inducción con la convergencia. Al principio
$d[v_0] = 0 = \delta(s, s)$. Al relajar $(v_{i-1}, v_i)$, $v_{i-1}$
ya es correcto, así que $v_i$ queda correcto, y lo que se haga en
medio no lo estropea, porque un valor correcto ya no se mueve. Esta
es la propiedad que hace funcionar a todos los algoritmos de la
sección: no hace falta saber cuál es el camino más corto, basta
garantizar que sus aristas se relajen en orden, aunque sea revueltas
entre muchas otras.

\begin{proposicion}[Subgrafo predecesor; Cormen, lema 24.17]
  \label{prop:subgrafo-pred}
  Si desde $s$ no se alcanza ningún ciclo negativo y en algún momento
  $d[v] = \delta(s, v)$ para todo $v \in V$, entonces $G_\pi$ es un
  árbol de caminos más cortos con raíz $s$.
\end{proposicion}

\textbf{Intuición.} Sin ciclos negativos, $G_\pi$ es siempre un árbol
con raíz $s$: una relajación no puede cerrar un ciclo de
predecesores. Al final, además, cada arista del árbol queda
\textbf{tensa}, $d[v] = d[\pi[v]] + w(\pi[v], v)$. Cuando se fijó
$\pi[v] = u$ quedó $d[v] = d[u] + w(u, v)$; desde entonces $d[v]$ no
cambió, porque si cambiara cambiaría $\pi[v]$, y $d[u]$ solo pudo
bajar, así que $d[v] \ge d[u] + w(u, v)$. La desigualdad del
triángulo da la otra mitad. Sumando las igualdades a lo largo del
camino del árbol, los $d$ se cancelan en telescopio y el camino pesa
$d[v] - d[s] = \delta(s, v)$.

\begin{idea}
  Un algoritmo de caminos más cortos desde una fuente es una política
  para escoger qué arista relajar. Si la política garantiza que las
  aristas de cada camino más corto se relajen en orden, la relajación
  de caminos da las distancias y la propiedad del subgrafo predecesor
  da el árbol. Bellman-Ford lo garantiza a fuerza de repetir, el
  algoritmo para grafos acíclicos con el orden topológico, y Dijkstra
  con la cola de prioridad.
\end{idea}

% =====================================================================

\subsection*{Bellman-Ford}

Bellman-Ford resuelve el caso general: pesos reales cualesquiera, con
la posibilidad de que haya ciclos negativos. Devuelve un booleano. Si
desde $s$ no se alcanza ningún ciclo negativo devuelve verdadero, y
entonces $d[v] = \delta(s, v)$ para todo $v$ y $G_\pi$ es un árbol de
caminos más cortos. Si se alcanza alguno devuelve falso, porque el
problema no tiene solución. La idea cabe en una línea: relajar todas
las aristas, $n - 1$ veces, y luego revisar.

\begin{algorithm}[H]
\caption{\textsc{Bellman-Ford}$(G, w, s)$}
\label{alg:bellman-ford}
\begin{algorithmic}[1]
\Procedure{Bellman-Ford}{$G, w, s$}
  \State \Call{Inicializar-Fuente}{$G, s$}
  \For{$i \gets 1$ \Hasta $n - 1$} \Comment{$n - 1$ pasadas}
    \ForAll{$(u, v) \in E$}
      \State \Call{Relajar}{$u, v, w$}
    \EndFor
  \EndFor
  \ForAll{$(u, v) \in E$} \Comment{revisión}
    \If{$d[v] > d[u] + w(u, v)$}
      \State \Return falso \Comment{hay un ciclo negativo alcanzable}
    \EndIf
  \EndFor
  \State \Return verdadero
\EndProcedure
\end{algorithmic}
\end{algorithm}

\begin{itemize}[itemsep=0.1em]
  \item \textbf{Línea 3, el $n - 1$.} Por la proposición
        \ref{prop:caminos-simples-fuente}, un camino más corto tiene a
        lo más $n - 1$ aristas, y cada pasada avanza al menos una
        arista de cada camino más corto. Es una cota de peor caso:
        cada relajación usa los valores que ya bajaron en la misma
        pasada, así que con un orden favorable una mejora recorre
        varias aristas de una vez.
  \item \textbf{Línea 4.} El orden de las aristas es arbitrario. No
        afecta el resultado, solo cuántas pasadas hacen algo de
        verdad.
  \item \textbf{Líneas 8 a 12.} Una pasada más que solo pregunta. Sin
        ciclos negativos todo $d$ ya es exacto y ninguna arista se
        puede relajar; si alguna todavía se puede, hay un ciclo
        negativo alcanzable desde $s$.
\end{itemize}

\begin{lema}[Cormen, lema 24.2]\label{lem:bellman-ford-pasadas}
  Si desde $s$ no se alcanza ningún ciclo negativo, al terminar las
  $n - 1$ pasadas de las líneas 3 a 7 se tiene $d[v] = \delta(s, v)$
  para todo $v$ alcanzable desde $s$.
\end{lema}

\begin{proof}
  Sea $v$ alcanzable y sea $p = \langle v_0, v_1, \dots, v_k\rangle$
  un camino más corto de $s = v_0$ a $v = v_k$; existe y tiene
  $k \le n - 1$ aristas por la proposición
  \ref{prop:caminos-simples-fuente}. Cada pasada relaja todas las
  aristas, así que la pasada $i$ relaja, entre otras, a
  $(v_{i-1}, v_i)$. Las aristas de $p$ quedan relajadas en orden, una
  en cada una de las pasadas $1, \dots, k$, y la relajación de caminos
  (proposición \ref{prop:relajacion-caminos}) da
  $d[v] = \delta(s, v)$.
\end{proof}

\begin{teorema}[Corrección de Bellman-Ford; Cormen, teorema 24.4]
  \label{teo:bellman-ford-correccion}
  Si desde $s$ no se alcanza ningún ciclo negativo,
  \textsc{Bellman-Ford}$(G, w, s)$ devuelve verdadero,
  $d[v] = \delta(s, v)$ para todo $v \in V$ y $G_\pi$ es un árbol de
  caminos más cortos con raíz $s$. Si desde $s$ se alcanza un ciclo
  negativo, devuelve falso.
\end{teorema}

\begin{proof}
  \textit{Sin ciclos negativos alcanzables.} Para los vértices
  alcanzables, $d[v] = \delta(s, v)$ es el lema
  \ref{lem:bellman-ford-pasadas}; para los demás, la propiedad sin
  camino. Con todas las distancias correctas, la propiedad del
  subgrafo predecesor da el árbol. Falta ver que la revisión no
  falla: para toda arista $(u, v)$, por la desigualdad del triángulo,
  \[
    d[v] = \delta(s, v) \le \delta(s, u) + w(u, v)
         = d[u] + w(u, v) ,
  \]
  así que la condición de la línea 9 nunca se cumple y el algoritmo
  devuelve verdadero.

  \textit{Con un ciclo negativo alcanzable.} Sea
  $c = \langle v_0, v_1, \dots, v_k\rangle$, con $v_0 = v_k$, un ciclo
  alcanzable desde $s$ con
  \[
    \sum_{i=1}^{k} w(v_{i-1}, v_i) < 0 .
  \]
  Supongamos que el algoritmo devuelve verdadero. Entonces ninguna
  arista del ciclo cumplió la condición de la línea 9:
  $d[v_i] \le d[v_{i-1}] + w(v_{i-1}, v_i)$ para $i = 1, \dots, k$.
  Sumando las $k$ desigualdades,
  \[
    \sum_{i=1}^{k} d[v_i] \le \sum_{i=1}^{k} d[v_{i-1}]
      + \sum_{i=1}^{k} w(v_{i-1}, v_i) .
  \]
  Como $v_0 = v_k$, cada vértice del ciclo aparece exactamente una vez
  en cada una de las dos sumas de $d$, así que son iguales. Además
  son finitas. Cada $v_i$ es alcanzable por un camino de a lo más
  $n - 1$ aristas, la pasada $j$ relaja la $j$-ésima arista de ese
  camino y $d$ solo baja, así que tras las $n - 1$ pasadas
  $d[v_i] < \infty$; y ningún $d$ vale $-\infty$, porque solo toma
  valores que son pesos de caminatas. Restando las dos sumas iguales
  queda $0 \le \sum_{i=1}^{k} w(v_{i-1}, v_i)$, que contradice que el
  ciclo es negativo.
\end{proof}

La segunda mitad es el corazón de la detección. Un ciclo negativo no
puede tener satisfechas las desigualdades de todas sus aristas: al
sumarlas alrededor del ciclo los $d$ se cancelan y queda
$0 \le w(c)$. Alguna tiene que fallar, y esa es la que la revisión
encuentra. El ejemplo \ref{ej:bellman-ford-ciclo-negativo} lo muestra
con números.

\textbf{Costo.} La línea 2 cuesta $\Tht{n}$. Las líneas 3 a 7 hacen
$n - 1$ pasadas, y cada pasada recorre las $m$ aristas con una
relajación de $\Ord{1}$: $(n - 1) \cdot \Tht{m} = \Tht{nm}$. La
revisión recorre las aristas otra vez, $\Ord{m}$. En total,
$\Tht{nm}$, con $\Tht{n}$ de memoria extra para $d$ y $\pi$. Dos
precisiones sobre la estructura de datos:
\begin{itemize}[itemsep=0.1em]
  \item ``Para toda arista'' con listas de adyacencia es ``para cada
        $u$, para cada $v$ en su lista'', que cuesta $\Tht{n + m}$
        por pasada y no $\Tht{m}$. El total es $\Tht{n(n + m)}$, que
        coincide con $\Tht{nm}$ cuando $m = \Omg{n}$, por ejemplo
        cuando todo vértice es alcanzable desde $s$ y por tanto
        $m \ge n - 1$. Con una lista plana de aristas cada pasada sí
        cuesta $\Tht{m}$.
  \item Con matriz de adyacencia recorrer las aristas cuesta
        $\Tht{n^2}$ por pasada y el total es $\Tht{n^3}$. En un grafo
        denso, $m = \Tht{n^2}$, Bellman-Ford cuesta $\Tht{n^3}$ aun
        con listas: lo mismo que Floyd para \emph{todos} los pares.
\end{itemize}

\begin{ejemplo}[Bellman-Ford sobre la figura 24.4 de Cormen]
  \label{ej:bellman-ford-24-4}
  El grafo tiene $V = \set{s, t, x, y, z}$, fuente $s$ y diez aristas:
  \[
    \begin{array}{lllll}
      (s, t) : 6 & (s, y) : 7 & (t, x) : 5 & (t, y) : 8
        & (t, z) : -4 \\
      (x, t) : -2 & (y, x) : -3 & (y, z) : 9 & (z, s) : 2
        & (z, x) : 7
    \end{array}
  \]
  Cada pasada relaja las aristas en el orden de Cormen,
  \[
    (t, x), (t, y), (t, z), (x, t), (y, x), (y, z), (z, x), (z, s),
    (s, t), (s, y) .
  \]
  Las aristas de $s$ van al final a propósito: es un orden malo, que
  obliga a usar las $n - 1 = 4$ pasadas. Cada celda es $d[v]$ con
  $\pi[v]$ entre paréntesis, al terminar la pasada.
  \begin{center}
  \small
  \begin{tabular}{cccccl}
    \toprule
    \textbf{Pasada} & $t$ & $x$ & $y$ & $z$
      & \textbf{Relajaciones que cambian algo} \\
    \midrule
    inicial & $\infty$ & $\infty$ & $\infty$ & $\infty$ & \\
    1 & $6\ (s)$ & $\infty$ & $7\ (s)$ & $\infty$
      & $(s, t)$: $\infty \to 6$; \ $(s, y)$: $\infty \to 7$ \\
    2 & $6\ (s)$ & $4\ (y)$ & $7\ (s)$ & $2\ (t)$
      & $(t, x)$: $\infty \to 11$; \ $(t, z)$: $\infty \to 2$;
        \ $(y, x)$: $11 \to 4$ \\
    3 & $2\ (x)$ & $4\ (y)$ & $7\ (s)$ & $2\ (t)$
      & $(x, t)$: $6 \to 2$ \\
    4 & $2\ (x)$ & $4\ (y)$ & $7\ (s)$ & $-2\ (t)$
      & $(t, z)$: $2 \to -2$ \\
    \bottomrule
  \end{tabular}
  \end{center}
  La misma ejecución en dibujos. En verde, los vértices cuya $d$ bajó
  en esa pasada; en azul, las aristas $(\pi[v], v)$.
  \begin{center}
  \begin{tikzpicture}[scale=0.8, transform shape,
      distancia/.style={font=\scriptsize\bfseries, text=acento,
        inner sep=1.5pt}]
    \foreach \p/\X/\Y/\dt/\dx/\dy/\dz/\tit in {
        a/0/0/\infty/\infty/\infty/\infty/{(a) inicial},
        b/6.2/0/6/\infty/7/\infty/{(b) tras la pasada 1},
        c/12.4/0/6/4/7/2/{(c) tras la pasada 2},
        d/3.1/-4.6/2/4/7/2/{(d) tras la pasada 3},
        e/9.3/-4.6/2/4/7/-2/{(e) tras la pasada 4}} {
      \begin{scope}[xshift=\X cm, yshift=\Y cm]
        \node[vert, label={[distancia]left:$0$}] (\p-s) at (0, 0)
          {$s$};
        \node[vert, label={[distancia]above:$\dt$}] (\p-t)
          at (1.8, 1.2) {$t$};
        \node[vert, label={[distancia]above:$\dx$}] (\p-x)
          at (4.2, 1.2) {$x$};
        \node[vert, label={[distancia]below:$\dy$}] (\p-y)
          at (1.8, -1.2) {$y$};
        \node[vert, label={[distancia]below:$\dz$}] (\p-z)
          at (4.2, -1.2) {$z$};
        \draw[dirigida] (\p-s) -- node[peso] {6} (\p-t);
        \draw[dirigida] (\p-s) -- node[peso] {7} (\p-y);
        \draw[dirigida] (\p-t) to[bend left=22] node[peso] {5}
          (\p-x);
        \draw[dirigida] (\p-x) to[bend left=22] node[peso] {$-2$}
          (\p-t);
        \draw[dirigida] (\p-t) -- node[peso] {8} (\p-y);
        \draw[dirigida] (\p-t) -- node[peso, pos=0.78] {$-4$}
          (\p-z);
        \draw[dirigida] (\p-y) -- node[peso, pos=0.78] {$-3$}
          (\p-x);
        \draw[dirigida] (\p-y) -- node[peso, pos=0.35] {9} (\p-z);
        \draw[dirigida] (\p-z) -- node[peso, pos=0.2] {2} (\p-s);
        \draw[dirigida] (\p-z) -- node[peso] {7} (\p-x);
        \node[font=\small] at (2.1, -2.4) {\tit};
      \end{scope}
    }
    % aristas de pi en cada instantanea
    \foreach \p in {b, c} {
      \draw[dirigidaArbol] (\p-s) -- node[peso] {6} (\p-t);
    }
    \foreach \p in {b, c, d, e} {
      \draw[dirigidaArbol] (\p-s) -- node[peso] {7} (\p-y);
    }
    \foreach \p in {c, d, e} {
      \draw[dirigidaArbol] (\p-y) -- node[peso, pos=0.78] {$-3$}
        (\p-x);
      \draw[dirigidaArbol] (\p-t) -- node[peso, pos=0.78] {$-4$}
        (\p-z);
    }
    \foreach \p in {d, e} {
      \draw[dirigidaArbol] (\p-x) to[bend left=22] node[peso] {$-2$}
        (\p-t);
    }
    % vertices cuya d bajo en esa pasada
    \foreach \p/\v in {b/t, b/y, c/x, c/z, d/t, e/z} {
      \node[vertOk] at (\p-\v) {$\v$};
    }
  \end{tikzpicture}
  \end{center}
  \begin{itemize}[itemsep=0.1em]
    \item \textbf{Pasada 1.} Solo cambian $t$ e $y$. Las ocho
          primeras aristas salen de vértices con $d = \infty$, y la
          regla $\infty + w = \infty$ hace que no muevan nada.
    \item \textbf{Pasada 2.} $(t, x)$ baja $x$ a $11$ y, más adelante
          en la misma pasada, $(y, x)$ lo baja a $4$. $(t, z)$ pone
          $z$ en $2$, que todavía no es su valor final.
    \item \textbf{Pasada 3.} $(x, t)$ da $4 - 2 = 2 < 6$: $t$ tenía
          $6$ por la arista directa y mejora por $s, y, x, t$, un
          camino de tres aristas.
    \item \textbf{Pasada 4.} $(t, z)$ da $2 - 4 = -2 < 2$, y $z$ se
          corrige por segunda vez.
    \item \textbf{Revisión.} Ninguna arista se puede relajar y el
          algoritmo devuelve verdadero. Las aristas del árbol quedan
          tensas, como $d[x] = 4 = d[y] + w(y, x)$; las demás, con
          holgura, como $(z, x)$: $4 \le -2 + 7$.
  \end{itemize}
  El árbol final es un solo camino, $s \to y \to x \to t \to z$, con
  $n - 1 = 4$ aristas: el peor caso de la cota, y con el orden de
  Cormen cada pasada avanza exactamente una arista de él. Con las
  aristas de $s$ al principio de la lista, la pasada 1 ya deja
  $t = 6$, $x = 4$, $y = 7$ y $z = 2$, porque $(s, y)$ y luego
  $(y, x)$ caen en la misma pasada; converge en tres pasadas y la
  cuarta no cambia nada. El orden cambia cuántas pasadas sirven, no
  el resultado: $\delta(s, \cdot) = 0, 2, 4, 7, -2$ para
  $s, t, x, y, z$.

  Un detalle más: el ciclo $s, y, x, t, z, s$ pesa
  $7 - 3 - 2 - 4 + 2 = 0$. Es un ciclo de peso cero, de los que no
  estorban, y por eso la arista $(z, s)$ también queda tensa,
  $d[s] = 0 = d[z] + 2$, sin estar en el árbol.
\end{ejemplo}

\begin{ejemplo}[El mismo grafo con un ciclo negativo]
  \label{ej:bellman-ford-ciclo-negativo}
  Se cambia un solo peso, $w(z, x)$, de $7$ a $4$ (Cormen, ejercicio
  24.1-1). Aparece el ciclo $z \to x \to t \to z$ de peso
  $4 - 2 - 4 = -2$, alcanzable desde $s$ por $(s, t)$. Como además
  hay arista $(z, s)$, desde el ciclo se llega a todos los vértices,
  incluida la fuente: $\delta(s, v) = -\infty$ para todo $v$, también
  para $v = s$.

  Las pasadas 1 a 3 son idénticas a las del ejemplo anterior: mientras
  $d[z] \ge 0$, la arista $(z, x)$ no mejora a $x$, valga $7$ o $4$.
  La diferencia sale en la pasada 4.
  \begin{center}
  \small
  \begin{tabular}{ccccccl}
    \toprule
    \textbf{Pasada} & $s$ & $t$ & $x$ & $y$ & $z$
      & \textbf{Qué pasa} \\
    \midrule
    3 & $0$ & $2\ (x)$ & $4\ (y)$ & $7\ (s)$ & $2\ (t)$
      & idéntico al grafo original \\
    4 & $0$ & $2\ (x)$ & $2\ (z)$ & $7\ (s)$ & $-2\ (t)$
      & $(t, z)$: $2 \to -2$; \ $(z, x)$: $4 \to 2$ \\
    \midrule
    \multicolumn{7}{l}{Revisión: $(x, t)$ da
      $d[t] = 2 > d[x] + w(x, t) = 2 - 2 = 0$, así que devuelve
      falso.} \\
    \midrule
    5 & $0$ & $0$ & $2$ & $7$ & $-2$ & si se siguiera relajando \\
    6 & $-2$ & $0$ & $0$ & $5$ & $-4$ & hasta $s$ baja \\
    8 & $-4$ & $-2$ & $-2$ & $3$ & $-6$
      & cada dos pasadas, $2$ menos \\
    \bottomrule
  \end{tabular}
  \end{center}
  En la pasada 4, $(t, z)$ baja $z$ a $-2$ y unas aristas después
  $(z, x)$ da $-2 + 4 = 2 < 4$; con el peso original daba
  $-2 + 7 = 5$ y no cambiaba nada. En la revisión, las tres
  desigualdades del ciclo quedan así:
  \begin{itemize}[itemsep=0.1em]
    \item $(z, x)$: $d[x] = 2 \le d[z] + 4 = 2$, se cumple;
    \item $(x, t)$: $d[t] = 2 \le d[x] - 2 = 0$, \textbf{falla};
    \item $(t, z)$: $d[z] = -2 \le d[t] - 4 = -2$, se cumple.
  \end{itemize}
  Es la prueba del teorema \ref{teo:bellman-ford-correccion} con
  números: la suma de los $d$ del ciclo es $2 + 2 - 2 = 2$, y si las
  tres se cumplieran quedaría $2 \le 2 + (-2)$. Alguna tiene que
  fallar, y falla $(x, t)$.
  \begin{center}
  \begin{tikzpicture}[scale=0.9, transform shape,
      distancia/.style={font=\scriptsize\bfseries, text=acento,
        inner sep=1.5pt}]
    \foreach \p/\X/\tit in {g/0/{(a) el ciclo $z, x, t, z$ pesa
                              $-2$},
                            h/7.4/{(b) tras la pasada 4}} {
      \begin{scope}[xshift=\X cm]
        \node[vert] (\p-s) at (0, 0) {$s$};
        \node[vert] (\p-t) at (1.8, 1.2) {$t$};
        \node[vert] (\p-x) at (4.2, 1.2) {$x$};
        \node[vert] (\p-y) at (1.8, -1.2) {$y$};
        \node[vert] (\p-z) at (4.2, -1.2) {$z$};
        \draw[dirigida] (\p-s) -- node[peso] {6} (\p-t);
        \draw[dirigida] (\p-s) -- node[peso] {7} (\p-y);
        \draw[dirigida] (\p-t) to[bend left=22] node[peso] {5}
          (\p-x);
        \draw[dirigida] (\p-x) to[bend left=22] node[peso] {$-2$}
          (\p-t);
        \draw[dirigida] (\p-t) -- node[peso] {8} (\p-y);
        \draw[dirigida] (\p-t) -- node[peso, pos=0.78] {$-4$}
          (\p-z);
        \draw[dirigida] (\p-y) -- node[peso, pos=0.78] {$-3$}
          (\p-x);
        \draw[dirigida] (\p-y) -- node[peso, pos=0.35] {9} (\p-z);
        \draw[dirigida] (\p-z) -- node[peso, pos=0.2] {2} (\p-s);
        \draw[dirigida] (\p-z) -- node[peso] {$\mathbf{4}$} (\p-x);
        \node[font=\small] at (2.1, -2.4) {\tit};
      \end{scope}
    }
    % (a) el ciclo negativo
    \draw[dirigidaResalte] (g-x) to[bend left=22] node[peso] {$-2$}
      (g-t);
    \draw[dirigidaResalte] (g-t) -- node[peso, pos=0.78] {$-4$}
      (g-z);
    \draw[dirigidaResalte] (g-z) -- node[peso] {$\mathbf{4}$} (g-x);
    % (b) estado tras la pasada 4: d y pi
    \foreach \v/\val/\lado in {s/0/left, t/2/above, x/2/above,
                               y/7/below, z/-2/below}
      \node[distancia, \lado=0.3cm] at (h-\v) {$\val$};
    \draw[dirigidaArbol] (h-s) -- node[peso] {7} (h-y);
    \draw[dirigidaArbol] (h-x) to[bend left=22] node[peso] {$-2$}
      (h-t);
    \draw[dirigidaArbol] (h-t) -- node[peso, pos=0.78] {$-4$} (h-z);
    \draw[dirigidaArbol] (h-z) -- node[peso] {$\mathbf{4}$} (h-x);
    \node[nota, align=left, anchor=west] at (12.0, 0.2)
      {revisión de $(x, t)$:\\ $d[t] = 2 > d[x] + w(x, t) = 0$,\\
       devuelve falso};
  \end{tikzpicture}
  \end{center}
  Dos cosas se ven en (b). La primera: $G_\pi$ ya no es un árbol.
  $\pi[t] = x$, $\pi[z] = t$ y $\pi[x] = z$ forman el ciclo
  $x \to t \to z \to x$, exactamente el ciclo negativo, y $t$ ya no
  cuelga de $s$. La segunda está en las últimas filas de la tabla: si
  se siguiera relajando, los valores bajarían sin fin, $2$ cada dos
  pasadas, el peso del ciclo, y hasta $d[s]$ deja de valer $0$. Por
  eso el algoritmo se detiene en $n - 1$ pasadas y pregunta, en lugar
  de esperar a que las $d$ se estabilicen.
\end{ejemplo}

\begin{ojo}
  \textbf{El infinito no es un entero grande.} En código es tentador
  representar $\infty$ con un centinela como $10^9$. Con pesos
  negativos, $10^9 + (-8) < 10^9$, y un vértice no alcanzable
  ``mejora'' a otro no alcanzable. En un grafo con $s$ aislado y el
  ciclo $h, i, j$ de la figura 24.1, Bellman-Ford sin protección
  devuelve falso, aunque ningún ciclo negativo sea alcanzable. La
  protección es relajar solo si $d[u] \neq \infty$, que es lo que la
  regla $\infty + a = \infty$ dice en papel. Si el centinela es el
  entero más grande del lenguaje, además, la suma se desborda.
\end{ojo}

\begin{ojo}
  \textbf{Qué detecta y qué no.} Bellman-Ford solo ve ciclos
  negativos alcanzables desde $s$. Uno al que no se llega no afecta
  nada y sus vértices se quedan en $\infty$, que es correcto. Para
  saber si el grafo tiene \emph{algún} ciclo negativo se agrega un
  vértice nuevo con aristas de peso $0$ hacia todos y se corre desde
  él; equivale a arrancar con todas las $d$ en $0$, y es el mismo
  truco de las restricciones de diferencia, al final de la sección.

  \smallskip
  \textbf{Cuando devuelve falso, las $d$ no significan nada.} Son
  números a media caída, $G_\pi$ puede tener ciclos y ni siquiera
  $d[s] = 0$ es confiable, como en el ejemplo anterior.

  \smallskip
  \textbf{Grafos no dirigidos.} Una arista $\set{u, v}$ es un par de
  arcos, y si su peso es negativo forman el ciclo $u, v, u$ de peso
  $2w(u, v) < 0$. Bellman-Ford devuelve falso en cuanto hay una
  arista negativa alcanzable. Con pesos no negativos funciona, pero
  ahí Dijkstra es más barato.
\end{ojo}

\begin{observacion}
  Bellman-Ford es programación dinámica sobre el número de aristas.
  Si $d^{(i)}[v]$ es el menor peso de una caminata de $s$ a $v$ con a
  lo más $i$ aristas,
  \[
    d^{(i)}[v] = \min\Bigl(d^{(i-1)}[v],\
      \min_{(u, v) \in E} \bigl(d^{(i-1)}[u] + w(u, v)\bigr)\Bigr) ,
  \]
  y la pasada $i$ calcula esa capa, solo que en el lugar: como usa
  valores ya actualizados, tras la pasada $i$ queda
  $d[v] \le d^{(i)}[v]$, que es todo lo que la prueba necesita. Para
  todos los pares, la misma recurrencia da el método por
  multiplicación de matrices (sección \ref{sec:todos-pares}). Cormen
  no lo presenta así en este capítulo; solo clasifica a Dijkstra como
  voraz y a Floyd-Warshall como programación dinámica.
\end{observacion}

\begin{ejercicio}
  Corre Bellman-Ford sobre el grafo de la figura 24.4 con fuente $z$ y
  el mismo orden de aristas (Cormen, ejercicio 24.1-1).

  \textit{Solución.} Para $s, t, x, y, z$, tras la pasada 1:
  $2\,(z)$, $8\,(s)$, $7\,(z)$, $9\,(s)$, $0$. En la pasada 2 bajan
  $t = 5\,(x)$ y $x = 6\,(y)$; en la pasada 3, $t = 4\,(x)$; la pasada
  4 no cambia nada y el algoritmo devuelve verdadero. Queda
  $\delta(z, \cdot) = 2, 4, 6, 9, 0$.
\end{ejercicio}

\begin{ejercicio}
  Sin ciclos negativos, si todo vértice tiene un camino más corto de a
  lo más $k$ aristas, haz que Bellman-Ford termine en $k + 1$ pasadas
  sin conocer $k$ (Cormen, ejercicio 24.1-3, donde $k$ se llama $m$).

  \textit{Solución.} Una bandera que diga si la pasada cambió algo.
  Si una pasada no cambia nada, la siguiente recibiría exactamente los
  mismos valores y tampoco cambiaría nada, así que se termina. Tras
  $k$ pasadas todo es correcto por la relajación de caminos, y la
  pasada $k + 1$ lo confirma. El peor caso sigue siendo $\Tht{nm}$.
\end{ejercicio}

\begin{ejercicio}
  Modifica Bellman-Ford para que ponga $d[v] = -\infty$ en todo $v$
  con un ciclo negativo en alguna caminata de $s$ a $v$, y para que,
  si devuelve falso, liste los vértices de un ciclo negativo (Cormen,
  ejercicios 24.1-4 y 24.1-6).

  \textit{Solución.} Tras las $n - 1$ pasadas, si una arista $(u, v)$
  todavía se puede relajar, $\delta(s, v) = -\infty$: un vértice con
  $\delta$ finito ya tiene su valor y no baja más. Y cada ciclo
  negativo alcanzable tiene al menos una arista así, por la prueba
  del teorema. Un BFS o DFS desde esos vértices $v$ marca con
  $-\infty$ todo lo alcanzado, y el costo sigue en $\Ord{nm}$. Para
  listar un ciclo se relaja una de esas aristas $(u, v)$ y desde $v$
  se sigue $\pi$ hacia atrás $n$ veces, lo que cae dentro de un ciclo
  de $G_\pi$; desde ahí se sigue $\pi$ hasta volver al mismo vértice.
  En el ejemplo \ref{ej:bellman-ford-ciclo-negativo}, desde $t$ sale
  el ciclo $x, t, z$.
\end{ejercicio}

% =====================================================================

\subsection*{Caminos más cortos en un grafo acíclico}

En un grafo dirigido acíclico no hace falta repetir pasadas. Si los
vértices se procesan en orden topológico, las aristas de cualquier
camino aparecen en orden, y una sola pasada cumple la hipótesis de la
relajación de caminos. Los pesos negativos están permitidos: sin
ciclos no hay ciclos negativos, y $\delta(s, v)$ siempre está bien
definido. El orden topológico se obtiene con DFS, listando los
vértices por tiempo de salida $\text{sale}[u]$ decreciente (sección
\ref{sec:aplicaciones-dfs}). En Cormen el procedimiento se llama
\textsc{DAG-Shortest-Paths}.

\begin{algorithm}[H]
\caption{\textsc{Caminos-DAG}$(G, w, s)$}
\label{alg:caminos-dag}
\begin{algorithmic}[1]
\Procedure{Caminos-DAG}{$G, w, s$}
  \State ordenar topológicamente los vértices de $G$
  \State \Call{Inicializar-Fuente}{$G, s$}
  \ForAll{$u \in V$, en orden topológico}
    \ForAll{$v$ en la lista de adyacencia de $u$}
      \State \Call{Relajar}{$u, v, w$}
    \EndFor
  \EndFor
\EndProcedure
\end{algorithmic}
\end{algorithm}

Cuando le toca a $u$, todos los vértices con aristas hacia $u$ ya se
procesaron, así que $d[u]$ ya es definitiva y sus aristas de salida se
relajan con el valor bueno. Es la misma idea que Dijkstra, fijar un
vértice y relajar lo que sale de él, solo que el orden se conoce de
antemano y no hace falta cola de prioridad.

\begin{teorema}[Cormen, teorema 24.5]\label{teo:caminos-dag-correccion}
  Si $G$ es dirigido, pesado y acíclico, al terminar
  \textsc{Caminos-DAG}$(G, w, s)$ se tiene $d[v] = \delta(s, v)$ para
  todo $v \in V$, y $G_\pi$ es un árbol de caminos más cortos con raíz
  $s$.
\end{teorema}

\begin{proof}
  Si $v$ no es alcanzable desde $s$, la propiedad sin camino da
  $d[v] = \infty = \delta(s, v)$. Si lo es, como no hay ciclos,
  $\delta(s, v)$ es un número real y lo alcanza un camino
  $p = \langle v_0, v_1, \dots, v_k\rangle$ con $v_0 = s$ y $v_k = v$
  (proposición \ref{prop:caminos-simples-fuente}). En el orden
  topológico cada $v_{i-1}$ va antes que $v_i$, porque
  $(v_{i-1}, v_i)$ es una arista, así que la línea 4 procesa
  $v_0, v_1, \dots, v_{k-1}$ en ese orden y las aristas de $p$ se
  relajan en el orden $(v_0, v_1), \dots, (v_{k-1}, v_k)$. La
  relajación de caminos da $d[v] = \delta(s, v)$. El árbol sale de la
  propiedad del subgrafo predecesor.
\end{proof}

\textbf{Costo.} El orden topológico con DFS cuesta $\Tht{n + m}$, e
\textsc{Inicializar-Fuente}, $\Tht{n}$. El ciclo de la línea 4 da $n$
vueltas y el de la línea 5 recorre la lista de $u$. Por análisis
agregado, como en BFS, las listas suman $m$ entradas, así que en total
se hacen exactamente $m$ relajaciones de $\Ord{1}$. Total:
$\Tht{n + m}$, lineal en el tamaño de las listas. Con matriz de
adyacencia sería $\Tht{n^2}$.

\begin{ejemplo}[Caminos en el DAG de la figura 24.5 de Cormen]
  \label{ej:caminos-dag-24-5}
  El grafo tiene $V = \set{r, s, t, x, y, z}$, ya en orden
  topológico (todas las aristas van de izquierda a derecha), fuente
  $s$ y diez aristas:
  \[
    \begin{array}{lllll}
      (r, s) : 5 & (r, t) : 3 & (s, t) : 2 & (s, x) : 6
        & (t, x) : 7 \\
      (t, y) : 4 & (t, z) : 2 & (x, y) : -1 & (x, z) : 1
        & (y, z) : -2
    \end{array}
  \]
  DFS desde $r$, con los vecinos en orden alfabético, da justamente el
  orden $r, s, t, x, y, z$ por $\text{sale}[u]$ decreciente. Una fila
  por vértice procesado, con las $d$ al terminar de relajar sus
  aristas:
  \begin{center}
  \small
  \begin{tabular}{cccccccl}
    \toprule
    $u$ & $r$ & $s$ & $t$ & $x$ & $y$ & $z$
      & \textbf{Relajaciones} \\
    \midrule
    inicial & $\infty$ & $0$ & $\infty$ & $\infty$ & $\infty$
      & $\infty$ & \\
    $r$ & $\infty$ & $0$ & $\infty$ & $\infty$ & $\infty$ & $\infty$
      & $d[r] = \infty$: $(r, s)$ y $(r, t)$ no cambian nada \\
    $s$ & $\infty$ & $0$ & $2$ & $6$ & $\infty$ & $\infty$
      & $t$: $\infty \to 2$; \ $x$: $\infty \to 6$ \\
    $t$ & $\infty$ & $0$ & $2$ & $6$ & $6$ & $4$
      & $x$: $2 + 7 = 9$, no; \ $y$: $\infty \to 6$;
        \ $z$: $\infty \to 4$ \\
    $x$ & $\infty$ & $0$ & $2$ & $6$ & $5$ & $4$
      & $y$: $6 - 1 = 5 < 6$; \ $z$: $6 + 1 = 7$, no \\
    $y$ & $\infty$ & $0$ & $2$ & $6$ & $5$ & $3$
      & $z$: $5 - 2 = 3 < 4$ \\
    $z$ & $\infty$ & $0$ & $2$ & $6$ & $5$ & $3$
      & sin aristas de salida \\
    \bottomrule
  \end{tabular}
  \end{center}
  En dibujos, en verde los vértices ya procesados y en azul las
  aristas $(\pi[v], v)$:
  \begin{center}
  \begin{tikzpicture}[scale=0.85, transform shape,
      distancia/.style={font=\scriptsize\bfseries, text=acento,
        inner sep=1.5pt}]
    \foreach \p/\Y/\dt/\dx/\dy/\dz/\tit in {
        b/0/2/6/\infty/\infty/{(a) tras $u = r$ y $u = s$},
        c/-2.9/2/6/6/4/{(b) tras $u = t$},
        d/-5.8/2/6/5/4/{(c) tras $u = x$},
        e/-8.7/2/6/5/3/{(d) tras $u = y$ y $u = z$: final}} {
      \begin{scope}[yshift=\Y cm]
        \node[vert, label={[distancia]above:$\infty$}] (\p-r)
          at (0, 0) {$r$};
        \node[vert, label={[distancia]above:$0$}] (\p-s) at (1.6, 0)
          {$s$};
        \node[vert, label={[distancia]above:$\dt$}] (\p-t)
          at (3.2, 0) {$t$};
        \node[vert, label={[distancia]above:$\dx$}] (\p-x)
          at (4.8, 0) {$x$};
        \node[vert, label={[distancia]above:$\dy$}] (\p-y)
          at (6.4, 0) {$y$};
        \node[vert, label={[distancia]above:$\dz$}] (\p-z)
          at (8.0, 0) {$z$};
        \draw[dirigida] (\p-r) -- node[peso] {5} (\p-s);
        \draw[dirigida] (\p-s) -- node[peso] {2} (\p-t);
        \draw[dirigida] (\p-t) -- node[peso] {7} (\p-x);
        \draw[dirigida] (\p-x) -- node[peso] {$-1$} (\p-y);
        \draw[dirigida] (\p-y) -- node[peso] {$-2$} (\p-z);
        \draw[dirigida] (\p-s) to[out=40, in=140, looseness=1.3]
          node[peso] {6} (\p-x);
        \draw[dirigida] (\p-x) to[out=40, in=140, looseness=1.3]
          node[peso] {1} (\p-z);
        \draw[dirigida] (\p-r) to[out=-40, in=-140, looseness=1.3]
          node[peso] {3} (\p-t);
        \draw[dirigida] (\p-t) to[out=-40, in=-140, looseness=1.3]
          node[peso] {4} (\p-y);
        \draw[dirigida] (\p-t) to[out=-50, in=-130]
          node[peso] {2} (\p-z);
        \node[font=\small, anchor=west] at (9.0, 0) {\tit};
      \end{scope}
    }
    % aristas de pi en cada instantanea
    \foreach \p in {b, c, d, e} {
      \draw[dirigidaArbol] (\p-s) -- node[peso] {2} (\p-t);
      \draw[dirigidaArbol] (\p-s) to[out=40, in=140, looseness=1.3]
        node[peso] {6} (\p-x);
    }
    \foreach \p in {c, d} {
      \draw[dirigidaArbol] (\p-t) to[out=-50, in=-130]
        node[peso] {2} (\p-z);
    }
    \draw[dirigidaArbol] (c-t) to[out=-40, in=-140, looseness=1.3]
      node[peso] {4} (c-y);
    \foreach \p in {d, e} {
      \draw[dirigidaArbol] (\p-x) -- node[peso] {$-1$} (\p-y);
    }
    \draw[dirigidaArbol] (e-y) -- node[peso] {$-2$} (e-z);
    % vertices ya procesados
    \foreach \p/\lista in {b/{r,s}, c/{r,s,t}, d/{r,s,t,x},
                           e/{r,s,t,x,y,z}} {
      \foreach \v in \lista {
        \node[vertOk] at (\p-\v) {$\v$};
      }
    }
  \end{tikzpicture}
  \end{center}
  \begin{itemize}[itemsep=0.1em]
    \item $r$ va antes que $s$ y no es alcanzable: sus aristas se
          relajan, pero con $d[r] = \infty$ no cambian nada.
    \item A $x$, la arista directa $(s, x)$ con $6$ le gana a
          $s, t, x$ con $9$.
    \item $y$ baja a $6$ por $t$ y luego a $5$ por $x$, gracias a la
          arista de peso $-1$; $z$ baja a $4$ y luego a $3$. Los pesos
          negativos no estorban: cuando le toca a cada vértice, ya le
          llegó todo lo que le podía llegar.
    \item Se hicieron exactamente $10 = m$ relajaciones, una por
          arista.
  \end{itemize}
  Resultado: $\delta(s, \cdot) = \infty, 0, 2, 6, 5, 3$ para
  $r, s, t, x, y, z$, con $\pi[t] = \pi[x] = s$, $\pi[y] = x$ y
  $\pi[z] = y$.
\end{ejemplo}

\subsubsection*{Ruta crítica}

En un \textbf{diagrama PERT} (\textit{program evaluation and review
technique}) las aristas son trabajos y sus pesos, la duración de cada
uno; los vértices son eventos. Si $(u, v)$ entra a $v$ y $(v, x)$ sale
de $v$, el trabajo $(u, v)$ tiene que terminar antes de que empiece
$(v, x)$. Una \textbf{ruta crítica} (\textit{critical path}) es un
camino \emph{más largo} del DAG: la sucesión de trabajos que más
tarda, y por tanto una cota inferior de la duración del proyecto
aunque todo lo demás se haga en paralelo. Hay dos formas de
calcularla, las dos en $\Tht{n + m}$:
\begin{itemize}[itemsep=0.1em]
  \item negar todos los pesos y correr \textsc{Caminos-DAG}: el
        camino más corto del grafo negado es el más largo del
        original;
  \item correr \textsc{Caminos-DAG} cambiando $\infty$ por $-\infty$
        en \textsc{Inicializar-Fuente} y $>$ por $<$ en
        \textsc{Relajar}.
\end{itemize}
Funciona porque en un DAG negar los pesos no crea ciclos negativos.
En un grafo con ciclos el camino más largo es NP-difícil: negar los
pesos de un ciclo positivo da uno negativo, y $\delta = -\infty$ deja
de decir algo sobre caminos sin repetir vértices.

\begin{ejemplo}[Ruta crítica de un proyecto]\label{ej:ruta-critica-pert}
  Un proyecto de investigación con seis eventos, de A (inicio) a F
  (fin). Cada arista es una tarea con su duración en días. A la
  derecha, la traza de \textsc{Caminos-DAG} con los pesos
  \emph{negados}, en el orden topológico A, B, C, D, E, F.
  \begin{center}
  \small
  \begin{tabular}{cll}
    \toprule
    \textbf{Arista} & \textbf{Tarea} & \textbf{Días} \\
    \midrule
    $A \to B$ & leer la literatura & $3$ \\
    $A \to C$ & conseguir los datos & $2$ \\
    $B \to D$ & implementar el modelo & $4$ \\
    $C \to D$ & limpiar los datos & $1$ \\
    $C \to E$ & etiquetar un conjunto de prueba & $6$ \\
    $D \to E$ & entrenar & $2$ \\
    $B \to F$ & escribir el estado del arte & $3$ \\
    $E \to F$ & evaluar y redactar resultados & $2$ \\
    \bottomrule
  \end{tabular}
  \hspace{0.6cm}
  \begin{tabular}{ccccccc}
    \toprule
    $u$ & $A$ & $B$ & $C$ & $D$ & $E$ & $F$ \\
    \midrule
    $A$ & $0$ & $-3$ & $-2$ & $\infty$ & $\infty$ & $\infty$ \\
    $B$ & $0$ & $-3$ & $-2$ & $-7$ & $\infty$ & $-6$ \\
    $C$ & $0$ & $-3$ & $-2$ & $-7$ & $-8$ & $-6$ \\
    $D$ & $0$ & $-3$ & $-2$ & $-7$ & $-9$ & $-6$ \\
    $E$ & $0$ & $-3$ & $-2$ & $-7$ & $-9$ & $-11$ \\
    $F$ & $0$ & $-3$ & $-2$ & $-7$ & $-9$ & $-11$ \\
    \bottomrule
  \end{tabular}
  \end{center}
  Al cambiar el signo, $-d$ es el día más temprano en que puede
  ocurrir cada evento: D no antes del día $7$ (hay que leer e
  implementar, $3 + 4$; los datos limpios están desde el día $3$), E
  no antes del día $9$, y el proyecto no termina antes del día $11$.
  \begin{center}
  \begin{tikzpicture}[distancia/.style={font=\scriptsize\bfseries,
      text=acento, inner sep=1.5pt}]
    \node[vert, label={[distancia]left:$0$}] (A) at (0, 0) {A};
    \node[vert, label={[distancia]above:$3$}] (B) at (2, 1.3) {B};
    \node[vert, label={[distancia]below:$2$}] (C) at (2, -1.3) {C};
    \node[vert, label={[distancia]below:$7$}] (D) at (4.2, 0) {D};
    \node[vert, label={[distancia]below:$9$}] (E) at (6.2, -1.3) {E};
    \node[vert, label={[distancia]right:$11$}] (F) at (8.2, 0) {F};
    \draw[dirigida] (A) -- node[peso] {2} (C);
    \draw[dirigida] (C) -- node[peso] {1} (D);
    \draw[dirigida] (C) -- node[peso] {6} (E);
    \draw[dirigida] (B) to[bend left=12] node[peso] {3} (F);
    \draw[dirigidaArbol] (A) -- node[peso] {3} (B);
    \draw[dirigidaArbol] (B) -- node[peso] {4} (D);
    \draw[dirigidaArbol] (D) -- node[peso] {2} (E);
    \draw[dirigidaArbol] (E) -- node[peso] {2} (F);
  \end{tikzpicture}
  \end{center}
  La ruta crítica es A, B, D, E, F: $3 + 4 + 2 + 2 = 11$ días.
  Retrasar cualquiera de esas cuatro tareas retrasa el proyecto. El
  camino A, C, E, F dura $10$, casi crítico: etiquetar puede atrasarse
  un día sin consecuencias. Lo que \emph{no} sirve es el camino más
  corto, A, B, F con $6$: no dice nada del proyecto, porque todas las
  tareas se tienen que hacer. Las cuatro rutas de A a F duran $11$,
  $6$, $7$ y $10$ días.
\end{ejemplo}

\begin{ojo}
  \textbf{El DAG tiene que ser DAG.} Si el grafo tiene un ciclo no hay
  orden topológico, pero el orden por $\text{sale}[u]$ decreciente se
  calcula igual sin quejarse, a menos que se revisen las aristas de
  retroceso, y \textsc{Caminos-DAG} devuelve basura sin avisar. A
  diferencia de Bellman-Ford, no tiene revisión final.

  \smallskip
  Del otro lado, dos ahorros válidos. Los vértices que van antes de
  $s$ en el orden no son alcanzables, porque todo camino desde $s$ va
  hacia adelante, así que el ciclo de la línea 4 puede arrancar en
  $s$. Y el último vértice del orden no tiene aristas de salida, así
  que tampoco hace falta procesarlo (Cormen, ejercicio 24.2-2).
\end{ojo}

\begin{ejercicio}
  Corre \textsc{Caminos-DAG} sobre el grafo de la figura 24.5 con
  fuente $r$ (Cormen, ejercicio 24.2-1).

  \textit{Solución.} Con el orden $r, s, t, x, y, z$: al procesar $r$
  quedan $s = 5$ y $t = 3$; al procesar $s$, $x = 11$; al procesar
  $t$, $x$ baja a $10$, $y = 7$ y $z = 5$; $x$, $y$ y $z$ ya no
  mejoran nada. Queda $\delta(r, \cdot) = 0, 5, 3, 10, 7, 5$ para
  $r, s, t, x, y, z$, con $\pi[s] = \pi[t] = r$ y
  $\pi[x] = \pi[y] = \pi[z] = t$.
\end{ejercicio}

\begin{ejercicio}
  Da un algoritmo de tiempo $\Tht{n + m}$ que cuente los caminos de un
  DAG (Cormen, ejercicio 24.2-4).

  \textit{Solución.} Sea $P(v)$ el número de caminos con al menos una
  arista que salen de $v$. Cada uno empieza con una arista $(v, u)$ y
  sigue con nada o con uno de los $P(u)$ caminos que salen de $u$, así
  que, en orden topológico inverso,
  \[
    P(v) = \sum_{(v, u) \in E} \bigl(1 + P(u)\bigr) .
  \]
  El total es $\sum_v P(v)$, y cada arista se mira una vez. En la
  figura 24.5 da $P = 21, 12, 7, 3, 1, 0$ para $r, s, t, x, y, z$, en
  total $44$ caminos.
\end{ejercicio}

% =====================================================================

\subsection*{Dijkstra, versión de Cormen}

La clase del 22 de octubre presentó Dijkstra con un ejemplo completo
sobre el grafo pesado (sección \ref{sec:grafos}), que sigue valiendo
tal cual: un grafo no dirigido se trata como un par de arcos opuestos
por arista, y la traza es la misma con el pseudocódigo de abajo. Aquí
se reescribe con \textsc{Inicializar-Fuente} y \textsc{Relajar}, como
lo hace Cormen, para que encaje en el esqueleto común y se pueda
probar con las seis propiedades. La hipótesis es la de siempre:
$w(u, v) \ge 0$ para toda arista.

\begin{algorithm}[H]
\caption{\textsc{Dijkstra}$(G, w, s)$, versión de Cormen}
\label{alg:dijkstra-cormen}
\begin{algorithmic}[1]
\Procedure{Dijkstra}{$G, w, s$}
  \State \Call{Inicializar-Fuente}{$G, s$}
  \State $S \gets \emptyset$
  \State $Q \gets V$ \Comment{cola de mínimos con llave $d$}
  \While{$Q \neq \emptyset$}
    \State $u \gets \Call{Extraer-Min}{Q}$
    \State $S \gets S \cup \set{u}$ \Comment{aquí $d[u] = \delta(s, u)$}
    \ForAll{$v$ en la lista de adyacencia de $u$}
      \State \Call{Relajar}{$u, v, w$} \Comment{si baja $d[v]$:
        \textsc{Decrementar-Llave}}
    \EndFor
  \EndWhile
\EndProcedure
\end{algorithmic}
\end{algorithm}

\begin{itemize}[itemsep=0.1em]
  \item \textbf{Línea 4.} $Q \gets V$ inserta los $n$ vértices de una
        vez. Nunca se inserta nada después y cada vértice sale una
        sola vez, así que el \textbf{mientras} da exactamente $n$
        vueltas. La primera vez sale $s$, el único con $d$ finita.
  \item \textbf{Línea 7.} El algoritmo no necesita $S$, que podría
        leer como $V - Q$; está para la prueba. El comentario marca el
        momento que el teorema de abajo garantiza.
  \item \textbf{Línea 9.} \textsc{Relajar} puede bajar la llave de
        $v$ en $Q$, así que en una implementación con montículo cada
        relajación exitosa es un \textsc{Decrementar-Llave}. Con
        $w \ge 0$ esa relajación nunca le pega a un vértice de $S$: si
        $v \in S$, por el teorema y la desigualdad del triángulo,
        $d[v] = \delta(s, v) \le \delta(s, u) + w(u, v)
        = d[u] + w(u, v)$, y la condición de \textsc{Relajar} falla.
        El decremento siempre actúa sobre un vértice que sigue en el
        montículo. Con pesos negativos esto deja de ser cierto
        (ejemplo \ref{ej:dijkstra-arista-negativa}).
\end{itemize}

\subsubsection*{Por qué funciona}

La clase del 22 de octubre dio el argumento en corto: cuando $u$ sale
con la $d$ mínima, el \emph{primer} vértice no fijado $y$ de un camino
más corto hacia $u$ ya tiene su $d$ exacta, porque su antecesor ya
fijado relajó la arista hacia él, y el tramo de $y$ a $u$ no resta.
Esa pieza es la que hace funcionar todo: sin ella, saber que
$d[y] \ge d[u]$ no dice nada del costo real de pasar por $y$. Aquí va
la prueba completa, apoyada en las propiedades de la relajación.

\begin{teorema}[Corrección de Dijkstra; Cormen, teorema 24.6]
  \label{teo:dijkstra-correccion}
  Si $G = (V, E)$ es dirigido, $w(u, v) \ge 0$ para toda arista y
  $s \in V$, \textsc{Dijkstra}$(G, w, s)$ termina con
  $d[u] = \delta(s, u)$ para todo $u \in V$.
\end{teorema}

\begin{proof}
  Se prueba el invariante: \emph{al inicio de cada vuelta del
  \textbf{mientras}, $d[v] = \delta(s, v)$ para todo $v \in S$}.

  \textit{Inicio.} $S = \emptyset$ y no hay nada que probar.

  \textit{Mantenimiento.} Basta ver que el vértice $u$ que entra a $S$
  en la línea 7 ya trae $d[u] = \delta(s, u)$; después no cambia, por
  la cota superior. Supongamos lo contrario, y sea $u$ el primer
  vértice que entra a $S$ con $d[u] \neq \delta(s, u)$.
  \begin{itemize}[itemsep=0.1em]
    \item $u \neq s$. El primero en salir es $s$, con
          $d[s] = 0 = \delta(s, s)$; como no hay pesos negativos,
          ninguna caminata cerrada pesa menos que $0$. En particular,
          justo antes de agregar $u$, $s \in S$.
    \item Hay un camino de $s$ a $u$. Si no, la propiedad sin camino
          daría $d[u] = \infty = \delta(s, u)$.
    \item Sea $p$ un camino más corto de $s$ a $u$; existe por la
          proposición \ref{prop:caminos-simples-fuente}, porque sin
          pesos negativos no hay ciclos negativos. $p$ empieza en
          $s \in S$ y termina en $u \notin S$. Sea $y$ el primer
          vértice de $p$ que no está en $S$ y $x \in S$ el anterior a
          él en $p$. Puede ser $x = s$ o $y = u$. El camino se parte
          en $s \overset{p_1}{\leadsto} x \to y
          \overset{p_2}{\leadsto} u$.
  \end{itemize}
  \begin{center}
  \begin{tikzpicture}
    \draw[draw=acento, fill=fondoIdea!60] (0, 0) ellipse (1.9 and 1.15);
    \node[nota, text=acento] at (-1.2, 0.75) {$S$};
    \node[vertOk] (s) at (-1.0, -0.1) {$s$};
    \node[vertOk] (x) at (0.9, -0.35) {$x$};
    \node[vert] (y) at (2.9, -0.35) {$y$};
    \node[vert] (u) at (5.2, 0.9) {$u$};
    \draw[dirigida, dashed] (s) to[out=-40, in=200]
      node[nota, below] {$p_1$} (x);
    \draw[dirigida] (x) -- (y);
    \draw[dirigida, dashed] (y) to[out=20, in=230]
      node[nota, below right] {$p_2$} (u);
    \node[nota, anchor=west] at (3.6, -0.8)
      {pesos $\ge 0$: $\delta(s, y) \le \delta(s, u)$};
  \end{tikzpicture}
  \end{center}
  Primero, $d[y] = \delta(s, y)$ en el momento en que $u$ entra a
  $S$. En efecto, $x$ entró a $S$ antes que $u$, así que por la
  elección de $u$ entró con $d[x] = \delta(s, x)$, y en esa vuelta se
  relajó $(x, y)$. Como $s \overset{p_1}{\leadsto} x \to y$ es un
  camino más corto (lema \ref{lem:subcaminos-cortos}), la convergencia
  da $d[y] = \delta(s, y)$ desde entonces.

  Segundo, $y$ va antes que $u$ en $p$ y $p_2$ no tiene pesos
  negativos, así que $\delta(s, y) = w(p_1) + w(x, y) \le w(p) =
  \delta(s, u)$. Con la cota superior,
  \[
    d[y] = \delta(s, y) \le \delta(s, u) \le d[u] .
  \]
  Pero cuando se eligió $u$ en la línea 6, $y$ también estaba en $Q$,
  así que $d[u] \le d[y]$. Todas las desigualdades son igualdades, y
  $d[u] = \delta(s, u)$, lo que contradice la elección de $u$.

  \textit{Terminación.} El ciclo acaba con $Q = \emptyset$, y como
  $Q = V - S$ al inicio de cada vuelta, $S = V$: todos los vértices
  tienen $d[v] = \delta(s, v)$.
\end{proof}

\begin{corolario}[Cormen, corolario 24.7]\label{cor:dijkstra-arbol}
  Con las mismas hipótesis, al terminar $G_\pi$ es un árbol de caminos
  más cortos con raíz $s$.
\end{corolario}

\begin{proof}
  Inmediato del teorema \ref{teo:dijkstra-correccion} y la propiedad
  del subgrafo predecesor; con $w \ge 0$ no hay ciclos negativos.
\end{proof}

La prueba usa la hipótesis $w \ge 0$ en un solo lugar, la desigualdad
$\delta(s, y) \le \delta(s, u)$. Si el tramo $p_2$ resta, $y$ puede
estar más lejos que $u$ aunque vaya antes en el camino, y el
argumento se cae justo ahí.

\subsubsection*{Costo según la cola de prioridad}

Fuera de las operaciones de la cola, todo cuesta $\Ord{1}$ por vértice
o por arista. Las operaciones se cuentan así:
\begin{itemize}[itemsep=0.1em]
  \item \textbf{Insertar}: $n$, todas en la línea 4.
  \item \textbf{Extraer-Min}: una por vuelta del \textbf{mientras},
        $n$ en total.
  \item \textbf{Decrementar-Llave}: a lo más una por relajación, y
        aquí entra el análisis agregado. Cada vértice entra a $S$ una
        sola vez, así que la lista de adyacencia de cada $u$ se
        recorre una sola vez, en la vuelta en que sale $u$. Las listas
        suman $m$ entradas: en toda la ejecución hay exactamente $m$
        relajaciones y a lo más $m$ decrementos. Contar por vuelta
        daría $n$ vueltas con hasta $n - 1$ decrementos cada una,
        $\Ord{n^2}$, una cota mucho peor en un grafo disperso.
\end{itemize}
El total es
$T = T_{\text{construir}} + n \cdot T_{\text{extraer}}
+ m \cdot T_{\text{decrementar}}$, y todo depende de la cola:
\begin{center}
\small
\begin{tabular}{lcccc}
  \toprule
  \textbf{Cola} & \textbf{Insertar} & \textbf{Extraer-Min}
    & \textbf{Decrementar} & \textbf{Total} \\
  & ($n$ veces) & ($n$ veces) & ($\le m$ veces) & \\
  \midrule
  arreglo por vértice & $\Ord{1}$ & $\Ord{n}$ & $\Ord{1}$
    & $\Ord{n^2 + m} = \Ord{n^2}$ \\
  montículo binario & $\Ord{n}$ en total & $\Ord{\log n}$
    & $\Ord{\log n}$ & $\Ord{(n + m) \log n}$ \\
  montículo de Fibonacci & $\Ord{1}$ & $\Ord{\log n}$ amort.
    & $\Ord{1}$ amort. & $\Ord{n \log n + m}$ \\
  \bottomrule
\end{tabular}
\end{center}
\begin{itemize}[itemsep=0.1em]
  \item \textbf{Arreglo indexado por vértice.} $d[v]$ vive en la
        casilla $v$, con una marca de si sigue en $Q$. Insertar y
        decrementar son asignaciones; extraer el mínimo obliga a
        recorrer el arreglo, $\Tht{n}$. Total $\Ord{n^2 + m}
        = \Ord{n^2}$, porque $m \le n^2$.
  \item \textbf{Montículo binario.} Construirlo cuesta $\Ord{n}$, como
        en la clase del 17 de septiembre; extraer y decrementar,
        $\Ord{\log n}$ cada una. Para decrementar hay que saber dónde
        está $v$ dentro del montículo: cada vértice guarda su
        posición y el montículo la actualiza en cada intercambio.
        Total $\Ord{(n + m) \log n}$, que es $\Ord{m \log n}$ si todos
        los vértices son alcanzables desde $s$, porque entonces
        $m \ge n - 1$.
  \item \textbf{Montículo de Fibonacci.} Extraer cuesta $\Ord{\log n}$
        amortizado y decrementar $\Ord{1}$ amortizado. Total
        $\Ord{n \log n + m}$. Cormen cuenta que se inventaron justo
        por esto: Dijkstra hace muchos más decrementos que
        extracciones, y abaratar el decremento sin encarecer la
        extracción mejora el total.
\end{itemize}

\textbf{El umbral de densidad.} El montículo binario le gana al
arreglo cuando $(n + m)\log n = \ord{n^2}$, o sea cuando
$m = \ord{n^2 / \log n}$. En un grafo denso, $m = \Tht{n^2}$, el
binario cuesta $\Tht{n^2 \log n}$ y el arreglo gana. Fibonacci nunca
pierde asintóticamente: en denso empata con el arreglo, en muy
disperso ($m = \Ord{n}$) empata con el binario, y en medio gana. Con
$n = 10^4$, sin constantes:
\begin{center}
\small
\begin{tabular}{lccc}
  \toprule
  $m$ & arreglo, $n^2$ & binario, $(n + m)\log_2 n$
    & Fibonacci, $n \log_2 n + m$ \\
  \midrule
  $5 \cdot 10^4$ (disperso) & $1.0 \cdot 10^8$ & $8.0 \cdot 10^5$
    & $1.8 \cdot 10^5$ \\
  $10^6 = n^{1.5}$ & $1.0 \cdot 10^8$ & $1.3 \cdot 10^7$
    & $1.1 \cdot 10^6$ \\
  $5 \cdot 10^7$ (denso) & $1.0 \cdot 10^8$ & $6.6 \cdot 10^8$
    & $5.0 \cdot 10^7$ \\
  \bottomrule
\end{tabular}
\end{center}
Aquí el umbral $n^2 / \log_2 n$ queda en unos $7.5 \cdot 10^6$
aristas. En la práctica casi siempre se usa el montículo binario: los
grafos reales suelen ser dispersos y las constantes del de Fibonacci
son grandes.

\subsubsection*{Parecidos con BFS y con Prim}

\textbf{BFS.} $S$ es el conjunto de vértices negros de BFS: en los
dos casos, lo que está ahí ya tiene su distancia final. Con todos los
pesos iguales a $1$, la cola de prioridad se comporta como la cola de
BFS y los vértices salen por capas de distancia. Cormen lo hace
preciso en el ejercicio 24.3-7: si cada arista de peso entero $w$ se
cambia por una cadena de $w$ aristas de peso $1$, BFS sobre el grafo
nuevo pinta de negro los vértices originales en el mismo orden en que
Dijkstra los extrae, salvo empates. Dijkstra es BFS con una cola que
sabe de pesos.

\textbf{Prim.} Los dos hacen crecer un conjunto, el árbol en Prim y
$S$ en Dijkstra, usan una cola de prioridad de mínimos para elegir el
vértice más ligero fuera de él, lo agregan y ajustan las llaves de
sus vecinos. Solo cambia la llave: $d[u] + w(u, v)$ en Dijkstra,
$w(u, v)$ en Prim (sección \ref{sec:mst}). Por eso cuestan lo mismo
con cada cola. Pero resuelven problemas distintos, y sus árboles
difieren, como mostró la clase del 22 de octubre al comparar el árbol
de Dijkstra con el de Kruskal.

\subsubsection*{Con una arista negativa}

\begin{ejemplo}[Dijkstra falla con una arista negativa]
  \label{ej:dijkstra-arista-negativa}
  Grafo dirigido con $A \to B$ de peso $2$, $A \to C$ de peso $3$,
  $C \to B$ de peso $-2$ y $B \to D$ de peso $2$, fuente A. Las
  distancias reales son $\delta(A, B) = 1$ por A, C, B y
  $\delta(A, D) = 3$ por A, C, B, D. Así corre el pseudocódigo de
  arriba; en negritas, los vértices de $S$.
  \begin{center}
  \small
  \begin{tabular}{cccccc}
    \toprule
    \textbf{Iter.} & \textbf{Sale} & A & B & C & D \\
    \midrule
    0 & & $0$ & $\infty$ & $\infty$ & $\infty$ \\
    1 & A & $\mathbf{0}$ & $2$ & $3$ & $\infty$ \\
    2 & B & $\mathbf{0}$ & $\mathbf{2}$ & $3$ & $4$ \\
    3 & C & $\mathbf{0}$ & $\mathbf{1}$ & $\mathbf{3}$ & $4$ \\
    4 & D & $\mathbf{0}$ & $\mathbf{1}$ & $\mathbf{3}$ & $\mathbf{4}$
      \\
    \midrule
    real & & $0$ & $1$ & $3$ & $3$ \\
    \bottomrule
  \end{tabular}
  \end{center}
  B sale en la iteración 2 con $d = 2$ y relaja $B \to D$: D queda en
  $4$. En la iteración 3 sale C y la relajación $C \to B$ baja $d[B]$
  a $3 - 2 = 1$, pero B ya está en $S$ y nadie vuelve a relajar
  $B \to D$. D sale con $4$ y la respuesta es incorrecta. En la
  prueba del teorema, cuando sale B el camino más corto hacia él es
  A, C, B; ahí $y = C$, con $d[C] = \delta(A, C) = 3$, pero
  $\delta(A, B) = 1 < 3$ porque $p_2 = (C, B)$ pesa $-2$.
  \begin{center}
  \begin{tikzpicture}[scale=0.9, transform shape]
    \foreach \px/\cap/\dA/\eA/\dB/\eB/\dC/\eC/\dD/\eD/\tAB/\tCB/\tBD
      in {
        0/{(a) sale A}/0/vertOk/2/vertMarca/3/vert/\infty/vert/%
          dirigidaArbol/dirigida/dirigida,
        5.4/{(b) sale B con $d = 2$}/0/vertOk/2/vertOk/3/vertMarca/%
          4/vert/dirigidaArbol/dirigida/dirigidaArbol,
        10.8/{(c) sale C: B baja, D no}/0/vertOk/1/vertOk/3/vertOk/%
          4/vertMarca/dirigida/dirigidaResalte/dirigidaArbol} {
      \begin{scope}[xshift=\px cm]
        \node[\eA, label={[nota, text=black]left:A}] (A) at (0, 0)
          {$\dA$};
        \node[\eB, label={[nota, text=black]above:B}] (B) at (1.9, 1)
          {$\dB$};
        \node[\eC, label={[nota, text=black]below:C}] (C)
          at (1.9, -1) {$\dC$};
        \node[\eD, label={[nota, text=black]above:D}] (D) at (3.8, 1)
          {$\dD$};
        \draw[\tAB] (A) -- node[peso] {2} (B);
        \draw[dirigidaArbol] (A) -- node[peso] {3} (C);
        \draw[\tCB] (C) -- node[peso] {$-2$} (B);
        \draw[\tBD] (B) -- node[peso] {2} (D);
        \node[font=\small] at (1.9, -2.1) {\cap};
      \end{scope}
    }
  \end{tikzpicture}
  \end{center}
  Dentro de cada vértice, su $d$. Verde: en $S$. Naranja: el mínimo de
  $Q$, que sale en la siguiente iteración. Azul: las aristas
  $(\pi[v], v)$. La arista resaltada en (c) es la relajación tardía:
  el árbol de predecesores ya es el correcto, A, C, B, D, pero
  $d[D] = 4 > d[B] + w(B, D) = 3$, porque el valor de D se calculó con
  el $d[B]$ viejo.
\end{ejemplo}

\begin{ojo}
  \textbf{Por qué hace falta D.} Con solo A, B y C, el pseudocódigo de
  Dijkstra que relaja todos los vecinos sin preguntar si ya salieron
  (el de la clase del 22 de octubre) termina con
  $d[B] = 1$, que es correcto: la relajación tardía de $C \to B$ sí
  baja $d[B]$, aunque B ya esté en $S$. Lo que falla es la afirmación
  de que $d[B]$ quedó fija al salir y, con montículo, el
  \textsc{Decrementar-Llave} sobre un vértice que ya no está en él.
  El error llega a la salida en cuanto algo depende de B, y para eso
  está D. Muchas implementaciones con montículo se saltan los vecinos
  que ya están en $S$; esa variante ni siquiera corrige a B y termina
  con $d = 0, 2, 3, 4$ para A, B, C, D. En ningún caso hay garantía.
\end{ojo}

\begin{ojo}
  \textbf{No siempre falla.} Si las únicas aristas negativas salen de
  $s$ y no hay ciclos negativos, Dijkstra es correcto (Cormen,
  ejercicio 24.3-10): en la prueba, $p_2$ va de $y \neq s$ a $u$ por
  un camino que no vuelve a pasar por $s$, así que no tiene aristas
  negativas. Con pesos negativos se pierde la garantía, no
  necesariamente la respuesta.

  \smallskip
  \textbf{Sumar una constante no arregla nada.} Es el arreglo que se
  le ocurre a todos: sumar a cada peso lo suficiente para que ninguno
  sea negativo. Falla porque castiga más a los caminos con más
  aristas. En el ejemplo, sumando $2$, $A \to B$ pesa $4$ y A, C, B
  pesa $5 + 0 = 5$, y el camino más corto a B cambia de A, C, B a la
  arista directa. La forma correcta de volver no negativos los pesos
  es la repesada de Johnson (sección \ref{sec:todos-pares}), que
  calcula sus potenciales con Bellman-Ford.

  \smallskip
  \textbf{Montículo sin decremento.} Bibliotecas como \texttt{heapq}
  de Python no tienen \textsc{Decrementar-Llave}. Lo usual es meter un
  par $(d[v], v)$ nuevo cada vez que baja $d[v]$ e ignorarlo al
  sacarlo si $v$ ya está en $S$. El montículo llega a tener hasta $m$
  entradas y el costo es $\Ord{m \log m} = \Ord{m \log n}$, porque
  $m \le n^2$.
\end{ojo}

\begin{ejercicio}
  Si el \textbf{mientras} se cambia por ``mientras $\abs{Q} > 1$'', el
  ciclo da $n - 1$ vueltas. ¿Sigue siendo correcto? (Cormen,
  ejercicio 24.3-3.)

  \textit{Solución.} Sí. El último vértice que queda en $Q$ es el
  mínimo de una cola de un elemento, y la prueba del teorema
  \ref{teo:dijkstra-correccion} muestra que el mínimo de $Q$ ya tiene
  $d = \delta$. Relajar sus aristas no cambiaría nada, porque todos
  los demás vértices están en $S$ con su distancia final y no bajan de
  ella.
\end{ejercicio}

\begin{ejercicio}
  Cada arista tiene una confiabilidad $r(u, v) \in [0, 1]$, la
  probabilidad de que no falle, independiente de las demás. Encuentra
  el camino más confiable entre dos vértices (Cormen, ejercicio
  24.3-6).

  \textit{Solución.} Maximizar el producto de las $r$ es minimizar la
  suma de los $-\log r(u, v)$, que son pesos no negativos porque
  $r \le 1$; las aristas con $r = 0$ se descartan. Se corre Dijkstra.
  Equivale a correr Dijkstra con producto y máximo en lugar de suma y
  mínimo: funciona porque multiplicar por $r \le 1$ nunca mejora, el
  análogo de $w \ge 0$.
\end{ejercicio}

\begin{ejercicio}
  Si los pesos son enteros en $\set{0, 1, \dots, W}$, haz que Dijkstra
  cueste $\Ord{Wn + m}$ (Cormen, ejercicio 24.3-8).

  \textit{Solución.} Toda distancia finita es a lo más $(n - 1)W$,
  así que la cola puede ser un arreglo de cubetas indexado por
  distancia, con un apuntador que solo avanza, porque las extracciones
  salen en orden no decreciente de $d$. Recorrer las cubetas cuesta
  $\Ord{Wn}$ en total y cada decremento mueve un vértice de cubeta en
  $\Ord{1}$.
\end{ejercicio}

% =====================================================================

\subsection*{Restricciones de diferencia}

Un \textbf{programa lineal} pide un vector $x$ que maximice una
función lineal $c^{T} x$ sujeto a un sistema de desigualdades
$Ax \le b$. A veces no hay función objetivo y solo interesa la
\textbf{factibilidad}: encontrar algún $x$ con $Ax \le b$ o decidir
que no existe. Hay algoritmos polinomiales para programación lineal en
general, pero muchos casos particulares se resuelven más rápido con
algoritmos de grafos. Este es uno.

\begin{definicion}[Sistema de restricciones de diferencia]
  \label{def:restricciones-diferencia}
  Un \textbf{sistema de restricciones de diferencia} (\textit{system
  of difference constraints}) es un programa lineal $Ax \le b$ en el
  que cada renglón de $A$ tiene un $1$, un $-1$ y ceros. Con
  incógnitas $x_1, \dots, x_p$ y $r$ restricciones, el renglón $k$
  dice
  \[
    x_j - x_i \le b_k , \qquad 1 \le i, j \le p,\ i \neq j,\
    1 \le k \le r .
  \]
\end{definicion}

Cormen llama $n$ al número de incógnitas y $m$ al de restricciones.
Aquí son $p$ y $r$, porque $n$ y $m$ se quedan para los vértices y las
aristas del grafo que se va a construir, que tendrá $n = p + 1$ y
$m = p + r$.

La aplicación natural es la calendarización. Si $x_i$ es el instante
en que empieza la tarea $i$, ``la tarea $j$ empieza a lo más $b$
unidades después de la $i$'' es $x_j - x_i \le b$, y ``la tarea $j$
empieza al menos $c$ unidades después de la $i$'' es
$x_j \ge x_i + c$, o sea $x_i - x_j \le -c$. Estas últimas, con lado
derecho negativo, son las que van a obligar a usar Bellman-Ford.

\begin{lema}[Traslación; Cormen, lema 24.8]
  \label{lem:traslacion-restricciones}
  Si $x = (x_1, \dots, x_p)$ es solución de un sistema de
  restricciones de diferencia y $a$ es cualquier constante,
  $x + a = (x_1 + a, \dots, x_p + a)$ también lo es.
\end{lema}

\begin{proof}
  $(x_j + a) - (x_i + a) = x_j - x_i$.
\end{proof}

Un sistema con solución tiene entonces una infinidad, y siempre se
puede fijar una incógnita en el valor que se quiera. En
calendarización: correr todo el calendario no cambia nada.

La observación clave es de forma. La restricción $x_j - x_i \le b_k$
se lee $x_j \le x_i + b_k$, que es la desigualdad del triángulo
$\delta(s, v_j) \le \delta(s, v_i) + w(v_i, v_j)$ con
$w(v_i, v_j) = b_k$. Si las incógnitas fueran distancias desde una
fuente común en un grafo con esas aristas, cumplirían todas las
restricciones solas.

\begin{definicion}[Grafo de restricciones]
  \label{def:grafo-restricciones}
  El \textbf{grafo de restricciones} (\textit{constraint graph}) del
  sistema es el grafo dirigido y pesado con vértices
  $V = \set{v_0, v_1, \dots, v_p}$, uno por incógnita y uno extra
  $v_0$, y aristas
  \begin{itemize}[itemsep=0.1em]
    \item $(v_i, v_j)$ con peso $b_k$ por cada restricción
          $x_j - x_i \le b_k$;
    \item $(v_0, v_i)$ con peso $0$, para $i = 1, \dots, p$.
  \end{itemize}
\end{definicion}

La arista va de la variable que resta a la que suma. El vértice $v_0$
es la fuente común: llega a todos, así que ningún $\delta(v_0, v_i)$
es $\infty$, y cualquier ciclo negativo del grafo es alcanzable desde
él.

\begin{teorema}[Cormen, teorema 24.9]\label{teo:restricciones-solucion}
  Sea $G$ el grafo de restricciones de un sistema de restricciones de
  diferencia. Si $G$ no tiene ciclos de peso negativo,
  \[
    x = \bigl(\delta(v_0, v_1), \delta(v_0, v_2), \dots,
             \delta(v_0, v_p)\bigr)
  \]
  es una solución. Si $G$ tiene un ciclo de peso negativo, el sistema
  no tiene solución.
\end{teorema}

\begin{proof}
  \textit{Sin ciclos negativos.} Todo $v_i$ es alcanzable desde $v_0$
  por la arista $(v_0, v_i)$, y no hay ciclos negativos, así que cada
  $\delta(v_0, v_i)$ es un número real (proposición
  \ref{prop:caminos-simples-fuente}). Para cada restricción
  $x_j - x_i \le b_k$ hay una arista $(v_i, v_j)$ de peso $b_k$, y la
  desigualdad del triángulo dice
  \[
    \delta(v_0, v_j) \le \delta(v_0, v_i) + b_k ,
  \]
  que con $x_i = \delta(v_0, v_i)$ es exactamente
  $x_j - x_i \le b_k$.

  \textit{Con un ciclo negativo.} Sea
  $c = \langle u_0, u_1, \dots, u_q\rangle$, con $u_0 = u_q$, un ciclo
  de peso negativo. No pasa por $v_0$, porque a $v_0$ no le entra
  ninguna arista, así que cada $u_l$ es el vértice de una incógnita,
  que se escribe $x(u_l)$. Cada arista $(u_{l-1}, u_l)$ del ciclo
  viene de una restricción
  \[
    x(u_l) - x(u_{l-1}) \le w(u_{l-1}, u_l) ,
    \qquad l = 1, \dots, q .
  \]
  Si $x$ fuera solución, cumpliría las $q$. Al sumarlas, a la
  izquierda cada incógnita del ciclo aparece una vez sumando y una vez
  restando, y queda $x(u_q) - x(u_0) = 0$; a la derecha queda
  $w(c)$. Entonces $0 \le w(c) < 0$, imposible.
\end{proof}

Para resolver el sistema basta construir el grafo y correr
Bellman-Ford desde $v_0$. Dijkstra no sirve, porque los $b_k$ pueden
ser negativos. Como $v_0$ alcanza a todos, el teorema
\ref{teo:bellman-ford-correccion} dice que Bellman-Ford devuelve falso
exactamente cuando el grafo tiene un ciclo negativo; si devuelve
verdadero, $d[v_i] = \delta(v_0, v_i)$. Cada restricción se da como
una terna $(i, j, b)$, que significa $x_j - x_i \le b$.

\begin{algorithm}[H]
\caption{\textsc{Resolver-Restricciones}$(p, R)$}
\label{alg:resolver-restricciones}
\begin{algorithmic}[1]
\Procedure{Resolver-Restricciones}{$p, R$}
  \State $V \gets \set{v_0, v_1, \dots, v_p}$, \quad $E \gets \emptyset$
  \For{$i \gets 1$ \Hasta $p$}
    \State agregar $(v_0, v_i)$ a $E$ con peso $0$
  \EndFor
  \ForAll{$(i, j, b) \in R$}
    \State agregar $(v_i, v_j)$ a $E$ con peso $b$
  \EndFor
  \If{\Call{Bellman-Ford}{$G, w, v_0$}}
    \State \Return $(d[v_1], d[v_2], \dots, d[v_p])$
  \Else
    \State \Return no hay solución \Comment{hay un ciclo negativo}
  \EndIf
\EndProcedure
\end{algorithmic}
\end{algorithm}

\textbf{Costo.} El grafo tiene $n = p + 1$ vértices y $m = p + r$
aristas, y construirlo cuesta $\Tht{p + r}$. Bellman-Ford cuesta
$\Tht{nm} = \Tht{(p + 1)(p + r)} = \Tht{p^2 + pr}$. El término $p^2$
lo ponen las $p$ aristas de $v_0$, que se relajan en cada una de las
$p$ pasadas sin aportar nada después de la primera; el ejercicio
\ref{ej:restricciones-rapido} lo quita.

\begin{ejemplo}[El sistema de la figura 24.8 de Cormen]
  \label{ej:restricciones-24-8}
  Cinco incógnitas y ocho restricciones:
  \[
    \begin{array}{llll}
      x_1 - x_2 \le 0 & x_1 - x_5 \le -1 & x_2 - x_5 \le 1
        & x_3 - x_1 \le 5 \\
      x_4 - x_1 \le 4 & x_4 - x_3 \le -1 & x_5 - x_3 \le -3
        & x_5 - x_4 \le -3
    \end{array}
  \]
  Bellman-Ford desde $v_0$, relajando primero las cinco aristas de
  $v_0$ y luego las ocho restricciones en el orden de arriba:
  \begin{center}
  \small
  \begin{tabular}{lccccc}
    \toprule
    & $d[v_1]$ & $d[v_2]$ & $d[v_3]$ & $d[v_4]$ & $d[v_5]$ \\
    \midrule
    inicio & $\infty$ & $\infty$ & $\infty$ & $\infty$ & $\infty$ \\
    pasada 1 & $-1$ & $0$ & $0$ & $-1$ & $-4$ \\
    pasada 2 & $-5$ & $-3$ & $0$ & $-1$ & $-4$ \\
    pasadas 3 a 5 & \multicolumn{5}{c}{sin cambios} \\
    \bottomrule
  \end{tabular}
  \end{center}
  Al final ninguna arista se puede relajar, así que Bellman-Ford
  devuelve verdadero, y la solución es $x = (-5, -3, 0, -1, -4)$.
  \begin{center}
  \begin{tikzpicture}[scale=1.05]
    \node[vertOk, label={[nota, text=black]left:$v_0$}] (v0)
      at (0, 0) {$0$};
    \node[vert, label={[nota, text=black]above:$v_1$}] (v1)
      at (3.5, 2.0) {$-5$};
    \node[vert, label={[nota, text=black]right:$v_2$}] (v2)
      at (5.6, 0.8) {$-3$};
    \node[vert, label={[nota, text=black]right:$v_3$}] (v3)
      at (4.7, -1.5) {$0$};
    \node[vert, label={[nota, text=black]below:$v_4$}] (v4)
      at (2.5, -1.5) {$-1$};
    \node[vert, label={[nota, text=black]above left:$v_5$}] (v5)
      at (1.9, 0.8) {$-4$};
    % aristas desde v0, todas de peso 0
    \draw[dirigida] (v0) -- node[peso] {0} (v5);
    \draw[dirigida] (v0) to[out=70, in=180] node[peso] {0} (v1);
    \draw[dirigida] (v0) .. controls (0.6, 3.6) and (5.6, 3.4) ..
      node[peso] {0} (v2);
    \draw[dirigida] (v0) to[out=-70, in=180] node[peso] {0} (v4);
    \draw[dirigidaArbol] (v0) .. controls (0.6, -3.0) and (4.0, -3.0)
      .. node[peso] {0} (v3);
    % una arista por restriccion: x_j - x_i <= b va de v_i a v_j
    \draw[dirigida] (v2) -- node[peso] {0} (v1);
    \draw[dirigidaArbol] (v5) -- node[peso] {$-1$} (v1);
    \draw[dirigidaArbol] (v5) -- node[peso, pos=0.25] {1} (v2);
    \draw[dirigida] (v1) -- node[peso, pos=0.7] {5} (v3);
    \draw[dirigida] (v1) -- node[peso, pos=0.75] {4} (v4);
    \draw[dirigidaArbol] (v3) -- node[peso] {$-1$} (v4);
    \draw[dirigida] (v3) -- node[peso, pos=0.3] {$-3$} (v5);
    \draw[dirigidaArbol] (v4) -- node[peso] {$-3$} (v5);
  \end{tikzpicture}
  \end{center}
  Dentro de cada vértice, $\delta(v_0, v_i)$. En azul, el árbol de
  caminos más cortos: $v_0 \to v_3 \to v_4 \to v_5$, que luego se
  abre hacia $v_1$ y $v_2$. Las aristas del árbol son restricciones
  justas, que se cumplen con igualdad; solo $x_1 - x_2 \le 0$ y
  $x_5 - x_3 \le -3$ tienen holgura, de $2$ y de $1$. Sumando $5$ a
  todo sale $x' = (0, 2, 5, 4, 1)$, que también es solución y tiene
  tiempos no negativos. Con las restricciones en el orden inverso,
  Bellman-Ford tarda una pasada más en estabilizarse y llega a lo
  mismo.
\end{ejemplo}

\begin{ejemplo}[Un sistema sin solución]
  \label{ej:restricciones-sin-solucion}
  Las restricciones $x_2 - x_1 \le 1$, $x_3 - x_2 \le -2$ y
  $x_1 - x_3 \le 0$ forman el ciclo $v_1 \to v_2 \to v_3 \to v_1$ de
  peso $1 - 2 + 0 = -1$, y sumarlas da $0 \le -1$. Bellman-Ford
  relaja primero las aristas de $v_0$ y luego las tres restricciones
  en ese orden.
  \begin{center}
  \begin{tikzpicture}[baseline=(current bounding box.center)]
    \node[vertOk, label={[nota, text=black]below:$v_0$}] (v0)
      at (0, 0) {$0$};
    \node[vert, label={[nota, text=black]above:$v_1$}] (v1)
      at (0, 1.8) {};
    \node[vert, label={[nota, text=black]right:$v_2$}] (v2)
      at (1.9, -1.1) {};
    \node[vert, label={[nota, text=black]left:$v_3$}] (v3)
      at (-1.9, -1.1) {};
    \foreach \v in {v1, v2, v3}
      \draw[dirigida] (v0) -- node[peso] {0} (\v);
    \draw[dirigidaFuera] (v1) to[bend left=25] node[peso] {1} (v2);
    \draw[dirigidaFuera] (v2) to[bend left=25] node[peso] {$-2$}
      (v3);
    \draw[dirigidaFuera] (v3) to[bend left=25] node[peso] {0} (v1);
  \end{tikzpicture}
  \hspace{1cm}
  \small
  \begin{tabular}{lccc}
    \toprule
    & $d[v_1]$ & $d[v_2]$ & $d[v_3]$ \\
    \midrule
    pasada 1 & $-2$ & $0$ & $-2$ \\
    pasada 2 & $-3$ & $-1$ & $-3$ \\
    pasada 3 & $-4$ & $-2$ & $-4$ \\
    \bottomrule
  \end{tabular}
  \end{center}
  Con $n = 4$ vértices, Bellman-Ford hace tres pasadas, y cada una
  baja todo en $1$, el peso del ciclo. La revisión encuentra que
  $(v_1, v_2)$ todavía relaja, $d[v_2] = -2 > d[v_1] + 1 = -3$, y
  devuelve falso. Los predecesores al final, $\pi[v_2] = v_1$,
  $\pi[v_3] = v_2$ y $\pi[v_1] = v_3$, recorren justo el ciclo: el
  propio $\pi$ es el certificado de que no hay solución.
\end{ejemplo}

¿Qué solución sale, de las infinitas que hay? Por la arista
$(v_0, v_i)$ de peso $0$, siempre $\delta(v_0, v_i) \le 0$ (Cormen,
ejercicio 24.4-3); si todos los $b_k$ son no negativos, $x = 0$ ya es
solución y es la que sale. Entre las soluciones con todas las
$x_i \le 0$, la de Bellman-Ford es la de suma máxima (ejercicio
24.4-8), y entre todas las soluciones minimiza
$\max_i x_i - \min_i x_i$ (ejercicio 24.4-9): en calendarización,
termina el proyecto en el menor intervalo posible.

\begin{observacion}
  El vértice $v_0$, con aristas de peso $0$ hacia todos, es el mismo
  truco del vértice extra del algoritmo de Johnson (sección
  \ref{sec:todos-pares}). Johnson lo agrega a un grafo cualquiera,
  corre Bellman-Ford desde él y usa $h(v) = \delta(v_0, v)$ para
  repesar cada arista como $w(u, v) + h(u) - h(v)$, que la desigualdad
  del triángulo hace no negativo. Visto desde aquí, $h$ es una
  solución del sistema $h(v) - h(u) \le w(u, v)$, con una restricción
  por arista: Johnson resuelve un sistema de restricciones de
  diferencia antes de correr Dijkstra.
\end{observacion}

\begin{ojo}
  \textbf{La dirección de la arista.} $x_j - x_i \le b$ da la arista
  $(v_i, v_j)$, no $(v_j, v_i)$. Si se invierten todas, lo que sale es
  $-x$ para una solución $x$ y hay que negar la respuesta; mezclar las
  dos convenciones da basura.

  \smallskip
  \textbf{Otras formas de restricción.} $x_j - x_i \ge b$ se reescribe
  $x_i - x_j \le -b$. Una igualdad $x_i = x_j + b$ son dos
  desigualdades, $x_i - x_j \le b$ y $x_j - x_i \le -b$. Las
  desigualdades estrictas no entran directamente; con incógnitas
  enteras, $x_j - x_i < b$ equivale a $x_j - x_i \le b - 1$.

  \smallskip
  \textbf{Restricciones de una sola variable}, como $x_i \le b$ o
  $-x_i \le b$ (Cormen, ejercicio 24.4-10). Es tentador usar $v_0$
  como el cero, pero sus aristas de peso $0$ agregarían a escondidas
  $x_i \le 0$, y con $-x_1 \le -5$ un sistema factible se volvería
  infactible. Lo correcto es una incógnita nueva $z$ que haga de cero,
  con $x_i - z \le b$ o $z - x_i \le b$, resolver el sistema ampliado
  y devolver $x_i - z$, que es solución por el lema
  \ref{lem:traslacion-restricciones}.
\end{ojo}

\begin{ejercicio}\label{ej:restricciones-rapido}
  Resuelve un sistema de $p$ incógnitas y $r$ restricciones en
  $\Ord{p + pr}$, sin construir $v_0$ (Cormen, ejercicios 24.4-5 y
  24.4-7).

  \textit{Solución.} Se inicializa $x_i \gets 0$ para todo $i$, que es
  lo que dejan las aristas de $v_0$ en cuanto se relajan. Después se
  hacen $p - 1$ pasadas sobre las $r$ restricciones, relajando
  $x_j \gets \min(x_j, x_i + b_k)$, y una pasada de revisión. Bastan
  $p - 1$: un camino desde $v_0$ tiene a lo más $p$ aristas, y la
  primera es una $(v_0, v_i)$, que la inicialización ya relajó. Una
  cadena $x_{i+1} - x_i \le -1$ listada de atrás hacia adelante sí
  necesita las $p - 1$ pasadas.
\end{ejercicio}

\begin{ejercicio}
  Encuentra una solución del sistema, o decide que no existe (Cormen,
  ejercicio 24.4-1):
  \[
    \begin{array}{lllll}
      x_1 - x_2 \le 1 & x_1 - x_4 \le -4 & x_2 - x_3 \le 2
        & x_2 - x_5 \le 7 & x_2 - x_6 \le 5 \\
      x_3 - x_6 \le 10 & x_4 - x_2 \le 2 & x_5 - x_1 \le -1
        & x_5 - x_4 \le 3 & x_6 - x_3 \le -8
    \end{array}
  \]

  \textit{Solución.} Bellman-Ford desde $v_0$ devuelve verdadero con
  $x = (-5, -3, 0, -1, -6, -8)$. Sumando $8$, $(3, 5, 8, 7, 2, 0)$ es
  una solución con valores no negativos.
\end{ejercicio}

\begin{ejercicio}
  Decide si tiene solución el sistema (Cormen, ejercicio 24.4-2)
  \[
    \begin{array}{lllll}
      x_1 - x_2 \le 4 & x_1 - x_5 \le 5 & x_2 - x_4 \le -6
        & x_3 - x_2 \le 1 & x_4 - x_1 \le 3 \\
      x_4 - x_3 \le 5 & x_4 - x_5 \le 10 & x_5 - x_3 \le -4
        & x_5 - x_4 \le -8 &
    \end{array}
  \]

  \textit{Solución.} No tiene. El grafo de restricciones tiene el
  ciclo $v_4 \to v_2 \to v_3 \to v_5 \to v_1 \to v_4$, de peso
  $-6 + 1 - 4 + 5 + 3 = -1$, y sumar sus cinco restricciones da
  $0 \le -1$. Bellman-Ford desde $v_0$ devuelve falso.
\end{ejercicio}

% =====================================================================

\subsection*{Resumen}

\begin{center}
\small
\begin{tabular}{l>{\raggedright\arraybackslash}p{2.5cm}%
                 >{\raggedright\arraybackslash}p{3.3cm}%
                 >{\raggedright\arraybackslash}p{2.6cm}l}
  \toprule
  \textbf{Algoritmo} & \textbf{Pesos} & \textbf{Qué relaja}
    & \textbf{Ciclos negativos} & \textbf{Costo} \\
  \midrule
  BFS & todos iguales & cada arista una vez, por capas & no hay
    & $\Tht{n + m}$ \\
  \addlinespace
  Bellman-Ford & cualesquiera & todas, $n - 1$ veces
    & los detecta & $\Tht{nm}$ \\
  \addlinespace
  \textsc{Caminos-DAG} & cualesquiera, grafo acíclico
    & cada arista una vez, en orden topológico & no hay
    & $\Tht{n + m}$ \\
  \addlinespace
  Dijkstra & $\ge 0$ & cada arista una vez, por $d$ creciente
    & no hay & $\Ord{(n + m)\log n}$ \\
  \bottomrule
\end{tabular}
\end{center}

El costo de Dijkstra es con montículo binario; con arreglo es
$\Ord{n^2}$ y con montículo de Fibonacci, $\Ord{n \log n + m}$. Un
sistema de restricciones de diferencia con $p$ incógnitas y $r$
restricciones se resuelve con Bellman-Ford desde el vértice extra
$v_0$, en $\Ord{p^2 + pr}$, o en $\Ord{p + pr}$ sin construir $v_0$.
Las seis propiedades de la relajación, que sostienen todas las
pruebas de esta sección, se demuestran en la sección
\ref{sec:caminos-pruebas}.

\newpage

% =====================================================================
%  Grafos - Demostraciones y problemas de caminos mas cortos
% =====================================================================

\clase{Demostraciones y problemas de caminos más cortos}%
  {Complemento: Cormen, cap.\ 24.5 y problemas}
\label{sec:caminos-pruebas}

\begin{quick_sum}
  Bellman-Ford, el algoritmo para DAG y Dijkstra hacen lo mismo:
  inicializan las estimaciones y relajan aristas. Solo cambian
  cuántas veces relajan cada arista y en qué orden. Por eso su
  corrección se parte en dos capas. Abajo hay seis propiedades de la
  relajación que valen para \emph{cualquier} secuencia de
  relajaciones; arriba, un argumento corto por algoritmo de que su
  orden cumple la hipótesis de alguna de ellas. La sección
  \ref{sec:caminos-fuente} enunció las seis; aquí se demuestran
  completas y en orden de dependencia, junto con el lema auxiliar que
  pide la convergencia y la prueba de que el subgrafo predecesor es
  siempre un árbol. Un ejemplo con relajaciones intercaladas muestra
  la propiedad central, la relajación de caminos, y una tabla resume
  qué propiedad usa cada algoritmo. La segunda mitad resuelve con esas
  mismas herramientas los problemas del capítulo 24 de Cormen: la
  mejora de Yen a Bellman-Ford, las cajas anidadas, el arbitraje entre
  monedas, el ciclo de peso medio mínimo de Karp y los caminos
  bitónicos.
\end{quick_sum}

% =====================================================================

\subsection*{Las reglas del juego}

Todo ocurre en un grafo dirigido $G = (V, E)$ con pesos
$w : E \to \R$ y una fuente $s$. Los algoritmos de la sección
\ref{sec:caminos-fuente} tocan las estimaciones solo de dos maneras.
\textsc{Inicializar-Fuente}$(G, s)$ (Algoritmo
\ref{alg:inicializar-fuente}) pone $d[v] = \infty$ y $\pi[v] =$ nulo
en todos los vértices y luego $d[s] = 0$.
\textsc{Relajar}$(u, v, w)$ (Algoritmo \ref{alg:relajar}) pregunta si
$d[v] > d[u] + w(u, v)$ y, si es así, asigna
$d[v] \gets d[u] + w(u, v)$ y $\pi[v] \gets u$ en el mismo renglón. De
ese pseudocódigo salen tres hechos que las pruebas usan sin parar:
\begin{enumerate}[label=(\roman*)]
  \item \textsc{Relajar} solo baja estimaciones: $d[v]$ nunca sube.
  \item $d[v]$ y $\pi[v]$ cambian juntos. Por eso, si $\pi[v] = u$, la
    última relajación que cambió $d[v]$ fue
    \textsc{Relajar}$(u, v, w)$.
  \item Con $d[u] = \infty$ no pasa nada, porque
    $\infty + w(u, v) = \infty$ y la prueba $d[v] > \infty$ es falsa.
    Por eso $\pi[v]$ solo recibe vértices con $d$ finita, y un vértice
    $v \ne s$ tiene $d[v]$ finita si y solo si $\pi[v]$ no es nulo:
    las dos cosas las asigna la misma relajación.
\end{enumerate}

Salvo la desigualdad del triángulo, que habla solo de $\delta$, todos
los resultados de esta parte suponen lo mismo: \emph{el grafo se
inicializó con \textsc{Inicializar-Fuente}$(G, s)$ y desde entonces
$d$ y $\pi$ solo cambiaron por llamadas a \textsc{Relajar} sobre
aristas de $G$}, en cualquier orden y cuantas veces se quiera. No se
supone nada sobre ese orden, y ahí está la gracia: lo que se pruebe
aquí vale para Bellman-Ford, para el DAG, para Dijkstra y para
cualquier algoritmo que se invente mañana, mientras solo relaje.

De la sección \ref{sec:caminos-fuente} se usa lo siguiente.
$\delta(s, v)$ es el menor peso de una caminata de $s$ a $v$: vale
$\infty$ si no hay ninguna y $-\infty$ si alguna toca un ciclo de peso
negativo, con $a + \infty = \infty$ para $a \ne -\infty$ y
$a + (-\infty) = -\infty$ para $a \ne \infty$. Cuando $\delta(s, v)$ es
finita la alcanza un camino, sin vértices repetidos y con a lo más
$n - 1$ aristas; un \emph{camino más corto} es un camino de peso
$\delta(s, v)$. El lema de subcaminos (Cormen, lema 24.1) dice que todo
subcamino de un camino más corto es un camino más corto entre sus
extremos. El subgrafo predecesor es $G_\pi = (V_\pi, E_\pi)$ con
\[
  V_\pi = \set{v \in V : \pi[v] \ne \text{nulo}} \cup \set{s},
  \qquad
  E_\pi = \set{(\pi[v], v) : v \in V_\pi - \set{s}} ,
\]
y un árbol de caminos más cortos con raíz $s$ es un subgrafo cuyos
vértices son (1) los alcanzables desde $s$, que (2) es un árbol con
raíz $s$ y en el que (3) el único camino de $s$ a cada $v$ es un
camino más corto en $G$.

La figura \ref{fig:caminos-dependencias} es el mapa de lo que sigue.
Cada flecha se lee ``se usa para probar''. En azul, las seis
propiedades; en blanco, los tres auxiliares.

% Mapa de dependencias entre las propiedades y sus auxiliares
\begin{figure}[htbp]
  \centering
  \begin{tikzpicture}[x=1cm, y=1cm,
      bloque/.style={draw=acento, fill=fondoIdea, rounded corners=2pt,
                     font=\scriptsize, align=center, inner sep=3pt,
                     minimum height=0.8cm, minimum width=2.2cm},
      bloqueAux/.style={bloque, draw=gray!70, fill=white}]
    \node[bloque] (tri) at (0, 0) {Desigualdad\\del triángulo};
    \node[bloque] (cota) at (3.3, 0) {Cota superior};
    \node[bloque] (sin) at (6.9, 1.3) {Sin camino};
    \node[bloque] (conv) at (6.9, 0) {Convergencia};
    \node[bloque] (rel) at (10.5, 0) {Relajación\\de caminos};
    \node[bloqueAux] (sub) at (10.5, 1.3) {Subcaminos};
    \node[bloqueAux] (aux) at (3.3, -1.4)
      {Tras relajar,\\$d[v] \le d[u] + w(u, v)$};
    \node[bloqueAux] (pred) at (3.3, -2.8)
      {Arista del predecesor,\\$d[v] \ge d[\pi[v]] + w(\pi[v], v)$};
    \node[bloque] (arbol) at (6.9, -2.8) {$G_\pi$ es un árbol};
    \node[bloque] (sp) at (10.5, -2.8) {Subgrafo\\predecesor};
    \draw[dirigida] (tri) -- (cota);
    \draw[dirigida] (cota) -- (sin);
    \draw[dirigida] (cota) -- (conv);
    \draw[dirigida] (aux) -- (conv);
    \draw[dirigida] (sub) -- (conv);
    \draw[dirigida] (sub) -- (rel);
    \draw[dirigida] (conv) -- (rel);
    \draw[dirigida] (cota.east) to[out=-10, in=100] (arbol.north);
    \draw[dirigida] (pred) -- (arbol);
    \draw[dirigida] (arbol) -- (sp);
    \draw[dirigida] (pred.south) to[out=-25, in=-155] (sp.south);
  \end{tikzpicture}
  \caption{De qué depende cada propiedad de la relajación.}
  \label{fig:caminos-dependencias}
\end{figure}

% =====================================================================

\subsection*{Triángulo, cota superior y sin camino}

\begin{lema}[Desigualdad del triángulo; Proposición~\ref{prop:triangulo}]
  \label{lem:prueba-triangulo}
  Para toda arista $(u, v) \in E$,
  \[
    \delta(s, v) \le \delta(s, u) + w(u, v) .
  \]
\end{lema}

\begin{proof}
  La idea cabe en una frase: ir a $u$ por el mejor camino y luego
  tomar $(u, v)$ es una manera de llegar a $v$, y la mejor manera no
  puede costar más que una cualquiera. Los infinitos obligan a separar
  tres casos según el valor de $\delta(s, u)$.

  \paso{$\delta(s, u) = \infty$.} El lado derecho vale $\infty$ y no
  hay nada que probar.

  \paso{$\delta(s, u)$ finito.} Sea $p$ un camino más corto de $s$ a
  $u$. Si se le pega la arista $(u, v)$, queda una caminata de $s$ a
  $v$ que pesa $\delta(s, u) + w(u, v)$. Puede repetir $v$, y no
  importa: $\delta(s, v)$ es el mínimo sobre todas las caminatas de
  $s$ a $v$ y no puede pasar del peso de esta en particular.

  \paso{$\delta(s, u) = -\infty$.} Hay caminatas de $s$ a $u$ de peso
  tan negativo como se quiera. Pegándoles $(u, v)$ salen caminatas a
  $v$ igual de negativas, salvo la constante $w(u, v)$; entonces
  $\delta(s, v) = -\infty$ y la desigualdad vale con igualdad.
\end{proof}

\begin{ojo}
  La desigualdad es sobre las distancias, no sobre los pesos. Las
  aristas no tienen que cumplir ninguna desigualdad del triángulo: en
  el ejemplo \ref{ej:caminos-relajacion}, más abajo, $w(s, b) = 6$ es
  mayor que $w(s, a) + w(a, b) = 5$. La que sí la cumple es $\delta$,
  en cualquier grafo y con pesos negativos o sin ellos.
\end{ojo}

\begin{lema}[Cota superior; Proposición~\ref{prop:cota-superior}]
  \label{lem:prueba-cota-superior}
  Después de \textsc{Inicializar-Fuente}$(G, s)$ se cumple
  $d[v] \ge \delta(s, v)$ para todo $v \in V$, y la desigualdad se
  mantiene tras cualquier secuencia de relajaciones. Además, en cuanto
  $d[v]$ alcanza $\delta(s, v)$, ya no cambia.
\end{lema}

\begin{proof}
  La primera parte es un invariante y se prueba por inducción sobre el
  número de relajaciones. Cuidado con la variable: no se induce sobre
  los vértices ni sobre las pasadas, sino sobre las llamadas a
  \textsc{Relajar}, una por una.

  \paso{Base.} Recién inicializado, $d[v] = \infty \ge \delta(s, v)$
  para todo $v \ne s$. Para la fuente, $d[s] = 0$ y
  $\delta(s, s) \le 0$, porque la caminata vacía de $s$ a $s$ pesa
  $0$. ($\delta(s, s)$ vale $0$ si ningún ciclo negativo pasa por $s$,
  y $-\infty$ si alguno pasa.)

  \paso{Paso inductivo.} Se llama \textsc{Relajar}$(u, v, w)$ y, por
  hipótesis de inducción, $d[x] \ge \delta(s, x)$ para todo $x$ antes
  de la llamada. La única estimación que puede cambiar es $d[v]$, y si
  cambia,
  \begin{align*}
    d[v] &= d[u] + w(u, v) \\
         &\ge \delta(s, u) + w(u, v)
           \razon{hipótesis de inducción} \\
         &\ge \delta(s, v)
           \razon{Lema \ref{lem:prueba-triangulo}} .
  \end{align*}
  Las demás estimaciones no se mueven, así que el invariante sigue en
  pie.

  \paso{Una vez alcanzada, no cambia.} Si $d[v] = \delta(s, v)$, no
  puede bajar, por el invariante que se acaba de probar, y no puede
  subir, por el hecho (i). Se queda donde está.
\end{proof}

La segunda mitad del lema parece un detalle y es la pieza que más se
usa. Es la que vuelve inofensivas las relajaciones de más: una vez que
un vértice tiene su distancia, nada de lo que pase después se la
quita.

\begin{observacion}
  Hay otra manera de ver la cota superior. Si $d[v]$ es finita, es el
  peso de una caminata de verdad de $s$ a $v$: la que fueron armando
  las relajaciones. Se ve por la misma inducción. $d[s] = 0$ es la
  caminata vacía, y cuando \textsc{Relajar} asigna
  $d[v] \gets d[u] + w(u, v)$, a la caminata que pesaba $d[u]$ se le
  pega $(u, v)$. Como ninguna caminata pesa menos que $\delta(s, v)$,
  $d[v] \ge \delta(s, v)$. Las estimaciones no son adivinanzas; son
  pesos de caminatas que existen.
\end{observacion}

\begin{corolario}[Sin camino; Proposición~\ref{prop:sin-camino}]
  \label{cor:prueba-sin-camino}
  Si no hay camino de $s$ a $v$, entonces
  $d[v] = \delta(s, v) = \infty$ después de inicializar y después de
  cualquier secuencia de relajaciones.
\end{corolario}

\begin{proof}
  Por la cota superior, $\infty = \delta(s, v) \le d[v]$; entonces
  $d[v] = \infty$.
\end{proof}

% =====================================================================

\subsection*{Convergencia y relajación de caminos}

La convergencia necesita saber qué deja una relajación, justo
después de hacerse. Es el lema auxiliar de Cormen (lema 24.13), y es
casi la definición de \textsc{Relajar}.

\begin{lema}[Lo que deja una relajación]
  \label{lem:prueba-tras-relajar}
  Para una arista $(u, v)$ con $u \ne v$, inmediatamente después de
  \textsc{Relajar}$(u, v, w)$ se cumple
  $d[v] \le d[u] + w(u, v)$.
\end{lema}

\begin{proof}
  Si antes de la llamada $d[v] > d[u] + w(u, v)$, la llamada asigna
  justo $d[v] = d[u] + w(u, v)$. Si no, ya se cumplía
  $d[v] \le d[u] + w(u, v)$, y la llamada no toca ni $d[u]$ ni
  $d[v]$. La condición $u \ne v$ solo excluye lazos: con un lazo de
  peso negativo, la asignación baja $d[u]$ junto con $d[v]$. Un camino
  nunca usa lazos, así que la convergencia no pierde nada.
\end{proof}

\begin{lema}[Convergencia; Proposición~\ref{prop:convergencia}]
  \label{lem:prueba-convergencia}
  Sea $s \leadsto u \to v$ un camino más corto de $s$ a $v$. Si
  $d[u] = \delta(s, u)$ en algún momento antes de llamar a
  \textsc{Relajar}$(u, v, w)$, entonces $d[v] = \delta(s, v)$ en todo
  momento después de esa llamada.
\end{lema}

\begin{proof}
  Se juntan tres piezas: el lema \ref{lem:prueba-tras-relajar} da una
  desigualdad, la cota superior da la otra y el lema de subcaminos
  pega las dos.

  Por la cota superior, una vez que $d[u] = \delta(s, u)$ ya no
  cambia; en particular sigue valiendo $\delta(s, u)$ cuando se relaja
  $(u, v)$. Justo después de esa llamada,
  \begin{align*}
    d[v] &\le d[u] + w(u, v)
           \razon{Lema \ref{lem:prueba-tras-relajar}} \\
         &= \delta(s, u) + w(u, v) \\
         &= \delta(s, v) \razon{subcaminos} .
  \end{align*}
  La última igualdad merece una línea. El tramo $s \leadsto u$ es un
  subcamino del camino más corto, así que pesa $\delta(s, u)$; el
  camino completo pesa entonces $\delta(s, u) + w(u, v)$, y por ser
  el más corto ese peso es $\delta(s, v)$. La cota superior da la
  desigualdad contraria, $d[v] \ge \delta(s, v)$. Entonces
  $d[v] = \delta(s, v)$, y por la segunda mitad de la cota superior se
  queda así.
\end{proof}

\begin{lema}[Relajación de caminos;
  Proposición~\ref{prop:relajacion-caminos}]
  \label{lem:prueba-relajacion-caminos}
  Sea $p = \langle v_0, v_1, \dots, v_k \rangle$ un camino más corto
  de $s = v_0$ a $v_k$. Si después de
  \textsc{Inicializar-Fuente}$(G, s)$ ocurre una secuencia de
  relajaciones que incluye, en este orden, a
  $(v_0, v_1), (v_1, v_2), \dots, (v_{k-1}, v_k)$, entonces
  $d[v_k] = \delta(s, v_k)$ al terminar esas relajaciones y en todo
  momento posterior. No importa qué otras relajaciones haya ni cómo
  se intercalen con las de $p$.
\end{lema}

\begin{proof}
  Se prueba por inducción sobre $i$ que, a partir de la relajación de
  $(v_{i-1}, v_i)$, se cumple $d[v_i] = \delta(s, v_i)$. Para $i = 0$
  ``a partir de'' significa desde el principio.

  \paso{Base, $i = 0$.} Por la inicialización, $d[v_0] = d[s] = 0$.
  Falta ver que $\delta(s, s) = 0$, y aquí hay un detalle que Cormen
  da por hecho. Si un ciclo negativo pasara por $s$, darle vueltas
  antes de seguir $p$ abarataría sin límite las caminatas a $v_k$, y
  no habría camino más corto a $v_k$. Como $p$ existe, ningún ciclo
  negativo pasa por $s$, y $\delta(s, s) = 0 = d[s]$. Por la cota
  superior, $d[s]$ ya no cambia.

  \paso{Paso inductivo.} Supóngase que $d[v_{i-1}] = \delta(s, v_{i-1})$
  desde que se relajó $(v_{i-2}, v_{i-1})$, o desde el principio si
  $i = 1$. Las relajaciones intercaladas no rompen esa igualdad, por
  la cota superior. Cuando le toca a $(v_{i-1}, v_i)$, el tramo
  $\langle v_0, \dots, v_{i-1}, v_i \rangle$ es un subcamino de $p$ y,
  por lo tanto, un camino más corto. La convergencia, con
  $u = v_{i-1}$ y $v = v_i$, da $d[v_i] = \delta(s, v_i)$ desde ese
  momento y para siempre.

  Con $i = k$ queda el enunciado.
\end{proof}

\begin{idea}
  La relajación de caminos dice que lo único que importa es una
  \emph{subsecuencia}: que las aristas de un camino más corto
  aparezcan en orden dentro de la secuencia de relajaciones. Lo que
  haya entre ellas no estorba, por la cota superior, y lo que se
  relaje antes de tiempo no cuenta. Un algoritmo no necesita conocer
  el camino más corto; le basta un orden que garantice esa
  subsecuencia para todos los caminos a la vez. Bellman-Ford la
  garantiza a la fuerza, con $n - 1$ pasadas completas; el algoritmo
  para DAG la obtiene gratis del orden topológico.
\end{idea}

\begin{ejemplo}[Relajación de caminos con relajaciones intercaladas]
  \label{ej:caminos-relajacion}
  Un grafo chico, con una arista negativa a propósito para que se vea
  que la propiedad no depende de Dijkstra. $V = \set{s, a, b, c}$ y
  $w(s, a) = 2$, $w(s, b) = 6$, $w(a, b) = 3$, $w(a, c) = 5$,
  $w(b, c) = -1$. Las distancias son $\delta(s, a) = 2$,
  $\delta(s, b) = 5$ y $\delta(s, c) = 4$, y el camino más corto a $c$
  es $p = \langle s, a, b, c \rangle$. Se relajan siete aristas en un
  orden cualquiera que contiene a las de $p$ en orden (pasos 4, 6
  y 7), mezcladas con otras:
  \begin{center}
  \small
  \begin{tabular}{cccccccl}
    \toprule
    \textbf{Paso} & \textbf{Arista} & $d[s]$ & $d[a]$ & $d[b]$
      & $d[c]$ & \textbf{¿Cambia?} & \textbf{Comentario} \\
    \midrule
    0 & & $0$ & $\infty$ & $\infty$ & $\infty$ & & inicialización \\
    1 & $(b, c)$ & $0$ & $\infty$ & $\infty$ & $\infty$ & no
      & $\infty + (-1) = \infty$ \\
    2 & $(s, b)$ & $0$ & $\infty$ & $6$ & $\infty$ & sí
      & $\pi[b] = s$ \\
    3 & $(b, c)$ & $0$ & $\infty$ & $6$ & $5$ & sí
      & $\pi[c] = b$; $5 > \delta(s, c)$ \\
    4 & $(s, a)$ & $0$ & $\mathbf{2}$ & $6$ & $5$ & sí
      & primera de $p$: $a$ converge \\
    5 & $(a, c)$ & $0$ & $2$ & $6$ & $5$ & no & $2 + 5 = 7 > 5$ \\
    6 & $(a, b)$ & $0$ & $2$ & $\mathbf{5}$ & $5$ & sí
      & segunda de $p$: $b$ converge, $\pi[b] = a$ \\
    7 & $(b, c)$ & $0$ & $2$ & $5$ & $\mathbf{4}$ & sí
      & tercera de $p$: $c$ converge \\
    \bottomrule
  \end{tabular}
  \end{center}
  En verde, los vértices con $d[v] = \delta(s, v)$; las flechas
  azules son las aristas $(\pi[v], v)$ de $G_\pi$, y el número junto a
  cada vértice es $d[v]$.
  \begin{center}
  \begin{tikzpicture}[x=1cm, y=1cm]
    % Tras el paso 3
    \begin{scope}
      \node[vertOk] (s) at (0, 0) {$s$};
      \node[vert] (a) at (1.5, 1.05) {$a$};
      \node[vert] (b) at (1.5, -1.05) {$b$};
      \node[vert] (c) at (3, 0) {$c$};
      \draw[dirigida] (s) -- node[peso] {2} (a);
      \draw[dirigidaArbol] (s) -- node[peso] {6} (b);
      \draw[dirigida] (a) -- node[peso] {3} (b);
      \draw[dirigida] (a) -- node[peso] {5} (c);
      \draw[dirigidaArbol] (b) -- node[peso] {$-1$} (c);
      \node[nota, left=1pt] at (s.west) {$0$};
      \node[nota, above=1pt] at (a.north) {$\infty$};
      \node[nota, below=1pt] at (b.south) {$6$};
      \node[nota, right=1pt] at (c.east) {$5$};
      \node[font=\small] at (1.5, -2.1) {tras el paso 3};
    \end{scope}
    % Tras el paso 6
    \begin{scope}[xshift=5cm]
      \node[vertOk] (s) at (0, 0) {$s$};
      \node[vertOk] (a) at (1.5, 1.05) {$a$};
      \node[vertOk] (b) at (1.5, -1.05) {$b$};
      \node[vert] (c) at (3, 0) {$c$};
      \draw[dirigidaArbol] (s) -- node[peso] {2} (a);
      \draw[dirigida] (s) -- node[peso] {6} (b);
      \draw[dirigidaArbol] (a) -- node[peso] {3} (b);
      \draw[dirigida] (a) -- node[peso] {5} (c);
      \draw[dirigidaArbol] (b) -- node[peso] {$-1$} (c);
      \node[nota, left=1pt] at (s.west) {$0$};
      \node[nota, above=1pt] at (a.north) {$2$};
      \node[nota, below=1pt] at (b.south) {$5$};
      \node[nota, right=1pt] at (c.east) {$5$};
      \node[font=\small] at (1.5, -2.1) {tras el paso 6};
    \end{scope}
    % Tras el paso 7
    \begin{scope}[xshift=10cm]
      \node[vertOk] (s) at (0, 0) {$s$};
      \node[vertOk] (a) at (1.5, 1.05) {$a$};
      \node[vertOk] (b) at (1.5, -1.05) {$b$};
      \node[vertOk] (c) at (3, 0) {$c$};
      \draw[dirigidaArbol] (s) -- node[peso] {2} (a);
      \draw[dirigida] (s) -- node[peso] {6} (b);
      \draw[dirigidaArbol] (a) -- node[peso] {3} (b);
      \draw[dirigida] (a) -- node[peso] {5} (c);
      \draw[dirigidaArbol] (b) -- node[peso] {$-1$} (c);
      \node[nota, left=1pt] at (s.west) {$0$};
      \node[nota, above=1pt] at (a.north) {$2$};
      \node[nota, below=1pt] at (b.south) {$5$};
      \node[nota, right=1pt] at (c.east) {$4$};
      \node[font=\small] at (1.5, -2.1) {tras el paso 7};
    \end{scope}
  \end{tikzpicture}
  \end{center}
  Tras el paso 3, $d[c] = 5$ es el peso de un camino real,
  $s \to b \to c$, pero no del más corto: la cota superior se cumple y
  nada más. Tras el paso 6, $b$ ya convergió, pero $c$ no se entera
  hasta que se vuelve a relajar $(b, c)$; la relajación del paso 3 se
  hizo cuando $d[b]$ todavía no era correcta y no cuenta. En el paso 7
  se completa la subsecuencia $(s, a), (a, b), (b, c)$ y $c$ queda en
  $4$ para siempre. Los pasos 1, 3 y 5 sobran, pero tampoco estorban.
\end{ejemplo}

El mismo grafo explica por qué Bellman-Ford hace $n - 1$ pasadas y
por qué al DAG le basta una. Si cada pasada relaja las aristas en el
orden $(b, c)$, $(a, b)$, $(s, a)$, $(a, c)$, $(s, b)$, que recorre
$p$ al revés, cada pasada avanza una sola arista de $p$:
\begin{center}
\small
\begin{tabular}{ccccc}
  \toprule
  \textbf{Tras la pasada} & $d[s]$ & $d[a]$ & $d[b]$ & $d[c]$ \\
  \midrule
  1 & $0$ & $\mathbf{2}$ & $6$ & $7$ \\
  2 & $0$ & $2$ & $\mathbf{5}$ & $5$ \\
  3 & $0$ & $2$ & $5$ & $\mathbf{4}$ \\
  \bottomrule
\end{tabular}
\end{center}
Son $3 = n - 1$ pasadas, el peor caso. El grafo no tiene ciclos y
$s, a, b, c$ es un orden topológico. \textsc{Caminos-DAG} relaja las
aristas agrupadas por origen en ese orden, $(s, a)$, $(s, b)$,
$(a, b)$, $(a, c)$, $(b, c)$, y termina en una pasada con
$d = (0, 2, 5, 4)$. No es suerte: en un orden topológico cada arista
de cualquier camino sale de un vértice que va antes que el origen de
la siguiente, así que las aristas de \emph{todo} camino se relajan en
orden. Aquí las de $p$ caen en las posiciones 1, 3 y 5.

% =====================================================================

\subsection*{El subgrafo predecesor}

Falta la otra mitad de la historia: los apuntadores $\pi$. Las dos
pruebas que siguen se apoyan en una desigualdad que vale en todo
momento para las aristas de $G_\pi$, la misma que en Cormen aparece
dentro de la prueba del lema 24.16.

\begin{lema}[Arista del predecesor]\label{lem:prueba-arista-pred}
  En todo momento, para cada vértice $v$ con $\pi[v]$ no nulo,
  $d[\pi[v]]$ es finita y
  \[
    d[v] \ge d[\pi[v]] + w(\pi[v], v) .
  \]
\end{lema}

\begin{proof}
  Sea $u = \pi[v]$. Por el hecho (ii), la última relajación que cambió
  $d[v]$ fue \textsc{Relajar}$(u, v, w)$, y dejó
  $d[v] = d[u] + w(u, v)$ con el valor que $d[u]$ tenía entonces, que
  era finito por el hecho (iii). Desde ese momento $d[v]$ no ha
  cambiado y $d[u]$ solo pudo bajar, por el hecho (i). El lado derecho
  solo pudo bajar y la desigualdad se mantiene.
\end{proof}

Con estimaciones correctas la desigualdad se vuelve igualdad sobre las
aristas del árbol; es lo que dice la última línea de la prueba del
lema \ref{lem:prueba-subgrafo-pred}.

\begin{lema}[$G_\pi$ es un árbol]\label{lem:prueba-gpi-arbol}
  Supóngase que $G$ no tiene ciclos de peso negativo alcanzables desde
  $s$. Después de \textsc{Inicializar-Fuente}$(G, s)$, el subgrafo
  predecesor $G_\pi$ es un árbol con raíz $s$, y cualquier secuencia
  de relajaciones mantiene esa propiedad.
\end{lema}

\begin{proof}
  Al inicio $V_\pi = \set{s}$ y $E_\pi = \emptyset$, un árbol de un
  solo vértice. Hay que ver dos cosas después de cada relajación: que
  $G_\pi$ no tiene ciclos y que de $s$ hay un único camino a cada
  vértice de $V_\pi$.

  \paso{$G_\pi$ no tiene ciclos.} Supóngase, buscando una
  contradicción, que alguna relajación cierra un ciclo
  $c = \langle v_0, v_1, \dots, v_k \rangle$ en $G_\pi$, con
  $v_k = v_0$ y $\pi[v_i] = v_{i-1}$ para $i = 1, \dots, k$.
  Numerando el ciclo desde donde convenga, la llamada que lo cierra es
  \textsc{Relajar}$(v_{k-1}, v_k, w)$. Se miran las estimaciones justo
  antes de esa llamada.

  Para $i = 1, \dots, k - 1$ el apuntador $\pi[v_i] = v_{i-1}$ ya
  estaba puesto, y el lema \ref{lem:prueba-arista-pred} da
  \[
    d[v_i] \ge d[v_{i-1}] + w(v_{i-1}, v_i) ,
    \qquad i = 1, \dots, k - 1 .
  \]
  En la arista que cierra el ciclo la desigualdad es estricta, porque
  la llamada sí cambia $\pi[v_k]$, y eso solo ocurre si pasa la prueba
  de \textsc{Relajar}:
  \[
    d[v_k] > d[v_{k-1}] + w(v_{k-1}, v_k) .
  \]
  Las $k$ estimaciones del ciclo son finitas: $d[v_{k-1}]$ porque la
  prueba pasó, y las otras porque cada $v_{i-1}$ con
  $1 \le i \le k - 1$ es el predecesor de $v_i$. Sumando las $k$
  desigualdades,
  \[
    \sum_{i=1}^{k} d[v_i] > \sum_{i=1}^{k} d[v_{i-1}]
      + \sum_{i=1}^{k} w(v_{i-1}, v_i) .
  \]
  Como $v_0 = v_k$, las dos sumas de estimaciones recorren los mismos
  vértices, una vez cada uno, y son el mismo número finito. Se cancelan
  y queda
  \[
    0 > \sum_{i=1}^{k} w(v_{i-1}, v_i) = w(c) .
  \]
  El ciclo $c$ es negativo. Además es alcanzable desde $s$, porque sus
  vértices tienen $d$ finita y, por la cota superior,
  $\delta(s, v_i) \le d[v_i] < \infty$. Eso contradice la hipótesis.

  \paso{Hay un único camino de $s$ a cada vértice de $V_\pi$.}
  Primero, $\pi[s]$ sigue siendo nulo: cambiarlo exigiría bajar $d[s]$
  de $0$, y la cota superior lo impide, porque sin ciclos negativos
  alcanzables $\delta(s, s) = 0$. Ahora tómese $v \in V_\pi$ y sígase
  la cadena $v, \pi[v], \pi[\pi[v]], \dots$. Por el lema
  \ref{lem:prueba-arista-pred} cada eslabón tiene $d$ finita, así que
  es $s$ o tiene $\pi$ no nulo (hecho (iii)): la cadena no sale de
  $V_\pi$. Tampoco repite vértices, porque $G_\pi$ no tiene ciclos.
  Entonces se detiene, y solo puede detenerse en un vértice de $V_\pi$
  con $\pi$ nulo, o sea en $s$. Leída al revés, es un camino de $s$ a
  $v$ en $G_\pi$.

  Es el único. A cada vértice $z \ne s$ de $V_\pi$ le entra en $G_\pi$
  exactamente una arista, $(\pi[z], z)$, y a $s$ ninguna. Un camino que
  sale de $s$ y llega a $v$ tiene que llegar desde $\pi[v]$, que a su
  vez tiene que venir de $\pi[\pi[v]]$, y así hasta $s$: está forzado a
  ser la cadena. Un grafo dirigido con un único camino de $s$ a cada
  vértice es un árbol con raíz $s$.
\end{proof}

\begin{lema}[Subgrafo predecesor; Proposición~\ref{prop:subgrafo-pred}]
  \label{lem:prueba-subgrafo-pred}
  Supóngase que $G$ no tiene ciclos negativos alcanzables desde $s$.
  Si después de \textsc{Inicializar-Fuente}$(G, s)$ una secuencia de
  relajaciones deja $d[v] = \delta(s, v)$ para todo $v \in V$,
  entonces $G_\pi$ es un árbol de caminos más cortos con raíz $s$.
\end{lema}

\begin{proof}
  Hay que comprobar las tres condiciones de la definición.

  \paso{(1) $V_\pi$ son los alcanzables.} Sin ciclos negativos
  alcanzables, $\delta(s, v)$ es finita si y solo si $v$ es alcanzable
  desde $s$. Por el hecho (iii), $V_\pi$ es justo el conjunto de
  vértices con $d$ finita ($s$ incluido, con $d[s] = 0$). Como
  $d[v] = \delta(s, v)$, $V_\pi = \set{v : \delta(s, v) < \infty}$,
  que son los alcanzables.

  \paso{(2) Es un árbol con raíz $s$.} Es el lema
  \ref{lem:prueba-gpi-arbol}.

  \paso{(3) Sus caminos son caminos más cortos.} Sea $v \in V_\pi$ y
  sea $p = \langle v_0, \dots, v_k \rangle$ el camino de $s = v_0$ a
  $v = v_k$ en $G_\pi$. Para cada $i$, $\pi[v_i] = v_{i-1}$, y el lema
  \ref{lem:prueba-arista-pred} con $d = \delta$ dice
  \[
    w(v_{i-1}, v_i) \le \delta(s, v_i) - \delta(s, v_{i-1}) .
  \]
  Al sumar sobre el camino, el lado derecho se telescopa:
  \begin{align*}
    w(p) = \sum_{i=1}^{k} w(v_{i-1}, v_i)
      &\le \sum_{i=1}^{k}
           \bigl( \delta(s, v_i) - \delta(s, v_{i-1}) \bigr) \\
      &= \delta(s, v_k) - \delta(s, v_0)
         \razon{suma telescópica} \\
      &= \delta(s, v) \razon{$\delta(s, s) = 0$} .
  \end{align*}
  Ningún camino de $s$ a $v$ pesa menos que $\delta(s, v)$, así que
  $w(p) = \delta(s, v)$ y $p$ es un camino más corto. De paso, todas
  las desigualdades de la suma eran igualdades: sobre las aristas del
  árbol el triángulo se cumple justo,
  $\delta(s, v_i) = \delta(s, v_{i-1}) + w(v_{i-1}, v_i)$.
\end{proof}

\begin{ojo}
  Las hipótesis de los lemas \ref{lem:prueba-gpi-arbol} y
  \ref{lem:prueba-subgrafo-pred} son distintas, y las dos hacen falta.
  \begin{itemize}
    \item Sin ciclos negativos, $G_\pi$ es árbol en todo momento, pero
      a media corrida no tiene por qué ser de caminos más cortos. Tras
      el paso 3 del ejemplo \ref{ej:caminos-relajacion} el árbol
      contiene $s \to b \to c$, que no es el camino más corto ni a $b$
      ni a $c$. Solo cuando todas las $d$ son correctas el árbol es el
      bueno.
    \item Con un ciclo negativo alcanzable, $G_\pi$ puede dejar de ser
      árbol: en el ejemplo \ref{ej:caminos-arbitraje} los apuntadores
      forman un ciclo, y es justo el ciclo negativo. La prueba de
      aciclicidad solo usó la hipótesis al final, para descartar el
      ciclo. Leída sin ella dice algo más fuerte: \emph{todo ciclo que
      aparezca en $G_\pi$ es negativo}. Eso es lo que permite recuperar
      un ciclo negativo siguiendo apuntadores.
  \end{itemize}
\end{ojo}

% =====================================================================

\subsection*{Quién usa qué}

Con las seis propiedades probadas, la corrección de cada algoritmo se
reduce a decir cuál de ellas aplica y por qué su orden de relajación
cumple la hipótesis. La tabla \ref{tab:caminos-quien-usa} lo resume
para los tres algoritmos de la sección \ref{sec:caminos-fuente} y para
los dos de los problemas de esta sección.

\begin{table}[htbp]
  \centering
  \small
  \begin{tabular}{>{\raggedright\arraybackslash}p{2.6cm}
                  >{\raggedright\arraybackslash}p{3.2cm}
                  >{\raggedright\arraybackslash}p{3.7cm}
                  >{\raggedright\arraybackslash}p{4.9cm}}
    \toprule
    \textbf{Algoritmo} & \textbf{Orden de relajación}
      & \textbf{Propiedades} & \textbf{Dónde entran} \\
    \midrule
    \textsc{Bellman-Ford} (Alg.~\ref{alg:bellman-ford})
      & $n - 1$ pasadas sobre todas las aristas
      & relajación de caminos, sin camino, triángulo, subgrafo
        predecesor
      & la pasada $i$ relaja la $i$-ésima arista de todo camino más
        corto; la revisión final pasa por el triángulo y falla con un
        ciclo negativo por la suma alrededor del ciclo \\
    \addlinespace
    \textsc{Caminos-DAG} (Alg.~\ref{alg:caminos-dag})
      & una pasada, orígenes en orden topológico
      & relajación de caminos, sin camino, subgrafo predecesor
      & el orden topológico pone en orden las aristas de cualquier
        camino \\
    \addlinespace
    \textsc{Dijkstra} (Alg.~\ref{alg:dijkstra-cormen})
      & cada arista una vez, orígenes en orden de $d$ creciente
      & cota superior, sin camino, convergencia, subgrafo predecesor
      & el primer vértice no fijado de un camino más corto ya tiene su
        distancia, por convergencia \\
    \addlinespace
    Yen (Alg.~\ref{alg:caminos-yen})
      & $\lceil n/2 \rceil$ pasadas de ida y regreso
      & relajación de caminos
      & cada pasada despacha dos corridas del camino \\
    \addlinespace
    Bitónicos (Alg.~\ref{alg:caminos-bitonicos})
      & cuatro pasadas en orden de peso
      & relajación de caminos
      & cada pasada despacha una corrida monótona \\
    \bottomrule
  \end{tabular}
  \caption{Qué propiedad de la relajación sostiene a cada algoritmo.}
  \label{tab:caminos-quien-usa}
\end{table}

\textbf{Bellman-Ford.} Sin ciclos negativos alcanzables, cada vértice
alcanzable $v$ tiene un camino más corto
$p = \langle v_0, \dots, v_k \rangle$ con $k \le n - 1$. La pasada $i$
relaja todas las aristas, entre ellas $(v_{i-1}, v_i)$; las de $p$
quedan relajadas en orden y la relajación de caminos da
$d[v] = \delta(s, v)$. Los no alcanzables se quedan en $\infty$ por
la propiedad sin camino, y el árbol sale del subgrafo predecesor. La
revisión final no encuentra nada que relajar porque
$d[v] = \delta(s, v) \le \delta(s, u) + w(u, v) = d[u] + w(u, v)$,
por el triángulo. Con un ciclo negativo alcanzable
$c = \langle v_0, \dots, v_k \rangle$, si ninguna de sus aristas
fallara la revisión, sumando $d[v_i] \le d[v_{i-1}] + w(v_{i-1}, v_i)$
alrededor de $c$ las estimaciones se cancelarían (son finitas, porque
tras $n - 1$ pasadas todo alcanzable tiene $d$ finita) y quedaría
$0 \le w(c)$. Es la misma suma de la prueba del lema
\ref{lem:prueba-gpi-arbol}, con la desigualdad volteada.

\textbf{Caminos-DAG.} En un orden topológico cada $v_{i-1}$ de un
camino va antes que $v_i$, así que al procesar los orígenes en ese
orden las aristas de cualquier camino se relajan en orden, cada una
exactamente una vez. Es el orden ``bueno'' del ejemplo
\ref{ej:caminos-relajacion}.

\textbf{Dijkstra.} Aquí la herramienta es la convergencia, porque el
orden no se conoce de antemano. Si $u$ fuera el primer vértice que se
fija con $d[u] \ne \delta(s, u)$, no sería $s$, y sería alcanzable
(si no, la propiedad sin camino lo dejaría correcto en $\infty$). Sea
$y$ el primer vértice no fijado de un camino más corto de $s$ a $u$ y
$x$ su predecesor en ese camino, ya fijado. Cuando $x$ se fijó tenía
$d[x] = \delta(s, x)$, porque $u$ es el primero que falla, y en ese
momento se relajó $(x, y)$: por convergencia, $d[y] = \delta(s, y)$.
Con pesos no negativos,
$\delta(s, y) \le \delta(s, u) \le d[u] \le d[y]$ (el tramo de $y$ a
$u$ no resta, cota superior, y $u$ se eligió con la $d$ mínima), así
que todo es igualdad y $d[u] = \delta(s, u)$, una contradicción. La
traza de Dijkstra de la clase del 22 de octubre ilustra además la
relajación de caminos: las aristas del camino A, D, E, F se relajan en
los pasos 1, 4 y 5, en orden, y F termina en $12$.

\subsubsection*{Ejercicios}

\begin{ejercicio}\label{ej:caminos-pi-fuente}
  Probar que si una secuencia de relajaciones deja $\pi[s]$ distinto
  de nulo, entonces un ciclo de peso negativo pasa por $s$. (Cormen,
  ejercicio 24.5-4.)

  \emph{Solución.} Para cambiar $\pi[s]$, una relajación tuvo que
  bajar $d[s]$ de su valor inicial $0$, así que $d[s] < 0$. Por la
  cota superior, $\delta(s, s) \le d[s] < 0$. Pero $\delta(s, s)$ solo
  puede valer $0$ o $-\infty$, y vale $-\infty$: un ciclo negativo pasa
  por $s$. Con la observación que sigue a la cota superior se ve
  directo: $d[s]$ es el peso de una caminata cerrada de $s$ a $s$, y
  una caminata cerrada de peso negativo contiene un ciclo negativo.
\end{ejercicio}

\begin{ejercicio}
  Con pesos no negativos, dar un grafo y una asignación de
  predecesores en la que cada $\pi[v]$ es el predecesor de $v$ en
  \emph{algún} camino más corto, pero $G_\pi$ tiene un ciclo. ¿Por qué
  las relajaciones nunca producen eso? (Cormen, ejercicio 24.5-5.)

  \emph{Solución.} $w(s, u) = w(s, v) = 1$ y $w(u, v) = w(v, u) = 0$.
  Tanto $s, u, v$ como $s, v, u$ son caminos más cortos (pesan $1$),
  así que $\pi[u] = v$ y $\pi[v] = u$ es una asignación legal en el
  sentido del enunciado, y cierra el ciclo $u, v, u$. Las relajaciones
  no lo permiten porque \textsc{Relajar} exige mejora estricta: el
  ciclo pesa $0$, y la prueba del lema \ref{lem:prueba-gpi-arbol}
  muestra que cerrar un ciclo en $G_\pi$ con relajaciones exige peso
  negativo.
\end{ejercicio}

\begin{ejercicio}
  Sin ciclos negativos alcanzables, probar que existe una secuencia de
  $n - 1$ relajaciones que deja $d[v] = \delta(s, v)$ para todo $v$.
  (Cormen, ejercicio 24.5-7.)

  \emph{Solución.} Tómese un árbol de caminos más cortos $T$ con raíz
  $s$ (existe: es el $G_\pi$ que deja Bellman-Ford, por el lema
  \ref{lem:prueba-subgrafo-pred}) y relájense sus aristas en un orden
  en que cada arista $(u, v)$ venga después de la que entra a $u$; por
  ejemplo, por niveles de BFS sobre $T$. Para cada alcanzable $v$, el
  camino de $s$ a $v$ en $T$ es un camino más corto y sus aristas se
  relajan en orden: relajación de caminos. Los no alcanzables ya
  valen $\infty$. $T$ tiene a lo más $n - 1$ aristas; si tiene menos,
  se completa con relajaciones cualesquiera, que no estropean nada.
\end{ejercicio}

\begin{ejercicio}
  Con un ciclo negativo alcanzable desde $s$, construir una secuencia
  infinita de relajaciones en la que cada una cambia alguna
  estimación. (Cormen, ejercicio 24.5-8.)

  \emph{Solución.} Primero se relajan en orden las aristas de un
  camino de $s$ a un vértice del ciclo $c$ y luego las del ciclo,
  omitiendo las que no cambien nada (si una no cambia, su destino ya
  tenía $d$ finita). Así todos los vértices de $c$ quedan con $d$
  finita. Desde ahí, si todas las aristas del ciclo cumplieran
  $d[v_i] \le d[v_{i-1}] + w(v_{i-1}, v_i)$, sumando alrededor de $c$
  saldría $0 \le w(c)$, que es falso. Siempre queda una arista del
  ciclo que se puede relajar; se relaja esa y se repite, para siempre.
  En el ejemplo \ref{ej:caminos-arbitraje} cada vuelta al ciclo baja
  las tres estimaciones en $0.0475$.
\end{ejercicio}

% =====================================================================

\subsection*{Problema: la mejora de Yen a Bellman-Ford}

\textbf{Planteamiento} (Cormen, problema 24-1). Antes de la primera
pasada se fija un orden lineal arbitrario $v_1, v_2, \dots, v_n$ de
los vértices y se parten las aristas en dos conjuntos:
\[
  E_f = \set{(v_i, v_j) \in E : i < j} , \qquad
  E_b = \set{(v_i, v_j) \in E : i > j} ,
\]
las que van hacia adelante en el orden y las que van hacia atrás. Se
supone que no hay lazos, para que toda arista caiga en uno de los dos.
Cada pasada recorre los vértices en el orden $v_1, \dots, v_n$
relajando sus aristas de $E_f$, y luego en el orden
$v_n, \dots, v_1$ relajando sus aristas de $E_b$. Yen mostró que así
bastan $\lceil n/2 \rceil$ pasadas en lugar de $n - 1$.

\textbf{Dos DAG escondidos.} $G_f = (V, E_f)$ no tiene ciclos y
$v_1, \dots, v_n$ es un orden topológico suyo: toda arista sube de
índice, y un ciclo tendría que bajar alguna vez para regresar.
Simétricamente, $G_b = (V, E_b)$ es acíclico con orden topológico
$v_n, \dots, v_1$. Cada mitad de pasada es exactamente
\textsc{Caminos-DAG} sobre uno de los dos DAG; Yen alterna los dos.

\textbf{Por qué bastan $\lceil n/2 \rceil$ pasadas.} Sea $v$
alcanzable y $p = \langle u_0, u_1, \dots, u_k \rangle$ un camino más
corto de $s = u_0$ a $v = u_k$, con $k \le n - 1$. Cada arista de $p$
es de ida o de regreso. Se parte $p$ en \emph{corridas}: bloques
máximos de aristas consecutivas del mismo tipo,
$R_1, R_2, \dots, R_r$, que alternan entre $E_f$ y $E_b$, con
$r \le k$. A lo largo de una corrida de $E_f$ los orígenes suben de
índice, y en la mitad de ida los vértices se procesan en índice
creciente; entonces la corrida entera se relaja en el orden del camino
dentro de una sola mitad de ida. Lo mismo pasa con las corridas de
$E_b$ en la mitad de regreso. Cada pasada despacha, entonces, una
corrida de $E_f$ y la de $E_b$ que le sigue. Si $p$ empieza con una
corrida de $E_b$, la primera mitad de ida se desperdicia. En el peor
caso hacen falta $\lceil (r + 1)/2 \rceil$ pasadas, y como
$r \le n - 1$, eso es a lo más $\lceil n/2 \rceil$. Las aristas de
$p$ quedaron relajadas en orden, y la relajación de caminos da
$d[v] = \delta(s, v)$.

\begin{algorithm}[H]
\caption{\textsc{Bellman-Ford-Yen}$(G, w, s)$}
\label{alg:caminos-yen}
\begin{algorithmic}[1]
\Procedure{Bellman-Ford-Yen}{$G, w, s$}
  \State fijar un orden lineal $v_1, v_2, \dots, v_n$ de $V$
  \State $E_f \gets \set{(v_i, v_j) \in E : i < j}$, \quad
         $E_b \gets \set{(v_i, v_j) \in E : i > j}$
  \State \Call{Inicializar-Fuente}{$G, s$}
  \For{$p \gets 1$ \Hasta $\lceil n / 2 \rceil$}
    \For{$i \gets 1, 2, \dots, n$} \Comment{mitad de ida}
      \ForAll{$(v_i, v_j) \in E_f$}
        \State \Call{Relajar}{$v_i, v_j, w$}
      \EndFor
    \EndFor
    \For{$i \gets n, n - 1, \dots, 1$} \Comment{mitad de regreso}
      \ForAll{$(v_i, v_j) \in E_b$}
        \State \Call{Relajar}{$v_i, v_j, w$}
      \EndFor
    \EndFor
  \EndFor
  \ForAll{$(u, v) \in E$}
    \If{$d[v] > d[u] + w(u, v)$}
      \State \Return falso \Comment{ciclo negativo alcanzable}
    \EndIf
  \EndFor
  \State \Return verdadero
\EndProcedure
\end{algorithmic}
\end{algorithm}

\textbf{Costo.} Las líneas 6 a 15 tocan cada arista exactamente una
vez por pasada: cada una está en $E_f$ o en $E_b$ y se visita desde su
origen. Con listas de adyacencia separadas para $E_f$ y $E_b$, que la
línea 3 arma una vez en $\Tht{n + m}$, una pasada cuesta
$\Tht{n + m}$. En total: $\Tht{n + m}$ de preparación,
$\lceil n/2 \rceil \cdot \Tht{n + m}$ de pasadas y $\Tht{m}$ de
revisión, o sea $\Tht{nm}$ cuando todos los vértices son alcanzables
($m \ge n - 1$). El orden de crecimiento no mejora; lo que se gana es
la mitad de las pasadas. La revisión de las líneas 17 a 21 sigue
siendo válida tras $\lceil n/2 \rceil$ pasadas: el argumento de
Bellman-Ford solo necesita que los vértices de un ciclo negativo
alcanzable tengan $d$ finita, y la tienen, porque el argumento de las
corridas vale para cualquier camino desde $s$, no solo para los más
cortos.

\begin{ejemplo}[Las pasadas de Yen]
  \label{ej:caminos-yen}
  Seis vértices en el orden lineal $v_1, \dots, v_6$ y el camino
  $v_1 \to v_4 \to v_2 \to v_5 \to v_3 \to v_6$, con todas las aristas
  de peso $1$ y fuente $s = v_1$. Arriba, en azul, las aristas de
  $E_f$; abajo, en naranja, las de $E_b$. Junto a cada arista va la
  pasada y la mitad en que se relaja con su origen ya correcto, y
  debajo de cada vértice su distancia final:
  \begin{center}
  \begin{tikzpicture}[x=1cm, y=1cm]
    \foreach \i in {1, ..., 6}
      \node[vert] (v\i) at ({1.8*(\i - 1)}, 0) {$v_\i$};
    \draw[dirigidaArbol] (v1) to[out=55, in=125, looseness=0.8]
      node[peso, pos=0.3] {1, ida} (v4);
    \draw[dirigidaArbol] (v2) to[out=70, in=110, looseness=1.0]
      node[peso, pos=0.5] {2, ida} (v5);
    \draw[dirigidaArbol] (v3) to[out=55, in=125, looseness=0.8]
      node[peso, pos=0.7] {3, ida} (v6);
    \draw[dirigidaResalte] (v4) to[out=-125, in=-55, looseness=0.9]
      node[peso, pos=0.6] {1, regreso} (v2);
    \draw[dirigidaResalte] (v5) to[out=-125, in=-55, looseness=0.9]
      node[peso, pos=0.4] {2, regreso} (v3);
    \foreach \i/\dist in {1/0, 2/2, 3/4, 4/1, 5/3, 6/5}
      \node[nota] at ({1.8*(\i - 1)}, -1.9) {$d = \dist$};
  \end{tikzpicture}
  \end{center}
  El camino tiene cinco corridas, tres de ida y dos de regreso. La
  pasada 1 resuelve $v_4$ en su mitad de ida y $v_2$ en la de regreso;
  la pasada 2, $v_5$ y $v_3$; la pasada 3, $v_6$. Son
  $3 = \lceil 6/2 \rceil$ pasadas. Bellman-Ford con las aristas en el
  orden inverso del camino necesita $5 = n - 1$.
\end{ejemplo}

\begin{ojo}
  Un lazo $(v, v)$ no está ni en $E_f$ ni en $E_b$. Si su peso es
  negativo es un ciclo negativo que el esquema no vería nunca, así que
  los lazos se revisan aparte; los de peso no negativo se pueden
  tirar, porque ningún camino más corto los usa.
\end{ojo}

% =====================================================================

\subsection*{Problema: cajas anidadas}

\textbf{Planteamiento} (Cormen, problema 24-2). Una caja de $k$
dimensiones $(x_1, \dots, x_k)$ \emph{anida} en otra
$(y_1, \dots, y_k)$ si existe una permutación $\sigma$ de
$\set{1, \dots, k}$ con $x_{\sigma(i)} < y_i$ para todo $i$: se puede
girar la primera para que quepa holgada dentro de la segunda. Cormen
llama $d$ a la dimensión; aquí es $k$, porque $d$ ya es la estimación.
Dadas $n$ cajas, se busca la sucesión más larga
$B_{i_1}, B_{i_2}, \dots$ en la que cada caja anida en la siguiente.

\textbf{Anidar es transitivo.} Si $x_{\sigma(i)} < y_i$ para todo $i$
y $y_{\tau(j)} < z_j$ para todo $j$, entonces
$x_{\sigma(\tau(j))} < y_{\tau(j)} < z_j$: la permutación compuesta
sirve de testigo.

\textbf{Cómo probar si anida sin recorrer permutaciones.} Se ordenan
las dimensiones de las dos cajas. Escribiendo $x_{(i)}$ para la
$i$-ésima más chica,
\[
  x \text{ anida en } y \iff x_{(i)} < y_{(i)}
  \text{ para todo } i .
\]
De derecha a izquierda, la permutación que ordena es el testigo. De
izquierda a derecha, las $i$ dimensiones más chicas de $y$ están
emparejadas por $\sigma$ con $i$ dimensiones distintas de $x$, cada
una menor que su pareja y por lo tanto menor que $y_{(i)}$; hay al
menos $i$ valores de $x$ por debajo de $y_{(i)}$, y entonces
$x_{(i)} < y_{(i)}$. Cuesta $\Ord{k \log k}$ ordenar y $\Ord{k}$
comparar. El criterio se comparó contra la fuerza bruta sobre todas
las permutaciones en 3000 pares aleatorios con $k \le 4$.

\textbf{La sucesión más larga es un camino más largo en un DAG.} Un
vértice por caja y una arista $B_i \to B_j$ si $B_i$ anida en $B_j$.
El grafo no tiene ciclos: un ciclo daría, por transitividad, una caja
que anida en sí misma, y eso es imposible con desigualdades
estrictas. La sucesión pedida es un camino con el mayor número de
vértices. Se resuelve con \textsc{Caminos-DAG} poniendo peso $-1$ a
cada arista, o directamente con la recurrencia
\[
  L(j) = 1 + \max \set{L(i) : B_i \to B_j} ,
\]
con $L(j) = 1$ si a $B_j$ no le entra ninguna arista, evaluada en orden
topológico. Hay un orden topológico gratis: ordenar las cajas por la
suma de sus dimensiones, porque si $B_i$ anida en $B_j$ su suma es
estrictamente menor. La sucesión se reconstruye guardando en cada caja
de dónde vino su máximo, como un $\pi$.

\textbf{Costo.} Ordenar las dimensiones de cada caja cuesta
$\Ord{n k \log k}$; comparar los $n(n - 1)$ pares ordenados,
$\Ord{n^2 k}$. El DAG tiene hasta $\Tht{n^2}$ aristas y el camino más
largo cuesta $\Ord{n^2}$. Total: $\Ord{n k \log k + n^2 k}$.

\begin{ejemplo}[Seis cajas de tres dimensiones]
  \label{ej:caminos-cajas}
  Las cajas, sus dimensiones ordenadas, la suma y el resultado de la
  recurrencia en el orden de las sumas:
  \begin{center}
  \small
  \begin{tabular}{cccccc}
    \toprule
    \textbf{Caja} & \textbf{Dimensiones} & \textbf{Ordenadas}
      & \textbf{Suma} & $L$ & \textbf{Viene de} \\
    \midrule
    $B_1$ & $(5, 1, 3)$ & $(1, 3, 5)$ & $9$ & $1$ & \\
    $B_2$ & $(2, 6, 4)$ & $(2, 4, 6)$ & $12$ & $2$ & $B_1$ \\
    $B_5$ & $(2, 2, 9)$ & $(2, 2, 9)$ & $13$ & $1$ & \\
    $B_3$ & $(7, 3, 5)$ & $(3, 5, 7)$ & $15$ & $3$ & $B_2$ \\
    $B_6$ & $(3, 7, 8)$ & $(3, 7, 8)$ & $18$ & $3$ & $B_2$ \\
    $B_4$ & $(6, 9, 4)$ & $(4, 6, 9)$ & $19$ & $\mathbf{4}$ & $B_3$ \\
    \bottomrule
  \end{tabular}
  \end{center}
  Comparando las columnas ordenadas salen ocho aristas. En azul, la
  sucesión más larga:
  \begin{center}
  \begin{tikzpicture}[x=1cm, y=1cm]
    \node[vert] (b1) at (0, 0) {$B_1$};
    \node[vert] (b2) at (2.4, 0) {$B_2$};
    \node[vert] (b3) at (4.8, 0.9) {$B_3$};
    \node[vert] (b6) at (4.8, -0.9) {$B_6$};
    \node[vert] (b4) at (7.2, 0) {$B_4$};
    \node[vert] (b5) at (9.6, 0) {$B_5$};
    \draw[dirigida] (b1) to[bend left=28] (b3);
    \draw[dirigida] (b1) to[bend right=28] (b6);
    \draw[dirigida] (b1) to[out=45, in=135, looseness=0.85] (b4);
    \draw[dirigida] (b2) -- (b6);
    \draw[dirigida] (b2) -- (b4);
    \draw[dirigidaArbol] (b1) -- (b2);
    \draw[dirigidaArbol] (b2) -- (b3);
    \draw[dirigidaArbol] (b3) -- (b4);
    \node[nota] at (9.6, -0.6) {no anida con nadie};
  \end{tikzpicture}
  \end{center}
  La respuesta es $B_1, B_2, B_3, B_4$, de longitud $4$. Ninguno de
  sus tres pasos se ve comparando coordenada a coordenada sin ordenar:
  $(5, 1, 3)$ contra $(2, 6, 4)$ falla en la primera, y sin embargo
  $B_1$ cabe en $B_2$ girándola. $B_5$ muestra que la suma solo sirve
  de orden: su suma queda entre las de $B_2$ y $B_3$, pero su
  dimensión $9$ le impide entrar en cualquier otra, y ninguna cabe en
  ella.
\end{ejemplo}

\begin{ojo}
  La desigualdad tiene que ser estricta. Dos cajas iguales no anidan;
  si se usara $\le$, cada par de cajas iguales formaría un ciclo y
  ``camino más largo'' dejaría de tener sentido.
\end{ojo}

% =====================================================================

\subsection*{Problema: arbitraje}

\textbf{Planteamiento} (Cormen, problema 24-3). \emph{Arbitraje} es
aprovechar tipos de cambio inconsistentes para convertir una unidad
de una moneda en más de una unidad de la misma moneda. Hay $n$ monedas
$c_1, \dots, c_n$ y una tabla $R$ en la que una unidad de $c_i$
compra $R[i, j]$ unidades de $c_j$. Se pide decidir si existe un ciclo
de monedas con
\[
  R[i_1, i_2] \cdot R[i_2, i_3] \cdots R[i_{k-1}, i_k] \cdot
  R[i_k, i_1] > 1
\]
y, si existe, imprimirlo.

\textbf{De productos a sumas.} El logaritmo convierte productos en
sumas y es creciente. Con $i_{k+1} = i_1$,
\[
  \prod_{j=1}^{k} R[i_j, i_{j+1}] > 1
  \iff \sum_{j=1}^{k} \ln R[i_j, i_{j+1}] > 0
  \iff \sum_{j=1}^{k} \bigl( -\ln R[i_j, i_{j+1}] \bigr) < 0 .
\]
Con pesos $w(c_i, c_j) = -\ln R[i, j]$, hay arbitraje si y solo si
hay un ciclo de peso negativo, y Bellman-Ford lo detecta. Un tipo de
cambio mayor que $1$ da un peso negativo, y eso es lo normal aquí:
Dijkstra no sirve. Si $R[i, j] = 0$ no hay arista, porque $\ln 0$ no
existe. Para que todo ciclo sea alcanzable aunque la tabla tenga
ceros, se agrega una fuente virtual $s$ con aristas de peso $0$ a
todas las monedas; equivale a inicializar todas las estimaciones en
$0$. El grafo queda con $n + 1$ vértices: $n$ es el número de monedas,
no de vértices.

El algoritmo completo queda así; la recuperación del ciclo, en las
líneas 12 a 21, se justifica después.

\begin{algorithm}[H]
\caption{\textsc{Arbitraje}$(R, n)$}
\label{alg:caminos-arbitraje}
\begin{algorithmic}[1]
\Procedure{Arbitraje}{$R, n$}
  \State $V \gets \set{s, c_1, \dots, c_n}$; aristas $(s, c_i)$ de
         peso $0$
  \State aristas $(c_i, c_j)$ de peso $-\ln R[i, j]$ para cada
         $i \ne j$ con $R[i, j] > 0$
  \State \Call{Inicializar-Fuente}{$G, s$}
  \For{$p \gets 1$ \Hasta $n$} \Comment{$\abs{V} - 1 = n$ pasadas}
    \ForAll{$(u, v) \in E$}
      \State \Call{Relajar}{$u, v, w$}
    \EndFor
  \EndFor
  \ForAll{$(u, v) \in E$}
    \If{$d[v] > d[u] + w(u, v) + \varepsilon$}
      \State \Call{Relajar}{$u, v, w$}
        \Comment{ahora $G_\pi$ tiene un ciclo}
      \State $x \gets v$
      \For{$i \gets 1$ \Hasta $n + 1$}
        \State $x \gets \pi[x]$ \Comment{al final, $x$ está en el ciclo}
      \EndFor
      \State $C \gets \langle x \rangle$, \quad $y \gets \pi[x]$
      \While{$y \neq x$}
        \State anteponer $y$ a $C$, \quad $y \gets \pi[y]$
      \EndWhile
      \State anteponer $x$ a $C$
      \State \Return $C$ \Comment{$c_{i_1}, \dots, c_{i_k}, c_{i_1}$}
    \EndIf
  \EndFor
  \State \Return no hay arbitraje
\EndProcedure
\end{algorithmic}
\end{algorithm}

\textbf{Cómo recuperar el ciclo.} Si tras las $n$ pasadas
($\abs{V} - 1 = n$) la arista $(u, v)$ todavía se puede relajar, se
relaja y se siguen los apuntadores desde $v$ durante $n + 1$ pasos. El
vértice al que se llega está en un ciclo de $G_\pi$, y ese ciclo es
negativo. La razón tiene tres partes.
\begin{itemize}
  \item Tras $i$ pasadas, cada $d[x]$ es a lo más el peso de cualquier
    caminata de $s$ a $x$ con a lo más $i$ aristas. Se ve por
    inducción sobre $i$: si la caminata termina en $(y, x)$, tras
    $i - 1$ pasadas $d[y]$ ya es a lo más el peso del resto, y en la
    pasada $i$ se relaja $(y, x)$ (lema \ref{lem:prueba-tras-relajar}).
    Entonces la relajación extra bajó $d[v]$ estrictamente por debajo
    del peso de toda caminata de $s$ a $v$ con a lo más $n$ aristas.
  \item Si la cadena $v, \pi[v], \pi[\pi[v]], \dots$ llegara a $s$
    sin repetir vértices, sería un camino de a lo más $n$ aristas.
    Sumando el lema \ref{lem:prueba-arista-pred} a lo largo de ella
    saldría $d[v] \ge d[s] + w(\text{cadena}) = w(\text{cadena})$,
    porque a la fuente virtual no le entra ninguna arista y $d[s]$ se
    queda en $0$. Eso contradice el punto anterior. Tampoco puede
    detenerse en otro vértice, porque todos los de la cadena tienen
    $d$ finita y $\pi$ no nulo. Entonces la cadena cae en un ciclo, y
    como solo hay $n + 1$ vértices, tras $n + 1$ pasos ya está dentro
    de él.
  \item Ese ciclo es negativo, por la versión fuerte del lema
    \ref{lem:prueba-gpi-arbol}: todo ciclo que aparece en $G_\pi$ lo
    es. Leído al derecho, es la sucesión de monedas.
\end{itemize}


Las líneas que no son obvias. En la línea 11, el $\varepsilon$ (por
ejemplo $10^{-12}$ relativo) evita que el redondeo de los logaritmos
invente un arbitraje donde el producto es exactamente $1$. La línea 12
relaja la arista antes de seguir apuntadores; es lo que garantiza el
ciclo en $G_\pi$. Las líneas 14 a 16 dan $n + 1 = \abs{V}$ saltos para
caer seguro en el ciclo y no en un vértice que cuelga de él. En las
líneas 17 a 21 se antepone porque los apuntadores van hacia atrás.

\textbf{Costo.} El grafo tiene $n + 1$ vértices y, con la tabla
llena, $n(n - 1) + n = n^2$ aristas. Construirlo (líneas 2 y 3) cuesta
$\Tht{n^2}$;
las $n$ pasadas, $n \cdot \Tht{n^2} = \Tht{n^3}$; la revisión,
$\Tht{n^2}$, y la recuperación, $\Tht{n}$. Total: $\Tht{n^3}$. Floyd
sobre los mismos pesos también cuesta $\Tht{n^3}$ y delata el ciclo
como una entrada negativa en la diagonal; es más fácil de programar,
aunque recuperar el ciclo pide guardar los intermedios.

\begin{ejemplo}[Arbitraje entre dólar, rupia y yen]
  \label{ej:caminos-arbitraje}
  Los tres tipos de cambio de Cormen ($1$ dólar compra $49$ rupias,
  $1$ rupia compra $2$ yenes y $1$ yen compra $0.0107$ dólares), más
  los de regreso, con el margen que cobra una casa de cambio. $R[i, j]$
  es cuántas unidades de la moneda $j$ se compran con una de la $i$:
  \[
    R = \begin{array}{c|ccc}
          & D & I & Y \\ \hline
        D & 1 & 49 & 92 \\
        I & 0.0200 & 1 & 2 \\
        Y & 0.0107 & 0.48 & 1
      \end{array}
    \qquad\qquad
    w = -\ln R = \begin{array}{c|ccc}
          & D & I & Y \\ \hline
        D & 0 & -3.892 & -4.522 \\
        I & 3.912 & 0 & -0.693 \\
        Y & 4.538 & 0.734 & 0
      \end{array}
  \]
  Un ciclo da ganancia si el producto de sus tipos pasa de $1$, o sea,
  si la suma de sus pesos es negativa. Hay cinco ciclos:
  \begin{center}
  \small
  \begin{tabular}{lcc}
    \toprule
    \textbf{Ciclo} & \textbf{Producto de $R$}
      & \textbf{Suma de $w$} \\
    \midrule
    $D \to I \to D$ & $0.9800$ & $0.0202$ \\
    $D \to Y \to D$ & $0.9844$ & $0.0157$ \\
    $I \to Y \to I$ & $0.9600$ & $0.0408$ \\
    $D \to I \to Y \to D$ & $1.0486$ & $-0.0475$ \\
    $D \to Y \to I \to D$ & $0.8832$ & $0.1242$ \\
    \bottomrule
  \end{tabular}
  \end{center}
  Solo $D \to I \to Y \to D$ gana: $49 \cdot 2 \cdot 0.0107 = 1.0486$,
  un $4.86\%$, y su peso es
  $-3.892 - 0.693 + 4.538 = -0.0475 = -\ln 1.0486$.
  \begin{center}
  \begin{tikzpicture}[x=1cm, y=1cm]
    \begin{scope}
      \node[vert] (D) at (0, 0) {$D$};
      \node[vert] (I) at (5, 0) {$I$};
      \node[vert] (Y) at (2.5, 3.5) {$Y$};
      \draw[dirigida] (D) to[bend left=18]
        node[peso, auto] {49} (I);
      \draw[dirigida] (I) to[bend left=18]
        node[peso, auto] {0.0200} (D);
      \draw[dirigida] (I) to[bend left=18]
        node[peso, auto, pos=0.42] {2} (Y);
      \draw[dirigida] (Y) to[bend left=18]
        node[peso, auto] {0.48} (I);
      \draw[dirigida] (Y) to[bend left=18]
        node[peso, auto, pos=0.58] {0.0107} (D);
      \draw[dirigida] (D) to[bend left=18]
        node[peso, auto] {92} (Y);
      \node[font=\small] at (2.5, -1.3) {tipos de cambio $R[i, j]$};
    \end{scope}
    \begin{scope}[xshift=7.6cm]
      \node[vert] (D) at (0, 0) {$D$};
      \node[vert] (I) at (5, 0) {$I$};
      \node[vert] (Y) at (2.5, 3.5) {$Y$};
      \draw[dirigidaResalte] (D) to[bend left=18]
        node[peso, auto] {$-3.892$} (I);
      \draw[dirigida] (I) to[bend left=18]
        node[peso, auto] {$3.912$} (D);
      \draw[dirigidaResalte] (I) to[bend left=18]
        node[peso, auto, pos=0.42] {$-0.693$} (Y);
      \draw[dirigida] (Y) to[bend left=18]
        node[peso, auto] {$0.734$} (I);
      \draw[dirigidaResalte] (Y) to[bend left=18]
        node[peso, auto, pos=0.58] {$4.538$} (D);
      \draw[dirigida] (D) to[bend left=18]
        node[peso, auto] {$-4.522$} (Y);
      \node[font=\small] at (2.5, -1.3)
        {pesos $-\ln R[i, j]$; ciclo de peso $-0.0475$};
    \end{scope}
  \end{tikzpicture}
  \end{center}
  Para seguirlo a mano se corre Bellman-Ford desde $D$ en lugar de la
  fuente virtual: la tabla está llena y $D$ alcanza a las otras dos.
  Son $\abs{V} - 1 = 2$ pasadas, con las aristas en el orden $(D, I)$,
  $(D, Y)$, $(I, D)$, $(I, Y)$, $(Y, D)$, $(Y, I)$. La columna
  \emph{candidato} es $d[u] + w(u, v)$:
  \begin{center}
  \small
  \begin{tabular}{cccccccl}
    \toprule
    \textbf{Pasada} & \textbf{Arista} & \textbf{Candidato}
      & \textbf{¿Cambia?} & $d[D]$ & $d[I]$ & $d[Y]$ & \\
    \midrule
    1 & $(D, I)$ & $-3.8918$ & sí & $0$ & $-3.8918$ & $\infty$
      & $\pi[I] = D$ \\
      & $(D, Y)$ & $-4.5218$ & sí & $0$ & $-3.8918$ & $-4.5218$
      & $\pi[Y] = D$ \\
      & $(I, D)$ & $0.0202$ & no & $0$ & $-3.8918$ & $-4.5218$ & \\
      & $(I, Y)$ & $-4.5850$ & sí & $0$ & $-3.8918$ & $-4.5850$
      & $\pi[Y] = I$ \\
      & $(Y, D)$ & $-0.0475$ & sí & $-0.0475$ & $-3.8918$
      & $-4.5850$ & $\pi[D] = Y$ \\
      & $(Y, I)$ & $-3.8510$ & no & $-0.0475$ & $-3.8918$
      & $-4.5850$ & \\
    \midrule
    2 & las seis & & 3 de 6 & $-0.0949$ & $-3.9393$ & $-4.6324$
      & bajan $D$, $I$ y $Y$ \\
    \midrule
    revisión & $(D, I)$ & $-3.9867$ & falla & & & &
      $-3.9393 > -3.9867$ \\
    \bottomrule
  \end{tabular}
  \end{center}
  Ya en la pasada 1, al relajar $(Y, D)$, la fuente recibe
  $\pi[D] = Y$ y $G_\pi$ queda con el ciclo $D \to I \to Y \to D$: es
  el ejercicio \ref{ej:caminos-pi-fuente} en acción, y el lema
  \ref{lem:prueba-gpi-arbol} deja de valer porque su hipótesis falla.
  Cada vuelta baja las tres estimaciones en $0.0475$ y nunca se
  estabilizan. La revisión encuentra que $(D, I)$ todavía se puede
  relajar: hay un ciclo negativo. Se relaja ($\pi[I]$ ya era $D$) y se
  retrocede $\abs{V} = 3$ veces desde $I$: $D$, $Y$, $I$. Desde $I$ se
  recorre $\pi$ hasta volver: $\pi[I] = D$, $\pi[D] = Y$,
  $\pi[Y] = I$. Al derecho es $I \to Y \to D \to I$, el mismo ciclo
  $D \to I \to Y \to D$ empezando en otra moneda. Con la fuente
  virtual y tres pasadas falla la misma arista y sale el mismo ciclo.
\end{ejemplo}

\begin{ojo}
  Los logaritmos se redondean. Un ciclo con producto exactamente $1$
  puede sumar $-10^{-16}$ y parecer arbitraje; por eso la línea 11
  compara con tolerancia. Y para recuperar el ciclo hay que relajar
  primero la arista que falla y luego retroceder $\abs{V}$ veces:
  retroceder menos puede dejar en un vértice que cuelga del ciclo sin
  estar en él.
\end{ojo}

% =====================================================================

\subsection*{Problema: ciclo de peso medio mínimo (Karp)}

\textbf{Planteamiento} (Cormen, problema 24-5). El \emph{peso medio}
de un ciclo $c$ con $\ell$ aristas es $\mu(c) = w(c)/\ell$, y
$\mu^* = \min_c \mu(c)$ sobre todos los ciclos dirigidos de $G$. Un
ciclo con $\mu(c) = \mu^*$ es un \emph{ciclo de peso medio mínimo}.
Se pide calcular $\mu^*$ en $\Ord{nm}$. Se supone que todo vértice es
alcanzable desde una fuente $s$, y se define
\[
  \delta_k(s, v) = \text{peso mínimo de una caminata de } s
  \text{ a } v \text{ con exactamente } k \text{ aristas} ,
\]
$\infty$ si no hay ninguna. Aquí sí importa que sean caminatas: pueden
repetir vértices, y con exactamente $n$ aristas siempre repiten.
La tabla de $\delta_k$ se llena por niveles, cada uno a partir del
anterior:
$\delta_k(s, v) = \min_{(u, v) \in E} \bigl( \delta_{k-1}(s, u) +
w(u, v) \bigr)$, con $\delta_0(s, s) = 0$ y $\delta_0(s, v) = \infty$
para $v \ne s$.

\begin{teorema}[Karp]\label{teo:caminos-karp}
  Si $G$ tiene al menos un ciclo y todo vértice es alcanzable desde
  $s$, entonces
  \[
    \mu^* = \min_{v \in V} \; \max_{0 \le k \le n - 1}
      \frac{\delta_n(s, v) - \delta_k(s, v)}{n - k} ,
  \]
  donde el mínimo omite los $v$ con $\delta_n(s, v) = \infty$ y el
  máximo omite los $k$ con $\delta_k(s, v) = \infty$.
\end{teorema}

\begin{proof}
  Sea $F(v)$ el máximo de adentro. La prueba sigue los incisos de
  Cormen.

  \paso{Paso 1: basta el caso $\mu^* = 0$.} Si se suma una constante
  $t$ al peso de cada arista, la media de todo ciclo sube $t$, así que
  $\mu^*$ sube $t$. Del otro lado, toda caminata de exactamente $k$
  aristas sube $kt$, así que $\delta_k(s, v)$ sube $kt$ y cada cociente
  sube $\frac{nt - kt}{n - k} = t$. Los dos lados se mueven igual; con
  $t = -\mu^*$ el caso general se reduce a $\mu^* = 0$. De aquí en
  adelante, $\mu^* = 0$.

  \paso{Paso 2: no hay ciclos negativos.} Un ciclo negativo tendría
  media negativa, menor que $\mu^*$. Entonces los caminos más cortos
  existen, no repiten vértices y tienen a lo más $n - 1$ aristas:
  \[
    \delta(s, v) = \min_{0 \le k \le n - 1} \delta_k(s, v) ,
  \]
  finito para todo $v$ porque todos son alcanzables.

  \paso{Paso 3: $F(v) \ge 0$.} Sea $v$ con $\delta_n(s, v)$ finita. Una
  caminata de $n$ aristas sigue siendo una caminata, así que
  $\delta_n(s, v) \ge \delta(s, v) = \delta_{k^*}(s, v)$ para el $k^*$
  que alcanza el mínimo del paso 2. El cociente de ese $k^*$ es
  $\ge 0$, y el máximo también.

  \paso{Paso 4: algún $v$ tiene $F(v) = 0$.} Sea $c$ un ciclo de peso
  $0$, que existe porque $\mu^* = 0$, y sea $u$ un vértice de $c$.
  Primero una observación: si el tramo de $c$ que va de $u$ a otro de
  sus vértices $v$ pesa $x$, entonces $\delta(s, v) = \delta(s, u) + x$.
  El resto del ciclo, de $v$ a $u$, pesa $-x$; pegar un camino más
  corto con un tramo del ciclo da una caminata, de donde
  $\delta(s, v) \le \delta(s, u) + x$ y
  $\delta(s, u) \le \delta(s, v) - x$, y las dos juntas dan la
  igualdad.

  Ahora se toma un camino más corto de $s$ a $u$, de $j \le n - 1$
  aristas, y se alarga dando vueltas a $c$ desde $u$ hasta completar
  exactamente $n$ aristas. Termina en algún vértice $v$ de $c$, y lo
  recorrido sobre $c$ pesa lo mismo que el tramo directo de $u$ a $v$,
  porque cada vuelta completa pesa $0$. Por la observación, esta
  caminata de $n$ aristas pesa $\delta(s, u) + x = \delta(s, v)$.
  Entonces $\delta_n(s, v) \le \delta(s, v)$ y, con el paso 3,
  $\delta_n(s, v) = \delta(s, v) \le \delta_k(s, v)$ para todo $k$.
  Todos los cocientes de $v$ son $\le 0$, así que $F(v) \le 0$, y con
  el paso 3, $F(v) = 0$.

  \paso{Paso 5.} Por el paso 3 el mínimo de los $F(v)$ es $\ge 0$, y
  por el paso 4 vale $0$, que es $\mu^*$. Los $v$ omitidos no cambian
  nada: con $\delta_n(s, v) = \infty$ su máximo sería $\infty$ y nunca
  ganaría el mínimo.
\end{proof}

\begin{ejemplo}[Ciclo de peso medio mínimo con Karp]
  \label{ej:caminos-karp}
  El caso interesante es cuando el ciclo de media mínima no es el de
  peso total mínimo. $V = \set{a, b, c, d}$, $n = 4$, fuente $s = a$
  y aristas $w(a, b) = 0$, $w(a, c) = 2$, $w(b, c) = 1$,
  $w(c, b) = 1$, $w(c, d) = 0$, $w(d, a) = 2$. Tiene tres ciclos:
  \begin{center}
  \begin{tikzpicture}[x=1cm, y=1cm]
    \node[vert] (a) at (0, 0) {$a$};
    \node[vert] (b) at (2.2, 1.3) {$b$};
    \node[vert] (c) at (4.4, 0) {$c$};
    \node[vert] (d) at (2.2, -1.3) {$d$};
    \draw[dirigidaResalte] (a) -- node[peso] {0} (b);
    \draw[dirigida] (a) -- node[peso] {2} (c);
    \draw[dirigidaResalte] (b) to[bend left=14] node[peso] {1} (c);
    \draw[dirigida] (c) to[bend left=14] node[peso] {1} (b);
    \draw[dirigidaResalte] (c) -- node[peso] {0} (d);
    \draw[dirigidaResalte] (d) -- node[peso] {2} (a);
    \node[nota, anchor=west] at (5.4, 0.6)
      {$b, c, b$: peso $2$, media $1$};
    \node[nota, anchor=west] at (5.4, 0)
      {$a, c, d, a$: peso $4$, media $4/3$};
    \node[nota, anchor=west] at (5.4, -0.6)
      {$a, b, c, d, a$: peso $3$, media $3/4$};
  \end{tikzpicture}
  \end{center}
  El ciclo más ligero es $b, c, b$, pero el de media mínima es el
  largo, en naranja: $\mu^* = 3/4$. La tabla de $\delta_k(a, v)$:
  \begin{center}
  \small
  \begin{tabular}{ccccc}
    \toprule
    $k$ & $a$ & $b$ & $c$ & $d$ \\
    \midrule
    0 & $0$ & $\infty$ & $\infty$ & $\infty$ \\
    1 & $\infty$ & $0$ & $2$ & $\infty$ \\
    2 & $\infty$ & $3$ & $1$ & $2$ \\
    3 & $4$ & $2$ & $4$ & $1$ \\
    4 & $3$ & $4$ & $3$ & $4$ \\
    \bottomrule
  \end{tabular}
  \end{center}
  Hay $\infty$ en renglones donde el vértice sí es alcanzable, pero no
  con ese número exacto de aristas: $a$ no se alcanza con $1$ ni con
  $2$. Por ejemplo, $\delta_2(a, b) = 3$ es $a \to c \to b$ y
  $\delta_4(a, a) = 3$ es la vuelta completa $a, b, c, d, a$. Para cada
  $v$ se evalúa $(\delta_4(a, v) - \delta_k(a, v)) / (4 - k)$ en los
  $k$ con $\delta_k(a, v)$ finita:
  \begin{center}
  \small
  \begin{tabular}{clc}
    \toprule
    $v$ & \textbf{Cocientes} & \textbf{Máximo} \\
    \midrule
    $a$ & $k = 0$: $\frac{3 - 0}{4} = \frac{3}{4}$; \
          $k = 3$: $\frac{3 - 4}{1} = -1$ & $\mathbf{3/4}$ \\[3pt]
    $b$ & $k = 1$: $\frac{4 - 0}{3} = \frac{4}{3}$; \
          $k = 2$: $\frac{4 - 3}{2} = \frac{1}{2}$; \
          $k = 3$: $\frac{4 - 2}{1} = 2$ & $2$ \\[3pt]
    $c$ & $k = 1$: $\frac{3 - 2}{3} = \frac{1}{3}$; \
          $k = 2$: $\frac{3 - 1}{2} = 1$; \
          $k = 3$: $\frac{3 - 4}{1} = -1$ & $1$ \\[3pt]
    $d$ & $k = 2$: $\frac{4 - 2}{2} = 1$; \
          $k = 3$: $\frac{4 - 1}{1} = 3$ & $3$ \\
    \bottomrule
  \end{tabular}
  \end{center}
  El mínimo de los máximos es $3/4$, en $v = a$, que está en el ciclo
  de media mínima, como predice el paso 4. Con los pesos corridos
  $w - 3/4$ del paso 1, la tabla da
  $\delta_4(a, a) = 0 = \delta_0(a, a)$: el vértice $a$ alcanza su
  distancia con exactamente $n$ aristas, como en la prueba.
\end{ejemplo}

El algoritmo llena esa tabla nivel por nivel y evalúa la fórmula del
teorema \ref{teo:caminos-karp}.

\begin{algorithm}[H]
\caption{\textsc{Karp}$(G, w, s)$}
\label{alg:caminos-karp}
\begin{algorithmic}[1]
\Procedure{Karp}{$G, w, s$}
  \State $D[k][v] \gets \infty$ para todo $0 \le k \le n$ y todo
         $v \in V$
  \State $D[0][s] \gets 0$
  \For{$k \gets 1$ \Hasta $n$}
    \ForAll{$(u, v) \in E$}
      \If{$D[k-1][u] + w(u, v) < D[k][v]$}
        \State $D[k][v] \gets D[k-1][u] + w(u, v)$
      \EndIf
    \EndFor
  \EndFor
  \State $\mu \gets \infty$
  \ForAll{$v \in V$ con $D[n][v] < \infty$}
    \State $M \gets -\infty$
    \For{$k \gets 0$ \Hasta $n - 1$}
      \If{$D[k][v] < \infty$}
        \State $M \gets \max\bigl(M,\ (D[n][v] - D[k][v]) / (n - k)
               \bigr)$
      \EndIf
    \EndFor
    \State $\mu \gets \min(\mu, M)$
  \EndFor
  \State \Return $\mu$ \Comment{$\infty$ si no hay ciclos}
\EndProcedure
\end{algorithmic}
\end{algorithm}

\textbf{Costo.} Las líneas 2 y 3 llenan $(n + 1) \cdot n$ casillas,
$\Tht{n^2}$. Las líneas 4 a 10 son $n$ niveles, y cada uno recorre las
$m$ aristas una vez: $\Tht{nm}$. Las líneas 11 a 20 evalúan $n$
cocientes para cada uno de los $n$ vértices: $\Tht{n^2}$. Como todos
son alcanzables desde $s$, $m \ge n - 1$ y $n^2 = \Ord{nm}$; el total
es $\Tht{nm}$. La tabla ocupa $\Tht{n^2}$ de memoria. En la línea 6
se lee $D[k-1]$ y se escribe $D[k]$: si se actualizara en el lugar,
sería Bellman-Ford y se perdería el ``exactamente $k$ aristas'' que
la fórmula necesita. Si ninguna fuente alcanza todo, se inicializa
$D[0][v] = 0$ para todo $v$, que equivale a una fuente virtual con
aristas de peso $0$; con esa inicialización la fórmula coincidió con
la fuerza bruta en 2000 grafos aleatorios.

\begin{ojo}
  Tres trampas al programarlo. La tabla no se actualiza en el lugar.
  Los $\infty$ se saltan: si $\delta_n(s, v) = \infty$ se omite el
  vértice entero, y si $\delta_k(s, v) = \infty$ se omite ese
  cociente. Y $\delta_k$ es sobre caminatas: restringida a caminos,
  que tienen a lo más $n - 1$ aristas, $\delta_n$ sería $\infty$ en
  todos lados y la fórmula no diría nada.
\end{ojo}

Por el paso 1, $\mu^*$ es el mayor $t$ tal que los pesos $w - t$ no
tienen ciclos negativos. Una búsqueda binaria sobre $t$ con
Bellman-Ford también lo encuentra, aunque más lento; Karp lo saca de
una sola tabla.

% =====================================================================

\subsection*{Problema: caminos bitónicos}

\textbf{Planteamiento} (Cormen, problema 24-6). Una sucesión es
\emph{bitónica} si crece y luego decrece, o si se vuelve así con un
corrimiento circular. $\langle 1, 4, 6, 8, 3, -2 \rangle$,
$\langle 9, 2, -4, -10, -5 \rangle$ y $\langle 1, 2, 3, 4 \rangle$
son bitónicas; $\langle 1, 3, 12, 4, 2, 10 \rangle$ no. Se da un grafo
con pesos todos distintos y la promesa de que, para cada $v$, los pesos
a lo largo de \emph{cualquier} camino más corto de $s$ a $v$ forman
una sucesión bitónica. Se piden los caminos más cortos lo más rápido
posible.

\begin{lema}\label{lem:caminos-bitonica-corridas}
  Una sucesión bitónica se parte en a lo más tres bloques
  consecutivos, cada uno monótono, con direcciones alternadas: sube,
  baja, sube, o baja, sube, baja (alguno puede faltar).
\end{lema}

\begin{proof}
  Algún corrimiento circular $r = \langle r_1, \dots, r_\ell \rangle$
  de la sucesión crece hasta un máximo $r_q$ y luego decrece. Sean
  $U = \langle r_1, \dots, r_q \rangle$, que sube, y
  $D = \langle r_{q+1}, \dots, r_\ell \rangle$, que baja. La sucesión
  original es $r$ cortada en algún punto $j$ y con las dos partes
  intercambiadas, $\langle r_{j+1}, \dots, r_\ell, r_1, \dots, r_j
  \rangle$. Si el corte cae dentro de $U$ ($j \le q$), se lee como el
  final de $U$ (sube), todo $D$ (baja) y el principio de $U$ (sube).
  Si cae dentro de $D$ ($j > q$), se lee como el final de $D$ (baja),
  todo $U$ (sube) y el principio de $D$ (baja).
\end{proof}

\textbf{Idea.} Se ordenan las aristas por peso y se hacen cuatro
pasadas de relajación: en orden creciente de peso, decreciente,
creciente y decreciente. Sea $p$ un camino más corto de $s$ a $v$. Un
bloque creciente de $p$ tiene pesos estrictamente crecientes, porque
son distintos; una pasada en orden creciente lo relaja entonces en el
orden del camino. Lo mismo con los bloques decrecientes y las pasadas
decrecientes. Por el lema \ref{lem:caminos-bitonica-corridas}, el
patrón de $p$ es sube-baja-sube, y lo cubren las pasadas 1, 2 y 3, o
baja-sube-baja, y lo cubren las pasadas 2, 3 y 4. En los dos casos las
aristas de $p$ quedan relajadas en orden, y la relajación de caminos
da $d[v] = \delta(s, v)$.

\begin{algorithm}[H]
\caption{\textsc{Caminos-Bitónicos}$(G, w, s)$}
\label{alg:caminos-bitonicos}
\begin{algorithmic}[1]
\Procedure{Caminos-Bitónicos}{$G, w, s$}
  \State $L \gets$ las aristas de $E$ ordenadas por peso creciente
  \State \Call{Inicializar-Fuente}{$G, s$}
  \ForAll{sentido en (creciente, decreciente, creciente,
          decreciente)}
    \ForAll{$(u, v)$ en $L$, recorrida en ese sentido}
      \State \Call{Relajar}{$u, v, w$}
    \EndFor
  \EndFor
\EndProcedure
\end{algorithmic}
\end{algorithm}

\textbf{Costo.} Ordenar cuesta $\Ord{m \log m} = \Ord{m \log n}$,
porque $m \le n^2$; inicializar, $\Tht{n}$, y las cuatro pasadas,
$\Tht{m}$. Total: $\Ord{n + m \log n}$, dominado por el
ordenamiento. Si los pesos son enteros en un rango chico y se ordenan
en tiempo lineal, queda $\Tht{n + m}$. En 1500 grafos aleatorios que
cumplen la promesa, las cuatro pasadas dieron siempre las distancias
de Bellman-Ford.

\begin{ejemplo}[Por qué hacen falta cuatro pasadas]
  \label{ej:caminos-bitonicos}
  El grafo más simple que lo muestra es un solo camino. Sus pesos,
  $\langle 3, 1, 2, 5, 4 \rangle$, forman una sucesión bitónica: con
  el corrimiento $\langle 1, 2, 5, 4, 3 \rangle$ sube y luego baja.
  Como es el único camino, también es el más corto, y la promesa se
  cumple. Junto a cada arista va la pasada en que se relaja con su
  origen ya correcto, y debajo de cada vértice su distancia:
  \begin{center}
  \begin{tikzpicture}[x=1cm, y=1cm]
    \foreach \v/\x in {s/0, a/1.9, b/3.8, c/5.7, d/7.6, t/9.5}
      \node[vert] (\v) at (\x, 0) {$\v$};
    \foreach \u/\v/\w/\p in {s/a/3/1, a/b/1/2, b/c/2/3, c/d/5/3,
                             d/t/4/4}
      \draw[dirigida] (\u) -- node[peso] {\w}
        node[nota, below=6pt] {pasada \p} (\v);
    \foreach \v/\dist in {s/0, a/3, b/4, c/6, d/11, t/15}
      \node[nota] at ([yshift=-1.0cm]\v) {$d = \dist$};
  \end{tikzpicture}
  \end{center}
  Tras las pasadas creciente y decreciente solo $s$, $a$ y $b$ tienen
  su distancia; la tercera pasada resuelve $c$ y $d$, y $t$ sigue en
  $\infty$ hasta la cuarta. El patrón es baja-sube-baja: el bloque
  $\langle 3 \rangle$, el bloque $\langle 1, 2, 5 \rangle$ y el bloque
  $\langle 4 \rangle$.
\end{ejemplo}

\begin{ojo}
  La solución que suele circular hace dos pasadas, creciente y luego
  decreciente, y solo cubre la forma sube-baja sin corrimiento. El
  enunciado de Cormen admite corrimientos circulares, y con ellos dos
  pasadas no bastan: en el ejemplo dejan mal a $c$, $d$ y $t$, y en
  las pruebas aleatorias fallaron en 162 de 1500 grafos. Tres pasadas
  tampoco: ningún orden de tres cubre a la vez sube-baja-sube y
  baja-sube-baja, y creciente, decreciente, creciente falló en 2. Los
  pesos distintos también son necesarios: con empates, dos aristas
  consecutivas del mismo peso pueden relajarse en el orden equivocado
  dentro de una pasada.
\end{ojo}

% Macros propias de la seccion 08. \paresCambio pinta en azul y
% negritas la casilla de una matriz que cambio en la iteracion;
% \paresNulo escribe el predecesor vacio en versalitas chicas para que
% las matrices Pi no se ensanchen.
\newcommand{\paresCambio}[1]{\textcolor{acento}{\boldsymbol{#1}}}
\newcommand{\paresNulo}{\text{\scriptsize\textsc{nulo}}}

\newpage

% =====================================================================
%  Grafos - Caminos mas cortos entre todos los pares
% =====================================================================

\clase{Caminos más cortos entre todos los pares}{Complemento: Cormen,
cap.\ 25}
\label{sec:todos-pares}

\begin{quick_sum}
  El problema es calcular $\delta(u, v)$ para todo par de vértices de
  un grafo dirigido con pesos, que pueden ser negativos siempre que no
  haya ciclos negativos, y entregar la respuesta como dos matrices:
  $D$ con los pesos y $\Pi$ con los predecesores que permiten
  reconstruir cada camino. Correr $n$ veces un algoritmo de fuente
  única es el punto de partida y la vara con que se mide lo demás.
  Siguen tres ideas. Una programación dinámica por número de aristas
  que resulta ser un producto de matrices con $\min$ en lugar de suma
  y $+$ en lugar de producto: $\Tht{n^4}$ tal cual y
  $\Tht{n^3 \log n}$ elevando al cuadrado. Floyd-Warshall, que cambia
  el estado a qué vértices se permiten como intermedios y baja a
  $\Tht{n^3}$; con él vienen $\Pi$, la versión en el lugar, la
  detección de ciclos negativos en la diagonal y la cerradura
  transitiva como caso booleano. Y Johnson, que repondera las aristas
  con un potencial calculado por Bellman-Ford para que ningún peso
  quede negativo y luego corre Dijkstra desde cada vértice; en grafos
  dispersos le gana a Floyd-Warshall. Una tabla final dice cuál
  conviene según la densidad.
\end{quick_sum}

\subsection*{El problema y su representación}

Sea $G = (V, E)$ un grafo dirigido con pesos $w : E \to \R$. Se quiere
el peso $\delta(u, v)$ de un camino más corto entre \textbf{todo} par
de vértices, y algo que permita reconstruir esos caminos. Los pesos
pueden ser negativos, pero se supone que no hay ciclos de peso
negativo: con uno alcanzable, $\delta$ vale $-\infty$ para algunos
pares y no existe camino más corto. Floyd-Warshall y Johnson, además,
detectan esa situación en lugar de devolver basura.

La clase del 22 de octubre (sección~\ref{sec:grafos}) presentó Floyd
con un ejemplo sobre el grafo dirigido de cuatro vértices; aquí va el
tratamiento completo, y ese ejemplo se cita en lugar de repetirse. A
diferencia de los algoritmos de fuente única, los de este capítulo
trabajan sobre la matriz de pesos y numeran los vértices de $1$ a $n$.
Cuando un grafo tenga letras, $A = 1$, $B = 2$, y así.

\begin{definicion}[Matriz de pesos]
  \label{def:matriz-pesos-pares}
  La matriz de pesos de $G$ es $W = (w_{ij})$, de $n \times n$, con
  \[
    w_{ij} = \begin{cases}
      0 & \text{si } i = j, \\
      w(i, j) & \text{si } i \neq j \text{ y } (i, j) \in E, \\
      \infty & \text{si } i \neq j \text{ y } (i, j) \notin E.
    \end{cases}
  \]
\end{definicion}

Es la matriz pesada de la clase del 22 de octubre: ceros en la
diagonal e $\infty$ donde no hay arista. Las matrices se escriben con
mayúscula y sus casillas con minúscula, $D = (d_{ij})$, y un
superíndice entre paréntesis marca la iteración, como en $D^{(k)}$.

\begin{definicion}[Distancias, predecesores y subgrafo de
  predecesores]
  \label{def:predecesores-pares}
  Al terminar, un algoritmo de todos los pares entrega $D = (d_{ij})$
  con $d_{ij} = \delta(i, j)$ y una \textbf{matriz de predecesores}
  $\Pi = (\pi_{ij})$: $\pi_{ij}$ es nulo si $i = j$ o si no hay camino
  de $i$ a $j$, y en otro caso es el predecesor de $j$ en algún camino
  más corto desde $i$. El \textbf{subgrafo de predecesores} de $i$ es
  $G_{\pi,i} = (V_{\pi,i}, E_{\pi,i})$, con
  \[
    V_{\pi,i} = \set{j \in V : \pi_{ij} \neq \text{nulo}} \cup \set{i},
    \qquad
    E_{\pi,i} = \set{(\pi_{ij}, j) : j \in V_{\pi,i} - \set{i}} .
  \]
\end{definicion}

Se pide que cada $G_{\pi,i}$ sea un árbol de caminos más cortos con
raíz $i$. Dicho de otro modo, el renglón $i$ de $\Pi$ es el arreglo
$\pi[\cdot]$ que dejaría una corrida de fuente única desde $i$
(sección~\ref{sec:caminos-fuente}), y con él se reconstruye cualquier
camino que salga de $i$.

\begin{algorithm}[H]
\caption{\textsc{Imprimir-Camino-Pares}$(\Pi, i, j)$}
\label{alg:imprimir-pares}
\begin{algorithmic}[1]
\Procedure{Imprimir-Camino-Pares}{$\Pi, i, j$}
  \If{$i = j$}
    \State imprimir $i$
  \ElsIf{$\pi_{ij} = \textsc{nulo}$}
    \State imprimir que no hay camino de $i$ a $j$
  \Else
    \State \Call{Imprimir-Camino-Pares}{$\Pi, i, \pi_{ij}$}
      \Comment{de $i$ al penúltimo}
    \State imprimir $j$
  \EndIf
\EndProcedure
\end{algorithmic}
\end{algorithm}

El procedimiento solo lee el renglón $i$: sube desde $j$ por los
predecesores hasta llegar a $i$ y escribe el camino de regreso. La
profundidad de la recursión es el número de aristas del camino, a lo
más $n - 1$, así que cuesta $\Ord{n}$. Supone que $G_{\pi,i}$ es un
árbol; con una $\Pi$ mal construida puede ciclarse para siempre, y más
adelante se verá cómo pasa eso con los empates.

\begin{ejemplo}[El grafo del capítulo]
  \label{ej:grafo-pares}
  Cormen usa un solo grafo en todo el capítulo (su figura 25.1), y
  esta sección también: cinco vértices, nueve aristas, dos de ellas
  negativas, y ningún ciclo negativo.
  \begin{center}
  \begin{tikzpicture}[baseline=(current bounding box.center)]
    \node[vert] (n1) at (-1.9, 0.4) {1};
    \node[vert] (n2) at (0, 1.8) {2};
    \node[vert] (n3) at (1.9, 0.4) {3};
    \node[vert] (n4) at (1.2, -1.8) {4};
    \node[vert] (n5) at (-1.2, -1.8) {5};
    \foreach \a/\b/\w/\p in {1/2/3/0.5, 1/3/8/0.78, 1/5/-4/0.5,
        2/4/1/0.72, 2/5/7/0.8, 3/2/4/0.5, 4/1/2/0.3, 4/3/-5/0.5,
        5/4/6/0.5}
      \draw[dirigida] (n\a) -- node[peso, pos=\p] {$\w$} (n\b);
  \end{tikzpicture}
  \hspace{1.2cm}
  $W = \begin{pmatrix}
    0 & 3 & 8 & \infty & -4 \\
    \infty & 0 & \infty & 1 & 7 \\
    \infty & 4 & 0 & \infty & \infty \\
    2 & \infty & -5 & 0 & \infty \\
    \infty & \infty & \infty & 6 & 0
  \end{pmatrix}$
  \end{center}
  Dos de sus ciclos: $1 \to 5 \to 4 \to 1$ pesa $-4 + 6 + 2 = 4$, y
  $2 \to 4 \to 3 \to 2$ pesa $1 - 5 + 4 = 0$. Ese ciclo de peso cero
  importa cuando se hable de empates en $\Pi$.
\end{ejemplo}

\subsubsection*{Correr $n$ veces un algoritmo de fuente única}

La solución obvia es correr un algoritmo de la
sección~\ref{sec:caminos-fuente} desde cada vértice. Cuánto cuesta
depende de cuál, y en Dijkstra, de la cola de prioridad:
\begin{itemize}[itemsep=0.1em]
  \item \textbf{Dijkstra} (algoritmo~\ref{alg:dijkstra-cormen}), solo
        con pesos no negativos. Con la cola en un arreglo, cada
        corrida hace $n$ extracciones de $\Tht{n}$ y hasta $m$
        decrementos de $\Ord{1}$, o sea $\Tht{n^2}$, y las $n$
        corridas suman $\Tht{n^3}$. Con montículo binario cada corrida
        cuesta $\Ord{(n + m)\log n}$ y el total es
        $\Ord{n(n + m)\log n}$, que Cormen escribe $\Ord{nm \log n}$
        suponiendo $m = \Omg{n}$. Con montículo de Fibonacci, extraer
        cuesta $\Ord{\log n}$ y decrementar $\Ord{1}$, ambos
        amortizados: $\Ord{n \log n + m}$ por corrida y
        $\Ord{n^2 \log n + nm}$ en total.
  \item \textbf{Bellman-Ford} (algoritmo~\ref{alg:bellman-ford}), con
        pesos de cualquier signo: $\Ord{nm}$ por corrida y
        $\Ord{n^2 m}$ en total, que en un grafo denso es $\Ord{n^4}$.
\end{itemize}
El resto de la sección persigue dos metas. Bajar ese $\Ord{n^4}$ a
$\Tht{n^3}$ sin renunciar a los pesos negativos, que es lo que logra
Floyd-Warshall, y alcanzar el costo de $n$ veces Dijkstra con
Fibonacci también con pesos negativos, que es lo que logra Johnson.

\begin{ojo}
  Al comparar con Floyd hay que comparar contra $n$ corridas. El costo
  de Dijkstra que se cita siempre, $\Ord{(n + m)\log n}$, es de una
  sola fuente; para todos los pares es $\Ord{n(n + m)\log n}$. En un
  grafo denso eso da $\Ord{n^3 \log n}$, peor que el $\Tht{n^3}$ de
  Floyd, y en uno disperso $\Ord{n^2 \log n}$, mucho mejor. La
  comparación completa está en la tabla del final de la sección.
\end{ojo}

% =====================================================================

\subsection*{Programación dinámica por número de aristas}

Las dos programaciones dinámicas del capítulo se apoyan en la misma
propiedad de estructura (Cormen, lema 24.1): \emph{todo subcamino de
un camino más corto es un camino más corto}. Si hubiera una forma más
barata de recorrer un tramo, cambiarla dentro del camino completo lo
abarataría, y entonces el camino no era el más corto. En particular,
si $p$ es un camino más corto de $i$ a $j \neq i$ y su última arista
es $(k, j)$, entonces $\delta(i, j) = \delta(i, k) + w_{kj}$.

\begin{definicion}[Caminatas con a lo más $r$ aristas]
  \label{def:aristas-pares}
  $l^{(r)}_{ij}$ es el peso mínimo de una caminata de $i$ a $j$ con a
  lo más $r$ aristas, o $\infty$ si no hay ninguna, y
  $L^{(r)} = (l^{(r)}_{ij})$.
\end{definicion}

Cormen usa la letra $m$ para este contador; aquí es $r$, porque $m$ ya
es el número de aristas. La definición habla de caminatas porque eso
es lo que construye la recurrencia de abajo, que puede repetir
vértices. Sin ciclos negativos da lo mismo: quitarle un ciclo a una
caminata le quita aristas y no la encarece, así que la mejor caminata
con a lo más $r$ aristas puede tomarse sin vértices repetidos.

Una caminata de $i$ a $j$ con a lo más $r$ aristas, o tiene a lo más
$r - 1$, o termina con una arista $(k, j)$ precedida por una caminata
de $i$ a $k$ con a lo más $r - 1$. Probar todos los penúltimos
vértices $k$ da la recurrencia
\[
  l^{(0)}_{ij} = \begin{cases} 0 & i = j \\ \infty & i \neq j
  \end{cases}
  \qquad
  l^{(r)}_{ij}
  = \min\Bigl( l^{(r-1)}_{ij},\
      \min_{1 \le k \le n} \set{l^{(r-1)}_{ik} + w_{kj}} \Bigr)
  = \min_{1 \le k \le n} \set{l^{(r-1)}_{ik} + w_{kj}} .
\]
La segunda igualdad usa $w_{jj} = 0$: el término $k = j$ vale
$l^{(r-1)}_{ij}$, que es justo la opción de no agregar arista. De
paso, $L^{(1)} = W$.

\begin{proposicion}[Cuándo parar]
  \label{prop:parar-pares}
  Si $G$ no tiene ciclos negativos, $l^{(r)}_{ij} = \delta(i, j)$ para
  todo $r \ge n - 1$.
\end{proposicion}

\begin{proof}
  Toda caminata de $i$ a $j$ pesa al menos $\delta(i, j)$: se le
  quitan ciclos, ninguno negativo, hasta dejar un camino, y el peso no
  sube. Por eso $l^{(r)}_{ij} \ge \delta(i, j)$ para todo $r$. Al
  revés, si $\delta(i, j) < \infty$ hay un camino más corto; no repite
  vértices, así que visita a lo más $n$ y tiene a lo más $n - 1$
  aristas, y cuenta en $l^{(r)}_{ij}$ para todo $r \ge n - 1$. Si
  $\delta(i, j) = \infty$ no hay caminata y los dos lados valen
  $\infty$.
\end{proof}

\subsubsection*{Extender caminos es multiplicar matrices}

\begin{algorithm}[H]
\caption{\textsc{Extender-Caminos}$(L, W)$}
\label{alg:extender-caminos}
\begin{algorithmic}[1]
\Procedure{Extender-Caminos}{$L, W$}
  \State $n \gets$ número de renglones de $L$
  \State $L' \gets$ matriz nueva de $n \times n$
  \For{$i \gets 1$ \Hasta $n$}
    \For{$j \gets 1$ \Hasta $n$}
      \State $l'_{ij} \gets \infty$ \Comment{neutro del mínimo}
      \For{$k \gets 1$ \Hasta $n$}
        \State $l'_{ij} \gets \min(l'_{ij},\ l_{ik} + w_{kj})$
          \Comment{$k$ es el penúltimo}
      \EndFor
    \EndFor
  \EndFor
  \State \Return $L'$
\EndProcedure
\end{algorithmic}
\end{algorithm}

\textsc{Extender-Caminos} recibe $L^{(r-1)}$ y $W$ y devuelve
$L^{(r)}$: calcula la recurrencia en todas las casillas. La línea~6
arranca en $\infty$ porque es el neutro del mínimo, como $0$ lo es de
la suma; en la línea~8, $k$ recorre los posibles penúltimos vértices.
Son tres ciclos anidados de $n$ vueltas con cuerpo $\Ord{1}$:
$\Tht{n^3}$.

El código es el producto de matrices de siempre,
$c_{ij} = \sum_k a_{ik} \cdot b_{kj}$, con dos sustituciones: la suma
se vuelve $\min$ y el producto se vuelve $+$. La tabla completa la
analogía.

\begin{center}
\small
\begin{tabular}{lll}
  \toprule
  & \textbf{Producto usual} & \textbf{Producto $(\min, +)$} \\
  \midrule
  Casilla & $c_{ij} = \sum_k a_{ik} \cdot b_{kj}$
          & $l'_{ij} = \min_k \, (l_{ik} + w_{kj})$ \\
  ``Suma'' & $+$ & $\min$ \\
  ``Producto'' & $\cdot$ & $+$ \\
  Neutro de la suma & $0$ & $\infty$ \\
  Neutro del producto & $1$ & $0$ \\
  Identidad & $I$: unos en la diagonal, ceros fuera
            & $L^{(0)}$: ceros en la diagonal, $\infty$ fuera \\
  \bottomrule
\end{tabular}
\end{center}

Con esas operaciones, $(\R \cup \set{\infty}, \min, +)$ es un
\textbf{semianillo}: $\min$ y $+$ son asociativas y tienen neutro, y
$+$ se distribuye sobre $\min$, $x + \min(y, z) = \min(x + y, x + z)$.
Lo que no tiene es resta: no hay manera de deshacer un mínimo. Se
escribe $A \cdot B$ para el producto $(\min, +)$ y $W^r$ para sus
potencias, de modo que $L^{(r)} = W^r$ y, por la
proposición~\ref{prop:parar-pares}, la respuesta es $W^{n-1}$.

\begin{proposicion}[El producto $(\min, +)$ es asociativo]
  \label{prop:asociativo-pares}
  Para matrices $A, B, C$ de $n \times n$,
  $(A \cdot B) \cdot C = A \cdot (B \cdot C)$.
\end{proposicion}

\begin{proof}
  Por la distributividad de $+$ sobre $\min$,
  \begin{align*}
    ((A \cdot B) \cdot C)_{ij}
    &= \min_{l} \Bigl( \min_{k} \, (a_{ik} + b_{kl}) + c_{lj} \Bigr)
     = \min_{k, l} \, (a_{ik} + b_{kl} + c_{lj}) \\
    &= \min_{k} \Bigl( a_{ik} + \min_{l} \, (b_{kl} + c_{lj}) \Bigr)
     = (A \cdot (B \cdot C))_{ij} . \qedhere
  \end{align*}
\end{proof}

La consecuencia práctica es $W^{2r} = W^r \cdot W^r$, que también se
lee directo: una caminata de a lo más $2r$ aristas se parte en dos de
a lo más $r$ cada una. No hace falta calcular las potencias una por
una.

\begin{algorithm}[H]
\caption{\textsc{Todos-Pares-Lento}$(W)$}
\label{alg:pares-lento}
\begin{algorithmic}[1]
\Procedure{Todos-Pares-Lento}{$W$}
  \State $n \gets$ número de renglones de $W$
  \State $L^{(1)} \gets W$
  \For{$r \gets 2$ \Hasta $n - 1$}
    \State $L^{(r)} \gets \Call{Extender-Caminos}{L^{(r-1)}, W}$
  \EndFor
  \State \Return $L^{(n-1)}$
\EndProcedure
\end{algorithmic}
\end{algorithm}

\begin{algorithm}[H]
\caption{\textsc{Todos-Pares-Rápido}$(W)$}
\label{alg:pares-rapido}
\begin{algorithmic}[1]
\Procedure{Todos-Pares-Rápido}{$W$}
  \State $n \gets$ número de renglones de $W$
  \State $L^{(1)} \gets W$, \quad $r \gets 1$
  \While{$r < n - 1$}
    \State $L^{(2r)} \gets \Call{Extender-Caminos}{L^{(r)}, L^{(r)}}$
      \Comment{elevar al cuadrado}
    \State $r \gets 2r$
  \EndWhile
  \State \Return $L^{(r)}$
\EndProcedure
\end{algorithmic}
\end{algorithm}

\textsc{Todos-Pares-Lento} hace $n - 2$ productos de $\Tht{n^3}$:
$\Tht{n^4}$, lo mismo que $n$ corridas de Bellman-Ford en un grafo
denso. Como está escrito guarda todas las $L^{(r)}$, pero cada una
solo necesita la anterior y dos matrices bastan.

\textsc{Todos-Pares-Rápido} calcula $W, W^2, W^4, \dots$ elevando al
cuadrado, la misma idea de la exponenciación rápida. Después de $t$
vueltas $r = 2^t$, y el ciclo se detiene en la primera $t$ con
$2^t \ge n - 1$, o sea $t = \lceil \log(n - 1) \rceil$. La última
matriz es $L^{(r)}$ con $n - 1 \le r < 2n - 2$, y por la
proposición~\ref{prop:parar-pares} pasarse de $n - 1$ no estropea
nada. Costo: $\lceil \log(n - 1) \rceil$ productos de $\Tht{n^3}$,
$\Tht{n^3 \log n}$, con dos matrices de memoria. Para $n = 1000$ la
versión lenta hace $998$ productos y la rápida $10$: del orden de
$10^{12}$ operaciones contra $10^{10}$. Floyd-Warshall, en la
siguiente subsección, hace $10^9$.

\begin{ejemplo}[Lento contra cuadrados repetidos]
  \label{ej:cuadrados-pares}
  Con $n = 5$ la versión lenta calcula $L^{(2)}, L^{(3)}, L^{(4)}$ y
  la rápida solo $L^{(2)}$ y $L^{(4)} = L^{(2)} \cdot L^{(2)}$.
  \begin{center}
  \begin{tikzpicture}[
    caja/.style={draw=acento, rounded corners=2pt, fill=fondoIdea,
                 font=\small, minimum width=1.1cm,
                 minimum height=0.6cm}]
    \node[font=\small, anchor=east] at (-0.7, 0) {Lenta:};
    \foreach \e/\x in {1/0, 2/2.2, 3/4.4, 4/6.6}
      \node[caja] (s\e) at (\x, 0) {$W^{\e}$};
    \foreach \a/\b in {1/2, 2/3, 3/4}
      \draw[dirigida] (s\a) -- node[nota, above] {$\cdot\, W$} (s\b);
    \node[font=\small, anchor=east] at (-0.7, -1.3) {Rápida:};
    \node[caja] (f1) at (0, -1.3) {$W^{1}$};
    \node[caja] (f2) at (2.2, -1.3) {$W^{2}$};
    \node[caja] (f4) at (6.6, -1.3) {$W^{4}$};
    \draw[dirigida] (f1) -- node[nota, above] {$(\cdot)^2$} (f2);
    \draw[dirigida] (f2) -- node[nota, above] {$(\cdot)^2$} (f4);
    \node[nota, anchor=west] at (7.4, 0) {$n - 2 = 3$ productos};
    \node[nota, anchor=west] at (7.4, -1.3)
      {$\lceil \log (n-1) \rceil = 2$ productos};
  \end{tikzpicture}
  \end{center}
  Sobre el grafo del ejemplo~\ref{ej:grafo-pares}, en azul las
  casillas que cambian respecto a la matriz anterior:
  {\small
  \setlength{\arraycolsep}{3.5pt}
  \[
  L^{(1)} = \begin{pmatrix}
    0 & 3 & 8 & \infty & -4 \\
    \infty & 0 & \infty & 1 & 7 \\
    \infty & 4 & 0 & \infty & \infty \\
    2 & \infty & -5 & 0 & \infty \\
    \infty & \infty & \infty & 6 & 0
  \end{pmatrix}
  \qquad
  L^{(2)} = \begin{pmatrix}
    0 & 3 & 8 & \paresCambio{2} & -4 \\
    \paresCambio{3} & 0 & \paresCambio{-4} & 1 & 7 \\
    \infty & 4 & 0 & \paresCambio{5} & \paresCambio{11} \\
    2 & \paresCambio{-1} & -5 & 0 & \paresCambio{-2} \\
    \paresCambio{8} & \infty & \paresCambio{1} & 6 & 0
  \end{pmatrix}
  \]
  \[
  L^{(3)} = \begin{pmatrix}
    0 & 3 & \paresCambio{-3} & 2 & -4 \\
    3 & 0 & -4 & 1 & \paresCambio{-1} \\
    \paresCambio{7} & 4 & 0 & 5 & 11 \\
    2 & -1 & -5 & 0 & -2 \\
    8 & \paresCambio{5} & 1 & 6 & 0
  \end{pmatrix}
  \qquad
  L^{(4)} = \begin{pmatrix}
    0 & \paresCambio{1} & -3 & 2 & -4 \\
    3 & 0 & -4 & 1 & -1 \\
    7 & 4 & 0 & 5 & \paresCambio{3} \\
    2 & -1 & -5 & 0 & -2 \\
    8 & 5 & 1 & 6 & 0
  \end{pmatrix}
  \]
  }%
  La rápida llega a la misma $L^{(4)}$, y $L^{(4)} \cdot W = L^{(4)}$,
  como pide la proposición~\ref{prop:parar-pares}. Una casilla baja en
  $L^{(r)}$ cuando el mejor camino con a lo más $r$ aristas usa
  exactamente $r$. La $(1, 2)$ vale $3$ hasta $L^{(3)}$ y baja a $1$
  en $L^{(4)}$ por el camino $1 \to 5 \to 4 \to 3 \to 2$, de cuatro
  aristas y peso $-4 + 6 - 5 + 4 = 1$. Las dos versiones llegan a ese
  $1$ por rutas distintas:
  \begin{itemize}[itemsep=0.1em]
    \item Lenta, $L^{(4)} = L^{(3)} \cdot W$: los candidatos
          $l^{(3)}_{1k} + w_{k2}$ son $3, 3, 1, \infty, \infty$ y gana
          $k = 3$, tres aristas hasta $3$ y luego $3 \to 2$.
    \item Rápida, $L^{(4)} = L^{(2)} \cdot L^{(2)}$: los candidatos
          $l^{(2)}_{1k} + l^{(2)}_{k2}$ son $3, 3, 12, 1, \infty$ y
          gana $k = 4$, dos aristas hasta $4$ ($1 \to 5 \to 4$) y dos
          más hasta $2$ ($4 \to 3 \to 2$).
  \end{itemize}
  En la rápida, $k$ ya no es el penúltimo vértice sino el punto medio
  del camino.
\end{ejemplo}

\begin{ojo}
  \begin{itemize}[itemsep=0.2em, leftmargin=1.2em]
    \item Pasarse de $n - 1$ solo es inofensivo sin ciclos negativos.
          Con uno, $L^{(r)}$ sigue bajando para siempre. Un ciclo
          negativo de $c$ aristas deja $l^{(r)}_{ii} < 0$ en sus
          vértices en cuanto $r \ge c$, y un ciclo tiene a lo más $n$
          aristas, así que basta revisar la diagonal de una $L^{(r)}$
          con $r \ge n$. Con $r = n - 1$ no alcanza: el ciclo
          $1 \to 2 \to 3 \to 1$ con pesos $-1$ deja la diagonal de
          $L^{(2)}$ en ceros y la de $L^{(3)}$ en $-3$.
    \item El algoritmo de Strassen no sirve para el producto
          $(\min, +)$: sus fórmulas restan, y el mínimo no tiene
          inverso.
  \end{itemize}
\end{ojo}

\begin{ejercicio}
  ¿Por qué hace falta $w_{ii} = 0$ en la recurrencia? ¿Qué papel juega
  $L^{(0)}$ en el producto $(\min, +)$? (Cormen, ejercicios 25.1-2 y
  25.1-3.)

  \emph{Solución.} Con $w_{jj} = 0$ el término $k = j$ es
  $l^{(r-1)}_{ij} + 0$, la opción de no agregar arista; sin eso,
  $L^{(r)}$ contaría caminatas de exactamente $r$ aristas y no de a lo
  más $r$. $L^{(0)}$, con ceros en la diagonal e $\infty$ fuera, es la
  identidad: $(L^{(0)} \cdot A)_{ij} = \min_k (l^{(0)}_{ik} + a_{kj})
  = a_{ij}$, porque el único término finito es el de $k = i$.
\end{ejercicio}

\begin{ejercicio}
  Sea $x$ el vector renglón con $0$ en la posición de una fuente $s$ e
  $\infty$ en las demás. ¿Qué es $x \cdot W^r$ en el producto
  $(\min, +)$, y a qué algoritmo de la
  sección~\ref{sec:caminos-fuente} corresponde calcular
  $x \cdot W, x \cdot W^2, \dots$? (Cormen, ejercicio 25.1-5.)

  \emph{Solución.} $(x \cdot W^r)_j$ es el peso mínimo de una caminata
  de $s$ a $j$ con a lo más $r$ aristas. Cada producto,
  $(x \cdot W)_j = \min_k (x_k + w_{kj})$, relaja todas las aristas
  que entran a $j$ con las estimaciones de la pasada anterior: es una
  pasada de Bellman-Ford, y $n - 1$ pasadas dan $\delta(s, \cdot)$. La
  programación dinámica por número de aristas es Bellman-Ford hecho
  matriz, para todas las fuentes a la vez.
\end{ejercicio}

% =====================================================================

\subsection*{Floyd-Warshall}

La clase del 22 de octubre trae la recurrencia, el pseudocódigo en el
lugar y un ejemplo sobre el grafo dirigido de cuatro vértices. Aquí se
justifica cada pieza: por qué la recurrencia es correcta y dónde entra
la hipótesis de no haber ciclos negativos, cómo se construye $\Pi$,
por qué basta una matriz y cómo se detectan los ciclos negativos. El
ejemplo grande usa pesos negativos, que es justo donde Floyd-Warshall
tiene algo que Dijkstra no.

La idea es cambiar el estado de la programación dinámica: en lugar de
cuántas aristas usa el camino, qué vértices puede usar como
intermedios.

\begin{definicion}[Vértice intermedio]
  \label{def:intermedio-floyd}
  Los \textbf{intermedios} de una caminata
  $p = \langle v_1, v_2, \dots, v_l \rangle$ son los vértices que
  aparecen en las posiciones $2$ a $l - 1$, es decir, todos menos los
  extremos. Si $G$ no tiene ciclos negativos, $d^{(k)}_{ij}$ es el
  peso mínimo de una caminata de $i$ a $j$ cuyos intermedios están
  todos en $\set{1, \dots, k}$, o $\infty$ si no hay ninguna, y
  $D^{(k)} = (d^{(k)}_{ij})$.
\end{definicion}

Con $k = 0$ no se permite ningún intermedio: la caminata es la arista
directa o la vacía, de ahí $d^{(0)}_{ij} = w_{ij}$. Con $k = n$ se
permite todo, así que $d^{(n)}_{ij} = \delta(i, j)$.

\begin{lema}[Estructura de Floyd-Warshall]
  \label{lem:estructura-floyd}
  Si $G$ no tiene ciclos negativos, para $k \ge 1$
  \[
    d^{(k)}_{ij} = \min\bigl( d^{(k-1)}_{ij},\
                   d^{(k-1)}_{ik} + d^{(k-1)}_{kj} \bigr) .
  \]
\end{lema}

\begin{proof}
  Primero $\le$. Una caminata con intermedios en
  $\set{1, \dots, k - 1}$ también los tiene en $\set{1, \dots, k}$,
  así que $d^{(k)}_{ij} \le d^{(k-1)}_{ij}$. Y pegar una caminata
  óptima de $i$ a $k$ con una de $k$ a $j$, las dos con intermedios en
  $\set{1, \dots, k - 1}$, da una caminata de $i$ a $j$ cuyos
  intermedios son los de ambas más el propio $k$:
  $d^{(k)}_{ij} \le d^{(k-1)}_{ik} + d^{(k-1)}_{kj}$.

  Ahora $\ge$, y aquí entra la hipótesis. Sea $p$ una caminata óptima
  de $i$ a $j$ con intermedios en $\set{1, \dots, k}$. Si $p$ repite
  un vértice, contiene un ciclo; ese ciclo no es negativo, y quitarlo
  deja una caminata con los mismos extremos, intermedios en el mismo
  conjunto y peso no mayor. Repitiendo, $p$ se puede tomar sin
  vértices repetidos. Hay dos casos.
  \begin{itemize}[itemsep=0.1em]
    \item $k$ no es intermedio de $p$. Entonces todos sus intermedios
          están en $\set{1, \dots, k - 1}$, así que
          $d^{(k)}_{ij} = w(p) \ge d^{(k-1)}_{ij}$.
    \item $k$ es intermedio de $p$. Aparece exactamente una vez,
          porque $p$ no repite vértices. Partiendo $p$ en $k$ quedan
          $p_1$ de $i$ a $k$ y $p_2$ de $k$ a $j$, y en ninguno de los
          dos $k$ es intermedio: es un extremo y no aparece en otra
          parte. Los dos usan solo $\set{1, \dots, k - 1}$, así que
          $d^{(k)}_{ij} = w(p_1) + w(p_2)
          \ge d^{(k-1)}_{ik} + d^{(k-1)}_{kj}$. \qedhere
  \end{itemize}
\end{proof}

La figura es la 25.3 de Cormen: el camino óptimo que pasa por $k$ se
parte en dos tramos que solo usan intermedios menores que $k$.

\begin{center}
\begin{tikzpicture}
  \node[vert] (i) at (0, 0) {$i$};
  \node[vertMarca] (k) at (4.5, 0.9) {$k$};
  \node[vert] (j) at (9, 0) {$j$};
  \draw[dirigida] (i) .. controls (1.3, 1.5) and (2.8, -0.5) .. (k)
    node[nota, pos=0.45, above=6pt] {$p_1$};
  \draw[dirigida] (k) .. controls (6.0, 2.1) and (7.2, -0.7) .. (j)
    node[nota, pos=0.5, above=6pt] {$p_2$};
  \draw[gray] (0, 1.9) -- (0, 2.05) -- (4.35, 2.05) -- (4.35, 1.9);
  \draw[gray] (4.65, 1.9) -- (4.65, 2.05) -- (9, 2.05) -- (9, 1.9);
  \node[nota, above] at (2.17, 2.05)
    {intermedios en $\set{1, \dots, k-1}$};
  \node[nota, above] at (6.82, 2.05)
    {intermedios en $\set{1, \dots, k-1}$};
  \draw[gray] (0, -0.6) -- (0, -0.75) -- (9, -0.75) -- (9, -0.6);
  \node[nota, below] at (4.5, -0.75)
    {$p$: intermedios en $\set{1, \dots, k}$, y $k$ aparece una sola
     vez};
\end{tikzpicture}
\end{center}

\begin{idea}
  La hipótesis de no haber ciclos negativos sirve para una sola cosa:
  poder tomar el óptimo sin vértices repetidos, y con eso que $k$
  aparezca a lo más una vez. Si un ciclo negativo pasara por $k$, la
  mejor caminata daría vueltas a ese ciclo sin límite, $k$ aparecería
  tantas veces como se quisiera y no habría manera de partirla en solo
  dos tramos. Ganada la hipótesis, cada casilla elige entre dos
  opciones, usar $k$ o no, en lugar de entre $n$ penúltimos vértices
  como en la subsección anterior. Las dos programaciones dinámicas
  tienen $n^3$ estados ($n^2$ pares por $n$ etapas); la diferencia
  entre $\Tht{n^4}$ y $\Tht{n^3}$ es lo que cuesta resolver cada uno.
\end{idea}

\subsubsection*{Los predecesores}

$\Pi$ se construye con una regla igual de simple. Al principio el
predecesor de $j$ desde $i$ es $i$ si hay arista, y nulo si no:
\[
  \pi^{(0)}_{ij} = \begin{cases}
    \text{nulo} & \text{si } i = j \text{ o } w_{ij} = \infty, \\
    i & \text{si } i \neq j \text{ y } w_{ij} < \infty .
  \end{cases}
\]
En la vuelta $k$, el predecesor solo cambia si pasar por $k$ mejora
estrictamente:
\[
  \pi^{(k)}_{ij} = \begin{cases}
    \pi^{(k-1)}_{ij} & \text{si } d^{(k-1)}_{ij} \le
                       d^{(k-1)}_{ik} + d^{(k-1)}_{kj}, \\
    \pi^{(k-1)}_{kj} & \text{si } d^{(k-1)}_{ij} >
                       d^{(k-1)}_{ik} + d^{(k-1)}_{kj}.
  \end{cases}
\]
Si el mejor camino pasa por $k$, su último tramo es un camino óptimo
de $k$ a $j$, así que el predecesor de $j$ es el que ya se tenía para
el par $(k, j)$. En un empate se conserva el predecesor viejo, y eso
no es un detalle de estilo: la otra elección rompe $\Pi$
(ejercicio~\ref{ej:empates-floyd}). Que cada $G_{\pi,i}$ salga un
árbol de caminos más cortos es el ejercicio 25.2-3 de Cormen; la clave
es mostrar que $\pi^{(k)}_{ij} = l$ implica
$d^{(k)}_{ij} \ge d^{(k)}_{il} + w_{lj}$, lo que impide ciclos en
$G_{\pi,i}$, y luego adaptar el argumento del subgrafo de
predecesores de fuente única
(proposición~\ref{prop:subgrafo-pred}).

\subsubsection*{El algoritmo completo}

La versión de abajo junta la recurrencia sin superíndices, que es la
de la clase, con el cálculo de $\Pi$ y la prueba de la diagonal.

\begin{algorithm}[H]
\caption{\textsc{Floyd-Warshall}$(W)$, con $\Pi$ y en el lugar}
\label{alg:floyd-warshall}
\begin{algorithmic}[1]
\Procedure{Floyd-Warshall}{$W$}
  \State $n \gets$ número de renglones de $W$; \quad $D \gets W$
  \For{$i \gets 1$ \Hasta $n$}
    \For{$j \gets 1$ \Hasta $n$}
      \State $\pi_{ij} \gets i$ si $i \neq j$ y $w_{ij} < \infty$; si
             no, $\pi_{ij} \gets \textsc{nulo}$
    \EndFor
  \EndFor
  \For{$k \gets 1$ \Hasta $n$} \Comment{$k$ va afuera, siempre}
    \For{$i \gets 1$ \Hasta $n$}
      \For{$j \gets 1$ \Hasta $n$}
        \If{$d_{ik} + d_{kj} < d_{ij}$}
          \State $d_{ij} \gets d_{ik} + d_{kj}$
          \State $\pi_{ij} \gets \pi_{kj}$
        \EndIf
      \EndFor
    \EndFor
  \EndFor
  \For{$i \gets 1$ \Hasta $n$}
    \If{$d_{ii} < 0$}
      \State \Return ``hay un ciclo de peso negativo''
    \EndIf
  \EndFor
  \State \Return $D, \Pi$
\EndProcedure
\end{algorithmic}
\end{algorithm}

\begin{itemize}[itemsep=0.1em]
  \item Líneas 3 a 7: $\Pi^{(0)}$, en $\Tht{n^2}$.
  \item Línea 8: $k$ tiene que ser el ciclo de afuera. La recurrencia
        necesita $D^{(k-1)}$ completa antes de empezar con $k$.
  \item Línea 11: la desigualdad es estricta. En un empate no se toca
        nada, que es el $\le$ de la regla de $\Pi$.
  \item Línea 13: $\pi_{kj}$ se lee de la misma matriz que se está
        modificando; la proposición~\ref{prop:floyd-lugar} dice por qué
        está bien.
  \item Líneas 18 a 22: el teorema~\ref{teo:floyd-diagonal}. Si hay un
        ciclo negativo, $D$ y $\Pi$ no significan nada y no deben
        usarse.
\end{itemize}

\textbf{Costo.} La inicialización es $\Tht{n^2}$, el triple ciclo
ejecuta su cuerpo de $\Ord{1}$ exactamente $n^3$ veces y la diagonal
cuesta $\Tht{n}$: $\Tht{n^3}$ en tiempo y $\Tht{n^2}$ en memoria. No
hay mejor caso. El costo depende solo de $n$, no de $m$ ni de los
pesos, y eso es una virtud en grafos densos y un defecto en los
dispersos.

\begin{proposicion}[Basta una matriz]
  \label{prop:floyd-lugar}
  Si $G$ no tiene ciclos negativos, durante la vuelta $k$ el renglón
  $k$ y la columna $k$ de $D$ y de $\Pi$ no cambian. Por eso la versión
  sin superíndices calcula exactamente $D^{(n)}$ y $\Pi^{(n)}$, con
  $\Tht{n^2}$ de memoria.
\end{proposicion}

\begin{proof}
  Primero, $d^{(k-1)}_{kk} = 0$: la caminata vacía de $k$ a $k$ pesa
  $0$, y ninguna caminata cerrada pesa menos porque no hay ciclos
  negativos. Entonces
  \[
    d^{(k)}_{ik} = \min\bigl(d^{(k-1)}_{ik},\
                   d^{(k-1)}_{ik} + d^{(k-1)}_{kk}\bigr)
                 = d^{(k-1)}_{ik} ,
  \]
  y lo mismo con $d^{(k)}_{kj}$. En $\Pi$, $\pi_{kj}$ solo cambiaría si
  $d_{kk} + d_{kj} < d_{kj}$, es decir, si $d_{kk} < 0$. En la vuelta
  $k$ las únicas casillas que se leen además de $d_{ij}$ son $d_{ik}$,
  $d_{kj}$ y $\pi_{kj}$, y valen lo mismo antes y después de que el
  ciclo pase por ellas. Leerlas ya actualizadas o todavía sin
  actualizar da igual, y por inducción sobre $k$ la matriz única
  coincide con $D^{(k)}$ y $\Pi^{(k)}$ al final de cada vuelta.
\end{proof}

Es la misma pregunta que surge en la mochila 0-1 con una sola fila,
donde el orden del recorrido sí importa porque la casilla que se lee
puede haber cambiado en la misma pasada. Aquí no importa porque las
casillas que se leen, el renglón y la columna $k$, son justo las que
no cambian.

\begin{teorema}[La diagonal delata los ciclos negativos]
  \label{teo:floyd-diagonal}
  Sin suponer nada sobre los ciclos, sea $D$ la matriz que deja
  Floyd-Warshall, por capas o en el lugar. $G$ tiene un ciclo de peso
  negativo si y solo si $d_{ii} < 0$ para algún $i$.
\end{teorema}

\begin{proof}
  Sin la hipótesis, $d^{(k)}_{ij}$ ya no se puede definir como un
  mínimo, así que la prueba habla de los números que calcula la
  recurrencia. Dos hechos, los dos por inducción sobre $k$:
  \begin{enumerate}[label=(\alph*), itemsep=0.1em]
    \item Cada valor finito es el peso de alguna caminata de $i$ a $j$.
          Para $k = 0$ es una arista o la caminata vacía, y cada
          actualización pega una caminata de $i$ a $k$ con una de $k$
          a $j$.
    \item $d^{(k)}_{ij} \le w(p)$ para todo camino $p$ de $i$ a $j$
          con intermedios en $\set{1, \dots, k}$. Si $k$ no está en
          $p$, vale porque $d^{(k)}_{ij} \le d^{(k-1)}_{ij}$; si está,
          aparece una vez y $p$ se parte en dos caminos que solo usan
          $\set{1, \dots, k - 1}$, como en el
          lema~\ref{lem:estructura-floyd}.
  \end{enumerate}
  La versión en el lugar cumple (a) por la misma razón, y sus valores
  nunca quedan arriba de los de la versión por capas, porque cada
  actualización toma el mínimo de números que ya son menores o
  iguales. Así que también cumple (b).

  ($\Leftarrow$) Si $d_{ii} < 0$, por (a) hay una caminata cerrada de
  peso negativo. Si no es un ciclo, repite un vértice a la mitad y se
  parte en dos caminatas cerradas más cortas cuyos pesos suman el
  total; una de ellas es negativa. Repitiendo se llega a un ciclo
  negativo.

  ($\Rightarrow$) Sea $C$ un ciclo negativo y $q$ su vértice de número
  más alto. Un lazo $(q, q)$ no cuenta, porque $w_{qq} = 0$ por
  definición (ver el ojo de abajo), así que $C$ tiene otro vértice
  $i$. Los dos tramos de $C$, de $i$ a $q$ y de $q$ a $i$, son caminos
  cuyos intermedios son otros vértices de $C$, todos menores que $q$.
  Por (b) con $k = q - 1$,
  \[
    d^{(q)}_{ii} \le d^{(q-1)}_{iq} + d^{(q-1)}_{qi} \le w(C) < 0 ,
  \]
  y como los valores nunca suben, $d^{(n)}_{ii} < 0$.
\end{proof}

\begin{ejemplo}[Floyd-Warshall sobre el grafo del capítulo]
  \label{ej:floyd-pares}
  Las seis parejas $D^{(k)}$, $\Pi^{(k)}$ sobre el grafo del
  ejemplo~\ref{ej:grafo-pares}; en azul, las casillas que cambian en
  la vuelta $k$. Coinciden con la figura 25.4 de Cormen.
  {\small
  \setlength{\arraycolsep}{3.5pt}
  \[
  D^{(0)} = \begin{pmatrix}
    0 & 3 & 8 & \infty & -4 \\
    \infty & 0 & \infty & 1 & 7 \\
    \infty & 4 & 0 & \infty & \infty \\
    2 & \infty & -5 & 0 & \infty \\
    \infty & \infty & \infty & 6 & 0
  \end{pmatrix}
  \qquad
  \Pi^{(0)} = \begin{pmatrix}
    \paresNulo & 1 & 1 & \paresNulo & 1 \\
    \paresNulo & \paresNulo & \paresNulo & 2 & 2 \\
    \paresNulo & 3 & \paresNulo & \paresNulo & \paresNulo \\
    4 & \paresNulo & 4 & \paresNulo & \paresNulo \\
    \paresNulo & \paresNulo & \paresNulo & 5 & \paresNulo
  \end{pmatrix}
  \]
  \[
  D^{(1)} = \begin{pmatrix}
    0 & 3 & 8 & \infty & -4 \\
    \infty & 0 & \infty & 1 & 7 \\
    \infty & 4 & 0 & \infty & \infty \\
    2 & \paresCambio{5} & -5 & 0 & \paresCambio{-2} \\
    \infty & \infty & \infty & 6 & 0
  \end{pmatrix}
  \qquad
  \Pi^{(1)} = \begin{pmatrix}
    \paresNulo & 1 & 1 & \paresNulo & 1 \\
    \paresNulo & \paresNulo & \paresNulo & 2 & 2 \\
    \paresNulo & 3 & \paresNulo & \paresNulo & \paresNulo \\
    4 & \paresCambio{1} & 4 & \paresNulo & \paresCambio{1} \\
    \paresNulo & \paresNulo & \paresNulo & 5 & \paresNulo
  \end{pmatrix}
  \]
  \[
  D^{(2)} = \begin{pmatrix}
    0 & 3 & 8 & \paresCambio{4} & -4 \\
    \infty & 0 & \infty & 1 & 7 \\
    \infty & 4 & 0 & \paresCambio{5} & \paresCambio{11} \\
    2 & 5 & -5 & 0 & -2 \\
    \infty & \infty & \infty & 6 & 0
  \end{pmatrix}
  \qquad
  \Pi^{(2)} = \begin{pmatrix}
    \paresNulo & 1 & 1 & \paresCambio{2} & 1 \\
    \paresNulo & \paresNulo & \paresNulo & 2 & 2 \\
    \paresNulo & 3 & \paresNulo & \paresCambio{2} & \paresCambio{2} \\
    4 & 1 & 4 & \paresNulo & 1 \\
    \paresNulo & \paresNulo & \paresNulo & 5 & \paresNulo
  \end{pmatrix}
  \]
  \[
  D^{(3)} = \begin{pmatrix}
    0 & 3 & 8 & 4 & -4 \\
    \infty & 0 & \infty & 1 & 7 \\
    \infty & 4 & 0 & 5 & 11 \\
    2 & \paresCambio{-1} & -5 & 0 & -2 \\
    \infty & \infty & \infty & 6 & 0
  \end{pmatrix}
  \qquad
  \Pi^{(3)} = \begin{pmatrix}
    \paresNulo & 1 & 1 & 2 & 1 \\
    \paresNulo & \paresNulo & \paresNulo & 2 & 2 \\
    \paresNulo & 3 & \paresNulo & 2 & 2 \\
    4 & \paresCambio{3} & 4 & \paresNulo & 1 \\
    \paresNulo & \paresNulo & \paresNulo & 5 & \paresNulo
  \end{pmatrix}
  \]
  \[
  D^{(4)} = \begin{pmatrix}
    0 & 3 & \paresCambio{-1} & 4 & -4 \\
    \paresCambio{3} & 0 & \paresCambio{-4} & 1 & \paresCambio{-1} \\
    \paresCambio{7} & 4 & 0 & 5 & \paresCambio{3} \\
    2 & -1 & -5 & 0 & -2 \\
    \paresCambio{8} & \paresCambio{5} & \paresCambio{1} & 6 & 0
  \end{pmatrix}
  \qquad
  \Pi^{(4)} = \begin{pmatrix}
    \paresNulo & 1 & \paresCambio{4} & 2 & 1 \\
    \paresCambio{4} & \paresNulo & \paresCambio{4} & 2 & \paresCambio{1} \\
    \paresCambio{4} & 3 & \paresNulo & 2 & \paresCambio{1} \\
    4 & 3 & 4 & \paresNulo & 1 \\
    \paresCambio{4} & \paresCambio{3} & \paresCambio{4} & 5 & \paresNulo
  \end{pmatrix}
  \]
  \[
  D^{(5)} = \begin{pmatrix}
    0 & \paresCambio{1} & \paresCambio{-3} & \paresCambio{2} & -4 \\
    3 & 0 & -4 & 1 & -1 \\
    7 & 4 & 0 & 5 & 3 \\
    2 & -1 & -5 & 0 & -2 \\
    8 & 5 & 1 & 6 & 0
  \end{pmatrix}
  \qquad
  \Pi^{(5)} = \begin{pmatrix}
    \paresNulo & \paresCambio{3} & 4 & \paresCambio{5} & 1 \\
    4 & \paresNulo & 4 & 2 & 1 \\
    4 & 3 & \paresNulo & 2 & 1 \\
    4 & 3 & 4 & \paresNulo & 1 \\
    4 & 3 & 4 & 5 & \paresNulo
  \end{pmatrix}
  \]
  }%
\end{ejemplo}

Qué pasa en cada vuelta. Ninguna casilla del renglón $k$ ni de la
columna $k$ cambia en la vuelta $k$, como dice la
proposición~\ref{prop:floyd-lugar}.
\begin{itemize}[itemsep=0.1em]
  \item $k = 1$: solo $4$ entra a $1$, así que solo el renglón $4$
        puede mejorar pasando por $1$: $d_{42} = d_{41} + d_{12} =
        2 + 3 = 5$ y $d_{45} = 2 + (-4) = -2$, las dos con predecesor
        $\pi_{12} = \pi_{15} = 1$.
  \item $k = 2$: a $2$ entran aristas desde $1$ y $3$. El renglón $1$
        gana $1 \to 2 \to 4$, con $d_{14} = 3 + 1 = 4$, y el $3$ gana
        $3 \to 2 \to 4$ y $3 \to 2 \to 5$, con $5$ y $11$.
  \item $k = 3$: un solo cambio, y es la mejora de un valor finito:
        $d_{42} = d_{43} + d_{32} = -5 + 4 = -1 < 5$ por la arista
        negativa $4 \to 3$, y $\pi_{42} \gets \pi_{32} = 3$.
  \item $k = 4$: la vuelta grande, nueve casillas. Casi todos los
        caminos buenos pasan por $4$: tiene las salidas $4 \to 1$ y
        $4 \to 3$, la negativa, y es la única salida de $5$, cuyo
        renglón se llena completo. Por ejemplo,
        $d_{52} = d_{54} + d_{42} = 6 + (-1) = 5$ con
        $\pi_{52} \gets \pi_{42} = 3$: el predecesor viene del renglón
        $4$, que ya sabía llegar a $2$ por $3$.
  \item $k = 5$: solo mejora el renglón $1$, gracias a la arista
        $1 \to 5$ de peso $-4$; por ejemplo,
        $d_{12} = d_{15} + d_{52} = -4 + 5 = 1$ y
        $\pi_{12} \gets \pi_{52} = 3$. Todo camino que sale de $5$
        empieza con $5 \to 4$, y los demás vértices ya llegaban a $4$
        más barato sin pasar por $5$.
\end{itemize}
Un detalle de la última vuelta: $d_{13}$ baja de $-1$ a $-3$ y
$\pi_{13}$ \emph{no} cambia. El camino nuevo,
$1 \to 5 \to 4 \to 3$, termina con la misma arista $4 \to 3$ que el
anterior, $1 \to 2 \to 4 \to 3$, así que $\pi_{53} = 4 = \pi_{13}$.
Cambiar $D$ no implica cambiar $\Pi$.

\textbf{Reconstrucción.} \textsc{Imprimir-Camino-Pares}$(\Pi^{(5)},
1, 2)$ sube por el renglón $1$: $\pi_{12} = 3$, $\pi_{13} = 4$,
$\pi_{14} = 5$, $\pi_{15} = 1$, e imprime $1, 5, 4, 3, 2$, con peso
$-4 + 6 - 5 + 4 = 1 = d_{12}$. De $3$ a $5$ sale $3, 2, 4, 1, 5$, con
peso $4 + 1 + 2 - 4 = 3$, y de $2$ a $1$ sale $2, 4, 1$, con peso $3$.
Los $20$ caminos que produce $\Pi^{(5)}$ pesan lo que dice $D^{(5)}$,
y $D^{(5)}$ es la misma $L^{(4)}$ del
ejemplo~\ref{ej:cuadrados-pares}.

\textbf{Cómo crece el árbol desde $1$.} El renglón $1$ de $\Pi^{(k)}$
define $G_{\pi,1}$ en cada etapa. En azul las aristas del árbol, en
naranja el vértice $k$ de la vuelta, y junto a cada vértice
$d^{(k)}_{1j}$:

% Cuatro instantaneas del arbol G_{pi,1}; las aristas grises claras
% son las del grafo que no estan en el arbol.
\begin{center}
\begin{tikzpicture}[x=0.55cm, y=0.6cm]
  \node[vert] (n1) at (-1.9, 0.4) {1};
  \node[vert] (n2) at (0, 1.8) {2};
  \node[vert] (n3) at (1.9, 0.4) {3};
  \node[vert] (n4) at (1.2, -1.8) {4};
  \node[vert] (n5) at (-1.2, -1.8) {5};
  \foreach \a/\b in {1/2, 1/3, 1/5, 2/4, 2/5, 3/2, 4/1, 4/3, 5/4}
    \draw[->, >=stealth, draw=gray!35] (n\a) -- (n\b);
  \foreach \a/\b in {1/2, 1/3, 1/5}
    \draw[dirigidaArbol] (n\a) -- (n\b);
  \node[nota, left=1pt] at (n1.west) {$0$};
  \node[nota, above=1pt] at (n2.north) {$3$};
  \node[nota, right=1pt] at (n3.east) {$8$};
  \node[nota, right=1pt] at (n4.east) {$\infty$};
  \node[nota, left=1pt] at (n5.west) {$-4$};
  \node[font=\small] at (0, -3.2) {$k = 0, 1$};
\end{tikzpicture}
\hspace{0.1cm}
\begin{tikzpicture}[x=0.55cm, y=0.6cm]
  \node[vert] (n1) at (-1.9, 0.4) {1};
  \node[vertMarca] (n2) at (0, 1.8) {2};
  \node[vert] (n3) at (1.9, 0.4) {3};
  \node[vert] (n4) at (1.2, -1.8) {4};
  \node[vert] (n5) at (-1.2, -1.8) {5};
  \foreach \a/\b in {1/2, 1/3, 1/5, 2/4, 2/5, 3/2, 4/1, 4/3, 5/4}
    \draw[->, >=stealth, draw=gray!35] (n\a) -- (n\b);
  \foreach \a/\b in {1/2, 1/3, 1/5, 2/4}
    \draw[dirigidaArbol] (n\a) -- (n\b);
  \node[nota, left=1pt] at (n1.west) {$0$};
  \node[nota, above=1pt] at (n2.north) {$3$};
  \node[nota, right=1pt] at (n3.east) {$8$};
  \node[nota, right=1pt] at (n4.east) {$4$};
  \node[nota, left=1pt] at (n5.west) {$-4$};
  \node[font=\small] at (0, -3.2) {$k = 2, 3$};
\end{tikzpicture}
\hspace{0.1cm}
\begin{tikzpicture}[x=0.55cm, y=0.6cm]
  \node[vert] (n1) at (-1.9, 0.4) {1};
  \node[vert] (n2) at (0, 1.8) {2};
  \node[vert] (n3) at (1.9, 0.4) {3};
  \node[vertMarca] (n4) at (1.2, -1.8) {4};
  \node[vert] (n5) at (-1.2, -1.8) {5};
  \foreach \a/\b in {1/2, 1/3, 1/5, 2/4, 2/5, 3/2, 4/1, 4/3, 5/4}
    \draw[->, >=stealth, draw=gray!35] (n\a) -- (n\b);
  \foreach \a/\b in {1/2, 4/3, 1/5, 2/4}
    \draw[dirigidaArbol] (n\a) -- (n\b);
  \node[nota, left=1pt] at (n1.west) {$0$};
  \node[nota, above=1pt] at (n2.north) {$3$};
  \node[nota, right=1pt] at (n3.east) {$-1$};
  \node[nota, right=1pt] at (n4.east) {$4$};
  \node[nota, left=1pt] at (n5.west) {$-4$};
  \node[font=\small] at (0, -3.2) {$k = 4$};
\end{tikzpicture}
\hspace{0.1cm}
\begin{tikzpicture}[x=0.55cm, y=0.6cm]
  \node[vert] (n1) at (-1.9, 0.4) {1};
  \node[vert] (n2) at (0, 1.8) {2};
  \node[vert] (n3) at (1.9, 0.4) {3};
  \node[vert] (n4) at (1.2, -1.8) {4};
  \node[vertMarca] (n5) at (-1.2, -1.8) {5};
  \foreach \a/\b in {1/2, 1/3, 1/5, 2/4, 2/5, 3/2, 4/1, 4/3, 5/4}
    \draw[->, >=stealth, draw=gray!35] (n\a) -- (n\b);
  \foreach \a/\b in {3/2, 4/3, 1/5, 5/4}
    \draw[dirigidaArbol] (n\a) -- (n\b);
  \node[nota, left=1pt] at (n1.west) {$0$};
  \node[nota, above=1pt] at (n2.north) {$1$};
  \node[nota, right=1pt] at (n3.east) {$-3$};
  \node[nota, right=1pt] at (n4.east) {$2$};
  \node[nota, left=1pt] at (n5.west) {$-4$};
  \node[font=\small] at (0, -3.2) {$k = 5$};
\end{tikzpicture}
\end{center}

Al principio el árbol son las aristas directas desde $1$. En $k = 2$
aparece $2 \to 4$, con $d_{14} = 3 + 1 = 4$. En $k = 4$ la arista
$4 \to 3$ reemplaza a $1 \to 3$, con $d_{13} = 4 - 5 = -1$. En $k = 5$
el árbol se vuelve un solo camino, $1 \to 5 \to 4 \to 3 \to 2$. A
diferencia de Dijkstra, ningún valor es definitivo antes de la última
vuelta: $d_{12}$ valió $3$ desde el principio y cambió hasta $k = 5$.

\begin{ejercicio}
  \label{ej:pi-clase-floyd}
  Calcula $\Pi$ para el ejemplo de Floyd de la clase del 22 de octubre
  y usa \textsc{Imprimir-Camino-Pares} para reconstruir los caminos de
  $3$ a $2$ y de $2$ a $1$.

  \emph{Solución.} Al terminar,
  \[
    \Pi^{(4)} = \begin{pmatrix}
      \paresNulo & 1 & 2 & 3 \\
      4 & \paresNulo & 2 & 3 \\
      4 & 1 & \paresNulo & 3 \\
      4 & 1 & 2 & \paresNulo
    \end{pmatrix} .
  \]
  De $3$ a $2$: $\pi_{32} = 1$, $\pi_{31} = 4$ y $\pi_{34} = 3$, así
  que se imprime $3, 4, 1, 2$, con peso $1 + 2 + 3 = 6 = d_{32}$. De
  $2$ a $1$ sale $2, 3, 4, 1$, con $2 + 1 + 2 = 5 = d_{21}$. Sin $\Pi$,
  el algoritmo solo sabe que $d_{32} = 6$; el camino se tendría que
  leer del dibujo.
\end{ejercicio}

\begin{ejemplo}[Caso borde: un ciclo negativo]
  \label{ej:floyd-ciclo-negativo}
  El grafo del ejemplo~\ref{ej:grafo-pares} con un solo cambio: la
  arista $4 \to 3$ pesa $-8$ en lugar de $-5$, y el ciclo
  $2 \to 4 \to 3 \to 2$ pasa a pesar $1 - 8 + 4 = -3$. La tabla da la
  diagonal de $D^{(k)}$ en cada vuelta de la versión por capas.
  \begin{center}
  \begin{tikzpicture}[baseline=(current bounding box.center)]
    \node[vert] (n1) at (-1.9, 0.4) {1};
    \node[vertMarca] (n2) at (0, 1.8) {2};
    \node[vertMarca] (n3) at (1.9, 0.4) {3};
    \node[vertMarca] (n4) at (1.2, -1.8) {4};
    \node[vert] (n5) at (-1.2, -1.8) {5};
    \foreach \a/\b/\w/\p in {1/2/3/0.5, 1/3/8/0.78, 1/5/-4/0.5,
        2/5/7/0.8, 4/1/2/0.3, 5/4/6/0.5}
      \draw[dirigida] (n\a) -- node[peso, pos=\p] {$\w$} (n\b);
    \foreach \a/\b/\w/\p in {2/4/1/0.72, 4/3/-8/0.5, 3/2/4/0.5}
      \draw[dirigidaResalte] (n\a) -- node[peso, pos=\p] {$\w$} (n\b);
    \node[nota] at (0, -2.7)
      {ciclo $2 \to 4 \to 3 \to 2$: $1 - 8 + 4 = -3$};
  \end{tikzpicture}
  \hspace{1cm}
  \small
  \begin{tabular}{cccccc}
    \toprule
    $k$ & $d_{11}$ & $d_{22}$ & $d_{33}$ & $d_{44}$ & $d_{55}$ \\
    \midrule
    $0$ & $0$ & $0$ & $0$ & $0$ & $0$ \\
    $1$ & $0$ & $0$ & $0$ & $0$ & $0$ \\
    $2$ & $0$ & $0$ & $0$ & $0$ & $0$ \\
    $3$ & $0$ & $0$ & $0$ & $-3$ & $0$ \\
    $4$ & $0$ & $-3$ & $-3$ & $-6$ & $0$ \\
    $5$ & $0$ & $-3$ & $-3$ & $-6$ & $0$ \\
    \bottomrule
  \end{tabular}
  \end{center}
  \begin{itemize}[itemsep=0.1em]
    \item La primera casilla negativa aparece en $k = 3$, no en
          $k = 4$. Para $i = 4$ el ciclo, leído como
          $4 \to 3 \to 2 \to 4$, solo tiene intermedios $3$ y $2$, y ya
          cuenta en $D^{(3)}$. Es la prueba del
          teorema~\ref{teo:floyd-diagonal} vista desde el vértice más
          alto del ciclo, que aquí es su propio extremo.
    \item En $k = 4$ los otros dos vértices del ciclo bajan a $-3$, y
          $d_{44}$ baja a $-6 = d_{44} + d_{44}$: la caminata da dos
          vueltas al ciclo. Con ciclos negativos los números ya no son
          pesos de caminos más cortos.
    \item $d_{11}$ y $d_{55}$ se quedan en $0$, aunque $1$ y $5$
          alcanzan el ciclo y se alcanzan desde él. La diagonal avisa
          que hay un ciclo, no qué pares están afectados: aquí los $25$
          pares tienen en realidad $\delta = -\infty$, mientras
          $D^{(5)}$ muestra números finitos como $d_{12} = -2$. Para
          marcarlos se pone $d_{ij} = -\infty$ cuando existe $k$ con
          $d_{kk} < 0$, $d_{ik} < \infty$ y $d_{kj} < \infty$: $i$
          llega a un ciclo negativo y de ahí a $j$.
    \item La versión en el lugar también lo detecta, con otra diagonal:
          $0, -5, -5, -8, -16$. La proposición~\ref{prop:floyd-lugar}
          deja de valer cuando $d_{kk} < 0$, pero el
          teorema~\ref{teo:floyd-diagonal} no.
  \end{itemize}
\end{ejemplo}

\begin{ojo}[breakable]
  \begin{itemize}[itemsep=0.2em, leftmargin=1.2em]
    \item \textbf{$k$ va afuera.} Con el ciclo de $k$ adentro, en el
          orden $i, j, k$, el algoritmo es otro y falla: en el grafo
          del capítulo da $d_{12} = 3$ y $d_{13} = 8$ en lugar de $1$
          y $-3$, y en $2000$ grafos aleatorios chicos falló en $721$.
          Curiosamente, en esas mismas pruebas repetir tres veces la
          versión mal ordenada dio el resultado correcto todas las
          veces. Es una observación experimental, no una razón para
          programarlo así.
    \item \textbf{Aritmética con $\infty$.} En código, $\infty$ suele
          ser un entero grande. Entonces $\infty + (-5) < \infty$
          inventa caminos que no existen, y $\infty + \infty$ puede
          desbordar. Hay que saltarse la actualización si
          $d_{ik} = \infty$ o $d_{kj} = \infty$.
    \item \textbf{Lazos negativos.} La matriz de pesos pone
          $w_{ii} = 0$ siempre, así que un lazo $(i, i)$ negativo, que
          es un ciclo negativo, desaparece y nadie lo detecta. Si el
          grafo puede tener lazos, hay que poner
          $w_{ii} = \min(0, w(i, i))$; con eso $d_{ii}$ sale negativo.
    \item \textbf{Grafos no dirigidos.} Una arista no dirigida es un
          par de arcos $(u, v)$ y $(v, u)$. Si pesa $-1$, el ciclo
          $u \to v \to u$ pesa $-2$: con una sola arista negativa ya
          hay ciclo negativo, y Floyd-Warshall lo reporta.
    \item \textbf{Los números explotan.} Con un ciclo negativo los
          valores bajan muy rápido: en el grafo dirigido completo de
          $n$ vértices con todos los pesos en $-1$, el mínimo de
          $D^{(n)}$ es $-2^n$ (verificado de $n = 2$ a $8$). Con pesos
          grandes eso desborda. Revisar $d_{kk} < 0$ al terminar cada
          vuelta y detenerse en cuanto aparezca cuesta $\Tht{n}$ extra
          por vuelta.
  \end{itemize}
\end{ojo}

\begin{ejercicio}
  \label{ej:empates-floyd}
  ¿Se puede romper el empate de la regla de $\Pi$ al revés, tomando
  $\pi^{(k-1)}_{kj}$ cuando
  $d^{(k-1)}_{ij} \ge d^{(k-1)}_{ik} + d^{(k-1)}_{kj}$? (Cormen,
  ejercicio 25.2-5.)

  \emph{Solución.} No. Con $k = j$ siempre hay empate, porque
  $d^{(k-1)}_{ij} = d^{(k-1)}_{ij} + d^{(k-1)}_{jj}$, y la regla toma
  $\pi^{(k-1)}_{jj}$, que es nulo: en la vuelta $j$ se borra toda la
  columna $j$ de $\Pi$. En el grafo del capítulo, $16$ de los $20$
  pares con $i \neq j$ terminan en nulo, y el ciclo de peso cero
  $2 \to 4 \to 3 \to 2$ hace además que $\pi_{44}$ deje de ser nulo.
  Excluir $k = i$ y $k = j$ de la regla nueva no basta, porque los
  ciclos de peso cero producen empates genuinos: en $3000$ grafos
  aleatorios con pesos en $\set{0, 1, 2}$ la variante falló en
  $1092$, con algún $\pi_{ii}$ distinto de nulo o con ciclos en
  $G_{\pi,i}$, donde \textsc{Imprimir-Camino-Pares} no termina. Con la
  regla original no falló ninguno.
\end{ejercicio}

% =====================================================================

\subsection*{Cerradura transitiva}

A veces no interesa cuánto cuesta ir de $i$ a $j$, sino si se puede.

\begin{definicion}[Cerradura transitiva]
  \label{def:cerradura-pares}
  La cerradura transitiva de $G$ es $G^* = (V, E^*)$ con
  $E^* = \set{(i, j) : \text{hay un camino de } i \text{ a } j
  \text{ en } G}$. Como el camino vacío cuenta, $E^*$ incluye todos
  los pares $(i, i)$.
\end{definicion}

Es el esquema de Floyd-Warshall con un camino que existe o no existe
en lugar de un peso. Sea $t^{(k)}_{ij} \in \set{0, 1}$ igual a $1$ si
hay un camino de $i$ a $j$ con intermedios en $\set{1, \dots, k}$.
Entonces $t^{(0)}_{ij} = 1$ si $i = j$ o $(i, j) \in E$, y $0$ si no, y
\[
  t^{(k)}_{ij} = t^{(k-1)}_{ij} \lor
                 \bigl( t^{(k-1)}_{ik} \land t^{(k-1)}_{kj} \bigr) ,
\]
que es el lema~\ref{lem:estructura-floyd} con $\lor$ en lugar de
$\min$ y $\land$ en lugar de $+$. Aquí no hace falta ninguna hipótesis
sobre ciclos: si hay alguna caminata de $i$ a $j$, hay un camino. Al
final, $(i, j) \in E^*$ si y solo si $t^{(n)}_{ij} = 1$.

\begin{algorithm}[H]
\caption{\textsc{Cerradura-Transitiva}$(G)$, en el lugar}
\label{alg:cerradura-transitiva}
\begin{algorithmic}[1]
\Procedure{Cerradura-Transitiva}{$G$}
  \For{$i \gets 1$ \Hasta $n$}
    \For{$j \gets 1$ \Hasta $n$}
      \State $t_{ij} \gets 1$ si $i = j$ o $(i, j) \in E$; si no,
             $t_{ij} \gets 0$
    \EndFor
  \EndFor
  \For{$k \gets 1$ \Hasta $n$}
    \For{$i \gets 1$ \Hasta $n$}
      \For{$j \gets 1$ \Hasta $n$}
        \State $t_{ij} \gets t_{ij} \lor (t_{ik} \land t_{kj})$
      \EndFor
    \EndFor
  \EndFor
  \State \Return $T$
\EndProcedure
\end{algorithmic}
\end{algorithm}

\textbf{Costo.} $\Tht{n^3}$, igual que Floyd-Warshall. La versión en
el lugar es correcta por el argumento de la
proposición~\ref{prop:floyd-lugar}, con $t_{kk} = 1$ en el papel de
$d_{kk} = 0$. La ganancia es práctica: operaciones lógicas sobre bits
en lugar de aritmética sobre palabras, y la memoria baja por un factor
del tamaño de la palabra. Si cada renglón de $T$ se guarda como un
vector de bits, la línea 10 para $i$ fija y todo $j$ se vuelve ``si
$t_{ik} = 1$, renglón $i \gets$ renglón $i$ OR renglón $k$'', y con
palabras de $b$ bits el total baja a $\Tht{n^3 / b}$ operaciones de
palabra. Hay dos alternativas. Poner peso $1$ a cada arista y correr
Floyd-Warshall: hay camino si y solo si $d_{ij} < n$. O, en grafos
dispersos, un BFS o DFS desde cada vértice
(sección~\ref{sec:busquedas}), en $\Tht{n(n + m)}$.

El nombre doble del algoritmo viene de aquí. Según las notas del
capítulo 25 de Cormen, Floyd basó su algoritmo de caminos más cortos
en un teorema de Warshall que describe cómo calcular la cerradura
transitiva de matrices booleanas.

\begin{ejemplo}[Cerradura transitiva, figura 25.5 de Cormen]
  \label{ej:cerradura-pares}
  Un grafo de cuatro vértices con aristas $2 \to 3$, $2 \to 4$,
  $3 \to 2$, $4 \to 1$ y $4 \to 3$. Se dibuja con $1$ al centro para
  que $G^*$ quede sin cruces; en naranja punteado, las aristas que
  agrega la cerradura, sin dibujar los pares $(i, i)$.
  \begin{center}
  \begin{tikzpicture}
    \node[vert] (n2) at (1.5, 2.6) {2};
    \node[vert] (n4) at (0, 0) {4};
    \node[vert] (n3) at (3, 0) {3};
    \node[vert] (n1) at (1.5, 0.95) {1};
    \draw[dirigida] (n4) -- (n1);
    \draw[dirigida] (n2) to[bend right=12] (n4);
    \draw[dirigida] (n2) to[bend left=12] (n3);
    \draw[dirigida] (n3) to[bend left=12] (n2);
    \draw[dirigida] (n4) to[bend right=12] (n3);
    \node[font=\small] at (1.5, -0.8) {$G$};
  \end{tikzpicture}
  \hspace{2cm}
  \begin{tikzpicture}
    \node[vert] (n2) at (1.5, 2.6) {2};
    \node[vert] (n4) at (0, 0) {4};
    \node[vert] (n3) at (3, 0) {3};
    \node[vert] (n1) at (1.5, 0.95) {1};
    \draw[dirigida] (n4) -- (n1);
    \draw[dirigida] (n2) to[bend right=12] (n4);
    \draw[dirigida] (n2) to[bend left=12] (n3);
    \draw[dirigida] (n3) to[bend left=12] (n2);
    \draw[dirigida] (n4) to[bend right=12] (n3);
    \draw[dirigidaFuera] (n3) to[bend right=12] (n4);
    \draw[dirigidaFuera] (n4) to[bend right=12] (n2);
    \draw[dirigidaFuera] (n2) -- (n1);
    \draw[dirigidaFuera] (n3) -- (n1);
    \node[font=\small] at (1.5, -0.8) {$G^*$};
  \end{tikzpicture}
  \end{center}
  {\small
  \setlength{\arraycolsep}{3.5pt}
  \[
  T^{(0)} = T^{(1)} = \begin{pmatrix}
    1 & 0 & 0 & 0 \\
    0 & 1 & 1 & 1 \\
    0 & 1 & 1 & 0 \\
    1 & 0 & 1 & 1
  \end{pmatrix}
  \quad
  T^{(2)} = \begin{pmatrix}
    1 & 0 & 0 & 0 \\
    0 & 1 & 1 & 1 \\
    0 & 1 & 1 & \paresCambio{1} \\
    1 & 0 & 1 & 1
  \end{pmatrix}
  \quad
  T^{(3)} = \begin{pmatrix}
    1 & 0 & 0 & 0 \\
    0 & 1 & 1 & 1 \\
    0 & 1 & 1 & 1 \\
    1 & \paresCambio{1} & 1 & 1
  \end{pmatrix}
  \quad
  T^{(4)} = \begin{pmatrix}
    1 & 0 & 0 & 0 \\
    \paresCambio{1} & 1 & 1 & 1 \\
    \paresCambio{1} & 1 & 1 & 1 \\
    1 & 1 & 1 & 1
  \end{pmatrix}
  \]
  }%
  \begin{itemize}[itemsep=0.1em]
    \item $k = 1$: nada cambia. $1$ no tiene aristas de salida, así que
          pasar por $1$ no lleva a ningún lado.
    \item $k = 2$: $t_{34} = t_{32} \land t_{24} = 1$, por
          $3 \to 2 \to 4$.
    \item $k = 3$: $t_{42} = t_{43} \land t_{32} = 1$, por
          $4 \to 3 \to 2$.
    \item $k = 4$: $t_{21} = t_{24} \land t_{41} = 1$ y
          $t_{31} = t_{34} \land t_{41} = 1$; el segundo usa el $t_{34}$
          que apareció en $k = 2$, por $3 \to 2 \to 4 \to 1$.
  \end{itemize}
  Al final $1$ solo se alcanza a sí mismo, y $2$, $3$ y $4$ alcanzan a
  todos.
\end{ejemplo}


% =====================================================================

\subsection*{Johnson: reponderar para poder usar Dijkstra}

Floyd-Warshall cuesta $\Tht{n^3}$ aunque el grafo tenga $m = 2n$
aristas. En un grafo disperso, $n$ corridas de Dijkstra harían del
orden de $n^2 \log n$, pero Dijkstra no admite pesos negativos, y $n$
corridas de Bellman-Ford cuestan $\Ord{n^2 m}$. El algoritmo de
Johnson (1977) junta lo mejor de los dos con una idea que vale más que
el algoritmo mismo: cambiar los pesos sin cambiar los caminos más
cortos.

\begin{definicion}[Reponderación]
  \label{def:reponderacion-johnson}
  Dada una función $h : V \to \R$ cualquiera, la
  \textbf{reponderación} de $w$ por $h$ es
  \[
    \hat{w}(u, v) = w(u, v) + h(u) - h(v) ,
  \]
  y $\hat{\delta}$ denota los pesos de caminos más cortos con
  $\hat{w}$.
\end{definicion}

A $h$ se le suele llamar \textbf{potencial}, y a $\hat{w}$,
\textbf{costos reducidos}; Cormen no le da nombre. El gorro sí es de
Cormen. Algunas notas escriben $w'$, pero aquí $G' = (V', E')$ va a
ser el grafo aumentado con un vértice nuevo, y $w'$ se leería como
``los pesos de $G'$'', que son otra cosa: en $G'$ los pesos siguen
siendo $w$.

\begin{lema}[La reponderación conserva los caminos más cortos]
  \label{lem:reponderacion-johnson}
  Sea $h : V \to \R$ cualquiera y $\hat{w}$ la reponderación por $h$.
  \begin{enumerate}[label=(\alph*), itemsep=0.1em]
    \item Para toda caminata
          $p = \langle v_0, v_1, \dots, v_k \rangle$,
          $\hat{w}(p) = w(p) + h(v_0) - h(v_k)$.
    \item $p$ es un camino más corto de $v_0$ a $v_k$ con $w$ si y
          solo si lo es con $\hat{w}$.
    \item $G$ tiene un ciclo negativo con $w$ si y solo si lo tiene
          con $\hat{w}$.
  \end{enumerate}
\end{lema}

\begin{proof}
  La suma es telescópica:
  \[
    \hat{w}(p)
    = \sum_{i=1}^{k} \bigl( w(v_{i-1}, v_i) + h(v_{i-1}) - h(v_i)
      \bigr)
    = w(p) + \sum_{i=1}^{k} \bigl( h(v_{i-1}) - h(v_i) \bigr)
    = w(p) + h(v_0) - h(v_k) ,
  \]
  porque cada $h(v_i)$ intermedio aparece una vez sumando y una vez
  restando. El término $h(v_0) - h(v_k)$ depende solo de los extremos:
  a todas las caminatas de $v_0$ a $v_k$ se les suma la misma
  cantidad, así que una pesa menos que otra con $w$ si y solo si pesa
  menos con $\hat{w}$, y eso es (b). Para un ciclo, $v_0 = v_k$ y (a)
  da $\hat{w}(c) = w(c)$, que es (c).
\end{proof}

\begin{corolario}
  \label{cor:deshacer-johnson}
  Si $G$ no tiene ciclos negativos,
  $\delta(u, v) = \hat{\delta}(u, v) + h(v) - h(u)$ para todo par, con
  $\infty$ en los dos lados cuando no hay camino.
\end{corolario}

Basta aplicar (a) a un camino más corto $p$ de $u$ a $v$, que por (b)
también lo es con $\hat{w}$. Para no equivocarse con el signo: con una
sola arista, $\hat{w} = w + h(u) - h(v)$, y despejando,
$w = \hat{w} + h(v) - h(u)$.

Lo que el lema no dice es que se conserve el orden entre destinos
distintos, porque cada destino tiene su propia $h(v)$. Dijkstra con
$\hat{w}$ puede extraer los vértices en un orden que no tiene nada que
ver con las distancias reales, y aun así las distancias salen bien; el
ejemplo~\ref{ej:johnson-cormen} lo muestra.

\subsubsection*{Cómo escoger $h$}

El lema vale para cualquier $h$. La gracia es encontrar una que deje
todos los pesos no negativos, para poder correr Dijkstra. Se agrega un
vértice nuevo $s$ con una arista de peso $0$ hacia cada vértice: el
\textbf{grafo aumentado} es $G' = (V', E')$, con $V' = V \cup \set{s}$,
$E' = E \cup \set{(s, v) : v \in V}$ y $w(s, v) = 0$. A $s$ no entra
ninguna arista, así que ningún ciclo pasa por $s$: $G'$ tiene un ciclo
negativo si y solo si $G$ lo tiene, y ningún camino más corto entre
vértices de $G$ pasa por $s$.

\begin{proposicion}[El potencial de Johnson]
  \label{prop:johnson-no-negativos}
  Si $G'$ no tiene ciclos negativos y $h(v) = \delta(s, v)$ para todo
  $v \in V'$, entonces $\hat{w}(u, v) \ge 0$ para toda
  $(u, v) \in E'$.
\end{proposicion}

\begin{proof}
  Todo $v$ es alcanzable desde $s$ por la arista $(s, v)$, así que
  $h(v)$ es finito, y de hecho $h(v) \le 0$. Por la desigualdad del
  triángulo (proposición~\ref{prop:triangulo}),
  $\delta(s, v) \le \delta(s, u) + w(u, v)$ para toda arista $(u, v)$,
  es decir, $h(v) \le h(u) + w(u, v)$, que reescrita es
  $\hat{w}(u, v) = w(u, v) + h(u) - h(v) \ge 0$. Las aristas de $s$
  quedan con $\hat{w}(s, v) = 0 + 0 - h(v) \ge 0$.
\end{proof}

La prueba dice algo más: $\hat{w}(u, v) = 0$ exactamente cuando la
desigualdad del triángulo es una igualdad, es decir, cuando la arista
está sobre algún camino más corto desde $s$. Y no hace falta que $h$
sea $\delta(s, \cdot)$: sirve cualquier $h$ con
$h(v) - h(u) \le w(u, v)$ en cada arista, que es una solución de un
sistema de restricciones de diferencia (Cormen, sección 24.4), y ese
sistema se resuelve con la misma construcción del vértice nuevo.

\begin{algorithm}[H]
\caption{\textsc{Johnson}$(G, w)$}
\label{alg:johnson}
\begin{algorithmic}[1]
\Procedure{Johnson}{$G, w$}
  \State $V' \gets V \cup \set{s}$, \quad
         $E' \gets E \cup \set{(s, v) : v \in V}$
         \Comment{$s$ es un vértice nuevo}
  \State $w(s, v) \gets 0$ para todo $v \in V$
  \If{\Call{Bellman-Ford}{$G', w, s$} devuelve falso}
    \State \Return ``$G$ tiene un ciclo de peso negativo''
  \EndIf
  \State $h(v) \gets d[v]$ para todo $v \in V$
    \Comment{$d[v] = \delta(s, v)$}
  \State $\hat{w}(u, v) \gets w(u, v) + h(u) - h(v)$ para toda
         $(u, v) \in E$ \Comment{$\ge 0$}
  \ForAll{$u \in V$}
    \State \Call{Dijkstra}{$G, \hat{w}, u$}
      \Comment{deja $d[v] = \hat{\delta}(u, v)$}
    \ForAll{$v \in V$}
      \State $d_{uv} \gets d[v] + h(v) - h(u)$
        \Comment{si $d[v] = \infty$, queda $\infty$}
    \EndFor
  \EndFor
  \State \Return $D$
\EndProcedure
\end{algorithmic}
\end{algorithm}

Bellman-Ford y Dijkstra son los algoritmos~\ref{alg:bellman-ford}
y~\ref{alg:dijkstra-cormen}; $d[v]$ es el arreglo de distancias que
dejan, y $d_{uv}$ la casilla de la matriz de salida.
\begin{itemize}[itemsep=0.1em]
  \item Línea 4: la única corrida de Bellman-Ford hace dos trabajos.
        Calcula $h$ y detecta ciclos negativos en todo $G$, porque
        desde $s$ se alcanza todo. Si se corriera desde un vértice de
        $G$, solo vería los ciclos alcanzables desde él y dejaría
        $h = \infty$ en los vértices que no alcanza
        (ejercicio~\ref{ej:fuente-johnson}).
  \item Línea 7: se copia $d[v]$ en $h(v)$ porque cada Dijkstra de la
        línea 10 sobrescribe $d$. Olvidar esa copia es el error de
        implementación más común.
  \item Línea 10: Dijkstra corre sobre $G$, no sobre $G'$. Daría lo
        mismo, porque desde $u \in V$ no se llega a $s$, pero no hace
        falta.
  \item Línea 12: si $d[v] = \infty$, la casilla queda en $\infty$. Con
        enteros y un centinela grande, sumarle $h(v) - h(u)$ desborda o
        da un número finito sin sentido; hay que revisar antes.
  \item Los predecesores salen gratis, aunque Cormen no los calcula.
        Por el lema~\ref{lem:reponderacion-johnson}, el árbol que
        Dijkstra construye con $\hat{w}$ desde $u$ es también árbol de
        caminos más cortos con $w$, y el renglón $u$ de $\Pi$ es el
        arreglo $\pi$ de esa corrida, sin corrección.
\end{itemize}

\begin{teorema}[Correctitud de Johnson]
  \label{teo:johnson-correcto}
  Si $G$ tiene un ciclo negativo, Johnson lo reporta. Si no,
  devuelve $D$ con $d_{uv} = \delta(u, v)$ para todo par.
\end{teorema}

\begin{proof}
  Bellman-Ford devuelve falso si y solo si hay un ciclo negativo
  alcanzable desde la fuente (sección~\ref{sec:caminos-fuente};
  Cormen, teorema 24.4). En $G'$ todo es alcanzable desde $s$, así que
  devuelve falso si y solo si $G'$ tiene un ciclo negativo, y eso pasa
  si y solo si $G$ lo tiene. Si devuelve verdadero, deja
  $d[v] = \delta(s, v) = h(v)$; por la
  proposición~\ref{prop:johnson-no-negativos}, $\hat{w} \ge 0$, así
  que cada Dijkstra calcula $\hat{\delta}(u, \cdot)$ correctamente, y
  el corolario~\ref{cor:deshacer-johnson} lo convierte en
  $\delta(u, \cdot)$.
\end{proof}

\textbf{Costo.} Construir $G'$ (líneas 2 y 3) cuesta $\Tht{n + m}$:
copiar las listas y agregar $n$ aristas. Bellman-Ford en $G'$
(línea 4) hace $n$ pasadas sobre $m + n$ aristas,
$\Ord{(n + 1)(m + n)} = \Ord{nm + n^2}$. Copiar $h$ (línea 7) cuesta
$\Tht{n}$, reponderar (línea 8) $\Tht{m}$, y deshacer la reponderación
(líneas 11 a 13) $\Tht{n^2}$, una suma por casilla de $D$. Falta lo
que domina, las $n$ corridas de Dijkstra de las líneas 9 y 10, y eso
depende de la cola de prioridad.
Una corrida de Dijkstra hace $n$ inserciones, $n$ extracciones del
mínimo y a lo más $m$ decrementos de llave, uno por relajación
exitosa. Con montículo de Fibonacci eso es
$n \cdot \Ord{1} + n \cdot \Ord{\log n} + m \cdot \Ord{1}
= \Ord{n \log n + m}$ por corrida y $\Ord{n^2 \log n + nm}$ las $n$
corridas. Bellman-Ford queda absorbido: su $nm$ por el de las
corridas y su $n^2$ por el $n^2 \log n$. El total es
\[
  \Ord{nm + n^2} + \Ord{n^2 \log n + nm} + \Tht{n^2}
  = \Ord{n^2 \log n + nm} .
\]
Con el montículo binario de la libreta son $\Ord{(n + m)\log n}$ por
corrida y $\Ord{n(n + m)\log n}$ en total, que Cormen escribe
$\Ord{nm \log n}$. Con un arreglo, $\Ord{n^3}$, y entonces no tiene
caso: se iguala a Floyd-Warshall con más trabajo. En los tres casos el
costo es el de $n$ veces Dijkstra con pesos no negativos; admitir
pesos negativos cuesta una corrida de Bellman-Ford, que no domina.

La memoria es $\Tht{n + m}$ de trabajo más los $\Tht{n^2}$ de la
salida. A diferencia de Floyd-Warshall, Johnson produce $D$ renglón
por renglón: si solo se necesitan algunos renglones, nunca hace falta
la matriz completa.

\begin{ejemplo}[El contraejemplo de Dijkstra, resuelto]
  \label{ej:johnson-libreta}
  La clase del 22 de octubre muestra que Dijkstra falla con las
  aristas $A \to B$ de peso $2$, $A \to C$ de peso $3$ y $C \to B$ de
  peso $-2$: desde A fija B con $d = 2$, pero A, C, B cuesta $1$. Es
  el ejemplo natural para Johnson.

  \textbf{Bellman-Ford en $G'$.} Con las aristas en el orden
  $(s, A)$, $(s, B)$, $(s, C)$, $(A, B)$, $(A, C)$, $(C, B)$, la
  primera pasada pone $d = 0$ en A, B y C por las aristas de $s$, y
  luego $(C, B)$ baja B a $0 - 2 = -2$. Las pasadas 2 y 3 no cambian
  nada y la verificación devuelve verdadero: $h(A) = 0$, $h(B) = -2$
  y $h(C) = 0$.

  \textbf{Reponderación.} $\hat{w}(A, B) = 2 + 0 - (-2) = 4$,
  $\hat{w}(A, C) = 3 + 0 - 0 = 3$ y
  $\hat{w}(C, B) = -2 + 0 - (-2) = 0$.

  \textbf{Dijkstra desde A con $\hat{w}$}, en el formato de la clase;
  en negritas, las $d$ ya fijas.
  \begin{center}
  \small
  \begin{tabular}{cccccl}
    \toprule
    \textbf{Paso} & \textbf{Sale} & A & B & C
      & \textbf{Relajaciones} \\
    \midrule
    0 & & $0$ & $\infty$ & $\infty$ & \\
    1 & A & $\mathbf{0}$ & $4$ & $3$ & B y C desde A \\
    2 & C & $\mathbf{0}$ & $3$ & $\mathbf{3}$ & B: $3 + 0 = 3 < 4$ \\
    3 & B & $\mathbf{0}$ & $\mathbf{3}$ & $\mathbf{3}$ & \\
    \midrule
    \multicolumn{2}{l}{$\delta(A, v)$} & $0$ & $3 + (-2) - 0 = 1$
      & $3 + 0 - 0 = 3$ & \\
    \bottomrule
  \end{tabular}
  \end{center}
  B empieza en $4$ por la arista directa y baja a $3$ al pasar por C,
  como debe. Desde C, $\hat{\delta}(C, B) = 0$ da
  $\delta(C, B) = 0 + (-2) - 0 = -2$; desde B no se alcanza nada y los
  $\infty$ se quedan. Las tres instantáneas son el grafo original, el
  aumentado con $h$ y el reponderado con el árbol de Dijkstra desde A,
  con $\hat{\delta}(A, v) / \delta(A, v)$ junto a cada vértice:
  \begin{center}
  \begin{tikzpicture}[baseline=(A.base)]
    \node[vert] (A) at (0, 0) {A};
    \node[vert] (B) at (2.8, 0) {B};
    \node[vert] (C) at (1.4, -1.5) {C};
    \draw[dirigida] (A) -- node[peso] {$2$} (B);
    \draw[dirigida] (A) -- node[peso] {$3$} (C);
    \draw[dirigida] (C) -- node[peso] {$-2$} (B);
    \node[font=\small] at (1.4, -3.75) {(a) $G$ con $w$};
  \end{tikzpicture}
  \hspace{0.9cm}
  \begin{tikzpicture}[baseline=(A.base)]
    \node[vert] (A) at (0, 0) {A};
    \node[vert] (B) at (2.8, 0) {B};
    \node[vert] (C) at (1.4, -1.5) {C};
    \node[vertNegro] (s) at (1.4, -2.9) {$s$};
    \draw[dirigida] (A) -- node[peso] {$2$} (B);
    \draw[dirigida] (A) -- node[peso] {$3$} (C);
    \draw[dirigida] (C) -- node[peso] {$-2$} (B);
    \draw[dirigida] (s) to[out=180, in=-100] node[peso] {$0$} (A);
    \draw[dirigida] (s) to[out=0, in=-80] node[peso] {$0$} (B);
    \draw[dirigida] (s) -- node[peso] {$0$} (C);
    \node[nota, above=1pt] at (A.north) {$h = 0$};
    \node[nota, above=1pt] at (B.north) {$h = -2$};
    \node[nota, right=1pt] at (C.east) {$h = 0$};
    \node[font=\small] at (1.4, -3.75) {(b) $G'$ y $h = \delta(s, \cdot)$};
  \end{tikzpicture}
  \hspace{0.9cm}
  \begin{tikzpicture}[baseline=(A.base)]
    \node[vertNegro] (A) at (0, 0) {A};
    \node[vert] (B) at (2.8, 0) {B};
    \node[vert] (C) at (1.4, -1.5) {C};
    \draw[dirigida] (A) -- node[peso] {$4$} (B);
    \draw[dirigidaArbol] (A) -- node[peso] {$3$} (C);
    \draw[dirigidaArbol] (C) -- node[peso] {$0$} (B);
    \node[nota, above=1pt] at (A.north) {$0/0$};
    \node[nota, above=1pt] at (B.north) {$3/1$};
    \node[nota, below=1pt] at (C.south) {$3/3$};
    \node[font=\small] at (1.4, -3.75) {(c) $\hat{w}$, Dijkstra desde A};
  \end{tikzpicture}
  \end{center}
  Floyd-Warshall, con $A = 1$, $B = 2$ y $C = 3$, llega a lo mismo:
  $d_{AB}$ baja de $2$ a $1$ en la vuelta $k = 3$, la de C.
\end{ejemplo}

\begin{ojo}
  Sumar una constante no sirve (Cormen, ejercicio 25.3-4). La
  tentación es restarle a cada arista el peso mínimo, $w^* = -2$, para
  que todas queden no negativas: $A \to B$ pasa a $4$, $A \to C$ a $5$
  y $C \to B$ a $0$. Dijkstra desde A elige ahora la arista directa,
  $4$ contra $5 + 0 = 5$, y responde que el mejor camino a B es
  $A \to B$, de peso original $2$. Está mal. Sumarle $\abs{w^*}$ a cada
  arista le suma $k\abs{w^*}$ a un camino de $k$ aristas, y castiga a
  los largos: A, C, B recibe $+4$ y la arista directa solo $+2$. La
  corrección de Johnson depende solo de los extremos, no de cuántas
  aristas tiene el camino.
  \begin{center}
  \begin{tikzpicture}
    \node[vertNegro] (A) at (0, 0) {A};
    \node[vert] (B) at (2.8, 0) {B};
    \node[vert] (C) at (1.4, -1.5) {C};
    \draw[dirigidaArbol] (A) -- node[peso] {$4$} (B);
    \draw[dirigidaArbol] (A) -- node[peso] {$5$} (C);
    \draw[dirigida] (C) -- node[peso] {$0$} (B);
    \node[nota, above=1pt] at (A.north) {$0$};
    \node[nota, above=1pt] at (B.north) {$4$};
    \node[nota, below=1pt] at (C.south) {$5$};
    \node[font=\small, anchor=west] at (3.6, -0.75)
      {$w(u, v) - w^*$, con $w^* = -2$: árbol equivocado};
  \end{tikzpicture}
  \end{center}
\end{ojo}

\begin{ejemplo}[La figura 25.6 de Cormen, completa]
  \label{ej:johnson-cormen}
  Es la corrida de Johnson sobre el grafo del
  ejemplo~\ref{ej:grafo-pares}.

  \textbf{Paso 1: Bellman-Ford en $G'$}, con las cinco aristas
  $(s, v)$ primero y luego las de $E$ en orden lexicográfico. Converge
  en dos pasadas, aunque hace las $\abs{V'} - 1 = 5$:
  \begin{center}
  \small
  \begin{tabular}{cccccccl}
    \toprule
    \textbf{Pasada} & $s$ & 1 & 2 & 3 & 4 & 5 & \textbf{Cambios} \\
    \midrule
    0 & $0$ & $\infty$ & $\infty$ & $\infty$ & $\infty$ & $\infty$ & \\
    1 & $0$ & $0$ & $0$ & $-5$ & $0$ & $-4$
      & las cinco $(s, v)$; $(1, 5)$: $-4$; $(4, 3)$: $-5$ \\
    2 & $0$ & $0$ & $-1$ & $-5$ & $0$ & $-4$ & $(3, 2)$: $-5 + 4 = -1$ \\
    3 a 5 & $0$ & $0$ & $-1$ & $-5$ & $0$ & $-4$ & ninguno \\
    \midrule
    $h$ & $0$ & $0$ & $-1$ & $-5$ & $0$ & $-4$
      & verificación: verdadero \\
    \bottomrule
  \end{tabular}
  \end{center}
  En la pasada 1, $(1, 5)$ con $d[1] = 0$ da $d[5] = -4$, y $(4, 3)$
  da $d[3] = -5$; en la pasada 2, $(3, 2)$ aprovecha ese $-5$:
  $d[2] = -5 + 4 = -1$. Todos los $h$ son $\le 0$, como debe ser.

  \textbf{Paso 2: reponderar}, con
  $\hat{w}(u, v) = w(u, v) + h(u) - h(v)$:
  \begin{center}
  \small
  \setlength{\tabcolsep}{4.5pt}
  \begin{tabular}{lccccccccc}
    \toprule
    $(u, v)$ & $(1, 2)$ & $(1, 3)$ & $(1, 5)$ & $(2, 4)$ & $(2, 5)$
      & $(3, 2)$ & $(4, 1)$ & $(4, 3)$ & $(5, 4)$ \\
    \midrule
    $w(u, v)$ & $3$ & $8$ & $-4$ & $1$ & $7$ & $4$ & $2$ & $-5$ & $6$ \\
    $h(u)$ & $0$ & $0$ & $0$ & $-1$ & $-1$ & $-5$ & $0$ & $0$ & $-4$ \\
    $h(v)$ & $-1$ & $-5$ & $-4$ & $0$ & $-4$ & $-1$ & $0$ & $-5$ & $0$ \\
    \midrule
    $\hat{w}(u, v)$ & $4$ & $13$ & $0$ & $0$ & $10$ & $0$ & $2$ & $0$
      & $2$ \\
    \bottomrule
  \end{tabular}
  \end{center}
  Las aristas de $s$ quedan con $\hat{w}(s, v) = -h(v)$, o sea
  $0, 1, 5, 0, 4$. Las de peso $0$ son las que están sobre algún
  camino más corto desde $s$. Los dos grafos, (a) con $w$ y $h$, y (b)
  con $\hat{w}$:
  \begin{center}
  \begin{tikzpicture}[scale=0.82, transform shape]
    \node[vertNegro] (s) at (-4.2, 0.4) {$s$};
    \node[vert] (n1) at (-1.9, 0.4) {1};
    \node[vert] (n2) at (0, 1.8) {2};
    \node[vert] (n3) at (1.9, 0.4) {3};
    \node[vert] (n4) at (1.2, -1.8) {4};
    \node[vert] (n5) at (-1.2, -1.8) {5};
    \foreach \a/\b/\w/\p in {1/2/3/0.5, 1/3/8/0.78, 1/5/-4/0.5,
        2/4/1/0.72, 2/5/7/0.8, 3/2/4/0.5, 4/1/2/0.3, 4/3/-5/0.5,
        5/4/6/0.5}
      \draw[dirigida] (n\a) -- node[peso, pos=\p] {$\w$} (n\b);
    \draw[dirigida] (s) -- node[peso] {$0$} (n1);
    \draw[dirigida] (s) to[out=75, in=180] node[peso] {$0$} (n2);
    \draw[dirigida] (s) to[out=85, in=90, looseness=1.25]
      node[peso] {$0$} (n3);
    \draw[dirigida] (s) to[out=-75, in=180] node[peso] {$0$} (n5);
    \draw[dirigida] (s) to[out=-85, in=-90, looseness=1.25]
      node[peso] {$0$} (n4);
    \node[nota, right=2pt] at (n2.north east) {$h = -1$};
    \node[nota, right=1pt] at (n3.east) {$h = -5$};
    \node[nota, right=1pt] at (n4.east) {$h = 0$};
    \node[nota, anchor=north east] at (n5.south west) {$h = -4$};
    \node[nota, anchor=south east] at (n1.north) {$h = 0$};
    \node[font=\normalsize] at (-1, -3.9) {(a) $G'$ con los pesos $w$};
  \end{tikzpicture}
  \hspace{0.3cm}
  \begin{tikzpicture}[scale=0.82, transform shape]
    \node[vertNegro] (s) at (-4.2, 0.4) {$s$};
    \node[vert] (n1) at (-1.9, 0.4) {1};
    \node[vert] (n2) at (0, 1.8) {2};
    \node[vert] (n3) at (1.9, 0.4) {3};
    \node[vert] (n4) at (1.2, -1.8) {4};
    \node[vert] (n5) at (-1.2, -1.8) {5};
    \foreach \a/\b/\w/\p in {1/2/4/0.5, 1/3/13/0.78, 1/5/0/0.5,
        2/4/0/0.72, 2/5/10/0.8, 3/2/0/0.5, 4/1/2/0.3, 4/3/0/0.5,
        5/4/2/0.5}
      \draw[dirigida] (n\a) -- node[peso, pos=\p] {$\w$} (n\b);
    \draw[dirigida] (s) -- node[peso] {$0$} (n1);
    \draw[dirigida] (s) to[out=75, in=180] node[peso] {$1$} (n2);
    \draw[dirigida] (s) to[out=85, in=90, looseness=1.25]
      node[peso] {$5$} (n3);
    \draw[dirigida] (s) to[out=-75, in=180] node[peso] {$4$} (n5);
    \draw[dirigida] (s) to[out=-85, in=-90, looseness=1.25]
      node[peso] {$0$} (n4);
    \node[font=\normalsize] at (-1, -3.9) {(b) $G'$ con $\hat{w}$};
  \end{tikzpicture}
  \end{center}

  \textbf{Paso 3: Dijkstra desde cada vértice con $\hat{w}$.} Desde
  el vértice 1:
  \begin{center}
  \small
  \begin{tabular}{cccccccl}
    \toprule
    \textbf{Paso} & \textbf{Sale} & 1 & 2 & 3 & 4 & 5
      & \textbf{Relajaciones} \\
    \midrule
    0 & & $0$ & $\infty$ & $\infty$ & $\infty$ & $\infty$ & \\
    1 & 1 & $\mathbf{0}$ & $4$ & $13$ & $\infty$ & $0$
      & 2, 3, 5 desde 1 \\
    2 & 5 & $\mathbf{0}$ & $4$ & $13$ & $2$ & $\mathbf{0}$
      & 4: $0 + 2$ \\
    3 & 4 & $\mathbf{0}$ & $4$ & $2$ & $\mathbf{2}$ & $\mathbf{0}$
      & 3: $2 + 0 = 2 < 13$ \\
    4 & 3 & $\mathbf{0}$ & $2$ & $\mathbf{2}$ & $\mathbf{2}$
      & $\mathbf{0}$ & 2: $2 + 0 = 2 < 4$ \\
    5 & 2 & $\mathbf{0}$ & $\mathbf{2}$ & $\mathbf{2}$ & $\mathbf{2}$
      & $\mathbf{0}$ & ninguna \\
    \midrule
    \multicolumn{2}{l}{$h(v) - h(1)$} & $0$ & $-1$ & $-5$ & $0$ & $-4$
      & \\
    \multicolumn{2}{l}{$\delta(1, v)$} & $0$ & $1$ & $-3$ & $2$ & $-4$
      & \\
    \bottomrule
  \end{tabular}
  \end{center}
  El vértice 2 empieza en $4$ por la arista directa y baja a $2$ por
  el camino $1, 5, 4, 3, 2$, que pesa $0 + 2 + 0 + 0 = 2$ con
  $\hat{w}$ y $-4 + 6 - 5 + 4 = 1$ con $w$. El último renglón aplica el
  corolario~\ref{cor:deshacer-johnson}; por ejemplo,
  $\delta(1, 3) = 2 + (-5) - 0 = -3$.

  Las cinco corridas, con el árbol de Dijkstra en azul, la fuente en
  negro y $\hat{\delta}(u, v) / \delta(u, v)$ junto a cada vértice:
  \begin{center}
  \begin{tikzpicture}[x=0.75cm, y=0.8cm]
    \node[vertNegro] (n1) at (-1.9, 0.4) {1};
    \node[vert] (n2) at (0, 1.8) {2};
    \node[vert] (n3) at (1.9, 0.4) {3};
    \node[vert] (n4) at (1.2, -1.8) {4};
    \node[vert] (n5) at (-1.2, -1.8) {5};
    \foreach \a/\b/\w/\p in {1/2/4/0.5, 1/3/13/0.78, 2/4/0/0.72,
        2/5/10/0.8, 4/1/2/0.3}
      \draw[dirigida] (n\a) -- node[peso, pos=\p] {$\w$} (n\b);
    \foreach \a/\b/\w/\p in {1/5/0/0.5, 5/4/2/0.5, 4/3/0/0.5,
        3/2/0/0.5}
      \draw[dirigidaArbol] (n\a) -- node[peso, pos=\p] {$\w$} (n\b);
    \node[nota, left=1pt] at (n1.west) {$0/0$};
    \node[nota, above=1pt] at (n2.north) {$2/1$};
    \node[nota, right=1pt] at (n3.east) {$2/{-3}$};
    \node[nota, right=1pt] at (n4.east) {$2/2$};
    \node[nota, left=1pt] at (n5.west) {$0/{-4}$};
    \node[font=\small] at (0, -2.9) {(c) fuente $1$};
  \end{tikzpicture}
  \hspace{0.2cm}
  \begin{tikzpicture}[x=0.75cm, y=0.8cm]
    \node[vert] (n1) at (-1.9, 0.4) {1};
    \node[vertNegro] (n2) at (0, 1.8) {2};
    \node[vert] (n3) at (1.9, 0.4) {3};
    \node[vert] (n4) at (1.2, -1.8) {4};
    \node[vert] (n5) at (-1.2, -1.8) {5};
    \foreach \a/\b/\w/\p in {1/2/4/0.5, 1/3/13/0.78, 2/5/10/0.8,
        3/2/0/0.5, 5/4/2/0.5}
      \draw[dirigida] (n\a) -- node[peso, pos=\p] {$\w$} (n\b);
    \foreach \a/\b/\w/\p in {2/4/0/0.72, 4/1/2/0.3, 4/3/0/0.5,
        1/5/0/0.5}
      \draw[dirigidaArbol] (n\a) -- node[peso, pos=\p] {$\w$} (n\b);
    \node[nota, left=1pt] at (n1.west) {$2/3$};
    \node[nota, above=1pt] at (n2.north) {$0/0$};
    \node[nota, right=1pt] at (n3.east) {$0/{-4}$};
    \node[nota, right=1pt] at (n4.east) {$0/1$};
    \node[nota, left=1pt] at (n5.west) {$2/{-1}$};
    \node[font=\small] at (0, -2.9) {(d) fuente $2$};
  \end{tikzpicture}
  \hspace{0.2cm}
  \begin{tikzpicture}[x=0.75cm, y=0.8cm]
    \node[vert] (n1) at (-1.9, 0.4) {1};
    \node[vert] (n2) at (0, 1.8) {2};
    \node[vertNegro] (n3) at (1.9, 0.4) {3};
    \node[vert] (n4) at (1.2, -1.8) {4};
    \node[vert] (n5) at (-1.2, -1.8) {5};
    \foreach \a/\b/\w/\p in {1/2/4/0.5, 1/3/13/0.78, 2/5/10/0.8,
        4/3/0/0.5, 5/4/2/0.5}
      \draw[dirigida] (n\a) -- node[peso, pos=\p] {$\w$} (n\b);
    \foreach \a/\b/\w/\p in {3/2/0/0.5, 2/4/0/0.72, 4/1/2/0.3,
        1/5/0/0.5}
      \draw[dirigidaArbol] (n\a) -- node[peso, pos=\p] {$\w$} (n\b);
    \node[nota, left=1pt] at (n1.west) {$2/7$};
    \node[nota, above=1pt] at (n2.north) {$0/4$};
    \node[nota, right=1pt] at (n3.east) {$0/0$};
    \node[nota, right=1pt] at (n4.east) {$0/5$};
    \node[nota, left=1pt] at (n5.west) {$2/3$};
    \node[font=\small] at (0, -2.9) {(e) fuente $3$};
  \end{tikzpicture}

  \vspace{0.3cm}
  \begin{tikzpicture}[x=0.75cm, y=0.8cm]
    \node[vert] (n1) at (-1.9, 0.4) {1};
    \node[vert] (n2) at (0, 1.8) {2};
    \node[vert] (n3) at (1.9, 0.4) {3};
    \node[vertNegro] (n4) at (1.2, -1.8) {4};
    \node[vert] (n5) at (-1.2, -1.8) {5};
    \foreach \a/\b/\w/\p in {1/2/4/0.5, 1/3/13/0.78, 2/4/0/0.72,
        2/5/10/0.8, 5/4/2/0.5}
      \draw[dirigida] (n\a) -- node[peso, pos=\p] {$\w$} (n\b);
    \foreach \a/\b/\w/\p in {4/1/2/0.3, 4/3/0/0.5, 3/2/0/0.5,
        1/5/0/0.5}
      \draw[dirigidaArbol] (n\a) -- node[peso, pos=\p] {$\w$} (n\b);
    \node[nota, left=1pt] at (n1.west) {$2/2$};
    \node[nota, above=1pt] at (n2.north) {$0/{-1}$};
    \node[nota, right=1pt] at (n3.east) {$0/{-5}$};
    \node[nota, right=1pt] at (n4.east) {$0/0$};
    \node[nota, left=1pt] at (n5.west) {$2/{-2}$};
    \node[font=\small] at (0, -2.9) {(f) fuente $4$};
  \end{tikzpicture}
  \hspace{0.2cm}
  \begin{tikzpicture}[x=0.75cm, y=0.8cm]
    \node[vert] (n1) at (-1.9, 0.4) {1};
    \node[vert] (n2) at (0, 1.8) {2};
    \node[vert] (n3) at (1.9, 0.4) {3};
    \node[vert] (n4) at (1.2, -1.8) {4};
    \node[vertNegro] (n5) at (-1.2, -1.8) {5};
    \foreach \a/\b/\w/\p in {1/2/4/0.5, 1/3/13/0.78, 1/5/0/0.5,
        2/4/0/0.72, 2/5/10/0.8}
      \draw[dirigida] (n\a) -- node[peso, pos=\p] {$\w$} (n\b);
    \foreach \a/\b/\w/\p in {5/4/2/0.5, 4/1/2/0.3, 4/3/0/0.5,
        3/2/0/0.5}
      \draw[dirigidaArbol] (n\a) -- node[peso, pos=\p] {$\w$} (n\b);
    \node[nota, left=1pt] at (n1.west) {$4/8$};
    \node[nota, above=1pt] at (n2.north) {$2/5$};
    \node[nota, right=1pt] at (n3.east) {$2/1$};
    \node[nota, right=1pt] at (n4.east) {$2/6$};
    \node[nota, left=1pt] at (n5.west) {$0/0$};
    \node[font=\small] at (0, -2.9) {(g) fuente $5$};
  \end{tikzpicture}
  \end{center}
  Desde el vértice 5, Dijkstra extrae 4, 3 y 2 empatados en
  $\hat{\delta} = 2$, pero sus distancias reales son $6$, $1$ y $5$: el
  orden de extracción no dice nada sobre $w$. Juntando los renglones,
  \[
    D = \begin{pmatrix}
      0 & 1 & -3 & 2 & -4 \\
      3 & 0 & -4 & 1 & -1 \\
      7 & 4 & 0 & 5 & 3 \\
      2 & -1 & -5 & 0 & -2 \\
      8 & 5 & 1 & 6 & 0
    \end{pmatrix} ,
  \]
  que es la $L^{(4)}$ del ejemplo~\ref{ej:cuadrados-pares} y la
  $D^{(5)}$ del ejemplo~\ref{ej:floyd-pares}.
\end{ejemplo}

\begin{ejemplo}[Caso borde: Johnson ante un ciclo negativo]
  \label{ej:johnson-ciclo-negativo}
  Al grafo del ejemplo~\ref{ej:johnson-libreta} se le agrega la arista
  $B \to A$ con peso $-2$. El ciclo A, C, B, A pesa $3 - 2 - 2 = -1$;
  el ciclo A, B, A pesa $0$ y no estorba.
  \begin{center}
  \begin{tikzpicture}
    \node[vertMarca] (A) at (0, 0) {A};
    \node[vertMarca] (B) at (2.8, 0) {B};
    \node[vertMarca] (C) at (1.4, -1.5) {C};
    \node[vertNegro] (s) at (1.4, -2.9) {$s$};
    \draw[dirigida] (A) to[bend left=12] node[peso] {$2$} (B);
    \draw[dirigidaResalte] (B) to[bend left=12] node[peso] {$-2$} (A);
    \draw[dirigidaResalte] (A) -- node[peso] {$3$} (C);
    \draw[dirigidaResalte] (C) -- node[peso] {$-2$} (B);
    \draw[dirigida] (s) to[out=180, in=-100] node[peso] {$0$} (A);
    \draw[dirigida] (s) to[out=0, in=-80] node[peso] {$0$} (B);
    \draw[dirigida] (s) -- node[peso] {$0$} (C);
    \node[nota, align=left, anchor=west] at (3.5, -1.2)
      {ciclo A, C, B, A:\\ $3 - 2 - 2 = -1$};
  \end{tikzpicture}
  \end{center}
  Bellman-Ford en $G'$, con las aristas en el orden $(s, A)$,
  $(s, B)$, $(s, C)$, $(A, B)$, $(A, C)$, $(C, B)$, $(B, A)$:
  \begin{center}
  \small
  \begin{tabular}{cccccl}
    \toprule
    \textbf{Pasada} & $s$ & A & B & C & \textbf{Cambios} \\
    \midrule
    0 & $0$ & $\infty$ & $\infty$ & $\infty$ & \\
    1 & $0$ & $-4$ & $-2$ & $0$
      & las tres $(s, v)$; $(C, B)$: $-2$; $(B, A)$: $-4$ \\
    2 & $0$ & $-5$ & $-3$ & $-1$ & $(A, C)$, $(C, B)$, $(B, A)$ \\
    3 & $0$ & $-6$ & $-4$ & $-2$ & $(A, C)$, $(C, B)$, $(B, A)$ \\
    \midrule
    \multicolumn{6}{l}{Verificación: $(A, C)$ da $-6 + 3 = -3 < -2$,
      devuelve falso} \\
    \bottomrule
  \end{tabular}
  \end{center}
  Cada pasada le da una vuelta más al ciclo, y los tres valores bajan
  $1$, el peso del ciclo. No hay punto fijo: después de las
  $\abs{V'} - 1 = 3$ pasadas todavía hay una arista que relaja, la
  verificación devuelve falso y Johnson se detiene en la línea 5 sin
  calcular ningún $h$. Si siguiera con estos $d$, la desigualdad del
  triángulo no valdría y saldrían pesos $\hat{w}$ negativos.
\end{ejemplo}

\begin{ojo}[breakable]
  \begin{itemize}[itemsep=0.2em, leftmargin=1.2em]
    \item \textbf{Deshacer la reponderación, y con el signo correcto.}
          Sin la corrección se devuelve $\hat{\delta}$: en el
          ejemplo~\ref{ej:johnson-cormen},
          $\hat{\delta}(1, 3) = 2$ pero $\delta(1, 3) = -3$. Se suma
          $h(v) - h(u)$, no al revés.
    \item \textbf{Pesos reales.} Con punto flotante, $\hat{w}$ puede
          salir $-10^{-16}$ por redondeo, y Dijkstra supone pesos no
          negativos. Conviene redondear a $0$ los negativos diminutos
          o trabajar con enteros.
    \item \textbf{Un ciclo negativo en cualquier parte tumba todo.}
          Johnson rechaza la entrada completa aunque muchos pares no
          estén afectados. Marcar con $-\infty$ solo los afectados ya
          es otro algoritmo.
    \item \textbf{Grafos no dirigidos.} Una arista negativa son dos
          arcos que forman un ciclo negativo de dos aristas, y Johnson
          rechaza el grafo. Con pesos no negativos, $h \equiv 0$ y
          Johnson se reduce a $n$ veces Dijkstra.
    \item \textbf{Empates.} $\hat{w}$ produce muchas aristas de peso
          $0$, y Dijkstra empata seguido. Las distancias salen bien
          con cualquier desempate, pero los árboles pueden salir
          distintos a los de Cormen.
  \end{itemize}
\end{ojo}

\begin{ejercicio}
  \label{ej:fuente-johnson}
  ¿Qué pasa si en lugar de agregar $s$ se corre Bellman-Ford desde un
  vértice de $G$? (Cormen, ejercicio 25.3-6.)

  \emph{Solución.} En general falla. Con $V = \set{a, b}$, una arista
  $(a, b)$ de peso $-1$ y la fuente en $b$, queda
  $h(a) = \delta(b, a) = \infty$ y
  $\hat{w}(a, b) = -1 + \infty - 0 = \infty$, y Dijkstra desde $a$
  responde que no hay camino a $b$, cuando $\delta(a, b) = -1$. Si $G$
  es fuertemente conexo sí funciona: todos los $h$ son finitos, la
  desigualdad del triángulo vale en toda arista y Bellman-Ford ve
  todos los ciclos.
\end{ejercicio}

\begin{ejercicio}
  Muestra que si $G$ tiene un ciclo $c$ de peso $0$, Johnson deja
  $\hat{w}(u, v) = 0$ en cada arista de $c$. (Cormen, ejercicio
  25.3-5.)

  \emph{Solución.} Por el lema~\ref{lem:reponderacion-johnson},
  $\hat{w}(c) = w(c) = 0$, y $\hat{w}(c)$ es una suma de términos
  $\ge 0$; la única forma de sumar $0$ es que todos valgan $0$. En el
  ejemplo~\ref{ej:johnson-cormen}, las tres aristas del ciclo
  $2 \to 4 \to 3 \to 2$ quedaron con $\hat{w} = 0$.
\end{ejercicio}

\begin{ejercicio}
  Corre Floyd-Warshall y Johnson sobre el grafo con vértices
  $1, \dots, 6$ y aristas $(1, 5)$ con peso $-1$, $(2, 1)$ con $1$,
  $(2, 4)$ con $2$, $(3, 2)$ con $2$, $(3, 6)$ con $-8$, $(4, 1)$ con
  $-4$, $(4, 5)$ con $3$, $(5, 2)$ con $7$, $(6, 2)$ con $5$ y
  $(6, 3)$ con $10$ (figura 25.2 de Cormen, ejercicios 25.2-1 y
  25.3-1).

  \emph{Respuesta.} Johnson obtiene $h = (-5, -3, 0, -1, -6, -8)$ y
  los pesos reponderados $\hat{w}(1, 5) = 0$, $\hat{w}(2, 1) = 3$,
  $\hat{w}(2, 4) = 0$, $\hat{w}(3, 2) = 5$, $\hat{w}(3, 6) = 0$,
  $\hat{w}(4, 1) = 0$, $\hat{w}(4, 5) = 8$, $\hat{w}(5, 2) = 4$,
  $\hat{w}(6, 2) = 0$ y $\hat{w}(6, 3) = 2$. Los dos llegan a
  \[
    D = \begin{pmatrix}
      0 & 6 & \infty & 8 & -1 & \infty \\
      -2 & 0 & \infty & 2 & -3 & \infty \\
      -5 & -3 & 0 & -1 & -6 & -8 \\
      -4 & 2 & \infty & 0 & -5 & \infty \\
      5 & 7 & \infty & 9 & 0 & \infty \\
      3 & 5 & 10 & 7 & 2 & 0
    \end{pmatrix} .
  \]
  En Floyd-Warshall la vuelta $k = 3$ no cambia nada y la diagonal
  queda en ceros. Los $\infty$ salen porque a $3$ y a $6$ solo se
  llega desde $3$ y $6$; sirven para practicar el manejo de infinitos
  de la línea 12 de \textsc{Johnson}.
\end{ejercicio}

% =====================================================================

\subsection*{Qué algoritmo usar}

La tabla junta todo lo de la sección. ``Sí'' en la columna de pesos
negativos quiere decir ``sí, siempre que no haya ciclos negativos'', y
todos los que admiten negativos también detectan esos ciclos:
Floyd-Warshall y los cuadrados repetidos con la diagonal, Bellman-Ford
y Johnson con la verificación final.

\begin{center}
\small
\begin{tabular}{lcccc}
  \toprule
  \textbf{Algoritmo} & \textbf{Pesos} & \textbf{General}
    & \textbf{Disperso} & \textbf{Denso} \\
  & \textbf{negativos} & & $m = \Tht{n}$ & $m = \Tht{n^2}$ \\
  \midrule
  Floyd-Warshall & sí & $\Tht{n^3}$ & $\Tht{n^3}$ & $\Tht{n^3}$ \\
  Cuadrados repetidos & sí & $\Tht{n^3 \log n}$ & $\Tht{n^3 \log n}$
    & $\Tht{n^3 \log n}$ \\
  \addlinespace
  $n$ veces Dijkstra, arreglo & no & $\Tht{n^3}$ & $\Tht{n^3}$
    & $\Tht{n^3}$ \\
  $n$ veces Dijkstra, binario & no & $\Ord{n(n + m)\log n}$
    & $\Ord{n^2 \log n}$ & $\Ord{n^3 \log n}$ \\
  $n$ veces Dijkstra, Fibonacci & no & $\Ord{n^2 \log n + nm}$
    & $\Ord{n^2 \log n}$ & $\Ord{n^3}$ \\
  \addlinespace
  $n$ veces Bellman-Ford & sí & $\Ord{n^2 m}$ & $\Ord{n^3}$
    & $\Ord{n^4}$ \\
  \addlinespace
  Johnson, binario & sí & $\Ord{n(n + m)\log n}$ & $\Ord{n^2 \log n}$
    & $\Ord{n^3 \log n}$ \\
  Johnson, Fibonacci & sí & $\Ord{n^2 \log n + nm}$
    & $\Ord{n^2 \log n}$ & $\Ord{n^3}$ \\
  \bottomrule
\end{tabular}
\end{center}

Para dimensionar, con $n = 10^4$ y solo el término dominante de cada
cota, sin constantes:
\begin{center}
\small
\begin{tabular}{lcc}
  \toprule
  \textbf{Cota} & \textbf{Disperso}, $m = 5 \cdot 10^4$
    & \textbf{Denso}, $m \approx 10^8$ \\
  \midrule
  Floyd-Warshall, $n^3$ & $10^{12}$ & $10^{12}$ \\
  Dijkstra o Johnson binario, $n(n + m)\log n$ & $8.0 \cdot 10^{9}$
    & $1.3 \cdot 10^{13}$ \\
  Dijkstra o Johnson Fibonacci, $n^2 \log n + nm$
    & $1.8 \cdot 10^{9}$ & $1.0 \cdot 10^{12}$ \\
  $n$ veces Bellman-Ford, $n^2 m$ & $5.0 \cdot 10^{12}$
    & $1.0 \cdot 10^{16}$ \\
  \bottomrule
\end{tabular}
\end{center}

\textbf{Dónde está el cruce.} Con montículo binario, Johnson le gana a
Floyd-Warshall cuando $m = \ord{n^2 / \log n}$, porque entonces
$n(n + m)\log n = \ord{n^3}$. Con Fibonacci, cuando $m = \ord{n^2}$; en el
caso denso empatan asintóticamente. En la práctica, en grafos densos
gana Floyd-Warshall: tres ciclos anidados sobre una matriz contigua
tienen constantes mucho menores que un montículo de Fibonacci, cuyas
constantes son famosamente grandes.

\begin{idea}
  Grafo denso o chico: Floyd-Warshall, por sencillez y constantes, con
  pesos de cualquier signo. Grafo grande y disperso con pesos no
  negativos: $n$ veces Dijkstra. Grafo grande y disperso con pesos
  negativos: Johnson, que cuesta lo mismo que $n$ veces Dijkstra más
  una corrida de Bellman-Ford. Si solo importa si hay camino:
  cerradura transitiva con bits, o un BFS por vértice si el grafo es
  disperso.
\end{idea}

\begin{observacion}
  En grafos con $m = \Tht{n^{1 + \varepsilon}}$ y
  $0 < \varepsilon \le 1$ fijo, se puede igualar a Fibonacci con una
  estructura sencilla (Cormen, problema 25-2). Un montículo $d$-ario
  con $d = n^{\varepsilon}$ tiene altura
  $\log_d n = 1/\varepsilon = \Ord{1}$: insertar y decrementar cuestan
  $\Ord{1}$ sin amortizar, y extraer, que compara con $d$ hijos por
  nivel, cuesta $\Ord{d \log_d n} = \Ord{n^{\varepsilon}}$. Las $n$
  extracciones de Dijkstra suman $\Ord{n^{1 + \varepsilon}} = \Ord{m}$
  y los $m$ decrementos, $\Ord{m}$: Dijkstra queda en $\Ord{m}$, y
  Johnson en $\Ord{nm}$. Es el mismo truco que motiva a Fibonacci:
  abaratar \textsc{Decrementar-Llave}, que es lo que Dijkstra más
  repite.
\end{observacion}

\newpage

% =====================================================================
%  Grafos - Flujo maximo
% =====================================================================

\clase{Flujo máximo: Ford-Fulkerson y Edmonds-Karp}
  {Complemento: Cormen, cap.\ 26.1 y 26.2}
\label{sec:flujo}

\begin{quick_sum}
  Una red de flujo es un grafo dirigido con una capacidad en cada
  arista, una fuente $s$ y un sumidero $t$; el problema es mandar de
  $s$ a $t$ tanto material como se pueda sin rebasar ninguna
  capacidad. La sección construye el método de Ford-Fulkerson:
  empezar con flujo cero y aumentar mientras exista un camino de $s$ a
  $t$ en la red residual, el grafo que registra cuánto cabe todavía en
  cada arista y cuánto se puede regresar. Regresar flujo es lo que
  separa al método de un voraz, que en la red de ejemplo se atora en
  $19$ cuando el máximo es $23$. El teorema de flujo máximo y corte
  mínimo prueba que, cuando ya no hay camino, el flujo es óptimo, y
  entrega un corte que lo certifica. Con capacidades enteras
  Ford-Fulkerson cuesta $\Ord{m \abs{f^*}}$, y una mala elección de
  caminos alcanza esa cota: dos millones de aumentos donde bastan dos.
  Edmonds-Karp escoge siempre el camino con menos aristas, con el BFS
  de la clase del 22 de octubre, y baja el costo a $\Ord{n m^2}$ sin
  importar las capacidades.
\end{quick_sum}

\subsection*{Redes de flujo}

Una compañía fabrica discos de hockey en Vancouver y los guarda en una
bodega en Winnipeg. Renta espacio en camiones que cubren rutas entre
ciudades, y cada ruta admite a lo más cierto número de cajas por día.
La pregunta es cuántas cajas diarias pueden llegar a Winnipeg. No
importa por dónde viajen ni cuánto cueste: importa la tasa. Ese es el
problema de flujo máximo, y el mismo modelo sirve para líquido en
tuberías, corriente en una red eléctrica o paquetes en una red de
datos. El ejemplo es de Cormen, y su red se usa en toda la sección.

\begin{definicion}[Red de flujo]
  Una \textbf{red de flujo} (\textit{flow network}) es un grafo
  dirigido $G = (V, E)$ en el que cada arista $(u, v)$ tiene una
  \textbf{capacidad} $c(u, v) \ge 0$, junto con dos vértices
  distinguidos: la \textbf{fuente} $s$ y el \textbf{sumidero} $t$. Por
  comodidad, $c(u, v) = 0$ si $(u, v) \notin E$. Se pide además que no
  haya lazos, que no haya \textbf{aristas antiparalelas} (si
  $(u, v) \in E$, entonces $(v, u) \notin E$) y que todo vértice esté
  en algún camino de $s$ a $t$.
\end{definicion}

La última condición tiene una consecuencia que se usa en todos los
análisis: todo vértice salvo $s$ tiene al menos una arista de entrada,
así que $m \ge n - 1$ y $\Ord{n + m} = \Ord{m}$.

\begin{definicion}[Flujo]
  Un \textbf{flujo} en $G$ es una función $f : V \times V \to \R$ que
  cumple dos propiedades:
  \begin{itemize}
    \item \textbf{Restricción de capacidad}: $0 \le f(u, v) \le
          c(u, v)$ para todos $u, v \in V$.
    \item \textbf{Conservación}: para todo $u \in V - \set{s, t}$,
          \[
            \sum_{v \in V} f(v, u) = \sum_{v \in V} f(u, v) .
          \]
  \end{itemize}
  El \textbf{valor} del flujo es lo que sale de la fuente menos lo que
  entra a ella:
  \[
    \abs{f} = \sum_{v \in V} f(s, v) - \sum_{v \in V} f(v, s) .
  \]
  El \textbf{problema de flujo máximo} pide, dados $G$, $s$ y $t$, un
  flujo de valor máximo. A un flujo así se le llama $f^*$.
\end{definicion}

La conservación dice que lo que entra a un vértice intermedio sale de
él; cuando el material es corriente eléctrica, es la ley de corrientes
de Kirchhoff. La restricción de capacidad obliga $f(u, v) = 0$ cuando
$(u, v) \notin E$, porque ahí $c(u, v) = 0$.

\begin{ojo}
  Las barras de $\abs{f}$ no son valor absoluto ni cardinalidad:
  denotan el valor del flujo. Y el término que resta lo que entra a
  $s$ no sobra. En una red de flujo normalmente no entra nada a la
  fuente, pero la definición se aplica también a flujos en redes
  residuales, que sí tienen aristas hacia $s$.
\end{ojo}

\begin{ejemplo}[La red de Cormen]\label{ej:flujo-red}
  La red tiene $V = \set{s, v_1, v_2, v_3, v_4, t}$ y nueve aristas:
  $n = 6$ y $m = 9$. Las ciudades son Vancouver ($s$), Edmonton
  ($v_1$), Calgary ($v_2$), Saskatoon ($v_3$), Regina ($v_4$) y
  Winnipeg ($t$). A la izquierda van las capacidades; a la derecha, un
  flujo, con cada arista etiquetada $f(u, v)/c(u, v)$. La diagonal
  separa los dos números, no los divide.
  \begin{center}
  \begin{tikzpicture}
    \foreach \n/\x/\y/\e in {s/0/0/s, v1/1.8/1.2/v_1,
        v2/1.8/-1.2/v_2, v3/4.2/1.2/v_3, v4/4.2/-1.2/v_4,
        t/6/0/t}
      \node[vert] (\n) at (\x, \y) {$\e$};
    \draw[dirigida] (s) -- node[peso] {16} (v1);
    \draw[dirigida] (s) -- node[peso] {13} (v2);
    \draw[dirigida] (v1) -- node[peso] {12} (v3);
    \draw[dirigida] (v2) -- node[peso] {4} (v1);
    \draw[dirigida] (v3) -- node[peso] {9} (v2);
    \draw[dirigida] (v2) -- node[peso] {14} (v4);
    \draw[dirigida] (v4) -- node[peso] {7} (v3);
    \draw[dirigida] (v3) -- node[peso] {20} (t);
    \draw[dirigida] (v4) -- node[peso] {4} (t);
    \node[nota] at (3, -2) {capacidades $c(u, v)$};
  \end{tikzpicture}
  \hspace{0.8cm}
  \begin{tikzpicture}
    \foreach \n/\x/\y/\e in {s/0/0/s, v1/1.8/1.2/v_1,
        v2/1.8/-1.2/v_2, v3/4.2/1.2/v_3, v4/4.2/-1.2/v_4,
        t/6/0/t}
      \node[vert] (\n) at (\x, \y) {$\e$};
    \draw[dirigida] (s) -- node[peso] {11/16} (v1);
    \draw[dirigida] (s) -- node[peso] {8/13} (v2);
    \draw[dirigida] (v1) -- node[peso] {12/12} (v3);
    \draw[dirigida] (v2) -- node[peso] {1/4} (v1);
    \draw[dirigida] (v3) -- node[peso] {4/9} (v2);
    \draw[dirigida] (v2) -- node[peso] {11/14} (v4);
    \draw[dirigida] (v4) -- node[peso] {7/7} (v3);
    \draw[dirigida] (v3) -- node[peso] {15/20} (t);
    \draw[dirigida] (v4) -- node[peso] {4/4} (t);
    \node[nota] at (3, -2) {flujo $f(u, v)/c(u, v)$, $\abs{f} = 19$};
  \end{tikzpicture}
  \end{center}
  El flujo respeta todas las capacidades. A $v_2$ le entran
  $f(s, v_2) + f(v_3, v_2) = 8 + 4 = 12$ unidades y le salen
  $f(v_2, v_1) + f(v_2, v_4) = 1 + 11 = 12$; en $v_1$, $v_3$ y $v_4$
  entran y salen $12$, $19$ y $11$, respectivamente. El valor es
  $\abs{f} = f(s, v_1) + f(s, v_2) = 11 + 8 = 19$. El máximo de esta
  red, se verá, es $23$.
\end{ejemplo}

\begin{observacion}
  Cormen escribe el flujo de una arista como atributo, $(u, v).f$, y
  dentro de la $\mathcal{O}$ abrevia $\abs{V}$ y $\abs{E}$ como $V$ y
  $E$: su $\Ord{VE^2}$ es el $\Ord{n m^2}$ de esta sección.
\end{observacion}

\subsubsection*{Aristas antiparalelas}

La definición prohíbe los pares $(u, v)$, $(v, u)$, y no por capricho.
La red residual, que viene enseguida, le asigna a cada par ordenado
una sola capacidad residual, y con un par antiparalelo no sabría si
$(u, v)$ es una arista original o el regreso de $(v, u)$. Pero los
pares aparecen solos en cualquier modelo: si la transportista ofrece
además $10$ cajas diarias de Edmonton a Calgary, la red tendría
$(v_1, v_2)$ junto a $(v_2, v_1)$.

La salida es partir una de las dos con un vértice nuevo $v'$: se
reemplaza $(v_1, v_2)$ por $(v_1, v')$ y $(v', v_2)$, las dos con la
capacidad original. Todo lo que pasaba por la arista vieja pasa ahora
por el camino de dos aristas, y al revés, así que las dos redes
admiten los mismos valores de flujo (Cormen, ejercicio 26.1-1).
Cuesta un vértice y una arista por cada par.

\begin{center}
\begin{tikzpicture}
  \node[vert] (v1) at (0, 0) {$v_1$};
  \node[vert] (v2) at (2.4, 0) {$v_2$};
  \draw[dirigida] (v1) to[bend left=25] node[peso] {10} (v2);
  \draw[dirigida] (v2) to[bend left=25] node[peso] {4} (v1);
  \node[nota] at (1.2, -1.2) {par antiparalelo: prohibido};
\end{tikzpicture}
\hspace{1.5cm}
\begin{tikzpicture}
  \node[vert] (v1) at (0, 0) {$v_1$};
  \node[vert] (v2) at (2.4, 0) {$v_2$};
  \node[vertMarca] (vp) at (1.2, 0.9) {$v'$};
  \draw[dirigida] (v1) -- node[peso] {10} (vp);
  \draw[dirigida] (vp) -- node[peso] {10} (v2);
  \draw[dirigida] (v2) -- node[peso] {4} (v1);
  \node[nota] at (1.2, -1.2) {se parte $(v_1, v_2)$ con $v'$};
\end{tikzpicture}
\end{center}

Con la red de Cormen más la arista $(v_1, v_2)$ de capacidad $10$, la
versión con el par y la versión partida tienen el mismo flujo máximo,
$23$, que además es el de la red sin la arista extra: más abajo se
verá por qué esa arista no ayuda.

\subsubsection*{Varias fuentes y varios sumideros}

Si hay fuentes $s_1, \dots, s_k$ y sumideros $t_1, \dots, t_\ell$, se
agrega una \textbf{superfuente} $s$ con aristas $(s, s_i)$ de
capacidad $\infty$ y un \textbf{supersumidero} $t$ con aristas
$(t_j, t)$ de capacidad $\infty$. La superfuente le da a cada $s_i$
todo lo que pida y el supersumidero recibe todo lo que llegue, así que
los flujos de la red original corresponden uno a uno con los de la red
nueva, con el mismo valor (Cormen, ejercicio 26.1-2). Cormen cuenta
las fuentes y los sumideros con $m$ y $n$; aquí esas letras ya son
$\abs{E}$ y $\abs{V}$, y por eso se usan $k$ y $\ell$.

En la figura, lo que se agrega va en naranja. La red original tiene
dos fuentes, dos sumideros y dos vértices intermedios, $x$ y $y$.

\begin{center}
\begin{tikzpicture}
  \foreach \n/\x/\y/\e in {s1/0/1/s_1, s2/0/-1/s_2, x/2.2/1/x,
      y/2.2/-1/y, t1/4.4/1/t_1, t2/4.4/-1/t_2}
    \node[vert] (\n) at (\x, \y) {$\e$};
  \node[vertMarca] (s) at (-2, 0) {$s$};
  \node[vertMarca] (t) at (6.4, 0) {$t$};
  \foreach \a/\b/\c in {s1/x/5, s2/x/3, s2/y/4, x/t1/6, y/t1/2,
      y/t2/5}
    \draw[dirigida] (\a) -- node[peso] {\c} (\b);
  \foreach \a/\b in {s/s1, s/s2, t1/t, t2/t}
    \draw[dirigidaFuera] (\a) -- node[peso] {$\infty$} (\b);
\end{tikzpicture}
\end{center}

En una implementación, $\infty$ se reemplaza por un número grande, por
ejemplo la suma de todas las capacidades originales. Como esas son
finitas, el flujo máximo también lo es (Cormen, ejercicio 26.2-5).

\begin{ejercicio}[Capacidades en los vértices; Cormen, ejercicio 26.1-7]
  Además de la capacidad de cada arista, cada vértice $v$ tiene un
  límite $l(v)$ de cuánto flujo puede pasar por él. Convierte esa red
  en una red de flujo ordinaria. ¿Cuántos vértices y aristas tiene?

  \textit{Solución.} Se parte cada vértice $v$ en $v_{\text{ent}}$ y
  $v_{\text{sal}}$, unidos por una arista $(v_{\text{ent}},
  v_{\text{sal}})$ de capacidad $l(v)$. Las aristas que entraban a $v$
  llegan a $v_{\text{ent}}$ y las que salían parten de
  $v_{\text{sal}}$. Todo el flujo que pasa por $v$ cruza esa arista,
  así que respeta $l(v)$. La red nueva tiene $2n$ vértices y $n + m$
  aristas. Es el truco de las antiparalelas al revés: allá se parte
  una arista con un vértice, aquí un vértice con una arista.
\end{ejercicio}

% =====================================================================

\subsection*{La red residual}

\subsubsection*{Por qué no basta un voraz}

La idea más directa es la de la clase del 8 de octubre: mientras haya
un camino de $s$ a $t$ cuyas aristas tengan holgura,
$f(u, v) < c(u, v)$, mandar por él todo lo que permita su arista más
angosta. Toma lo mejor que ve en el momento y nunca revisa una
decisión ya tomada. En esta red eso no alcanza.

\begin{ejemplo}[El voraz se atora en $19$]\label{ej:flujo-voraz}
  El voraz empieza con el camino $s, v_1, v_3, v_2, v_4, t$, que
  admite $\min(16, 12, 9, 14, 4) = 4$ unidades. Después busca con BFS
  entre las aristas con holgura y encuentra $s, v_1, v_3, t$, con $8$,
  y $s, v_2, v_4, v_3, t$, con $7$. Lleva $19$. A la izquierda está ese
  flujo: $(v_1, v_3)$, $(v_4, v_3)$ y $(v_4, t)$, en naranja,
  quedaron saturadas, y desde $s$ solo se llega a $v_1$, $v_2$ y $v_4$
  por aristas con holgura. No hay más camino y el voraz se detiene.
  \begin{center}
  \begin{tikzpicture}
    \foreach \n/\x/\y/\e in {s/0/0/s, v1/1.8/1.2/v_1,
        v2/1.8/-1.2/v_2, v3/4.2/1.2/v_3, v4/4.2/-1.2/v_4,
        t/6/0/t}
      \node[vert] (\n) at (\x, \y) {$\e$};
    \draw[dirigida] (s) -- node[peso] {12/16} (v1);
    \draw[dirigida] (s) -- node[peso] {7/13} (v2);
    \draw[dirigidaResalte] (v1) -- node[peso] {12/12} (v3);
    \draw[dirigida] (v2) -- node[peso] {0/4} (v1);
    \draw[dirigida] (v3) -- node[peso] {4/9} (v2);
    \draw[dirigida] (v2) -- node[peso] {11/14} (v4);
    \draw[dirigidaResalte] (v4) -- node[peso] {7/7} (v3);
    \draw[dirigida] (v3) -- node[peso] {15/20} (t);
    \draw[dirigidaResalte] (v4) -- node[peso] {4/4} (t);
    \node[nota] at (3, -2) {flujo del voraz, $\abs{f} = 19$};
  \end{tikzpicture}
  \hspace{0.8cm}
  \begin{tikzpicture}
    \foreach \n/\x/\y/\e in {s/0/0/s, v1/1.8/1.2/v_1,
        v2/1.8/-1.2/v_2, v3/4.2/1.2/v_3, v4/4.2/-1.2/v_4,
        t/6/0/t}
      \node[vert] (\n) at (\x, \y) {$\e$};
    \draw[dirigida] (s) to[bend left=16] node[peso, auto] {4} (v1);
    \draw[dirigida] (v1) to[bend left=16] node[peso, auto] {12} (s);
    \draw[dirigida] (v2) to[bend left=16] node[peso, auto] {7} (s);
    \draw[dirigida] (v3) -- node[peso] {12} (v1);
    \draw[dirigida] (v2) -- node[peso] {4} (v1);
    \draw[dirigida] (v3) to[bend left=16] node[peso, auto] {5} (v2);
    \draw[dirigida] (v2) to[bend left=16] node[peso, auto] {3} (v4);
    \draw[dirigida] (v4) to[bend left=16] node[peso, auto] {11} (v2);
    \draw[dirigida] (v3) -- node[peso] {7} (v4);
    \draw[dirigida] (t) to[bend left=16] node[peso, auto] {15} (v3);
    \draw[dirigida] (t) -- node[peso] {4} (v4);
    \draw[dirigidaArbol] (s) to[bend left=16] node[peso, auto] {6} (v2);
    \draw[dirigidaArbol] (v2) to[bend left=16] node[peso, auto] {4} (v3);
    \draw[dirigidaArbol] (v3) to[bend left=16] node[peso, auto] {5} (t);
    \node[nota] at (3, -2) {su red residual $G_f$};
  \end{tikzpicture}
  \end{center}
  Pero $19$ no es el máximo. Estorban las $4$ unidades que el primer
  camino mandó de $v_3$ a $v_2$: ocupan capacidad de $(v_1, v_3)$ que
  hace falta para llegar a $t$ por arriba. A la derecha está la red
  residual de ese flujo, que se define enseguida. En ella sí hay
  camino, $s, v_2, v_3, t$, y su arista de en medio, $(v_2, v_3)$, no
  existe en $G$: es el regreso de las $4$ unidades de $(v_3, v_2)$.
  Mandar $4$ por ese camino deja $f(s, v_2) = 11$, $f(v_3, v_2) = 0$ y
  $f(v_3, t) = 19$, de valor $23$. Dicho en cajas: $v_2$ recibe ahora
  de $s$ lo que antes le llegaba de $v_3$, y $v_3$ manda esas $4$
  unidades directo a $t$.
\end{ejemplo}

El voraz no siempre falla: si empieza por $s, v_1, v_3, t$, llega a
$23$. Falla cuando una decisión temprana le estorba a una posterior y
no tiene forma de deshacerla. La red residual es justo esa forma.

\subsubsection*{Capacidad residual}

\begin{definicion}[Red residual]\label{def:red-residual}
  Sea $f$ un flujo en $G$. La \textbf{capacidad residual}
  (\textit{residual capacity}) de un par $u, v \in V$ es
  \[
    c_f(u, v) =
    \begin{cases}
      c(u, v) - f(u, v) & \text{si } (u, v) \in E , \\
      f(v, u)           & \text{si } (v, u) \in E , \\
      0                 & \text{en otro caso.}
    \end{cases}
  \]
  La \textbf{red residual} de $G$ inducida por $f$ es
  $G_f = (V, E_f)$, con
  \[
    E_f = \set{(u, v) \in V \times V : c_f(u, v) > 0} .
  \]
  Sus aristas se llaman \textbf{aristas residuales}.
\end{definicion}

El primer renglón es el espacio que queda en la arista. El segundo es
el \textbf{regreso de flujo}: si de $v$ a $u$ van $f(v, u)$ unidades,
se puede cancelar cualquier parte de ellas, y visto desde el resultado
eso equivale a mandar flujo de $u$ a $v$. Como no hay aristas
antiparalelas, a cada par ordenado le toca exactamente un renglón.
Cada arista residual es una arista de $E$ o la inversa de una, así que
$\abs{E_f} \le 2m$. $G_f$ no es una red de flujo, porque puede tener
pares antiparalelos, pero un flujo en $G_f$ se define igual que en una
red, con las capacidades $c_f$.

\begin{ejemplo}[Una arista y sus dos residuales]\label{ej:flujo-arista}
  Si $c(u, v) = 16$ y $f(u, v) = 11$, caben $5$ unidades más,
  $c_f(u, v) = 5$, y se pueden regresar hasta $11$, $c_f(v, u) = 11$.
  \begin{center}
  \begin{tikzpicture}
    \node[vert] (u) at (0, 0) {$u$};
    \node[vert] (v) at (2.4, 0) {$v$};
    \draw[dirigida] (u) -- node[peso] {11/16} (v);
    \node[nota] at (1.2, -0.8) {en $G$: $f(u, v)/c(u, v)$};
  \end{tikzpicture}
  \hspace{1.5cm}
  \begin{tikzpicture}
    \node[vert] (u) at (0, 0) {$u$};
    \node[vert] (v) at (2.4, 0) {$v$};
    \draw[dirigida] (u) to[bend left=25] node[peso] {5} (v);
    \draw[dirigida] (v) to[bend left=25] node[peso] {11} (u);
    \node[nota] at (1.2, -1.0) {en $G_f$: $5$ de ida, $11$ de regreso};
  \end{tikzpicture}
  \end{center}
  Una arista saturada solo deja la residual de regreso, y una vacía
  solo la de ida. En las figuras de esta sección las aristas
  residuales de capacidad $0$ no se dibujan.
\end{ejemplo}

\subsubsection*{Aumento}

\begin{definicion}[Aumento]
  Si $f$ es un flujo en $G$ y $f'$ uno en $G_f$, el \textbf{aumento}
  de $f$ por $f'$ es la función $f \uparrow f'$ dada por
  \[
    (f \uparrow f')(u, v) =
    \begin{cases}
      f(u, v) + f'(u, v) - f'(v, u) & \text{si } (u, v) \in E , \\
      0 & \text{en otro caso.}
    \end{cases}
  \]
\end{definicion}

Se suma lo que $f'$ manda por $(u, v)$ y se resta lo que manda por la
arista de regreso, $(v, u)$. Empujar por la de regreso se llama
\textbf{cancelación}: mandar $5$ cajas de $u$ a $v$ y $2$ de $v$ a $u$
es, visto desde el resultado, mandar $3$ de $u$ a $v$.

\begin{lema}[Cormen, lema 26.1]\label{lem:flujo-aumento}
  Si $f$ es un flujo en $G$ y $f'$ es un flujo en $G_f$, entonces
  $f \uparrow f'$ es un flujo en $G$ y
  $\abs{f \uparrow f'} = \abs{f} + \abs{f'}$.
\end{lema}

\begin{proof}
  Sea $h = f \uparrow f'$. Primero la capacidad. Fuera de $E$, $h$
  vale $0$. Si $(u, v) \in E$, como $f'(v, u) \le c_f(v, u) =
  f(u, v)$,
  \[
    h(u, v) \ge f(u, v) + f'(u, v) - f(u, v) = f'(u, v) \ge 0 ,
  \]
  y como $f'(v, u) \ge 0$ y $f'(u, v) \le c_f(u, v) = c(u, v) -
  f(u, v)$,
  \[
    h(u, v) \le f(u, v) + f'(u, v) \le c(u, v) .
  \]

  Para la conservación y el valor conviene mirar el flujo en neto.
  Para cualquier función $g$ en $V \times V$ sea
  $\bar{g}(u, v) = g(u, v) - g(v, u)$, lo que $g$ manda de $u$ a $v$
  descontando lo que manda de $v$ a $u$. La conservación en $u$ dice
  $\sum_{v} \bar{g}(u, v) = 0$, y el valor es
  $\abs{g} = \sum_{v} \bar{g}(s, v)$. Basta probar que para todo par
  \[
    \bar{h}(u, v) = \bar{f}(u, v) + \bar{f'}(u, v) , \tag{$*$}
  \]
  porque al sumar sobre $v$ la conservación de $f$ y la de $f'$ dan la
  de $h$, y en $u = s$ queda $\abs{h} = \abs{f} + \abs{f'}$. Hay tres
  casos. Si $(u, v) \in E$, entonces $(v, u) \notin E$, así que
  $f(v, u) = 0$ y $h(v, u) = 0$, y
  $\bar{h}(u, v) = f(u, v) + f'(u, v) - f'(v, u) = \bar{f}(u, v) +
  \bar{f'}(u, v)$. Si $(v, u) \in E$, es el caso anterior con los
  papeles cambiados, y los dos lados de $(*)$ solo cambian de signo.
  Si ninguna de las dos está en $E$, entonces $f$, $h$ y $f'$ valen
  $0$ en los dos sentidos ($f'$ porque ahí $c_f = 0$), y los dos lados
  son $0$.
\end{proof}

Cormen prueba la parte del valor separando los vecinos de $s$ en los
de salida y los de entrada, y deja un paso de las sumas como ejercicio
(26.2-1). El flujo neto $\bar{g}$ resuelve la conservación y el valor
de un golpe.

\subsubsection*{Caminos aumentantes}

\begin{definicion}[Camino aumentante]
  Un \textbf{camino aumentante} (\textit{augmenting path}) es un
  camino $p$ de $s$ a $t$ en la red residual $G_f$. Su
  \textbf{capacidad residual} es
  \[
    c_f(p) = \min\set{c_f(u, v) : (u, v) \text{ está en } p} ,
  \]
  lo más que se puede mandar por todas sus aristas a la vez. El flujo
  $f_p$ en $G_f$ vale $c_f(p)$ en las aristas de $p$ y $0$ en las
  demás.
\end{definicion}

\begin{lema}[Cormen, lema 26.2 y corolario 26.3]\label{lem:flujo-camino}
  Si $p$ es un camino aumentante, $f_p$ es un flujo en $G_f$ con
  $\abs{f_p} = c_f(p) > 0$, y $f \uparrow f_p$ es un flujo en $G$ con
  $\abs{f \uparrow f_p} = \abs{f} + c_f(p) > \abs{f}$.
\end{lema}

\begin{proof}
  En cada arista de $p$, $f_p$ vale $c_f(p)$, que no pasa de
  $c_f(u, v)$ porque es el mínimo; fuera de $p$ vale $0$. Como $p$ no
  repite vértices, a cada vértice intermedio de $p$ le entra
  exactamente una arista de $p$ y le sale exactamente una, las dos con
  $c_f(p)$; los vértices fuera de $p$ no tocan ninguna. De $s$ sale una
  arista de $p$ y no entra ninguna, así que $\abs{f_p} = c_f(p)$, que
  es positivo porque toda arista residual tiene capacidad positiva. La
  segunda parte es el lema~\ref{lem:flujo-aumento} con $f' = f_p$.
\end{proof}

Cormen deja la primera parte como ejercicio (26.2-7). Con este lema el
método ya está completo: mientras haya un camino aumentante, aumentar
por él, y cada vuelta sube el valor. Falta saber que cuando no queda
ningún camino el flujo es máximo, y para eso hacen falta los cortes.

% =====================================================================

\subsection*{Cortes y el teorema de flujo máximo y corte mínimo}

\begin{definicion}[Corte]
  Un \textbf{corte} (\textit{cut}) $(S, T)$ de la red es una
  partición de $V$ con $s \in S$ y $t \in T = V - S$. El
  \textbf{flujo neto} que lo cruza es
  \[
    f(S, T) = \sum_{u \in S} \sum_{v \in T} f(u, v)
            - \sum_{u \in S} \sum_{v \in T} f(v, u) ,
  \]
  y su \textbf{capacidad} es
  \[
    c(S, T) = \sum_{u \in S} \sum_{v \in T} c(u, v) .
  \]
  Un \textbf{corte mínimo} es uno de capacidad mínima entre todos los
  cortes de la red.
\end{definicion}

La asimetría es deliberada. El flujo neto cuenta lo que va de $S$ a
$T$ y resta lo que regresa de $T$ a $S$; la capacidad solo cuenta las
aristas de $S$ a $T$, porque una arista de regreso no le agrega
espacio al cruce: lo mejor que puede hacer es no llevar nada. Es la
misma idea de partir los vértices en dos lados que sostiene a Kruskal
(sección~\ref{sec:mst}), con tres diferencias: aquí el grafo es
dirigido, $s$ y $t$ quedan en lados opuestos, y la capacidad mide un
solo sentido.

\begin{ejemplo}[Un corte de la red de Cormen]\label{ej:flujo-corte}
  Con el flujo de valor $19$ del ejemplo~\ref{ej:flujo-red}, tómese
  $S = \set{s, v_1, v_2}$ y $T = \set{v_3, v_4, t}$ (es el de la
  izquierda en la figura del ejemplo~\ref{ej:flujo-corte-minimo}). De
  $S$ a $T$ cruzan $(v_1, v_3)$ y $(v_2, v_4)$; de $T$ a $S$,
  $(v_3, v_2)$. Entonces
  \[
    f(S, T) = 12 + 11 - 4 = 19 , \qquad c(S, T) = 12 + 14 = 26 .
  \]
  La capacidad de $(v_3, v_2)$, $9$, no se suma a $c(S, T)$, pero su
  flujo, $4$, sí se resta en $f(S, T)$. Confundir las dos cuentas
  rompe la prueba del corolario~\ref{cor:flujo-cota}.
\end{ejemplo}

Que el flujo neto haya salido igual a $\abs{f}$ no es casualidad. Todo
lo que sale de $s$ tiene que cruzar el corte para llegar a $t$, y lo
que cruza de regreso se descuenta.

\begin{lema}[Cormen, lema 26.4]\label{lem:flujo-neto-corte}
  Para todo flujo $f$ y todo corte $(S, T)$, $f(S, T) = \abs{f}$.
\end{lema}

\begin{proof}
  Para cada $u \in S - \set{s}$, la conservación dice
  $\sum_{v} f(u, v) - \sum_{v} f(v, u) = 0$ ($u \ne t$ porque
  $t \in T$). Sumar todos esos ceros a la definición de $\abs{f}$ no
  la cambia:
  \[
    \abs{f} = \sum_{u \in S} \Bigl( \sum_{v \in V} f(u, v)
              - \sum_{v \in V} f(v, u) \Bigr) .
  \]
  Se parte cada suma interior en $v \in S$ y $v \in T$:
  \[
    \abs{f} = \sum_{u \in S} \sum_{v \in S} f(u, v)
            - \sum_{u \in S} \sum_{v \in S} f(v, u)
            + \sum_{u \in S} \sum_{v \in T} f(u, v)
            - \sum_{u \in S} \sum_{v \in T} f(v, u) .
  \]
  Las dos primeras sumas tienen exactamente los mismos términos,
  $f(x, y)$ con $x, y \in S$, y se cancelan. Las dos últimas son
  $f(S, T)$.
\end{proof}

En la red de Cormen hay $2^4 = 16$ cortes, uno por cada subconjunto de
$\set{v_1, v_2, v_3, v_4}$ que acompaña a $s$, y con el flujo de valor
$19$ los $16$ dan $f(S, T) = 19$.

\begin{corolario}[Cormen, corolario 26.5]\label{cor:flujo-cota}
  Para todo flujo $f$ y todo corte $(S, T)$, $\abs{f} \le c(S, T)$.
\end{corolario}

\begin{proof}
  \[
    \abs{f} = f(S, T)
    = \sum_{u \in S} \sum_{v \in T} f(u, v)
      - \sum_{u \in S} \sum_{v \in T} f(v, u)
    \le \sum_{u \in S} \sum_{v \in T} f(u, v)
    \le \sum_{u \in S} \sum_{v \in T} c(u, v) = c(S, T) .
  \]
  La primera desigualdad tira términos no negativos que se restaban;
  la segunda es la restricción de capacidad.
\end{proof}

Cualquier corte es una cota superior para cualquier flujo. En la red
de Cormen, el corte del ejemplo~\ref{ej:flujo-corte} dice que ningún
flujo pasa de $26$, y el corte $(\set{s}, V - \set{s})$ dice que
ninguno pasa de $16 + 13 = 29$. La mejor de esas cotas la da un corte
mínimo, y el teorema central dice que esa cota se alcanza.

\begin{teorema}[Flujo máximo y corte mínimo; Cormen,
  teorema 26.6]\label{teo:flujo-corte}
  Si $f$ es un flujo en $G$, son equivalentes:
  \begin{enumerate}
    \item $f$ es un flujo máximo en $G$.
    \item La red residual $G_f$ no tiene caminos aumentantes.
    \item $\abs{f} = c(S, T)$ para algún corte $(S, T)$ de $G$.
  \end{enumerate}
\end{teorema}

\begin{proof}
  (1) $\Rightarrow$ (2). Si hubiera un camino aumentante $p$, por el
  lema~\ref{lem:flujo-camino} el flujo $f \uparrow f_p$ tendría valor
  mayor que $\abs{f}$, y $f$ no sería máximo.

  (2) $\Rightarrow$ (3). Sea
  \[
    S = \set{v \in V : \text{hay un camino de } s \text{ a } v
             \text{ en } G_f}
  \]
  y $T = V - S$. Es un corte: $s \in S$, y $t \notin S$ porque en
  $G_f$ no hay camino de $s$ a $t$. Sean $u \in S$ y $v \in T$. Si
  $(u, v) \in E$, entonces $f(u, v) = c(u, v)$: de lo contrario
  $c_f(u, v) > 0$, y $v$ sería alcanzable desde $s$ pasando por $u$.
  Si $(v, u) \in E$, entonces $f(v, u) = 0$: de lo contrario
  $c_f(u, v) = f(v, u) > 0$, y otra vez $v$ estaría en $S$. Si ninguna
  de las dos está en $E$, las dos valen $0$ en $f$ y en $c$. Así
  \[
    f(S, T) = \sum_{u \in S} \sum_{v \in T} c(u, v)
            - \sum_{u \in S} \sum_{v \in T} 0 = c(S, T) ,
  \]
  y por el lema~\ref{lem:flujo-neto-corte}, $\abs{f} = f(S, T) =
  c(S, T)$.

  (3) $\Rightarrow$ (1). Por el corolario~\ref{cor:flujo-cota},
  $\abs{f'} \le c(S, T)$ para todo flujo $f'$. Si $f$ alcanza esa
  cota, ningún flujo lo supera.
\end{proof}

Del teorema salen dos cosas. La primera, que el valor de un flujo
máximo es igual a la capacidad de un corte mínimo. La segunda, que ese
corte se obtiene con una búsqueda desde $s$ sobre la red residual
final, en $\Ord{n + m} = \Ord{m}$. El corte es un
\textbf{certificado}: cualquiera puede comprobar que un flujo es
máximo sumando unas cuantas capacidades, sin volver a correr nada.

\begin{idea}
  Las dos mitades del teorema hacen trabajos distintos. El corolario
  dice que ningún flujo pasa de ningún corte: los cortes son cotas
  superiores. La construcción de (2) $\Rightarrow$ (3) dice que,
  cuando ya no hay camino aumentante, lo alcanzable en $G_f$ forma un
  corte que el flujo llena exactamente: aristas de ida saturadas,
  aristas de vuelta vacías. Flujo y corte se encuentran, y ninguno de
  los dos se puede mejorar.
\end{idea}

% =====================================================================

\subsection*{El método de Ford-Fulkerson}

El método junta todo lo anterior: empezar con flujo cero y, mientras
la red residual tenga un camino aumentante, aumentar por él. Cormen lo
llama \emph{método} y no algoritmo porque no dice cómo se encuentra el
camino. Cada manera de buscarlo es una implementación distinta, con su
propio costo.

\begin{algorithm}[H]
\caption{\textsc{Ford-Fulkerson}$(G, s, t)$}
\label{alg:ford-fulkerson}
\begin{algorithmic}[1]
\Procedure{Ford-Fulkerson}{$G, s, t$}
  \ForAll{$(u, v) \in E$}
    \State $f(u, v) \gets 0$
  \EndFor
  \While{existe un camino $p$ de $s$ a $t$ en la red residual $G_f$}
    \State $c_f(p) \gets \min\set{c_f(u, v) : (u, v) \text{ está en }
           p}$
    \ForAll{$(u, v)$ en $p$}
      \If{$(u, v) \in E$}
        \State $f(u, v) \gets f(u, v) + c_f(p)$
          \Comment{empuja por una arista de $G$}
      \Else
        \State $f(v, u) \gets f(v, u) - c_f(p)$
          \Comment{cancela flujo de $(v, u)$}
      \EndIf
    \EndFor
  \EndWhile
  \State \Return $f$
\EndProcedure
\end{algorithmic}
\end{algorithm}

Las líneas que no son obvias:
\begin{itemize}[beginpenalty=10000]
  \item \textbf{Línea 5.} ``Existe un camino'' es cualquier búsqueda
        desde $s$ sobre $G_f$, BFS o DFS. Si llega a $t$ devuelve el
        camino; si no, el ciclo termina, y por el
        teorema~\ref{teo:flujo-corte} el flujo es máximo.
  \item \textbf{Línea 6.} El cuello de botella del camino.
  \item \textbf{Líneas 8 a 12.} Cada arista residual de $p$ es una
        arista de $G$ o la inversa de una, nunca las dos cosas, porque
        $G$ no tiene antiparalelas. En el primer caso se empuja flujo;
        en el segundo se cancela flujo que iba en sentido contrario.
        Es el aumento $f \uparrow f_p$ escrito arista por arista.
\end{itemize}

\textbf{Estructura de datos.} $G_f$ nunca se construye. Se guarda
$G' = (V, E')$, con
\[
  E' = \set{(u, v) : (u, v) \in E \text{ o } (v, u) \in E} ,
\]
así que la lista de cada vértice trae a sus vecinos de salida y
también a los de entrada. Cada entrada sabe si va hacia adelante o
de regreso y apunta a la arista original, que guarda $c$ y $f$. Así
$c_f(u, v)$ cuesta $\Ord{1}$: $c - f$ hacia adelante, $f$ de regreso.
Las aristas de $G_f$ son las de $G'$ con $c_f > 0$, y se filtran
durante la búsqueda. Sin antiparalelas, $\abs{E'} = 2m$.

\textbf{Costo, término por término.}
\begin{itemize}[beginpenalty=10000]
  \item Líneas 2 a 4: $\Tht{m}$.
  \item Línea 5: una búsqueda sobre $G'$ cuesta $\Ord{n + 2m}$, y
        como $m \ge n - 1$, eso es $\Ord{m}$.
  \item Líneas 6 a 13: $p$ no repite vértices, así que tiene a lo más
        $n - 1$ aristas: $\Ord{n} \subseteq \Ord{m}$.
  \item Número de iteraciones. Si las capacidades son enteras, por
        inducción $f$ y todas las $c_f$ son enteras: empiezan en $0$ y
        cada actualización suma o resta el entero $c_f(p)$. Entonces
        $c_f(p) \ge 1$ y, por el lema~\ref{lem:flujo-camino},
        $\abs{f}$ sube al menos $1$ en cada iteración. Como nunca pasa
        de $\abs{f^*}$, hay a lo más $\abs{f^*}$ iteraciones.
\end{itemize}

\begin{resultado}
  Con capacidades enteras, Ford-Fulkerson hace a lo más $\abs{f^*}$
  aumentos y cuesta
  \[
    \Ord{m} + \abs{f^*} \cdot \Ord{m} = \Ord{m \abs{f^*}} .
  \]
  El costo depende del valor del flujo y no solo del tamaño de la red:
  es pseudopolinomial, porque $\abs{f^*}$ puede ser exponencial en el
  número de bits de la entrada.
\end{resultado}

\begin{corolario}[Integralidad]\label{cor:flujo-integralidad}
  Si todas las capacidades son enteras, Ford-Fulkerson devuelve un
  flujo máximo con $f(u, v)$ entero en cada arista.
\end{corolario}

\begin{proof}
  Es la misma inducción del costo: $f$ empieza en $0$ y cada
  actualización suma o resta un entero. Al terminar, el
  teorema~\ref{teo:flujo-corte} dice que el flujo es máximo.
\end{proof}

Cormen lo enuncia hasta la sección 26.3 (teorema 26.10), y es lo que
hace funcionar el emparejamiento bipartito de la
sección~\ref{sec:emparejamiento}. Con capacidades racionales todo se
multiplica por el común denominador y el argumento sigue igual. Con
irracionales el método puede no terminar, y Cormen advierte que el
valor puede ni siquiera converger al máximo.

\begin{ejemplo}[La ejecución de la figura 26.6 de
  Cormen]\label{ej:flujo-ff266}
  Ford-Fulkerson sobre la red de Cormen, desde flujo cero, con los
  caminos que escoge Cormen. Primero el resumen; en la última columna,
  el valor después de cada aumento.
  \begin{center}
  \small
  \begin{tabular}{clcll}
    \toprule
    & \textbf{Camino aumentante $p$} & $c_f(p)$
      & \textbf{Arista crítica} & $\abs{f}$ \\
    \midrule
    (a) & $s, v_1, v_3, v_2, v_4, t$ & $\min(16, 12, 9, 14, 4) = 4$
        & $(v_4, t)$ & $4$ \\
    (b) & $s, v_2, v_1, v_3, t$ & $\min(13, 4, 8, 20) = 4$
        & $(v_2, v_1)$ & $8$ \\
    (c) & $s, v_1, v_2, v_3, t$ & $\min(12, 4, 4, 16) = 4$
        & $(v_1, v_2)$ y $(v_2, v_3)$ & $12$ \\
    (d) & $s, v_2, v_4, v_3, t$ & $\min(9, 10, 7, 12) = 7$
        & $(v_4, v_3)$ & $19$ \\
    (e) & $s, v_1, v_3, t$ & $\min(8, 4, 5) = 4$
        & $(v_1, v_3)$ & $23$ \\
    (f) & no hay camino & & & $23$ \\
    \bottomrule
  \end{tabular}
  \end{center}
  En cada renglón de figuras, a la izquierda está $G_f$ con el camino
  aumentante en azul y a la derecha el flujo que resulta del aumento.

  \noindent\begin{minipage}{\linewidth}
    \textbf{(a)} El flujo es cero, así que $G_f = G$ y cada arista
    residual lleva la capacidad completa. El camino tiene cinco
    aristas, y la más angosta es $(v_4, t)$.
    \begin{center}
    \begin{tikzpicture}
      \foreach \n/\x/\y/\e in {s/0/0/s, v1/1.8/1.2/v_1,
          v2/1.8/-1.2/v_2, v3/4.2/1.2/v_3, v4/4.2/-1.2/v_4,
          t/6/0/t}
        \node[vert] (\n) at (\x, \y) {$\e$};
      \draw[dirigida] (s) -- node[peso] {13} (v2);
      \draw[dirigida] (v2) -- node[peso] {4} (v1);
      \draw[dirigida] (v4) -- node[peso] {7} (v3);
      \draw[dirigida] (v3) -- node[peso] {20} (t);
      \draw[dirigidaArbol] (s) -- node[peso] {16} (v1);
      \draw[dirigidaArbol] (v1) -- node[peso] {12} (v3);
      \draw[dirigidaArbol] (v3) -- node[peso] {9} (v2);
      \draw[dirigidaArbol] (v2) -- node[peso] {14} (v4);
      \draw[dirigidaArbol] (v4) -- node[peso] {4} (t);
    \end{tikzpicture}
    \hspace{0.8cm}
    \begin{tikzpicture}
      \foreach \n/\x/\y/\e in {s/0/0/s, v1/1.8/1.2/v_1,
          v2/1.8/-1.2/v_2, v3/4.2/1.2/v_3, v4/4.2/-1.2/v_4,
          t/6/0/t}
        \node[vert] (\n) at (\x, \y) {$\e$};
      \draw[dirigida] (s) -- node[peso] {4/16} (v1);
      \draw[dirigida] (s) -- node[peso] {0/13} (v2);
      \draw[dirigida] (v1) -- node[peso] {4/12} (v3);
      \draw[dirigida] (v2) -- node[peso] {0/4} (v1);
      \draw[dirigida] (v3) -- node[peso] {4/9} (v2);
      \draw[dirigida] (v2) -- node[peso] {4/14} (v4);
      \draw[dirigida] (v4) -- node[peso] {0/7} (v3);
      \draw[dirigida] (v3) -- node[peso] {0/20} (t);
      \draw[dirigida] (v4) -- node[peso] {4/4} (t);
    \end{tikzpicture}
    \end{center}
  \end{minipage}

  \noindent\begin{minipage}{\linewidth}
    \textbf{(b)} Ya aparecen aristas de regreso, una por cada arista
    con flujo: $(v_1, s)$, $(v_3, v_1)$, $(v_2, v_3)$, $(v_4, v_2)$ y
    $(t, v_4)$. El camino satura $(v_2, v_1)$.
    \begin{center}
    \begin{tikzpicture}
      \foreach \n/\x/\y/\e in {s/0/0/s, v1/1.8/1.2/v_1,
          v2/1.8/-1.2/v_2, v3/4.2/1.2/v_3, v4/4.2/-1.2/v_4,
          t/6/0/t}
        \node[vert] (\n) at (\x, \y) {$\e$};
      \draw[dirigida] (s) to[bend left=16] node[peso, auto] {12} (v1);
      \draw[dirigida] (v1) to[bend left=16] node[peso, auto] {4} (s);
      \draw[dirigida] (v3) to[bend left=16] node[peso, auto] {4} (v1);
      \draw[dirigida] (v3) to[bend left=16] node[peso, auto] {5} (v2);
      \draw[dirigida] (v2) to[bend left=16] node[peso, auto] {4} (v3);
      \draw[dirigida] (v2) to[bend left=16] node[peso, auto] {10} (v4);
      \draw[dirigida] (v4) to[bend left=16] node[peso, auto] {4} (v2);
      \draw[dirigida] (v4) -- node[peso] {7} (v3);
      \draw[dirigida] (t) -- node[peso] {4} (v4);
      \draw[dirigidaArbol] (s) -- node[peso] {13} (v2);
      \draw[dirigidaArbol] (v1) to[bend left=16] node[peso, auto] {8} (v3);
      \draw[dirigidaArbol] (v2) -- node[peso] {4} (v1);
      \draw[dirigidaArbol] (v3) -- node[peso] {20} (t);
    \end{tikzpicture}
    \hspace{0.8cm}
    \begin{tikzpicture}
      \foreach \n/\x/\y/\e in {s/0/0/s, v1/1.8/1.2/v_1,
          v2/1.8/-1.2/v_2, v3/4.2/1.2/v_3, v4/4.2/-1.2/v_4,
          t/6/0/t}
        \node[vert] (\n) at (\x, \y) {$\e$};
      \draw[dirigida] (s) -- node[peso] {4/16} (v1);
      \draw[dirigida] (s) -- node[peso] {4/13} (v2);
      \draw[dirigida] (v1) -- node[peso] {8/12} (v3);
      \draw[dirigida] (v2) -- node[peso] {4/4} (v1);
      \draw[dirigida] (v3) -- node[peso] {4/9} (v2);
      \draw[dirigida] (v2) -- node[peso] {4/14} (v4);
      \draw[dirigida] (v4) -- node[peso] {0/7} (v3);
      \draw[dirigida] (v3) -- node[peso] {4/20} (t);
      \draw[dirigida] (v4) -- node[peso] {4/4} (t);
    \end{tikzpicture}
    \end{center}
  \end{minipage}

  \noindent\begin{minipage}{\linewidth}
    \textbf{(c)} El camino usa dos aristas de regreso: $(v_1, v_2)$,
    inversa de $(v_2, v_1)$, y $(v_2, v_3)$, inversa de $(v_3, v_2)$.
    El aumento cancela las $4$ unidades de cada una, y $f(v_2, v_1)$
    y $f(v_3, v_2)$ vuelven a $0$ (en naranja). El flujo de valor $12$
    se lee ahora como $8$ unidades por $s, v_1, v_3, t$ y $4$ por
    $s, v_2, v_4, t$: el método deshizo el rodeo de (a) y (b).
    \begin{center}
    \begin{tikzpicture}
      \foreach \n/\x/\y/\e in {s/0/0/s, v1/1.8/1.2/v_1,
          v2/1.8/-1.2/v_2, v3/4.2/1.2/v_3, v4/4.2/-1.2/v_4,
          t/6/0/t}
        \node[vert] (\n) at (\x, \y) {$\e$};
      \draw[dirigida] (v1) to[bend left=16] node[peso, auto] {4} (s);
      \draw[dirigida] (s) to[bend left=16] node[peso, auto] {9} (v2);
      \draw[dirigida] (v2) to[bend left=16] node[peso, auto] {4} (s);
      \draw[dirigida] (v1) to[bend left=16] node[peso, auto] {4} (v3);
      \draw[dirigida] (v3) to[bend left=16] node[peso, auto] {8} (v1);
      \draw[dirigida] (v3) to[bend left=16] node[peso, auto] {5} (v2);
      \draw[dirigida] (v2) to[bend left=16] node[peso, auto] {10} (v4);
      \draw[dirigida] (v4) to[bend left=16] node[peso, auto] {4} (v2);
      \draw[dirigida] (v4) -- node[peso] {7} (v3);
      \draw[dirigida] (t) to[bend left=16] node[peso, auto] {4} (v3);
      \draw[dirigida] (t) -- node[peso] {4} (v4);
      \draw[dirigidaArbol] (s) to[bend left=16] node[peso, auto] {12} (v1);
      \draw[dirigidaArbol] (v1) -- node[peso] {4} (v2);
      \draw[dirigidaArbol] (v2) to[bend left=16] node[peso, auto] {4} (v3);
      \draw[dirigidaArbol] (v3) to[bend left=16] node[peso, auto] {16} (t);
    \end{tikzpicture}
    \hspace{0.8cm}
    \begin{tikzpicture}
      \foreach \n/\x/\y/\e in {s/0/0/s, v1/1.8/1.2/v_1,
          v2/1.8/-1.2/v_2, v3/4.2/1.2/v_3, v4/4.2/-1.2/v_4,
          t/6/0/t}
        \node[vert] (\n) at (\x, \y) {$\e$};
      \draw[dirigida] (s) -- node[peso] {8/16} (v1);
      \draw[dirigida] (s) -- node[peso] {4/13} (v2);
      \draw[dirigida] (v1) -- node[peso] {8/12} (v3);
      \draw[dirigidaResalte] (v2) -- node[peso] {0/4} (v1);
      \draw[dirigidaResalte] (v3) -- node[peso] {0/9} (v2);
      \draw[dirigida] (v2) -- node[peso] {4/14} (v4);
      \draw[dirigida] (v4) -- node[peso] {0/7} (v3);
      \draw[dirigida] (v3) -- node[peso] {8/20} (t);
      \draw[dirigida] (v4) -- node[peso] {4/4} (t);
    \end{tikzpicture}
    \end{center}
  \end{minipage}

  \noindent\begin{minipage}{\linewidth}
    \textbf{(d)} Un camino de cuatro aristas, sin regresos, que
    satura $(v_4, v_3)$.
    \begin{center}
    \begin{tikzpicture}
      \foreach \n/\x/\y/\e in {s/0/0/s, v1/1.8/1.2/v_1,
          v2/1.8/-1.2/v_2, v3/4.2/1.2/v_3, v4/4.2/-1.2/v_4,
          t/6/0/t}
        \node[vert] (\n) at (\x, \y) {$\e$};
      \draw[dirigida] (s) to[bend left=16] node[peso, auto] {8} (v1);
      \draw[dirigida] (v1) to[bend left=16] node[peso, auto] {8} (s);
      \draw[dirigida] (v2) to[bend left=16] node[peso, auto] {4} (s);
      \draw[dirigida] (v1) to[bend left=16] node[peso, auto] {4} (v3);
      \draw[dirigida] (v3) to[bend left=16] node[peso, auto] {8} (v1);
      \draw[dirigida] (v2) -- node[peso] {4} (v1);
      \draw[dirigida] (v3) -- node[peso] {9} (v2);
      \draw[dirigida] (v4) to[bend left=16] node[peso, auto] {4} (v2);
      \draw[dirigida] (t) to[bend left=16] node[peso, auto] {8} (v3);
      \draw[dirigida] (t) -- node[peso] {4} (v4);
      \draw[dirigidaArbol] (s) to[bend left=16] node[peso, auto] {9} (v2);
      \draw[dirigidaArbol] (v2) to[bend left=16] node[peso, auto] {10} (v4);
      \draw[dirigidaArbol] (v4) -- node[peso] {7} (v3);
      \draw[dirigidaArbol] (v3) to[bend left=16] node[peso, auto] {12} (t);
    \end{tikzpicture}
    \hspace{0.8cm}
    \begin{tikzpicture}
      \foreach \n/\x/\y/\e in {s/0/0/s, v1/1.8/1.2/v_1,
          v2/1.8/-1.2/v_2, v3/4.2/1.2/v_3, v4/4.2/-1.2/v_4,
          t/6/0/t}
        \node[vert] (\n) at (\x, \y) {$\e$};
      \draw[dirigida] (s) -- node[peso] {8/16} (v1);
      \draw[dirigida] (s) -- node[peso] {11/13} (v2);
      \draw[dirigida] (v1) -- node[peso] {8/12} (v3);
      \draw[dirigida] (v2) -- node[peso] {0/4} (v1);
      \draw[dirigida] (v3) -- node[peso] {0/9} (v2);
      \draw[dirigida] (v2) -- node[peso] {11/14} (v4);
      \draw[dirigida] (v4) -- node[peso] {7/7} (v3);
      \draw[dirigida] (v3) -- node[peso] {15/20} (t);
      \draw[dirigida] (v4) -- node[peso] {4/4} (t);
    \end{tikzpicture}
    \end{center}
  \end{minipage}

  \noindent\begin{minipage}{\linewidth}
    \textbf{(e)} El último aumento satura $(v_1, v_3)$. El flujo de la
    derecha vale $23$.
    \begin{center}
    \begin{tikzpicture}
      \foreach \n/\x/\y/\e in {s/0/0/s, v1/1.8/1.2/v_1,
          v2/1.8/-1.2/v_2, v3/4.2/1.2/v_3, v4/4.2/-1.2/v_4,
          t/6/0/t}
        \node[vert] (\n) at (\x, \y) {$\e$};
      \draw[dirigida] (v1) to[bend left=16] node[peso, auto] {8} (s);
      \draw[dirigida] (s) to[bend left=16] node[peso, auto] {2} (v2);
      \draw[dirigida] (v2) to[bend left=16] node[peso, auto] {11} (s);
      \draw[dirigida] (v3) to[bend left=16] node[peso, auto] {8} (v1);
      \draw[dirigida] (v2) -- node[peso] {4} (v1);
      \draw[dirigida] (v3) -- node[peso] {9} (v2);
      \draw[dirigida] (v2) to[bend left=16] node[peso, auto] {3} (v4);
      \draw[dirigida] (v4) to[bend left=16] node[peso, auto] {11} (v2);
      \draw[dirigida] (v3) -- node[peso] {7} (v4);
      \draw[dirigida] (t) to[bend left=16] node[peso, auto] {15} (v3);
      \draw[dirigida] (t) -- node[peso] {4} (v4);
      \draw[dirigidaArbol] (s) to[bend left=16] node[peso, auto] {8} (v1);
      \draw[dirigidaArbol] (v1) to[bend left=16] node[peso, auto] {4} (v3);
      \draw[dirigidaArbol] (v3) to[bend left=16] node[peso, auto] {5} (t);
    \end{tikzpicture}
    \hspace{0.8cm}
    \begin{tikzpicture}
      \foreach \n/\x/\y/\e in {s/0/0/s, v1/1.8/1.2/v_1,
          v2/1.8/-1.2/v_2, v3/4.2/1.2/v_3, v4/4.2/-1.2/v_4,
          t/6/0/t}
        \node[vert] (\n) at (\x, \y) {$\e$};
      \draw[dirigida] (s) -- node[peso] {12/16} (v1);
      \draw[dirigida] (s) -- node[peso] {11/13} (v2);
      \draw[dirigida] (v1) -- node[peso] {12/12} (v3);
      \draw[dirigida] (v2) -- node[peso] {0/4} (v1);
      \draw[dirigida] (v3) -- node[peso] {0/9} (v2);
      \draw[dirigida] (v2) -- node[peso] {11/14} (v4);
      \draw[dirigida] (v4) -- node[peso] {7/7} (v3);
      \draw[dirigida] (v3) -- node[peso] {19/20} (t);
      \draw[dirigida] (v4) -- node[peso] {4/4} (t);
    \end{tikzpicture}
    \end{center}
  \end{minipage}

  \noindent\begin{minipage}{\linewidth}
    \textbf{(f)} La red residual final. Desde $s$ se llega a $v_1$ y a
    $v_2$; desde $v_2$, a $v_4$; desde $v_1$ y $v_4$ solo se regresa.
    $t$ no es alcanzable y el ciclo termina con $\abs{f^*} = 23$. En
    verde, los vértices alcanzables desde $s$.
    \begin{center}
    \begin{tikzpicture}
      \node[vertOk] (s) at (0, 0) {$s$};
      \node[vertOk] (v1) at (1.8, 1.2) {$v_1$};
      \node[vertOk] (v2) at (1.8, -1.2) {$v_2$};
      \node[vert] (v3) at (4.2, 1.2) {$v_3$};
      \node[vertOk] (v4) at (4.2, -1.2) {$v_4$};
      \node[vert] (t) at (6, 0) {$t$};
      \draw[dirigida] (s) to[bend left=16] node[peso, auto] {4} (v1);
      \draw[dirigida] (v1) to[bend left=16] node[peso, auto] {12} (s);
      \draw[dirigida] (s) to[bend left=16] node[peso, auto] {2} (v2);
      \draw[dirigida] (v2) to[bend left=16] node[peso, auto] {11} (s);
      \draw[dirigida] (v3) -- node[peso] {12} (v1);
      \draw[dirigida] (v2) -- node[peso] {4} (v1);
      \draw[dirigida] (v3) -- node[peso] {9} (v2);
      \draw[dirigida] (v2) to[bend left=16] node[peso, auto] {3} (v4);
      \draw[dirigida] (v4) to[bend left=16] node[peso, auto] {11} (v2);
      \draw[dirigida] (v3) -- node[peso] {7} (v4);
      \draw[dirigida] (v3) to[bend left=16] node[peso, auto] {1} (t);
      \draw[dirigida] (t) to[bend left=16] node[peso, auto] {19} (v3);
      \draw[dirigida] (t) -- node[peso] {4} (v4);
    \end{tikzpicture}
    \end{center}
  \end{minipage}
\end{ejemplo}

\begin{ojo}
  Los caminos de la figura 26.6 no salen de un BFS. Cormen no dice
  cómo los eligió, y el primero tiene cinco aristas cuando hay caminos
  de tres, como $s, v_1, v_3, t$. Quien intente reproducir la figura
  con un BFS obtiene otra ejecución, la de Edmonds-Karp, que aparece
  más adelante. Estos caminos sirven para que se vea una cancelación,
  y cualquier elección de caminos es una ejecución válida del método.
\end{ojo}

\begin{ejemplo}[El corte mínimo]\label{ej:flujo-corte-minimo}
  Los vértices verdes de la red residual final forman
  $S = \set{s, v_1, v_2, v_4}$, con $T = \set{v_3, t}$. A la derecha
  está ese corte sobre el flujo máximo; a la izquierda, para
  contrastar, el corte del ejemplo~\ref{ej:flujo-corte} sobre el flujo
  de valor $19$. En verde, los vértices de $S$.
  \begin{center}
  \begin{tikzpicture}
    \foreach \n/\x/\y/\e in {s/0/0/s, v1/1.8/1.2/v_1,
        v2/1.8/-1.2/v_2}
      \node[vertOk] (\n) at (\x, \y) {$\e$};
    \foreach \n/\x/\y/\e in {v3/4.2/1.2/v_3, v4/4.2/-1.2/v_4, t/6/0/t}
      \node[vert] (\n) at (\x, \y) {$\e$};
    \foreach \a/\b/\f/\c in {s/v1/11/16, s/v2/8/13, v2/v1/1/4,
        v4/v3/7/7, v3/t/15/20, v4/t/4/4}
      \draw[dirigida] (\a) -- node[peso] {\f/\c} (\b);
    % las tres que cruzan el corte llevan la etiqueta lejos de la linea
    \foreach \a/\b/\f/\c/\p in {v1/v3/12/12/0.65, v3/v2/4/9/0.35,
        v2/v4/11/14/0.65}
      \draw[dirigida] (\a) -- node[peso, pos=\p] {\f/\c} (\b);
    \draw[corte] (2.5, 1.9) -- (2.5, -1.9);
    \node[nota] at (2.2, -2.1) {$S$};
    \node[nota] at (2.8, -2.1) {$T$};
  \end{tikzpicture}
  \hspace{0.8cm}
  \begin{tikzpicture}
    \foreach \n/\x/\y/\e in {s/0/0/s, v1/1.8/1.2/v_1,
        v2/1.8/-1.2/v_2, v4/4.2/-1.2/v_4}
      \node[vertOk] (\n) at (\x, \y) {$\e$};
    \foreach \n/\x/\y/\e in {v3/4.2/1.2/v_3, t/6/0/t}
      \node[vert] (\n) at (\x, \y) {$\e$};
    \foreach \a/\b/\f/\c in {s/v1/12/16, s/v2/11/13, v2/v1/0/4,
        v3/v2/0/9, v2/v4/11/14, v4/v3/7/7, v3/t/19/20, v4/t/4/4}
      \draw[dirigida] (\a) -- node[peso] {\f/\c} (\b);
    \draw[dirigida] (v1) -- node[peso, pos=0.4] {12/12} (v3);
    \draw[corte] plot[smooth] coordinates {(3.65, 1.95) (3.65, 1.2)
      (3.8, 0.6) (4.2, 0.45) (4.9, 0.15) (5.45, -0.37) (5.75, -1.0)
      (5.8, -1.7)};
    \node[nota] at (3.3, 1.9) {$S$};
    \node[nota] at (4.0, 1.9) {$T$};
  \end{tikzpicture}
  \end{center}
  De $S$ a $T$ cruzan $(v_1, v_3)$ con $12/12$, $(v_4, v_3)$ con $7/7$
  y $(v_4, t)$ con $4/4$: todas saturadas, como exige la prueba de
  (2) $\Rightarrow$ (3). La única de $T$ a $S$ es $(v_3, v_2)$, con
  $0/9$: vacía. Entonces
  \[
    f(S, T) = 12 + 7 + 4 - 0 = 23 = c(S, T) .
  \]
  En el corte de la izquierda, en cambio, sobra holgura en
  $(v_2, v_4)$, con $11/14$, y regresan $4$ unidades por
  $(v_3, v_2)$. La red tiene $16$ cortes; los seis más baratos son
  \begin{center}
  \small
  \begin{tabular}{lclc}
    \toprule
    $S$ & $c(S, T)$ & $S$ & $c(S, T)$ \\
    \midrule
    $\set{s, v_1, v_2, v_4}$ & $23$ & $\set{s, v_1, v_2}$ & $26$ \\
    $\set{s, v_1, v_2, v_3, v_4}$ & $24$ & $\set{s}$ & $29$ \\
    $\set{s, v_1}$ & $25$ & $\set{s, v_2, v_4}$ & $31$ \\
    \bottomrule
  \end{tabular}
  \end{center}
  y el mínimo es único. La lectura física: a Saskatoon solo pueden
  entrar $12 + 7 = 19$ cajas diarias, aunque la ruta a Winnipeg
  aguante $20$, y por Regina pasan directo $4$ más. No hay manera de
  enrutar que pase de $23$. También se ve por qué la arista extra
  $(v_1, v_2)$ del apartado de antiparalelas no ayuda: sus dos
  extremos están en $S$, así que no cruza el corte mínimo.
\end{ejemplo}

\begin{idea}
  Ford-Fulkerson parece un voraz: toma un camino cualquiera y lo
  llena hasta su cuello de botella. La diferencia está en las aristas
  de regreso. La decisión de la iteración (a), mandar $4$ unidades por
  $(v_3, v_2)$, se deshizo en (c), y el voraz del
  ejemplo~\ref{ej:flujo-voraz}, que no podía deshacerla, se atoró en
  $19$. Dijkstra y Kruskal nunca revisan una decisión, y no les hace
  falta: Dijkstra porque con pesos no negativos el vértice que fija ya
  tiene su distancia definitiva, Kruskal porque la arista más ligera
  que cruza un corte siempre es segura. En flujo no hay una propiedad
  así para los caminos: ninguna elección local es segura. Las aristas
  de regreso vuelven provisional cada decisión, y el
  teorema~\ref{teo:flujo-corte} garantiza que, cuando ya no hay
  camino, no queda nada que corregir.
\end{idea}

% =====================================================================

\subsection*{Cuando la elección del camino importa}

Ford-Fulkerson no dice cómo escoger el camino, y eso decide el costo.
La cota $\Ord{m \abs{f^*}}$ no es pesimismo del análisis: hay redes
donde se alcanza.

\begin{ejemplo}[El caso patológico; Cormen, figura
  26.7]\label{ej:flujo-patologico}
  La red tiene $V = \set{s, u, v, t}$, con
  $c(s, u) = c(s, v) = c(u, t) = c(v, t) = 10^6$ y $c(u, v) = 1$. El
  flujo máximo vale $2 \cdot 10^6$: $10^6$ por $s, u, t$ y $10^6$ por
  $s, v, t$, y el corte $(\set{s}, \set{u, v, t})$ tiene capacidad
  $2 \cdot 10^6$. Supóngase que Ford-Fulkerson elige $s, u, v, t$ en
  las iteraciones impares y $s, v, u, t$ en las pares. Cada camino
  pasa por $(u, v)$ o por su inversa, de capacidad residual $1$, así
  que el flujo sube de uno en uno.
  \begin{center}
  \small
  \begin{tabular}{clccc}
    \toprule
    $i$ & \textbf{Camino} & $c_f(p)$ & $f(u, v)$ después & $\abs{f}$ \\
    \midrule
    $1$ & $s, u, v, t$ & $1$ & $1$ & $1$ \\
    $2$ & $s, v, u, t$ & $1$ & $0$ & $2$ \\
    $3$ & $s, u, v, t$ & $1$ & $1$ & $3$ \\
    $\vdots$ & & & & $\vdots$ \\
    $2 \cdot 10^6$ & $s, v, u, t$ & $1$ & $0$ & $2 \cdot 10^6$ \\
    \bottomrule
  \end{tabular}
  \end{center}
  Las tres instantáneas son la red original con el primer camino, la
  residual después del primer aumento con el segundo camino, y la
  residual después del segundo. En la tercera, $(u, v)$ está otra vez
  disponible con capacidad $1$ y las cuatro aristas de los lados
  perdieron una unidad: el ciclo se repite.
  \begin{center}
  \begin{tikzpicture}
    \node[vert] (s) at (0, 0) {$s$};
    \node[vert] (u) at (2.2, 1.3) {$u$};
    \node[vert] (v) at (2.2, -1.3) {$v$};
    \node[vert] (t) at (4.4, 0) {$t$};
    \draw[dirigida] (s) -- node[peso] {$10^6$} (v);
    \draw[dirigida] (u) -- node[peso] {$10^6$} (t);
    \draw[dirigidaArbol] (s) -- node[peso] {$10^6$} (u);
    \draw[dirigidaArbol] (u) -- node[peso] {1} (v);
    \draw[dirigidaArbol] (v) -- node[peso] {$10^6$} (t);
  \end{tikzpicture}
  \hspace{0.4cm}
  \begin{tikzpicture}
    \node[vert] (s) at (0, 0) {$s$};
    \node[vert] (u) at (2.2, 1.3) {$u$};
    \node[vert] (v) at (2.2, -1.3) {$v$};
    \node[vert] (t) at (4.4, 0) {$t$};
    \draw[dirigida] (s) to[bend left=18] node[peso, auto] {$10^6{-}1$} (u);
    \draw[dirigida] (u) to[bend left=18] node[peso, auto] {1} (s);
    \draw[dirigida] (v) to[bend right=18]
      node[peso, auto, swap] {$10^6{-}1$} (t);
    \draw[dirigida] (t) to[bend right=18] node[peso, auto, swap] {1} (v);
    \draw[dirigidaArbol] (s) -- node[peso] {$10^6$} (v);
    \draw[dirigidaArbol] (v) -- node[peso] {1} (u);
    \draw[dirigidaArbol] (u) -- node[peso] {$10^6$} (t);
  \end{tikzpicture}
  \hspace{0.4cm}
  \begin{tikzpicture}
    \node[vert] (s) at (0, 0) {$s$};
    \node[vert] (u) at (2.2, 1.3) {$u$};
    \node[vert] (v) at (2.2, -1.3) {$v$};
    \node[vert] (t) at (4.4, 0) {$t$};
    \draw[dirigida] (s) to[bend left=18] node[peso, auto] {$10^6{-}1$} (u);
    \draw[dirigida] (u) to[bend left=18] node[peso, auto] {1} (s);
    \draw[dirigida] (s) to[bend right=18]
      node[peso, auto, swap] {$10^6{-}1$} (v);
    \draw[dirigida] (v) to[bend right=18] node[peso, auto, swap] {1} (s);
    \draw[dirigida] (u) -- node[peso] {1} (v);
    \draw[dirigida] (u) to[bend left=18] node[peso, auto] {$10^6{-}1$} (t);
    \draw[dirigida] (t) to[bend left=18] node[peso, auto] {1} (u);
    \draw[dirigida] (v) to[bend right=18]
      node[peso, auto, swap] {$10^6{-}1$} (t);
    \draw[dirigida] (t) to[bend right=18] node[peso, auto, swap] {1} (v);
  \end{tikzpicture}
  \end{center}
  Nada en el método impide esa elección; la línea 5 solo pide que
  exista un camino. Son $2 \cdot 10^6$ iteraciones de $\Ord{m}$ cada
  una, y la cota $\Tht{m \abs{f^*}}$ se alcanza. Con $C$ en lugar de
  $10^6$ la estrategia alternada hace $2C$ aumentos; se comprobó
  simulando con $C = 1, 2, 3, 10, 1000$ y $10^6$. Mientras tanto, los
  caminos $s, u, t$ y $s, v, t$, de dos aristas cada uno, resuelven la
  red en dos aumentos.
\end{ejemplo}

% =====================================================================

\subsection*{Edmonds-Karp}

Edmonds-Karp fija la elección: el camino aumentante es siempre un
camino de $s$ a $t$ con el menor número de aristas en $G_f$, y lo
encuentra un BFS. En el ejemplo patológico eso ya basta. Los caminos
más cortos son $s, u, t$ y $s, v, t$, de dos aristas, y el de tres que
pasa por $(u, v)$ nunca se toca: dos aumentos contra dos millones.

El BFS es el de la clase del 22 de octubre
(sección~\ref{sec:grafos}; a fondo en la sección~\ref{sec:busquedas})
con dos cambios: recorre las listas de $G'$, que incluyen a los
vecinos de entrada, y solo cruza $(u, v)$ si $c_f(u, v) > 0$. Recibe
$f$ porque con él calcula cada $c_f(u, v)$ en $\Ord{1}$, con la
estructura $G'$ de Ford-Fulkerson; la red residual sigue sin
construirse. Su $d[v]$ es la \textbf{distancia residual}
$\delta_f(s, v)$: el número de aristas de un camino más corto de $s$ a
$v$ en $G_f$, o $\infty$ si no hay camino.

\begin{algorithm}[H]
\caption{\textsc{BFS-Residual}$(G, f, s)$}
\label{alg:flujo-bfs-residual}
\begin{algorithmic}[1]
\Procedure{BFS-Residual}{$G, f, s$}
  \ForAll{$v \in V$}
    \State $\text{visto}[v] \gets$ falso, \quad $d[v] \gets \infty$,
           \quad $\pi[v] \gets \textsc{nulo}$
  \EndFor
  \State $\text{visto}[s] \gets$ verdadero, \quad $d[s] \gets 0$
  \State $Q \gets$ cola vacía; \ \textsc{Encolar}$(Q, s)$
  \While{$Q$ no está vacía}
    \State $u \gets \textsc{Desencolar}(Q)$
    \ForAll{$v$ con $(u, v) \in E$ o $(v, u) \in E$}
      \Comment{lista de $u$ en $G'$}
      \If{\textbf{no} $\text{visto}[v]$ \textbf{y} $c_f(u, v) > 0$}
        \State $\text{visto}[v] \gets$ verdadero, \quad
               $d[v] \gets d[u] + 1$, \quad $\pi[v] \gets u$
        \State \textsc{Encolar}$(Q, v)$
      \EndIf
    \EndFor
  \EndWhile
\EndProcedure
\end{algorithmic}
\end{algorithm}

\begin{algorithm}[H]
\caption{\textsc{Edmonds-Karp}$(G, s, t)$}
\label{alg:edmonds-karp}
\begin{algorithmic}[1]
\Procedure{Edmonds-Karp}{$G, s, t$}
  \ForAll{$(u, v) \in E$}
    \State $f(u, v) \gets 0$
  \EndFor
  \Loop
    \State \textsc{BFS-Residual}$(G, f, s)$
      \Comment{llena visto, $d$ y $\pi$}
    \If{\textbf{no} $\text{visto}[t]$}
      \State \Return $f$ \Comment{no hay camino: $f$ es máximo}
    \EndIf
    \State $\Delta \gets \infty$, \quad $v \gets t$
    \While{$v \neq s$} \Comment{primera pasada: cuello de botella}
      \State $\Delta \gets \min(\Delta, c_f(\pi[v], v))$, \quad
             $v \gets \pi[v]$
    \EndWhile
    \State $v \gets t$
    \While{$v \neq s$} \Comment{segunda pasada: aumento}
      \State $u \gets \pi[v]$
      \If{$(u, v) \in E$}
        \State $f(u, v) \gets f(u, v) + \Delta$
      \Else
        \State $f(v, u) \gets f(v, u) - \Delta$
      \EndIf
      \State $v \gets u$
    \EndWhile
  \EndLoop
\EndProcedure
\end{algorithmic}
\end{algorithm}

Las líneas que no son obvias en Edmonds-Karp:
\begin{itemize}[beginpenalty=10000]
  \item \textbf{Línea 6.} Después del BFS, $d[v] = \delta_f(s, v)$ y
        siguiendo $\pi$ desde $t$ se recorre un camino más corto, al
        revés.
  \item \textbf{Líneas 10 a 13.} Primera pasada de $t$ a $s$ para
        saber el cuello de botella; al salir, $\Delta = c_f(p)$. Hacen
        falta dos pasadas porque $c_f(p)$ solo se conoce después de
        recorrer todo el camino.
  \item \textbf{Líneas 14 a 23.} Segunda pasada: la misma
        actualización que las líneas 7 a 13 de Ford-Fulkerson.
  \item El BFS puede detenerse en cuanto descubre $t$. Ahorra trabajo
        y no cambia la cota.
\end{itemize}

\begin{ejemplo}[Edmonds-Karp sobre la red de Cormen; ejercicio
  26.2-3]\label{ej:flujo-ek}
  Misma red, pero ahora cada camino lo elige un BFS que revisa los
  vecinos en el orden $s, v_1, v_2, v_3, v_4, t$. La tabla da las
  distancias residuales antes de cada aumento.
  \begin{center}
  \small
  \begin{tabular}{cccccclccc}
    \toprule
    & \multicolumn{5}{c}{$\delta_f(s, \cdot)$ antes del aumento}
    & & & & \\
    \cmidrule(lr){2-6}
    \textbf{Iter.} & $v_1$ & $v_2$ & $v_3$ & $v_4$ & $t$
    & \textbf{Camino} & $c_f(p)$ & \textbf{Crítica} & $\abs{f}$ \\
    \midrule
    1 & 1 & 1 & 2 & 2 & 3 & $s, v_1, v_3, t$ & 12 & $(v_1, v_3)$
      & 12 \\
    2 & 1 & 1 & 3 & 2 & 3 & $s, v_2, v_4, t$ & 4 & $(v_4, t)$ & 16 \\
    3 & 1 & 1 & 3 & 2 & 4 & $s, v_2, v_4, v_3, t$ & 7 & $(v_4, v_3)$
      & 23 \\
    fin & 1 & 1 & $\infty$ & 2 & $\infty$ & no hay camino & & & \\
    \bottomrule
  \end{tabular}
  \end{center}
  Las redes residuales de las iteraciones 1 y 2, con su camino en
  azul:
  \begin{center}
  \begin{tikzpicture}
    \foreach \n/\x/\y/\e in {s/0/0/s, v1/1.8/1.2/v_1,
        v2/1.8/-1.2/v_2, v3/4.2/1.2/v_3, v4/4.2/-1.2/v_4,
        t/6/0/t}
      \node[vert] (\n) at (\x, \y) {$\e$};
    \draw[dirigida] (s) -- node[peso] {13} (v2);
    \draw[dirigida] (v2) -- node[peso] {4} (v1);
    \draw[dirigida] (v3) -- node[peso] {9} (v2);
    \draw[dirigida] (v2) -- node[peso] {14} (v4);
    \draw[dirigida] (v4) -- node[peso] {7} (v3);
    \draw[dirigida] (v4) -- node[peso] {4} (t);
    \draw[dirigidaArbol] (s) -- node[peso] {16} (v1);
    \draw[dirigidaArbol] (v1) -- node[peso] {12} (v3);
    \draw[dirigidaArbol] (v3) -- node[peso] {20} (t);
  \end{tikzpicture}
  \hspace{0.8cm}
  \begin{tikzpicture}
    \foreach \n/\x/\y/\e in {s/0/0/s, v1/1.8/1.2/v_1,
        v2/1.8/-1.2/v_2, v3/4.2/1.2/v_3, v4/4.2/-1.2/v_4,
        t/6/0/t}
      \node[vert] (\n) at (\x, \y) {$\e$};
    \draw[dirigida] (s) to[bend left=16] node[peso, auto] {4} (v1);
    \draw[dirigida] (v1) to[bend left=16] node[peso, auto] {12} (s);
    \draw[dirigida] (v3) -- node[peso] {12} (v1);
    \draw[dirigida] (v2) -- node[peso] {4} (v1);
    \draw[dirigida] (v3) -- node[peso] {9} (v2);
    \draw[dirigida] (v4) -- node[peso] {7} (v3);
    \draw[dirigida] (v3) to[bend left=16] node[peso, auto] {8} (t);
    \draw[dirigida] (t) to[bend left=16] node[peso, auto] {12} (v3);
    \draw[dirigidaArbol] (s) -- node[peso] {13} (v2);
    \draw[dirigidaArbol] (v2) -- node[peso] {14} (v4);
    \draw[dirigidaArbol] (v4) -- node[peso] {4} (t);
  \end{tikzpicture}
  \end{center}
  y la de la iteración 3:
  \begin{center}
  \begin{tikzpicture}
    \foreach \n/\x/\y/\e in {s/0/0/s, v1/1.8/1.2/v_1,
        v2/1.8/-1.2/v_2, v3/4.2/1.2/v_3, v4/4.2/-1.2/v_4,
        t/6/0/t}
      \node[vert] (\n) at (\x, \y) {$\e$};
    \draw[dirigida] (s) to[bend left=16] node[peso, auto] {4} (v1);
    \draw[dirigida] (v1) to[bend left=16] node[peso, auto] {12} (s);
    \draw[dirigida] (v2) to[bend left=16] node[peso, auto] {4} (s);
    \draw[dirigida] (v3) -- node[peso] {12} (v1);
    \draw[dirigida] (v2) -- node[peso] {4} (v1);
    \draw[dirigida] (v3) -- node[peso] {9} (v2);
    \draw[dirigida] (v4) to[bend left=16] node[peso, auto] {4} (v2);
    \draw[dirigida] (t) to[bend left=16] node[peso, auto] {12} (v3);
    \draw[dirigida] (t) -- node[peso] {4} (v4);
    \draw[dirigidaArbol] (s) to[bend left=16] node[peso, auto] {9} (v2);
    \draw[dirigidaArbol] (v2) to[bend left=16] node[peso, auto] {10} (v4);
    \draw[dirigidaArbol] (v4) -- node[peso] {7} (v3);
    \draw[dirigidaArbol] (v3) to[bend left=16] node[peso, auto] {8} (t);
  \end{tikzpicture}
  \end{center}
  La red residual final es la misma del renglón (f) del
  ejemplo~\ref{ej:flujo-ff266}, y el flujo final es idéntico al de
  (e). Son tres aumentos en lugar de cinco, y ninguno usa una arista
  de regreso. Las distancias ilustran el lema que sigue:
  $\delta_f(s, v_3)$ pasa de $2$ a $3$ y luego a $\infty$,
  $\delta_f(s, t)$ va de $3$ a $3$ a $4$, y ninguna baja; las
  longitudes de los caminos, $3, 3, 4$, tampoco. Cada aumento tiene
  exactamente una arista crítica, y ninguna se repite. Hay un empate:
  $s, v_1, v_3, t$ y $s, v_2, v_4, t$ miden $3$, y si el BFS revisara
  $v_2$ antes que $v_1$ saldrían en orden inverso.
\end{ejemplo}

\subsubsection*{Por qué Edmonds-Karp es polinomial}

Todo el análisis descansa en un hecho sobre las distancias
residuales: con caminos más cortos, ningún vértice se acerca a $s$.

\begin{lema}[Distancias residuales; Cormen,
  lema 26.7]\label{lem:flujo-distancias}
  Durante Edmonds-Karp, para todo $v \in V$, la distancia residual
  $\delta_f(s, v)$ nunca baja de un aumento al siguiente.
\end{lema}

\begin{proof}
  Supóngase que sí baja alguna, y tómese el primer aumento que baja
  una distancia: $f$ es el flujo antes de él y $f'$ el de después.
  Entre los vértices cuya distancia bajó, sea $v$ el de
  $\delta_{f'}(s, v)$ más chica. Así
  $\delta_{f'}(s, v) < \delta_f(s, v)$, y $v \ne s$ porque
  $\delta(s, s) = 0$ siempre. Sea $u$ el vértice anterior a $v$ en un
  camino más corto de $s$ a $v$ en $G_{f'}$: $(u, v) \in E_{f'}$ y
  $\delta_{f'}(s, u) = \delta_{f'}(s, v) - 1$. Como $u$ quedó más
  cerca que $v$, por la elección de $v$ su distancia no bajó:
  \[
    \delta_{f'}(s, u) \ge \delta_f(s, u) .
  \]

  Si $(u, v)$ ya estaba en $E_f$, la desigualdad del triángulo para
  distancias en aristas (propiedad~\ref{prop:triangulo}, con todos los
  pesos iguales a $1$) daría
  \[
    \delta_f(s, v) \le \delta_f(s, u) + 1 \le \delta_{f'}(s, u) + 1
    = \delta_{f'}(s, v) ,
  \]
  contra la hipótesis de que la distancia de $v$ bajó. Entonces
  $(u, v) \notin E_f$ y $(u, v) \in E_{f'}$: la arista apareció con el
  aumento. La capacidad residual de $(u, v)$ solo sube si el flujo
  entre $u$ y $v$ cambia a favor de $u$, es decir, si el camino
  aumentante usó la arista residual $(v, u)$. Ese camino es un camino
  más corto en $G_f$, así que $\delta_f(s, u) = \delta_f(s, v) + 1$,
  y
  \[
    \delta_f(s, v) = \delta_f(s, u) - 1 \le \delta_{f'}(s, u) - 1
    = \delta_{f'}(s, v) - 2 ,
  \]
  otra contradicción con $\delta_{f'}(s, v) < \delta_f(s, v)$.
\end{proof}

Cormen enuncia el lema para $v \in V - \set{s, t}$, pero la prueba no
usa esa restricción; en $t$ la distancia tampoco baja, como se vio en
el ejemplo.

\begin{ojo}
  El lema dice ``nunca baja'', no ``siempre sube''. En el
  ejemplo~\ref{ej:flujo-ek}, $\delta_f(s, v_1)$ y $\delta_f(s, v_2)$
  se quedan en $1$ todo el tiempo. Y la garantía es del BFS: con DFS
  los caminos no son los más cortos y el lema deja de valer. Un
  Dijkstra que use las capacidades como pesos tampoco es
  Edmonds-Karp: a lo largo de un camino las capacidades no se suman,
  se toma su mínimo.
\end{ojo}

\begin{definicion}[Arista crítica]
  Una arista residual $(u, v)$ de un camino aumentante $p$ es
  \textbf{crítica} si $c_f(u, v) = c_f(p)$: es un cuello de botella.
\end{definicion}

Todo camino aumentante tiene al menos una arista crítica, y después
del aumento esa arista sale de $G_f$: si era una arista de $E$, quedó
saturada; si era la inversa de una, el flujo de esa arista volvió a
$0$.

\begin{teorema}[Cormen, teorema 26.8]\label{teo:flujo-ek-aumentos}
  Edmonds-Karp hace a lo más $n m$ aumentos.
\end{teorema}

\begin{proof}
  Sea $(u, v)$ crítica cuando el flujo es $f$. Como el camino es uno
  de los más cortos, su tramo hasta $v$ también lo es, y
  $\delta_f(s, v) = \delta_f(s, u) + 1$. Después del aumento $(u, v)$
  sale de $G_f$, y para volver a ser crítica primero tiene que
  reaparecer. Eso solo pasa si un aumento posterior, con flujo $f'$,
  empuja por $(v, u)$; ese camino también es de los más cortos, así
  que $\delta_{f'}(s, u) = \delta_{f'}(s, v) + 1$. Por el
  lema~\ref{lem:flujo-distancias}, $\delta_{f'}(s, v) \ge
  \delta_f(s, v)$, y entonces
  \[
    \delta_{f'}(s, u) = \delta_{f'}(s, v) + 1 \ge \delta_f(s, v) + 1
    = \delta_f(s, u) + 2 .
  \]
  Como las distancias nunca bajan, entre dos veces seguidas que
  $(u, v)$ es crítica la distancia de $s$ a $u$ sube al menos $2$.
  Además, cada vez que es crítica, $u$ está en un camino aumentante
  antes que $t$, y el tramo de $s$ a $u$ de ese camino no repite
  vértices ni pasa por $t$: tiene a lo más $n - 1$ vértices, y
  $\delta_f(s, u) \le n - 2$. Empezando en al menos $0$ y subiendo de
  $2$ en $2$ sin pasar de $n - 2$, $(u, v)$ es crítica a lo más
  $(n - 2)/2 + 1 = n/2$ veces. Las aristas residuales posibles son las
  de $E$ y sus inversas, a lo más $2m$, así que hay a lo más
  $2m \cdot n/2 = nm$ eventos críticos. Cada aumento tiene al menos
  uno.
\end{proof}

Cormen cuenta ``$\Ord{E}$ pares de vértices'' sin precisar; la cuenta
exacta, $2m$ aristas residuales posibles con $n/2$ eventos cada una,
da la cota $nm$.

\textbf{Costo.} Cada iteración hace un BFS sobre $G'$, de
$\Ord{n + 2m} = \Ord{m}$ con listas, por la misma cuenta del BFS de la
libreta, y dos pasadas por el camino, de $\Ord{n}$. Son a lo más $nm$
iteraciones, más un último BFS que no encuentra camino.

\begin{resultado}
  Edmonds-Karp hace $\Ord{nm}$ aumentos y cuesta
  \[
    \Ord{nm} \cdot \Ord{m} + \Ord{m} = \Ord{n m^2} ,
  \]
  sin importar las capacidades.
\end{resultado}

La cota no usa que las capacidades sean enteras: ni el lema ni el
teorema lo suponen, así que Edmonds-Karp termina con capacidades
reales cualesquiera. Y como sigue siendo una implementación de
Ford-Fulkerson, con capacidades enteras también cumple
$\Ord{m \abs{f^*}}$; vale la menor de las dos cotas. En la red de
Cormen las cotas de aumentos son $nm = 54$ y $\abs{f^*} = 23$, y bastan
$3$. En el caso patológico son $nm = 20$ contra $2 \cdot 10^6$, y
bastan $2$.

Con matriz de adyacencia cada BFS cuesta $\Tht{n^2}$ y el total sube a
$\Ord{n^3 m}$; la tabla de matriz contra listas de la clase del 22 de
octubre aplica tal cual. En un grafo denso, $m = \Tht{n^2}$ y
$\Ord{n m^2} = \Ord{n^5}$. Eso motiva los algoritmos de empujar y
reetiquetar de la sección~\ref{sec:push-relabel}, que no buscan
caminos aumentantes y llegan a $\Ord{n^2 m}$ y a $\Ord{n^3}$.

% =====================================================================

\subsection*{Errores comunes}

\begin{ojo}
  \begin{itemize}
    \item \textbf{Antiparalelas en la entrada.} La definición de $c_f$
          supone que no hay. Un grafo no dirigido, donde cada arista
          se modela con dos dirigidas, las tiene siempre. Hay dos
          salidas: partir cada par con un vértice $v'$, o usar
          $c_f(u, v) = c(u, v) - f(u, v) + f(v, u)$. Una
          implementación con matriz de capacidades residuales
          $r[u][v]$, que al aumentar resta en $r[u][v]$ y suma en
          $r[v][u]$, hace lo segundo sin pensarlo.
    \item \textbf{El flujo máximo no es único; su valor sí.} Aumentar
          el flujo de valor $19$ del ejemplo~\ref{ej:flujo-red} por
          $s, v_2, v_3, t$ da un flujo de valor $23$ con
          $f(v_2, v_1) = 1$; el de la figura 26.6 vale $23$ con
          $f(v_2, v_1) = 0$. También puede haber varios cortes
          mínimos, y la búsqueda final entrega uno. En la red de
          Cormen el corte mínimo es único.
    \item \textbf{La red es dirigida.} El flujo máximo de $s$ a $t$
          no es el de $t$ a $s$. En la red de Cormen no sale ninguna
          arista de $t$, y el flujo de $t$ a $s$ vale $0$.
    \item \textbf{Vértices fuera de todo camino de $s$ a $t$.} La
          definición los excluye. Si aparecen, hay un flujo máximo que
          no pasa por ellos (Cormen, ejercicio 26.1-3), y se pueden
          borrar antes de empezar.
  \end{itemize}
\end{ojo}

% =====================================================================

\subsection*{Ejercicios}

\begin{ejercicio}[Cormen, ejercicio 26.2-2]
  Con el flujo de valor $19$ del ejemplo~\ref{ej:flujo-red}, calcula
  el flujo neto y la capacidad del corte
  $(\set{s, v_2, v_4}, \set{v_1, v_3, t})$.

  \textit{Solución.} De $S$ a $T$ cruzan $(s, v_1)$, $(v_2, v_1)$,
  $(v_4, v_3)$ y $(v_4, t)$; de $T$ a $S$, $(v_3, v_2)$. Entonces
  $f(S, T) = 11 + 1 + 7 + 4 - 4 = 19$, como exige el
  lema~\ref{lem:flujo-neto-corte}, y
  $c(S, T) = 16 + 4 + 7 + 4 = 31$.
\end{ejercicio}

\begin{ejercicio}
  Corre el voraz sin aristas de regreso del
  ejemplo~\ref{ej:flujo-voraz}, pero empezando por
  $s, v_2, v_1, v_3, t$, el camino (b) de la figura 26.6, y después
  siempre con BFS entre las aristas con holgura. ¿Dónde se atora?
  Encuentra un camino aumentante en la red residual de ese flujo.

  \textit{Solución.} Toma $s, v_2, v_1, v_3, t$ con $4$,
  $s, v_1, v_3, t$ con $8$, $s, v_2, v_4, t$ con $4$ y
  $s, v_2, v_4, v_3, t$ con $5$, y se atora en $21$: $(s, v_2)$ y
  $(v_1, v_3)$ quedaron saturadas, y desde $s$ solo se llega a $v_1$.
  En la red residual el único camino aumentante es
  $s, v_1, v_2, v_4, v_3, t$, con
  $c_f(p) = \min(8, 4, 5, 2, 3) = 2$. Su segunda arista es la inversa
  de $(v_2, v_1)$: cancela $2$ de las $4$ unidades del primer camino,
  y el valor llega a $23$, la capacidad del corte mínimo.
\end{ejercicio}

\begin{ejercicio}
  Encuentra el flujo máximo de la red con superfuente y supersumidero
  del apartado de varias fuentes, y un corte que lo certifique.

  \textit{Solución.} El flujo máximo vale $10$. Un corte que lo
  certifica es $S = \set{s, s_1, s_2, x}$: de $S$ a $T$ solo cruzan
  $(x, t_1)$ y $(s_2, y)$, con $6 + 4 = 10$. Un flujo de ese valor:
  $5$ por $s_1, x, t_1$, $1$ por $s_2, x, t_1$, $2$ por $s_2, y, t_1$
  y $2$ por $s_2, y, t_2$. Las aristas de capacidad $\infty$ nunca
  cruzan un corte mínimo, porque su capacidad lo volvería infinito.
\end{ejercicio}

\begin{ejercicio}[Conectividad de aristas; Cormen, ejercicio 26.2-11]
  La conectividad de aristas de un grafo no dirigido conexo es el
  mínimo número de aristas que hay que quitar para desconectarlo.
  Muestra cómo calcularla con a lo más $n - 1$ problemas de flujo
  máximo, cada uno sobre una red de $n$ vértices y $\Ord{m}$ aristas.

  \textit{Solución.} Se fija un vértice $u$ y, para cada otro vértice
  $v$, se calcula el flujo máximo de $u$ a $v$ en la red dirigida que
  pone las dos direcciones de cada arista con capacidad $1$. La
  respuesta es el mínimo de esos $n - 1$ valores: en cualquier
  conjunto mínimo de aristas que desconecta el grafo, $u$ queda de un
  lado y algún $v$ del otro, y ese flujo máximo es igual al corte
  mínimo entre $u$ y $v$. Los pares antiparalelos se manejan con la
  capacidad residual combinada de los errores comunes; partirlos con
  vértices nuevos subiría el número de vértices a $\Ord{n + m}$.
\end{ejercicio}

% =====================================================================

\subsection*{Resumen}

\begin{center}
\small
\begin{tabular}{lll}
  \toprule
  & \textbf{Ford-Fulkerson} & \textbf{Edmonds-Karp} \\
  \midrule
  Camino aumentante & cualquiera & el de menos aristas, con BFS \\
  Número de aumentos & $\le \abs{f^*}$ con capacidades enteras
    & $\le nm$ \\
  Costo & $\Ord{m \abs{f^*}}$ con capacidades enteras
    & $\Ord{n m^2}$ \\
  Capacidades irracionales & puede no terminar & termina \\
  Red de Cormen & $5$ aumentos con los caminos de Cormen
    & $3$ aumentos \\
  Caso patológico & hasta $2 \cdot 10^6$ aumentos & $2$ aumentos \\
  \bottomrule
\end{tabular}
\end{center}

Lo que se lleva de la sección: el flujo se mejora por caminos
aumentantes en la red residual, cuyas aristas de regreso permiten
deshacer decisiones; cuando ya no hay camino, el
teorema~\ref{teo:flujo-corte} entrega a la vez la optimalidad y un
corte que la certifica; y la regla para escoger el camino es lo que
separa un algoritmo pseudopolinomial de uno polinomial.

\newpage

% =====================================================================
%  Grafos - Emparejamiento bipartito
% =====================================================================

\clase{Emparejamiento bipartito máximo}{Complemento: Cormen, cap.\ 26.3}
\label{sec:emparejamiento}

\begin{quick_sum}
  Un emparejamiento es un conjunto de aristas que no comparten
  extremos: cada máquina hace a lo más una tarea, cada alumno toma a
  lo más un proyecto. En un grafo bipartito, el emparejamiento más
  grande se encuentra con flujo máximo: se agregan una fuente y un
  sumidero, todas las capacidades valen $1$, y los flujos enteros de
  esa red resultan ser exactamente los emparejamientos. Como
  Ford-Fulkerson solo produce flujos enteros, el problema queda
  resuelto en $\Ord{nm}$. Leído sobre el grafo original, cada camino
  aumentante es una cadena de reasignaciones que termina en un
  vértice libre, y el teorema de Berge dice que un emparejamiento es
  máximo justo cuando ya no queda ninguna. Después vienen el teorema
  de Hall, que dice cuándo no hay emparejamiento perfecto y entrega
  un testigo que se revisa a mano, el de König, la idea de
  Hopcroft-Karp, que baja el costo a $\Ord{\sqrt{n}\, m}$, y la razón
  por la que en grafos no bipartitos todo esto se queda corto.
\end{quick_sum}

\subsection*{El problema}

Cormen lo plantea con máquinas y tareas. Hay un conjunto $L$ de
máquinas y un conjunto $R$ de tareas que se ejecutan al mismo
tiempo; la arista $\set{u, v}$ dice que la máquina $u$ sabe hacer la
tarea $v$. Cada máquina hace a lo más una tarea y cada tarea la hace
a lo más una máquina. La pregunta es a cuántas máquinas se les puede
dar trabajo a la vez. Con alumnos y proyectos es lo mismo: cada
alumno acepta algunos proyectos, cada proyecto admite un alumno, y se
quiere acomodar a tantos alumnos como se pueda.

El grafo es bipartito: $V = L \cup R$ con $L \cap R = \emptyset$, y
toda arista une un vértice de $L$ con uno de $R$. Aquí los lados se
llaman $L$ (izquierda) y $R$ (derecha), como en Cormen; son los $V_1$
y $V_2$ de la sección~\ref{sec:hamilton-arboles}.

\begin{definicion}[Emparejamiento]
  \label{def:emparejamiento}
  En un grafo no dirigido $G = (V, E)$, un \textbf{emparejamiento}
  (\textit{matching}) es un conjunto $M \subseteq E$ tal que cada
  vértice es incidente en a lo más una arista de $M$. Un vértice
  incidente en una arista de $M$ está \textbf{emparejado}; si no,
  está \textbf{libre}. Si $\set{u, v} \in M$, $v$ es la
  \textbf{pareja} de $u$ y $u$ la de $v$.
\end{definicion}

\begin{definicion}[Máximo, maximal y perfecto]
  \label{def:emparejamiento-maximo}
  Un emparejamiento $M$ es
  \begin{itemize}
    \item \textbf{máximo} (\textit{maximum}) si
      $\abs{M} \ge \abs{M'}$ para todo emparejamiento $M'$ de $G$;
    \item \textbf{maximal} (\textit{maximal}) si no está contenido
      propiamente en otro emparejamiento, es decir, si ninguna arista
      de $E - M$ se le puede agregar;
    \item \textbf{perfecto} (\textit{perfect}) si todo vértice de $G$
      está emparejado.
  \end{itemize}
\end{definicion}

Todo máximo es maximal, pero no al revés, y esa diferencia es todo
el problema. Un perfecto es máximo, y en un bipartito exige
$\abs{L} = \abs{R}$. El \textbf{problema del emparejamiento
bipartito máximo} pide, dado $G$ bipartito, un emparejamiento máximo.
Cormen supone además que todo vértice tiene al menos una arista; de
ahí sale $m \ge n/2$, que se usa en las cotas de tiempo. Un vértice
aislado no cambia la respuesta y se quita en $\Tht{n}$ antes de
empezar.

\begin{ejemplo}[Figura 26.8 de Cormen: maximal contra máximo]
  \label{ej:emparejamiento-cormen}
  Los vértices se numeran como en el ejercicio 26.3-1 de Cormen:
  $L = \set{1, 2, 3, 4, 5}$ y $R = \set{6, 7, 8, 9}$, de arriba
  abajo, con
  \[
    E = \set{\set{1,6}, \set{2,6}, \set{2,8}, \set{3,7}, \set{3,8},
             \set{3,9}, \set{4,8}, \set{5,8}} ,
  \]
  $n = 9$ y $m = 8$. En azul van las aristas del emparejamiento y en
  verde los vértices emparejados.
  \begin{center}
  \begin{tikzpicture}
    \foreach \pre/\x/\lst/\ok in
      {a/0/{1/6, 3/8}/{1,3,6,8},
       b/4.2/{2/6, 3/7, 5/8}/{2,3,5,6,7,8}} {
      \begin{scope}[xshift=\x cm]
        \foreach \i/\y in {1/0, 2/-0.9, 3/-1.8, 4/-2.7, 5/-3.6}
          \node[vert] (\pre\i) at (0, \y) {\i};
        \foreach \i/\y in {6/-0.45, 7/-1.35, 8/-2.25, 9/-3.15}
          \node[vert] (\pre\i) at (1.8, \y) {\i};
        \foreach \i in \ok
          \node[vertOk] at (\pre\i) {\i};
        \foreach \a/\b in {1/6, 2/6, 2/8, 3/7, 3/8, 3/9, 4/8, 5/8}
          \draw[arista] (\pre\a) -- (\pre\b);
        \foreach \a/\b in \lst
          \draw[aristaArbol] (\pre\a) -- (\pre\b);
        \node[nota] at (0, -4.2) {$L$};
        \node[nota] at (1.8, -4.2) {$R$};
      \end{scope}
    }
    \node[font=\small] at (0.9, -4.75) {(a) maximal, $\abs{M_a} = 2$};
    \node[font=\small] at (5.1, -4.75) {(b) máximo, $\abs{M_b} = 3$};
  \end{tikzpicture}
  \end{center}
  En (a), $M_a = \set{\set{1,6}, \set{3,8}}$ es maximal: cada una de
  las otras seis aristas toca a $1$, $3$, $6$ u $8$. En (b),
  $M_b = \set{\set{2,6}, \set{3,7}, \set{5,8}}$ es máximo, y el
  argumento es corto: $7$ y $9$ solo tienen de vecino a $3$, así que
  a lo más uno de los dos queda emparejado, y de $R$ se emparejan a
  lo más tres vértices. Cormen no usa la palabra \textit{maximal}; el
  pie de su figura dice solo que (a) tiene cardinalidad $2$.
\end{ejemplo}

La forma obvia de construir un emparejamiento es voraz: recorrer las
aristas y tomar cada una que no choque con las ya tomadas. Lo que
sale es maximal por construcción, pero su tamaño depende del orden.
En la figura 26.8, con las aristas en el orden en que se listaron
arriba, el voraz toma $\set{1,6}$, $\set{2,8}$ y $\set{3,7}$, que es
máximo; si empieza por $\set{1,6}$ y $\set{3,8}$, se atora en $M_a$.
Lo único que garantiza es la mitad.

\begin{proposicion}
  \label{prop:emparejamiento-maximal}
  Si $M$ es un emparejamiento maximal y $M^*$ uno máximo, entonces
  $\abs{M^*} \le 2\abs{M}$.
\end{proposicion}

\begin{proof}
  Cada arista de $M^*$ comparte al menos un extremo con alguna arista
  de $M$; si no, se le podría agregar a $M$, que no sería maximal. A
  cada arista de $M^*$ se le asigna uno de esos extremos compartidos.
  Los vértices asignados son distintos, porque dos aristas de $M^*$
  no comparten extremos, y todos son extremos de aristas de $M$, que
  en total tienen $2\abs{M}$. Entonces $\abs{M^*} \le 2\abs{M}$.
\end{proof}

La cota se alcanza: en el camino $a$--$b$--$c$--$d$, el voraz que
empieza por la arista $\set{b, c}$ ya no puede agregar nada y se
queda con $1$, mientras que $\set{\set{a,b}, \set{c,d}}$ mide $2$.
Cormen no incluye esta proposición en 26.3; está aquí porque dice
exactamente cuánto puede perder el voraz.

\begin{ojo}
  Tres confusiones que conviene separar desde el principio. Maximal
  no es máximo: $M_a$ es maximal y mide $2$, y el máximo mide $3$. El
  máximo no es único: la figura 26.8 tiene $10$ emparejamientos de
  tamaño $3$, y distintas reglas de desempate llegan a distintos
  máximos, todos del mismo tamaño. Y máximo no es perfecto: ahí
  $\abs{L} = 5$ y $\abs{R} = 4$, así que algún vértice de $L$ se
  queda libre siempre.
\end{ojo}

Para pasar de $M_a$ a un máximo hay que deshacer algo: sacar la
arista $\set{3,8}$ y meter dos nuevas. La pregunta es cómo encontrar
esos reacomodos sin probar todas las combinaciones. Cormen responde
con flujo máximo; leída sobre el grafo original, la misma respuesta
son los caminos aumentantes.

% =====================================================================

\subsection*{Reducción a flujo máximo}

\begin{definicion}[Red de flujo correspondiente]
  \label{def:red-emparejamiento}
  Sea $G = (V, E)$ bipartito con $V = L \cup R$. Su \textbf{red de
  flujo correspondiente} es $G' = (V', E')$ con
  $V' = V \cup \set{s, t}$ y
  \[
    E' = \set{(s, u) : u \in L}
         \cup \set{(u, v) : u \in L,\ v \in R,\ \set{u, v} \in E}
         \cup \set{(v, t) : v \in R} ,
  \]
  con capacidad $c(u, v) = 1$ en todas las aristas de $E'$.
\end{definicion}

Cada arista de $G$ se orienta de $L$ a $R$, la fuente $s$ apunta a
todo $L$ y todo $R$ apunta al sumidero $t$. La red tiene
$\abs{E'} = m + n$ aristas, y con $m \ge n/2$ queda
$m \le \abs{E'} \le 3m$, o sea $\abs{E'} = \Tht{m}$. Todo el modelado
está en las capacidades de $(s, u)$ y $(v, t)$, que son capacidades
de vértice disfrazadas. A $u \in L$ solo le entra la arista $(s, u)$,
así que le entra a lo más una unidad, y por conservación esa unidad
sale por una sola arista hacia $R$. Del lado de $R$ pasa lo mismo con
$(v, t)$.

Para la figura 26.8 la red tiene $11$ vértices y
$\abs{E'} = 5 + 8 + 4 = 17$ aristas; es el panel (c) de esa figura.
En azul va el flujo que corresponde a $M_b$: una unidad por cada uno
de los caminos $s, 2, 6, t$, \ $s, 3, 7, t$ \ y \ $s, 5, 8, t$, de
valor $\abs{f} = 3$.

\begin{center}
\begin{tikzpicture}
  \node[vert] (s) at (-2, -1.8) {$s$};
  \node[vert] (t) at (4.2, -1.8) {$t$};
  \foreach \i/\y in {1/0, 2/-0.9, 3/-1.8, 4/-2.7, 5/-3.6}
    \node[vert] (\i) at (0, \y) {\i};
  \foreach \i/\y in {6/-0.45, 7/-1.35, 8/-2.25, 9/-3.15}
    \node[vert] (\i) at (2.2, \y) {\i};
  \foreach \a/\b in {s/1, s/4, 1/6, 2/8, 3/8, 3/9, 4/8, 9/t}
    \draw[dirigida] (\a) -- (\b);
  \foreach \a/\b in {s/2, s/3, s/5, 2/6, 3/7, 5/8, 6/t, 7/t, 8/t}
    \draw[dirigidaArbol] (\a) -- (\b);
  \node[nota] at (0, -4.2) {$L$};
  \node[nota] at (2.2, -4.2) {$R$};
  \node[nota, align=center] at (1.1, 0.75)
    {todas las capacidades valen $1$; en azul, $f = 1$};
\end{tikzpicture}
\end{center}

\begin{algorithm}[H]
\caption{\textsc{Emparejamiento-Bipartito}$(G, L, R)$}
\label{alg:emparejamiento-flujo}
\begin{algorithmic}[1]
\Procedure{Emparejamiento-Bipartito}{$G, L, R$}
  \State $V' \gets V \cup \set{s, t}$, \quad $E' \gets \emptyset$
  \ForAll{$u \in L$}
    \State $E' \gets E' \cup \set{(s, u)}$
  \EndFor
  \ForAll{$\set{u, v} \in E$ con $u \in L$, $v \in R$}
    \State $E' \gets E' \cup \set{(u, v)}$ \Comment{solo de $L$ a $R$}
  \EndFor
  \ForAll{$v \in R$}
    \State $E' \gets E' \cup \set{(v, t)}$
  \EndFor
  \State $c(e) \gets 1$ para toda $e \in E'$
  \State $f \gets \textsc{Ford-Fulkerson}(G', s, t)$
  \State \Return $M = \set{\set{u, v} : u \in L,\ v \in R,\
         f(u, v) = 1}$
\EndProcedure
\end{algorithmic}
\end{algorithm}

La línea 13 es el algoritmo~\ref{alg:ford-fulkerson} de la
sección~\ref{sec:flujo}, con cualquier forma de buscar el camino; con
BFS es Edmonds-Karp (algoritmo~\ref{alg:edmonds-karp}). La línea 14
lee el emparejamiento en el flujo, y que eso funcione es justo lo que
hay que demostrar: que los flujos enteros son emparejamientos y que
el flujo que devuelve Ford-Fulkerson es entero.

\begin{lema}[Emparejamientos y flujos enteros]
  \label{lem:emparejamiento-flujo}
  Si $M$ es un emparejamiento de $G$, existe un flujo entero $f$ en
  $G'$ con $\abs{f} = \abs{M}$. Recíprocamente, si $f$ es un flujo
  entero en $G'$, entonces
  $M = \set{\set{u, v} : u \in L,\ v \in R,\ f(u, v) = 1}$ es un
  emparejamiento de $G$ con $\abs{M} = \abs{f}$.
\end{lema}

\begin{proof}
  \textit{Ida.} Por cada $\set{u, v} \in M$, con $u \in L$, se pone
  $f(s, u) = f(u, v) = f(v, t) = 1$, y $0$ en las demás aristas. Las
  capacidades se respetan: cada $u \in L$ está en a lo más una arista
  de $M$, así que $(s, u)$ recibe a lo más una unidad, y lo mismo
  $(v, t)$ del lado de $R$. Se conserva el flujo: a un $u \in L$
  emparejado le entra $1$ por $(s, u)$ y le sale $1$ por $(u, v)$, y a
  uno libre no le entra ni le sale nada; en $R$ igual. A $s$ no le
  entra ninguna arista, así que
  $\abs{f} = \sum_{u \in L} f(s, u)$, que cuenta los vértices
  emparejados de $L$: $\abs{f} = \abs{M}$.

  \textit{Vuelta.} Como $f$ es entero y $0 \le f \le 1$, cada arista
  lleva $0$ o $1$. A $u \in L$ solo le entra $(s, u)$, así que le
  entra $0$ o $1$; por conservación le sale lo mismo, y como cada
  arista de salida lleva $0$ o $1$, a lo más una arista $(u, v)$
  lleva $1$. Del lado de $R$, a $v$ solo le sale $(v, t)$, y el mismo
  argumento dice que a lo más una arista $(u, v)$ que le entra lleva
  $1$. Entonces ningún vértice está en dos aristas de $M$, y $M$ es
  un emparejamiento. Por último, sumando la conservación en cada
  $u \in L$,
  $\abs{f} = \sum_{u \in L} f(s, u)
           = \sum_{u \in L} \sum_{v \in R} f(u, v) = \abs{M}$.
\end{proof}

\begin{ojo}
  El flujo tiene que ser entero. En $K_{2,2}$, con $L = \set{a, b}$ y
  $R = \set{x, y}$, poner $1/2$ en las cuatro aristas de en medio y
  $1$ en las de $s$ y de $t$ es un flujo válido de valor $2$, máximo,
  y de él no se lee ningún emparejamiento: no hay arista con
  $f = 1$. Un algoritmo de flujo máximo cualquiera podría devolver
  ese flujo. Lo que salva a la reducción es el siguiente teorema, que
  vale para los métodos de caminos aumentantes.
\end{ojo}

\begin{teorema}[Integralidad]
  \label{teo:emparejamiento-integralidad}
  Si todas las capacidades son enteras, el flujo máximo $f$ que
  devuelve Ford-Fulkerson tiene $\abs{f}$ entero y $f(u, v)$ entero
  en cada arista.
\end{teorema}

\begin{proof}
  Por inducción sobre las iteraciones. Al principio $f = 0$. Si $f$
  es entero, todas las capacidades residuales son enteras, porque son
  de la forma $c(u, v) - f(u, v)$ o $f(v, u)$
  (definición~\ref{def:red-residual}). Entonces
  $c_f(p) = \min\set{c_f(u, v) : (u, v) \text{ está en } p}$ es un
  entero, y el aumento suma o resta ese entero en cada arista del
  camino: $f$ sigue siendo entero. Es el mismo argumento que acota el
  número de iteraciones de Ford-Fulkerson en la
  sección~\ref{sec:flujo}.
\end{proof}

\begin{corolario}
  \label{cor:emparejamiento-flujo}
  El tamaño de un emparejamiento máximo de $G$ es igual al valor de un
  flujo máximo en $G'$, y el emparejamiento que devuelve
  \textup{\textsc{Emparejamiento-Bipartito}} es máximo.
\end{corolario}

\begin{proof}
  Sea $f^*$ el flujo máximo que devuelve Ford-Fulkerson. Es entero
  por el teorema~\ref{teo:emparejamiento-integralidad}, así que el
  lema~\ref{lem:emparejamiento-flujo} lo convierte en un
  emparejamiento $M$ con $\abs{M} = \abs{f^*}$. Si $M'$ es cualquier
  otro emparejamiento, el mismo lema da un flujo $f'$ con
  $\abs{f'} = \abs{M'}$, y como $f^*$ es máximo,
  $\abs{M'} = \abs{f'} \le \abs{f^*} = \abs{M}$.
\end{proof}

\subsubsection*{Cuánto cuesta}

\begin{itemize}
  \item Construir $G'$ (líneas 2 a 12) cuesta $\Tht{n + m}$: una
    pasada por $L$, una por $E$ y una por $R$.
  \item Cada iteración de Ford-Fulkerson busca un camino de $s$ a $t$
    en $G'_f$ con BFS o DFS. La residual tiene $n + 2$ vértices y a
    lo más $2\abs{E'} = 2(m + n)$ aristas, así que la búsqueda cuesta
    $\Ord{n + m} = \Ord{m}$, usando $n \le 2m$. Actualizar el flujo a
    lo largo del camino es $\Ord{n}$ y queda absorbido.
  \item Cada aumento sube $\abs{f}$ en $c_f(p)$, que aquí vale
    exactamente $1$ porque todas las capacidades residuales son $0$ o
    $1$. El flujo máximo es
    $\abs{M^*} \le \min(\abs{L}, \abs{R}) \le n/2$, así que hay a lo
    más $n/2$ aumentos y una búsqueda final que falla.
\end{itemize}

En total, $\Ord{n} \cdot \Ord{m} = \Ord{nm}$. La cota no depende de
cómo se elija el camino: vale con BFS y con DFS. Es la cota
$\Ord{m\,\abs{f^*}}$ de Ford-Fulkerson con $\abs{f^*} \le n/2$; la de
Edmonds-Karp en general, $\Ord{nm^2}$, aquí es peor y no hace falta.
Cormen escribe $\min(L, R)$ para la cota del flujo; se lee
$\min(\abs{L}, \abs{R})$.

\begin{ojo}
  Las aristas de $G'$ van solo de $L$ a $R$. Si se dejaran en los dos
  sentidos, un camino de flujo podría cruzar de $L$ a $R$, regresar a
  $L$ por una arista al revés y volver a cruzar, y una unidad de flujo
  ya no correspondería a una arista del emparejamiento: la prueba del
  lema se cae en el paso que dice que a $u \in L$ solo le entra
  $(s, u)$. Además, las redes de flujo de la
  sección~\ref{sec:flujo} no admiten aristas antiparalelas, y con la
  orientación de $L$ a $R$ no aparece ninguna.
\end{ojo}

% =====================================================================

\subsection*{Caminos aumentantes en el grafo original}

Ford-Fulkerson trabaja sobre $G'$, pero lo que hace se entiende mejor
en $G$. La correspondencia que sigue no la escribe Cormen en 26.3 (la
deja implícita en el problema 26-6); es la clave para leer cualquier
traza del algoritmo y para escribir una versión que no construye
$G'$.

\begin{definicion}[Camino aumentante respecto a $M$]
  \label{def:camino-aumentante-emparejamiento}
  Sea $M$ un emparejamiento de $G$ bipartito. Un camino es
  \textbf{alternante} si sus aristas están, una sí y una no, en
  $E - M$ y en $M$. Un \textbf{camino aumentante} respecto a $M$
  (\textit{augmenting path}) es un camino alternante que empieza en
  un vértice libre de $L$ y termina en un vértice libre de $R$.
\end{definicion}

La primera y la última arista de un camino aumentante están fuera de
$M$, así que tiene $k$ aristas de $M$ y $k + 1$ fuera de $M$, y su
longitud $2k + 1$ es impar. La \textbf{diferencia simétrica} de dos
conjuntos es $A \oplus B = (A - B) \cup (B - A)$: lo que está en uno
solo de los dos.

\begin{proposicion}[La residual de un emparejamiento]
  \label{prop:emparejamiento-residual}
  Sea $M$ un emparejamiento de $G$ y $f$ el flujo entero que le
  corresponde en $G'$ por el lema~\ref{lem:emparejamiento-flujo}. La
  red residual $G'_f$ tiene exactamente estas aristas, con $u \in L$
  y $v \in R$:
  \begin{itemize}
    \item $(s, u)$ si $u$ está libre, y $(u, s)$ si está emparejado;
    \item $(u, v)$ si $\set{u, v} \in E - M$, y $(v, u)$ si
      $\set{u, v} \in M$;
    \item $(v, t)$ si $v$ está libre, y $(t, v)$ si está emparejado.
  \end{itemize}
  En consecuencia, los caminos simples de $s$ a $t$ en $G'_f$ son
  exactamente los $s, u_0, v_0, u_1, v_1, \dots, u_k, v_k, t$ tales
  que $P = u_0, v_0, \dots, u_k, v_k$ es un camino aumentante
  respecto a $M$, y aumentar el flujo por uno de ellos produce el
  flujo de $M \oplus P$.
\end{proposicion}

\begin{proof}
  En $G'$ todas las capacidades valen $1$ y $f$ solo toma los valores
  $0$ y $1$. Por la definición de red residual
  (definición~\ref{def:red-residual}), una arista $(x, y) \in E'$
  aparece en $G'_f$ hacia adelante si $f(x, y) = 0$ y aparece
  invertida, como $(y, x)$, si $f(x, y) = 1$. Como $f$ vale $1$
  exactamente en las $(s, u)$ y $(v, t)$ de vértices emparejados y en
  las $(u, v)$ con $\set{u, v} \in M$, salen las tres líneas de la
  lista.

  Sea ahora un camino simple de $s$ a $t$ en $G'_f$. No usa aristas
  que entran a $s$ ni que salen de $t$, porque no regresa a $s$ y
  termina en $t$. Su primera arista es entonces $(s, u_0)$ con $u_0$
  libre. Desde un vértice de $L$ las únicas salidas que quedan son
  aristas $(u, v)$ con $\set{u, v} \notin M$. Desde un vértice
  $v \in R$ hay una sola salida: hacia $t$ si $v$ está libre, y hacia
  su pareja si está emparejado. Así, el camino alterna una arista
  fuera de $M$, de $L$ a $R$, con una de $M$, de $R$ a $L$, hasta
  llegar a un vértice libre de $R$, desde donde sale a $t$. Como el
  camino es simple, los vértices no se repiten, y lo que queda entre
  $s$ y $t$ es un camino aumentante. Al revés, todo camino aumentante
  con $s$ al principio y $t$ al final es un camino en $G'_f$.

  Su capacidad residual es $1$. Aumentar pone en $1$ el flujo de
  $(s, u_0)$, de $(v_k, t)$ y de cada $(u_i, v_i)$, que estaban en
  $0$, y regresa a $0$ el de cada arista $(u_{i+1}, v_i)$ de $M$, que
  el camino recorre al revés. Las aristas de $M$ que están en $P$
  salen y las de $P$ fuera de $M$ entran: es el flujo de
  $M \oplus P$.
\end{proof}

La frase clave de la prueba es que un vértice emparejado de $R$ tiene
una sola salida en la residual, hacia su pareja. Por eso los caminos
alternan solos, sin que nadie se lo pida al algoritmo.

\begin{lema}
  \label{lem:emparejamiento-diferencia}
  Si $P$ es un camino aumentante respecto a $M$, entonces
  $M \oplus P$ es un emparejamiento y
  $\abs{M \oplus P} = \abs{M} + 1$.
\end{lema}

\begin{proof}
  $P$ tiene $k$ aristas en $M$ y $k + 1$ fuera; $M \oplus P$ quita las
  primeras y agrega las segundas, así que mide $\abs{M} + 1$. Falta
  ver que sigue siendo emparejamiento. Un vértice interior de $P$ toca
  dos aristas de $P$, una de $M$ y otra no, y ninguna otra arista de
  $M$: cambia una por la otra y se queda con exactamente una. Los
  extremos de $P$ estaban libres y ganan una sola arista, la primera
  o la última de $P$. Los vértices fuera de $P$ no cambian.
\end{proof}

\begin{ejemplo}[Del maximal al máximo con un camino aumentante]
  \label{ej:emparejamiento-aumento}
  Se vuelve a la figura 26.8. Respecto a $M_a$ los vértices libres
  son $2$, $4$ y $5$ en $L$, y $7$ y $9$ en $R$. Hay seis caminos
  aumentantes, todos de longitud $3$: de $2$, $4$ o $5$ a $8$, de $8$
  a su pareja $3$, y de $3$ a $7$ o a $9$. Todos pasan por la arista
  $\set{3,8}$ de $M_a$, la que estorba: ocupa a la vez al único
  vecino de $7$ y $9$ y al único vecino de $4$ y $5$. Con
  $P = 4$--$8$--$3$--$7$:
  \begin{center}
  \small
  \begin{tabular}{ccc}
    \toprule
    \textbf{Arista de $P$} & \textbf{¿Está en $M_a$?}
      & \textbf{En $M_a \oplus P$} \\
    \midrule
    $\set{4,8}$ & no & entra \\
    $\set{3,8}$ & sí & sale \\
    $\set{3,7}$ & no & entra \\
    \bottomrule
  \end{tabular}
  \end{center}
  y $M_a \oplus P = \set{\set{1,6}, \set{3,7}, \set{4,8}}$, de tamaño
  $3$. En la figura, el trazo naranja es $P$ dibujado sobre $M_a$.
  \begin{center}
  \begin{tikzpicture}[caminoAum/.style={draw=orange!60,
      line width=6pt, line cap=round, opacity=0.55}]
    \foreach \pre/\x/\lst/\ok/\cam in
      {a/0/{1/6, 3/8}/{1,3,6,8}/{4/8, 8/3, 3/7},
       b/4.2/{1/6, 3/7, 4/8}/{1,3,4,6,7,8}/{}} {
      \begin{scope}[xshift=\x cm]
        \foreach \i/\y in {1/0, 2/-0.9, 3/-1.8, 4/-2.7, 5/-3.6}
          \coordinate (\pre\i) at (0, \y);
        \foreach \i/\y in {6/-0.45, 7/-1.35, 8/-2.25, 9/-3.15}
          \coordinate (\pre\i) at (1.8, \y);
        \foreach \a/\b in \cam
          \draw[caminoAum] (\pre\a) -- (\pre\b);
        \foreach \i in {1,...,9}
          \node[vert] (\pre\i) at (\pre\i) {\i};
        \foreach \i in \ok
          \node[vertOk] at (\pre\i) {\i};
        \foreach \a/\b in {1/6, 2/6, 2/8, 3/7, 3/8, 3/9, 4/8, 5/8}
          \draw[arista] (\pre\a) -- (\pre\b);
        \foreach \a/\b in \lst
          \draw[aristaArbol] (\pre\a) -- (\pre\b);
      \end{scope}
    }
    \node[font=\small] at (0.9, -4.3) {$M_a$ y el camino $P$};
    \node[font=\small] at (5.1, -4.3) {$M_a \oplus P$};
    \draw[dirigida] (2.5, -1.8) -- node[nota, above] {$\oplus$}
      (3.6, -1.8);
  \end{tikzpicture}
  \end{center}
  Respecto a $M_b$ no hay ningún camino aumentante, como debe ser
  para un máximo.
\end{ejemplo}

\begin{teorema}[Berge]
  \label{teo:emparejamiento-berge}
  Un emparejamiento $M$ es máximo si y solo si no existe ningún
  camino aumentante respecto a $M$.
\end{teorema}

\begin{proof}
  Si hay un camino aumentante $P$, el
  lema~\ref{lem:emparejamiento-diferencia} da un emparejamiento
  $M \oplus P$ más grande, y $M$ no es máximo.

  Al revés, supóngase que $M$ no es máximo y sea $M^*$ un
  emparejamiento con $\abs{M^*} > \abs{M}$. En el grafo
  $H = (V, M \oplus M^*)$ cada vértice tiene grado a lo más $2$: una
  arista de $M$ y una de $M^*$, como mucho. Un grafo así se parte en
  caminos y ciclos, y en cada uno las aristas alternan entre
  $M - M^*$ y $M^* - M$, porque dos aristas del mismo emparejamiento
  no comparten vértice. Un ciclo alternante tiene tantas aristas de un
  emparejamiento como del otro. Como $M \oplus M^*$ tiene más aristas
  de $M^*$ que de $M$, algún camino de $H$ tiene más aristas de $M^*$
  que de $M$; alternando, empieza y termina con aristas de $M^* - M$.
  Sus dos extremos están libres en $M$. Si un extremo $x$ tuviera una
  arista $e \in M$, esa arista no estaría en $M^*$, porque la única
  arista de $M^*$ que toca a $x$ es la del camino, que no es de $M$;
  entonces $e$ estaría en $H$ y $x$ tendría grado $2$, y no sería
  extremo. Ese camino es aumentante respecto a $M$: tiene longitud
  impar, así que un extremo está en $L$ y el otro en $R$, y se lee
  desde el de $L$.
\end{proof}

El nombre viene de Berge (1957); Cormen lo deja como parte del
problema 26-6 sin nombrarlo. Con la
proposición~\ref{prop:emparejamiento-residual}, el teorema de Berge para
bipartitos es la equivalencia entre flujo máximo y ausencia de
caminos aumentantes del teorema~\ref{teo:flujo-corte}. La prueba de
arriba, en cambio, no usa flujo, y usa la bipartición solo para
nombrar los extremos del camino. Por eso sigue valiendo en cualquier
grafo, cosa que se retoma al final de la sección.

\subsubsection*{El algoritmo escrito sobre $G$}

Con la proposición~\ref{prop:emparejamiento-residual}, Ford-Fulkerson con
BFS sobre $G'$ se escribe sin construir $G'$. Es el BFS de la clase
del 22 de octubre con dos cambios: arranca con todos los vértices
libres de $L$ en la cola, y al llegar a un vértice emparejado de $R$
salta directo a su pareja. El arreglo $\text{pareja}[v]$ guarda la
pareja de cada vértice, o nulo si está libre.

\begin{algorithm}[H]
\caption{\textsc{Camino-Aumentante}$(G, M)$}
\label{alg:emparejamiento-camino}
\begin{algorithmic}[1]
\Procedure{Camino-Aumentante}{$G, M$}
  \ForAll{$v \in V$}
    \State $\text{visto}[v] \gets$ falso, \quad $\pi[v] \gets$ nulo
  \EndFor
  \State $Q \gets$ cola vacía
  \ForAll{$u \in L$ libre en $M$}
    \State $\text{visto}[u] \gets$ verdadero; \ \textsc{Encolar}$(Q, u)$
  \EndFor
  \While{$Q$ no está vacía}
    \State $u \gets \textsc{Desencolar}(Q)$ \Comment{siempre $u \in L$}
    \ForAll{$v$ adyacente a $u$ con $\set{u, v} \notin M$}
      \If{\textbf{no} $\text{visto}[v]$}
        \State $\text{visto}[v] \gets$ verdadero, \quad $\pi[v] \gets u$
        \If{$v$ está libre en $M$}
          \State \Return el camino de $v$ hacia atrás por $\pi$
        \EndIf
        \State $w \gets \text{pareja}[v]$
          \Comment{$\set{w, v} \in M$}
        \State $\text{visto}[w] \gets$ verdadero, \quad
               $\pi[w] \gets v$
        \State \textsc{Encolar}$(Q, w)$
      \EndIf
    \EndFor
  \EndWhile
  \State \Return nulo
\EndProcedure
\end{algorithmic}
\end{algorithm}

Las líneas que no son obvias en \textsc{Camino-Aumentante}:
\begin{itemize}
  \item Líneas 6 a 8: todos los libres de $L$ son raíces a la vez. En
    $G'$ son el primer nivel del BFS desde $s$, que aquí ya no
    aparece.
  \item Línea 11: desde $L$ solo se avanza por aristas fuera de $M$,
    que son las aristas $(u, v)$ de la residual.
  \item Líneas 14 a 16: el primer vértice libre de $R$ que se descubre
    cierra el camino; en $G'_f$ es la arista $(v, t)$.
  \item Líneas 17 a 19: desde un $v$ emparejado solo se puede seguir a
    su pareja. No hace falta preguntar si $w$ ya fue visto: $w$ está
    emparejado, así que no fue raíz, y solo se llega a él desde $v$,
    que se acaba de visitar por primera vez. Por eso en la cola solo
    hay vértices de $L$.
\end{itemize}

\begin{algorithm}[H]
\caption{\textsc{Emparejamiento-Máximo}$(G, L, R)$}
\label{alg:emparejamiento-maximo}
\begin{algorithmic}[1]
\Procedure{Emparejamiento-Máximo}{$G, L, R$}
  \State $M \gets \emptyset$
  \Repeat
    \State $P \gets \textsc{Camino-Aumentante}(G, M)$
    \If{$P \neq$ nulo}
      \State $M \gets M \oplus P$ \Comment{crece en uno}
    \EndIf
  \Until{$P =$ nulo}
  \State \Return $M$
\EndProcedure
\end{algorithmic}
\end{algorithm}

\textsc{Emparejamiento-Máximo} termina con un máximo por el teorema
de Berge: \textsc{Camino-Aumentante} devuelve nulo solo cuando $t$ no
es alcanzable en $G'_f$, es decir, cuando no hay camino aumentante.
Una llamada visita cada vértice una vez y revisa cada arista una vez
desde su extremo en $L$: $\Tht{n + m}$. Hay a lo más
$\abs{M^*} + 1 \le n/2 + 1$ llamadas, y $M \oplus P$ cuesta $\Ord{n}$
actualizando $\text{pareja}[\,]$. En total
$\Ord{n(n + m)} = \Ord{nm}$ sin vértices aislados, igual que la
reducción, y no es casualidad: es el mismo algoritmo. Con los vecinos
en el mismo orden, elige exactamente los mismos caminos que
Edmonds-Karp sobre $G'$.

\begin{ojo}
  Cormen advierte que un camino aumentante respecto a $M$ está
  ``relacionado con, pero es distinto de'' un camino aumentante en
  una red de flujo. El de $G$ va de un libre de $L$ a un libre de $R$
  y alterna; el de $G'_f$ va de $s$ a $t$. Se corresponden uno a uno
  por la proposición~\ref{prop:emparejamiento-residual}, pero no son el
  mismo objeto. Y no cualquier camino alternante sirve: uno que
  termina en un vértice emparejado no es aumentante, y aplicarle
  $\oplus$ no hace crecer a $M$.
\end{ojo}

% =====================================================================

\subsection*{Ejemplo: alumnos y proyectos}

\begin{ejemplo}[El emparejamiento crece de $1$ a $4$]
  \label{ej:emparejamiento-alumnos}
  Cuatro alumnos, Ana, Beto, Caro y Dani, forman
  $L = \set{A, B, C, D}$, y cuatro proyectos forman
  $R = \set{1, 2, 3, 4}$. Cada alumno acepta:
  \begin{center}
  \small
  \begin{tabular}{ll}
    \toprule
    \textbf{Alumno} & \textbf{Proyectos que acepta} \\
    \midrule
    A (Ana) & 1, 2 \\
    B (Beto) & 3, 4 \\
    C (Caro) & 1, 3 \\
    D (Dani) & 1 \\
    \bottomrule
  \end{tabular}
  \end{center}
  Se abrevia la arista $\set{A, 1}$ como $A1$. Entonces
  $E = \set{A1, A2, B3, B4, C1, C3, D1}$, $n = 8$, $m = 7$ y
  $\abs{E'} = 15$. El proyecto 1 es el más pedido y Dani no tiene
  alternativa. Se corre \textsc{Emparejamiento-Máximo}, que es
  Edmonds-Karp sobre $G'$, con los vecinos en orden alfabético y
  numérico.
  \begin{center}
  \small
  \begin{tabular}{cllll}
    \toprule
    \textbf{It.} & \textbf{Camino en $G'_f$} & \textbf{En $G$}
      & \textbf{$M$ después} & $\abs{M}$ \\
    \midrule
    1 & $s, A, 1, t$ & A--1 & $\set{A1}$ & 1 \\
    2 & $s, B, 3, t$ & B--3 & $\set{A1, B3}$ & 2 \\
    3 & $s, C, 1, A, 2, t$ & C--1--A--2 & $\set{A2, B3, C1}$ & 3 \\
    4 & $s, D, 1, C, 3, B, 4, t$ & D--1--C--3--B--4
      & $\set{A2, B4, C3, D1}$ & 4 \\
    \bottomrule
  \end{tabular}
  \end{center}
  En cada instantánea, el trazo naranja es el camino que se aplica en
  esa iteración, dibujado sobre el $M$ de antes. Los vértices se
  dibujan en el orden A, D, C, B y 2, 1, 3, 4 para que ninguna arista
  se cruce; el BFS los recorre en orden alfabético.
  \begin{center}
  \begin{tikzpicture}[caminoAum/.style={draw=orange!60,
      line width=6pt, line cap=round, opacity=0.55}]
    \foreach \k/\lst/\ok/\cam/\txt in
      {1/{}/{}/{A/1}/{it.\ 1: A--1},
       2/{A/1}/{A,1}/{B/3}/{it.\ 2: B--3},
       3/{A/1, B/3}/{A,B,1,3}/{C/1, 1/A, A/2}/{it.\ 3: C--1--A--2},
       4/{A/2, B/3, C/1}/{A,B,C,1,2,3}/{D/1, 1/C, C/3, 3/B, B/4}%
         /{it.\ 4: D--1--C--3--B--4},
       5/{A/2, B/4, C/3, D/1}/{A,B,C,D,1,2,3,4}/{}/{final}} {
      \begin{scope}[xshift={3.3cm*(\k-1)}, name prefix=p\k-]
        \foreach \v/\y in {A/0, D/-0.9, C/-1.8, B/-2.7}
          \coordinate (\v) at (0, \y);
        \foreach \v/\y in {2/0, 1/-0.9, 3/-1.8, 4/-2.7}
          \coordinate (\v) at (1.5, \y);
        \foreach \a/\b in \cam
          \draw[caminoAum] (\a) -- (\b);
        \foreach \v in {A, B, C, D, 1, 2, 3, 4}
          \node[vert] (\v) at (\v) {\v};
        \foreach \v in \ok
          \node[vertOk] at (\v) {\v};
        \foreach \a/\b in {A/1, A/2, B/3, B/4, C/1, C/3, D/1}
          \draw[arista] (\a) -- (\b);
        \foreach \a/\b in \lst
          \draw[aristaArbol] (\a) -- (\b);
        \node[font=\scriptsize] at (0.75, -3.35) {\txt};
      \end{scope}
    }
  \end{tikzpicture}
  \end{center}
  Las dos primeras iteraciones son voraces. En la tercera, Caro
  quiere el 1, que es de Ana, o el 3, que es de Beto; el BFS
  encuentra primero que Ana puede mudarse al 2, que está libre. En la
  cuarta, Dani solo acepta el 1, que ahora es de Caro: Caro se muda
  al 3, que es de Beto, y Beto al 4. Cada camino aumentante es una
  cadena de reasignaciones que termina en un proyecto libre, y el
  final es perfecto.
\end{ejemplo}

\subsubsection*{Las búsquedas de las iteraciones 3 y 4}

Con $M = \set{A1, B3}$, la parte de $G'_f$ que importa tiene $(s, C)$
y $(s, D)$, porque C y D están libres; $(C, 1)$, $(C, 3)$, $(D, 1)$,
$(A, 2)$ y $(B, 4)$, aristas fuera de $M$; $(1, A)$ y $(3, B)$, las
aristas de $M$ invertidas; y $(2, t)$ y $(4, t)$. El BFS desde $s$:
\begin{center}
\small
\begin{tabular}{cll}
  \toprule
  \textbf{Sale} & \textbf{Descubre} & \textbf{Cola después} \\
  \midrule
  $s$ & C, D & [C, D] \\
  C & 1, 3 & [D, 1, 3] \\
  D & nada: 1 ya visto & [1, 3] \\
  1 & A, por la arista invertida de $A1$ & [3, A] \\
  3 & B, por la arista invertida de $B3$ & [A, B] \\
  A & 2 & [B, 2] \\
  B & 4 & [2, 4] \\
  2 & $t$: se detiene & \\
  \bottomrule
\end{tabular}
\end{center}
Las distancias quedan $d[C] = d[D] = 1$, $d[1] = d[3] = 2$,
$d[A] = d[B] = 3$, $d[2] = d[4] = 4$ y $d[t] = 5$. Siguiendo $\pi$
desde $t$ se recupera $t, 2, A, 1, C, s$, o sea el camino
C--1--A--2. En \textsc{Camino-Aumentante} es la misma búsqueda sin
$s$ ni $t$ y con el paso de cada vértice de $R$ a su pareja hecho en
el acto: la cola pasa por [C, D], [D, A, B] y [A, B], y al sacar A se
descubre 2, que está libre.

En la iteración 4, con $M = \set{A2, B3, C1}$, la residual deja una
sola ruta, $s, D, 1, C, 3, B, 4, t$, con $d[t] = 7$. A y 2 ni
siquiera se alcanzan: a A solo se entra desde 2, por la arista
invertida de $A2$, y a 2 solo desde $t$. Las aristas de $M$ aparecen
invertidas, de $R$ a $L$; son las que permiten deshacer una
asignación, la cancelación de flujo de la sección~\ref{sec:flujo}.

\begin{center}
\begin{tikzpicture}[caminoAum/.style={draw=orange!60,
    line width=6pt, line cap=round, opacity=0.55}]
  \node[vert] (s) at (-2, -1.35) {$s$};
  \node[vert] (t) at (4, -1.35) {$t$};
  \foreach \v/\y in {A/0, D/-0.9, C/-1.8, B/-2.7}
    \node[vert] (\v) at (0, \y) {\v};
  \foreach \v/\y in {2/0, 1/-0.9, 3/-1.8, 4/-2.7}
    \node[vert] (\v) at (2, \y) {\v};
  \foreach \a/\b in {s/D, D/1, 1/C, C/3, 3/B, B/4, 4/t}
    \draw[caminoAum] (\a.center) -- (\b.center);
  \foreach \v in {A, B, C, D, 1, 2, 3, 4}
    \node[vert] at (\v) {\v};
  \node[vert] at (s) {$s$};
  \node[vert] at (t) {$t$};
  \foreach \a/\b in {s/D, A/1, D/1, C/3, B/4, 4/t}
    \draw[dirigida] (\a) -- (\b);
  \foreach \a/\b in {2/A, 1/C, 3/B}
    \draw[dirigidaArbol] (\a) -- (\b);
  \node[nota, align=center] at (1, 0.95)
    {gris: arista de $E - M$, de $L$ a $R$;
     azul: arista de $M$, invertida \\
     (se omiten las que entran a $s$ y las que salen de $t$)};
\end{tikzpicture}
\end{center}

% =====================================================================

\subsection*{Cuando no hay emparejamiento perfecto: Hall y König}

La \textbf{vecindad} de $X \subseteq V$ es
$N(X) = \set{y \in V : \set{x, y} \in E \text{ para algún }
x \in X}$. Si un conjunto $X$ de alumnos acepta entre todos menos
proyectos que alumnos tiene, $\abs{N(X)} < \abs{X}$, es imposible
acomodarlos a todos. Lo notable es que esa es la única obstrucción.

\begin{teorema}[Hall]
  \label{teo:emparejamiento-hall}
  Sea $G$ bipartito con $V = L \cup R$. Existe un emparejamiento que
  cubre todo $L$ si y solo si $\abs{A} \le \abs{N(A)}$ para todo
  $A \subseteq L$. En particular, si $\abs{L} = \abs{R}$, $G$ tiene
  un emparejamiento perfecto si y solo si se cumple esa condición.
\end{teorema}

Cormen lo pide en el ejercicio 26.3-4, en la versión con
$\abs{L} = \abs{R}$.

\begin{proof}
  Si un emparejamiento cubre $L$, cada $x \in A$ tiene una pareja
  distinta, y todas están en $N(A)$: $\abs{A} \le \abs{N(A)}$.

  La vuelta se prueba por contrapositiva, y el flujo la hace
  constructiva. Supóngase que el emparejamiento máximo $M$ que
  devuelve \textsc{Emparejamiento-Bipartito} no cubre $L$. Sea $f$
  su flujo, $S$ el conjunto de vértices alcanzables desde $s$ en
  $G'_f$, el mismo que aparece en la prueba del
  teorema~\ref{teo:flujo-corte}, y $X = L \cap S$. Con la
  proposición~\ref{prop:emparejamiento-residual} se leen tres hechos:
  \begin{enumerate}
    \item Todo vértice libre de $L$ está en $X$, por la arista
      $(s, u)$.
    \item $N(X) \subseteq R \cap S$. Sean $x \in X$ y
      $\set{x, v} \in E$. Si la arista no está en $M$, $(x, v)$ está
      en la residual y $v$ es alcanzable. Si está en $M$, la única
      arista que entra a $x$ en la residual es $(v, x)$, así que a
      $x$ se llegó desde $v$, y $v \in S$.
    \item Todo $v \in R \cap S$ está emparejado, porque si no,
      $(v, t)$ haría alcanzable a $t$ y habría camino aumentante. Y
      su pareja está en $X$, por la arista $(v, u)$.
  \end{enumerate}
  Por 2 y 3, la pareja es una biyección entre $R \cap S$ y los
  vértices emparejados de $X$, y además $N(X) = R \cap S$. Por 1,
  \[
    \abs{N(X)} = \abs{X} - (\text{número de libres de } L)
               < \abs{X} ,
  \]
  y $X$ viola la condición.
\end{proof}

\begin{observacion}
  La misma cuenta da el tamaño del máximo:
  \[
    \abs{M^*} = \abs{L} - \max_{A \subseteq L}
                \bigl(\abs{A} - \abs{N(A)}\bigr) .
  \]
  Para cualquier $A$, a lo más $\abs{N(A)}$ de sus vértices quedan
  emparejados, así que al menos $\abs{A} - \abs{N(A)}$ quedan libres;
  el $X$ de la prueba alcanza la igualdad. Con $A = \emptyset$ se ve
  que el máximo de la derecha nunca es negativo.
\end{observacion}

\begin{ejemplo}[Tres alumnos para dos proyectos]
  \label{ej:emparejamiento-hall}
  Los mismos cuatro alumnos, con otras preferencias:
  \begin{center}
  \small
  \begin{tabular}{ll}
    \toprule
    \textbf{Alumno} & \textbf{Proyectos que acepta} \\
    \midrule
    A (Ana) & 1, 2 \\
    B (Beto) & 2, 3, 4 \\
    C (Caro) & 1 \\
    D (Dani) & 1, 2 \\
    \bottomrule
  \end{tabular}
  \end{center}
  Así $E = \set{A1, A2, B2, B3, B4, C1, D1, D2}$, con $n = 8$ y
  $m = 8$. \textsc{Emparejamiento-Máximo}, con el mismo orden de
  vecinos, encuentra $A1$, después $B2$, y en la tercera iteración el
  camino D--2--B--3: Dani toma el 2 y Beto se muda al 3. Queda
  $M = \set{A1, B3, D2}$. La cuarta búsqueda arranca del único libre
  de $L$, que es Caro, y falla:
  \begin{center}
  \small
  \begin{tabular}{cll}
    \toprule
    \textbf{Sale} & \textbf{Descubre} & \textbf{Cola después} \\
    \midrule
    C & 1, que es de A; se encola A & [A] \\
    A & 2, que es de D; se encola D & [D] \\
    D & nada: 1 ya visto y 2 es su pareja & [\,] \\
    \bottomrule
  \end{tabular}
  \end{center}
  Los alcanzables son $S = \set{s, C, 1, A, 2, D}$, así que
  $X = \set{A, C, D}$ y $N(X) = \set{1, 2}$:
  $\abs{N(X)} = 2 < 3 = \abs{X}$. Tres alumnos compiten por dos
  proyectos, y ningún algoritmo puede acomodar a los cuatro. La
  deficiencia $\abs{X} - \abs{N(X)} = 1$ es justo
  $\abs{L} - \abs{M} = 4 - 3$. El proyecto 4 queda libre: Beto podría
  tomarlo, pero eso no ayuda a nadie de $X$. El subconjunto se llama
  $X$ y no $A$ porque Ana ya es A.
  \begin{center}
  \begin{tikzpicture}
    \draw[zona] (-0.5, 0.45) rectangle (0.5, -2.25);
    \draw[zona] (1.3, 0.45) rectangle (2.3, -1.35);
    \foreach \v/\y in {C/0, A/-0.9, D/-1.8, B/-2.7}
      \node[vert] (\v) at (0, \y) {\v};
    \foreach \v/\y in {1/0, 2/-0.9, 3/-1.8, 4/-2.7}
      \node[vert] (\v) at (1.8, \y) {\v};
    \foreach \v in {A, D, B, 1, 2, 3}
      \node[vertOk] at (\v) {\v};
    \foreach \a/\b in {C/1, A/1, A/2, D/1, D/2, B/2, B/3, B/4}
      \draw[arista] (\a) -- (\b);
    \foreach \a/\b in {A/1, D/2, B/3}
      \draw[aristaArbol] (\a) -- (\b);
    \node[nota, left=6pt] at (-0.5, -0.9) {$X$};
    \node[nota, right=6pt] at (2.3, -0.45) {$N(X)$};
    \begin{scope}[xshift=5.6cm]
      \node[vert] (tC) at (0, 0) {C};
      \node[vertOk] (t1) at (0, -0.9) {1};
      \node[vertOk] (tA) at (0, -1.8) {A};
      \node[vertOk] (t2) at (0, -2.7) {2};
      \node[vertOk] (tD) at (0, -3.6) {D};
      \draw[arista] (tC) -- (t1);
      \draw[aristaArbol] (t1) -- (tA);
      \draw[arista] (tA) -- (t2);
      \draw[aristaArbol] (t2) -- (tD);
      \draw[aristaFuera] (tD) to[out=0, in=0, looseness=1.3] (t1);
      \node[nota, anchor=west] at (1.5, -2.25) {D--1: 1 ya visto};
      \node[nota, above=2pt] at (tC.north) {raíz libre};
    \end{scope}
  \end{tikzpicture}
  \end{center}
  A la izquierda, $X$ y $N(X)$ sombreados. A la derecha, el árbol de
  la búsqueda que falla: alterna una arista fuera de $M$ y una de $M$,
  y se agota sin tocar un libre de $R$.
\end{ejemplo}

Del mismo conjunto $S$ sale un segundo certificado. Una
\textbf{cubierta de vértices} de $G$ es un conjunto $C \subseteq V$
que toca todas las aristas.

\begin{teorema}[König]
  \label{teo:emparejamiento-konig}
  En un grafo bipartito, el tamaño de un emparejamiento máximo es
  igual al de una cubierta de vértices mínima.
\end{teorema}

\begin{proof}
  Para cualquier emparejamiento $M'$ y cualquier cubierta $C'$,
  $\abs{M'} \le \abs{C'}$: las aristas de $M'$ no comparten extremos,
  y cada una necesita su propio vértice de $C'$. Basta entonces
  exhibir una cubierta del tamaño del máximo $M$. Los tres hechos de
  la prueba de Hall no usan que $M$ deje libres en $L$, así que valen
  aquí. Sea $C = (L - S) \cup (R \cap S)$. Toca todas las aristas: si
  $\set{u, v} \in E$ con $u \in L$, la toca $u$ cuando $u \notin S$, y
  cuando $u \in S$ el hecho 2 pone a $v$ en $R \cap S$. Y mide
  $\abs{M}$: todo vértice de $C$ está emparejado (los libres de $L$
  están en $S$ por el hecho 1, y los de $R \cap S$ están emparejados
  por el hecho 3), y cada arista de $M$ tiene exactamente un extremo
  en $C$, porque por los hechos 2 y 3 sus dos extremos están en $S$ o
  los dos fuera.
\end{proof}

Cormen no lo menciona, pero vale la pena: para convencer a alguien de
que $M$ es máximo basta enseñarle $\abs{M}$ vértices que tocan todas
las aristas, y eso se revisa a mano. En la figura 26.8,
$\set{3, 6, 8}$ toca las ocho aristas, así que ningún emparejamiento
pasa de $3$. En el ejemplo~\ref{ej:emparejamiento-hall}, la cubierta es
$\set{B} \cup \set{1, 2}$. La cubierta es el corte mínimo leído en
$G$: las aristas de $G'$ que van de $S$ a $V' - S$ son exactamente las
$(s, u)$ con $u \in L - S$ y las $(v, t)$ con $v \in R \cap S$.

% =====================================================================

\subsection*{Hopcroft-Karp: varios caminos por fase}

\textsc{Emparejamiento-Máximo} hace un BFS completo para agregar una
sola arista. Hopcroft y Karp (problema 26-6 de Cormen) aprovechan
cada BFS para varios caminos. Sea $\ell$ la longitud del camino
aumentante más corto respecto a $M$. Una \textbf{fase} encuentra un
conjunto \emph{maximal} de caminos aumentantes de longitud $\ell$,
ajenos en vértices, y los aplica todos juntos.

\begin{algorithm}[H]
\caption{\textsc{Hopcroft-Karp}$(G)$}
\label{alg:emparejamiento-hopcroft-karp}
\begin{algorithmic}[1]
\Procedure{Hopcroft-Karp}{$G$}
  \State $M \gets \emptyset$
  \Repeat
    \State $\ell \gets$ longitud del camino aumentante más corto
           respecto a $M$
    \State $\mathcal{P} \gets$ un conjunto maximal de caminos
           aumentantes de longitud $\ell$, ajenos en vértices
    \State $M \gets M \oplus (P_1 \cup \dots \cup P_k)$, donde
           $\mathcal{P} = \set{P_1, \dots, P_k}$
  \Until{$\mathcal{P} = \emptyset$}
    \Comment{no queda camino aumentante}
  \State \Return $M$
\EndProcedure
\end{algorithmic}
\end{algorithm}

Que se puedan aplicar juntos sale del
lema~\ref{lem:emparejamiento-diferencia}: como los $P_i$ no comparten
vértices, aplicar uno no toca a los otros, y
$M \oplus (P_1 \cup \dots \cup P_k)$ es un emparejamiento de tamaño
$\abs{M} + k$. Una fase cuesta $\Ord{m}$. Un BFS desde todos los
libres de $L$ a la vez, como en \textsc{Camino-Aumentante}, arma las
capas y da $\ell$; se detiene al completar la capa donde aparece el
primer libre de $R$. Después, un DFS que solo baja de una capa a la
siguiente saca los caminos uno por uno y borra cada vértice que usa o
que resulta un callejón sin salida, así que cada arista se revisa una
sola vez en toda la fase.

El número de fases descansa en tres hechos, que Cormen deja como
partes del problema:
\begin{enumerate}
  \item Después de una fase, el camino aumentante más corto mide más
    de $\ell$. Las longitudes son impares, así que $\ell$ sube al
    menos $2$ por fase. Esta es la parte larga del problema y su
    prueba se omite.
  \item Si el más corto mide $\ell$, entonces
    $\abs{M^*} \le \abs{M} + n/(\ell + 1)$. En la prueba del teorema
    de Berge, cada componente de $M \oplus M^*$ tiene a lo más una
    arista más de $M^*$ que de $M$, así que hay al menos
    $\abs{M^*} - \abs{M}$ componentes que son caminos aumentantes, y
    son ajenos en vértices. Cada uno tiene al menos $\ell$ aristas, o
    sea $\ell + 1$ vértices, y no hay más de $n$ vértices.
  \item Después de $\sqrt{n}$ fases, $\ell > \sqrt{n}$, y por el
    hecho 2 faltan menos de $n / \sqrt{n} = \sqrt{n}$ aristas. Cada
    fase agrega al menos una, así que quedan menos de $\sqrt{n}$
    fases: a lo más $2\sqrt{n}$ en total.
\end{enumerate}
El costo total es $\Ord{\sqrt{n}\, m}$, contra $\Ord{nm}$ de
\textsc{Emparejamiento-Máximo}.

\begin{ejemplo}[Hopcroft-Karp sobre alumnos y proyectos]
  \label{ej:emparejamiento-hopcroft-karp}
  En el ejemplo~\ref{ej:emparejamiento-alumnos} los caminos que aplicó
  Edmonds-Karp medían $1$, $1$, $3$ y $5$, y eso ya sugiere las
  fases. El DFS recorre los vecinos en orden alfabético.
  \begin{center}
  \small
  \begin{tabular}{cclll}
    \toprule
    \textbf{Fase} & $\ell$ & \textbf{Caminos más cortos}
      & \textbf{Conjunto aplicado} & $\abs{M}$ \\
    \midrule
    1 & 1 & las 7 aristas & A--1, B--3 & 2 \\
    2 & 3 & C--1--A--2, C--3--B--4, D--1--A--2 & C--1--A--2 & 3 \\
    3 & 5 & D--1--C--3--B--4 & el mismo & 4 \\
    \bottomrule
  \end{tabular}
  \end{center}
  En la fase 1, A toma el 1 y B el 3, y C y D quedan bloqueados.
  $\ell$ crece estrictamente, $1$, $3$, $5$, y hay tres fases en lugar
  de cuatro búsquedas. La fase 2 enseña la diferencia entre maximal y
  máximo para conjuntos de caminos: $\set{\text{C--1--A--2}}$ es
  maximal, porque los otros dos caminos de longitud $3$ chocan con
  él, pero C--3--B--4 y D--1--A--2 son ajenos y habrían terminado el
  trabajo en esa misma fase.
\end{ejemplo}

\begin{ojo}
  Hopcroft-Karp pide en cada fase un conjunto maximal de caminos, no
  uno máximo. Encontrar el máximo número de caminos ajenos sería otro
  problema de flujo; el maximal sale con un DFS voraz en $\Ord{m}$ y
  basta para la cota. La primera fase, con $M = \emptyset$ y
  $\ell = 1$, es exactamente el voraz del principio de la sección:
  construye un emparejamiento maximal.
\end{ojo}

% =====================================================================

\subsection*{Grafos que no son bipartitos}

La reducción a flujo necesita la partición en $L$ y $R$ para orientar
las aristas. En un triángulo el emparejamiento máximo mide $1$, pero
no hay manera de repartir sus vértices en dos lados. El teorema de
Berge sigue valiendo, con camino aumentante definido como un camino
alternante entre dos vértices libres: su prueba usó la bipartición
solo para nombrar los extremos. Lo que se rompe es la búsqueda.

\begin{ejemplo}[Una búsqueda alternante que se pierde]
  \label{ej:emparejamiento-flor}
  El grafo es el ciclo $a, b, c, d, e$ de longitud $5$ más la arista
  $\set{e, y}$, con $M = \set{\set{b,c}, \set{d,e}}$. Los libres son
  $a$ e $y$, y $P = a$--$b$--$c$--$d$--$e$--$y$ es un camino
  aumentante: $M \oplus P = \set{\set{a,b}, \set{c,d}, \set{e,y}}$
  es perfecto.
  \begin{center}
  \begin{tikzpicture}[caminoAum/.style={draw=orange!60,
      line width=6pt, line cap=round, opacity=0.55}]
    \foreach \pre/\x/\lst/\cam in
      {i/0/{b/c, d/e}/{a/b, b/c, c/d, d/e, e/y},
       d/5.4/{a/b, c/d, e/y}/{}} {
      \begin{scope}[xshift=\x cm]
        \coordinate (\pre a) at (0, 1.1);
        \coordinate (\pre b) at (1.05, 0.34);
        \coordinate (\pre c) at (0.65, -0.89);
        \coordinate (\pre d) at (-0.65, -0.89);
        \coordinate (\pre e) at (-1.05, 0.34);
        \coordinate (\pre y) at (-2.3, 0.34);
        \foreach \p/\q in \cam
          \draw[caminoAum] (\pre\p) -- (\pre\q);
        \foreach \p in {a, b, c, d, e, y}
          \node[vert] (\pre\p) at (\pre\p) {$\p$};
        \foreach \p/\q in {a/b, b/c, c/d, d/e, e/a, e/y}
          \draw[arista] (\pre\p) -- (\pre\q);
        \foreach \p/\q in \lst
          \draw[aristaArbol] (\pre\p) -- (\pre\q);
      \end{scope}
    }
    \node[font=\small] at (-0.6, -1.6) {$M$ y el camino $P$};
    \node[font=\small] at (4.8, -1.6) {$M \oplus P$};
  \end{tikzpicture}
  \end{center}
  Una búsqueda como \textsc{Camino-Aumentante}, que desde un vértice
  emparejado salta a su pareja, sale de $a$ y descubre $b$ y $e$ por
  aristas fuera de $M$; de ahí salta a sus parejas, $c$ y $d$. Desde
  $c$ y $d$ la única arista fuera de $M$ es $\set{c, d}$, que une dos
  vértices ya vistos, y la búsqueda se detiene sin encontrar nada. El
  problema es $e$: se le marcó como visto al llegar desde $a$, en el
  papel de un vértice de $R$, y desde ahí solo se siguió a su pareja.
  El camino aumentante necesita llegar a $e$ por su arista de $M$,
  desde $d$, y salir por $\set{e, y}$. En un bipartito cada vértice
  tiene un solo papel posible; aquí $e$ tiene dos, porque el ciclo
  $a, b, c, d, e$ es impar.
\end{ejemplo}

La arista $\set{c, d}$ une dos vértices a los que se llegó por su
arista de $M$ y cierra el ciclo impar; por el
teorema~\ref{teo:bipartito-ciclos}, eso no puede pasar en un
bipartito. Edmonds (1965) llamó \textbf{flor} (\textit{blossom}) a un
ciclo impar así, y su algoritmo la contrae a un solo vértice, sigue
buscando en el grafo más chico y, cuando encuentra un camino
aumentante, lo expande por el lado correcto de la flor. En el
ejemplo, contraer $a, b, c, d, e$ deja un vértice libre unido a $y$:
un camino aumentante de longitud $1$, que expandido es
$a$--$b$--$c$--$d$--$e$--$y$. El algoritmo de flores corre en tiempo
polinomial y queda fuera del capítulo 26 de Cormen.

Por eso, antes de aplicar la reducción hay que verificar que el grafo
es bipartito y obtener $L$ y $R$, con el
algoritmo~\ref{alg:bipartito-bfs} de la
sección~\ref{sec:hamilton-arboles}, en $\Tht{n + m}$. El grafo $G$ de
la libreta, por ejemplo, no es bipartito: tiene el ciclo
A, B, E, F, C de longitud $5$.

% =====================================================================

\subsection*{Ejercicios}

\begin{ejercicio}[Cormen 26.3-1]
  Correr Ford-Fulkerson sobre la red de la figura 26.8(c), tomando en
  cada iteración el camino lexicográficamente menor, y dar el corte
  mínimo final.

  \textit{Solución.} Tres aumentos: $s, 1, 6, t$; luego
  $s, 2, 8, t$; luego $s, 3, 7, t$; la cuarta búsqueda falla. En la
  segunda iteración, $s, 2, 6, \dots$ sería menor, pero desde $6$ solo
  se puede volver a $1$ por la arista invertida, y desde $1$ solo a
  $s$: callejón. Queda $M = \set{\set{1,6}, \set{2,8}, \set{3,7}}$,
  máximo y distinto de $M_b$. Al final los alcanzables son
  $S = \set{s, 1, 2, 4, 5, 6, 8}$; el corte lo cruzan $(s, 3)$,
  $(6, t)$ y $(8, t)$, de capacidad $3$, y la cubierta de König es
  $\set{3, 6, 8}$.
\end{ejercicio}

\begin{ejercicio}[Cormen 26.3-3]
  Acotar la longitud de un camino aumentante en $G'$.

  \textit{Solución.} Por la proposición~\ref{prop:emparejamiento-residual}
  el camino es $s, u_0, v_0, \dots, u_k, v_k, t$, con los $u_i$
  distintos en $L$ y los $v_i$ distintos en $R$. Entonces
  $k + 1 \le \min(\abs{L}, \abs{R})$, y la longitud, $2(k + 1) + 1$,
  es a lo más $2\min(\abs{L}, \abs{R}) + 1$.
\end{ejercicio}

\begin{ejercicio}[Cormen 26.3-5]
  Un grafo es $d$-regular si todo vértice tiene grado $d \ge 1$.
  Probar que todo bipartito $d$-regular tiene un emparejamiento
  perfecto.

  \textit{Solución.} Contando las aristas desde cada lado,
  $d\abs{L} = m = d\abs{R}$, así que $\abs{L} = \abs{R}$. Con Hall:
  para $A \subseteq L$, las $d\abs{A}$ aristas que salen de $A$ caen
  en $N(A)$, cuyos vértices reciben $d$ aristas cada uno en total, así
  que $d\abs{A} \le d\abs{N(A)}$. Con flujo: poner $1/d$ en cada
  arista de en medio y $1$ en las de $s$ y de $t$ es un flujo
  fraccionario de valor $\abs{L}$, que es lo más que permite el corte
  alrededor de $s$. Entonces el flujo máximo vale $\abs{L}$, y por el
  teorema~\ref{teo:emparejamiento-integralidad} Ford-Fulkerson lo alcanza con un
  flujo entero, que es un emparejamiento perfecto.
\end{ejercicio}

\begin{ejercicio}[Cupos]
  Cada proyecto $v$ admite hasta $k_v$ alumnos y cada alumno toma un
  solo proyecto. ¿Qué cambia en la reducción?

  \textit{Solución.} Solo la capacidad de $(v, t)$, que pasa a ser
  $k_v$. Las capacidades siguen siendo enteras, así que el
  teorema~\ref{teo:emparejamiento-integralidad} da un flujo máximo entero: cada
  alumno manda $0$ o $1$ unidad por una sola arista y a cada proyecto
  le llegan a lo más $k_v$. Lo mismo vale para máquinas que hacen
  hasta $k$ tareas, cambiando la capacidad de $(s, u)$. Si además los
  alumnos tienen preferencias con peso, el problema ya es otro, el de
  asignación, que se resuelve con el algoritmo húngaro.
\end{ejercicio}

\begin{ejercicio}[Bipartitos completos]
  ¿Cuánto mide el emparejamiento máximo de $K_{p,q}$? ¿Cuántos
  emparejamientos perfectos tiene $K_{r,r}$?

  \textit{Solución.} $\min(p, q)$: cualquier emparejamiento deja
  libres al menos $\abs{p - q}$ vértices del lado grande, y en
  $K_{p,q}$ cualquier vértice del lado chico se puede unir con
  cualquiera del grande. En $K_{r,r}$ el primer vértice de $L$ escoge
  pareja entre $r$, el segundo entre los $r - 1$ que quedan, y así:
  $r!$ emparejamientos perfectos, uno por cada permutación. Aquí se
  usa $r$ y no $n$ porque $K_{r,r}$ tiene $n = 2r$ vértices.
\end{ejercicio}

\newpage

% =====================================================================
%  Grafos - Empujar y reetiquetar
% =====================================================================

\clase{Empujar y reetiquetar}{Complemento: Cormen, cap.\ 26.4 y 26.5}
\label{sec:push-relabel}

\begin{quick_sum}
  Ford-Fulkerson mejora un flujo buscando caminos completos de $s$ a
  $t$. Los algoritmos de empujar y reetiquetar trabajan al revés:
  saturan de golpe todo lo que sale de la fuente y después cada
  vértice resuelve, mirando solo a sus vecinos, qué hacer con lo que le
  sobra. Lo empuja cuesta abajo, o sube de altura cuando no tiene a
  dónde empujarlo. Durante la ejecución no hay un flujo sino un
  preflujo, y una función de altura garantiza desde el primer paso que
  no existe camino aumentante; cuando ningún vértice interno guarda
  exceso, el preflujo es un flujo y es máximo. El algoritmo genérico
  cuesta $\Ord{n^2 m}$ sin importar las capacidades. Reetiquetar al
  frente fija el orden de las operaciones con una lista que siempre es
  orden topológico de las aristas por las que se puede empujar, y baja
  a $\Ord{n^3}$. La sección cierra con una tabla de todos los
  algoritmos de la parte de grafos.
\end{quick_sum}

% =====================================================================

\subsection*{La idea: agua que corre cuesta abajo}

El método de Ford-Fulkerson (algoritmo~\ref{alg:ford-fulkerson})
mantiene un flujo válido en todo momento y, en cada iteración, busca en
la red residual un camino completo de $s$ a $t$. Para mover unas
cuantas unidades mira la red entera. Los algoritmos de
\textbf{empujar y reetiquetar} (\textit{push-relabel}, de Goldberg y de
Goldberg y Tarjan) son locales: operan sobre un vértice a la vez y solo
miran las aristas residuales que salen de él. A cambio, durante la
ejecución renuncian a la conservación de flujo. A un vértice le puede
entrar más de lo que sale, y lo que sobra se queda guardado en él.

Cormen lo explica con fluidos. Las aristas son tuberías y los vértices
son uniones de tuberías. Cada vértice tiene además un depósito donde
se acumula el agua que entró y todavía no sale. Todo el sistema está
montado sobre plataformas de altura ajustable, y el agua solo corre
cuesta abajo. La fuente está fija a altura $n$ y el sumidero a altura
$0$; los demás vértices empiezan en $0$ y solo suben.

\begin{itemize}
  \item Al inicio se abre la llave por completo: todas las tuberías que
        salen de $s$ se llenan a su capacidad y el agua cae en los
        depósitos de los vecinos de $s$.
  \item Un vértice con agua en su depósito la empuja hacia un vecino
        que esté exactamente un nivel abajo, por una tubería que
        todavía tenga holgura.
  \item Si todas sus tuberías con holgura llevan a vecinos de su nivel
        o más altos, el vértice se levanta: se \textbf{reetiqueta} a un
        nivel por encima del más bajo de esos vecinos, y entonces sí
        puede vaciarse.
\end{itemize}

Lo que cabe hasta $t$ llega a $t$. Cuando las tuberías que entran a
$t$ se saturan, el agua que sobra se queda atorada en los depósitos;
esos vértices siguen subiendo, rebasan la altura de $s$ y el sobrante
regresa a la fuente por las mismas tuberías por las que llegó. Al
final ningún depósito interno tiene agua y lo que llegó a $t$ es un
flujo máximo.

\begin{idea}
  Ford-Fulkerson conserva la \textbf{factibilidad} (siempre tiene un
  flujo) y trabaja hasta alcanzar la \textbf{optimalidad} (que no
  quede camino aumentante). Empujar y reetiquetar conserva la
  optimalidad desde el primer paso (nunca hay camino de $s$ a $t$ en
  la red residual) y trabaja hasta alcanzar la factibilidad (que se
  vacíen los depósitos). Los dos terminan en el mismo punto, un flujo
  sin camino aumentante, que es máximo por el teorema de flujo máximo
  y corte mínimo (teorema~\ref{teo:flujo-corte}).
\end{idea}

\begin{ejemplo}[Caminos contra tuberías]
  \label{ej:push-tuberias}
  La misma red de cuatro vértices vista por los dos métodos. A la
  izquierda, lo que hace Ford-Fulkerson: un camino aumentante de $s$ a
  $t$, un chorro sin ramas. A la derecha, la red como tuberías justo
  después del primer reetiquetado del ejemplo~\ref{ej:push-generico}:
  $s$ en lo alto, las dos tuberías que salen de $s$ saturadas, los
  depósitos de $a$ y $b$ con su exceso, y $a$ ya levantado un nivel
  para vaciarse hacia $b$ y $t$. $b$ sigue al nivel de $t$ y todavía no
  puede empujar.
  \begin{center}
  \begin{tikzpicture}[nivel/.style={draw=gray!25, thin}]
    % --- izquierda: Ford-Fulkerson ---
    \begin{scope}
      \node[vert] (s) at (0, 0) {$s$};
      \node[vert] (a) at (2, 1.2) {$a$};
      \node[vert] (b) at (2, -1.2) {$b$};
      \node[vert] (t) at (4, 0) {$t$};
      \draw[dirigida] (s) -- node[peso] {2} (b);
      \draw[dirigida] (a) -- node[peso] {1} (t);
      \draw[dirigidaArbol] (s) -- node[peso] {4} (a);
      \draw[dirigidaArbol] (a) -- node[peso] {2} (b);
      \draw[dirigidaArbol] (b) -- node[peso] {3} (t);
      \node[nota, align=center] at (2, -2.3)
        {Ford-Fulkerson: un camino $s, a, b, t$\\
         en $G_f$, un chorro sin ramas};
    \end{scope}
    % --- derecha: la misma red como tuberias ---
    \begin{scope}[xshift=6.4cm, yshift=-1.5cm, y=0.8cm]
      \foreach \k in {0, ..., 4} {
        \draw[nivel] (-0.5, \k) -- (7.0, \k);
        \node[nota] at (-0.8, \k) {\k};
      }
      \node[nota] at (-0.8, 4.75) {$h$};
      % plataformas
      \foreach \x/\y in {0/4, 2.2/0, 4.2/1, 6.4/0}
        \draw[gray!60!black, line width=2pt]
          (\x - 0.45, \y - 0.42) -- (\x + 0.45, \y - 0.42);
      \node[vert] (s) at (0, 4) {$s$};
      \node[vert] (b) at (2.2, 0) {$b$};
      \node[vert] (a) at (4.2, 1) {$a$};
      \node[vert] (t) at (6.4, 0) {$t$};
      % depositos encima de cada vertice; 4 unidades miden 1.0
      \foreach \x/\y/\e in {4.2/1/4, 2.2/0/2} {
        \fill[acento!35] (\x - 0.17, \y + 0.42)
          rectangle (\x + 0.17, \y + 0.42 + 0.25 * \e);
        \draw[gray!60!black] (\x - 0.17, \y + 1.5)
          -- (\x - 0.17, \y + 0.42) -- (\x + 0.17, \y + 0.42)
          -- (\x + 0.17, \y + 1.5);
        \node[nota, anchor=west] at (\x + 0.2, \y + 1.2) {$e = \e$};
      }
      \draw[dirigidaArbol] (s) -- node[peso, pos=0.4] {4/4} (a);
      \draw[dirigidaArbol] (s) -- node[peso, pos=0.45] {2/2} (b);
      \draw[dirigida] (a) -- node[peso] {0/2} (b);
      \draw[dirigida] (a) -- node[peso] {0/1} (t);
      \draw[dirigida] (b) to[bend right=22]
        node[peso, pos=0.6] {0/3} (t);
      \node[nota, align=center] at (3.1, -1.6)
        {Empujar y reetiquetar: $s$ fijo a altura $n = 4$. $a$ ya
         subió a $1$\\y puede vaciar su depósito cuesta abajo; $b$
         sigue al nivel de $t$};
    \end{scope}
  \end{tikzpicture}
  \end{center}
  En esta edición Cormen describe las tuberías solo en prosa; la
  figura es propia. En todas las figuras que siguen la altura de cada
  vértice es su coordenada vertical, con líneas grises en cada nivel;
  dentro del círculo va el exceso y en cada arista el par $f/c$.
\end{ejemplo}

% =====================================================================

\subsection*{Preflujo, exceso y alturas}

Se trabaja con una red de flujo $G = (V, E)$ como la de la
sección~\ref{sec:flujo}: dirigida, con capacidades $c(u, v) \ge 0$,
fuente $s$, sumidero $t$, sin aristas antiparalelas y con cada vértice
en algún camino de $s$ a $t$, de donde $m \ge n - 1$. La capacidad
residual $c_f$ y la red residual $G_f = (V, E_f)$ son las de la
definición~\ref{def:red-residual}. Esa definición solo usa que
$0 \le f(u, v) \le c(u, v)$, así que vale tal cual para lo que sigue.

\begin{definicion}[Preflujo]
  \label{def:preflujo}
  Un \textbf{preflujo} (\textit{preflow}) es una función
  $f : V \times V \to \R$ que cumple la restricción de capacidad,
  $0 \le f(u, v) \le c(u, v)$, y la conservación relajada
  \[
    \sum_{v \in V} f(v, u) - \sum_{v \in V} f(u, v) \ \ge\ 0
    \qquad \text{para toda } u \in V - \set{s} .
  \]
  El \textbf{exceso} (\textit{excess flow}) de $u$ es lo que entra
  menos lo que sale,
  \[
    e(u) = \sum_{v \in V} f(v, u) - \sum_{v \in V} f(u, v) .
  \]
  Un vértice $u \in V - \set{s, t}$ con $e(u) > 0$ está
  \textbf{desbordado} (\textit{overflowing}).
\end{definicion}

En un preflujo $e(u) \ge 0$ para todo $u \neq s$. Como cada $f(x, y)$
suma en $e(y)$ y resta en $e(x)$, los excesos suman cero y por lo tanto
$e(s) \le 0$. Por definición $s$ y $t$ nunca están desbordados, aunque
$t$ acumule exceso: ahí es justamente donde debe terminar el agua. Un
preflujo sin vértices desbordados cumple conservación en todo
$V - \set{s, t}$, es decir, es un flujo.

\begin{definicion}[Función de altura]
  \label{def:altura-push}
  Una \textbf{función de altura} (\textit{height function}) es una
  función $h : V \to \N$ con
  \[
    h(s) = n , \qquad h(t) = 0 , \qquad
    h(u) \le h(v) + 1 \ \text{ para toda arista residual }
    (u, v) \in E_f .
  \]
\end{definicion}

La condición dice que una arista residual puede subir cuanto quiera,
pero baja a lo más un nivel. Leída al revés: si $h(u) > h(v) + 1$,
entonces $(u, v)$ no es arista residual (Cormen, lema 26.12). En la
literatura a $h$ se le llama también \textit{distance function}: si en
$G_f$ hay un camino de $u$ a $t$ con $k$ aristas, sumando la condición
a lo largo de él sale $h(u) \le k$, así que $h(u)$ nunca supera la
distancia en aristas de $u$ a $t$. La última parte de la sección
vuelve sobre esto.

\begin{observacion}
  Cormen usa atributos en el pseudocódigo ($u.h$, $u.e$, $(u, v).f$,
  $u.N$, $u.\mathit{current}$) y funciones en los lemas ($h(u)$,
  $e(u)$, $f(u, v)$). Aquí se usa la forma funcional en todas partes:
  $h(u)$, $e(u)$, $f(u, v)$, $N(u)$ y $\text{actual}(u)$.
\end{observacion}

% =====================================================================

\subsection*{Las dos operaciones básicas}

\subsubsection*{Empujar}

\textsc{Empujar}$(u, v)$ aplica cuando $u$ está desbordado,
$c_f(u, v) > 0$ y $h(u) = h(v) + 1$: hay holgura de $u$ a $v$ y $v$
está exactamente un nivel abajo. Mueve
$\Delta = \min\bigl(e(u), c_f(u, v)\bigr)$ unidades de $u$ a $v$. Es
un \textbf{empuje saturante} si después $c_f(u, v) = 0$, y entonces
$(u, v)$ desaparece de la red residual; si no, es \textbf{no
saturante}.

\begin{algorithm}[H]
\caption{\textsc{Empujar}$(u, v)$ \ (\textsc{Push} en Cormen)}
\label{alg:empujar}
\begin{algorithmic}[1]
\Procedure{Empujar}{$u, v$}
  \Comment{aplica si $e(u) > 0$, $c_f(u, v) > 0$ y $h(u) = h(v) + 1$}
  \State $\Delta \gets \min\bigl(e(u),\ c_f(u, v)\bigr)$
  \If{$(u, v) \in E$}
    \State $f(u, v) \gets f(u, v) + \Delta$ \Comment{arista original}
  \Else
    \State $f(v, u) \gets f(v, u) - \Delta$ \Comment{arista de regreso}
  \EndIf
  \State $e(u) \gets e(u) - \Delta$, \quad $e(v) \gets e(v) + \Delta$
\EndProcedure
\end{algorithmic}
\end{algorithm}

\begin{itemize}
  \item \textbf{Línea 2.} $\Delta$ es lo más que se puede mover sin
        dejar $e(u)$ negativo y sin rebasar la capacidad residual.
  \item \textbf{Líneas 3 a 7.} La arista residual $(u, v)$ puede ser
        original o de regreso. Si es original, se aumenta su flujo. Si
        la original es $(v, u)$, empujar de $u$ a $v$ cancela flujo que
        iba de $v$ a $u$. Por eso Cormen prohíbe las antiparalelas: con
        $(u, v)$ y $(v, u)$ en $E$, la pregunta de la línea 3 no decide
        nada.
  \item El procedimiento no revisa alturas; $h(u) = h(v) + 1$ es una
        precondición que garantiza quien lo llama. Cuesta $\Ord{1}$ si
        $c_f$ se calcula en tiempo constante a partir de $c$ y $f$.
\end{itemize}

\begin{lema}[Cormen, lema 26.13]
  \label{lem:push-no-saturante}
  Después de un empuje no saturante de $u$ a $v$, $u$ ya no está
  desbordado.
\end{lema}

\begin{proof}
  Si el empuje no saturó, después queda $c_f(u, v) - \Delta > 0$, así
  que $\Delta \neq c_f(u, v)$ y el mínimo lo dio $e(u)$. Se empujó
  todo el exceso y $e(u)$ queda en $0$.
\end{proof}

\subsubsection*{Reetiquetar}

\textsc{Reetiquetar}$(u)$ aplica cuando $u$ está desbordado y
$h(u) \le h(v)$ para toda arista residual $(u, v)$: tiene agua y no
hay nada abajo hacia dónde empujarla.

\begin{algorithm}[H]
\caption{\textsc{Reetiquetar}$(u)$ \ (\textsc{Relabel} en Cormen)}
\label{alg:reetiquetar}
\begin{algorithmic}[1]
\Procedure{Reetiquetar}{$u$}
  \Comment{aplica si $e(u) > 0$ y $h(u) \le h(v)$ para toda
           $(u, v) \in E_f$}
  \State $h(u) \gets 1 + \min\set{h(v) : (u, v) \in E_f}$
\EndProcedure
\end{algorithmic}
\end{algorithm}

La nueva altura es la mayor que permite la condición
$h(u) \le h(v) + 1$ sobre las aristas residuales que salen de $u$. Con
ella queda al menos una arista por la que se puede empujar: la que va
al vecino que dio el mínimo. El conjunto sobre el que se toma el
mínimo no es vacío. Como $e(u) > 0$, a $u$ le entra flujo por alguna
arista $(v, u)$, y entonces $c_f(u, v) \ge f(v, u) > 0$. Recorrer las
aristas residuales de $u$ cuesta $\Ord{\text{gr}(u)}$ con listas.

\subsubsection*{Inicializar el preflujo y el algoritmo genérico}

\begin{algorithm}[H]
\caption{\textsc{Inicializar-Preflujo}$(G, s)$}
\label{alg:inicializar-preflujo}
\begin{algorithmic}[1]
\Procedure{Inicializar-Preflujo}{$G, s$}
  \ForAll{$v \in V$}
    \State $h(v) \gets 0$, \quad $e(v) \gets 0$
  \EndFor
  \ForAll{$(u, v) \in E$}
    \State $f(u, v) \gets 0$
  \EndFor
  \State $h(s) \gets n$
  \ForAll{$v$ adyacente a $s$}
    \State $f(s, v) \gets c(s, v)$ \Comment{satura todo lo que sale de $s$}
    \State $e(v) \gets c(s, v)$, \quad $e(s) \gets e(s) - c(s, v)$
  \EndFor
\EndProcedure
\end{algorithmic}
\end{algorithm}

\begin{algorithm}[H]
\caption{\textsc{Empujar-Reetiquetar-Genérico}$(G)$
         \ (\textsc{Generic-Push-Relabel} en Cormen)}
\label{alg:empujar-reetiquetar-generico}
\begin{algorithmic}[1]
\Procedure{Empujar-Reetiquetar-Genérico}{$G$}
  \State \Call{Inicializar-Preflujo}{$G, s$}
  \While{existe una operación \textsc{Empujar} o \textsc{Reetiquetar}
         que aplique}
    \State elegir una cualquiera y ejecutarla
  \EndWhile
\EndProcedure
\end{algorithmic}
\end{algorithm}

\begin{itemize}
  \item \textbf{Líneas 9 a 11 de la inicialización.} Saturan todas las
        aristas que salen de $s$; es mandar la capacidad completa del
        corte $(\set{s}, V - \set{s})$, y por eso $e(s)$ queda en menos
        esa capacidad.
  \item \textbf{$h(s) = n$ es una altura válida} justamente porque esas
        aristas quedaron saturadas. Las únicas aristas con
        $h(u) > h(v) + 1$ son las que salen de $s$, y ninguna es
        residual. Las de regreso $(v, s)$ suben, y subir está permitido.
  \item \textbf{El genérico no fija orden.} Cualquier secuencia de
        operaciones válidas termina en un flujo máximo. El orden cambia
        cuántas operaciones se hacen y cuál flujo máximo se obtiene,
        nunca su valor.
\end{itemize}

\begin{ojo}
  \begin{itemize}
    \item \textbf{Se empuja exactamente un nivel.} Empujar exige
          $h(u) = h(v) + 1$, no ``cualquier vecino más bajo''. De todos
          modos no existen aristas residuales que bajen más de un
          nivel.
    \item \textbf{$s$ y $t$ nunca se empujan ni se reetiquetan.}
          $h(s) = n$ y $h(t) = 0$ quedan fijas. $e(t)$ solo crece.
    \item \textbf{A media ejecución $f$ no es un flujo}, así que
          hablar de su valor $\abs{f}$ no tiene sentido. El valor del
          flujo final es $e(t) = -e(s)$.
  \end{itemize}
\end{ojo}

% =====================================================================

\subsection*{El genérico en una red de cuatro vértices}

\begin{ejemplo}[Empujar y reetiquetar genérico]
  \label{ej:push-generico}
  La red tiene $V = \set{s, a, b, t}$, $n = 4$, $m = 5$ y capacidades
  $c(s, a) = 4$, $c(s, b) = 2$, $c(a, b) = 2$, $c(a, t) = 1$,
  $c(b, t) = 3$. Su flujo máximo vale $4$: el corte
  $(\set{s, a, b}, \set{t})$ tiene capacidad $1 + 3 = 4$.
  \begin{center}
  \begin{tikzpicture}
    \node[vert] (s) at (0, 0) {$s$};
    \node[vert] (a) at (2, 1.2) {$a$};
    \node[vert] (b) at (2, -1.2) {$b$};
    \node[vert] (t) at (4, 0) {$t$};
    \draw[dirigida] (s) -- node[peso] {4} (a);
    \draw[dirigida] (s) -- node[peso] {2} (b);
    \draw[dirigida] (a) -- node[peso] {2} (b);
    \draw[dirigida] (a) -- node[peso] {1} (t);
    \draw[dirigida] (b) -- node[peso] {3} (t);
  \end{tikzpicture}
  \end{center}
  Como el genérico no fija orden, hace falta una regla para que la
  traza sea reproducible: \emph{se atiende el desbordado que va primero
  en el orden $a, b$; si tiene una arista admisible (residual y un
  nivel abajo) se empuja por la primera en el orden
  $\langle s, b, t\rangle$ para $a$ y $\langle s, a, t\rangle$ para
  $b$; si no tiene, se reetiqueta.}
  \begin{center}
  \small
  \begin{tabular}{rlcccccc}
    \toprule
    \textbf{Paso} & \textbf{Operación} & $h(a)$ & $h(b)$
      & $e(s)$ & $e(a)$ & $e(b)$ & $e(t)$ \\
    \midrule
    0 & \textsc{Inicializar-Preflujo} & 0 & 0 & $-6$ & 4 & 2 & 0 \\
    1 & \textsc{Reetiquetar}$(a)$: $0 \to 1$ & \textbf{1} & 0
      & $-6$ & 4 & 2 & 0 \\
    2 & \textsc{Empujar}$(a, b)$: 2, saturante & 1 & 0
      & $-6$ & 2 & 4 & 0 \\
    3 & \textsc{Empujar}$(a, t)$: 1, saturante & 1 & 0
      & $-6$ & 1 & 4 & 1 \\
    4 & \textsc{Reetiquetar}$(a)$: $1 \to 5$ & \textbf{5} & 0
      & $-6$ & 1 & 4 & 1 \\
    5 & \textsc{Empujar}$(a, s)$: 1, no saturante & 5 & 0
      & $-5$ & 0 & 4 & 1 \\
    6 & \textsc{Reetiquetar}$(b)$: $0 \to 1$ & 5 & \textbf{1}
      & $-5$ & 0 & 4 & 1 \\
    7 & \textsc{Empujar}$(b, t)$: 3, saturante & 5 & 1
      & $-5$ & 0 & 1 & 4 \\
    8 & \textsc{Reetiquetar}$(b)$: $1 \to 5$ & 5 & \textbf{5}
      & $-5$ & 0 & 1 & 4 \\
    9 & \textsc{Empujar}$(b, s)$: 1, no saturante & 5 & 5
      & $-4$ & 0 & 0 & 4 \\
    \bottomrule
  \end{tabular}
  \end{center}
  Paso a paso:
  \begin{itemize}
    \item \textbf{Paso 0.} $h(s) = 4$ y se saturan $(s, a)$ y
          $(s, b)$. Los dos internos quedan desbordados.
    \item \textbf{Paso 1.} $a$ está a altura $0$ y sus aristas
          residuales son $(a, s)$ de regreso con $c_f = 4$, $(a, b)$
          con $2$ y $(a, t)$ con $1$, hacia alturas $4$, $0$ y $0$.
          Ninguna baja, así que se reetiqueta a
          $1 + \min\set{4, 0, 0} = 1$.
    \item \textbf{Pasos 2 y 3.} Ahora $(a, b)$ y $(a, t)$ son
          admisibles. Empuja $\min(4, 2) = 2$ a $b$ y $\min(2, 1) = 1$
          a $t$; los dos empujes saturan.
    \item \textbf{Paso 4.} A $a$ le queda $1$ y su única arista
          residual es $(a, s)$, hacia altura $4$. Se reetiqueta a $5$:
          queda \emph{arriba de la fuente}.
    \item \textbf{Paso 5.} Devuelve $1$ a $s$. No satura, porque
          $c_f(a, s)$ baja de $4$ a $3$, y como dice el
          lema~\ref{lem:push-no-saturante}, $a$ se vacía.
    \item \textbf{Pasos 6 a 9.} Lo mismo con $b$, que tiene $4$. Sube a
          $1$, manda $3$ a $t$ saturando $(b, t)$, sube a
          $1 + \min\set{h(s), h(a)} = 1 + \min\set{4, 5} = 5$ y
          devuelve $1$ a $s$.
  \end{itemize}
  Cuatro instantáneas, cada vértice a su altura y, en naranja, los
  desbordados:
  \begin{center}
  \begin{tikzpicture}[y=0.45cm, nivel/.style={draw=gray!25, thin},
      pr/.style={minimum size=0.8cm, align=center, font=\scriptsize}]
    \foreach \px/\titulo in {0/{Paso 0}, 8.2/{Paso 3}} {
      \begin{scope}[xshift=\px cm]
        \foreach \k in {0, ..., 5} {
          \draw[nivel] (-0.5, \k) -- (6.0, \k);
          \node[nota] at (-0.8, \k) {\k};
        }
        \node[font=\small] at (2.7, 6.9) {\titulo};
      \end{scope}
    }
    % ----- paso 0 -----
    \begin{scope}
      \node[vert, pr] (s) at (0, 4) {$s$\\$-6$};
      \node[vertMarca, pr] (a) at (1.8, 0) {$a$\\$4$};
      \node[vertMarca, pr] (b) at (3.6, 0) {$b$\\$2$};
      \node[vert, pr] (t) at (5.4, 0) {$t$\\$0$};
      \draw[dirigida] (s) -- node[peso] {4/4} (a);
      \draw[dirigida] (s) to[bend left=25] node[peso] {2/2} (b);
      \draw[dirigida] (a) -- node[peso] {0/2} (b);
      \draw[dirigida] (a) to[bend right=35] node[peso] {0/1} (t);
      \draw[dirigida] (b) -- node[peso] {0/3} (t);
    \end{scope}
    % ----- paso 3 -----
    \begin{scope}[xshift=8.2cm]
      \node[vert, pr] (s) at (0, 4) {$s$\\$-6$};
      \node[vertMarca, pr] (a) at (1.8, 1) {$a$\\$1$};
      \node[vertMarca, pr] (b) at (3.6, 0) {$b$\\$4$};
      \node[vert, pr] (t) at (5.4, 0) {$t$\\$1$};
      \draw[dirigida] (s) -- node[peso] {4/4} (a);
      \draw[dirigida] (s) to[bend left=25] node[peso] {2/2} (b);
      \draw[dirigida] (a) -- node[peso, above right=-1pt] {2/2} (b);
      \draw[dirigida] (a) to[bend right=50]
        node[peso, pos=0.6] {1/1} (t);
      \draw[dirigida] (b) -- node[peso] {0/3} (t);
    \end{scope}
  \end{tikzpicture}
  \end{center}
  \begin{center}
  \begin{tikzpicture}[y=0.45cm, nivel/.style={draw=gray!25, thin},
      pr/.style={minimum size=0.8cm, align=center, font=\scriptsize}]
    \foreach \px/\titulo in {0/{Paso 5}, 8.2/{Paso 9 (final)}} {
      \begin{scope}[xshift=\px cm]
        \foreach \k in {0, ..., 5} {
          \draw[nivel] (-0.5, \k) -- (6.0, \k);
          \node[nota] at (-0.8, \k) {\k};
        }
        \node[font=\small] at (2.7, 6.9) {\titulo};
      \end{scope}
    }
    % ----- paso 5 -----
    \begin{scope}
      \node[vert, pr] (s) at (0, 4) {$s$\\$-5$};
      \node[vert, pr] (a) at (1.8, 5) {$a$\\$0$};
      \node[vertMarca, pr] (b) at (3.6, 0) {$b$\\$4$};
      \node[vert, pr] (t) at (5.4, 0) {$t$\\$1$};
      \draw[dirigida] (s) -- node[peso] {3/4} (a);
      \draw[dirigida] (s) to[bend right=20] node[peso] {2/2} (b);
      \draw[dirigida] (a) -- node[peso, pos=0.4] {2/2} (b);
      \draw[dirigida] (a) to[bend left=30] node[peso] {1/1} (t);
      \draw[dirigida] (b) -- node[peso] {0/3} (t);
    \end{scope}
    % ----- paso 9 -----
    \begin{scope}[xshift=8.2cm]
      \node[vert, pr] (s) at (0, 4) {$s$\\$-4$};
      \node[vert, pr] (a) at (1.8, 5) {$a$\\$0$};
      \node[vert, pr] (b) at (3.6, 5) {$b$\\$0$};
      \node[vert, pr] (t) at (5.4, 0) {$t$\\$4$};
      \draw[dirigida] (s) -- node[peso] {3/4} (a);
      \draw[dirigida] (s) to[bend right=40]
        node[peso, pos=0.3] {1/2} (b);
      \draw[dirigida] (a) -- node[peso] {2/2} (b);
      \draw[dirigida] (a) -- node[peso, pos=0.45] {1/1} (t);
      \draw[dirigida] (b) -- node[peso] {3/3} (t);
    \end{scope}
  \end{tikzpicture}
  \end{center}
  El resultado es $f(s, a) = 3$, $f(s, b) = 1$, $f(a, b) = 2$,
  $f(a, t) = 1$, $f(b, t) = 3$, con valor $e(t) = -e(s) = 4$.
\end{ejemplo}

Edmonds-Karp (algoritmo~\ref{alg:edmonds-karp}) resuelve la misma red
con tres caminos completos: $s, a, t$ con $1$, $s, b, t$ con $2$ y
$s, a, b, t$ con $1$. Empujar y reetiquetar nunca construyó un camino.
Mandó $6$ unidades hacia adelante, dejó que cada vértice resolviera
con sus vecinos qué hacer con lo suyo y devolvió a la fuente las $2$
que no cabían.

Las alturas finales también dicen dónde está el corte mínimo:
$h(t) = 0$, $h(s) = 4$ y $h(a) = h(b) = 5$. Ningún vértice quedó a
altura $1$, y ninguna arista residual puede bajar de un nivel mayor
que $1$ a uno menor que $1$ sin pasar por él, así que
$\set{v : h(v) > 1} = \set{s, a, b}$ es el lado de $s$ de un corte
mínimo. El ejercicio~\ref{ej:push-corte} lo generaliza.

\begin{ojo}
  \textbf{Las alturas sí pasan de $n$.} Es tentador pensar que nadie
  sube más que la fuente, y no es así. El exceso que no puede llegar a
  $t$ solo regresa a $s$ si los vértices que lo guardan rebasan $h(s)$:
  aquí $a$ y $b$ terminan en $5$ con $n = 4$. El tope real es
  $2n - 1$, y se prueba más abajo.
\end{ojo}

% =====================================================================

\subsection*{Por qué el resultado es un flujo máximo}

La corrección descansa en dos hechos: a un vértice desbordado siempre
se le puede hacer algo, y una función de altura impide que haya camino
de $s$ a $t$ en la red residual.

\begin{lema}[Cormen, lema 26.14]
  \label{lem:push-siempre-aplica}
  Si $f$ es un preflujo, $h$ una función de altura y $u$ está
  desbordado, entonces a $u$ le aplica un empuje o un reetiquetado.
\end{lema}

\begin{proof}
  Para toda arista residual $(u, v)$ se tiene $h(u) \le h(v) + 1$. Si
  ningún empuje aplica, ninguna cumple la igualdad, así que
  $h(u) < h(v) + 1$, es decir, $h(u) \le h(v)$, para todas ellas. Esa
  es la condición de \textsc{Reetiquetar}. El mínimo se toma sobre un
  conjunto no vacío porque $e(u) > 0$ obliga a que entre flujo por
  alguna $(v, u)$, y entonces $(u, v) \in E_f$.
\end{proof}

\begin{lema}[Cormen, lemas 26.15 y 26.16]
  \label{lem:push-altura-valida}
  Durante el algoritmo genérico las alturas nunca bajan, cada
  reetiquetado sube $h(u)$ en al menos $1$, y $h$ es una función de
  altura en todo momento.
\end{lema}

\begin{proof}[Idea de la prueba]
  Solo \textsc{Reetiquetar} cambia alturas. Antes de reetiquetar,
  $h(u) \le h(v)$ para todo vecino residual, así que
  $h(u) < 1 + \min h(v)$: la altura sube. Que $h$ sigue siendo
  función de altura se prueba por inducción en el número de
  operaciones; tras la inicialización ya lo es.
  \begin{itemize}
    \item \textbf{Reetiquetar $u$.} Las aristas residuales que salen de
          $u$ cumplen la condición por la elección del mínimo. Para una
          que entra, $(w, u)$, ya se tenía $h(w) \le h(u) + 1$, y
          $h(u)$ solo subió.
    \item \textbf{Empujar de $u$ a $v$.} Puede crear la arista residual
          $(v, u)$, que cumple porque $h(v) = h(u) - 1 < h(u) + 1$, y
          puede quitar $(u, v)$, lo que solo elimina una restricción.
          \qedhere
  \end{itemize}
\end{proof}

\begin{lema}[Cormen, lema 26.17]
  \label{lem:push-sin-camino}
  Si $f$ es un preflujo y $h$ una función de altura, no existe camino
  de $s$ a $t$ en la red residual $G_f$.
\end{lema}

\begin{proof}
  Supóngase que hay uno, $\langle v_0, v_1, \dots, v_k\rangle$ con
  $v_0 = s$ y $v_k = t$. Se puede tomar simple, así que tiene a lo más
  $n$ vértices y $k \le n - 1$ aristas. Cada $(v_i, v_{i+1})$ es arista
  residual, luego $h(v_i) \le h(v_{i+1}) + 1$. Sumando las $k$
  desigualdades, los términos intermedios se cancelan y queda
  \[
    h(s) \ \le\ h(t) + k \ =\ k \ \le\ n - 1 ,
  \]
  lo que contradice $h(s) = n$.
\end{proof}

La prueba es una cuenta de escalones: de $s$ a $t$ hay que bajar $n$
niveles, cada arista residual baja a lo más uno y un camino simple no
tiene $n$ aristas.

\begin{teorema}[Corrección del genérico; Cormen, teorema 26.18]
  \label{teo:push-correccion}
  Si el algoritmo genérico (algoritmo~\ref{alg:empujar-reetiquetar-generico})
  termina, el preflujo que calcula es un flujo máximo.
\end{teorema}

\begin{proof}
  El invariante es ``$f$ es un preflujo''. Vale tras la inicialización.
  \textsc{Reetiquetar} no toca $f$. \textsc{Empujar} lo conserva:
  $\Delta \le c_f(u, v)$ respeta la capacidad y $\Delta \le e(u)$ deja
  $e(u) \ge 0$. Al terminar no queda ningún vértice desbordado, porque
  si hubiera uno le aplicaría alguna operación
  (lema~\ref{lem:push-siempre-aplica}) y el ciclo no habría terminado.
  Entonces $e(u) = 0$ en todo $V - \set{s, t}$ y $f$ es un flujo. Por
  el lema~\ref{lem:push-altura-valida}, $h$ sigue siendo función de
  altura, y por el lema~\ref{lem:push-sin-camino} no hay camino
  aumentante en $G_f$. Por el teorema~\ref{teo:flujo-corte}, $f$ es un
  flujo máximo.
\end{proof}

Falta ver que termina. Eso sale de las cotas de la siguiente
subsección, que no dependen de las capacidades: el genérico termina
siempre, incluso con capacidades irracionales, caso en el que
Ford-Fulkerson con caminos arbitrarios puede no terminar.

% =====================================================================

\subsection*{Cuánto cuesta}

Todo el análisis cuelga de un tope para las alturas, y el tope sale de
que todo exceso vino de $s$ y tiene por dónde regresar.

\begin{lema}[Cormen, lema 26.19]
  \label{lem:push-camino-a-s}
  Si $x$ está desbordado, hay un camino simple de $x$ a $s$ en $G_f$.
\end{lema}

\begin{proof}
  Sea $U$ el conjunto de vértices alcanzables desde $x$ en $G_f$ y
  supóngase $s \notin U$. Al sumar los excesos de los vértices de $U$,
  cada $f(x', y')$ con los dos extremos en $U$ suma en uno y resta en
  otro, y queda
  \[
    \sum_{u \in U} e(u)
    = \sum_{u \in U} \sum_{v \notin U} f(v, u)
    - \sum_{u \in U} \sum_{v \notin U} f(u, v) .
  \]
  El lado izquierdo es positivo, porque $e(x) > 0$ y ningún vértice de
  $U$ tiene exceso negativo (el único que puede tenerlo es $s$). Como
  el segundo término de la derecha no es negativo, el primero es
  positivo: hay $u' \in U$ y $v' \notin U$ con $f(v', u') > 0$. Pero
  entonces $c_f(u', v') \ge f(v', u') > 0$, la arista de regreso
  $(u', v')$ es residual y $v'$ sería alcanzable desde $x$. Esto
  contradice $v' \notin U$.
\end{proof}

\begin{lema}[Tope de las alturas; Cormen, lema 26.20]
  \label{lem:push-tope-alturas}
  En todo momento $h(u) \le 2n - 1$ para todo $u \in V$.
\end{lema}

\begin{proof}
  $h(s) = n$ y $h(t) = 0$ no cambian. Un vértice interno $u$ solo
  cambia de altura al reetiquetarse, y en ese momento está desbordado.
  Reetiquetar no cambia $f$, así que después sigue desbordado y, por el
  lema~\ref{lem:push-camino-a-s}, hay un camino simple
  $\langle u = v_0, \dots, v_k = s\rangle$ en $G_f$ con $k \le n - 1$
  aristas. Sumando $h(v_i) \le h(v_{i+1}) + 1$ a lo largo del camino,
  igual que en el lema~\ref{lem:push-sin-camino},
  \[
    h(u) \ \le\ h(s) + k \ \le\ n + (n - 1) \ =\ 2n - 1 . \qedhere
  \]
\end{proof}

\begin{corolario}[Reetiquetados; Cormen, corolario 26.21]
  \label{cor:push-reetiquetados}
  Cada vértice se reetiqueta a lo más $2n - 1$ veces, y hay a lo más
  $(2n - 1)(n - 2) < 2n^2$ reetiquetados en total.
\end{corolario}

\begin{proof}[Idea de la prueba]
  Solo los $n - 2$ vértices internos se reetiquetan. Cada uno empieza
  en $0$, sube al menos $1$ por reetiquetado y no pasa de $2n - 1$.
\end{proof}

\begin{lema}[Empujes saturantes; Cormen, lema 26.22]
  \label{lem:push-saturantes}
  El algoritmo genérico hace menos de $2nm$ empujes saturantes.
\end{lema}

\begin{proof}[Idea de la prueba]
  Se cuentan juntos los empujes saturantes entre $u$ y $v$ en los dos
  sentidos; para que haya alguno, $(u, v)$ o $(v, u)$ está en $E$. Tras
  uno de $u$ a $v$ queda $h(v) = h(u) - 1$ y la arista $(u, v)$ deja de
  ser residual. Para volver a empujar de $u$ a $v$, primero tiene que
  empujarse algo de $v$ a $u$, lo que exige $h(v) = h(u) + 1$. Como
  $h(u)$ no baja, $h(v)$ tuvo que subir al menos $2$. Lo mismo vale
  para $h(u)$ entre dos saturantes de $v$ a $u$. Así, entre dos
  saturantes consecutivos del par, uno de los dos vértices subió al
  menos $2$ niveles. Con alturas entre $0$ y $2n - 1$, cada vértice
  sube de dos en dos menos de $n$ veces, y el par admite menos de $2n$
  saturantes. Los pares que cuentan son las $m$ aristas de $E$.
\end{proof}

El ejemplo~\ref{ej:push-ping-pong} muestra esta escalera en acción.
El último conteo es el que domina, y su prueba usa una herramienta que
vale la pena ver completa: una \textbf{función potencial}, un número
que nunca es negativo, que algunas operaciones suben y otras bajan.
Si se sabe cuánto puede subir en total y cuánto baja cada operación de
cierto tipo, se acota cuántas de esas operaciones hay.

\begin{lema}[Empujes no saturantes; Cormen, lema 26.23]
  \label{lem:push-no-saturantes}
  El algoritmo genérico hace menos de $4n^2(n + m)$ empujes no
  saturantes.
\end{lema}

\begin{proof}
  Sea
  \[
    \Phi \ =\ \sum_{v \text{ desbordado}} h(v) ,
  \]
  la suma de las alturas de los vértices desbordados. Tras la
  inicialización los desbordados son los vecinos de $s$, todos a altura
  $0$, así que $\Phi = 0$; y como las alturas son naturales,
  $\Phi \ge 0$ siempre. Se revisa cómo cambia $\Phi$ con cada tipo de
  operación.
  \begin{itemize}
    \item \textbf{Reetiquetar $u$.} $u$ está desbordado antes y después,
          el conjunto de desbordados no cambia y solo cambia $h(u)$.
          $\Phi$ sube exactamente lo que sube $h(u)$. La altura de cada
          interno empieza en $0$, nunca baja y no pasa de $2n - 1$
          (lema~\ref{lem:push-tope-alturas}), así que la suma de todas
          estas subidas en toda la ejecución es a lo más
          $(n - 2)(2n - 1) < 2n^2$.
    \item \textbf{Empuje saturante de $u$ a $v$.} Ninguna altura cambia.
          $u$ puede dejar de estar desbordado, lo que solo resta, y $v$
          puede volverse desbordado, lo que suma
          $h(v) \le 2n - 1$. Cada uno sube $\Phi$ menos de $2n$, y por
          el lema~\ref{lem:push-saturantes} hay menos de $2nm$, así que
          en total suben $\Phi$ menos de $4n^2 m$.
    \item \textbf{Empuje no saturante de $u$ a $v$.} Ninguna altura
          cambia. Por el lema~\ref{lem:push-no-saturante}, $u$ deja de
          estar desbordado y su término $h(u)$ sale de la suma. A lo
          más entra $v$, si antes no estaba desbordado y no es $s$ ni
          $t$, con $h(v) = h(u) - 1$. El cambio es a lo más
          $-h(u) + (h(u) - 1) = -1$: cada no saturante baja $\Phi$ en
          al menos $1$.
  \end{itemize}
  Sea $N$ el número de empujes no saturantes. $\Phi$ empieza en $0$,
  sube en total menos de $2n^2 + 4n^2 m$, baja al menos $N$ y termina
  en un valor no negativo:
  \[
    0 \ \le\ \Phi_{\text{final}} \ <\ 2n^2 + 4n^2 m - N ,
    \qquad\text{luego}\qquad
    N \ <\ 2n^2 + 4n^2 m \ \le\ 4n^2(n + m) . \qedhere
  \]
\end{proof}

Cormen acota la subida por reetiquetados con $2n \cdot 2n^2$ en lugar
de $2n^2$, y por eso enuncia la cota como $4n^2(n + m)$; la cuenta de
arriba es un poco más fina, pero la conclusión asintótica es la misma.

\begin{teorema}[Cormen, teorema 26.24 y corolario 26.25]
  \label{teo:push-costo}
  El algoritmo genérico hace $\Ord{n^2 m}$ operaciones básicas y se
  puede implementar en tiempo $\Ord{n^2 m}$.
\end{teorema}

\begin{proof}[Derivación]
  Se cuentan las operaciones y después cuánto cuesta cada una con la
  estructura adecuada.
  \begin{enumerate}
    \item \emph{Reetiquetados.} Menos de $2n^2$
          (corolario~\ref{cor:push-reetiquetados}). Cada uno recorre los
          vecinos de $u$ en $\Ord{\text{gr}(u)}$. Como $u$ se reetiqueta
          a lo más $2n - 1$ veces, el total es
          $\sum_u (2n - 1)\,\text{gr}(u) = (2n - 1) \cdot 2m = \Ord{nm}$,
          con la misma suma de grados que da el $\Tht{n + m}$ de BFS.
    \item \emph{Empujes saturantes.} Menos de $2nm$, a $\Ord{1}$ cada
          uno: $\Ord{nm}$.
    \item \emph{Empujes no saturantes.} Menos de $4n^2(n + m)$, a
          $\Ord{1}$ cada uno. Como $m \ge n - 1$, esto es
          $\Ord{n^2 m}$.
    \item \emph{Elegir la operación.} Los desbordados se guardan en una
          pila o en una cola, y para el vértice de turno se usa un
          apuntador al vecino que se está revisando, como en
          \textsc{Descargar} (algoritmo~\ref{alg:descargar}). El
          apuntador recorre la lista de $u$ una vez entre dos
          reetiquetados de $u$: otra vez $\Ord{nm}$ en total. Meter y
          sacar de la pila cuesta $\Ord{1}$ por operación.
  \end{enumerate}
  Total: $\Ord{nm + nm + n^2 m + nm} = \Ord{n^2 m}$. Domina el conteo
  de empujes no saturantes, y ninguna cuenta depende de las
  capacidades.
\end{proof}

\begin{ojo}
  El paso de $4n^2(n + m)$ a $\Ord{n^2 m}$ usa $m \ge n - 1$, que sale
  del supuesto de que todo vértice está en un camino de $s$ a $t$;
  Cormen no lo dice en la prueba. Si hay vértices que no llegan a $t$,
  el algoritmo funciona igual: lo que les entra sube y regresa a $s$.
  Lo único que se pierde es poder escribir la cota como $\Ord{n^2 m}$;
  queda $\Ord{n^2(n + m)}$.
\end{ojo}

\begin{ejemplo}[La misma red, empezando por $b$: ping-pong]
  \label{ej:push-ping-pong}
  Misma red y misma regla que en el ejemplo~\ref{ej:push-generico},
  pero atendiendo primero a $b$. Ahora son $15$ operaciones
  ($7$ reetiquetados, $4$ saturantes y $4$ no saturantes) en lugar de
  $9$, y el valor sigue siendo $4$.
  \begin{center}
  \small
  \begin{tabular}{rlcccccc}
    \toprule
    \textbf{Paso} & \textbf{Operación} & $h(a)$ & $h(b)$
      & $e(s)$ & $e(a)$ & $e(b)$ & $e(t)$ \\
    \midrule
    0 & \textsc{Inicializar-Preflujo} & 0 & 0 & $-6$ & 4 & 2 & 0 \\
    1 & \textsc{Reetiquetar}$(b)$: $0 \to 1$ & 0 & \textbf{1}
      & $-6$ & 4 & 2 & 0 \\
    2 & \textsc{Empujar}$(b, t)$: 2, no saturante & 0 & 1
      & $-6$ & 4 & 0 & 2 \\
    3 & \textsc{Reetiquetar}$(a)$: $0 \to 1$ & \textbf{1} & 1
      & $-6$ & 4 & 0 & 2 \\
    4 & \textsc{Empujar}$(a, t)$: 1, saturante & 1 & 1
      & $-6$ & 3 & 0 & 3 \\
    5 & \textsc{Reetiquetar}$(a)$: $1 \to 2$ & \textbf{2} & 1
      & $-6$ & 3 & 0 & 3 \\
    6 & \textsc{Empujar}$(a, b)$: 2, saturante & 2 & 1
      & $-6$ & 1 & 2 & 3 \\
    7 & \textsc{Empujar}$(b, t)$: 1, saturante & 2 & 1
      & $-6$ & 1 & 1 & 4 \\
    8 & \textsc{Reetiquetar}$(b)$: $1 \to 3$ & 2 & \textbf{3}
      & $-6$ & 1 & 1 & 4 \\
    9 & \textsc{Empujar}$(b, a)$: 1, no saturante & 2 & 3
      & $-6$ & 2 & 0 & 4 \\
    10 & \textsc{Reetiquetar}$(a)$: $2 \to 4$ & \textbf{4} & 3
      & $-6$ & 2 & 0 & 4 \\
    11 & \textsc{Empujar}$(a, b)$: 1, saturante & 4 & 3
      & $-6$ & 1 & 1 & 4 \\
    12 & \textsc{Reetiquetar}$(b)$: $3 \to 5$ & 4 & \textbf{5}
      & $-6$ & 1 & 1 & 4 \\
    13 & \textsc{Empujar}$(b, s)$: 1, no saturante & 4 & 5
      & $-5$ & 1 & 0 & 4 \\
    14 & \textsc{Reetiquetar}$(a)$: $4 \to 5$ & \textbf{5} & 5
      & $-5$ & 1 & 0 & 4 \\
    15 & \textsc{Empujar}$(a, s)$: 1, no saturante & 5 & 5
      & $-4$ & 0 & 0 & 4 \\
    \bottomrule
  \end{tabular}
  \end{center}
  Tres instantáneas del ping-pong entre $a$ y $b$; solo se dibujan
  $a$, $b$ y $t$, y en naranja grueso el empuje de ese paso.
  \begin{center}
  \begin{tikzpicture}[y=0.45cm, nivel/.style={draw=gray!25, thin},
      pr/.style={minimum size=0.8cm, align=center, font=\scriptsize}]
    \foreach \px/\titulo in {0/{Paso 6: $a$ empuja a $b$},
        5.6/{Paso 9: $b$ devuelve a $a$},
        11.2/{Paso 11: $a$ empuja a $b$}} {
      \begin{scope}[xshift=\px cm]
        \foreach \k in {0, ..., 5} {
          \draw[nivel] (-0.4, \k) -- (4.6, \k);
          \node[nota] at (-0.7, \k) {\k};
        }
        \node[font=\small] at (2.1, 6.4) {\titulo};
      \end{scope}
    }
    \begin{scope}
      \node[vertMarca, pr] (a) at (1.0, 2) {$a$\\$1$};
      \node[vertMarca, pr] (b) at (2.6, 1) {$b$\\$2$};
      \node[vert, pr] (t) at (4.0, 0) {$t$\\$3$};
      \draw[dirigidaResalte]
        (a) -- node[peso, above right=-1pt] {2/2} (b);
      \draw[dirigida] (b) -- node[peso, above right=-1pt] {2/3} (t);
    \end{scope}
    \begin{scope}[xshift=5.6cm]
      \node[vertMarca, pr] (a) at (1.0, 2) {$a$\\$2$};
      \node[vert, pr] (b) at (2.6, 3) {$b$\\$0$};
      \node[vert, pr] (t) at (4.0, 0) {$t$\\$4$};
      \draw[dirigida] (a) -- node[peso, below right=-1pt] {1/2} (b);
      \draw[dirigidaResalte]
        (b.north west) to[bend right=45] node[peso, above] {1} (a.north);
      \draw[dirigida] (b) -- node[peso, right=1pt] {3/3} (t);
    \end{scope}
    \begin{scope}[xshift=11.2cm]
      \node[vertMarca, pr] (a) at (1.0, 4) {$a$\\$1$};
      \node[vertMarca, pr] (b) at (2.6, 3) {$b$\\$1$};
      \node[vert, pr] (t) at (4.0, 0) {$t$\\$4$};
      \draw[dirigidaResalte]
        (a) -- node[peso, above right=-1pt] {2/2} (b);
      \draw[dirigida] (b) -- node[peso, right=1pt] {3/3} (t);
    \end{scope}
  \end{tikzpicture}
  \end{center}
  $a$ satura $(a, b)$ en el paso 6 y otra vez en el 11. Entre los dos,
  $b$ tuvo que devolverle $1$ por la arista de regreso (paso 9, que
  baja $f(a, b)$ de $2$ a $1$), y para eso subió de $1$ a $3$; después
  $a$ tuvo que subir de $2$ a $4$ para volver a empujar. Cada ida y
  vuelta le cuesta dos niveles a cada uno: es exactamente el argumento
  del lema~\ref{lem:push-saturantes}. Como las alturas están topadas
  en $2n - 1$, el ping-pong no puede durar para siempre; al final los
  dos rebasan $h(s) = 4$ y el sobrante regresa a la fuente.
\end{ejemplo}

En este caso hasta el flujo final sale igual que con el otro orden,
aunque la red tiene otro flujo máximo entero, con $f(a, b) = 1$ y
$f(s, a) = f(s, b) = 2$. En general el orden sí puede cambiar el flujo
que se obtiene, nunca su valor. Los conteos de los dos ejemplos quedan
muy por debajo de las cotas, que con $n = 4$ y $m = 5$ son $14$
reetiquetados, $40$ saturantes y $576$ no saturantes: las cotas son
para el peor caso y en redes chicas son holgadas.

% =====================================================================

\subsection*{Reetiquetar al frente}

El genérico deja libre el orden de las operaciones, y el orden importa:
en la misma red un orden hizo $9$ operaciones y otro $15$.
\textbf{Reetiquetar al frente} (\textit{relabel-to-front}) fija uno
con dos ideas. La primera es atender un vértice hasta vaciarlo, en vez
de alternar entre vértices. La segunda es el orden en que se atienden:
una lista $L$ de los vértices internos que se recorre de principio a
fin, y cada vez que un vértice sube de altura pasa al frente de la
lista y el recorrido sigue detrás de él. Para ver por qué ese orden
funciona hay que mirar qué aristas sirven para empujar.

\subsubsection*{Aristas admisibles}

\begin{definicion}[Arista admisible]
  \label{def:arista-admisible}
  Con un preflujo $f$ y una función de altura $h$, una arista $(u, v)$
  es \textbf{admisible} (\textit{admissible}) si $c_f(u, v) > 0$ y
  $h(u) = h(v) + 1$; las demás son \textbf{inadmisibles}. La
  \textbf{red admisible} es $G_{f,h} = (V, E_{f,h})$, con $E_{f,h}$ el
  conjunto de aristas admisibles.
\end{definicion}

Las aristas admisibles son exactamente aquellas por las que se puede
empujar. Tres hechos sobre ellas sostienen todo el algoritmo.

\begin{lema}[Cormen, lema 26.26]
  \label{lem:push-admisible-aciclica}
  La red admisible $G_{f,h}$ es acíclica.
\end{lema}

\begin{proof}
  Supóngase un ciclo $\langle v_0, v_1, \dots, v_k\rangle$ de aristas
  admisibles, con $v_0 = v_k$ y $k > 0$. Cada arista cumple
  $h(v_{i-1}) = h(v_i) + 1$. Sumando las $k$ igualdades,
  \[
    \sum_{i=1}^{k} h(v_{i-1}) = \sum_{i=1}^{k} h(v_i) + k .
  \]
  Las dos sumas recorren los mismos vértices del ciclo, cada uno una
  vez, y son iguales. Queda $0 = k$, una contradicción.
\end{proof}

\begin{lema}[Cormen, lemas 26.27 y 26.28]
  \label{lem:push-tras-reetiquetar}
  \begin{enumerate}
    \item Empujar de $u$ a $v$ no crea aristas admisibles, aunque puede
          quitar $(u, v)$.
    \item Si $u$ está desbordado y no tiene aristas admisibles de
          salida, \textsc{Reetiquetar}$(u)$ aplica. Después, $u$ tiene
          al menos una arista admisible de salida y ninguna de entrada.
  \end{enumerate}
\end{lema}

\begin{proof}
  (1) La única arista residual que puede aparecer al empujar es
  $(v, u)$, y no es admisible porque $h(v) = h(u) - 1$. Las alturas
  no cambian.

  (2) Sin admisibles de salida, ningún empuje aplica a $u$ y, por el
  lema~\ref{lem:push-siempre-aplica}, aplica el reetiquetado. Después,
  la arista hacia el vecino que dio el mínimo es admisible. Si hubiera
  una admisible de entrada $(v, u)$, se tendría $h(v) = h(u) + 1$ con
  la altura nueva de $u$. Como $h(u)$ subió al menos $1$, antes era
  $h(v) > h(u) + 1$ y $(v, u)$ no era residual. Reetiquetar no cambia
  la red residual, así que tampoco lo es después.
\end{proof}

\subsubsection*{Listas de vecinos y la operación Descargar}

La \textbf{lista de vecinos} $N(u)$ de un vértice interno contiene
cada $v$ con $(u, v) \in E$ o $(v, u) \in E$: los posibles extremos de
una arista residual que sale de $u$. Su orden es arbitrario pero fijo
durante todo el algoritmo. El apuntador $\text{actual}(u)$ señala el
vecino que se está revisando; empieza en el primero y vale nulo cuando
se pasa del último.

\textbf{Descargar} un vértice desbordado es empujar todo su exceso
por aristas admisibles, reetiquetándolo cuando haga falta.

\begin{algorithm}[H]
\caption{\textsc{Descargar}$(u)$ \ (\textsc{Discharge} en Cormen)}
\label{alg:descargar}
\begin{algorithmic}[1]
\Procedure{Descargar}{$u$}
  \While{$e(u) > 0$}
    \State $v \gets \text{actual}(u)$
    \If{$v = \text{nulo}$} \Comment{se recorrió toda $N(u)$}
      \State \Call{Reetiquetar}{$u$}
      \State $\text{actual}(u) \gets$ primer vértice de $N(u)$
    \ElsIf{$c_f(u, v) > 0$ \textbf{y} $h(u) = h(v) + 1$}
      \State \Call{Empujar}{$u, v$} \Comment{$(u, v)$ es admisible}
    \Else
      \State $\text{actual}(u) \gets$ vértice que sigue a $v$ en $N(u)$
    \EndIf
  \EndWhile
\EndProcedure
\end{algorithmic}
\end{algorithm}

Cada vuelta del mientras hace exactamente una de tres cosas:
\begin{enumerate}
  \item \textbf{$v$ nulo (líneas 4 a 6).} Se acabó la lista: todas las
        aristas de salida resultaron inadmisibles en esta pasada. Se
        reetiqueta $u$ y el apuntador regresa al inicio.
  \item \textbf{$(u, v)$ admisible (líneas 7 y 8).} Se empuja, y el
        apuntador \emph{no} avanza. Si el empuje no saturó, $e(u) = 0$
        y la llamada termina; si saturó, en la vuelta siguiente la
        arista aparece como inadmisible y entonces sí avanza.
  \item \textbf{$(u, v)$ inadmisible (línea 10).} Se avanza al
        siguiente vecino.
\end{enumerate}
Ni reetiquetar ni avanzar cambian $e(u)$, así que la última acción de
toda llamada sobre un vértice desbordado es un empuje. Y
$\text{actual}(u)$ se conserva entre llamadas: la siguiente descarga
de $u$ sigue donde se quedó la anterior.

\begin{lema}[Cormen, lema 26.29]
  \label{lem:push-descargar-valido}
  Cuando \textsc{Descargar}$(u)$ llama a \textsc{Reetiquetar}$(u)$ en
  la línea 5, el reetiquetado aplica.
\end{lema}

\begin{proof}[Idea de la prueba]
  Los empujes están protegidos por la condición de la línea 7. Para el
  reetiquetado, basta ver que $u$ no tiene aristas admisibles de salida
  (lema~\ref{lem:push-tras-reetiquetar}). En la pasada completa por
  $N(u)$ que acaba de terminar, y que pudo repartirse entre varias
  llamadas porque el apuntador no se reinicia, cada arista de salida se
  encontró inadmisible en algún momento. Ninguna pudo volverse
  admisible después: empujar no crea admisibles, $u$ no se reetiquetó
  en esa pasada, y un vértice $v$ que se reetiquetó no tiene
  admisibles de entrada, en particular no tiene $(u, v)$.
\end{proof}

\subsubsection*{El algoritmo}

\begin{algorithm}[H]
\caption{\textsc{Reetiquetar-al-Frente}$(G, s, t)$
         \ (\textsc{Relabel-To-Front} en Cormen)}
\label{alg:reetiquetar-al-frente}
\begin{algorithmic}[1]
\Procedure{Reetiquetar-al-Frente}{$G, s, t$}
  \State \Call{Inicializar-Preflujo}{$G, s$}
  \State $L \gets V - \set{s, t}$ en cualquier orden
  \ForAll{$u \in V - \set{s, t}$}
    \State $\text{actual}(u) \gets$ primer vértice de $N(u)$
  \EndFor
  \State $u \gets$ primer vértice de $L$
  \While{$u \neq \text{nulo}$}
    \State $h_{\text{antes}} \gets h(u)$
    \State \Call{Descargar}{$u$}
    \If{$h(u) > h_{\text{antes}}$} \Comment{$u$ se reetiquetó}
      \State mover $u$ al frente de $L$
    \EndIf
    \State $u \gets$ vértice que sigue a $u$ en $L$
  \EndWhile
\EndProcedure
\end{algorithmic}
\end{algorithm}

\begin{itemize}
  \item \textbf{Línea 3.} $L$ tiene solo los vértices que pueden
        desbordarse; $s$ y $t$ quedan fuera.
  \item \textbf{Líneas 9 a 13.} $h_{\text{antes}}$ sirve para saber si
        \textsc{Descargar} reetiquetó a $u$. Si lo hizo, $u$ pasa al
        frente de $L$.
  \item \textbf{Línea 14.} Se sigue con el sucesor de $u$ en su posición
        \emph{actual}. Si $u$ acaba de pasar al frente, el recorrido
        vuelve a pasar por todos los que quedaron detrás de él.
\end{itemize}

\begin{teorema}[Invariante y corrección de reetiquetar al frente]
  \label{teo:reetiquetar-frente}
  Cada vez que se evalúa la condición de la línea 8, $L$ es un orden
  topológico de la red admisible $G_{f,h}$ y ningún vértice anterior a
  $u$ en $L$ tiene exceso. En consecuencia, al terminar no hay
  vértices desbordados y $f$ es un flujo máximo.
\end{teorema}

\begin{proof}[Idea de la prueba]
  Que $G_{f,h}$ tenga orden topológico lo garantiza el
  lema~\ref{lem:push-admisible-aciclica}. Al inicio no hay aristas
  admisibles: $h(s) = n \ge 2$ y todos los demás están en $0$, así que
  cualquier orden sirve. Empujar no crea admisibles. Reetiquetar sí,
  pero solo de salida del vértice reetiquetado, que además se queda
  sin admisibles de entrada (lema~\ref{lem:push-tras-reetiquetar}):
  ponerlo al frente de $L$ conserva el orden topológico. Para la
  segunda parte, \textsc{Descargar}$(u)$ empuja solo por aristas
  admisibles, que en $L$ van hacia adelante (o a $s$ o $t$, que no
  están en $L$); los vértices anteriores a $u$ no reciben nada y
  siguen sin exceso, y $u$ queda vacío. Si $u$ pasó al frente, no hay
  nadie antes de él. Cuando el recorrido se sale de la lista
  ($u = \text{nulo}$), ningún vértice de $L$ tiene exceso y el
  teorema~\ref{teo:push-correccion} da la conclusión.
\end{proof}

\subsubsection*{Por qué cuesta $\Ord{n^3}$}

Se llama \textbf{fase} al tramo de ejecución entre dos reetiquetados
consecutivos (Cormen, teorema 26.30).
\begin{enumerate}
  \item \emph{Llamadas a Descargar.} Hay menos de $2n^2$
        reetiquetados, luego $\Ord{n^2}$ fases. Dentro de una fase, una
        descarga que no reetiqueta deja a $u$ en su lugar y la
        siguiente es sobre un vértice más adelante en $L$; una que
        reetiqueta termina la fase. Como $L$ tiene $n - 2$ vértices,
        una fase tiene a lo más $n - 1$ llamadas, y hay $\Ord{n^3}$ en
        total. Eso también acota el trabajo del mientras de la línea 8
        fuera de \textsc{Descargar}.
  \item \emph{Reetiquetados dentro de Descargar.} $\Ord{nm}$ en total,
        con la misma cuenta que en el teorema~\ref{teo:push-costo}.
  \item \emph{Avances del apuntador.} Entre dos reetiquetados de $u$,
        $\text{actual}(u)$ recorre $N(u)$ una vez:
        $\Ord{\text{gr}(u)}$ por reetiquetado, $\Ord{n\,\text{gr}(u)}$
        por vértice y $\Ord{nm}$ en total.
  \item \emph{Empujes saturantes.} $\Ord{nm}$
        (lema~\ref{lem:push-saturantes}).
  \item \emph{Empujes no saturantes.} Uno no saturante deja
        $e(u) = 0$ y termina la llamada, así que hay a lo más uno por
        llamada: $\Ord{n^3}$.
\end{enumerate}
El total es $\Ord{n^3 + nm}$, que es $\Ord{n^3}$ porque sin
antiparalelas $m \le n(n - 1)/2$. Comparado con el genérico, como
$m \ge n - 1$, $\Ord{n^3}$ nunca es peor que $\Ord{n^2 m}$: empatan en
grafos dispersos y reetiquetar al frente gana en densos. La mejora está
toda en el quinto renglón. El orden de $L$ topa los empujes no
saturantes en uno por llamada, en lugar de los $\Ord{n^2 m}$ que
permite la función potencial.

\begin{ojo}
  En reetiquetar al frente el apuntador $\text{actual}(u)$
  \textbf{no se reinicia} entre llamadas a \textsc{Descargar}. Si se
  reiniciara en cada llamada el algoritmo seguiría siendo correcto,
  pero se perdería la cuenta de avances que da el $\Ord{nm}$: cada
  llamada podría volver a recorrer la lista completa.
\end{ojo}

\subsubsection*{La figura 26.10 de Cormen}

\begin{ejemplo}[Reetiquetar al frente paso a paso]
  \label{ej:push-figura-26-10}
  La red tiene $V = \set{s, x, y, z, t}$, $n = 5$ y capacidades
  $c(s, x) = 12$, $c(s, y) = 14$, $c(x, y) = 5$, $c(y, z) = 8$,
  $c(z, x) = 7$, $c(z, t) = 10$, $c(x, t) = 16$. Su flujo máximo vale
  $20$. Las listas de vecinos son las de Cormen:
  $N(x) = \langle s, y, z, t\rangle$, $N(y) = \langle s, x, z\rangle$,
  $N(z) = \langle x, y, t\rangle$, y la lista inicial es
  $L = \langle x, y, z\rangle$. En las instantáneas, los desbordados
  van en naranja.
  \begin{center}
  \begin{tikzpicture}[y=0.45cm, nivel/.style={draw=gray!25, thin},
      pr/.style={minimum size=0.8cm, align=center, font=\scriptsize}]
    \foreach \px/\titulo in {0/{(a) inicio, $L = \langle x, y, z \rangle$},
        8.4/{(b) tras \textsc{Descargar}$(x)$}} {
      \begin{scope}[xshift=\px cm]
        \foreach \k in {0, ..., 6} {
          \draw[nivel] (-0.5, \k) -- (6.9, \k);
          \node[nota] at (-0.8, \k) {\k};
        }
        \node[font=\small] at (3.2, 7.4) {\titulo};
      \end{scope}
    }
    % (a)
    \begin{scope}
      \node[vert, pr] (s) at (0, 5) {$s$\\$-26$};
      \node[vertMarca, pr] (x) at (1.6, 0) {$x$\\$12$};
      \node[vertMarca, pr] (y) at (3.2, 0) {$y$\\$14$};
      \node[vert, pr] (z) at (4.8, 0) {$z$\\$0$};
      \node[vert, pr] (t) at (6.4, 0) {$t$\\$0$};
      \draw[dirigida] (s) -- node[peso, pos=0.55, left=1pt] {12/12} (x);
      \draw[dirigida] (s) -- node[peso, pos=0.4, above right=0pt]
        {14/14} (y);
      \draw[dirigida] (x) -- node[peso, above=1pt] {0/5} (y);
      \draw[dirigida] (y) -- node[peso, above=1pt] {0/8} (z);
      \draw[dirigida] (z) -- node[peso, above=1pt] {0/10} (t);
      \draw[dirigida] (z) to[bend left=30] node[peso, pos=0.3] {0/7} (x);
      \draw[dirigida] (x) to[bend right=45]
        node[peso, pos=0.7] {0/16} (t);
    \end{scope}
    % (b)
    \begin{scope}[xshift=8.4cm]
      \node[vert, pr] (s) at (0, 5) {$s$\\$-26$};
      \node[vert, pr] (x) at (1.6, 1) {$x$\\$0$};
      \node[vertMarca, pr] (y) at (3.2, 0) {$y$\\$19$};
      \node[vert, pr] (z) at (4.8, 0) {$z$\\$0$};
      \node[vert, pr] (t) at (6.4, 0) {$t$\\$7$};
      \draw[dirigida] (s) -- node[peso, pos=0.5, left=1pt] {12/12} (x);
      \draw[dirigida] (s) -- node[peso, pos=0.4, above right=0pt]
        {14/14} (y);
      \draw[dirigida] (x) -- node[peso, above=2pt, pos=0.35] {5/5} (y);
      \draw[dirigida] (y) -- node[peso, above=1pt] {0/8} (z);
      \draw[dirigida] (z) -- node[peso, above=1pt] {0/10} (t);
      \draw[dirigida] (z) to[out=-150, in=-70, looseness=1.2]
        node[peso, pos=0.62] {0/7} (x);
      \draw[dirigida] (x) to[bend right=45]
        node[peso, pos=0.7] {7/16} (t);
    \end{scope}
  \end{tikzpicture}
  \end{center}
  \textbf{De (a) a (b).} Se descarga $x$. No tiene admisibles, sube a
  $1$, empuja $5$ a $y$ (satura $(x, y)$) y $7$ a $t$. Se reetiquetó,
  así que pasa al frente, donde ya estaba; sigue $y$.
  \begin{center}
  \begin{tikzpicture}[y=0.45cm, nivel/.style={draw=gray!25, thin},
      pr/.style={minimum size=0.8cm, align=center, font=\scriptsize}]
    \foreach \px/\titulo in {0/{(c) tras \textsc{Descargar}$(y)$,
        $L = \langle y, x, z \rangle$},
        8.4/{(d) tras \textsc{Descargar}$(x)$}} {
      \begin{scope}[xshift=\px cm]
        \foreach \k in {0, ..., 6} {
          \draw[nivel] (-0.5, \k) -- (6.9, \k);
          \node[nota] at (-0.8, \k) {\k};
        }
        \node[font=\small] at (3.2, 7.4) {\titulo};
      \end{scope}
    }
    % (c)
    \begin{scope}
      \node[vert, pr] (s) at (0, 5) {$s$\\$-20$};
      \node[vertMarca, pr] (x) at (1.6, 1) {$x$\\$5$};
      \node[vert, pr] (y) at (3.2, 6) {$y$\\$0$};
      \node[vertMarca, pr] (z) at (4.8, 0) {$z$\\$8$};
      \node[vert, pr] (t) at (6.4, 0) {$t$\\$7$};
      \draw[dirigida] (s) -- node[peso, pos=0.5, left=1pt] {12/12} (x);
      \draw[dirigida] (s) -- node[peso] {8/14} (y);
      \draw[dirigida] (x) -- node[peso, pos=0.45] {0/5} (y);
      \draw[dirigida] (y) -- node[peso, pos=0.4] {8/8} (z);
      \draw[dirigida] (z) -- node[peso, above=1pt] {0/10} (t);
      \draw[dirigida] (z) to[bend left=30] node[peso, pos=0.3] {0/7} (x);
      \draw[dirigida] (x) to[bend right=45]
        node[peso, pos=0.7] {7/16} (t);
    \end{scope}
    % (d)
    \begin{scope}[xshift=8.4cm]
      \node[vert, pr] (s) at (0, 5) {$s$\\$-20$};
      \node[vert, pr] (x) at (1.6, 1) {$x$\\$0$};
      \node[vert, pr] (y) at (3.2, 6) {$y$\\$0$};
      \node[vertMarca, pr] (z) at (4.8, 0) {$z$\\$8$};
      \node[vert, pr] (t) at (6.4, 0) {$t$\\$12$};
      \draw[dirigida] (s) -- node[peso, pos=0.5, left=1pt] {12/12} (x);
      \draw[dirigida] (s) -- node[peso] {8/14} (y);
      \draw[dirigida] (x) -- node[peso, pos=0.45] {0/5} (y);
      \draw[dirigida] (y) -- node[peso, pos=0.4] {8/8} (z);
      \draw[dirigida] (z) -- node[peso, above=1pt] {0/10} (t);
      \draw[dirigida] (z) to[bend left=30] node[peso, pos=0.3] {0/7} (x);
      \draw[dirigida] (x) to[bend right=45]
        node[peso, pos=0.7] {12/16} (t);
    \end{scope}
  \end{tikzpicture}
  \end{center}
  \textbf{De (b) a (c).} Se descarga $y$ con $19$ unidades: sube a
  $1$, manda $8$ a $z$, sube a $2$, devuelve $5$ a $x$ cancelando el
  flujo de $(x, y)$, sube a $6$ (arriba de $s$) y devuelve $6$ a la
  fuente. Pasa al frente: $L = \langle y, x, z\rangle$.

  \textbf{De (c) a (d).} Sigue $x$, que recibió $5$. Su apuntador se
  había quedado en $t$ desde su descarga anterior, y $(x, t)$ es
  admisible: empuja $5$ a $t$ sin reetiquetarse y se queda en su lugar.
  \begin{center}
  \begin{tikzpicture}[y=0.45cm, nivel/.style={draw=gray!25, thin},
      pr/.style={minimum size=0.8cm, align=center, font=\scriptsize}]
    \foreach \px/\titulo in {0/{(e) tras \textsc{Descargar}$(z)$,
        $L = \langle z, y, x \rangle$}} {
      \begin{scope}[xshift=\px cm]
        \foreach \k in {0, ..., 6} {
          \draw[nivel] (-0.5, \k) -- (6.9, \k);
          \node[nota] at (-0.8, \k) {\k};
        }
        \node[font=\small] at (3.2, 7.4) {\titulo};
      \end{scope}
    }
    % (e)
    \begin{scope}
      \node[vert, pr] (s) at (0, 5) {$s$\\$-20$};
      \node[vert, pr] (x) at (1.6, 1) {$x$\\$0$};
      \node[vert, pr] (y) at (3.2, 6) {$y$\\$0$};
      \node[vert, pr] (z) at (4.8, 1) {$z$\\$0$};
      \node[vert, pr] (t) at (6.4, 0) {$t$\\$20$};
      \draw[dirigida] (s) -- node[peso, pos=0.5, left=1pt] {12/12} (x);
      \draw[dirigida] (s) -- node[peso] {8/14} (y);
      \draw[dirigida] (x) -- node[peso, pos=0.45] {0/5} (y);
      \draw[dirigida] (y) -- node[peso, pos=0.4] {8/8} (z);
      \draw[dirigida] (z) -- node[peso, above=3pt, pos=0.3] {8/10} (t);
      \draw[dirigida] (z) -- node[peso, above=1pt] {0/7} (x);
      \draw[dirigida] (x) to[bend right=40]
        node[peso, pos=0.6] {12/16} (t);
    \end{scope}
  \end{tikzpicture}
  \end{center}
  \textbf{De (d) a (e).} Sigue $z$ con $8$. Sube a $1$ y manda $8$ a
  $t$. Pasa al frente: $L = \langle z, y, x\rangle$. Después se
  recorren $y$ y $x$ sin exceso y la lista se acaba. El flujo final
  tiene $f(s, x) = 12$, $f(s, y) = 8$, $f(x, t) = 12$, $f(y, z) = 8$,
  $f(z, t) = 8$ y $f(x, y) = f(z, x) = 0$, con valor $20$. Hubo $5$
  reetiquetados, $3$ empujes saturantes y $4$ no saturantes.

  El estado después de cada llamada a \textsc{Descargar}:
  \begin{center}
  \small
  \setlength{\tabcolsep}{4.5pt}
  \begin{tabular}{clcccccccccc}
    \toprule
    & \textbf{Descarga} & \textbf{Reetiq.} & $L$ después
      & $h(x)$ & $h(y)$ & $h(z)$ & $e(s)$ & $e(x)$ & $e(y)$ & $e(z)$
      & $e(t)$ \\
    \midrule
    (a) & inicio & & $x, y, z$ & 0 & 0 & 0 & $-26$ & 12 & 14 & 0 & 0 \\
    (b) & $x$ & sí, $0 \to 1$ & $x, y, z$ & 1 & 0 & 0 & $-26$ & 0 & 19
        & 0 & 7 \\
    (c) & $y$ & sí, $0 \to 6$ & $y, x, z$ & 1 & 6 & 0 & $-20$ & 5 & 0
        & 8 & 7 \\
    (d) & $x$ & no & $y, x, z$ & 1 & 6 & 0 & $-20$ & 0 & 0 & 8 & 12 \\
    (e) & $z$ & sí, $0 \to 1$ & $z, y, x$ & 1 & 6 & 1 & $-20$ & 0 & 0
        & 0 & 20 \\
    & $y$, luego $x$ & no & $z, y, x$ & 1 & 6 & 1 & $-20$ & 0 & 0 & 0
        & 20 \\
    \bottomrule
  \end{tabular}
  \end{center}
  La descarga de $y$ es la figura 26.9 de Cormen. Al empezar,
  $h(s) = 5$, $h(x) = 1$ y $h(z) = 0$; la columna de acción dice a qué
  vecino apunta $\text{actual}(y)$ y qué se hace en cada vuelta del
  mientras.
  \begin{center}
  \small
  \begin{tabular}{clcc}
    \toprule
    \textbf{Vuelta} & \textbf{$\text{actual}(y)$ y acción}
      & $h(y)$ & $e(y)$ \\
    \midrule
    1 & $s$: $h(y) = 0 \neq h(s) + 1$, avanza & 0 & 19 \\
    2 & $x$: $h(x) = 1$, avanza & 0 & 19 \\
    3 & $z$: $h(z) = 0 = h(y)$, avanza & 0 & 19 \\
    4 & nulo: \textsc{Reetiquetar}, $1 + \min\set{5, 1, 0}$ & 1 & 19 \\
    5 & $s$: avanza & 1 & 19 \\
    6 & $x$: $h(x) = h(y)$, avanza & 1 & 19 \\
    7 & $z$: \textbf{empuja 8}, satura $(y, z)$ & 1 & 11 \\
    8 & $z$: $c_f(y, z) = 0$, avanza & 1 & 11 \\
    9 & nulo: \textsc{Reetiquetar}, $1 + \min\set{5, 1}$ & 2 & 11 \\
    10 & $s$: avanza & 2 & 11 \\
    11 & $x$: \textbf{empuja 5} de regreso, $f(x, y)$ vuelve a $0$
       & 2 & 6 \\
    12 & $x$: $c_f(y, x) = 0$, avanza & 2 & 6 \\
    13 & $z$: $c_f(y, z) = 0$, avanza & 2 & 6 \\
    14 & nulo: \textsc{Reetiquetar}, $1 + \min\set{5}$ & 6 & 6 \\
    15 & $s$: \textbf{empuja 6} de regreso a la fuente & 6 & 0 \\
    \bottomrule
  \end{tabular}
  \end{center}
  Las vueltas 7 y 11 muestran por qué el apuntador no avanza tras un
  empuje: en la vuelta siguiente la arista ya saturada se descubre
  inadmisible y entonces sí avanza. La vuelta 11 es un empuje por
  arista de regreso: $(x, y) \in E$ y $y$ le devuelve a $x$ las $5$
  unidades que $x$ le había mandado.
\end{ejemplo}

% =====================================================================

\subsection*{Comparación con Ford-Fulkerson y Edmonds-Karp}

Los dos enfoques llegan al mismo certificado de optimalidad por lados
opuestos. La primera tabla compara lo que mantiene cada uno; la
segunda, cuánto cuesta.

\begin{center}
\small
\begin{tabular}{l>{\raggedright\arraybackslash}p{5.6cm}%
                 >{\raggedright\arraybackslash}p{5.6cm}}
  \toprule
  & \textbf{Ford-Fulkerson y Edmonds-Karp}
    & \textbf{Empujar y reetiquetar} \\
  \midrule
  Mantiene & un flujo: conserva en cada vértice
           & un preflujo: a cada vértice le entra al menos lo que le
             sale \\
  \addlinespace
  Paso básico & aumentar por un camino completo de $s$ a $t$ en $G_f$
              & empujar o reetiquetar en un vértice, mirando solo a
                sus vecinos \\
  \addlinespace
  Invariante & $f$ es un flujo factible
             & no hay camino de $s$ a $t$ en $G_f$ \\
  \addlinespace
  Termina cuando & ya no hay camino aumentante
                 & ya no hay vértices desbordados \\
  \addlinespace
  Certificado & el corte de los alcanzables desde $s$ en $G_f$
              & el mismo corte, o el que da un hueco en las alturas \\
  \bottomrule
\end{tabular}
\end{center}

\begin{center}
\small
\begin{tabular}{l>{\raggedright\arraybackslash}p{3.4cm}cccc}
  \toprule
  \textbf{Algoritmo} & \textbf{Qué elige} & \textbf{Costo}
    & \textbf{Denso} & \textbf{Disperso} & \textbf{Usa $c$} \\
  \midrule
  Ford-Fulkerson & cualquier camino aumentante
    & $\Ord{m \abs{f^*}}$ & \multicolumn{2}{c}{depende de $\abs{f^*}$}
    & sí \\
  \addlinespace
  Edmonds-Karp & el camino aumentante más corto (BFS)
    & $\Ord{n m^2}$ & $\Ord{n^5}$ & $\Ord{n^3}$ & no \\
  \addlinespace
  Genérico & cualquier operación que aplique
    & $\Ord{n^2 m}$ & $\Ord{n^4}$ & $\Ord{n^3}$ & no \\
  \addlinespace
  Reetiquetar al frente & descargar en el orden de $L$
    & $\Ord{n^3}$ & $\Ord{n^3}$ & $\Ord{n^3}$ & no \\
  \bottomrule
\end{tabular}
\end{center}

``Denso'' quiere decir $m = \Tht{n^2}$ y ``disperso'' $m = \Tht{n}$.
La cota de Ford-Fulkerson vale con capacidades enteras y depende del
valor $\abs{f^*}$ del flujo máximo, no del tamaño de la red; la
columna ``Usa $c$'' dice si la cota depende de las capacidades. En
grafos densos la diferencia entre $n^5$ y $n^3$ es la razón de ser de
esta sección.

\subsubsection*{En la práctica}

Las notas del capítulo de Cormen dicen que las variantes de empujar y
reetiquetar con buenas heurísticas son los algoritmos de flujo máximo
más rápidos en la práctica. Hay dos órdenes alternativos con la misma
cota $\Ord{n^3}$: atender los desbordados en una cola FIFO (Goldberg y
Tarjan) y descargar siempre el desbordado más alto (Cheriyan y
Maheshwari), que según Cherkassky y Goldberg es la mejor variante. Y
dos heurísticas que cambian mucho el tiempo real:
\begin{itemize}
  \item \textbf{Reetiquetado global.} $h(u)$ nunca supera la distancia
        en aristas de $u$ a $t$ en $G_f$, y si $h(u) \ge n$, $h(u) - n$
        no supera la distancia de $u$ a $s$ (ejercicio 26.4-8). De vez
        en cuando conviene fijar las alturas en esas distancias exactas,
        calculadas con un BFS desde $t$ (y otro desde $s$ para los que
        ya no llegan a $t$) sobre la red residual con las aristas
        volteadas: es el $d[v]$ del BFS de la
        sección~\ref{sec:busquedas}.
  \item \textbf{Brecha} (\textit{gap}). Si ninguna altura vale $k$, con
        $0 < k < n$, los vértices de altura mayor que $k$ ya no pueden
        mandar nada a $t$, y se les puede subir de golpe por encima de
        $n$ (ejercicio~\ref{ej:push-brecha}).
\end{itemize}

% =====================================================================

\subsection*{Ejercicios}

\begin{ejercicio}[Cormen, 26.4-1]
  Probar que tras \textsc{Inicializar-Preflujo} se tiene
  $e(s) \le -\abs{f^*}$, con $f^*$ un flujo máximo.

  \textit{Solución.} La inicialización satura todo lo que sale de $s$,
  así que $e(s) = -c(\set{s}, V - \set{s})$, menos la capacidad de un
  corte. Ningún flujo supera la capacidad de un corte
  (sección~\ref{sec:flujo}), luego
  $c(\set{s}, V - \set{s}) \ge \abs{f^*}$. En el
  ejemplo~\ref{ej:push-generico}, $e(s) = -6 \le -4$.
\end{ejercicio}

\begin{ejercicio}[Cormen, 26.4-4]
  \label{ej:push-corte}
  Dado el flujo máximo que calculó empujar y reetiquetar, encontrar un
  corte mínimo en $\Ord{n}$ usando las alturas finales.

  \textit{Solución.} Hay $n - 2$ vértices internos y $n - 1$ alturas
  posibles en $\set{1, \dots, n - 1}$, así que alguna altura $k$ de ese
  rango no la tiene nadie. Sea $S = \set{v : h(v) > k}$. Contiene a
  $s$, porque $h(s) = n > k$, y no a $t$. Si $(u, v)$ es residual con
  $u \in S$, entonces $h(v) \ge h(u) - 1 \ge k$, y como nadie está en
  $k$, $h(v) > k$: ninguna arista residual sale de $S$. Por el mismo
  argumento que en el teorema~\ref{teo:flujo-corte}, las aristas de
  $S$ a $V - S$ van saturadas y las de regreso vacías, así que
  $c(S, V - S) = \abs{f}$ y el corte es mínimo. En el
  ejemplo~\ref{ej:push-figura-26-10} las alturas finales son
  $h(x) = h(z) = 1$ y $h(y) = 6$; nadie está en $2$, $S = \set{s, y}$ y
  $c(S, V - S) = c(s, x) + c(y, z) = 12 + 8 = 20$. La alternativa sin
  alturas es un BFS desde $s$ en $G_f$, en $\Ord{n + m}$.
\end{ejercicio}

\begin{ejercicio}[Cormen, 26.5-3]
  ¿Sigue funcionando el genérico si \textsc{Reetiquetar} hace solo
  $h(u) \gets h(u) + 1$?

  \textit{Solución.} Sí. Antes de reetiquetar $h(u) \le h(v)$ para
  todo vecino residual, así que sumar $1$ conserva la función de altura,
  las alturas siguen topadas en $2n - 1$ y las cotas no cambian. Lo
  que se pierde es que tras reetiquetar haya ya una arista admisible:
  en reetiquetar al frente, \textsc{Descargar} puede reetiquetar varias
  veces seguidas, una por cada pasada inútil por $N(u)$. Cada pasada
  extra se paga con un reetiquetado y el total sigue en $\Ord{n^3}$.
  En \textsc{Descargar}, el mínimo de la versión original salta de una
  vez a la altura a la que llegarían esas subidas de uno en uno.
\end{ejercicio}

\begin{ejercicio}[Heurística de brecha; Cormen, 26.5-5]
  \label{ej:push-brecha}
  Si en algún momento ninguna altura vale $k$, con $0 < k \le n - 1$,
  probar que desde ningún vértice con $h(v) > k$ hay camino a $t$ en
  $G_f$, y que subir cada $v \neq s$ con $h(v) > k$ a
  $\max(h(v), n + 1)$ deja una función de altura.

  \textit{Solución.} Sea $S = \set{v : h(v) > k}$. Como en el
  ejercicio~\ref{ej:push-corte}, ninguna arista residual sale de $S$, y
  $t \notin S$: desde $S$ no se llega a $t$. Sea $h'$ la altura después
  de la subida y $(u, v)$ una arista residual. Si $u \notin S$, la
  altura de $u$ no cambió y la de $v$ solo pudo subir, así que la
  condición se conserva. Si $u \in S$, también $v \in S$. Cuando $u$
  subió a $n + 1$, $h'(v) \ge n$ (vale $n$ si $v = s$ y al menos
  $n + 1$ si no), y $h'(u) \le h'(v) + 1$. Cuando $u$ no cambió,
  $h'(u) = h(u) \le h(v) + 1 \le h'(v) + 1$. La heurística ahorra
  trabajo porque el exceso atorado en $S$ solo puede regresar a $s$, y
  para eso esos vértices tienen que rebasar $n$ de todos modos: la
  brecha los sube de una vez en lugar de reetiquetarlos nivel por
  nivel. Las notas del capítulo dicen que es crucial en la práctica.
\end{ejercicio}

\begin{ejercicio}[Cormen, 26.5-1]
  Correr \textsc{Reetiquetar-al-Frente} sobre la red de la
  sección~\ref{sec:flujo} (figura 26.1 de Cormen) con
  $L = \langle v_1, v_2, v_3, v_4\rangle$ y listas
  $N(v_1) = \langle s, v_2, v_3\rangle$,
  $N(v_2) = \langle s, v_1, v_3, v_4\rangle$,
  $N(v_3) = \langle v_1, v_2, v_4, t\rangle$ y
  $N(v_4) = \langle v_2, v_3, t\rangle$.

  \textit{Solución.} El flujo máximo vale $23$, igual que con
  Edmonds-Karp. Se hacen $11$ reetiquetados, $3$ empujes saturantes y
  $10$ no saturantes, y las alturas finales son
  $h(v_1) = h(v_2) = 7$, $h(v_3) = 1$ y $h(v_4) = 6$. Que dos vértices
  terminen en $7 > h(s) = 6$ es el exceso que no cupo regresando a la
  fuente.
\end{ejercicio}

\subsection*{Problemas del capítulo 26}

Cada uno en una línea; los dos últimos conectan con la
sección~\ref{sec:emparejamiento}.
\begin{itemize}
  \item \textbf{26-1, problema de escape.} En una cuadrícula con varios
        puntos de inicio, decidir si hay caminos ajenos en vértices
        desde todos ellos hasta la frontera; se reduce a flujo
        partiendo cada vértice en uno de entrada y otro de salida,
        unidos por una arista de capacidad $1$.
  \item \textbf{26-2, cobertura mínima por caminos.} En un DAG, partir
        los vértices en el menor número de caminos ajenos; se arma un
        bipartito con una copia de salida y una de entrada por vértice
        y la respuesta es $n$ menos el emparejamiento máximo.
  \item \textbf{26-3, consultoría algorítmica.} Escoger expertos con
        costo y trabajos con ganancia para maximizar la ganancia neta;
        la respuesta es la suma de ganancias menos la capacidad de un
        corte mínimo.
  \item \textbf{26-4, actualizar el flujo máximo.} Si una capacidad
        entera sube o baja en $1$, recuperar un flujo máximo en
        $\Ord{n + m}$ con a lo más una búsqueda de camino aumentante.
  \item \textbf{26-5, flujo máximo por escalamiento.} Aumentar solo por
        caminos de capacidad residual al menos $K$, con $K$ una
        potencia de $2$ que se va partiendo a la mitad; con $C$ la
        capacidad máxima, cuesta $\Ord{m^2 \log C}$.
  \item \textbf{26-6, Hopcroft-Karp.} Emparejamiento bipartito máximo
        en $\Ord{\sqrt{n}\, m}$, aumentando en cada fase por un
        conjunto maximal de caminos aumentantes más cortos y ajenos en
        vértices.
\end{itemize}

El ejercicio 26.4-5 pide resolver con empujar y reetiquetar el
emparejamiento bipartito de la sección~\ref{sec:emparejamiento}. Sobre
la red de capacidades $1$ de esa sección, reetiquetar al frente lo
resuelve sin cambiar nada, en $\Ord{n^3}$.

% =====================================================================

\subsection*{Resumen de la parte de grafos}

La tabla junta todos los algoritmos de la parte, en el orden en que
aparecen, y leída de arriba abajo cuenta la historia de la parte. Los
recorridos lineales, BFS y DFS, son la base de casi todo lo demás: el
orden topológico, las componentes, los puentes, el coloreo bipartito,
los caminos en un DAG, la cerradura transitiva por recorridos y la
búsqueda de caminos aumentantes. Donde aparece un logaritmo hay una
cola de prioridad o un ordenamiento; donde hay pesos negativos aparece
Bellman-Ford; y el único renglón exponencial, Held-Karp, es el de un
problema NP-difícil.

Tres convenciones de la tabla. ``Red de flujo'' quiere decir grafo
dirigido con capacidades $c \ge 0$, fuente y sumidero, sin aristas
antiparalelas (sección~\ref{sec:flujo}). En restricciones de
diferencia, $p$ es el número de incógnitas y $r$ el de restricciones,
para no confundirlos con $n$ y $m$. Los costos son con listas de
adyacencia salvo que se diga otra cosa.

\begingroup
\small
\setlength{\tabcolsep}{4pt}
\begin{center}
\begin{tabular}{@{}>{\raggedright\arraybackslash}p{2.5cm}%
                >{\raggedright\arraybackslash}p{3.3cm}%
                >{\raggedright\arraybackslash}p{3.1cm}%
                >{\raggedright\arraybackslash}p{2.7cm}%
                >{\raggedright\arraybackslash}p{2.9cm}c@{}}
  \toprule
  \textbf{Algoritmo} & \textbf{Problema} & \textbf{Requisitos}
    & \textbf{Estructura} & \textbf{Costo} & \textbf{Secc.} \\
  \midrule
  \multicolumn{6}{@{}l}{\textit{Recorridos y estructura del grafo}} \\
  \addlinespace[2pt]
  BFS & recorrido por capas; distancias en aristas
      & ninguno; sin pesos
      & cola & $\Ord{n + m}$; $\Ord{n^2}$ con matriz
      & \ref{sec:busquedas} \\
  \addlinespace
  DFS & recorrido en profundidad; tiempos y tipos de arista
      & ninguno & pila o recursión & $\Tht{n + m}$
      & \ref{sec:busquedas} \\
  \addlinespace
  Orden topológico con DFS & toda arista hacia adelante
      & dirigido acíclico (DAG)
      & DFS e inserción al frente & $\Tht{n + m}$
      & \ref{sec:aplicaciones-dfs} \\
  \addlinespace
  Kahn & lo mismo, quitando vértices sin entrada
      & dirigido; detecta ciclos
      & cola y arreglo $\text{gr}^{-}$ & $\Tht{n + m}$
      & \ref{sec:aplicaciones-dfs} \\
  \addlinespace
  Componentes fuertemente conexas & clases de alcanzabilidad mutua
      & dirigido & dos DFS y grafo transpuesto & $\Tht{n + m}$
      & \ref{sec:aplicaciones-dfs} \\
  \addlinespace
  Puntos de articulación y puentes & vértices y aristas que
        desconectan
      & no dirigido, conexo
      & DFS con $\text{entra}[v]$ y $\text{low}[v]$ & $\Tht{n + m}$
      & \ref{sec:aplicaciones-dfs} \\
  \addlinespace
  Hierholzer & circuito o camino euleriano
      & conexo; grados pares, o dos impares y salir de uno
      & pila y apuntador por lista & $\Tht{n + m}$
      & \ref{sec:teoria-grafos} \\
  \addlinespace
  Coloreo bipartito con BFS & decidir si es bipartito; las dos partes
      & no dirigido & BFS con color por vértice & $\Tht{n + m}$
      & \ref{sec:hamilton-arboles} \\
  \addlinespace
  Held-Karp & agente viajero: ciclo hamiltoniano mínimo
      & completo y pesado; $n$ chico
      & tabla $C[S][j]$ por subconjuntos
      & $\Tht{n^2 2^n}$ tiempo, $\Tht{n 2^n}$ memoria
      & \ref{sec:hamilton-arboles} \\
  \midrule
  \multicolumn{6}{@{}l}{\textit{Árboles de expansión mínima}} \\
  \addlinespace[2pt]
  Kruskal & árbol de expansión de peso mínimo
      & no dirigido, conexo, pesado
      & ordenamiento y conjuntos disjuntos & $\Ord{m \log n}$
      & \ref{sec:mst} \\
  \addlinespace
  Prim & lo mismo, creciendo un solo árbol
      & no dirigido, conexo, pesado
      & montículo con decrementar llave
      & $\Ord{m \log n}$ binario; $\Ord{m + n \log n}$ Fibonacci;
        $\Tht{n^2}$ arreglo
      & \ref{sec:mst} \\
  \bottomrule
\end{tabular}
\end{center}

\begin{center}
\begin{tabular}{@{}>{\raggedright\arraybackslash}p{2.5cm}%
                >{\raggedright\arraybackslash}p{3.3cm}%
                >{\raggedright\arraybackslash}p{3.1cm}%
                >{\raggedright\arraybackslash}p{2.7cm}%
                >{\raggedright\arraybackslash}p{2.9cm}c@{}}
  \toprule
  \textbf{Algoritmo} & \textbf{Problema} & \textbf{Requisitos}
    & \textbf{Estructura} & \textbf{Costo} & \textbf{Secc.} \\
  \midrule
  \multicolumn{6}{@{}l}{\textit{Caminos más cortos}} \\
  \addlinespace[2pt]
  Bellman-Ford & desde una fuente; detecta ciclos negativos
      & dirigido, pesos de cualquier signo
      & arreglos $d$ y $\pi$; $n - 1$ pasadas
      & $\Ord{nm}$ & \ref{sec:caminos-fuente} \\
  \addlinespace
  Caminos en un DAG & desde una fuente
      & DAG, pesos de cualquier signo
      & orden topológico & $\Tht{n + m}$
      & \ref{sec:caminos-fuente} \\
  \addlinespace
  Dijkstra & desde una fuente
      & pesos no negativos
      & montículo de mínimos
      & $\Ord{(n + m) \log n}$ binario; $\Ord{n^2}$ arreglo;
        $\Ord{n \log n + m}$ Fibonacci
      & \ref{sec:caminos-fuente} \\
  \addlinespace
  Restricciones de diferencia & resolver $x_j - x_i \le b$ o probar
        que no hay solución
      & dos incógnitas por restricción
      & Bellman-Ford en el grafo de restricciones
      & $\Ord{p^2 + pr}$ & \ref{sec:caminos-fuente} \\
  \addlinespace
  Floyd-Warshall & entre todos los pares
      & sin ciclos negativos
      & matriz $n \times n$ & $\Tht{n^3}$ & \ref{sec:todos-pares} \\
  \addlinespace
  Cerradura transitiva & alcanzabilidad entre todos los pares
      & dirigido & matriz booleana
      & $\Tht{n^3}$; $\Tht{n(n + m)}$ con $n$ recorridos
      & \ref{sec:todos-pares} \\
  \addlinespace
  Johnson & entre todos los pares, grafos dispersos
      & sin ciclos negativos
      & Bellman-Ford y $n$ veces Dijkstra
      & $\Ord{nm \log n}$ binario; $\Ord{n^2 \log n + nm}$ Fibonacci
      & \ref{sec:todos-pares} \\
  \midrule
  \multicolumn{6}{@{}l}{\textit{Flujo y emparejamiento}} \\
  \addlinespace[2pt]
  Ford-Fulkerson & flujo máximo
      & red de flujo; capacidades enteras
      & red residual; BFS o DFS & $\Ord{m \abs{f^*}}$
      & \ref{sec:flujo} \\
  \addlinespace
  Edmonds-Karp & flujo máximo & red de flujo
      & red residual y BFS & $\Ord{n m^2}$ & \ref{sec:flujo} \\
  \addlinespace
  Emparejamiento bipartito por flujo & emparejamiento máximo
      & bipartito
      & red de capacidades $1$ y Ford-Fulkerson & $\Ord{nm}$
      & \ref{sec:emparejamiento} \\
  \addlinespace
  Hopcroft-Karp & emparejamiento máximo & bipartito
      & BFS por capas y DFS
      & $\Ord{\sqrt{n}\, m}$ & \ref{sec:emparejamiento} \\
  \addlinespace
  Empujar y reetiquetar genérico & flujo máximo & red de flujo
      & alturas, excesos y pila de desbordados & $\Ord{n^2 m}$
      & \ref{sec:push-relabel} \\
  \addlinespace
  Reetiquetar al frente & flujo máximo & red de flujo
      & lista $L$ y listas de vecinos & $\Ord{n^3}$
      & \ref{sec:push-relabel} \\
  \bottomrule
\end{tabular}
\end{center}
\endgroup

\end{document}

%########################## END OF GRAFOSCOMPLETO.TEX ##########################
%###############################################################################
